\documentclass[11pt,a4paper]{book}
\usepackage[utf8]{inputenc}
\usepackage[T1]{fontenc}
\usepackage[shorthands=off,english,french,german]{babel}
\usepackage{amsmath,amssymb,amsthm}
\usepackage{geometry}
\usepackage{graphicx}
\usepackage{fancyhdr}
\usepackage{titlesec}
\usepackage{enumitem}
\usepackage{array}
\usepackage{longtable}
\usepackage{multirow}
\usepackage{booktabs}
\usepackage{url}
\usepackage[protrusion=true,expansion=false]{microtype}
\usepackage[bookmarks=true,hidelinks]{hyperref}

\geometry{paperwidth=6.5in,paperheight=9.5in,margin=1in,footskip=0.5in}

\pagestyle{fancy}
\fancyhf{}
\fancyhead[LE,RO]{\thepage}
\fancyhead[RE]{\nouppercase{\leftmark}}
\fancyhead[LO]{\nouppercase{\rightmark}}
\renewcommand{\headrulewidth}{0pt}

\titleformat{\section}{\normalfont\Large\bfseries}{\thesection}{1em}{}
\titleformat{\subsection}{\normalfont\large\bfseries}{\thesubsection}{1em}{}

\newtheorem*{theorem}{Theorem}
\newtheorem*{lemma}{Lemma}
\newtheorem*{proposition}{Proposition}
\newtheorem*{corollary}{Corollary}
\newtheorem*{definition}{Definition}
\newtheorem*{remark}{Remark}

\setlength{\emergencystretch}{3em}
\tolerance=2000
\hbadness=10000
\hfuzz=10pt
\vbadness=10000

% Fallback macros for commands used in source chunks
\providecommand{\modd}{\ \mathrm{mod}\ }
\providecommand{\tome}[1]{t.~#1}
\providecommand{\pages}[1]{pp.~#1}
\providecommand{\page}[1]{p.~#1}
\providecommand{\vols}[1]{vols.~#1}
\providecommand{\vol}[1]{vol.~#1}
\providecommand{\Hessian}{\mathrm{Hess}}
\providecommand{\Jacobian}{\mathrm{Jac}}
\providecommand{\Jac}{\mathrm{Jac}}
\providecommand{\Pf}{\mathrm{Pf}}
\providecommand{\nominaltimes}{\cdot}
\providecommand{\nthroot}[2]{\sqrt[#1]{#2}}
\providecommand{\half}{\tfrac{1}{2}}
\providecommand{\third}{\tfrac{1}{3}}
\providecommand{\diff}{\,d}
\providecommand{\dd}{\,d}
\providecommand{\sn}{\mathrm{sn}}
\providecommand{\cn}{\mathrm{cn}}
\providecommand{\dn}{\mathrm{dn}}
\providecommand{\qform}[1]{(#1)}
\providecommand{\V}[1]{#1}
\providecommand{\uth}{\textsuperscript{th}}
\providecommand{\fA}{\mathfrak{A}}
\providecommand{\tocentry}[2]{#1\dotfill #2\par}
\providecommand{\IndexEntry}[2]{#1\dotfill #2\par}
\providecommand{\paperentry}[3]{#1.\ #2\dotfill #3\par}
\providecommand{\mypar}{\par}
\providecommand{\wrule}{\noindent\rule{\linewidth}{0.5pt}\par}
\providecommand{\art}[1]{Art.~#1}
\providecommand{\arto}[1]{Art.~#1}
\providecommand{\asterism}{\par\smallskip\centerline{$\ast\,\ast\,\ast$}\smallskip\par}
\providecommand{\Bigast}{$\ast$}
\providecommand{\canont}{\mathop{\mathrm{canon}}\nolimits}
\providecommand{\Canon}{\mathop{\mathrm{Canon}}\nolimits}
\providecommand{\Disc}{\mathop{\mathrm{Disc}}\nolimits}
\providecommand{\Res}{\mathop{\mathrm{Res}}\nolimits}
\providecommand{\between}{,}
\providecommand{\nthcomma}[1]{,}
% Pre-declare counters that some chunks expect
\makeatletter
\@ifundefined{c@art}{\newcounter{art}}{}
\makeatother

\title{The Collected Mathematical Papers of Arthur Cayley\\Vol. VII}
\author{Arthur Cayley}
\date{Modern \LaTeX{} reconstruction}

\begin{document}
\maketitle


% =========================================================
% Chunk: cayley_vol07_pages_001_050.tex  (orig pages 1--50)
% =========================================================
% ============================================================
% TITLE PAGE
% ============================================================
\thispagestyle{empty}
\begin{center}
\vspace*{2cm}
{\Huge\scshape The Collected Mathematical Papers}\\[1.5em]
{\LARGE\scshape of}\\[1.5em]
{\Huge Arthur Cayley}\\[2em]
{\LARGE\scshape Vol.\ VII}\\[3em]
{\large Cambridge University Press}\\[1em]
{1894}
\end{center}
\vfill
\clearpage

% ============================================================
% INTRODUCTORY NOTE
% ============================================================
\thispagestyle{empty}
\section*{Introductory Note}
\addcontentsline{toc}{section}{Introductory Note}

The present volume contains 68 articles numbered 417 to 484 (the last four
being posthumously published papers).  The period covered is from 1866 to
1877, but several papers of earlier date are here included on account of their
having been inserted by the author in the volumes of the \textit{Proceedings of the
London Mathematical Society} for 1865--1876.

\clearpage

% ============================================================
% TABLE OF CONTENTS
% ============================================================

\clearpage

% ============================================================
% PAPER 417
% ============================================================
\section{[417]}
\subsection{On the Locus of the Foci of the Conics which pass through Four
Given Points.}

\begin{center}
\textit{[From the Philosophical Magazine, vol.~\textsc{xxxii}. (1866), pp.~362--365.]}
\end{center}

\textsc{The} curve which is the locus of the foci of the conics which pass through four
given points is, as appears from a general theorem of M.~Chasles, a sextic curve
having a double point at each of the circular points at infinity; and Professor
Sylvester, in his ``Supplemental Note on the Analogues in Space to the Cartesian
Ovals \textit{in plano}'' (Phil.\ Mag., Oct.\ 1862), thus gives the analytical investigation.
Taking the four points as the intersections of the conics $U=0$, $V=0$, the locus
of the foci is at once deduced from the two conics of the series $U+\lambda V=0$
which pass through the circular points at infinity $J$, $J'$ respectively; it is in fact
the locus of the intersections of the tangents from $J$ and the tangents from $J'$ to any
conic of the series; which curve, if $J$, $J'$ are the circular points at infinity, is the
required curve of foci.  Taking $U+\lambda V=0$ for the equation of a conic of the series,
the pair of tangents from $J$ is given by an equation of the form
\[
(\lambda, 1 \,\,\check{a}\, x, y, z)^2 = 0,
\]
and the pair of tangents from $J'$ by an equation of the like form
\[
(\lambda, 1 \,\,\check{a}'\, x, y, z)^2 = 0;
\]
and by eliminating $\lambda$ from these equations, we have the required locus.

I propose to show how, by an obvious transformation, the same sextic curve
presents itself as a unicursal curve which can be easily obtained in a tridynamical
form.  Take the coordinates of the four points to be
\[
(a,b,1), \quad (a,-b,1), \quad (-a,b,1), \quad (-a,-b,1),
\]
so that the equations of the two conics are
\begin{align*}
U &= \tfrac{x^2}{a^2} + \tfrac{y^2}{b^2} - z^2 = 0, \\
V &= x^2 - y^2 - (a^2-b^2)z^2 = 0.
\end{align*}
Then the general conic of the series is $U + \lambda V = 0$, viz.
\[
\left(\frac{1}{a^2}+\lambda\right)x^2 + \left(\frac{1}{b^2}-\lambda\right)y^2 - (1+\lambda(a^2-b^2))z^2 = 0,
\]
and we thence find without difficulty
\begin{align*}
\lambda : 1 : \mu : m &=
a(y - \beta z)(y + \gamma z)(y - \delta z) \\
&: -b(x - \alpha z)(x + \gamma z)(x - \delta z) \\
&: a(y + \alpha z)(y - \beta z)(y - \delta z) \\
&: -b(x + \alpha z)(x - \gamma z)(x - \delta z),
\end{align*}
where $\alpha$, $\beta$, $\gamma$, $\delta$ are linear functions of $a$, $b$; and thence also
\[
\lambda : \mu = a(y - \beta z)(y + \gamma z) : a(y + \alpha z)(y - \beta z),
\]
which however it is not necessary to write down.  The resulting curve of the
order~6 is (not one of a group of curves, but the very curve) the locus of the foci
of the conics through the four points.  And the determination of the ratios
$\lambda : \mu : \nu : m$ is in fact quite simple; for writing
\begin{align*}
A &= (x-a)^2 + (y-b)^2 \\
  &= \rho^2 - 2(ax + by) + \&\text{c.} \quad (\text{if } \rho^2 = a^2 + b^2),
\end{align*}
and therefore $\sqrt{A} = \rho \sqrt{1 - \dots}$, with similar values for $\sqrt{B}$, $\sqrt{C}$, $\sqrt{D}$,
it is easy to see that, considering $\lambda$, $\mu$, $\nu$, $m$ as standing for $+\sqrt{A}$, $+\sqrt{B}$,
$+\sqrt{C}$, $+\sqrt{D}$ respectively, the equation
\[
\lambda\sqrt{A} + \mu\sqrt{B} + \nu\sqrt{C} + m\sqrt{D} = 0
\]
is satisfied identically.  In fact, if for a moment
\[
\lambda\sqrt{A} + \mu\sqrt{B} = M, \quad \nu\sqrt{C} + m\sqrt{D} = N,
\]
then $M+N=0$ is the equation of a conic passing through the four points
$(a,b)$, $(a,-b)$, $(-a,b)$, $(-a,-b)$ and having for its foci the points
$(x,y)$ and $(x',y')$; and the equation $M+N=0$ is satisfied identically
whenever the foci of the conic coincide, that is, whenever the conic is a circle,
or in the limit whenever the conic reduces to a point-pair.  Hence the required
locus is given by the equation
\[
\lambda\sqrt{A} + \mu\sqrt{B} + \nu\sqrt{C} + m\sqrt{D} = 0.
\]

\textsc{Ps.}  In general the curve
$\lambda\sqrt{A} + \mu\sqrt{B} + \nu\sqrt{C} + m\sqrt{D} = 0$
has (exclusively of multiple points at infinity) six double points; viz.\ these
are situate at the intersections of the pairs of circles,
\begin{align*}
&\lambda\sqrt{A} + \mu\sqrt{B} = 0, \quad \nu\sqrt{C} + m\sqrt{D} = 0, \\
&\lambda\sqrt{A} + \nu\sqrt{C} = 0, \quad \mu\sqrt{B} + m\sqrt{D} = 0, \\
&\lambda\sqrt{A} + m\sqrt{D} = 0, \quad \mu\sqrt{B} + \nu\sqrt{C} = 0.
\end{align*}
In the case of the curve of foci, the first, second, and third pairs of circles
intersect respectively in the points $(AB.CD)$, $(AC.BD)$, $(AD.BC)$, which,
as mentioned above, are double points on the sextic curve.

\clearpage

% ============================================================
% PAPER 418
% ============================================================
\section{[418]}
\subsection{A Remark on Differential Equations.}

\begin{center}
\textit{[From the Philosophical Magazine, vol.~\textsc{xxxii}. (1866), pp.~379--381.]}
\end{center}

\textsc{Consider} a differential equation $f(x,y,p)=0$, of the first order, but of the
degree~$n$, where $f$ is a rational and integral function of $(x,y,p)$ not rationally
decomposable into factors: the integral equation contains an arbitrary constant~$c$,
and represents therefore a system of curves, for any one of which curves the
differential equation is satisfied: the differential equation is in fact obtained from
the integral equation by eliminating~$c$.

Conversely, if the integral equation is $\phi(x,y,c)=0$, then the differential
equation is obtained by the elimination of~$c$ from the equations
\[
\phi(x,y,c)=0, \quad \frac{d\phi}{dx} + \frac{d\phi}{dy}\,p = 0,
\]
where, as usual, $p = \dfrac{dy}{dx}$.

But the equation $\phi(x,y,c)=0$ is not of necessity the equation of an undecomposable
curve, and the undecomposable curve which constitutes the proper solution of
the differential equation cannot always be represented by an equation of the form in
question.  For although $\phi$ regarded as a function of $(x,y,c)$ is not rationally
decomposable into factors, yet it may very well happen that $\phi$ regarded as a
function of $(x,y)$ is rationally decomposable into factors; that is, the equation
$\phi=0$ may represent, not a single curve, but a system of curves, each of them
varying with the arbitrary constant~$c$.  The system of curves represented by
$\phi=0$ is, it is clear, the complete solution of the differential equation, and it
comprises as a partial solution the undecomposable curve which is the proper solution
of the differential equation.

\clearpage

% ============================================================
% PAPER 419
% ============================================================
\section{[419]}
\subsection{A Theorem on Differential Operators.}

\begin{center}
\textit{[From a paper by Prof.\ Sylvester, ``Note on the Test Operators which occur in the
Calculus of Invariants, \&c.,'' Philosophical Magazine, vol.~\textsc{xxxii}. (1866),
pp.~461--472, see p.~471]}
\end{center}

\textsc{The} paper concludes with an Observation from Professor Cayley as follows:

``In the case of two variables, if
\[
P_1 = \left(ax + by\right)\frac{d}{dx} + \left(cx + dy\right)\frac{d}{dy},
\]
then in the notation of matrices,
\[
\begin{pmatrix} a & b \\ c & d \end{pmatrix}
\begin{pmatrix} \dfrac{d}{dx} \\[6pt] \dfrac{d}{dy} \end{pmatrix}
=
\begin{pmatrix} P_1 \\ P_2 \end{pmatrix},
\]
and therefore
\[
P_1^2 = \left(ax+by,\ cx+dy\,\,\check{a}\,\frac{d}{dx},\,\frac{d}{dy}\right)^2,
\]
whence also
\[
P_1^2 = (ax+by)^2\frac{d^2}{dx^2} + 2(ax+by)(cx+dy)\frac{d^2}{dx\,dy} + (cx+dy)^2\frac{d^2}{dy^2},
\]
and so on.  And generally for any number of variables, if
\[
P = (a_{11}x_1 + a_{12}x_2 + \cdots + a_{1n}x_n)\frac{d}{dx_1} + \cdots,
\]
then $P^k$ is obtained by the like rule.''

\clearpage

% ============================================================
% PAPER 420
% ============================================================
\section{[420]}
\subsection{On Riccati's Equation.}

\begin{center}
\textit{[From the Philosophical Magazine, vol.~\textsc{xxxvi}. (1868), pp.~348--351.]}
\end{center}

\textsc{The} following is, it appears to me, the proper form in which to present the
solution of Riccati's equation.

The equation may be written
\[
\frac{dy}{dx} + y^2 = x^{2q-2},
\]
which is integrable by algebraic and exponential functions if $(2i+1)q = \pm 1$,
$i$~being zero, or a positive integer.  To effect the integration, writing $y = \dfrac{1}{u}\dfrac{du}{dx}$,
we have
\[
\frac{d^2u}{dx^2} = x^{2q-2}u.
\]

The peculiar advantage of this well-known transformation has not (so far as I am
aware) been explicitly stated; it puts in evidence the form under which the
sought-for function $y$ contains the constant of integration.  In fact if $u=P$, $u=Q$
be two particular solutions of the equation in~$u$, then the general solution is
$u = CP + DQ$; and denoting by $P'$, $Q'$ the derived functions, the value of $y$ is
\[
y = \frac{CP' + DQ'}{CP + DQ},
\]
showing the form under which the constant of integration $C \div D$ is contained in~$y$.
To complete the solution, assume
\[
u = ze^{\frac{1}{q}x^q};
\]
we find
\[
\frac{d^2z}{dx^2} + 2x^{q-1}\frac{dz}{dx} + (q-1)x^{q-2}\,z = 0;
\]
considering first the particular integral of the form
\[
z = A + Bx^q + Cx^{2q} + Dx^{3q} + \&\text{c.},
\]
we find that the equation will be satisfied if
\begin{align*}
(q-1)A + q(q-1)\,B &= 0, \\
(3q-1)B + 2q(2q-1)\,C &= 0, \\
(5q-1)C + 3q(3q-1)\,D &= 0, \\
(7q-1)D + 4q(4q-1)\,E &= 0, \\
&\&\text{c.}
\end{align*}

If $A=1$, this is
\begin{align*}
A &= \phantom{-}1, \\
B &= -\frac{q-1}{q(q-1)}, \\
C &= +\frac{(q-1)(3q-1)}{q(q-1)\,2q(2q-1)}, \\
D &= -\frac{(q-1)(3q-1)(5q-1)}{q(q-1)\,2q(2q-1)\,3q(3q-1)}, \\
&\&\text{c.},
\end{align*}
where it is to be noticed that the series may be considered to stop so soon as
there is in the numerator a factor $=0$.  For instance, if $5q-1=0$, then if the
particular integral had been assumed to be $z = A + Bx^q + Cx^{2q}$, the only
conditions to be satisfied by the coefficients are the first and second equations
giving the foregoing values of $A$, $B$, $C$.  It is immaterial that the analytical
expressions of $F$ and the subsequent coefficients contain in the denominators the
evanescent factor $5q-1$; the coefficients after $C$ do not ever come into consideration.

Thus if $(2i+1)q = +1$, the series terminates, and we have for $u$ the finite
particular solution
\[
u = P = \left(1 - \frac{q-1}{q(q-1)}x^q + \frac{(q-1)(3q-1)}{q(q-1)\,2q(2q-1)}x^{2q} - \&\text{c.}\right)
e^{\frac{1}{q}x^q}:
\]
and it is easy to see that we may herein change the sign of~$x^q$, thereby obtaining
another finite particular solution,
\[
u = Q = \left(1 + \frac{q-1}{q(q-1)}x^q + \frac{(q-1)(3q-1)}{q(q-1)\,2q(2q-1)}x^{2q}
+ \&\text{c.}\right) e^{-\frac{1}{q}x^q}.
\]

Reverting to the equation in~$z$, we have next a particular solution of the form
\[
z = Ax + Bx^{q+1} + Cx^{2q+1} + Dx^{3q+1} + \&\text{c.},
\]
giving between the coefficients the relation
\begin{align*}
(q+1)A + (q+1)\,q\,B &= 0, \\
(3q+1)B + (2q+1)\,2q\,C &= 0, \\
(5q+1)C + (3q+1)\,3q\,D &= 0, \\
(7q+1)D + (4q+1)\,4q\,E &= 0, \\
&\&\text{c.}
\end{align*}

If $A=1$, we have
\begin{align*}
A &= \phantom{-}1, \\
B &= -\frac{(q+1)}{(q+1)\,q}, \\
C &= +\frac{(q+1)(3q+1)}{(q+1)\,q\,(2q+1)\,2q}, \\
D &= -\frac{(q+1)(3q+1)(5q+1)}{(q+1)\,q\,(2q+1)\,2q\,(3q+1)\,3q}, \\
&\&\text{c.},
\end{align*}
where, as in the former case, the series is considered to terminate as soon as there
is an evanescent factor in the numerator, without any regard to the subsequent
coefficients which contain in the denominators the same evanescent factor.
[In particular, $q=-1$, we have the solution $z=x$.]

Hence if we have $(2i+1)q=-1$, the series terminates, and we have for~$u$ the
finite particular solution,
\[
u = P = x\left(1 - \frac{q+1}{(q+1)\,q}x^q + \frac{(q+1)(3q+1)}{(q+1)\,q\,(2q+1)\,2q}x^{2q}
- \&\text{c.}\right)e^{\frac{1}{q}x^q}:
\]
from which, changing the sign of~$x^q$, we deduce the other finite particular solution,
\[
u = Q = x\left(1 + \frac{q+1}{(q+1)\,q}x^q + \frac{(q+1)(3q+1)}{(q+1)\,q\,(2q+1)\,2q}x^{2q}
+ \&\text{c.}\right)e^{-\frac{1}{q}x^q}.
\]

Hence, in the equation
\[
\frac{dy}{dx} + y^2 = x^{2q-2},
\]
where $q(2i+1) = \pm 1$, we have (writing $D=1$)
\[
y = \frac{CP' + Q'}{CP + Q},
\]
where $C$ is the constant of integration, $P$, $Q$ are finite series as above, and
$P'$, $Q'$ are the derived functions of $P$ and~$Q$.  Writing successively $i=0$, $i=1$,
$i=2$, \&\text{c.}, we may tabulate the solutions
\begin{align*}
\frac{dy}{dx} + y^2 = 1,       &\qquad P = e^x,            && Q = e^{-x}, \\
\frac{dy}{dx} + y^2 = x^{-4},  &\qquad P = xe^{-\frac{1}{x}}, && Q = xe^{\frac{1}{x}}, \\
\frac{dy}{dx} + y^2 = x^{-\frac{8}{3}}, &\qquad
P = (1 - 3x^{\frac{1}{3}})\,e^{3x^{\frac{1}{3}}}, &&
Q = (1 + 3x^{\frac{1}{3}})\,e^{-3x^{\frac{1}{3}}}, \\
\frac{dy}{dx} + y^2 = x^{-\frac{8}{5}}, &\qquad
P = x(1 + 3x^{-\frac{1}{3}})\,e^{-3x^{-\frac{1}{3}}}, &&
Q = x(1 - 3x^{-\frac{1}{3}})\,e^{3x^{-\frac{1}{3}}}, \\
\frac{dy}{dx} + y^2 = x^{-\frac{8}{5}}, &\qquad
P = (1 - 5x^{\frac{1}{5}} + \tfrac{25}{3}x^{\frac{2}{5}})\,e^{5x^{\frac{1}{5}}}, &&
Q = (1 + 5x^{\frac{1}{5}} + \tfrac{25}{3}x^{\frac{2}{5}})\,e^{-5x^{\frac{1}{5}}}, \\
&\&\text{c.}
\end{align*}

It is hardly necessary to make the final step of calculating $P'$ and~$Q'$ and
substituting in~$y$; but, as an example, take the above equation
$\dfrac{dy}{dx} + y^2 = x^{-\frac{8}{3}}$: we have
\[
y = \frac{-3x^{-\frac{1}{3}}(Ce^{3x^{\frac{1}{3}}} + e^{-3x^{\frac{1}{3}}})}
{C(1-3x^{\frac{1}{3}})\,e^{3x^{\frac{1}{3}}} + (1+3x^{\frac{1}{3}})\,e^{-3x^{\frac{1}{3}}}},
\]
which is readily identified with the solution, p.~98 of Boole's \textit{Differential
Equations} (Cambridge, 1859).

\begin{flushright}
\textit{Cambridge, September 29, 1868.}
\end{flushright}

\clearpage

% ============================================================
% PAPER 421
% ============================================================
\section{[421]}
\subsection{Note on the Solvibility of Equations by Means of Radicals.}

\begin{center}
\textit{[From the Philosophical Magazine, vol.~\textsc{xxxvi}. (1868), pp.~386, 387.]}
\end{center}

\textsc{In} regard to the theorem that the general quintic equation of the $n$th order is
not solvable by radicals, I believe that the proofs which have been given depend, or
at any rate that a proof may be given that shall depend, on the following two lemmas:

\medskip\noindent
\textsc{I.}  A one-valued (or symmetrical) function of $n$ letters is a perfect $k$th power,
only when the $k$th root is a one-valued function of the $n$ letters.

There is an exception in the case $k=2$, whatever be the value of~$n$: viz.\ the
product of the squares of the differences is a one-valued function, a perfect square;
but its square root, or the product of the simple differences, is a two-valued function.
It is in virtue of this exception that a quadric equation is solvable by radicals; we
have the one-valued function $(x_1-x_2)^2$, the square of a two-valued function
$x_1-x_2$, and thence the two roots are each expressible in the form
\[
\tfrac{1}{2}\{x_1+x_2 + \sqrt{(x_1-x_2)^2}\}.
\]

\medskip\noindent
\textsc{II.}  A two-valued function of $n$ letters is a perfect $k$th power, only when the
$k$th root is a two-valued function of the $n$ letters.

There is an exception in the case $k=3$, when $n=3$ or $4$: viz.\ for $n=3$ we have
$(x_1 + \omega x_2 + \omega^2 x_3)^3$ ($\omega$ an imaginary cube root of unity) a two-valued function,
and a perfect cube; whereas its cube root is the six-valued function
$x_1 + \omega x_2 + \omega^2 x_3$.  And similarly for $n=4$ we have, for instance,
\[
\{x_1x_2 + x_3x_4 + \omega(x_1x_3 + x_2x_4) + \omega^2(x_1x_4 + x_2x_3)\}^3
\]
a two-valued function, and a perfect cube, whereas its cube root is a six-valued
function.  And it is in virtue of this exception that a cubic or a quartic equation
is solvable by radicals.  But I assume that for $n>4$ the lemma is true without
exception.

The course of demonstration would be something as follows: Imagine, if possible,
the root of an equation expressed, by means of radicals, in terms of the coefficients;
the expression cannot contain any radical such as $\sqrt[p]{X}$, $p>2$, where $X$ is a one-valued
(or rational) function of the coefficients, not a perfect $p$th power, for the reason that,
expressing the coefficients in terms of the roots, such function $\sqrt[p]{X}$ is not a rational
function of the roots; if it were so, by lemma~I.\ it would be a one-valued (that is,
a symmetrical) function of the roots; consequently a rational function of the coefficients,
or $X$ expressed in terms of the coefficients, would be a perfect $p$th power.

The expression may however contain a radical $\sqrt{X}$, $X$ a one-valued (or rational)
function of the coefficients, not a perfect square: viz.\ $X$ may be any square function
multiplied into that function of the coefficients which is equal to the product of the
squared differences of the roots, or, say, multiplied into the discriminant; that is, we
may have $X = Q^2\nabla$, or $\sqrt{X} = Q\sqrt{\nabla}$.

We have next to consider whether the expression can contain any radical
$\sqrt[p]{X}$, where $X$, not being a rational function of the coefficients, is a function expressible by
radicals.  But the foregoing reasoning shows that if this be so, $X$ cannot contain any
radical other than the radical $\sqrt{Q^2\nabla}$ or $Q\sqrt{\nabla}$, as above; that is, $X$ must be
$= P + Q\sqrt{\nabla}$, where $P$ and $Q$ are rational functions of the coefficients, and where we
may assume that $P + Q\sqrt{\nabla}$ is not a perfect $p$th power of a function of the like form
$P' + Q'\sqrt{\nabla}$.  But then, expressing the coefficients in terms of the roots, we have
$P + Q\sqrt{\nabla}$, a (rational) two-valued function of the roots; and there is no radical
$\sqrt[p]{P + Q\sqrt{\nabla}}$, which is a rational function of the roots; for by lemma~II., if such
radical existed we should have $\sqrt[p]{P + Q\sqrt{\nabla}}$ a (rational) two-valued function of the
roots; that is, it would be $= P' + Q'\sqrt{\nabla}$, $P'$ and $Q'$ one-valued (symmetrical) functions
of the roots, consequently rational functions of the coefficients; or $P + Q\sqrt{\nabla}$ would
be a perfect $p$th power $(P' + Q'\sqrt{\nabla})^p$.

The conclusion is that for $n>4$ there is not (besides the function $P + Q\sqrt{\nabla}$)
any function of the coefficients, expressible by means of radicals, which, when the
coefficients are expressed in terms of the roots, will be a rational function of the
roots, and consequently there is no possibility of expressing the roots in terms of the
coefficients by means of radicals.

\begin{flushright}
\textit{Cambridge, October 1, 1868.}
\end{flushright}

\clearpage

% ============================================================
% PAPER 422
% ============================================================
\section{[422]}
\subsection{On the Geodesic Lines on an Oblate Spheroid.}

\begin{center}
\textit{[From the Philosophical Magazine, vol.~\textsc{xl}. (1870), pp.~329--340.]}
\end{center}

\textsc{The} theory of the geodesic lines on an oblate spheroid of any excentricity
whatever was investigated by Legendre\footnote{\textit{M\'{e}m.\ de l'Inst.}\ 1806; see also the
\textit{Exer.\ de Calcul Int\'{e}gral}, t.~\textsc{i}. (1811), p.~178, and the \textit{Trait\'{e} des Fonctions
Elliptiques}, t.~\textsc{i}. (1825), p.~360.}; and the general course of them is well known,
viz.\ each geodesic line undulates between two parallels equidistant from the equator
(being thus either a closed curve, or a curve of indefinite length, according to the
distance between the two parallels): at a point of contact with the parallel the curve
is, of course, at right angles to the meridian; say this is~$V$, a vertex of the geodesic
line, and let the meridian through~$V$ meet the equator in~$A$; the geodesic line proceeds
from~$V$ to meet the equator in a point~$N$, the node, where $AN$ is at most $= 90^\circ$;
and the undulations are obtained by the repetition of this portion $VN$ of the geodesic
line alternately on each side of the equator and of the meridian.

I consider in the present paper the series of geodesic lines which cut at right
angles a given meridian~$AC$, or, say, a series of geodesic normals.  It may be
remarked that as $V$ passes from the position~$A$ on the equator to the pole~$C$, the
angular distance $AN$ increases from a certain determinate value (equal, as will appear,
to $\dfrac{C}{A}\,90^\circ$, if $C$, $A$ are the polar and equatorial axes respectively) up to the value~$90^\circ$;
and it thus appears that, attending only to their course after they first meet the
equator, the geodesic normals have an envelope resembling in its general appearance
the evolute of an ellipse (see fig.~1 and also fig.~2), the centre hereof being the point
$B$ at the distance $BA = 90^\circ$, and the axes coinciding in direction with the equator
$BA$ and meridian~$BC$: this is in fact a real geodesic evolute of the meridian~$CA$.
The point~$\alpha$ is, it is clear, the intersection of the equator by the geodesic line for
which $V$ is consecutive to the point~$A$ (so that $\angle BOA = \left(1 - \dfrac{C}{A}\right)90^\circ$); and the
point $\gamma$ is the intersection of the meridian~$CB$ by the geodesic line for which $V$ is
consecutive to the point~$C$; and its position will be in this way presently determined.
I was anxious, with a view to the construction of a drawing and a model, to obtain
some numerical results in relation to a spheroid of considerable excentricity, and I
selected that for which $\dfrac{C}{A} = \tfrac{1}{2}$ (polar axis $= \tfrac{1}{2}$ equatorial).

Before proceeding further, I remark that Legendre's expression ``reduced latitude''
is used in what is not, I think, the ordinary sense; and I propose to substitute the
term ``parametric latitude'': viz., in fig.~3, referring the point~$P$ on the ellipse by means
of the ordinate $MPQ$ to a point~$Q$ on the circle, radius $OK\,(=OA$, fig.~1), and drawing
the normal~$PT$, then we have for the point~$P$ the three latitudes,
\begin{align*}
\lambda  &= \angle PTK, \quad \text{normal latitude}, \\
\lambda'' &= \angle POK, \quad \text{central latitude}, \\
\lambda'  &= \angle QOK, \quad \text{parametric latitude};
\end{align*}
viz.\ $\lambda'$ is the parameter most convenient for the expression of the values of the
coordinates $x$, $y$ ($x = A\cos\lambda'$, $y = C\sin\lambda'$) of a point~$P$ on the ellipse.  The relations
between the three latitudes are
\[
\tan\lambda'' = \frac{C}{A}\tan\lambda' = \frac{C^2}{A^2}\tan\lambda,
\]
so that $\lambda''$, $\lambda'$, $\lambda$ are in the order of increasing magnitude.  I use in like manner
$l$, $l'$, $l''$ in regard to the vertex~$V$.  The course of a geodesic line is determined by the
equation
\[
\cos\lambda'\sin\alpha = \text{const.},
\]
where $\lambda'$ is the reduced latitude of any point~$P$ on the geodesic line, and $\alpha$ is at
this point the azimuth of the geodesic line, or its inclination to the meridian.  Hence,
if $l'$ be the parametric latitude of the vertex~$V$, the equation is
\[
\cos\lambda'\sin\alpha = \cos l'
\]
(whence also, when $\lambda'=0$, $\alpha = 90^\circ-l'$; that is, the geodesic line cuts the equator at
an angle $=l'$, the parametric latitude of the vertex).  The equation in question,
$\cos\lambda'\sin\alpha = \cos l'$, leads at once to Legendre's other equations: viz.\ taking,
as above, $A$, $C$ for the equatorial and polar semiaxes respectively, and $\delta$ for the
excentricity, $\delta = \sqrt{1 - \dfrac{C^2}{A^2}}$; and to determine the position of~$P$ on the meridian,
using (instead of the parametric latitude~$\lambda'$) the angle~$\phi$ determined by the equation
\[
\cos\phi = \frac{\sin\lambda'}{\sin l'},
\]
and writing, moreover, $s$ to denote the geodesic distance~$VP$, and $\Lambda$ to denote the
longitude of~$P$ measured from the meridian~$CA$ which passes through the vertex~$V$,
these are
\begin{align*}
ds &= \phantom{\frac{\cos l'}{A}\,} d\phi\,\sqrt{C^2 + A^2\delta^2\sin^2 l'\cos^2\phi}, \\
d\Lambda &= \frac{\cos l'}{A}\,
\frac{d\phi\,\sqrt{C^2 + A^2\delta^2\sin^2 l'\cos^2\phi}}{1 - \sin^2 l'\cos^2\phi};
\end{align*}
which differential expressions are to be integrated from $\phi=0$; and the equations then
determine $\lambda'$, $s$, and $\Lambda$, all in terms of the angle~$\phi$,---that is, virtually $s$ and~$\Lambda$, the
length and longitude, in terms of the parametric latitude~$\lambda'$.

\clearpage

Writing, with Legendre,
\begin{align*}
c^2 &= \frac{A^2\delta^2\sin^2 l'}{C^2 + A^2\delta^2\sin^2 l'},
    \quad= \phantom{1-{}\,}\delta^2\sin^2 l, \\
b^2 &= 1 - c^2,
    \;= \frac{C^2}{C^2 + A^2\delta^2\sin^2 l'},
    \;= 1 - \delta^2\sin^2 l;
\end{align*}
also
\[
n = \tan^2 l', \quad M = \frac{C}{bA\cos l'} = \frac{C}{A\cos l'},
\]
then the formul\ae\ become
\begin{align*}
ds &= \frac{C}{b}\,d\phi\,\sqrt{1 - c^2\sin^2\phi}, \\
d\Lambda &= M\,\frac{d\phi\,\sqrt{1 - c^2\sin^2\phi}}{1 + n\sin^2\phi}.
\end{align*}
Hence integrating from $\phi=0$, and using the notations $F$, $E$, $\Pi$ of elliptic functions,
we have
\begin{align*}
s &= \frac{C}{b}\,E(c,\phi), \\
\Lambda &= \frac{M}{n}\left\{(n + c^2)\,\Pi\,(n,\,c,\,\phi) - c^2F(c,\phi)\right\};
\end{align*}
viz.\ these belong to any point~$P$ whatever on the geodesic line, parametric latitude
of vertex $=l'$; and if we write herein $\phi = 90^\circ$, then they will refer to the node~$N$,
or point of intersection with the equator.

The position of the point~$\alpha$ is at once obtained by writing $l'=0$: viz.\ this gives
$c=0$, $b=1$, $M = \dfrac{C}{A}$, $n=0$: the differential expressions are $ds = C\,d\phi$, $d\Lambda = \dfrac{C}{A}\,d\phi$.  Or
integrating from $\phi=0$ to $\phi = \tfrac{1}{2}\pi$, we have $s = A \cdot \dfrac{C}{A} \cdot \tfrac{1}{2}\pi$, $\Lambda = \dfrac{C}{A} \cdot \tfrac{1}{2}\pi$, agreeing with each
other, and giving longitude of $\alpha = \dfrac{C}{A} \cdot \tfrac{1}{2}\pi$; or, what is the same thing,
$\angle\alpha OB = \tfrac{1}{2}\pi\left(1 - \dfrac{C}{A}\right)$.

Writing in the formul\ae\ $l' = 90^\circ$, we have $c=\delta$, $b = \dfrac{C}{A}$, $\dfrac{M}{n} = 0$; whence $d\Lambda = 0$, or
$\Lambda = \text{const.} = \tfrac{1}{2}\pi$, since the geodesic line here coincides with the meridian~$CB$; and
moreover $s = AE(\delta,\phi)$; viz.\ this is merely the expression of the distance from~$C$ of
a point~$P$ on the meridian~$CB$.  But we do not thus obtain the position of the point~$\gamma$.

To find it we must consider a position of~$V$ consecutive to~$C$, say, $l' = \tfrac{1}{2}\pi - \epsilon$,
where $\epsilon$ is indefinitely small; $n$ is thus indefinitely large, and the integral $\Pi\,(n,\,c,\,\phi)$ is not
conveniently dealt with.  But it may be replaced by an expression depending on
$\Pi\left(\dfrac{c^2}{n},\,c,\,\phi\right)$, where $\dfrac{c^2}{n}$ is indefinitely small; viz.\ (Legendre, \textit{Fonct.\ Ellipt.}\ vol.~I. p.~69)
we have
\[
\Pi\,(n,\,c,\,\phi) = F(c,\phi) + \frac{1}{\sqrt{\alpha}}\tan^{-1}
\frac{\sqrt{\alpha}\,\tan\phi}{\sqrt{1-c^2\sin^2\phi}}
- \Pi\left(\frac{c^2}{n},\,c,\,\phi\right),
\]
where
\[
\alpha = (1+n)\left(1+\frac{c^2}{n}\right).
\]
We thus have
\[
\Lambda = \frac{M}{n}\left\{nF(c,\phi) + \frac{\sqrt{n}\,\sqrt{c^2+n}}{\sqrt{1+n}}
\tan^{-1}\frac{\sqrt{\alpha}\,\tan\phi}{\sqrt{1-c^2\sin^2\phi}}
- (c^2+n)\,\Pi\left(\frac{c^2}{n},\,c,\,\phi\right)\right\},
\]
where, $\dfrac{c^2}{n}$ being small,
\[
\Pi\left(\frac{c^2}{n},\,c,\,\phi\right) =
\int\!\frac{d\phi}{\left(1+\dfrac{c^2}{n}\sin^2\phi\right)\sqrt{1-c^2\sin^2\phi}},
\]
\[
= \int\!\frac{\left(1-\dfrac{c^2}{n}\sin^2\phi\right)d\phi}{\sqrt{1-c^2\sin^2\phi}},
= \left(1-\frac{1}{n}\right)F(c,\phi) + \frac{1}{n}\,E(c,\phi).
\]

And expanding also the $\tan^{-1}$ term, we thus have
\begin{align*}
\Lambda &= \frac{M}{n}\left\{nF(c,\phi) + \frac{\sqrt{n}\,\sqrt{c^2+n}}{\sqrt{1+n}}
\left[\tfrac{1}{2}\pi - \frac{\sqrt{1-c^2\sin^2\phi}}{\tan\phi}
\frac{\sqrt{n}}{\sqrt{(1+n)(c^2+n)}}\right]\right. \\
&\qquad\qquad\qquad\qquad
\left. - (c^2+n)\left[\left(1-\frac{1}{n}\right)F(c,\phi) + \frac{1}{n}\,E(c,\phi)\right]\right\}, \\
&= \frac{M}{n}\left\{\left(b^2+\frac{c^2}{n}\right)F(c,\phi)
- \left(1+\frac{c^2}{n}\right)E(c,\phi)
+ \frac{\sqrt{n}\,\sqrt{c^2+n}}{\sqrt{1+n}} \cdot \tfrac{1}{2}\pi
- \frac{n}{n+1}\cot\phi\,\sqrt{1-c^2\sin^2\phi}\right\},
\end{align*}
which, in the term in $\{\}$ neglecting negative powers of~$n$, becomes
\[
\Lambda = \frac{M}{n}\left\{\sqrt{n} \cdot \tfrac{1}{2}\pi + b^2F(c,\phi) - E(c,\phi)
- \cot\phi\,\sqrt{1-c^2\sin^2\phi}\right\}.
\]
We may moreover write $c=\delta$, $b=\dfrac{C}{A}$, $\phi = 90^\circ-\lambda'$,
$n = \dfrac{1}{\epsilon^2}$, $M=\epsilon$, and therefore $\dfrac{M}{n}=\epsilon$, so
that the formula is
\begin{align*}
\Lambda &= \epsilon\left\{\frac{1}{\epsilon} \cdot \tfrac{1}{2}\pi
+ b^2F(c,\,90^\circ-\lambda') - E(c,\,90^\circ-\lambda')
- \tan\lambda'\,\sqrt{1-c^2\cos^2\lambda'}\right\}, \\
&= \tfrac{1}{2}\pi - \epsilon\left\{\tan\lambda'\,\sqrt{1-c^2\cos^2\lambda'}
+ E(c,\,90^\circ-\lambda') - b^2F(c,\,90^\circ-\lambda')\right\},
\end{align*}
where I retain $c$, $b$ as standing for $\sqrt{1-\dfrac{C^2}{A^2}}$, $\dfrac{C}{A}$ respectively.

\clearpage

Writing herein $\lambda'=0$, we have
\[
\Lambda = \tfrac{1}{2}\pi - \epsilon\,(E,\,c - b^2F,\,c),
\]
where the coefficient $E,\,c - b^2F,\,c$ is
\[
= \int_0^{\frac{1}{2}\pi}\!d\theta\left(\sqrt{1-c^2\sin^2\theta} - \frac{1-c^2}{\sqrt{1-c^2\sin^2\theta}}\right)
= c^2\!\int_0^{\frac{1}{2}\pi}\!\frac{\cos^2\theta\,d\theta}{\sqrt{1-c^2\sin^2\theta}}
\]
consequently positive; that is, $\Lambda$, the longitude of the node, is less than~$90^\circ$, as it
should be.  Hence in order that $\Lambda$ may be $= 90^\circ$, we must have $\lambda'$ negative, say,
$\lambda' = -\mu'$, where $\mu'$ is positive; and, observing that we may under the signs $E$, $F$
write $90^\circ-\mu'$ instead of $90^\circ+\mu'$, we thus have
\[
\tfrac{1}{2}\pi = \tfrac{1}{2}\pi + \epsilon\left\{\sqrt{1-c^2\cos^2\mu'}\,\tan\mu'
- E(c,\,90^\circ-\mu') + b^2F(c,\,90^\circ-\mu')\right\};
\]
that is, we must have
\[
\tan\mu'\,\sqrt{1-c^2\cos^2\mu'} = E(c,\,90^\circ-\mu') - b^2F(c,\,90^\circ-\mu');
\]
viz.\ $\mu'$ is here the parametric latitude (south) of the intersection of the meridian~$CB$
with the consecutive geodesic line---that is, of the point~$\gamma$.  As $\mu'$ increases from
$0$ to~$90^\circ$, the left-hand side increases from $0$ to~$\infty$; and the right-hand side, beginning
from a positive value and either attaining a maximum or not, ultimately decreases
to~$0$; there is consequently a real root, which is easily found by trial.

Thus $\dfrac{C}{A} = \tfrac{1}{2}$, $C = \tfrac{1}{2}\sqrt{3}$ (angle of modulus $= 60^\circ$), $b = \tfrac{1}{2}$; or the equation is
\[
\tan\mu'\,\sqrt{1 - \tfrac{3}{4}\cos^2\mu'}
= E\,(90^\circ-\mu') - \tfrac{1}{4}F\,(90^\circ-\mu').
\]
Using Legendre's Table~IX., we have
\begin{center}
\begin{tabular}{c|c|c|c|c|c}
\toprule
$\mu'$ & $90^\circ-\mu'$ & $E$. & $F$. & $E-\tfrac{1}{4}F$. & $\tan\mu'\sqrt{1-\tfrac{3}{4}\cos^2\mu'}$ \\
\midrule
 $0^\circ$  & $90^\circ$ & $1.21105$ & $2.15651$ & $\cdot 6719$ & $\cdot 0$ \\
$10$  & $80$  & $1.12248$ & $1.81252$ & $\cdot 6693$ & \\
$20$  & $70$  & $1.02663$ & $1.49441$ & $\cdot 6530$ & \\
$30$  & $60$  & $\cdot 91839$ & $1.21253$ & $\cdot 6153$ & $\cdot 3819$ \\
$40$  & $50$  & $\cdot 79538$ & $\cdot 96465$ & $\cdot 5542$ & $\cdot 6278$ \\
\bottomrule
\end{tabular}
\end{center}
so that we see the required value is between $30^\circ$ and $40^\circ$; and a rough interpolation
gives the value $\mu' = 37^\circ\,40'$.  But repeating the calculation with the values $37^\circ$ and
$38^\circ$, we have
\begin{center}
\begin{tabular}{c|c|c|c|c|c}
\toprule
$\mu'$ & $90^\circ-\mu'$ & $E$. & $F$. & $E-\tfrac{1}{4}F$. & $\tan\mu'\sqrt{1-\tfrac{3}{4}\cos^2\mu'}$ \\
\midrule
 $37^\circ$ & $53^\circ$ & $\cdot 833879$ & $1.035870$ & $\cdot 57419$ & $\cdot 54425$ \\
$38$  & $52$  & $\cdot 821197$ & $1.011849$ & $\cdot 56823$ & $\cdot 57108$ \\
\bottomrule
\end{tabular}
\end{center}
whence, interpolating, $\mu' = 37^\circ\,55'$.

\clearpage

The semiaxes of the geodesic evolute, measured according to their longitude and
parametric latitude respectively, are thus $B\alpha$, long.\ of $\alpha = 45^\circ$; $B\gamma$, param.\ lat.
$= 37^\circ\,55'$.  But measuring them according to their geodesic distance, the equatorial
radius~$A$ being taken $=1$, we have
\begin{align*}
B\alpha &= \tfrac{1}{4}\pi &&= \cdot 78540, \\
B\gamma &= \left(\frac{C}{b}-1\right)\{E, - E(52^\circ\,5')\} = 1.21106 - 0.82225 = \cdot 38881.
\end{align*}

Reverting to the general formul\ae\ for $s$, $\Lambda$, but writing therein $A=1$, and therefore
$C = \sqrt{1-\delta^2}$; writing also $\phi=90^\circ$ (that is, making the formul\ae\ to refer to the node
$N$ of the geodesic line), we have
\begin{align*}
s &= \frac{\sqrt{1-\delta^2}}{\sqrt{1-\delta^2\sin^2 l}}\,E,\,c, \\
\Lambda &= \frac{\sqrt{1-\delta^2}}{n\cos l}\left\{(n+c^2)\,\Pi\,(n,\,c) - c^2F,\,c\right\};
\end{align*}
but for the calculation of the second of these formul\ae\ by means of Legendre's Tables
it is necessary to express $\Pi\,(n,\,c)$ in terms of the functions $E$, $F$.

The proper formula is given in \textit{Fonct.\ Ellipt.}\ vol.~I. p.~137; viz.\ this is
\[
\frac{\Delta\,(b,\,\theta)}{\sin\theta\cos\theta}\,\Pi\,(n,\,c)
= \tfrac{1}{2}\pi + \frac{\sin\theta}{\cos\theta}\,\Delta\,(b,\,\theta)\,F,\,c
+ F,\,c\,F\,(b,\,\theta) - F,\,c\,E\,(b,\,\theta) - E,\,c\,F\,(b,\,\theta),
\]
where $\Delta\,(b,\,\theta) = \sqrt{1-b^2\sin^2\theta}$.  $\theta$ is an angle given by the equation $\cot\theta = \sqrt{n}$; we
have $n = \tan^2 l'$; consequently $\theta = 90^\circ - l'$.  Substituting this value, except that for
shortness I retain $E\,(b,\,\theta)$, $F\,(b,\,\theta)$ in place of $E\,(b,\,90^\circ-l')$, $F\,(b,\,90^\circ-l')$, we have
\begin{align*}
\Delta\,(b,\,\theta) &= \sqrt{1 - b^2\cos^2 l'}, \\
&= \sqrt{1 - (1-\delta^2\sin^2 l)\cos^2 l'}, \;= \sin l;
\end{align*}
and thence
\[
\tan\theta\,\Delta\,(b,\,\theta) = \cot l\sin l = \frac{\cos l}{\sqrt{1-\delta^2}};
\]
whence
\[
\Pi\,(n,\,c) = \frac{\sin l'\cos l'}{\sin l}
\left\{\tfrac{1}{2}\pi + F,\,c\left[\frac{\cos l}{\sqrt{1-\delta^2}}
+ F\,(b,\,\theta) - E\,(b,\,\theta)\right] - E,\,c\,F\,(b,\,\theta)\right\}.
\]
Also
\[
n + c^2 = \tan^2 l' + \delta^2\sin^2 l = \sin^2 l\sec^2 l.
\]
Hence
\[
(n+c^2)\,\Pi\,(n,\,c) - c^2F,\,c
= \sin^2 l\left\{\sec^2 l'\,\Pi\,(n,\,c) - \delta^2F,\,c\right\};
\]
and multiplying this by
\[
\frac{\sqrt{1-\delta^2}}{n\cos l}, \;= \frac{\sqrt{1-\delta^2}\cos l}{\tan^2 l'\cos^2 l'},
\]
the exterior factor is
\[
\frac{\sqrt{1-\delta^2}\cos l\tan^2 l'}{\tan^2 l'}, \;= \frac{\cos l}{\sqrt{1-\delta^2}},
\]
and we have
\[
\Lambda = \frac{\cos l}{\sqrt{1-\delta^2}}
\left\{\sec^2 l'\,\Pi\,(n,\,c) - \delta^2F,\,c\right\},
\]
which is the formula I used in the calculations.  It would, however, have been better
to reduce a step further; viz.\ we have
\begin{align*}
\sec^2 l'\,\Pi\,(n,\,c)
&= \frac{\tan l'}{\tan l\cos l}\,\{\}, \\
&= \frac{\sqrt{1-\delta^2}}{\cos l}
\left\{\tfrac{1}{2}\pi + F,\,c\left[\frac{\cos l}{\sqrt{1-\delta^2}}
+ F\,(b,\,\theta) - E\,(b,\,\theta)\right] - E,\,c\,F\,(b,\,\theta)\right\}, \\
&= \frac{\sqrt{1-\delta^2}}{\cos l}
\left\{\tfrac{1}{2}\pi + F,\,c\left[F\,(b,\,\theta) - E\,(b,\,\theta)\right]
- E,\,c\,F\,(b,\,\theta)\right\} + F,\,c,
\end{align*}
and thence
\begin{align*}
\sec^2 l'\,\Pi\,(n,\,c) - \delta^2F,\,c
&= \frac{\sqrt{1-\delta^2}}{\cos l}
\left\{\tfrac{1}{2}\pi + F,\,c\left[\sqrt{1-\delta^2}\cos l
\right.
\right. \\
&\qquad\left.\left. + F\,(b,\,\theta) - E\,(b,\,\theta)\right]
- E,\,c\,F\,(b,\,\theta)\right\};
\end{align*}
or, finally,
\[
\Lambda = \tfrac{1}{2}\pi + F,\,c\,F\,(b,\,\theta) - F,\,c\,E\,(b,\,\theta)
- E,\,c\,F\,(b,\,\theta) + \sqrt{1-\delta^2}\cos l\,F,\,c.
\]

It is easy with this expression of~$\Lambda$ to obtain the results already found for the
extreme values $l'= 0^\circ$, $l'= 90^\circ$.

As Legendre's Tables have for argument, not the modulus~$c$, but the angle of the
modulus, say~$\chi$ (that is, $\sin\chi = c = \delta\sin l$), it is convenient to replace
$\sqrt{1-\delta^2\sin^2 l}$ by its value $\cos\chi$; and the formul\ae\ thus are
\begin{align*}
s &= \frac{\sqrt{1-\delta^2}}{\cos\chi}\,E,\,c, \\
\Lambda &= \tfrac{1}{2}\pi + F,\,c\left[\sqrt{1-\delta^2}\cos l
+ F\,(b,\,\theta) - E\,(b,\,\theta)\right] - E,\,c\,F\,(b,\,\theta),
\end{align*}
where
\[
c = \sin\chi = \delta\sin l, \quad
\tan l' = \sqrt{1-\delta^2}\tan l, \quad
\theta = 90^\circ - l';
\]
and in the case intended to be numerically discussed, $\delta = \tfrac{1}{2}\sqrt{3}$,
$\sqrt{1-\delta^2} = \tfrac{1}{2}$.  I take $l'$ as the argument, giving it the values  $0^\circ$, $10^\circ$, \dots\  $90^\circ$,
and perform the calculation as shown in the Table.

\clearpage

The Table is as follows (the values are for $\delta = \tfrac{1}{2}\sqrt{3}$, $\sqrt{1-\delta^2} = \tfrac{1}{2}$):

\begin{center}
\small
\begin{tabular}{c|c|c|c|c|c|c|c|c|c}
\toprule
\parbox{0.6cm}{\centering $l'$} &
\parbox{0.6cm}{\centering $\theta$} &
\parbox{1cm}{\centering $l$} &
\parbox{1cm}{\centering $\chi = c$} &
\parbox{1.2cm}{\centering $90^\circ-\chi=l$} &
\parbox{0.8cm}{\centering $\tfrac{1}{2}\cos l$} &
\parbox{1cm}{\centering $F,\,c$} &
\parbox{1cm}{\centering $E,\,c$} &
\parbox{1.2cm}{\centering $F\,(b,\,\theta)$} &
\parbox{1.2cm}{\centering $E\,(b,\,\theta)$} \\
\midrule
 $0^\circ$  & $90^\circ$  & $0^\circ$     & $60.0$    & $0^\circ$    & ---       & $1.21105$     & $0.83162$   & ---       & ---       \\
$10^\circ$  & $80^\circ$  & $19^\circ 26'$ & $16.75$   & $73.25$     & $\cdot47151$ & $0.20548$     & $1.8686$    & $1.6050$  & $1.5376$  \\
$20^\circ$  & $70^\circ$  & $36^\circ 3'$  & $30.63$   & $59.37$     & $\cdot40425$ & $0.22814$     & $1.6532$    & $1.6910$  & $1.4633$  \\
$30^\circ$  & $60^\circ$  & $49^\circ 6'$  & $40.90$   & $49.10$     & $\cdot32737$ & $0.25478$     & $1.4167$    & $1.7980$  & $1.3857$  \\
$40^\circ$  & $50^\circ$  & $59^\circ 13'$ & $48.07$   & $41.93$     & $\cdot25589$ & $0.27626$     & $1.2163$    & $1.8891$  & $1.3232$  \\
$50^\circ$  & $40^\circ$  & $67^\circ 14'$ & $52.98$   & $37.02$     & $\cdot19349$ & $0.29935$     & $1.0645$    & $1.9923$  & $1.2777$  \\
$60^\circ$  & $30^\circ$  & $73^\circ 54'$ & $56.32$   & $33.68$     & $\cdot13866$ & $0.31479$     & $0.9557$    & $2.0644$  & $1.2461$  \\
$70^\circ$  & $20^\circ$  & $79^\circ 41'$ & $58.43$   & $31.57$     & $\cdot08954$ & $0.32540$     & $0.8850$    & $2.1154$  & $1.2260$  \\
$80^\circ$  & $10^\circ$  & $84^\circ 58'$ & $59.62$   & $30.38$     & $\cdot04387$ & $0.33169$     & $0.8446$    & $2.1463$  & $1.2147$  \\
$90^\circ$  & $0^\circ$   & $90^\circ$     &           & $60.0$      &           &               & $0.8316$    &           &           \\
\bottomrule
\end{tabular}
\end{center}

\begin{center}
\small
\begin{tabular}{c|c|c|c|c|c|c|c|c|c}
\toprule
\parbox{0.8cm}{\centering $\tfrac{1}{2}\cos l + F\,(b,\,\theta) - E\,(b,\,\theta)$} &
\parbox{0.8cm}{\centering Do.\ log.} &
\parbox{0.8cm}{\centering Add log $F,\,c$.} &
\parbox{0.8cm}{\centering (1) Nat.} &
\parbox{1cm}{\centering $\log F\,(b,\,\theta)$} &
\parbox{0.8cm}{\centering Add log $E,\,c$.} &
\parbox{0.8cm}{\centering (2) Nat.} &
\parbox{1cm}{\centering $\tfrac{1}{2}\pi+(1)-(2)$} &
\parbox{0.8cm}{\centering Do.\ in $(^\circ)$.} &
\parbox{1cm}{\centering $\log E,\,c$ $-\log\cos\chi$} \\
\midrule
$\cdot 7854$ & $45^\circ$ & & & & & & $\cdot 7854$ & $45^\circ$ & $\cdot 7854$ \\
$\cdot 8022$ & $45$  & $58$ & $20569$ & $1.6058$ & $\cdot 8029$ & & & & \\
$\cdot 8525$ & $48$  & $51$ & $23075$ & $1.7012$ & $\cdot 8506$ & & & & \\
$\cdot 9257$ & $53$  & $2$  & $26323$ & $1.8333$ & $\cdot 9166$ & & & & \\
$1.0110$    & $57$  & $56$ & $29668$ & $1.9801$ & $\cdot 9900$ & & & & \\
$1.1166$    & $63$  & $59$ & $32682$ & $2.1224$ & $1.0612$ & & & & \\
$1.2250$    & $70$  & $11$ & $35178$ & $2.2479$ & $1.1239$ & & & & \\
$1.3378$    & $76$  & $39$ & $36939$ & $2.3410$ & $1.1705$ & & & & \\
$1.4536$    & $83$  & $17$ & $38454$ & $2.4240$ & $1.2120$ & & & & \\
$1.5708$    & $90$  &      &         &          & $1.2111$ & & & & \\
\bottomrule
\end{tabular}
\end{center}

$\chi$, $90^\circ-\chi$ in degrees and decimals of a degree, to correspond with Legendre's Tables.

where the columns marked with an $*$ show respectively the longitude of the node,
and the length (or distance of node from vertex), for the geodesic lines belonging to
the different values of the argument~$l'$.

\clearpage

The remarks which follow have reference to the stereographic projection of the
figure on the plane of the equator, the centre of projection being the pole (say the
South Pole) of the spheroid.  It is to be remarked that if a point~$P$ of the spheroid
is projected as above, by means of an ordinate into the point~$Q$ of the sphere radius
$OK\,(=OA)$, then projecting stereographically as to the spheroid and the sphere from
the south poles thereof respectively, the points $P$ and $Q$ have the same projection.
And it is hence easy to show that an azimuth~$\alpha$ at a point of the meridian (parametric
latitude~$\lambda'$, normal latitude~$\lambda$, and therefore $\tan\lambda' = \dfrac{C}{A}\tan\lambda$) is projected
into an angle~$(\alpha)$ such that
\[
\tan{(\alpha)} = \frac{\sin\lambda'}{\sin\lambda}\tan\alpha.
\]

In fact in fig.~3, if we take therein $OK$, $OC$ for the axes of $x$, $z$ respectively,
and the axis of~$y$ at right angles to the plane of the paper, and if we have at~$P$
on the surface of the spheroid an element of length~$PR$ at the inclination~$\alpha$ to the
meridian~$PK$, then if $x$, $y$, $z$ are the coordinates of~$P$, and $x+\delta x$, $y+\delta y$, $z+\delta z$ those
of~$R$, we have
\begin{alignat*}{3}
\delta x &= \phantom{-}\rho\cos\alpha\sin\lambda, &\qquad
\delta z &= -\rho\cos\alpha\cos\lambda, &\qquad
\delta y &= \phantom{-}\rho\sin\alpha,
\end{alignat*}
and thence
\[
\tan\alpha = \frac{\delta y}{\sqrt{\delta x^2 + \delta z^2}}.
\]
Now, if the meridian and the points $P$, $R$ are referred by lines parallel to $Oz$ to the
surface of the sphere radius~$OA$, the only difference is that the ordinates~$z$ are
increased in the ratio $C:A$; so that if the projected angle be~$(\alpha)$, we have
\[
\tan{(\alpha)} = \frac{\delta y}{\sqrt{\delta x^2 + \dfrac{A^2}{C^2}\delta z^2}};
\]
and then projecting the sphere stereographically from its south pole, the angle in the
projection is $=(\alpha)$.  And according to the foregoing remark, the angle~$(\alpha)$ thus obtained
is also the projection of~$\alpha$ from the south pole of the spheroid.  We have thus
\[
\frac{\tan{(\alpha)}}{\tan\alpha}
= \frac{\sqrt{\delta x^2 + \delta z^2}}{\sqrt{\delta x^2 + \dfrac{A^2}{C^2}\delta z^2}},
= \frac{\sqrt{\sin^2\lambda + \cos^2\lambda}}
{\sqrt{\sin^2\lambda + \dfrac{A^2}{C^2}\cos^2\lambda}},
= \sqrt{\frac{1+\cot^2\lambda}{1+\cot^2\lambda'}},
= \frac{\sin\lambda'}{\sin\lambda},
\]
which is the required relation.

The foregoing equations,
\[
\cos\lambda'\sin\alpha = \cos l', \quad
\tan\lambda' = \frac{C}{A}\tan\lambda, \quad
\tan{(\alpha)} = \frac{\sin\lambda'}{\sin\lambda}\tan\alpha,
\]
determine in the stereographic projection the inclination~$(\alpha)$ to the radius, or projection
of the meridian, of the geodesic line (parametric latitude of vertex $=l'$) at the point
the parametric latitude of which is $=\lambda'$; viz.\ they enable the construction (in the
projection) of the direction of the successive elements of the geodesic line.  There
would be no difficulty in performing the construction geometrically; but it would,
I think, be more convenient to calculate~$(\alpha)$ numerically for a given value of~$l'$ and
for the successive values of~$\lambda'$.  Observe that for $\lambda'=0$ we have (as above) $90^\circ-\alpha=l'$,
and then $\dfrac{\sin\lambda'}{\sin\lambda} = \dfrac{\tan\lambda'}{\tan\lambda} = \dfrac{C}{A}$, consequently $\tan{(\alpha)} = \dfrac{C}{A}\cot l'$: but we have also
$\cot l' = \dfrac{A}{C}\cot l$, so that this equation becomes $\tan{(\alpha)} = \cot l$, or we have $90^\circ-(\alpha)=l$;
viz.\ in the projection, the geodesic line cuts the equator at an angle $l=$ the normal
latitude of the vertex of the geodesic line.

The preceding formul\ae\ and results have enabled me to construct a drawing, on
a large scale, of the stereographic projection of the geodesic lines for the spheroid,
polar axis $= \tfrac{1}{2}$ equatorial axis.

\clearpage

% ============================================================
% PAPER 423
% ============================================================
\section{[423]}
\subsection{On the Plane Representation of a Solid Figure.}

\begin{center}
\textit{[From the Philosophical Magazine, vol.~\textsc{xli}. (1871), pp.~286--290.]}
\end{center}

\textsc{We} represent \textit{in plano} the position of a point~$P$ whose coordinates in space are
$(x,\,y,\,z)$ by drawing these coordinates, on the same scale or on different scales, and
in given directions from a fixed origin in the plane; $OM = x$, $MP' = y$, $P'P'' = z$.  But
observe that the point~$P''$ \textit{alone} does not completely represent the point~$P$; in fact
$P''$ represents a whole series of points lying in a line; any one such point is the
point whose coordinates are $Om$, $mp'$, $p'P''$.  For the complete representation of~$P$ we
require the \textit{two points} $P'$, $P''$: these might be distinguished as the projection~$P''$, and
the foot-point~$P'$.  The two points $P'$, $P''$ are obviously such that the line joining them
is in a given direction.

The preceding is, of course, the ordinary method of orthogonal projection, or
geometrical delineation of a solid figure: it may be used under various forms; for
example, the coordinates $x$, $y$, $z$ may be taken on the same scale and in directions
inclined to each other at angles of  $120^\circ$ (isometrical projection); or the coordinates $x$, $y$
may be drawn on the same scale and at their actual inclination,  $90^\circ$, to each other;
and the coordinate~$z$ on the same or an altered scale in any given direction; the
points $P'$ then give a true ground-plan of the solid figure, and the lengths of the
lines $P'P''$ give the altitudes of the several points~$P$: this is also a method in
ordinary use.

But it is to be observed that the points $P'$, $P''$ are both of them \textit{projections},
and that the general theory is as follows: we represent the position of the point~$P$
by means of its projections $P'$, $P''$, from two fixed points $\Omega'$, $\Omega''$ respectively; the
line joining these points passes, it is clear, through a fixed point~$\Omega$ which is the
intersection of the plane of projection by the line which joins the two points $\Omega'$, $\Omega''$.

Hence we say that a point~$P$ in space is represented \textit{in plano} by any two points
$P'$, $P''$ which are such that the line joining them passes through a fixed point~$\Omega$.
And we have thus a \textit{system of constructive geometry} which is the more simple on
account of the generality of its basis, and which is at once applicable to any of the
special projections above referred to.  I establish the fundamental notions of such a
geometry, and by way of illustration apply it to the solution of the well-known pro-
blem of finding the lines which meet four given lines in space.

A point~$P$ (as already mentioned) is given by its projections $P'$, $P''$, which are
points such that the line joining them passes through the fixed point~$\Omega$.

A line~$L$ is given by its projections $L'$, $L''$, which are any two lines in the plane.
We speak of the point $(P',\,P'')$, meaning the point~$P$ whose projections are $P'$, $P''$;
and similarly of the line $(L',\,L'')$, meaning the line whose projections are $L'$, $L''$.

If $P'$, $P''$ coincide, then the point~$P$ is in the plane of projection; and so if
$L'$, $L''$ coincide, then the line~$L$ is in the plane of projection.

If through~$\Omega$ we draw a line meeting $L'$, $L''$ in the points $P'$, $P''$ respectively,
these are the projections of a point~$P$ on the line~$L$.  In particular the intersection
of $L'$, $L''$ (considered as two coincident points) represents the intersection of the line
$L$ with the plane of projection.

\clearpage

The line through the points $(P',\,P'')$ and $(Q',\,Q'')$ has for its projections the lines
$P'Q'$ and $P''Q''$.

Two lines $(L',\,L'')$ and $(M',\,M'')$ intersect each other if only the intersections $L'M'$
and $L''M''$ are the projections of a point~$P$---that is, if the line through the points
$L'M'$ and $L''M''$ passes through~$\Omega$.  And then clearly $P$ is the intersection of the two
lines.

A plane~$\Pi$ is conveniently given by means of its trace~$\Theta$ on the plane of pro-
jection, and of the projections $(P',\,P'')$ of a point on the plane; or, say, by means
of the trace~$\Theta$, and of a point~$P$ on the plane.

Suppose, however, that a plane is given by means of a line~$L$ and a point~$P$
on the plane.  The trace~$\Theta$ passes through the point of intersection of the line~$L$ with
the plane of projection---that is, through the point of intersection of the projections
$L'$, $L''$.  To find another point on the trace, we have only to imagine on the line~$L$
a point~$Q$, and, joining this with~$P$, to suppose the line~$PQ$ produced to meet the
plane of projection.  The construction is obvious; but by way of illustration I give it
in full.  Through $\Omega$ draw a line meeting $L'$, $L''$ in $Q'$, $Q''$ respectively (then these are
the projections of a point~$Q$ on the line~$L$); the lines $P'Q'$ and $P''Q''$ are the pro-
jections of the line~$PQ$, and the intersection of $P'Q'$ and $P''Q''$ is therefore the required
point on the trace~$\Theta$.

The line of intersection of two planes passes through the point of intersection of
their traces $\Theta_1$, $\Theta_2$; whence, if the planes have in common a point~$P$, the line of
intersection is the line joining $P$ with the intersection of the traces $\Theta_1$, $\Theta_2$.

In what precedes we have the solution of the following problem:---``Given a point
$P$, and two lines $L_1$, $L_2$, to find a line through~$P$ meeting the two lines $L_1$, $L_2$.''  The
required line is in fact the line of intersection of the planes $(P,\,L_1)$ and $(P,\,L_2)$; we
have seen how to construct the traces $\Theta_1$ and $\Theta_2$ of these planes respectively; and
the required line is the line joining $P$ with the intersection of $\Theta_1$ and $\Theta_2$.

I proceed now to the problem to find the two lines, each of them meeting four
given lines, $L_1$, $L_2$, $L_3$, $L_4$ (these being, of course, given by means of their projections
$(L_1',\,L_1'')$ \&\text{c.}).  The question is in effect to find on the line~$L_1$ a point~$P$ such that,
drawing from it a line to meet $L_2$, $L_3$, and also a line to meet $L_3$, $L_4$, these shall
be one and the same line.

Now, considering in the first instance $P$ as an arbitrary point on the line~$L_1$,
the line from~$P$ to meet $L_2$, $L_3$ is any line whatever meeting the lines $L_1$, $L_2$, $L_3$:
say it is a generating line of the hyperboloid whose directrices are $L_1$, $L_2$, $L_3$, or of
the hyperboloid $L_1L_2L_3$.  Hence projecting from any point $\Omega'$ whatever, the generating
lines and directrices are projected into tangents of one and the same conic.  We know
the projections $L_1'$, $L_2'$, $L_3'$ of the directrices; to find two other tangents of the conic,
we take two arbitrary positions of $P$ on the line~$L_1$, and construct as above the pro-
jections $M'$, $N'$ of the lines from these to meet the lines $L_2$, $L_3$.  The conic is then
determined as the conic touching the five lines $L_1'$, $L_2'$, $L_3'$, $M'$, $N'$: instead of
$L_1'$ we may take the lines $L_2'$, $L_3'$, $M'$, $N'$ and the line $\Omega'$; and similarly the conic
$\Sigma''$ is determined as the conic touching the five lines $L_2''$, $L_3''$, $M''$, $N''$ and the line
$\Omega''$.  Each tangent $T'$ of $\Sigma'$, combined with the corresponding tangent $T''$ of $\Sigma''$,
represents a line $T$ meeting $L_1$, $L_2$, $L_3$; to establish the correspondence, observe that,
inasmuch as the line $T$ meets $L_1$, the intersections of $T'$, $L_1'$ and of $T''$, $L_1''$ must lie
in a line with $\Omega$; if $T'$ be given, the point $(T'', L_1'')$ is thus uniquely determined,
and therefore also $T''$ (since $L_1''$ is a tangent of $\Sigma''$); and similarly if $T''$ be given,
$T'$ is uniquely determined: the correspondence $T'$, $T''$ is thus, as it should be, a
$(1, 1)$ correspondence.

Considering in like manner the lines $L_1''$, $L_2''$, $L_3''$, $L_4''$, $M''$, $N''$ we have touching
these a conic $\Sigma''$; each tangent $T'$ of $\Sigma'$, combined with the corresponding tangent
$T''$ of $\Sigma''$, represents a line $T$ meeting $L_1$, $L_2$, $L_3$, $L_4$, the correspondence being
a $(1, 1)$ correspondence such as in the former case.  The conics $\Sigma'$, $\Sigma'$ both touch
$L_1'$, $L_2'$ and have thus with each other and be a common two tangents.  Say one of
these is $T' = T'$, the corresponding common tangent of $\Sigma''$, $\Sigma''$ (these conics both
touch $L_1''$, $L_2''$, and have thus in common two tangents).  We have thus $T' = T'$,
and $T'' = T''$, as the common tangents of $\Sigma'$, $\Sigma'$, and $\Sigma''$, $\Sigma''$; and taking the
other common tangents we have the projections of the other line meeting $L_1$, $L_2$,
$L_3$, $L_4$.

The whole process is:---Construct $M'$, $M''$ and $N'$, $N''$ each of them the projections
of a line through a point $P$ of $L_1$, which meets $L_2$, $L_3$; and $M'$, $M''$ and $N'$, $N''$ each
of them the projections of a line through a point $P$ of $L_1$, which meets $L_3$, $L_4$; we
have then the conics $\Sigma'$, $\Sigma''$ touching $L_1'$, $L_2'$, $L_3'$, $M'$, $N'$, and $L_1''$, $L_2''$, $L_3''$,
$M''$, $N''$ respectively, and then the projections of each of the required lines are
$T' = T'$, a common tangent of $\Sigma'$, $\Sigma'$, and $T'' = T''$, the corresponding common
tangent of $\Sigma''$, $\Sigma''$.

It is material to remark how the construction is simplified when there is given
one of the lines, say, $M$, which meets $L_1$, $L_2$, $L_3$, $L_4$.  Here if $M$ is a common
directrix of the two hyperboloids we may for the hyperbolas $\Sigma'$ and $\Sigma''$ consider,
instead of $L_1$, $L_2$, $L_3$ and two new generating lines, the lines $L_1$, $L_2$, $L_3$, $M$, and a
single new generating line $N$; and similarly for the hyperbolas $\Sigma''$, $\Sigma''$ the lines
$L_1''$, $L_2''$, $L_3''$, $M''$, and a single new generating line $N$.  The conics $\Sigma'$, $\Sigma'$ have
the three tangents $L_1'$, $L_2'$, $M'$, and therefore only a single other common tangent,
$T' = T'$; and similarly $\Sigma''$, $\Sigma''$ have the three tangents $L_1''$, $L_2''$, $M''$, and therefore
only a single other common tangent, $T'' = T''$; and we have thus the other line
meeting $L_1$, $L_2$, $L_3$, $L_4$, and cutting the four given lines.

I take the opportunity of mentioning the following theorem: ``If in a given
triangle we inscribe a variable triangle of given form, the envelope of each side of
the variable triangle is a conic touching the two sides (of the given triangle) which
contain the extremities of the variable side in question.''

We have thence a solution of the problem (Principia, Book~I. Sect.~V. Lemma
XXVII.), in a given quadrilateral to inscribe a quadrangle of given form.  The
question in effect is: in the triangle $ABC$ to inscribe a triangle $\alpha\beta\gamma$ of given form;
and in the triangle $ADE$ a triangle $\alpha'\beta'\gamma'$ of given form, in such wise that the sides
$\alpha\gamma$, $\alpha'\gamma'$ may be coincident.  The envelope of $\alpha\gamma$ is a conic touching $AD$, $AE$,
and the envelope of $\alpha'\gamma'$ a conic also touching $AD$, $AE$; of which two common
tangents, one may be taken for the position of the side $\alpha\gamma = \alpha'\gamma'$; there are thus
two other common tangents, either of which may be the side $\alpha\gamma = \alpha'\gamma'$; and the
problem admits accordingly of two solutions.
\clearpage


% =========================================================
% Chunk: cayley_vol07_pages_051_100.tex  (orig pages 51--100)
% =========================================================
% No title page; direct transcription begins
\setcounter{page}{51}

% ============================================================
% PAPER 433: SUR LES SURFACES TETRAEDRALES (continued)
% ============================================================
\selectlanguage{french}

\noindent\paperref{433}\quad
\begin{center}
\textbf{SUR LES SURFACES TETRAEDRALES.}
\end{center}

\noindent
(ce qui donne quatre syst\`emes de valeurs de $\lambda : \mu : \nu$), et puis
\[
\Omega = af\frac{\nu - \mu}{\lambda} + bg\frac{\lambda - \nu}{\mu} + ch\frac{\mu - \lambda}{\nu}:
\]
la surface devient une surface r\'egl\'ee, savoir, la \textit{quadricuspidale} de M.~de~la~Gournerie; et, en supposant
\[
(af)^{\frac{1}{3}} + (bg)^{\frac{1}{3}} + (ch)^{\frac{1}{3}} = 0,
\]
et
\[
\lambda : \mu : \nu = (af)^{\frac{1}{3}} : (bg)^{\frac{1}{3}} : (ch)^{\frac{1}{3}},
\]
ce qui donne
\[
\Omega = \Bigl[(af)^{\frac{1}{3}} - (bg)^{\frac{1}{3}}\Bigr]\,
\Bigl[(bg)^{\frac{1}{3}} - (ch)^{\frac{1}{3}}\Bigr]\,
\Bigl[(ch)^{\frac{1}{3}} - (af)^{\frac{1}{3}}\Bigr]:
\]
la surface devient d\'eveloppable.

\begin{flushright}
\textit{Cambridge, 18 Octobre 1866.}
\end{flushright}

\begin{center}
\textit{Deuxi\`eme M\'emoire. Notes pp.~279---283.}
\end{center}

\smallskip\noindent
\textbf{Note I.}\quad
\textsc{Sur la d\'ecomposition du lieu des g\'en\'eratrices en surfaces t\'etra\'edrales distinctes.}

\smallskip
Il me semble qu'une de vos conclusions a besoin d'\^etre modifi\'ee. Ainsi la surface t\'etra\'edrale d\'eriv\'ee de deux courbes triangulaires \`a exposant $\dfrac{1}{m}$ ($m$ \'etant un entier positif), laquelle, selon un de vos th\'eor\`emes, serait de l'ordre $2m^2$, para\^it se d\'ecomposer en $m$ surfaces chacune de l'ordre $2m$. Il y a pour cela une raison \`a \textit{priori}; en effet, pour deux triangulaires de la forme en question, en employant votre construction, on peut \'etablir une correspondance non-seulement entre les deux syst\`emes de points

% [Illustration: Diagram showing Plan $z=0$ and Plan $w=0$ with curves]

\newpage
\setcounter{page}{52}

\noindent
$A, B, \ldots$, et $A', B', \ldots$, mais aussi entre chaque point $A$ et un seul point correspondant $A'$, car l'\'equation de la premi\`ere courbe \'etant de la forme
\[
fx^{\frac{1}{m}} + gy^{\frac{1}{m}} + hz^{\frac{1}{m}} = 0,
\]
on satisfait \`a cette \'equation en \'ecrivant
\[
x : y : z = a(\theta + \alpha)^m : b(\theta + \beta)^m : c(\theta + \gamma)^m,
\]
o\`u $\theta$ est un param\`etre variable. De m\^eme, l'\'equation de la seconde courbe \'etant
\[
fx^{\frac{1}{m}} + gy^{\frac{1}{m}} + k'w^{\frac{1}{m}} = 0,
\]
on satisfait \`a cette \'equation en \'ecrivant
\[
x : y : w = a'(\theta' + \alpha')^m : b'(\theta' + \beta')^m : d'(\theta' + \delta')^m,
\]
o\`u $\theta'$ est aussi un param\`etre variable.

Pour la droite $OP$, on a
\[
\frac{x}{y} = \frac{a\,(\theta + \alpha)^m}{b\,(\theta + \beta)^m},
\]
et pour la droite $O'P'$
\[
\frac{x}{y} = \frac{a'\,(\theta' + \alpha')^m}{b'\,(\theta' + \beta')^m};
\]
donc, la condition pour la correspondance des droites est
\[
\frac{a\,(\theta + \alpha)^m}{b\,(\theta + \beta)^m}
= \lambda\,\frac{a'\,(\theta' + \alpha')^m}{b'\,(\theta' + \beta')^m},
\]
ce qui donne $m$ valeurs diff\'erentes pour $\theta'$ en termes de $\theta$. Mais chacune de ces valeurs est de la forme
\[
\theta' = \frac{A\theta + B}{C\theta + D},
\]
et, en ne faisant attention qu'\`a une seule valeur de $\theta'$, on a le point
\[
x : y : z = a\,(\theta + \alpha)^m : b\,(\theta + \beta)^m : c\,(\theta + \gamma)^m,
\]
qui correspond \`a un point unique
\[
x : y : w = a'\,(\theta' + \alpha')^m : b'\,(\theta' + \beta')^m : d'\,(\theta' + \delta')^m.
\]

Pour le cas de l'exposant $\frac{1}{3}$, on a, de cette mani\`ere, une surface de l'ordre 6. J'ai v\'erifi\'e cela dans le cas particulier de la surface d\'eveloppable. Il est tr\`es-singulier (c'est M.~Salmon qui m'a fait cette remarque) qu'en \'ecrivant dans cette \'equation $(x^2, y^2, z^2, w^2)$ au lieu de $(x, y, z, w)$, on obtient l'\'equation d'une surface du douzi\`eme ordre, lieu des centres de courbure d'un ellipso"ide.

\begin{flushright}
\textit{Cambridge, 15 F\'evrier 1866.}
\end{flushright}

\newpage
\setcounter{page}{53}

\noindent\paperref{433}\quad
\begin{center}
\textbf{SUR LES SURFACES TETRAEDRALES.}
\end{center}

\smallskip\noindent
\textbf{Note II.}\quad
\textsc{A l'occasion de l'ordre des surfaces t\'etra\'edrales.}

\smallskip
Je crois que j'ai n\'eglig\'e de vous faire conna\^itre un th\'eor\`eme assez g\'en\'eral au sujet de l'ordre de ces surfaces. En consid\'erant dans l'espace deux courbes (planes ou \`a double courbure) des ordres $m$ et $m'$ respectivement, et en supposant qu'il y ait entre les points de ces deux courbes une correspondance $(a, a')$, (c'est-\`a-dire qu'\`a un point donn\'e de la courbe $m$ correspondent $a'$ points sur la courbe $m'$, et \`a un point donn\'e de la courbe $m'$ correspondent $a$ points sur la courbe $m$), alors la surface r\'egl\'ee d\'eriv\'ee est de l'ordre $am' + a'm$. De m\^eme, en consid\'erant dans l'espace deux courbes des ordres $n$ et $n'$ avec une correspondance $(b, b')$, la surface d\'eriv\'ee est de l'ordre $bn' + b'n$; et si l'on a en m\^eme temps les deux syst\`emes de courbes avec les correspondances $(a, a')$ et $(b, b')$, la surface t\'etra\'edrale d\'eriv\'ee est de l'ordre
\[
(am' + a'm)(bn' + b'n) + (an' + a'n)(bm' + b'm).
\]
C'est ce que l'on voit sans peine au moyen des \'equations de la surface t\'etra\'edrale.

En effet, si les \'equations des deux courbes du premier syst\`eme sont
\[
(x, y, z)^m = 0, \quad (x, y, w)^{m'} = 0,
\]
et que la condition de correspondance entre deux points $(x, y, z)$ et $(x, y, w)$ soit exprim\'ee par l'\'equation
\[
(x, y, z, w)^{a, a'} = 0,
\]
o\`u $(x, y, z, w)^{a, a'}$ d\'enote une fonction de l'ordre $a$ en $x, y, z$ et de l'ordre $a'$ en $x, y, w$; alors en \'eliminant $x$ et $y$ entre ces trois \'equations, on obtient pour la surface r\'egl\'ee une \'equation de l'ordre $am' + a'm$ en $z$ et $w$. De m\^eme pour le second syst\`eme de courbes, la surface r\'egl\'ee est de l'ordre $bn' + b'n$. Et pour la surface t\'etra\'edrale, en \'eliminant $x$ et $y$ entre les quatre \'equations, on obtient une \'equation dont l'ordre en $z$ et $w$ est
\[
(am' + a'm)(bn' + b'n) + (an' + a'n)(bm' + b'm),
\]
ce qui est le r\'esultat \'enonc\'e.

\newpage
\setcounter{page}{54}

\noindent
\textbf{Note III.}\quad
\textsc{Sur la surface compl\'ementaire.}

\smallskip
Je profite de cette occasion pour rectifier une erreur que j'ai commise dans ma note sur votre M\'emoire des surfaces t\'etra\'edrales. J'ai dit que la surface compl\'ementaire d'une surface t\'etra\'edrale \'etait une surface r\'egl\'ee; cela est vrai pour certaines surfaces t\'etra\'edrales, mais non pas en g\'en\'eral. Le th\'eor\`eme correct est le suivant:

\textit{La surface compl\'ementaire d'une surface t\'etra\'edrale est la surface t\'etra\'edrale d\'eriv\'ee des deux syst\`emes de courbes, chacun pris deux fois.}

Ainsi, si la surface t\'etra\'edrale est d\'eriv\'ee des courbes $(m, m')$ et $(n, n')$, la surface compl\'ementaire est d\'eriv\'ee des courbes $(m, m)$ et $(n', n')$; et si l'on suppose que la correspondance soit dans chaque cas $(1, 1)$, la surface t\'etra\'edrale sera de l'ordre $2mm' \cdot 2nn' = 4mm'nn'$, et la surface compl\'ementaire sera de l'ordre $2m^2 \cdot 2n'^2 = 4m^2n'^2$; ce qui donne pour la somme des deux surfaces l'ordre $4mm'nn' + 4m^2n'^2$.

\begin{flushright}
\textit{Cambridge, 15 F\'evrier 1866.}
\end{flushright}

\bigskip
\noindent\rule{\textwidth}{0.4pt}

% ============================================================
% PAPER 434: ON CERTAIN SKEW SURFACES, OTHERWISE SCROLLS
% ============================================================
\selectlanguage{english}

\newpage
\setcounter{page}{55}

\noindent\paperref{434}\quad
\begin{center}
\textbf{434.}\\[6pt]
\textbf{ON CERTAIN SKEW SURFACES, OTHERWISE SCROLLS.}\\[4pt]
\small
\textit{[From the \textit{Transactions of the Cambridge Philosophical Society}, vol.~\textsc{xi}. Part~\textsc{ii}. (1869),\\
pp.~277---289. Read Nov.~11, 1867.]}
\end{center}

\normalsize
\setlength{\parindent}{1em}
The investigations contained in the present Memoir were suggested to me by the Memoirs of M.~De~la~Goumerie, presented by him to the Academy of Sciences in 1865 and 1866, published in extract in the Comptes Rendus, and reproduced in his work ``Recherches sur les surfaces reglees tetraedrales symetriques, par Jules De la Goumerie, avec des notes par Arthur Cayley,'' 8vo. Paris, 1867. Although the results or the greater part of them are virtually contained in the formulae established by M.~De~la~Goumerie, yet the special theory of the tetraedral surfaces derived from two pairs of curves of the orders $(m, m')$ and $(n, n')$ respectively, does not appear to have been explicitly developed; and I was led to enter upon it by a very elegant theorem communicated to me by M.~Chasles in September 1865, and subsequently published in the Comptes Rendus, as a portion of the theorem ``Sur les courbes \`a double courbure et sur les surfaces r\'egl\'ees,'' t.~\textsc{lxi}. p.~859 (1865 Oct.~9), viz. if in a tetraedron $ABCD$ we inscribe a variable triangle the sides whereof meet the edges in fixed anharmonic ratios, then the three generating lines of the three hyperboloids $(AB, CD)$, $(AC, BD)$, $(AD, BC)$ respectively, meet in a point; and conversely, if they meet in a point, the anharmonic ratios are equal; or say the variable triangle is inscribed in the tetraedron with a given form; the locus of the point of intersection is a skew surface or \textit{scroll}, the order of which was to be determined.

The particular case of Chasles's theorem, where the variable triangle reduces itself to a point, gives the theorem that the locus of a point such that its projections from the four edges on the opposite faces are generating lines of the hyperboloids, is a sextic surface, the locus of the point of intersection of the three generating lines: this surface was mentioned as a sextic scroll by M.~Chasles, and the same surface was considered by me in the paper ``On Tetrahedral Surfaces,'' Phil.~Trans. vol.~\textsc{cliii}. (1863), pp.~417--436, see p.~420.

\newpage
\setcounter{page}{56}

\noindent\paperref{434}\quad
\textbf{ON CERTAIN SKEW SURFACES, OTHERWISE SCROLLS.}

\smallskip
The Memoir is divided into three Parts.

\textsc{Part I.}\quad Article Nos. 1 to 9. On the scrolls in general.

\textsc{Part II.}\quad Article Nos. 10 to 14. On the particular class of scrolls.

\textsc{Part III.}\quad Article Nos. 15 to 29. The Curves $U$, $U'$ are henceforward ``triangular'' curves.

\smallskip
\textbf{Article Nos.~1 to 9. On the scrolls in general.}

\smallskip
\textit{Generation of the scroll.}

1.~Consider any two curves $U$, $U'$ (in space, planes or skew); let their orders be $m$, $m'$ respectively, and let there be between the points of the two curves a correspondence $(\alpha, \alpha')$, viz. to each point of $U$ correspond $\alpha'$ points of $U'$, and to each point of $U'$ correspond $\alpha$ points of $U$; the number of united points (or points corresponding to each other) on $U$ is $= \alpha + \alpha'$; and in like manner the number on $U'$ is $= \alpha + \alpha'$.

2.~In particular let $U$, $U'$ be the sections of a scroll $S$ by two fixed planes; the scroll is generated by the lines which join the corresponding points of the two curves; the class of generation is $(\alpha, \alpha')$; and the order of the scroll is $= m\alpha' + m'\alpha$. This includes as a particular case the generation of a scroll by means of two director curves and a director line: viz. if the director line meet the planes of $U$, $U'$ in the points $O$, $O'$ respectively, and if $OP$, $O'P'$ are corresponding generating lines, then the points $P$, $P'$ on the director curves have a correspondence $(1, 1)$, and the order of the scroll is $= m + m'$.

3.~This conclusion may be otherwise established by means of the section of the scroll by an arbitrary plane: viz. taking the section by the plane of $U$, this is made up of the curve $U$ taken $\alpha'$ times (or, what is the same thing, the complete section is a curve of the order $m\alpha'$), together with the $\alpha$ lines through the several intersections of the plane with the director line; the total order of the section is $= m\alpha' + \alpha$; but this is equal to the order of the scroll into the number of generating lines through an arbitrary point of the plane of $U$: an arbitrary point of the plane has with the curve $U$ a $(1, m)$ correspondence, and therefore with the curve $U'$ a $(\alpha, m\alpha')$ correspondence; hence through the point pass $m\alpha'$ generating lines of the scroll, and the order of the scroll is thus $= m\alpha' + m'\alpha$.

\newpage
\setcounter{page}{57}

\noindent\paperref{434}\quad
\textbf{ON CERTAIN SKEW SURFACES, OTHERWISE SCROLLS.}

4.~There are on the scroll certain curves $V$ of the order $m$, viz. the section by the plane of $U$ consists of the curve $U$ ($= m$) together with the generating lines through the $m$ intersections of the plane with the director line; each of these lines meets $U$ in a single point, and the $m$ points thus obtained are points of a curve $V$ of the order $m$. There are in like manner curves $V'$ of the order $m'$.

5.~The number of these curves $V$ is $= \alpha$; and the like number of curves $V'$ is $= \alpha'$.

6.~In particular the curves $U$, $U'$ may reduce themselves each to a point: we have then a scroll of the order 2, viz. this is a quadric surface, which may be considered as a scroll generated by a variable line which meets each of two given lines, and also each of two given points; or say the scroll has for its director lines two lines, and for its director points two points.

\textit{Section of the scroll by the plane of $U$ or of $U'$.}

7.~By what precedes, the scroll $S$ ($m\alpha' + m'\alpha$) meets the plane of $U$ in the curve $U$ taken $\alpha'$ times (order $m\alpha'$) together with $\alpha m'$ generating lines; and it meets the plane of $U'$ in the curve $U'$ taken $\alpha$ times (order $m'\alpha$) together with $\alpha' m$ generating lines.

\textit{Section of the scroll by the tangent plane at a point of $U$ or of $U'$.}

8.~The foregoing results will subsist if for the arbitrary plane we substitute the tangent plane at a point $A$ of the curve $U$: the section by the tangent plane at $A$ is made up of the curve $U$ ($= m$) taken $\alpha'$ times, order $m\alpha'$, together with the generating lines through the $m'$ intersections (other than $A$) of the tangent plane with the director line. But each of these generating lines meets $U$ in $\alpha$ points (besides the point $A$), and there are thus $\alpha m' - \alpha$ of these points, giving rise to $\alpha m' - \alpha$ generating lines which lie in the tangent plane; the complete section by the tangent plane is therefore of the order $m\alpha' + \alpha m' - \alpha$, that is, $= (m + m')\alpha' - \alpha$, or say $= m\alpha' + m'\alpha - \alpha$, which is the order of the scroll, and the section is therefore the complete section.

\newpage
\setcounter{page}{58}

\noindent\paperref{434}\quad
\textbf{ON CERTAIN SKEW SURFACES, OTHERWISE SCROLLS.}

9.~The section by the tangent plane through $A$ is made up of $(m - \alpha)(m' - \alpha') - \alpha\alpha'$ generating lines (viz. these are, the line $V'\!A$ $\alpha'(m - \omega - \alpha)$ times, the line $V\!A$ $\alpha(m' - \omega' - \alpha')$ times, and the line $AA'$ $\alpha\alpha'$ times, if $\omega$, $\omega'$ are certain integers), and of the curve $U$ taken $\alpha'$ times (order $m\alpha'$), and the curve $U'$ taken $\alpha$ times (order $m'\alpha$).

\bigskip
\textbf{Article Nos.~10 to 14. On the particular class of scrolls.}

\smallskip
\textit{10. Geometrical Construction of a class of Scrolls.}

1.~Consider any two curves $\Omega$, $\Omega'$ of the same deficiency, and such that to each point of $\Omega$ there corresponds a single tangent of $\Omega'$, and to each tangent of $\Omega'$ a single tangent of $\Omega$ (this assumes that the curves $\Omega$, $\Omega'$ are rational transformations one of the other, and that they have consequently the same Deficiency). This being so, let the points of $U$ which lie on any tangent of $\Omega$ and the points of $U'$ which lie on the corresponding tangent of $\Omega'$ be taken to be corresponding points of $U$, $U'$.

The correspondence is then $(\mu\mu', \mu'\mu)$; in fact through a given point of $U$ there pass $\mu$ tangents of $\Omega$, and the corresponding $\mu'$ tangents of $\Omega'$ each intersect $U'$ in a single point $P'$; that is, to a given point $P$ correspond $\mu'$ points $P'$; and similarly to a given point $P'$ correspond $\mu$ points $P$.

And this being so, the number of united points, that is, points of $U$ through which pass corresponding tangents of $\Omega$, $\Omega'$, is $= \mu + \mu'$.

\newpage
\setcounter{page}{59}

\noindent\paperref{434}\quad
\textbf{ON CERTAIN SKEW SURFACES, OTHERWISE SCROLLS.}

2.~The generation of the scroll is obvious: to each point $P$ of $U$ we have the corresponding tangent of $\Omega$, and taking on this the $\alpha'$ corresponding points of $U'$, we have through $P$ the $\alpha'$ generating lines of the scroll; and the order of the scroll is $= m\alpha' + m'\alpha$.

3.~This conclusion may be otherwise established as follows: the scroll is generated by the lines which join corresponding points of $U$ and $U'$; these lines meet each of the planes of $U$, $U'$ in a single point, and therefore meet any plane whatever in a single point; that is, they are generating lines of a scroll. And the order of the scroll is $= m\alpha' + m'\alpha$, as appears by the section made by the plane of $U$.

4.~There are on the scroll curves $V$ of the order $m$, and curves $V'$ of the order $m'$; the number of curves $V$ is $= \alpha$, and the number of curves $V'$ is $= \alpha'$.

5.~The section by the plane of $U$ is made up of the curve $U$ taken $\alpha'$ times (order $m\alpha'$), together with $\alpha m'$ generating lines.

6.~In particular the curves $\Omega$, $\Omega'$ may reduce themselves each to a point: the scroll is then a quadric surface, which may be considered as generated by a variable line meeting each of two given lines and each of two given points.

7.~By what precedes, the scroll $S$ ($m\alpha' + m'\alpha$) meets the plane of $U$ in the curve $U$ taken $\alpha'$ times (order $m\alpha'$) together with $\alpha m'$ generating lines; and it meets the plane of $U'$ in the curve $U'$ taken $\alpha$ times (order $m'\alpha$) together with $\alpha' m$ generating lines.

8.~The section by the tangent plane through $A$ is made up of the curve $U$ taken $\alpha'$ times (order $m\alpha'$), together with $\alpha m' - \alpha$ generating lines, and is of the order $m\alpha' + \alpha m' - \alpha = m\alpha' + m'\alpha - \alpha$, which is the order of the scroll.

\newpage
\setcounter{page}{60}

\noindent\paperref{434}\quad
\textbf{ON CERTAIN SKEW SURFACES, OTHERWISE SCROLLS.}

9.~The section by the tangent plane through $A$ is made up of $(m - \alpha)(m' - \alpha') - \alpha\alpha'$ generating lines (viz. these are, the line $V'\!A$ $\alpha'(m - \omega - \alpha)$ times, the line $V\!A$ $\alpha(m' - \omega' - \alpha')$ times, and the line $AA'$ $\alpha\alpha'$ times), and of the curve $U$ taken $\alpha'$ times (order $m\alpha'$), and the curve $U'$ taken $\alpha$ times (order $m'\alpha$).

\bigskip
\textbf{Article Nos.~15 to 29. The Curves $U$, $U'$ are henceforward ``triangular'' curves.}

\smallskip
\textit{15.}
Let $r = \pm\dfrac{p}{q}$, where $p$, $q$ are positive integers prime to each other, and let the given sections be
\begin{align*}
z &= 0, \quad Ax^r + By^r \quad\; + Dw^r = 0,\\
w &= 0, \quad A'x^r + B'y^r + C'z^r \quad\; = 0,
\end{align*}
where it is to be observed that $r$ being $= +\dfrac{p}{q}$, the two given sections are of the order $pq$, the order of the scroll is $= 2p^2q^2$; each of the given sections is a $pq$-tuple line on the scroll, and the plane thereof meets the scroll in the section taken $pq$ times, and in the $pq$ generating lines: but $r$ being $= -\dfrac{p}{q}$, the two given sections are each of the order $2pq$, with three $pq$-tuple points ($\omega = \alpha = \beta = pq$, $\omega' = \alpha' = \beta' = pq$), and thence the order of the scroll is $\frac{1}{2}(2pq)^2 = 2p^2q^2$; each of the sections is a $pq$-tuple line on the scroll, and the plane meets the scroll only in the section taken $pq$ times. But in either case, if $q$ be $> 1$, that is, if $r$ be fractional, it will presently appear that the scroll of the order $2p^2q^2$ breaks up into $q$ scrolls each of the order $2p^2q$.

\newpage
\setcounter{page}{61}

\noindent\paperref{434}\quad
\textbf{ON CERTAIN SKEW SURFACES, OTHERWISE SCROLLS.}

\textit{16. Analytical Theory of the scrolls.}

11.~It will be convenient to consider first the case $r = +\dfrac{p}{q}$; and to simplify the formulae I write
\[
\xi = x^{\frac{1}{q}}, \quad \eta = y^{\frac{1}{q}}, \quad \zeta = z^{\frac{1}{q}}, \quad \omega = w^{\frac{1}{q}}.
\]
The equations of the two curves $U$, $U'$ then are
\begin{align*}
\zeta^q &= 0, \quad A\xi^p + B\eta^p + D\omega^p = 0,\\
\omega^q &= 0, \quad A'\xi^p + B'\eta^p + C'\zeta^p = 0;
\end{align*}
and the equations of any two corresponding points thereof may be taken to be
\begin{alignat*}{2}
(x, y, z, w) &= (\alpha(\theta + \alpha)^q,\; &&\beta(\theta + \beta)^q,\; \gamma(\theta + \gamma)^q,\; 0),\\
(x, y, z, w) &= (\alpha'(\theta' + \alpha')^q,\; &&\beta'(\theta' + \beta')^q,\; 0,\; \delta'(\theta' + \delta')^q),
\end{alignat*}
respectively; whence the equations of the line joining the two points are
\begin{align*}
\frac{x}{\alpha(\theta + \alpha)^q} &= \frac{y}{\beta(\theta + \beta)^q} = \frac{z}{\gamma(\theta + \gamma)^q},\\
\frac{x}{\alpha'(\theta' + \alpha')^q} &= \frac{y}{\beta'(\theta' + \beta')^q} = \frac{w}{\delta'(\theta' + \delta')^q};
\end{align*}
that is,
\[
\frac{x}{\alpha(\theta + \alpha)^q} = \frac{y}{\beta(\theta + \beta)^q},
\qquad
\frac{z}{\gamma(\theta + \gamma)^q} = 0,
\]
and
\[
\frac{x}{\alpha'(\theta' + \alpha')^q} = \frac{y}{\beta'(\theta' + \beta')^q},
\qquad
\frac{w}{\delta'(\theta' + \delta')^q} = 0.
\]

\newpage
\setcounter{page}{62}

\noindent\paperref{434}\quad
\textbf{ON CERTAIN SKEW SURFACES, OTHERWISE SCROLLS.}

12.~Or, what is the same thing, the line is the intersection of the two planes
\begin{align*}
\frac{x}{\alpha(\theta + \alpha)^q} - \frac{y}{\beta(\theta + \beta)^q} &= 0,
& \frac{z}{\gamma(\theta + \gamma)^q} &= 0,\\
\frac{x}{\alpha'(\theta' + \alpha')^q} - \frac{y}{\beta'(\theta' + \beta')^q} &= 0,
& \frac{w}{\delta'(\theta' + \delta')^q} &= 0.
\end{align*}

13.~Substituting the values $X = \theta Y$, $X' = \theta' Y'$, we have
\[
M(\theta Y, Y, W) = 0, \quad M(\theta' Y', Y', W') = 0,
\]
and thence eliminating $\theta$, $\theta'$, we obtain the equation of the scroll.

14.~We may find the sections of the scroll by the coordinate planes, and the sections by the planes through the coordinate axes, without obtaining the equation of the scroll itself. Thus the section by the plane $z = 0$ is found by eliminating $\theta$, $\theta'$ between the three equations
\[
\frac{x}{\alpha(\theta + \alpha)^q} = \frac{y}{\beta(\theta + \beta)^q},
\qquad
\frac{x}{\alpha'(\theta' + \alpha')^q} = \frac{y}{\beta'(\theta' + \beta')^q},
\qquad
\frac{w}{\delta'(\theta' + \delta')^q} = 0.
\]

\newpage
\setcounter{page}{63}

\noindent\paperref{434}\quad
\textbf{ON CERTAIN SKEW SURFACES, OTHERWISE SCROLLS.}

15.~The section by the plane $z = 0$ is thus a curve of the order $2pq$, viz. it is composed of the two curves
\begin{align*}
w &= 0, \quad Ax^r + By^r = 0,\\
w &= 0, \quad A'x^r + B'y^r = 0,
\end{align*}
each taken $q$ times.

16.~In like manner the section by the plane $w = 0$ is a curve of the order $2pq$, composed of the two curves
\begin{align*}
z &= 0, \quad Ax^r + By^r = 0,\\
z &= 0, \quad A'x^r + B'y^r = 0,
\end{align*}
each taken $q$ times.

17.~The section by the plane $x = 0$ is found by eliminating $\theta$, $\theta'$ between the three equations
\[
\frac{1}{\alpha(\theta + \alpha)^q} = \frac{y}{\beta(\theta + \beta)^q},
\qquad
\frac{1}{\alpha'(\theta' + \alpha')^q} = \frac{y}{\beta'(\theta' + \beta')^q},
\qquad
\frac{z}{\gamma(\theta + \gamma)^q} = 0,\;
\frac{w}{\delta'(\theta' + \delta')^q} = 0.
\]

18.~The section by the plane $x = 0$ is a curve of the order $2pq$, composed of the curves
\begin{align*}
y &= 0, \quad C'z^r + Dw^r = 0,\\
y &= 0, \quad Cz^r + D'w^r = 0,
\end{align*}
each taken $q$ times.

\newpage
\setcounter{page}{64}

\noindent\paperref{434}\quad
\textbf{ON CERTAIN SKEW SURFACES, OTHERWISE SCROLLS.}

19.~The section by the plane $y = 0$ is a curve of the order $2pq$, composed of the curves
\begin{align*}
x &= 0, \quad C'z^r + Dw^r = 0,\\
x &= 0, \quad Cz^r + D'w^r = 0,
\end{align*}
each taken $q$ times.

20.~The section by the plane through the axis $z$, $w$ (that is, the plane $x = 0$, $y = 0$) is a curve of the order $2pq$, composed of the two lines $z = 0$, $w = 0$ each taken $pq$ times.

21.~The section by the plane through the axis $x$, $y$ (that is, the plane $z = 0$, $w = 0$) is a curve of the order $2pq$, composed of the two lines $x = 0$, $y = 0$ each taken $pq$ times.

22.~The section by the plane through the axis $x$, $z$ (that is, the plane $y = 0$, $w = 0$) is a curve of the order $2pq$, composed of the two lines $x = 0$, $z = 0$ each taken $pq$ times.

23.~The section by the plane through the axis $y$, $w$ (that is, the plane $x = 0$, $z = 0$) is a curve of the order $2pq$, composed of the two lines $y = 0$, $w = 0$ each taken $pq$ times.

24.~It may be remarked that the sections by the planes $z = 0$ and $w = 0$ each consist of two curves of the order $pq$, each taken $q$ times; and that the sections by the planes $x = 0$ and $y = 0$ each consist of two curves of the order $pq$, each taken $q$ times.

25.~The sections by the planes through the coordinate axes each consist of two lines, each taken $pq$ times.

26.~We see that the scroll is a surface of the order $2p^2q^2$, having the curves $U$ and $U'$ each as $pq$-tuple lines; and the planes of these curves meet the scroll in these curves taken $pq$ times, and in $pq$ generating lines.

\newpage
\setcounter{page}{65}

\noindent\paperref{434}\quad
\textbf{ON CERTAIN SKEW SURFACES, OTHERWISE SCROLLS.}

27.~If $q > 1$, that is, if $r$ be fractional, the scroll breaks up into $q$ scrolls, each of the order $2p^2q$; and each of these scrolls has the curves $U$ and $U'$ each as $p$-tuple lines, and the planes of these curves meet each scroll in these curves taken $p$ times, and in $p$ generating lines.

28.~The case $r = -\dfrac{p}{q}$ is in like manner obtained by writing in the formulae for $r = +\dfrac{p}{q}$, $x^{-1}$, $y^{-1}$, $z^{-1}$, $w^{-1}$ in place of $x$, $y$, $z$, $w$ respectively. The resulting scroll is of the same order $2p^2q^2$, but the nature of the singularities is different.

29.~The scrolls of the order $2p^2q$ are generated by lines which join corresponding points of two curves $U$, $U'$; and the correspondence is established by means of two curves $\Omega$, $\Omega'$ of the same deficiency, such that to each point of $\Omega$ corresponds a single tangent of $\Omega'$, and to each tangent of $\Omega'$ corresponds a single tangent of $\Omega$.

\bigskip
\begin{flushright}
\textit{Cambridge, Nov.~11, 1867.}
\end{flushright}

\bigskip
\noindent\rule{\textwidth}{0.4pt}

% ============================================================
% PAPER 435: ON THE SIX COORDINATES OF A LINE
% ============================================================
\newpage
\setcounter{page}{66}

\noindent\paperref{435}\quad
\begin{center}
\textbf{435.}\\[6pt]
\textbf{ON THE SIX COORDINATES OF A LINE.}\\[4pt]
\small
\textit{[From the \textit{Transactions of the Cambridge Philosophical Society}, vol.~\textsc{xi}. Part~\textsc{ii}. (1871),\\
pp.~231---238. Read March~9, 1868.]}
\end{center}

\normalsize
\setlength{\parindent}{1em}
\textit{Introduction. Article Nos.~1 to 8.}

1.~In the geometry of three dimensions, a line may be determined by means of six coordinates $(a, b, c, f, g, h)$, which are not independent, but are connected by the equation $af + bg + ch = 0$. These coordinates may be regarded as the coefficients of a bilinear form, or as the elements of an anti-symmetric matrix; and they serve to express the condition that two lines may intersect, or that a line may belong to a given complex, congruence, or regulus. The coordinates were first introduced by Grassmann in 1844, and independently by Pl"ucker in 1846; and they have been extensively employed by Cayley, Salmon, and other writers.

2.~If $(x, y, z, w)$ are the coordinates of a point, and $(X, Y, Z, W)$ the coordinates of a plane, then the line joining the two points $(x, y, z, w)$ and $(x', y', z', w')$ has for its six coordinates
\begin{align*}
a &= y z' - y' z, & b &= z x' - z' x, & c &= x y' - x' y,\\
f &= x w' - x' w, & g &= y w' - y' w, & h &= z w' - z' w;
\end{align*}
and these satisfy the relation $af + bg + ch = 0$.

\newpage
\setcounter{page}{67}

\noindent\paperref{435}\quad
\textbf{ON THE SIX COORDINATES OF A LINE.}

3.~The six coordinates may also be expressed in terms of the coordinates of two planes. If $(X, Y, Z, W)$ and $(X', Y', Z', W')$ are the coordinates of two planes, then their line of intersection has for its six coordinates
\begin{align*}
a &= X W' - X' W, & b &= Y W' - Y' W, & c &= Z W' - Z' W,\\
f &= Y Z' - Y' Z, & g &= Z X' - Z' X, & h &= X Y' - X' Y;
\end{align*}
and these also satisfy the relation $af + bg + ch = 0$.

4.~The condition that two lines $(a, b, c, f, g, h)$ and $(a_1, b_1, c_1, f_1, g_1, h_1)$ may intersect is
\[
a f_1 + b g_1 + c h_1 + f a_1 + g b_1 + h c_1 = 0.
\]
This may be written in the form
\[
\begin{vmatrix}
a & b & c & f & g & h \\
a_1 & b_1 & c_1 & f_1 & g_1 & h_1
\end{vmatrix}
= 0,
\]
where the expression on the left denotes the sum of the products of corresponding elements, taken with the proper signs.

5.~If the two lines coincide, we have the conditions
\[
\frac{a}{a_1} = \frac{b}{b_1} = \frac{c}{c_1} = \frac{f}{f_1} = \frac{g}{g_1} = \frac{h}{h_1},
\]
or, what is the same thing,
\[
a : b : c : f : g : h = a_1 : b_1 : c_1 : f_1 : g_1 : h_1.
\]

\newpage
\setcounter{page}{68}

\noindent\paperref{435}\quad
\textbf{ON THE SIX COORDINATES OF A LINE.}

6.~The condition that the line $(a, b, c, f, g, h)$ may belong to the linear complex
\[
AX + BY + CZ + DW + EU + FV = 0
\]
is
\[
Aa + Bb + Cc + Df + Eg + Fh = 0.
\]
The coordinates of the line satisfy this linear equation; and conversely, any six quantities satisfying this equation may be taken as the coordinates of a line belonging to the complex.

7.~The condition that the line $(a, b, c, f, g, h)$ may belong to the linear congruence
\begin{align*}
A_1 X + B_1 Y + C_1 Z + D_1 W + E_1 U + F_1 V &= 0,\\
A_2 X + B_2 Y + C_2 Z + D_2 W + E_2 U + F_2 V &= 0,
\end{align*}
is that its coordinates satisfy both of the linear equations
\begin{align*}
A_1 a + B_1 b + C_1 c + D_1 f + E_1 g + F_1 h &= 0,\\
A_2 a + B_2 b + C_2 c + D_2 f + E_2 g + F_2 h &= 0.
\end{align*}

8.~The condition that the line $(a, b, c, f, g, h)$ may belong to the regulus
\[
\Phi(X, Y, Z, W, U, V) = 0
\]
is that its coordinates satisfy the quadratic equation
\[
\Phi(a, b, c, f, g, h) = 0.
\]

\newpage
\setcounter{page}{69}

\noindent\paperref{435}\quad
\textbf{ON THE SIX COORDINATES OF A LINE.}

formulae of transformation to a new set of coordinate planes $x_0 = 0$, $y_0 = 0$, $z_0 = 0$, $w_0 = 0$. And I shall also show in what manner the absolute magnitudes of the coordinates may be fixed. But deferring the consideration of these questions, I consider the planes $x = 0$, $y = 0$, $z = 0$, $w = 0$ as given planes, and take the six coordinates $(a, b, c, f, g, h)$ of a line to be determined as above in reference to these given planes, the absolute values of these coordinates remaining indeterminate, and their ratios only being attended to. And I proceed to consider the various questions which present themselves in the geometry of the line, considered as thus determined by means of its six coordinates $(a, b, c, f, g, h)$.

\smallskip
\textbf{Article, Nos.~9 to 18. (Various Sub-headings.) Elementary Theorems.}

\smallskip
\textit{Condition that a line may be in a given plane.}

9.~Taking the line to be $(a, b, c, f, g, h)$, the equation of the given plane to be
\[
(A, B, C, D \,\backslash\!\!\backslash\, x, y, z, w) = 0;
\]
then if $(\alpha, \beta, \gamma, \delta)$, $(\alpha', \beta', \gamma', \delta')$ are the coordinates of any two points on the line, we have the system of equations \textit{ante}, No.~2, and substituting therein for $\beta\gamma' - \beta'\gamma$, \&c. the values $(a, b, c, f, g, h)$, we find
\[
\begin{vmatrix}
0, & c, & -b, & f \\
-c, & 0, & a, & g \\
b, & -a, & 0, & h \\
-f, & -g, & -h, & 0
\end{vmatrix}
(A, B, C, D) = 0;
\]
which equations, equivalent to a twofold relation, are the required condition. It may be remarked that, treating $(A, B, C, D)$ as current plane coordinates, each equation of the system is that of a point lying in the line.

\newpage
\setcounter{page}{70}

\noindent\paperref{435}\quad
\textbf{ON THE SIX COORDINATES OF A LINE.}

\textit{Condition that a line may pass through a given point.}

10.~The coordinates of the given point are taken to be $(\alpha, \beta, \gamma, \delta)$. If
\[
(A, B, C, D \,\backslash\!\!\backslash\, x, y, z, w) = 0, \quad
(A', B', C', D' \,\backslash\!\!\backslash\, x, y, z, w) = 0,
\]
are the equations of any two planes through the line, then we have the system of equations \textit{ante} No.~3, and substituting therein for $AB' - A'B$, \&c. their values in terms of the coordinates $(a, b, c, f, g, h)$ of the line, we have
\[
\begin{vmatrix}
0, & h, & -g, & a \\
-h, & 0, & f, & b \\
g, & -f, & 0, & c \\
-a, & -b, & -c, & 0
\end{vmatrix}
(\alpha, \beta, \gamma, \delta) = 0;
\]
which equations, equivalent to a twofold relation, are the required condition. It is obvious that, treating $(\alpha, \beta, \gamma, \delta)$ as current point coordinates, each equation of the system is the equation of a plane through the given line.

\textit{Condition that two lines may intersect.}

11.~For two lines $(a, b, c, f, g, h)$ and $(a_1, b_1, c_1, f_1, g_1, h_1)$ the condition of intersection is, as already mentioned,
\[
a f_1 + b g_1 + c h_1 + f a_1 + g b_1 + h c_1 = 0.
\]

\newpage
\setcounter{page}{71}

\noindent\paperref{435}\quad
\textbf{ON THE SIX COORDINATES OF A LINE.}

It may be remarked that the expression on the left-hand side is the determinant of the matrix
\[
\begin{pmatrix}
a & b & c & f & g & h \\
a_1 & b_1 & c_1 & f_1 & g_1 & h_1
\end{pmatrix},
\]
and that the vanishing of this determinant expresses that the two lines intersect. The determinant may be written in the form
\[
\begin{vmatrix}
a & h & -g & a_1 & h_1 & -g_1 \\
-h & b & f & -h_1 & b_1 & f_1 \\
g & -f & c & g_1 & -f_1 & c_1 \\
a & b & c & a_1 & b_1 & c_1
\end{vmatrix}
= 0,
\]
or in other equivalent forms.

\textit{Condition that a line may meet each of four given lines.}

12.~If the four given lines are $(a_1, b_1, c_1, f_1, g_1, h_1)$, $(a_2, b_2, c_2, f_2, g_2, h_2)$, $(a_3, b_3, c_3, f_3, g_3, h_3)$, $(a_4, b_4, c_4, f_4, g_4, h_4)$, then the conditions that the line $(a, b, c, f, g, h)$ may meet each of them are
\begin{align*}
a f_1 + b g_1 + c h_1 + f a_1 + g b_1 + h c_1 &= 0,\\
a f_2 + b g_2 + c h_2 + f a_2 + g b_2 + h c_2 &= 0,\\
a f_3 + b g_3 + c h_3 + f a_3 + g b_3 + h c_3 &= 0,\\
a f_4 + b g_4 + c h_4 + f a_4 + g b_4 + h c_4 &= 0.
\end{align*}
These are four linear equations in the six coordinates $(a, b, c, f, g, h)$, which together with the quadratic relation $af + bg + ch = 0$ determine the ratios of these coordinates. The number of solutions is in general $= \infty^1$, that is, the lines meeting four given lines form a regulus.

\newpage
\setcounter{page}{72}

\noindent\paperref{435}\quad
\textbf{ON THE SIX COORDINATES OF A LINE.}

\textit{Condition that a line may belong to a given complex of the second order.}

13.~The general equation of a complex of the second order is
\[
(A, B, C, F, G, H \,\backslash\!\!\backslash\, a, b, c, f, g, h)^2 = 0,
\]
where $(A, B, C, F, G, H \,\backslash\!\!\backslash\, a, b, c, f, g, h)^2$ denotes a quadratic function of the six coordinates. The condition that the line $(a, b, c, f, g, h)$ may belong to the complex is that its coordinates satisfy this equation.

14.~The complex of the second order may be represented as the assemblage of lines which satisfy the equation
\[
\lambda(a, b, c, f, g, h)^2 + \mu(a, b, c, f, g, h)(a_1, b_1, c_1, f_1, g_1, h_1) + \nu(a_1, b_1, c_1, f_1, g_1, h_1)^2 = 0,
\]
where $\lambda$, $\mu$, $\nu$ are parameters, and $(a, b, c, f, g, h)$, $(a_1, b_1, c_1, f_1, g_1, h_1)$ are the coordinates of two fixed lines.

\textit{Condition that a line may belong to a given congruence of the second order.}

15.~The general equation of a congruence of the second order is given by the intersection of two complexes of the second order; the coordinates of the lines belonging to the congruence satisfy two quadratic equations.

\newpage
\setcounter{page}{73}

\noindent\paperref{435}\quad
\textbf{ON THE SIX COORDINATES OF A LINE.}

16.~The congruence of the second order may be represented as the assemblage of lines which satisfy the equations
\begin{align*}
&(A_1, B_1, C_1, F_1, G_1, H_1 \,\backslash\!\!\backslash\, a, b, c, f, g, h)^2 = 0,\\
&(A_2, B_2, C_2, F_2, G_2, H_2 \,\backslash\!\!\backslash\, a, b, c, f, g, h)^2 = 0,
\end{align*}
where the two quadratic forms are independent.

\textit{Condition that a line may belong to a given regulus.}

17.~The regulus is a special case of the congruence; it is the assemblage of lines which meet each of three given lines. The condition that the line $(a, b, c, f, g, h)$ may belong to a regulus is that its coordinates satisfy three linear equations, which express that the line meets each of the three given lines.

18.~Conversely, any three linear equations in the six coordinates determine a regulus; for they express that the line meets each of three given lines, and the lines meeting three given lines form a regulus.

\newpage
\setcounter{page}{74}

\noindent\paperref{435}\quad
\textbf{ON THE SIX COORDINATES OF A LINE.}

\smallskip
\textbf{Article, Nos.~19 to 35. Various Formulae.}

\smallskip
\textit{Expression for the coordinates of a line in terms of the coordinates of two points.}

19.~If $(x, y, z, w)$ and $(x', y', z', w')$ are two points on the line, then the six coordinates of the line are given by the formulae
\begin{align*}
a &= y z' - y' z, & b &= z x' - z' x, & c &= x y' - x' y,\\
f &= x w' - x' w, & g &= y w' - y' w, & h &= z w' - z' w;
\end{align*}
and conversely, if $(a, b, c, f, g, h)$ are the coordinates of a line, and $(x, y, z, w)$ is a point on the line, then the coordinates of any other point on the line may be expressed in terms of these quantities.

20.~If $(X, Y, Z, W)$ and $(X', Y', Z', W')$ are the coordinates of two planes through the line, then the six coordinates of the line are given by the formulae
\begin{align*}
a &= X W' - X' W, & b &= Y W' - Y' W, & c &= Z W' - Z' W,\\
f &= Y Z' - Y' Z, & g &= Z X' - Z' X, & h &= X Y' - X' Y.
\end{align*}

\newpage
\setcounter{page}{75}

\noindent\paperref{435}\quad
\textbf{ON THE SIX COORDINATES OF A LINE.}

\textit{Expression for the coordinates of a line in terms of the coordinates of two planes.}

21.~If $(X, Y, Z, W)$ and $(X', Y', Z', W')$ are the coordinates of two planes through the line, then the six coordinates of the line may be expressed in the form
\begin{align*}
a &= \begin{vmatrix} X & W \\ X' & W' \end{vmatrix},
& b &= \begin{vmatrix} Y & W \\ Y' & W' \end{vmatrix},
& c &= \begin{vmatrix} Z & W \\ Z' & W' \end{vmatrix},\\
f &= \begin{vmatrix} Y & Z \\ Y' & Z' \end{vmatrix},
& g &= \begin{vmatrix} Z & X \\ Z' & X' \end{vmatrix},
& h &= \begin{vmatrix} X & Y \\ X' & Y' \end{vmatrix}.
\end{align*}
These formulae are analogous to those of No.~19, and they also satisfy the relation $af + bg + ch = 0$.

22.~The formulae of Nos.~19 and 20 are connected by the principle of duality. The coordinates of a line may be regarded as point-coordinates or as plane-coordinates; and the formulae for the one case are transformed into those for the other case by the substitution of plane-coordinates for point-coordinates.

23.~If $(x, y, z, w)$ is a point on the line $(a, b, c, f, g, h)$, then the coordinates satisfy the equations
\begin{align*}
c y - b z + g w &= 0,\\
-c x + a z - f w &= 0,\\
b x - a y + h w &= 0,\\
-g x + f y - h z &= 0.
\end{align*}
These are the equations which express that the point $(x, y, z, w)$ lies on the line $(a, b, c, f, g, h)$.

\newpage
\setcounter{page}{76}

\noindent\paperref{435}\quad
\textbf{ON THE SIX COORDINATES OF A LINE.}

24.~If $(X, Y, Z, W)$ is a plane through the line $(a, b, c, f, g, h)$, then the coordinates satisfy the equations
\begin{align*}
c Y - b Z + g W &= 0,\\
-c X + a Z - f W &= 0,\\
b X - a Y + h W &= 0,\\
-g X + f Y - h Z &= 0.
\end{align*}
These are the equations which express that the plane $(X, Y, Z, W)$ passes through the line $(a, b, c, f, g, h)$.

25.~The six coordinates $(a, b, c, f, g, h)$ may be expressed in terms of four arbitrary quantities. If we write
\[
a = p q' - p' q, \quad b = r s' - r' s, \quad c = t u' - t' u,
\]
\[
f = p r' - p' r, \quad g = q s' - q' s, \quad h = t v' - t' v,
\]
where $p$, $q$, $r$, $s$, $t$, $u$, $v$, $p'$, $q'$, $r'$, $s'$, $t'$, $u'$, $v'$ are arbitrary quantities subject to the condition $af + bg + ch = 0$, then these expressions give the general solution of the equation $af + bg + ch = 0$.

26.~The general solution may also be written in the form
\begin{align*}
a &= \lambda_1 \mu_2 - \lambda_2 \mu_1, & b &= \lambda_1 \nu_2 - \lambda_2 \nu_1, & c &= \lambda_1 \rho_2 - \lambda_2 \rho_1,\\
f &= \mu_1 \nu_2 - \mu_2 \nu_1, & g &= \nu_1 \rho_2 - \nu_2 \rho_1, & h &= \rho_1 \mu_2 - \rho_2 \mu_1,
\end{align*}
where $\lambda_1$, $\mu_1$, $\nu_1$, $\rho_1$, $\lambda_2$, $\mu_2$, $\nu_2$, $\rho_2$ are arbitrary quantities.

\newpage
\setcounter{page}{77}

\noindent\paperref{435}\quad
\textbf{ON THE SIX COORDINATES OF A LINE.}

\textit{Expression for the moment of two lines.}

27.~If $(a, b, c, f, g, h)$ and $(a_1, b_1, c_1, f_1, g_1, h_1)$ are the coordinates of two lines, the moment of the two lines is defined as the product of the perpendicular distance between them into the sine of the angle between them. The expression for the moment is
\[
M = a f_1 + b g_1 + c h_1 + f a_1 + g b_1 + h c_1.
\]
This expression is symmetrical with respect to the two lines, and it vanishes when the lines intersect.

28.~The moment may also be expressed in the form
\[
M = \sqrt{(af + bg + ch)(a_1 f_1 + b_1 g_1 + c_1 h_1)} \cdot \sin\theta,
\]
where $\theta$ is the angle between the lines. But since $af + bg + ch = 0$ and $a_1 f_1 + b_1 g_1 + c_1 h_1 = 0$, this expression is indeterminate; and the proper expression is that given in No.~27.

29.~The square of the moment is given by the formula
\[
M^2 = (a f_1 + b g_1 + c h_1 + f a_1 + g b_1 + h c_1)^2.
\]
This may be expressed as a determinant; but the expression is complicated, and it is generally more convenient to use the formula of No.~27.

\newpage
\setcounter{page}{78}

\noindent\paperref{435}\quad
\textbf{ON THE SIX COORDINATES OF A LINE.}

\textit{Expression for the distance of a point from a line.}

30.~If $(x, y, z, w)$ is a point, and $(a, b, c, f, g, h)$ is a line, then the distance of the point from the line is given by the formula
\[
d = \frac{\sqrt{(c y - b z + g w)^2 + (-c x + a z - f w)^2 + (b x - a y + h w)^2}}{\sqrt{a^2 + b^2 + c^2 + f^2 + g^2 + h^2}}.
\]
This formula may be simplified when the coordinates are normalized so that $a^2 + b^2 + c^2 + f^2 + g^2 + h^2 = 1$.

31.~The distance may also be expressed in terms of the moment of the line and a line through the point parallel to the given line. If $(a_1, b_1, c_1, f_1, g_1, h_1)$ is a line through the point $(x, y, z, w)$ parallel to the line $(a, b, c, f, g, h)$, then the distance of the point from the line is equal to the moment of the two lines divided by the square root of the sum of the squares of the coordinates.

\textit{Expression for the angle between two lines.}

32.~If $(a, b, c, f, g, h)$ and $(a_1, b_1, c_1, f_1, g_1, h_1)$ are two lines, the angle $\theta$ between them is given by the formula
\[
\cos\theta = \frac{a a_1 + b b_1 + c c_1 + f f_1 + g g_1 + h h_1}{\sqrt{a^2 + b^2 + c^2 + f^2 + g^2 + h^2}\,\sqrt{a_1^2 + b_1^2 + c_1^2 + f_1^2 + g_1^2 + h_1^2}}.
\]

\newpage
\setcounter{page}{79}

\noindent\paperref{435}\quad
\textbf{ON THE SIX COORDINATES OF A LINE.}

\textit{Case of the fivefold relation.}

33.~The fivefold relation
\[
\begin{Vmatrix}
a, & b, & c, & f, & g, & h \\
a_1, & b_1, & c_1, & f_1, & g_1, & h_1 \\
a_2, & b_2, & c_2, & f_2, & g_2, & h_2 \\
a_3, & b_3, & c_3, & f_3, & g_3, & h_3 \\
a_4, & b_4, & c_4, & f_4, & g_4, & h_4
\end{Vmatrix}
= 0
\]
expresses that the line $(a, b, c, f, g, h)$ belongs to the linear complex determined by the four lines $(a_1, b_1, \ldots)$, $(a_2, b_2, \ldots)$, $(a_3, b_3, \ldots)$, $(a_4, b_4, \ldots)$. This relation may be regarded as the condition that the five lines have a common transversal; or that they belong to the same linear congruence.

34.~The fivefold relation may be written in the form of a determinant of the fifth order; but it is generally more convenient to use the form given in No.~33. The relation is equivalent to five linear equations in the coordinates $(a, b, c, f, g, h)$, which express that the line belongs to the linear complex determined by the four given lines.

\textit{Case of the fourfold relation.}

35.~The fourfold relation
\[
\begin{Vmatrix}
a, & b, & c, & f, & g, & h \\
a_1, & b_1, & c_1, & f_1, & g_1, & h_1 \\
a_2, & b_2, & c_2, & f_2, & g_2, & h_2 \\
a_3, & b_3, & c_3, & f_3, & g_3, & h_3
\end{Vmatrix}
= 0
\]
expresses that the line $(a, b, c, f, g, h)$ belongs to the regulus determined by the three lines $(a_1, b_1, \ldots)$, $(a_2, b_2, \ldots)$, $(a_3, b_3, \ldots)$. This relation is equivalent to four linear equations in the coordinates $(a, b, c, f, g, h)$, which express that the line meets each of the three given lines.

\newpage
\setcounter{page}{80}

\noindent\paperref{435}\quad
\textbf{ON THE SIX COORDINATES OF A LINE.}

\smallskip
\textbf{Article, Nos.~36 to 58. Further Developments.}

\smallskip
\textit{Expression for the coordinates of the line of intersection of two planes.}

36.~If $(X, Y, Z, W)$ and $(X', Y', Z', W')$ are the coordinates of two planes, the six coordinates of their line of intersection are given by the formulae
\begin{align*}
a &= X W' - X' W, & b &= Y W' - Y' W, & c &= Z W' - Z' W,\\
f &= Y Z' - Y' Z, & g &= Z X' - Z' X, & h &= X Y' - X' Y;
\end{align*}
and these satisfy the relation $af + bg + ch = 0$.

37.~The coordinates of the line joining two points $(x, y, z, w)$ and $(x', y', z', w')$ are given by the formulae
\begin{align*}
a &= y z' - y' z, & b &= z x' - z' x, & c &= x y' - x' y,\\
f &= x w' - x' w, & g &= y w' - y' w, & h &= z w' - z' w.
\end{align*}

38.~The condition that the line $(a, b, c, f, g, h)$ may meet the line $(a_1, b_1, c_1, f_1, g_1, h_1)$ is
\[
a f_1 + b g_1 + c h_1 + f a_1 + g b_1 + h c_1 = 0.
\]
This condition may be written in the form of a determinant; but the form given above is generally more convenient.

\newpage
\setcounter{page}{81}

\noindent\paperref{435}\quad
\textbf{ON THE SIX COORDINATES OF A LINE.}

\textit{Expression for the coordinates of a line in terms of six arbitrary quantities.}

39.~The six coordinates $(a, b, c, f, g, h)$ of a line are connected by the single equation $af + bg + ch = 0$; they may therefore be expressed in terms of five independent quantities. If we introduce a sixth quantity $k$, and write
\begin{align*}
a &= k \alpha, & b &= k \beta, & c &= k \gamma,\\
f &= k \phi, & g &= k \psi, & h &= k \chi,
\end{align*}
where $\alpha$, $\beta$, $\gamma$, $\phi$, $\psi$, $\chi$ are quantities satisfying the equation $\alpha\phi + \beta\psi + \gamma\chi = 0$, then the ratios $a : b : c : f : g : h$ are determined by the five quantities $\alpha : \beta : \gamma : \phi : \psi : \chi$.

40.~The absolute values of the coordinates may be fixed by imposing an additional condition; for example, we may assume that $a^2 + b^2 + c^2 + f^2 + g^2 + h^2 = 1$, or that one of the coordinates is equal to unity. But it is generally more convenient to leave the absolute values indeterminate, and to consider only the ratios of the coordinates.

41.~The coordinates of a line may be normalized so that $a^2 + b^2 + c^2 + f^2 + g^2 + h^2 = 1$; and when this is done, the distance of a point from the line, and the angle between two lines, may be expressed in simpler forms. But the normalization is not essential, and it is often more convenient to use the unnormalized coordinates.

\newpage
\setcounter{page}{82}

\noindent\paperref{435}\quad
\textbf{ON THE SIX COORDINATES OF A LINE.}

\textit{Transformation of coordinates.}

42.~If we change the system of coordinate planes, the six coordinates of a line undergo a linear transformation. Let the new coordinate planes be $x_0 = 0$, $y_0 = 0$, $z_0 = 0$, $w_0 = 0$, where
\begin{align*}
x_0 &= a_1 x + b_1 y + c_1 z + d_1 w,\\
y_0 &= a_2 x + b_2 y + c_2 z + d_2 w,\\
z_0 &= a_3 x + b_3 y + c_3 z + d_3 w,\\
w_0 &= a_4 x + b_4 y + c_4 z + d_4 w.
\end{align*}
Then if $(a, b, c, f, g, h)$ are the coordinates of a line referred to the old system, and $(a_0, b_0, c_0, f_0, g_0, h_0)$ the coordinates of the same line referred to the new system, we have
\begin{align*}
a_0 &= (a_2 b_3 - a_3 b_2) a + (a_3 b_1 - a_1 b_3) b + (a_1 b_2 - a_2 b_1) c\\
&\quad + (a_2 c_3 - a_3 c_2) f + (a_3 c_1 - a_1 c_3) g + (a_1 c_2 - a_2 c_1) h,\\
&\ \,\vdots
\end{align*}
with similar formulae for $b_0$, $c_0$, $f_0$, $g_0$, $h_0$.

43.~These formulae of transformation are linear in the coordinates $(a, b, c, f, g, h)$, and they preserve the relation $af + bg + ch = 0$. They may be written in the matrix form
\[
\begin{pmatrix} a_0 \\ b_0 \\ c_0 \\ f_0 \\ g_0 \\ h_0 \end{pmatrix}
= M
\begin{pmatrix} a \\ b \\ c \\ f \\ g \\ h \end{pmatrix},
\]
where $M$ is a $6 \times 6$ matrix whose elements are the minors of the matrix of transformation of point-coordinates.

\newpage
\setcounter{page}{83}

\noindent\paperref{435}\quad
\textbf{ON THE SIX COORDINATES OF A LINE.}

44.~The matrix $M$ satisfies the condition that it preserves the bilinear form $af + bg + ch$; that is, if $(a, b, c, f, g, h)$ satisfy $af + bg + ch = 0$, then $(a_0, b_0, c_0, f_0, g_0, h_0)$ also satisfy $a_0 f_0 + b_0 g_0 + c_0 h_0 = 0$. This condition is expressed by the equation
\[
M^T J M = J,
\]
where $J$ is the matrix
\[
J = \begin{pmatrix}
0 & 0 & 0 & 1 & 0 & 0 \\
0 & 0 & 0 & 0 & 1 & 0 \\
0 & 0 & 0 & 0 & 0 & 1 \\
1 & 0 & 0 & 0 & 0 & 0 \\
0 & 1 & 0 & 0 & 0 & 0 \\
0 & 0 & 1 & 0 & 0 & 0
\end{pmatrix}.
\]

45.~The transformation formulae may also be expressed in terms of the coordinates of two planes. If $(X, Y, Z, W)$ and $(X', Y', Z', W')$ are the coordinates of two planes, and $(X_0, Y_0, Z_0, W_0)$, $(X'_0, Y'_0, Z'_0, W'_0)$ are their coordinates referred to the new system, then the coordinates of the line of intersection are transformed by the same formulae as above.

46.~The absolute magnitudes of the coordinates may be fixed by imposing the condition that the determinant of the matrix of transformation shall be equal to unity. When this is done, the coordinates are said to be \textit{normalized}; and the formulae of transformation are then completely determined.

\newpage
\setcounter{page}{84}

\noindent\paperref{435}\quad
\textbf{ON THE SIX COORDINATES OF A LINE.}

\textit{Condition that four lines may belong to the same regulus.}

47.~If $(a_1, b_1, c_1, f_1, g_1, h_1)$, $(a_2, b_2, c_2, f_2, g_2, h_2)$, $(a_3, b_3, c_3, f_3, g_3, h_3)$, $(a_4, b_4, c_4, f_4, g_4, h_4)$ are four lines, the condition that they may belong to the same regulus is that the determinant
\[
\begin{vmatrix}
a_1 & b_1 & c_1 & f_1 & g_1 & h_1 \\
a_2 & b_2 & c_2 & f_2 & g_2 & h_2 \\
a_3 & b_3 & c_3 & f_3 & g_3 & h_3 \\
a_4 & b_4 & c_4 & f_4 & g_4 & h_4
\end{vmatrix}
\]
shall be of rank not exceeding 3; that is, all the determinants of the fourth order formed from it shall vanish.

48.~This condition is equivalent to the three conditions that the four lines shall have a common transversal; or that they shall belong to the same linear congruence. It may also be expressed as the condition that the six coordinates of any one of the lines can be expressed as linear functions of the coordinates of the other three.

49.~The condition that four lines may belong to the same regulus may be written in the form
\[
\begin{vmatrix}
a_1 f_2 + b_1 g_2 + c_1 h_2 + f_1 a_2 + g_1 b_2 + h_1 c_2 & \cdots \\ \vdots & \ddots
\end{vmatrix}
= 0,
\]
where the determinant is of the fourth order, and its elements are the moments of pairs of lines.

\newpage
\setcounter{page}{85}

\noindent\paperref{435}\quad
\textbf{ON THE SIX COORDINATES OF A LINE.}

\textit{Condition that five lines may have a common transversal.}

50.~If $(a_1, b_1, c_1, f_1, g_1, h_1)$, $(a_2, b_2, c_2, f_2, g_2, h_2)$, $(a_3, b_3, c_3, f_3, g_3, h_3)$, $(a_4, b_4, c_4, f_4, g_4, h_4)$, $(a_5, b_5, c_5, f_5, g_5, h_5)$ are five lines, the condition that they may have a common transversal is that the determinant
\[
\begin{vmatrix}
a_1 & b_1 & c_1 & f_1 & g_1 & h_1 \\
a_2 & b_2 & c_2 & f_2 & g_2 & h_2 \\
a_3 & b_3 & c_3 & f_3 & g_3 & h_3 \\
a_4 & b_4 & c_4 & f_4 & g_4 & h_4 \\
a_5 & b_5 & c_5 & f_5 & g_5 & h_5
\end{vmatrix}
= 0,
\]
where the determinant is of the fifth order. This condition is equivalent to the condition that the five lines belong to the same linear complex; or that they have $\infty^1$ common transversals.

51.~The condition may also be expressed as the vanishing of the determinant of the matrix whose rows are the coordinates of the five lines. This determinant is of the fifth order, and its vanishing expresses that the five lines are linearly dependent.

\textit{Condition that six lines may have two common transversals.}

52.~If $(a_1, b_1, c_1, f_1, g_1, h_1)$, $\ldots$, $(a_6, b_6, c_6, f_6, g_6, h_6)$ are six lines, the condition that they may have two common transversals is that the determinant
\[
\begin{vmatrix}
a_1 & b_1 & c_1 & f_1 & g_1 & h_1 \\
\vdots & \vdots & \vdots & \vdots & \vdots & \vdots \\
a_6 & b_6 & c_6 & f_6 & g_6 & h_6
\end{vmatrix}
= 0,
\]
where the determinant is of the sixth order. This condition expresses that the six lines belong to the same linear congruence; or that they have $\infty^2$ common transversals.

\newpage
\setcounter{page}{86}

\noindent\paperref{435}\quad
\textbf{ON THE SIX COORDINATES OF A LINE.}

53.~The condition that six lines may have two common transversals may also be expressed as the condition that the six lines are linearly dependent; that is, that there exist quantities $\lambda_1$, $\lambda_2$, $\ldots$, $\lambda_6$, not all zero, such that
\[
\lambda_1 (a_1, b_1, c_1, f_1, g_1, h_1) + \cdots + \lambda_6 (a_6, b_6, c_6, f_6, g_6, h_6) = (0, 0, 0, 0, 0, 0).
\]
This condition is equivalent to the vanishing of the determinant of the sixth order formed from the coordinates of the six lines.

\textit{Expression for the coordinates of the common transversal of four lines.}

54.~If four lines are given, the coordinates of their common transversal may be found by solving the four linear equations which express that the transversal meets each of the given lines. Let $(a, b, c, f, g, h)$ be the coordinates of the transversal; then the conditions that it meets the four given lines are
\begin{align*}
a f_1 + b g_1 + c h_1 + f a_1 + g b_1 + h c_1 &= 0,\\
a f_2 + b g_2 + c h_2 + f a_2 + g b_2 + h c_2 &= 0,\\
a f_3 + b g_3 + c h_3 + f a_3 + g b_3 + h c_3 &= 0,\\
a f_4 + b g_4 + c h_4 + f a_4 + g b_4 + h c_4 &= 0.
\end{align*}
These four equations, together with $af + bg + ch = 0$, determine the ratios of the six coordinates.

55.~The coordinates of the transversal may be expressed as determinants of the fourth order formed from the coordinates of the four given lines. If we write
\[
A = \begin{vmatrix} f_1 & g_1 & h_1 & a_1 \\ \vdots & & & \\ f_4 & g_4 & h_4 & a_4 \end{vmatrix},
\quad \text{etc.},
\]
then the coordinates of the transversal are proportional to $(A, B, C, F, G, H)$.

\newpage
\setcounter{page}{87}

\noindent\paperref{435}\quad
\textbf{ON THE SIX COORDINATES OF A LINE.}

56.~The coordinates of the transversal may also be found by the following construction. Let the four given lines be $L_1$, $L_2$, $L_3$, $L_4$; through $L_1$ and $L_2$ draw two arbitrary planes, and let their line of intersection be $M_{12}$; through $L_3$ and $L_4$ draw two arbitrary planes, and let their line of intersection be $M_{34}$. Then the common transversal of $L_1$, $L_2$, $L_3$, $L_4$ is the line of intersection of the two planes which contain $M_{12}$ and $M_{34}$ respectively, and which pass through the line of intersection of the planes of $L_1$, $L_2$ and $L_3$, $L_4$.

57.~This construction gives a geometrical method of finding the common transversal of four lines; and it may be extended to the case of five or more lines. The construction is based on the principle that the common transversal of four lines is the line of intersection of two planes, each of which contains two of the given lines.

\textit{Expression for the coordinates of the line of shortest distance between two lines.}

58.~If $(a, b, c, f, g, h)$ and $(a_1, b_1, c_1, f_1, g_1, h_1)$ are two lines, the coordinates of their common perpendicular (or line of shortest distance) may be found as follows. Let $(\lambda, \mu, \nu, \rho, \sigma, \tau)$ be the coordinates of the common perpendicular; then the conditions that it meets each of the two lines are
\begin{align*}
\lambda f + \mu g + \nu h + \rho a + \sigma b + \tau c &= 0,\\
\lambda f_1 + \mu g_1 + \nu h_1 + \rho a_1 + \sigma b_1 + \tau c_1 &= 0.
\end{align*}
And the condition that it is perpendicular to each of the two lines is that its direction-cosines are proportional to the differences of the direction-cosines of the two lines.

\newpage
\setcounter{page}{88}

\noindent\paperref{435}\quad
\textbf{ON THE SIX COORDINATES OF A LINE.}

\smallskip
\textbf{Article, Nos.~59 to 75. Statical and Kinematical Applications.}

\smallskip
The coordinates $(a, b, c, f, g, h)$, as last defined, are peculiarly convenient in kinematical and mechanical questions, as will appear from the following investigations.

\textit{Rotations.}

59.~Using the term rotation to denote an infinitesimal rotation, I say first that a rotation $\lambda$ round the line $(a, b, c, f, g, h)$ produces in the point $(x, y, z)$ rigidly connected with this line the displacements
\begin{align*}
\delta x &= \lambda(\phantom{hx} - hy + gz - a),\\
\delta y &= \lambda(\phantom{hx} hx \phantom{+} - fz - b),\\
\delta z &= \lambda(-gx + fy \phantom{+ gz} - c).
\end{align*}

\textit{Forces.}

60.~A force $R$ acting along the line $(a, b, c, f, g, h)$ has for its components along the axes
\begin{align*}
X &= R a, & Y &= R b, & Z &= R c,
\end{align*}
and its moments about the axes are
\begin{align*}
L &= R f, & M &= R g, & N &= R h.
\end{align*}
These six quantities satisfy the relation $a f + b g + c h = 0$, which is the condition that the force may be a single force acting along a definite line.

\newpage
\setcounter{page}{89}

\noindent\paperref{435}\quad
\textbf{ON THE SIX COORDINATES OF A LINE.}

\textit{Couples.}

61.~A couple $G$ in the plane $(l, m, n, p)$ has for its components
\begin{align*}
L &= G l, & M &= G m, & N &= G n,
\end{align*}
and its forces along the axes are
\begin{align*}
X &= G p, & Y &= G q, & Z &= G r,
\end{align*}
where $p$, $q$, $r$ are determined by the condition that the system is equivalent to a couple.

62.~The six components of a couple satisfy the relation $a f + b g + c h = 0$, which is the condition that the system may be reducible to a single force or a couple. In fact, the general system of forces acting on a rigid body may be reduced to a single force acting along a certain line, together with a couple in a plane perpendicular to that line.

\textit{Wrenches.}

63.~A \textit{wrench} is a force $R$ acting along a line $(a, b, c, f, g, h)$, together with a couple $G$ in a plane perpendicular to the line. The six components of the wrench are
\begin{align*}
X &= R a, & Y &= R b, & Z &= R c,\\
L &= R f + G a, & M &= R g + G b, & N &= R h + G c.
\end{align*}
These six quantities are connected by the relation
\[
a L + b M + c N - f X - g Y - h Z = 0,
\]
which is the condition that the system may be a wrench.

\newpage
\setcounter{page}{90}

\noindent\paperref{435}\quad
\textbf{ON THE SIX COORDINATES OF A LINE.}

64.~The pitch of the wrench is defined as the ratio $p = G/R$ of the couple to the force. The six components may then be written
\begin{align*}
X &= R a, & Y &= R b, & Z &= R c,\\
L &= R(f + pa), & M &= R(g + pb), & N &= R(h + pc).
\end{align*}
And the relation between them becomes
\[
a L + b M + c N - f X - g Y - h Z = p(a X + b Y + c Z) = 0,
\]
since $a X + b Y + c Z = R(a^2 + b^2 + c^2)$.

65.~The axis of the wrench is the line $(a, b, c, f, g, h)$; and the pitch is the ratio of the couple to the force. The wrench is completely determined by its axis and its pitch; and conversely, any line and any pitch determine a wrench.

66.~If the pitch is zero, the wrench reduces to a single force; if the pitch is infinite, the wrench reduces to a couple. The general wrench is therefore intermediate between a force and a couple.

\textit{Screws.}

67.~A \textit{screw} is a line associated with a pitch. The six coordinates of a screw are the six coordinates of the line, together with the pitch. The screw may be regarded as a wrench of unit force; and its six components are
\begin{align*}
X &= a, & Y &= b, & Z &= c,\\
L &= f + pa, & M &= g + pb, & N &= h + pc.
\end{align*}

\newpage
\setcounter{page}{91}

\noindent\paperref{435}\quad
\textbf{ON THE SIX COORDINATES OF A LINE.}

68.~The theory of screws was developed by Sir Robert Ball in his ``Theory of Screws'' (1876). A screw is determined by its axis $(a, b, c, f, g, h)$ and its pitch $p$; and it may be regarded as a geometrical entity consisting of a line and a scalar quantity associated with it.

69.~Two screws are said to be \textit{reciprocal} when the virtual coefficient of the one with respect to the other vanishes. If $(a, b, c, f, g, h, p)$ and $(a_1, b_1, c_1, f_1, g_1, h_1, p_1)$ are two screws, the condition of reciprocity is
\[
a f_1 + b g_1 + c h_1 + f a_1 + g b_1 + h c_1 + p(a a_1 + b b_1 + c c_1) + p_1(a_1 a + b_1 b + c_1 c) = 0.
\]

70.~A system of screws is said to be a \textit{co-reciprocal system} when every pair of screws in the system are reciprocal. Such a system contains at most six screws; and when it contains exactly six, the screws are said to form a \textit{screw-six}.

\textit{Applications to Mechanics.}

71.~The six coordinates of a line are of great importance in the theory of the equilibrium and motion of rigid bodies. The condition of equilibrium of a rigid body under the action of any system of forces is that the sum of the components and the sum of the moments about any line shall vanish.

\newpage
\setcounter{page}{92}

\noindent\paperref{435}\quad
\textbf{ON THE SIX COORDINATES OF A LINE.}

72.~If $(X_i, Y_i, Z_i, L_i, M_i, N_i)$ are the components of the $i$th force, then the conditions of equilibrium are
\begin{align*}
\sum X_i &= 0, & \sum Y_i &= 0, & \sum Z_i &= 0,\\
\sum L_i &= 0, & \sum M_i &= 0, & \sum N_i &= 0.
\end{align*}
These six equations express that the resultant force and the resultant couple vanish.

73.~The six coordinates of a line are also useful in the theory of virtual work. If a rigid body undergoes a small displacement, the work done by the forces acting on it is equal to the sum of the products of the components of the forces into the corresponding displacements.

74.~If $(\delta x, \delta y, \delta z)$ is the displacement of a point, and $(\delta\theta_x, \delta\theta_y, \delta\theta_z)$ is the rotation of the body, then the virtual work is
\[
\sum (X_i \delta x + Y_i \delta y + Z_i \delta z + L_i \delta\theta_x + M_i \delta\theta_y + N_i \delta\theta_z) = 0.
\]
This equation expresses the principle of virtual work for a rigid body.

75.~The six coordinates of a line are also applicable to the theory of attractions and potentials. If a body is attracted by a system of forces, the potential of the system may be expressed in terms of the six coordinates of the lines of action of the forces.

\newpage
\setcounter{page}{93}

\noindent\paperref{435}\quad
\textbf{ON THE SIX COORDINATES OF A LINE.}

\smallskip
\textbf{Article, Nos.~76 to 99. Further Analytical Developments.}

\smallskip
\textit{Expression for the condition that six lines may be generators of the same quadric surface.}

76.~If six lines are generators of the same quadric surface, their coordinates satisfy a certain condition. Let the equation of the quadric surface be
\[
Ax^2 + By^2 + Cz^2 + Dw^2 + 2Fyz + 2Gzx + 2Hxy + 2Lxw + 2Myw + 2Nzw = 0.
\]
Then the condition that the line $(a, b, c, f, g, h)$ may be a generator of the quadric is
\[
A a^2 + B b^2 + C c^2 + D f^2 + E g^2 + F h^2 + 2F bc + 2G ca + 2H ab + 2L af + 2M bg + 2N ch = 0.
\]
This is a quadratic equation in the coordinates of the line; and the condition that six lines may be generators of the same quadric is that their coordinates satisfy this equation.

77.~The condition may be written in the form
\[
\sum_{i,j} Q_{ij} a_i a_j = 0,
\]
where $Q_{ij}$ are the coefficients of the quadric, and $a_i$ are the coordinates of the line. This is a quadratic condition; and the six lines satisfy it when they are generators of the same quadric.

\newpage
\setcounter{page}{94}

\noindent\paperref{435}\quad
\textbf{ON THE SIX COORDINATES OF A LINE.}

\textit{Expression for the condition that a line may touch a given quadric surface.}

78.~If the quadric surface is given by the equation
\[
Ax^2 + By^2 + Cz^2 + Dw^2 + 2Fyz + 2Gzx + 2Hxy + 2Lxw + 2Myw + 2Nzw = 0,
\]
then the condition that the line $(a, b, c, f, g, h)$ may touch the quadric is that the discriminant of the quadratic form vanish. This condition is
\[
\begin{vmatrix}
A & H & G & L \\
H & B & F & M \\
G & F & C & N \\
L & M & N & D
\end{vmatrix}
= 0.
\]
This is a quartic condition in the coordinates of the line.

79.~The condition may also be expressed as the condition that the line is a tangent to the quadric; that is, that it meets the quadric in two coincident points. This condition is equivalent to the condition that the line is a generator of the quadric, or that it lies altogether on the quadric.

\textit{Expression for the condition that a line may belong to a given linear complex.}

80.~The condition that the line $(a, b, c, f, g, h)$ may belong to the linear complex
\[
AX + BY + CZ + DW + EU + FV = 0
\]
is
\[
Aa + Bb + Cc + Df + Eg + Fh = 0.
\]
This is a linear condition in the coordinates of the line; and the lines belonging to a linear complex form a threefold system.

\newpage
\setcounter{page}{95}

\noindent\paperref{435}\quad
\textbf{ON THE SIX COORDINATES OF A LINE.}

81.~The linear complex may be represented as the assemblage of lines which satisfy the linear equation
\[
Aa + Bb + Cc + Df + Eg + Fh = 0.
\]
The coefficients $(A, B, C, D, E, F)$ are the coordinates of the complex; and the equation expresses that the line $(a, b, c, f, g, h)$ belongs to the complex.

82.~The condition that a line may belong to a given linear congruence is that its coordinates satisfy two linear equations
\begin{align*}
A_1 a + B_1 b + C_1 c + D_1 f + E_1 g + F_1 h &= 0,\\
A_2 a + B_2 b + C_2 c + D_2 f + E_2 g + F_2 h &= 0.
\end{align*}
The lines belonging to a linear congruence form a twofold system; they are the common lines of two linear complexes.

83.~The condition that a line may belong to a given regulus is that its coordinates satisfy three linear equations. The lines belonging to a regulus form a onefold system; they are the generators of a quadric surface.

\textit{Expression for the condition that four lines may have a common transversal.}

84.~If $(a_1, b_1, c_1, f_1, g_1, h_1)$, $(a_2, b_2, c_2, f_2, g_2, h_2)$, $(a_3, b_3, c_3, f_3, g_3, h_3)$, $(a_4, b_4, c_4, f_4, g_4, h_4)$ are four lines, the condition that they may have a common transversal is that the determinant
\[
\begin{vmatrix}
a_1 & b_1 & c_1 & f_1 & g_1 & h_1 \\
a_2 & b_2 & c_2 & f_2 & g_2 & h_2 \\
a_3 & b_3 & c_3 & f_3 & g_3 & h_3 \\
a_4 & b_4 & c_4 & f_4 & g_4 & h_4
\end{vmatrix}
= 0.
\]

\newpage
\setcounter{page}{96}

\noindent\paperref{435}\quad
\textbf{ON THE SIX COORDINATES OF A LINE.}

This condition is equivalent to the condition that the four lines belong to the same linear congruence; or that the six coordinates of their common transversal satisfy the four linear equations which express that the transversal meets each of the given lines.

85.~The condition may also be expressed as the vanishing of the determinant of the matrix whose rows are the coordinates of the four lines. This determinant is of the fourth order in the elements of each line; and its vanishing expresses that the four lines are linearly dependent.

86.~The condition that four lines may have a common transversal is equivalent to the condition that the four lines belong to the same regulus; for if four lines have a common transversal, they belong to the linear congruence determined by the transversal and any line meeting it; and conversely, if four lines belong to the same regulus, they have a common transversal.

\textit{Expression for the condition that five lines may belong to the same linear complex.}

87.~If $(a_1, b_1, c_1, f_1, g_1, h_1)$, $\ldots$, $(a_5, b_5, c_5, f_5, g_5, h_5)$ are five lines, the condition that they may belong to the same linear complex is that the determinant
\[
\begin{vmatrix}
a_1 & b_1 & c_1 & f_1 & g_1 & h_1 \\
\vdots & \vdots & \vdots & \vdots & \vdots & \vdots \\
a_5 & b_5 & c_5 & f_5 & g_5 & h_5
\end{vmatrix}
= 0.
\]

\newpage
\setcounter{page}{97}

\noindent\paperref{435}\quad
\textbf{ON THE SIX COORDINATES OF A LINE.}

This condition is equivalent to the condition that the five lines are linearly dependent; that is, that there exist quantities $\lambda_1$, $\ldots$, $\lambda_5$, not all zero, such that
\[
\lambda_1 (a_1, b_1, c_1, f_1, g_1, h_1) + \cdots + \lambda_5 (a_5, b_5, c_5, f_5, g_5, h_5) = (0, 0, 0, 0, 0, 0).
\]

88.~The condition may also be expressed as the vanishing of the determinant of the fifth order formed from the coordinates of the five lines. This determinant is a skew-symmetric determinant of the fifth order; and its vanishing expresses that the five lines belong to the same linear complex.

89.~The condition that five lines may belong to the same linear complex is equivalent to the condition that they have $\infty^1$ common transversals; for if five lines belong to a linear complex, any line of the complex meets them all; and conversely, if five lines have $\infty^1$ common transversals, these transversals form a linear complex, and the five lines belong to it.

\textit{Expression for the condition that six lines may belong to the same linear congruence.}

90.~If $(a_1, b_1, c_1, f_1, g_1, h_1)$, $\ldots$, $(a_6, b_6, c_6, f_6, g_6, h_6)$ are six lines, the condition that they may belong to the same linear congruence is that the determinant
\[
\begin{vmatrix}
a_1 & b_1 & c_1 & f_1 & g_1 & h_1 \\
\vdots & \vdots & \vdots & \vdots & \vdots & \vdots \\
a_6 & b_6 & c_6 & f_6 & g_6 & h_6
\end{vmatrix}
= 0.
\]

\newpage
\setcounter{page}{98}

\noindent\paperref{435}\quad
\textbf{ON THE SIX COORDINATES OF A LINE.}

This condition is equivalent to the condition that the six lines are linearly dependent; that is, that there exist quantities $\lambda_1$, $\ldots$, $\lambda_6$, not all zero, such that
\[
\lambda_1 (a_1, b_1, c_1, f_1, g_1, h_1) + \cdots + \lambda_6 (a_6, b_6, c_6, f_6, g_6, h_6) = (0, 0, 0, 0, 0, 0).
\]

91.~The condition may also be expressed as the vanishing of the determinant of the sixth order formed from the coordinates of the six lines. This determinant is a skew-symmetric determinant of the sixth order; and its vanishing expresses that the six lines belong to the same linear congruence.

92.~The condition that six lines may belong to the same linear congruence is equivalent to the condition that they have $\infty^2$ common transversals; for if six lines belong to a linear congruence, any line of the congruence meets them all; and conversely, if six lines have $\infty^2$ common transversals, these transversals form a linear congruence, and the six lines belong to it.

\textit{Expression for the condition that seven lines may have a common transversal.}

93.~If seven lines are given, the condition that they may have a common transversal is that the coordinates of the lines satisfy a certain system of equations. This condition is equivalent to the condition that the seven lines belong to the same linear complex; or that they have $\infty^1$ common transversals.

\newpage
\setcounter{page}{99}

\noindent\paperref{435}\quad
\textbf{ON THE SIX COORDINATES OF A LINE.}

94.~The condition that seven lines may have a common transversal may be expressed as the vanishing of certain determinants formed from the coordinates of the lines. If $(a_1, b_1, c_1, f_1, g_1, h_1)$, $\ldots$, $(a_7, b_7, c_7, f_7, g_7, h_7)$ are the seven lines, then the condition is that all the determinants of the sixth order formed from the matrix
\[
\begin{pmatrix}
a_1 & b_1 & c_1 & f_1 & g_1 & h_1 \\
\vdots & \vdots & \vdots & \vdots & \vdots & \vdots \\
a_7 & b_7 & c_7 & f_7 & g_7 & h_7
\end{pmatrix}
\]
shall vanish. This condition is equivalent to seven independent equations; and when it is satisfied, the seven lines have a common transversal.

95.~The condition may also be expressed as the condition that the seven lines are linearly dependent; that is, that there exist quantities $\lambda_1$, $\ldots$, $\lambda_7$, not all zero, such that
\[
\lambda_1 (a_1, b_1, c_1, f_1, g_1, h_1) + \cdots + \lambda_7 (a_7, b_7, c_7, f_7, g_7, h_7) = (0, 0, 0, 0, 0, 0).
\]

\bigskip
\noindent\rule{\textwidth}{0.4pt}

% ============================================================
% PAPER 436: ON A CERTAIN SEXTIC TORSE
% ============================================================
\newpage
\setcounter{page}{99}

\noindent\paperref{436}\quad
\begin{center}
\textbf{436.}\\[6pt]
\textbf{ON A CERTAIN SEXTIC TORSE.}\\[4pt]
\small
\textit{[From the \textit{Transactions of the Cambridge Philosophical Society}, vol.~\textsc{xi}. Part~\textsc{iii}. (1871),\\
pp.~507---523. Read Nov.~8, 1869.]}
\end{center}

\normalsize
\setlength{\parindent}{1em}
The torse (developable surface) intended to be considered is that which has for its edge of regression an excubo-quartic curve, or say a \textit{unicursal} quartic curve. I call to mind that (excluding the plane quartic) a quartic curve is either a quadriquadric, viz. it is the complete intersection of two quadric surfaces; or else it is an excubo-quartic, viz. there is through the curve only one quadric surface, and the curve is the partial intersection of this quadric surface with a cubic surface through two generating lines (of the same kind) of the quadric surface. Returning to the quadriquadric curve, this may be general, nodal, or cuspidal; viz. if the two quadric surfaces have an ordinary contact, the curve of intersection is a nodal quadriquadric; if they have a stationary contact, the curve is a cuspidal quadriquadric.

The unicursal quartic is a curve such that the coordinates $(x, y, z, w)$ of any point thereof are proportional to rational and integral quartic functions $(*\!\!\backslash\!\!\backslash\, \theta, 1)^4$ of a variable parameter $\theta$; and the general unicursal quartic is in fact the excubo-quartic; but included as particular cases of the unicursal curve (although not as cases of the excubo-quartic as above defined) we have the nodal quadriquadric and the cuspidal quadriquadric. The torse having for its edge of regression a unicursal curve is a sextic torse; and this is in fact the order of the torse derived from the excubo-quartic, and from the nodal quadriquadric; but for the cuspidal quadriquadric, there is a depression of one, and the torse becomes a quintic torse. The equations have been obtained of (1) the sextic torse derived from the nodal quadriquadric, (2) the quintic torse derived from the cuspidal quadriquadric, (3) the sextic torse derived from a certain special excubo-quartic; but the equation of the torse derived from the general unicursal quartic has not yet been found. To show at the outset what the analytical problem is, I

\newpage
\setcounter{page}{100}

\noindent\paperref{436}\quad
\textbf{ON A CERTAIN SEXTIC TORSE.}

remark that the unicursal quartic may be represented by the equations
\begin{align*}
x : y : z : w &= (\theta - \alpha)(\theta - \beta)(\theta - \gamma)(\theta - \delta) \\
& : (\theta - \alpha')(\theta - \beta')(\theta - \gamma')(\theta - \delta') \\
& : (\theta - \alpha'')(\theta - \beta'')(\theta - \gamma'')(\theta - \delta'') \\
& : (\theta - \alpha''')(\theta - \beta''')(\theta - \gamma''')(\theta - \delta'''),
\end{align*}
where $\alpha$, $\beta$, $\gamma$, $\delta$, $\alpha'$, $\beta'$, $\gamma'$, $\delta'$, $\alpha''$, $\beta''$, $\gamma''$, $\delta''$, $\alpha'''$, $\beta'''$, $\gamma'''$, $\delta'''$ are arbitrary constants; or, what is the same thing, the curve is the partial intersection of the quadric surface
\[
\frac{x}{(\theta - \alpha)(\theta - \beta)(\theta - \gamma)(\theta - \delta)}
= \frac{y}{(\theta - \alpha')(\theta - \beta')(\theta - \gamma')(\theta - \delta')}
\]
with the cubic surface which is the locus of the point of intersection of the three planes
\begin{align*}
\frac{x}{(\theta - \alpha)(\theta - \beta)(\theta - \gamma)(\theta - \delta)} &=
\frac{z}{(\theta - \alpha'')(\theta - \beta'')(\theta - \gamma'')(\theta - \delta'')},\\
\frac{y}{(\theta - \alpha')(\theta - \beta')(\theta - \gamma')(\theta - \delta')} &=
\frac{w}{(\theta - \alpha''')(\theta - \beta''')(\theta - \gamma''')(\theta - \delta''')},\\
\frac{x}{(\theta - \alpha)(\theta - \beta)(\theta - \gamma)(\theta - \delta)} &=
\frac{w}{(\theta - \alpha''')(\theta - \beta''')(\theta - \gamma''')(\theta - \delta''')}.
\end{align*}
The torse is the envelope of the osculating plane of the curve; or, what is the same thing, it is the locus of the tangent lines to the curve. The equation of the torse is to be obtained by eliminating $\theta$ between the equations of the tangent line and the condition that the line is tangent to the curve.
\clearpage


% =========================================================
% Chunk: cayley_vol07_pages_101_150.tex  (orig pages 101--150)
% =========================================================
\begin{titlepage}
\begin{center}
\vspace*{2cm}
{\Large\textsc{The Collected Mathematical Papers of}}\\[0.5em]
{\LARGE\textbf{Arthur Cayley, Sc.D., F.R.S.}}\\[0.3em]
{\large Sadlerian Professor of Pure Mathematics in the University of Cambridge}\\[2cm]
{\Huge\textbf{Vol.\ VII}}\\[1cm]
{\large Cambridge: At the University Press. 1894}\\[2cm]
\rule{\textwidth}{0.4pt}\\[1cm]
{\Large\textbf{Transcription of Pages 101--150}}\\[0.5em]
{\normalsize Papers [436] through [445]}\\[1cm]
\rule{\textwidth}{0.4pt}
\end{center}
\end{titlepage}


\newpage

% ============================================================
% PAPER [436] -- ON A CERTAIN SEXTIC TORSE
% Pages 101--115 (PDF pages 122--136)
% ============================================================
\section*{[436]}
\addcontentsline{toc}{section}{[436] On a Certain Sextic Torse}

\begin{center}
\textbf{\Large ON A CERTAIN SEXTIC TORSE.}
\end{center}

\begin{flushright}
\textit{[From the Philosophical Transactions of the Royal Society of London,}\\
\textit{for the year 1868 (Part\ I), pp.\ 377--396.]}
\end{flushright}

\vspace{1em}

\subsection*{Standard Equation of the Unicursal Quartic.}

2. The coordinates $(x, y, z, w)$ being originally taken to be proportional to any four given quartic functions $(\ast\between\theta, 1)^4$ of the parameter $\theta$, then forming a linear function of the coordinates, we have four sets of values of the multipliers, each reducing the function of $\theta$ to a perfect fourth power; that is, writing $(X, Y, Z, W)$ for the linear functions of the original coordinates, and taking $(X, Y, Z, W)$ as coordinates, it appears that the unicursal quartic may be represented by the equations
\[
X : Y : Z : W = (\theta + \alpha)^4 : (\theta + \beta)^4 : (\theta + \gamma)^4 : (\theta + \delta)^4.
\]

\subsection*{Tangent Line, and Osculating Plane of the Unicursal Quartic.}

3. The equations of the tangent line at the point $(\theta)$ (that is, the point the coordinates whereof are as $(\theta + \alpha)^4 : (\theta + \beta)^4 : (\theta + \gamma)^4 : (\theta + \delta)^4$) are at once seen to be
\[
\begin{vmatrix}
X, & Y, & Z, & W \\
(\theta + \alpha)^4, & (\theta + \beta)^4, & (\theta + \gamma)^4, & (\theta + \delta)^4 \\
(\theta + \alpha)^3, & (\theta + \beta)^3, & (\theta + \gamma)^3, & (\theta + \delta)^3
\end{vmatrix} = 0,
\]
and that of the osculating plane to be
\[
\begin{vmatrix}
X, & Y, & Z, & W \\
(\theta + \alpha)^4, & (\theta + \beta)^4, & (\theta + \gamma)^4, & (\theta + \delta)^4 \\
(\theta + \alpha)^3, & (\theta + \beta)^3, & (\theta + \gamma)^3, & (\theta + \delta)^3 \\
(\theta + \alpha)^2, & (\theta + \beta)^2, & (\theta + \gamma)^2, & (\theta + \delta)^2
\end{vmatrix} = 0.
\]

Writing as in the sequel
\begin{align*}
a &= \beta - \gamma, & f &= \alpha - \delta, \\
b &= \gamma - \alpha, & g &= \beta - \delta, \\
c &= \alpha - \beta, & h &= \gamma - \delta,
\end{align*}
the equations of the tangent line become
\begin{align*}
\frac{hY}{(\theta + \beta)^3} - \frac{gZ}{(\theta + \gamma)^3} + \frac{aW}{(\theta + \delta)^3} &= 0, \\
-\frac{hX}{(\theta + \alpha)^3} \qquad\quad + \frac{fZ}{(\theta + \gamma)^3} + \frac{bW}{(\theta + \delta)^3} &= 0, \\
\frac{gX}{(\theta + \alpha)^3} - \frac{fY}{(\theta + \beta)^3} \qquad\quad + \frac{cW}{(\theta + \delta)^3} &= 0, \\
-\frac{aX}{(\theta + \alpha)^3} - \frac{bY}{(\theta + \beta)^3} - \frac{cZ}{(\theta + \gamma)^3} \qquad\quad &= 0,
\end{align*}
(equivalent of course to two equations), and the equation of the osculating plane becomes
\[
\frac{ahg\,X}{(\theta + \alpha)^2} + \frac{hbf\,Y}{(\theta + \beta)^2} + \frac{gfc\,Z}{(\theta + \gamma)^2} + \frac{abc\,W}{(\theta + \delta)^2} = 0.
\]

\subsection*{Modification of the foregoing notation, and final form for the Unicursal Quartic.}

4. If instead of the coordinates $(X, Y, Z, W)$ we introduce the coordinates $(x, y, z, w)$ connected therewith by the relations
\[
ahgX : hbfY : gfcZ : abcW = x : y : z : w,
\]
or, what is the same thing,
\[
X : Y : Z : W = bcfx : cagy : abhz : fghw,
\]
then the curve is given by the equations
\[
x : y : z : w = ahg(\theta + \alpha)^4 : hbf(\theta + \beta)^4 : cfg(\theta + \gamma)^4 : abc(\theta + \delta)^4.
\]

The equations of the tangent line are
\begin{align*}
\frac{cy}{(\theta + \beta)^3} - \frac{bz}{(\theta + \gamma)^3} + \frac{fw}{(\theta + \delta)^3} &= 0, \\
-\frac{cx}{(\theta + \alpha)^3} \qquad\quad + \frac{az}{(\theta + \gamma)^3} + \frac{gw}{(\theta + \delta)^3} &= 0, \\
\frac{bx}{(\theta + \alpha)^3} - \frac{ay}{(\theta + \beta)^3} \qquad\quad + \frac{hw}{(\theta + \delta)^3} &= 0, \\
-\frac{fx}{(\theta + \alpha)^3} - \frac{gy}{(\theta + \beta)^3} - \frac{hz}{(\theta + \gamma)^3} \qquad\quad &= 0,
\end{align*}
and the equation of the osculating plane is
\[
\frac{x}{(\theta + \alpha)^2} + \frac{y}{(\theta + \beta)^2} + \frac{z}{(\theta + \gamma)^2} + \frac{w}{(\theta + \delta)^2} = 0.
\]

\subsection*{Determination of the Sextic Torse.}

5. Starting from the equation of the osculating plane written under the form
\begin{multline*}
x(\theta + \beta)^2(\theta + \gamma)^2(\theta + \delta)^2
+ y(\theta + \gamma)^2(\theta + \delta)^2(\theta + \alpha)^2 \\
+ z(\theta + \delta)^2(\theta + \alpha)^2(\theta + \beta)^2
+ w(\theta + \alpha)^2(\theta + \beta)^2(\theta + \gamma)^2 = 0,
\end{multline*}
the equation of the torse is obtained by equating to zero the discriminant of the sextic function. Writing as before
\begin{align*}
a &= \beta - \gamma, & f &= \alpha - \delta, \\
b &= \gamma - \alpha, & g &= \beta - \delta, \\
c &= \alpha - \beta, & h &= \gamma - \delta,
\end{align*}
equations which give
\begin{alignat*}{3}
h &- g &&+ a &&= 0, \\
-h &\quad &&+ f + b &&= 0, \\
g &- f &&\quad + c &&= 0, \\
-a &- b &&- c &&= 0,
\end{alignat*}
\begin{alignat*}{3}
h\beta &- g\gamma &&+ a\delta &&= 0, \\
-h\alpha &\quad &&+ f\gamma + b\delta &&= 0, \\
g\alpha &- f\beta &\quad + c\delta &&= 0, \\
-a\alpha &- b\beta &- c\gamma &\quad\quad &= 0,
\end{alignat*}
and also
\[
af + bg + ch = 0,
\]
then the discriminant is a function of $(x, y, z, w)$, $(a, b, c, f, g, h)$ of the degree 10 in $(x, y, z, w)$ and the degree 30 in $(a, b, c, f, g, h)$. But the equation in $\theta$ has two equal roots, or the discriminant vanishes, if any one of the quantities $(x, y, z, w)$ is $= 0$; and again, if any one of the differences $\alpha - \beta$, \&c.\ (that is any one of the quantities $a$, $b$, $c$, $f$, $g$, $h$) is $= 0$: the discriminant thus contains the factors $xyzw$ and $(abcfgh)^2$, and throwing these out, we have an equation of the form
\[
\Lambda = (a, b, c, f, g, h)^{12}(x, y, z, w)^6 = 0,
\]
which is the equation of the sextic torse.

\subsection*{Principal Sections of the Torse.}

6. Consider for instance the section by the plane $w = 0$. Writing $w = 0$, the equation of the osculating plane is
\[
(\theta + \delta)^2\bigl[x(\theta + \beta)^2(\theta + \gamma)^2 + y(\theta + \gamma)^2(\theta + \alpha)^2 + z(\theta + \alpha)^2(\theta + \beta)^2\bigr] = 0.
\]
The discriminant of the sextic function vanishes identically in virtue of the double factor $(\theta + \delta)^2$. But omitting this factor, the equation becomes
\[
x(\theta + \beta)^2(\theta + \gamma)^2 + y(\theta + \gamma)^2(\theta + \alpha)^2 + z(\theta + \alpha)^2(\theta + \beta)^2 = 0.
\]
The discriminant of this quartic function of $\theta$ is a function of $x$, $y$, $z$, $a$, $b$, $c$ of the degree 6 in $(x, y, z)$ and 12 in $(a, b, c)$; it contains however the factors $xyz$, $a^2b^2c^2$, and the remaining factor is of the form
\[
\Phi = (a, b, c)^8 (x, y, z)^3 = 0,
\]
which is the equation of a cubic curve; that is, the section by the plane $w = 0$ is made up of the three lines $x = 0$, $y = 0$, $z = 0$ (which are the tangents of the quartic at the points $\alpha$, $\beta$, $\gamma$ respectively), and of a cubic curve, the intersection of the plane $w = 0$ with the torse.

7. More generally, for the section by any plane $lx + my + nz + pw = 0$, writing for convenience $lx + my + nz + pw = 0$ in place of $w = 0$ in the foregoing equation of the osculating plane, the section of the torse is made up of the six lines which are the tangents of the unicursal quartic at the six points given by the equation $l(\theta + \alpha)^2 + m(\theta + \beta)^2 + n(\theta + \gamma)^2 + p(\theta + \delta)^2 = 0$, and of a sextic curve which is the proper intersection of the plane with the torse.

\subsection*{Equation of the Sextic Torse.}

8. The complete expression of the discriminant of the sextic function is very complicated. Using a known formula for the discriminant of a binary sextic $(\ast\between x, y)^6$, or proceeding directly, we may write the sextic equation under the form
\[
(a, b, c, d, e, f, g\between x, y)^6 = 0,
\]
and express the discriminant as a function of the coefficients $(a, b, c, d, e, f, g)$: the calculations are however very laborious. I have in fact obtained the equation of the torse by a different method, which I proceed to explain.

\subsection*{Analytical Theory of the Correspondence of Points on the Unicursal Quartic.}

9. Consider two points of the quartic, the parameters of which are $\theta$, $\theta_1$ respectively. The line joining these two points is a generating line of the torse, if the osculating planes at the two points coincide, that is, if the two points have the same osculating plane. The condition for this may be obtained as follows.

The equation of the osculating plane at the point $\theta$ is
\[
\frac{x}{(\theta + \alpha)^2} + \frac{y}{(\theta + \beta)^2} + \frac{z}{(\theta + \gamma)^2} + \frac{w}{(\theta + \delta)^2} = 0.
\]
Hence the osculating planes at the points $\theta$, $\theta_1$ will be the same, if
\[
\frac{x}{(\theta + \alpha)^2} + \frac{y}{(\theta + \beta)^2} + \frac{z}{(\theta + \gamma)^2} + \frac{w}{(\theta + \delta)^2} = 0,
\]
and
\[
\frac{x}{(\theta_1 + \alpha)^2} + \frac{y}{(\theta_1 + \beta)^2} + \frac{z}{(\theta_1 + \gamma)^2} + \frac{w}{(\theta_1 + \delta)^2} = 0,
\]
for the same values of $(x, y, z, w)$; that is, if the two fractions
\[
\frac{1}{(\theta + \alpha)^2(\theta_1 + \alpha)^2},\quad
\frac{1}{(\theta + \beta)^2(\theta_1 + \beta)^2},\quad
\frac{1}{(\theta + \gamma)^2(\theta_1 + \gamma)^2},\quad
\frac{1}{(\theta + \delta)^2(\theta_1 + \delta)^2}
\]
are proportional to the same quantities $(x, y, z, w)$. The condition is therefore
\[
\begin{vmatrix}
\frac{1}{(\theta + \alpha)^2(\theta_1 + \alpha)^2}, &
\frac{1}{(\theta + \beta)^2(\theta_1 + \beta)^2}, &
\frac{1}{(\theta + \gamma)^2(\theta_1 + \gamma)^2}, &
\frac{1}{(\theta + \delta)^2(\theta_1 + \delta)^2} \\
\frac{1}{(\theta + \alpha)^3(\theta_1 + \alpha)^2}, &
\frac{1}{(\theta + \beta)^3(\theta_1 + \beta)^2}, &
\frac{1}{(\theta + \gamma)^3(\theta_1 + \gamma)^2}, &
\frac{1}{(\theta + \delta)^3(\theta_1 + \delta)^2} \\
\frac{1}{(\theta + \alpha)^2(\theta_1 + \alpha)^3}, &
\frac{1}{(\theta + \beta)^2(\theta_1 + \beta)^3}, &
\frac{1}{(\theta + \gamma)^2(\theta_1 + \gamma)^3}, &
\frac{1}{(\theta + \delta)^2(\theta_1 + \delta)^3}
\end{vmatrix} = 0,
\]
which is the condition for the coincidence of the osculating planes at the two points $\theta$, $\theta_1$.

10. This determinant may be simplified. In fact, multiplying the several rows and columns by proper factors, and using the identity
\[
\frac{1}{(\theta + \alpha)^2(\theta_1 + \alpha)^2} - \frac{1}{(\theta + \beta)^2(\theta_1 + \beta)^2} = \frac{(\alpha - \beta)(\theta + \theta_1 + \alpha + \beta)}{(\theta + \alpha)^2(\theta + \beta)^2(\theta_1 + \alpha)^2(\theta_1 + \beta)^2},
\]
\&c., the condition may be expressed in a more manageable form. Writing for shortness
\[
\theta + \theta_1 = 2\sigma,\quad \theta\theta_1 = \pi,
\]
the condition of correspondence of the two points $\theta$, $\theta_1$ on the unicursal quartic assumes the form of an equation of the sixth order in $(\sigma, \pi)$, which may be written
\[
\Psi(\sigma, \pi) = 0.
\]

\subsection*{Equation of the Torse derived from the Correspondence.}

11. For any given point $\theta$ of the quartic, the equation $\Psi(\sigma, \pi) = 0$ gives six values of the pair $(\sigma, \pi)$, that is, six points $\theta_1$ which are corresponding points of $\theta$; and the line joining $\theta$ to any one of these points $\theta_1$ is a generating line of the torse. The equation of the torse may be obtained by eliminating $\theta$ between the equations of the line and the condition of correspondence. The resulting equation in $(x, y, z, w)$ is the required equation of the sextic torse.

12. To obtain this in a more explicit form, observe that for a generating line of the torse, the coordinates $(x, y, z, w)$ satisfy the equations of the line joining two corresponding points $\theta$, $\theta_1$. These equations may be written
\[
\frac{x}{(\theta + \alpha)^2(\theta_1 + \alpha)^2} =
\frac{y}{(\theta + \beta)^2(\theta_1 + \beta)^2} =
\frac{z}{(\theta + \gamma)^2(\theta_1 + \gamma)^2} =
\frac{w}{(\theta + \delta)^2(\theta_1 + \delta)^2}.
\]
Writing these equal ratios $= \rho$, and expressing that the coordinates thus determined satisfy the equation of the osculating plane at the point $\theta$, we obtain an equation which, combined with $\Psi(\sigma, \pi) = 0$, determines the generating lines of the torse which pass through a given point of the quartic.

13. The actual calculation of the equation of the torse by this method is still complicated, but it is more symmetrical than the direct calculation of the discriminant of the sextic function. I do not propose to give the complete expansion, but only to verify that the form of the equation is as stated, and to obtain some further properties of the torse.

\subsection*{Verification of the form $\Lambda = (a, b, c, f, g, h)^{12}(x, y, z, w)^6 = 0$.}

14. We have seen that the discriminant of the sextic function is of the degree 10 in $(x, y, z, w)$ and 30 in $(a, b, c, f, g, h)$, and that it contains the factors $xyzw$ and $(abcfgh)^2$. The remaining factor is of the degree 6 in $(x, y, z, w)$ and 12 in $(a, b, c, f, g, h)$, which agrees with the form $\Lambda = 0$.

15. To verify that the coefficients of the equation are of the degree 12 in $(a, b, c, f, g, h)$, observe that if we change $(\alpha, \beta, \gamma, \delta)$ into $(k\alpha, k\beta, k\gamma, k\delta)$, the quantities $(a, b, c, f, g, h)$ are all multiplied by $k$, while the equation of the osculating plane remains unalter if at the same time we change $\theta$ into $k\theta$. The equation of the torse must therefore be homogeneous of the same degree in $(a, b, c, f, g, h)$ and in the coefficients of the equation of the osculating plane; that is, of the degree 12 in $(a, b, c, f, g, h)$, since each coefficient of the sextic function is of the degree 2 in these quantities. This completes the verification.

\subsection*{Further Properties of the Torse.}

16. The torse is a developable surface of the sixth order; it has for its edge of regression a curve of the sixteenth order, which is the locus of the points of intersection of four consecutive generating lines of the torse. The cuspidal curve of the torse is the unicursal quartic itself; in fact, the osculating plane at any point of the quartic touches the torse along the entire generating line which passes through that point, and the point itself is a cuspidal point on the section of the torse by the osculating plane.

17. The nodal curve of the torse is of a complicated nature. Any point on the nodal curve is the intersection of two non-consecutive generating lines of the torse; the nodal curve is therefore the locus of the points of intersection of pairs of corresponding points on the quartic, for all possible pairs. It may be shown that the nodal curve is of the order 18, and that it breaks up into a number of distinct curves, corresponding to the different kinds of correspondence between points of the quartic.

18. The complete intersection of the torse with a plane is a curve of the sixth order, which, as we have seen, breaks up into six lines (the tangents of the quartic at six points) and a proper sextic curve. The six lines are of course concurrent at the point where the plane meets the quartic; they form a degenerate part of the section, the remaining part being the proper sextic curve, which is the partial intersection of the plane with the torse.

19. It may be remarked that the torse here considered is a particular case of a more general class of surfaces, which may be called the \textit{osculating torses} of a curve. For any curve in space of any order, the locus of the tangent lines which meet the curve in three consecutive points (or, what is the same thing, the envelope of the osculating planes) is a developable surface, which may be called the osculating torse of the curve. For a curve of the order $n$, the osculating torse is in general of the order $3n - 6$; thus for a cubic curve it is of the order 3 (the developable surface generated by the tangents of the cubic), and for a quartic curve it is of the order 6, as in the present case.

20. The theory may be extended to curves in space of any number of dimensions. For a curve in space of $r$ dimensions, the osculating torse is the envelope of the osculating hyperplanes; its order and other properties may be investigated by similar methods.

\newpage
% ============================================================
% PAPER [437] -- DEMONSTRATION NOUVELLE DU THEOREME DE M. CASEY
% Page 115 (PDF page 137)
% ============================================================
\section*{[437]}
\addcontentsline{toc}{section}{[437] Demonstration nouvelle du Theoreme de M. Casey par rapport aux Cercles qui touchent a trois Cercles donnes}

\begin{center}
\textbf{\Large DEMONSTRATION NOUVELLE DU THEOREME DE M. CASEY}\\[0.3em]
\textbf{\Large PAR RAPPORT AUX CERCLES QUI TOUCHENT A TROIS}\\[0.3em]
\textbf{\Large CERCLES DONNES.}
\end{center}

\begin{flushright}
\textit{[From the \textit{Annali di Matematica pura ed applicata},}\\
\textit{torn.\ \textsc{i}.\ (1867), pp.\ 132--134.]}
\end{flushright}

\vspace{1em}

\noindent
This is in fact the investigation contained in the paper 414, ``On Polyzomal Curves otherwise the curves $\sqrt{P} + \sqrt{Q} + \sqrt{R} = 0$,'' Annex II.\ pp.\ 568--573, ``On Casey's theorem for the circle which touches three given circles,'' viz.\ it is based on the identity of the two problems
\begin{enumerate}[label=$\circ$]
\item to find a circle touching three given circles,
\item to find a cone-sphere (sphere of radius zero) passing through three given points in space.
\end{enumerate}

\newpage
% ============================================================
% PAPER [438] -- NOTE SUR QUELQUES TORSES SEXTIQUES
% Pages 116--118 (PDF pages 138--140)
% ============================================================
\section*{[438]}
\addcontentsline{toc}{section}{[438] Note sur quelques Torses Sextiques}

\begin{center}
\textbf{\Large NOTE SUR QUELQUES TORSES SEXTIQUES.}
\end{center}

\begin{flushright}
\textit{[From the \textit{Annali di Matematica pura ed applicata}, tom.\ \textsc{ii}.\ (1868), pp.\ 99, 100.]}
\end{flushright}

\vspace{1em}

\selectlanguage{french}
\noindent
Je d\'{e}sire d'appeler attention aux surfaces d\'{e}veloppables, ou torses, donn\'{e}es par l\'{e}quation
\[
(ae - 4bd + 3c^2)^3 - 27\,(ace - ad^2 - b^2e + 2bcd - c^3)^2 = 0.
\]

Dans cette \'{e}quation $(a, b, c, d, e)$ sont des fonctions lin\'{e}aires quelconques des quatre coordonn\'{e}es $(x, y, z, t)$; ces quantit\'{e}s sont donc li\'{e}es par une \'{e}quation lin\'{e}aire
\[
Aa + 4Bb + 6Cc + 4Dd + Ee = 0,
\]
et je remarque que la classification des torses comprises sous l\'{e}quation mentionn\'{e}e d\'{e}pend des propri\'{e}t\'{e}s invariantives de la fonction $(A, B, C, D, E \between \tau, 1)^4$.

En effet la torse a une courbe cuspidale, ou ar\^{e}te de rebroussement, donn\'{e}e par les \'{e}quations
\[
ae - 4bd + 3c^2 = 0,\quad ace - ad^2 - b^2e + 2bcd - c^3 = 0,
\]
et une courbe nodale donn\'{e}e par les \'{e}quations
\[
\frac{ac - b^2}{a} = \frac{ad - bc}{2b} = \frac{ae + 2bd - 3c^2}{6c} = \frac{be - cd}{2d} = \frac{ce - d^2}{e}:
\]
ces deux courbes se rencontrent dans les points donn\'{e}s par les \'{e}quations
\[
\frac{a}{b} = \frac{b}{c} = \frac{c}{d} = \frac{d}{e},
\]
lesquels sont des points stationnaires de la courbe cuspidale. Pour trouver ces points, en \'{e}crivant $a : b : c : d : e = \tau^4 : \tau^3 : \tau^2 : \tau : 1$, on obtient pour le param\`{e}tre $\tau$ l\'{e}quation
\[
(A, B, C, D, E \between \tau, 1)^4 = 0.
\]

Et l'on voit ainsi qu'il y a un rapport entre la th\'{e}orie de la surface et cette \'{e}quation; \`{a} toute particularit\'{e} invariantive de l'\'{e}quation, il y correspondra quelque particularit\'{e} de la torse.

Les cas \`{a} consid\'{e}rer sont:
\begin{enumerate}[label=$\circ$]
\item Racines in\'{e}gales, sans aucune relation invariantive. C'est le cas g\'{e}n\'{e}ral; je l'ai consid\'{e}r\'{e} dans le M\'{e}moire, ``On a certain Sextic Torse'' [436].
\end{enumerate}
\selectlanguage{english}

\newpage
% ============================================================
% PAPER [439] -- ADDITION A LA NOTE SUR QUELQUES TORSES SEXTIQUES
% Pages 119--120 (PDF pages 140--142)
% ============================================================
\section*{[439]}
\addcontentsline{toc}{section}{[439] Addition a la Note sur quelques Torses Sextiques}

\begin{center}
\textbf{\Large ADDITION A LA NOTE SUR QUELQUES\\TORSES SEXTIQUES.}
\end{center}

\begin{flushright}
\textit{[From the \textit{Annali di Matematica pura ed applicata}, tom.\ \textsc{ii}.\ (1868), pp.\ 219--221.]}
\end{flushright}

\vspace{1em}

\selectlanguage{french}
\noindent
Je viens de trouver ce que signifie la condition $J = 0$.

Consid\'{e}rons deux surfaces quadriques qui se touchent (d'un contact ordinaire). Les \'{e}quations tangentielles peuvent s'\'{e}crire sous la forme
\begin{align*}
a\alpha^2 + b\beta^2 + 2n\gamma\epsilon &= 0, \\
2a\alpha^2 - 2b\beta^2 + c\epsilon^2 + 2n\gamma\delta\epsilon &= 0, \\
a'\alpha^2 + b'\beta^2 + 2n'\gamma\epsilon &= 0, \\
2a'\alpha^2 - 2b'\beta^2 + c'\epsilon^2 + 2n'\gamma\delta\epsilon &= 0,
\end{align*}
et l'on satisfait \`{a} ces \'{e}quations par des valeurs de $\xi$, $\eta$, $\zeta$, $\omega$ qui contiennent un param\`{e}tre arbitraire $\theta$, en \'{e}crivant
\[
\alpha : \beta : \gamma : \delta : \epsilon = 1 : \theta : \theta^2 : \theta^3 : \theta^4,
\]
en effet cela donne
\begin{align*}
a\alpha^2 + b\beta^2 + 2n\gamma\epsilon &= 0, \\
2a\alpha^2 - 2b\beta^2 + c\epsilon^2 + 2n\gamma\delta\epsilon &= 0, \\
a'\alpha^2 + b'\beta^2 + 2n'\gamma\epsilon &= 0, \\
2a'\alpha^2 - 2b'\beta^2 + c'\epsilon^2 + 2n'\gamma\delta\epsilon &= 0,
\end{align*}
ce qui d\'{e}termine les valeurs de $\alpha : \beta : \gamma : \delta : \epsilon$.

L'\'{e}quation du plan tangent commun sera donc
\[
\gamma\epsilon + 4\delta\epsilon + \epsilon^2 = 0,
\]
ou, en multipliant par $12\gamma\epsilon$, cette \'{e}quation sera
\[
(a, b, c, d, e \between \theta, 1)^4 = 0,
\]
avec $\gamma^2 + 4\delta\epsilon + \epsilon^2 = 0$,
les valeurs des coefficients \'{e}tant
\[
a = 12\beta\omega,\quad c = 2(\delta\omega + \epsilon z),\quad e = 12\gamma\omega.
\]

En repr\'{e}sentant par $Aa + 4Bb + 6Cc + 4Dd + Ee = 0$ l'\'{e}quation qui lie les fonctions lin\'{e}aires $a$, $b$, $c$, $d$, $e$, cette \'{e}quation sera $a - e = 0$; on a donc $A = -E = 1$, $B = C = D = 0$; l'invariant $J$ de la fonction $(A, B, C, D, E \between \tau, 1)^4$ est donc $= 0$.

Nous arrivons ainsi \`{a} la conclusion que la torse sextique
\[
(ae - 4bd + 3c^2)^3 - 27(ace - ad^2 - b^2e - c^3 + 2bcd)^2 = 0,
\]
o\`{u} les fonctions lin\'{e}aires $a$, $b$, $c$, $d$, $e$ sont li\'{e}es par une \'{e}quation $Aa + 4Bb + 6Cc + 4Dd + Ee = 0$, telle que l'invariant
\[
J = ACE - AD^2 - EB^2 - C^3 - 2BCD
\]
de la fonction $(A, B, C, D, E \between \tau, 1)^4$ est $= 0$ (cas $7^\circ$ de la Note), est la torse envelopp\'{e}e par le plan tangent commun de deux surfaces quadriques qui se touchent d'un contact ordinaire.

J'ai trouv\'{e} l'\'{e}quation de cette torse dans le M\'{e}moire, ``On the Developable Surfaces which arise from two Surfaces of the Second Order,'' Camb.\ and Dubl.\ Math.\ Jour., t.\ v.\ (1850), pp.\ 46--57, voir p.\ 56, [85]: $a$, $b$, $c$, $n$, $a'$, $b'$, $c'$, $n'$ y d\'{e}notent les m\^{e}mes coefficients comme \`{a} pr\'{e}sent, et en \'{e}crivant
\begin{align*}
bc' - b'c &= f, & an' - a'n &= p, \\
ca' - c'a &= g, & bn' - b'n &= q, \\
ab' - a'b &= h, & cn' - c'n &= r,
\end{align*}
(et de l\`{a} $pf + qg + rh = 0$), l'\'{e}quation trouv\'{e}e est
\[
f^2 g^2 h^2 H^2 \ldots + 4p^2 q^2 r^2 (q\alpha^2 + p\beta^2)^2 = 0.
\]

Je vais v\'{e}rifier ces termes. Partant de l'\'{e}quation $(a, b, c, d, e \between \theta, 1)^4 = 1$, l'\'{e}quation de la torse, en y introduisant pour commodity le facteur $-\tfrac{1}{108}pqr$, sera
\begin{multline*}
-\tfrac{1}{108}pqr \bigl[\{(h\mu + \epsilon z)^2 + 12\gamma\omega^2 - 3(\alpha\nu^2 - \beta\gamma^2)\}^2 \\
- \{2(\delta\omega + \epsilon z)^2 - 9(\delta\omega + \epsilon z)(\alpha\nu^2 - \beta\gamma^2) - 27(h\mu + \epsilon z)\eta^2\omega^2 + 54(\alpha\mu\nu^2 + \beta\gamma^3)\gamma\omega\}\bigr] = 0.
\end{multline*}

En prenant $\gamma\epsilon = ab' - a'b = h$, on obtient pour $\alpha$, $\beta$, $\gamma$, $\delta$, $\epsilon$ les valeurs
\begin{align*}
\alpha' &= -\tfrac{1}{2}A\sqrt{\gamma}, & \beta' &= \frac{\sqrt{\gamma}\cdot\sqrt{p - 9qy}}{2pqr}, \\
\gamma' &= -\tfrac{1}{2}A\sqrt{\gamma}, & \delta' &= \frac{\sqrt{\gamma}\cdot\sqrt{p + 9qy}}{2pqr}, & \epsilon' &= \frac{4\gamma = 1}{2pqr}.
\end{align*}

Les termes en $\nu\zeta$ et $(\alpha\nu^2, \beta\gamma^2)^2$ sont
\[
-\tfrac{1}{4}h M'' \{[4\gamma^2 (\delta^2 + 12\xi)^2 - (2\alpha\nu^2 - 72\xi\beta\gamma)]\nu\zeta - 108(\alpha^2\nu^4 - \beta^2\gamma^4)\gamma^2\}.
\]
ces termes sont donc
\begin{align*}
&= -\tfrac{1}{4} h M'' \{432\gamma^2(\delta^2 - 4\xi\gamma^2\gamma'\omega\nu\zeta - 108(\alpha^2\nu^4 - \beta^2\gamma^4)\gamma^2\}; \\
&= -\tfrac{1}{4} h M'' \{-432g^2 \cdot \tfrac{q^2 r^2}{h^2} \omega\nu^2 + 864(\beta\gamma^4 + \ldots)\},
\end{align*}
comme cela doit \^{e}tre.

Cambridge, le 22 septembre 1868.
\selectlanguage{english}

\newpage
% ============================================================
% PAPER [440] -- NOTE SUR UNE TRANSFORMATION GEOMETRIQUE
% Pages 121--122 (PDF pages 143--144)
% ============================================================
\section*{[440]}
\addcontentsline{toc}{section}{[440] Note sur une Transformation Geometrique}

\begin{center}
\textbf{\Large NOTE SUR UNE TRANSFORMATION GEOMETRIQUE.}
\end{center}

\begin{flushright}
\textit{[From the \textit{Journal f\'{u}r die reine und angewandte Mathematik} (Crelle),}\\
\textit{torn.\ \textsc{lxvii}.\ (1867), pp.\ 95, 96.]}
\end{flushright}

\vspace{1em}

\selectlanguage{french}
\noindent
La lecture de la Note de M.~Hesse, ``Ein Uebertragungsprincip'' (t.\ LXVI.\ p.\ 15 de ce Journal) m'a sugg\'{e}r\'{e} les remarques suivantes.

Soient $(a_1, b_1, c_1, d_1)$, $(a_2, b_2, c_2, d_2)$, $(a_3, b_3, c_3, d_3)$ des constantes donn\'{e}es, on peut supposer que les coordonn\'{e}es $(x, y)$ d'un point quelconque dans un plan soient exprim\'{e}es en fonctions des param\`{e}tres variables $(u, v)$ par les \'{e}quations
\[
x = \frac{a_1 + b_1 u + c_1 v + d_1 uv}{a_3 + b_3 u + c_3 v + d_3 uv},\quad
y = \frac{a_2 + b_2 u + c_2 v + d_2 uv}{a_3 + b_3 u + c_3 v + d_3 uv}.
\]

En introduisant une nouvelle ind\'{e}termin\'{e}e $s$, ces \'{e}quations peuvent \^{e}tre \'{e}crites dans la forme
\begin{align*}
sx &= a_1 + b_1 u + c_1 v + d_1 uv, \\
sy &= a_2 + b_2 u + c_2 v + d_2 uv, \\
s &= a_3 + b_3 u + c_3 v + d_3 uv;
\end{align*}
pour des valeurs donn\'{e}es des coordonn\'{e}es $(x, y)$ la quantit\'{e} $s$ est en g\'{e}n\'{e}ral d\'{e}termin\'{e}e par une \'{e}quation quadratique, et les param\`{e}tres $u$ et $v$ sont des fonctions lin\'{e}aires donn\'{e}es de $s$; il y a cependant deux cas particuliers qu'il convient de distinguer.

$1^\circ$. L'\'{e}quation quadratique en $s$ peut avoir la racine $s = 0$ et, d\'{e}barrass\'{e}e de ce facteur, se r\'{e}duire par cons\'{e}quent \`{a} une \'{e}quation lin\'{e}aire; ce cas particulier a lieu si
\[
(abc)(bcd) = (abd)(acd),
\]
est remplie, o\`{u} la notation $(abc)$ d\'{e}signe le d\'{e}terminant
\[
\begin{vmatrix}
a_1, & b_1, & c_1 \\
a_2, & b_2, & c_2 \\
a_3, & b_3, & c_3
\end{vmatrix}.
\]
Dans ce cas $u$ et $v$ sont des fonctions rationnelles de $(x, y)$ et la transformation a la signification g\'{e}om\'{e}trique suivante:

En consid\'{e}rant deux droites quelconques $L$, $M$ dans l'espace et en menant par le point donn\'{e} $(x, y)$ la droite unique qui rencontre ces deux droites, on peut supposer que $u$ et $v$ soient des param\`{e}tres qui d\'{e}terminent les positions des points de rencontre sur les deux droites respectivement; c.\`{a}.d.\ que $u$ soit la distance d'un point fixe sur la droite $L$ au point de rencontre avec la droite $G$, et de m\^{e}me que $v$ soit la distance d'un point fixe sur la droite $M$ au point de rencontre avec la droite $G$.

$2^\circ$. Supposons $b_1 : c_1 = b_2 : c_2 = b_3 : c_3$, ou ce qui est au fond la m\^{e}me chose $b_1 - c_1 = 0$, $b_2 - c_2 = 0$, $b_3 - c_3 = 0$; alors $s$ est d\'{e}termin\'{e}e par une \'{e}quation simple, mais $u$ et $v$ ne sont plus des fonctions rationnelles de $s$; on voit que dans ce cas $u + v$ et $uv$ sont des fonctions rationnelles de $(x, y)$, et que par cons\'{e}quent $u$ et $v$ sont les racines d'une \'{e}quation quadratique qui contient $(x, y)$ lin\'{e}airement.

On peut supposer que $u$ et $v$ soient les param\`{e}tres de deux points sur une droite donn\'{e}e, c.\`{a}.d.\ que $u$ et $v$ soient les distances de ces deux points respectivement \`{a} un point fixe situ\'{e} sur la droite donn\'{e}e; on a ainsi la transformation de M.~Hesse.

Je n'ai pas cherch\'{e} la signification g\'{e}om\'{e}trique des formules g\'{e}n\'{e}rales.

Cambridge, 10 octobre 1866.
\selectlanguage{english}

\newpage
% ============================================================
% PAPER [441] -- NOTE SUR L'ALGORITHME DES TANGENTES DOUBLES
% Pages 123--125 (PDF pages 145--147)
% ============================================================
\section*{[441]}
\addcontentsline{toc}{section}{[441] Note sur l'Algorithme des Tangentes Doubles d'une Courbe du Quatrieme Ordre}

\begin{center}
\textbf{\Large NOTE SUR L'ALGORITHME DES TANGENTES DOUBLES D'UNE\\COURBE DU QUATRI\`{E}ME ORDRE.}
\end{center}

\begin{flushright}
\textit{[From the \textit{Journal f\'{u}r die reine und angewandte Mathematik} (Crelle),}\\
\textit{torn.\ \textsc{lxviii}.\ (1868), pp.\ 176--179.]}
\end{flushright}

\vspace{1em}

\selectlanguage{french}
\noindent
On n'a pas, je crois, assez fait attention \`{a} l'algorithme (tir\'{e} de la consid\'{e}ration d'une figure dans l'espace) qu'a trouv\'{e} M.~Hesse (dans le m\'{e}moire ``Ueber die Doppeltangenten der Curven vierter Ordnung,'' t.\ XLIX de ce Journal, 1855) pour d\'{e}noter les tangentes doubles (ou bitangentes) d'une courbe du quatri\`{e}me ordre. Voici en quoi cet algorithme consiste.

En employant les huit symboles $1, 2, 3, \ldots 8$, les 28 bitangentes sont repr\'{e}sent\'{e}es par les combinaisons binaires $12, 13, 14, \ldots 78$. Cela pos\'{e}, consid\'{e}rons une expression quelconque $12.13.14$, ou $12.34$,\ldots\ ou disons un ``terme'' qui repr\'{e}sente un syst\`{e}me d'une seule ou de plusieurs des bitangentes. On peut op\'{e}rer sur ce terme avec deux esp\`{e}ces de substitutions; la substitution ordinaire qui consiste \`{a} changer l'arrangement $12345678$ des huit symboles en un autre arrangement quelconque; et la substitution ``bifide'' repr\'{e}sent\'{e}e par un symbole tel que $1234.5678$, lequel d\'{e}note qu'il faut enterechanger les combinaisons $12$ et $34$, $13$ et $24$, $14$ et $23$, $56$ et $78$, $57$ et $68$, $58$ et $67$, en ne changeant pas les autres combinaisons. Par exemple en op\'{e}rant avec $1234.5678$ sur $34.45.56.17$ on obtient $12.45.78.17$.

Le nombre de ces substitutions bifides est 35, ou en comptant la substitution, unit\'{e}, qui ne change aucune des combinaisons, ce nombre est 36. Appelons ``homotypiques'' deux termes qui se d\'{e}rivent l'un de l'autre par une substitution ordinaire; ``syntypiques'' qui se d\'{e}rivent l'un de l'autre par une substitution ordinaire ou bifide; ``sous-groupe'' le syst\`{e}me entier des termes homotypiques \`{a} un terme donn\'{e}: ``groupe'' le syst\`{e}me entier des termes syntypiques \`{a} un terme donn\'{e}. Un groupe peut contenir un seul sous-groupe, ou plusieurs sous-groupes; mais il importe de remarquer que la notion du sous-groupe n'a pas de signification g\'{e}om\'{e}trique, et ne sert que comme moyen de former les termes du groupe.

Cela \'{e}tant, le th\'{e}or\`{e}me g\'{e}om\'{e}trique est celui-ci; ``les syst\`{e}mes de bitangentes repr\'{e}sent\'{e}es par des termes syntypiques (ou autrement dit, par des termes qui appartiennent au m\^{e}me groupe) ont les m\^{e}mes propri\'{e}t\'{e}s g\'{e}om\'{e}triques.''

Par exemple, en consid\'{e}rant les bitangentes deux \`{a} deux, on a deux sous-groupes, l'un compos\'{e} de termes homotypiques \`{a} $12.13$; l'autre, de termes homotypiques \`{a} $12.34$---ou disons, le sous-groupe $12.13$ de 168 termes et le sous-groupe $12.34$ de 210 termes; mais ces deux sous-groupes ne forment qu'un seul groupe: pour montrer cela il suffit d'op\'{e}rer sur $12.13$, par exemple avec la substitution $1245.3678$, ce qui donne $45.13$, terme homotypique \`{a} $12.34$. Cela veut dire qu'il n'y a pas de combinaison de deux bitangentes qui se distingue d'une mani\`{e}re quelconque de toute autre combinaison de deux bitangentes.

Mais en combinant les bitangentes trois \`{a} trois, on a les deux sous-groupes $12.34.56$ (420 termes) et $12.23.34$ (840 termes) qui forment un groupe de 1260 termes; les trois bitangentes repr\'{e}sent\'{e}es par un quelconque des 1260 termes ont leurs six points de contact sur une m\^{e}me conique. Les trois autres sous-groupes $12.23.31$ (56 termes), $12.23.45$ (1680 termes) et $12.13.14$ (280 termes) forment un groupe de 2016 termes, et pour trois bitangentes repr\'{e}sent\'{e}es par un terme quelconque de ce groupe, les six points de contact ne sont pas situ\'{e}s sur une m\^{e}me conique.

Comme un autre exemple j'explique la constitution des 63 ``groupes'' de Steiner (voir le m\'{e}moire de Steiner, ``Eigenschaften der Curven vierten Grades n\"{u}cksichtlich ihrer Doppeltangenten,'' t.\ XLIX de ce journal, 1855) ou (pour \'{e}viter l'emploi de ce mot groupe dans une nouvelle signification) disons les 63 termes $G$ de Steiner, chaque terme compos\'{e} de 6 paires de bitangentes. On a ici un sous-groupe de 35 termes $G$, de la forme
\[
12.34;\ 13.42;\ 14.23;\ 56.78;\ 57.86;\ 58.67
\]
(pour abr\'{e}ger on peut d\'{e}noter ce terme par $1234.5678$), et un sous-groupe de 28 termes $G$, de la forme
\[
13.32;\ 14.42;\ 15.52;\ 16.62;\ 17.72;\ 18.82
\]
(pour abr\'{e}ger on peut de m\^{e}me d\'{e}noter ce terme par $12.345678$), les deux sous-groupes forment le groupe des 63 termes $G$.

Steiner a de plus consid\'{e}r\'{e} les ``syst\`{e}mes'' ou disons les termes $S_1$, $S_2$, compos\'{e}s chacun de trois termes $G$; savoir 315 termes $S_1$ et 336 termes $S_2$. Les 315 termes $S_1$ sont ici un groupe compos\'{e} d'un sous-groupe de 105 termes $3G$, de la forme
\[
1234.5678;\ 1256.3478;\ 1278.3456
\]
et un sous-groupe de 210 termes $2G_1 + G_2$ de la forme
\[
12.345678;\ 34.125678 \quad\text{et}\quad 1234.5678.
\]
Et de m\^{e}me les 336 termes $S_2$ sont un groupe compos\'{e} d'un sous-groupe de 280 termes $2G_1 + G_2$ de la forme
\[
1234.5678;\ 5234.1678 \quad\text{et}\quad 15.234678
\]
et un sous-groupe de 56 termes $3G_2$ de la forme
\[
12.34.3678;\ 13.24.5678;\ 31.245678.
\]

Il va sans dire que je me suis servi de l'abreviation $1234.5678$ pour d\'{e}noter le terme $12.34$; $13.42$; $14.23$; $56.78$; $57.86$; $58.67$; et de m\^{e}me pour les autres termes $G_1$ ou $G_2$.

M.~Aronhold (dans le m\'{e}moire ``Ueber den gegenseitigen Zusammenhang der 28 Doppeltangenten einer allgemeinen Curve vierten Grades,'' Berl.\ Monatsber.\ Juli 1864), partant de 7 bitangentes donn\'{e}es, a trouv\'{e} une construction pour les autres 21 bitangentes. Les bitangentes donn\'{e}es doivent \^{e}tre ind\'{e}pendantes; savoir pour trois quelconques de ces 7 bitangentes, les six points de contact ne sont pas situ\'{e}s sur une m\^{e}me conique. Les bitangentes repr\'{e}sent\'{e}es par les termes $12, 13, 14, 15, 16, 17, 18$ sont un tel syst\`{e}me de bitangentes ind\'{e}pendantes; et en d\'{e}notant de cette mani\`{e}re les 7 bitangentes donn\'{e}es, la bitangente construite par le moyen de la conique qui touche cinq de ces droites, par exemple les droites $38, 48, 58, 68, 78$, (ou conique $34567$) peut \^{e}tre d\'{e}not\'{e}e par $12$, et de m\^{e}me pour les autres bitangentes cherch\'{e}es; on a ainsi le syst\`{e}me entier des bitangentes d\'{e}not\'{e}es comme auparavant par $12, 13, 14, \ldots 78$.

J'ajoute que le groupe qui contient $18, 28, 38, 48, 58, 68, 78$ est compos\'{e} d'un sous-groupe $18, 28, 38, 48, 58, 68, 78$ de 8 termes, et d'un sous-groupe $12, 23, 31, 48, 58, 68, 78$ de 280 termes; le groupe contient donc 288 termes; savoir il y a ce nombre 288 de syst\`{e}mes de sept bitangentes ind\'{e}pendantes qui peuvent chacun servir \`{a} trouver par la construction d'Aronhold les autres 21 bitangentes.

\textit{P.S.} J'ai trouv\'{e} \`{a} propos de la m\'{e}thode de M.~Aronhold une forme commode pour l'\'{e}quation de la conique qui touche cinq droites donn\'{e}es; en supposant que l'on ait identiquement $x + y + z + w = 0$, et que les droites donn\'{e}es soient $x = 0$, $y = 0$, $z = 0$, $w = 0$, et $ax + by + cz + dw = 0$, l'\'{e}quation de la conique est
\begin{multline*}
(a - d)^2 (b - c)^2 (xw + yz) + (b - d)^2 (c - a)^2 (yw + zx) \\
+ (c - d)^2 (a - b)^2 (zw + xy) = 0.
\end{multline*}
J'ajoute qu'en \'{e}crivant pour abr\'{e}ger
\[
\alpha : \beta : \gamma = (a - d)(b - c) : (b - d)(c - a) : (c - d)(a - b)
\]
(d'o\`{u} $\alpha + \beta + \gamma = 0$) les coordonn\'{e}es $(x, y, z, w)$ des points de contact avec les droites $x = 0$, $y = 0$, $z = 0$, $w = 0$ sont $(0, \gamma, \beta, \alpha)$, $(\gamma, 0, \alpha, \beta)$, $(\beta, \alpha, 0, \gamma)$, $(\alpha, \beta, \gamma, 0)$ respectivement; et que les coordonn\'{e}es du point de contact avec la droite $ax + by + cz + dw = 0$ sont
\[
x : y : z : w = (bcd) : -(cda) : (dab) : -(abc)
\]
o\`{u}, pour abr\'{e}ger, $(bcd)$ d\'{e}note $(b - c)(c - d)(d - b)$, et de m\^{e}me pour $(cda)$, $(dab)$, $(abc)$.

Cambridge, le 23 septembre 1867.
\selectlanguage{english}

\newpage
% ============================================================
% PAPER [442] -- NOTE SUR LA SURFACE DU QUATRIEME ORDRE
% Pages 126--128 (PDF pages 148--150)
% ============================================================
\section*{[442]}
\addcontentsline{toc}{section}{[442] Note sur la Surface du Quatrieme Ordre douee de seize Points Singuliers et de seize Plans Singuliers}

\begin{center}
\textbf{\Large NOTE SUR LA SURFACE DU QUATRI\`{E}ME ORDRE DOU\'{E}E DE\\SEIZE POINTS SINGULIERS ET DE SEIZE PLANS SINGULIERS.}
\end{center}

\begin{flushright}
\textit{[From the \textit{Journal f\'{u}r die reine und angewandte Mathematik} (Crelle),}\\
\textit{tom.\ \textsc{lxxii}.\ (1871), pp.\ 292--293.]}
\end{flushright}

\vspace{1em}

\selectlanguage{french}
\noindent
L'\'{e}quation de M.~Kummer se transforme sans difficult\'{e} en celle-ci
\[
\sqrt{\alpha x \left(\gamma'\gamma'' y - \beta'\beta'' z - \frac{w}{\alpha}\right)}
+ \sqrt{\beta y \left(\alpha'\alpha'' z - \gamma'\gamma'' x - \frac{w}{\beta}\right)}
+ \sqrt{\gamma z \left(\beta'\beta'' x - \alpha'\alpha'' y - \frac{w}{\gamma}\right)} = 0,
\]
o\`{u}
\[
\alpha + \beta + \gamma = 0,\quad \alpha' + \beta' + \gamma' = 0,\quad \alpha'' + \beta'' + \gamma'' = 0.
\]

Or cette \'{e}quation rendue rationnelle prend, apr\`{e}s toutes les r\'{e}ductions n\'{e}cessaires, la forme suivante:
\begin{multline*}
w^2(x^2 + y^2 + z^2 - 2yz - 2zx - 2xy) \\
+ 2w(\alpha\alpha'\alpha''(y^2z - yz^2) + \beta\beta'\beta''(z^2x - zx^2) + \gamma\gamma'\gamma''(x^2y - xy^2) + \theta xyz) \\
+ (\alpha\alpha'\alpha'' yz + \beta\beta'\beta'' zx + \gamma\gamma'\gamma'' xy)^2 = 0,
\end{multline*}
o\`{u}, pour abr\'{e}ger, l'on a \'{e}crit
\begin{align*}
\theta &= (\beta - \gamma)\alpha'\alpha'' + (\gamma - \alpha)\beta'\beta'' + (\alpha - \beta)\gamma'\gamma'', \\
&= (\beta' - \gamma')\alpha''\alpha + (\gamma' - \alpha')\beta''\beta + (\alpha' - \beta')\gamma''\gamma, \\
&= (\beta'' - \gamma'')\alpha\alpha' + (\gamma'' - \alpha'')\beta\beta' + (\alpha'' - \beta'')\gamma\gamma', \\
&= -\tfrac{1}{3}\{(\beta - \gamma)(\beta' - \gamma')(\beta'' - \gamma'') + (\gamma - \alpha)(\gamma' - \alpha')(\gamma'' - \alpha'') + (\alpha - \beta)(\alpha' - \beta')(\alpha'' - \beta'')\},
\end{align*}
l'identit\'{e} de ces diff\'{e}rentes valeurs de $\theta$ \'{e}tant facile \`{a} v\'{e}rifier.

En repr\'{e}sentant par $Aw^2 + 2Bw + C = 0$ la forme rationnelle de l'\'{e}quation de la surface, on trouve pour le discriminant $AC - B^2$ de cette \'{e}quation du second degr\'{e} en $w$ la valeur
\[
AC - B^2 = Wa^2\alpha\beta\gamma - 27(\ldots).
\]

L'\'{e}quation de la surface rendue rationnelle est sym\'{e}trique par rapport aux trois syst\`{e}mes de quantit\'{e}s $(\alpha, \beta, \gamma)$, $(\alpha', \beta', \gamma')$, $(\alpha'', \beta'', \gamma'')$; la forme rationnelle de la m\^{e}me \'{e}quation peut donc \^{e}tre pr\'{e}sent\'{e}e de trois mani\`{e}res diff\'{e}rentes, savoir:
\begin{align*}
\sqrt{\alpha x (\gamma'\gamma'' y - \beta'\beta'' z - \tfrac{w}{\alpha})} + \sqrt{\beta y (\alpha'\alpha'' z - \gamma'\gamma'' x - \tfrac{w}{\beta})} + \sqrt{\gamma z (\beta'\beta'' x - \alpha'\alpha'' y - \tfrac{w}{\gamma})} &= 0, \\
\sqrt{\alpha' x (\gamma\gamma'' y - \beta\beta'' z - \tfrac{w}{\alpha'})} + \sqrt{\beta' y (\alpha\alpha'' z - \gamma\gamma'' x - \tfrac{w}{\beta'})} + \sqrt{\gamma' z (\beta\beta'' x - \alpha\alpha'' y - \tfrac{w}{\gamma'})} &= 0, \\
\sqrt{\alpha'' x (\gamma\gamma' y - \beta\beta' z - \tfrac{w}{\alpha''})} + \sqrt{\beta'' y (\alpha\alpha' z - \gamma\gamma' x - \tfrac{w}{\beta''})} + \sqrt{\gamma'' z (\beta\beta' x - \alpha\alpha' y - \tfrac{w}{\gamma''})} &= 0,
\end{align*}
et l'on voit de plus que les \'{e}quations des seize plans singuliers sont
\begin{align*}
x = 0,\quad y = 0,\quad z = 0,\quad w &= 0, \\
\frac{x}{\alpha} + \frac{y}{\beta} + \frac{z}{\gamma} = 0,\quad
\frac{x}{\alpha} + \frac{y}{\beta} + \frac{z}{\gamma} = 0,\quad
\frac{x}{\alpha} + \frac{y}{\beta} + \frac{z}{\gamma} &= 0, \\
\frac{\gamma'\gamma'' y - \beta'\beta'' z}{\alpha} - \frac{w}{\alpha^2} = 0,\quad
\frac{\alpha'\alpha'' z - \gamma'\gamma'' x}{\beta} - \frac{w}{\beta^2} = 0,\quad
\frac{\beta'\beta'' x - \alpha'\alpha'' y}{\gamma} - \frac{w}{\gamma^2} &= 0, \\
\frac{\gamma''\gamma y - \beta''\beta z}{\alpha'} - \frac{w}{\alpha'^2} = 0,\quad
\frac{\alpha''\alpha z - \gamma''\gamma x}{\beta'} - \frac{w}{\beta'^2} = 0,\quad
\frac{\beta''\beta x - \alpha''\alpha y}{\gamma'} - \frac{w}{\gamma'^2} &= 0, \\
\frac{\gamma\gamma' y - \beta\beta' z}{\alpha''} - \frac{w}{\alpha''^2} = 0,\quad
\frac{\alpha\alpha' z - \gamma\gamma' x}{\beta''} - \frac{w}{\beta''^2} = 0,\quad
\frac{\beta\beta' x - \alpha\alpha' y}{\gamma''} - \frac{w}{\gamma''^2} &= 0,
\end{align*}
les quantit\'{e}s $\alpha$, $\beta$, $\gamma$, etc. \'{e}tant li\'{e}es entre elles par les trois \'{e}quations
\[
\alpha + \beta + \gamma = 0,\quad \alpha' + \beta' + \gamma' = 0,\quad \alpha'' + \beta'' + \gamma'' = 0.
\]
Voil\`{a} ce me semble la forme la plus simple pour l'\'{e}quation de cette surface.

Cambridge, le 23 f\'{e}vrier 1871.
\selectlanguage{english}

\newpage
% ============================================================
% PAPER [443] -- ON THE SOLUTION OF A QUARTIC EQUATION
% Pages 129--130 (PDF pages 150--151)
% ============================================================
\section*{[443]}
\addcontentsline{toc}{section}{[443] On the Solution of a Quartic Equation}

\begin{center}
\textbf{\Large ON THE SOLUTION OF A QUARTIC EQUATION.}
\end{center}

\begin{flushright}
\textit{[From the \textit{Mathematische Annalen}, vol.\ \textsc{i}.\ (1869), pp.\ 54, 55.]}
\end{flushright}

\vspace{1em}

\noindent
If $U$ denote the quartic function $(a, b, c, d, e\between x, y)^4$, $H$ its Hessian $= (ac - b^2, \ldots \between x, y)^2$, and $\Phi$ the sexticovariant $= (a^2d - 3abc + 2b^3, \ldots \between x, y)^6$, then the equation $U = 0$ can be solved by means of the equation $\Phi = 0$. In fact, writing
\[
X = \frac{x}{(ac - b^2)x + (ad - bc)y},\quad
Y = \frac{y}{(ac - b^2)x + (ad - bc)y},
\]
and similarly
\[
X' = \frac{x}{(ad - bc)x + (ae - c^2)y},\quad
Y' = \frac{y}{(ad - bc)x + (ae - c^2)y},
\]
we have identically
\[
X^3 + Y^3 + Z^3 = 0,
\]
so that the expression $\alpha X + \beta Y + \gamma Z$ will be a square if only $\alpha^3 + \beta^3 + \gamma^3 = 0$. (To see this observe that in virtue of the equation $X^3 + Y^3 + Z^3 = 0$, we have $X + iY$, $X - iY$ each of them a square, and thence $\alpha X + \beta Y + \gamma Z = \Box$.) Hence, assuming $\alpha^3 + \beta^3 + \gamma^3 = 0$, we have an equation of the form
\[
(\alpha X + \beta Y + \gamma Z)^2 = \lambda U + \mu H,
\]
and the equation $U = 0$ will be satisfied if $\alpha X + \beta Y + \gamma Z = 0$ and $H = 0$; that is, the values of $(x, y)$ which satisfy $U = 0$ are given by the intersection of the quartic curve $U = 0$ with the conic $H = 0$, and by the condition $\alpha X + \beta Y + \gamma Z = 0$.

The conic $H = 0$ meets the quartic $U = 0$ in eight points, which are the four double points of the quartic (counted twice) and four other points. The condition $\alpha X + \beta Y + \gamma Z = 0$ determines which of these points satisfy the original equation.

\newpage
% ============================================================
% PAPER [444] -- ON THE CENTRO-SURFACE OF AN ELLIPSOID
% Pages 131--133 (PDF pages 152--154)
% ============================================================
\section*{[444]}
\addcontentsline{toc}{section}{[444] On the Centro-surface of an Ellipsoid}

\begin{center}
\textbf{\Large ON THE CENTRO-SURFACE OF AN ELLIPSOID.}
\end{center}

\begin{flushright}
\textit{[From the \textit{Proceedings of the London Mathematical Society},}\\
\textit{vol.\ \textsc{iii}.\ (1869--1871), pp.\ 16--18.]}
\end{flushright}

\vspace{1em}

\noindent
The President [Prof.~Cayley] gave an account of his investigations on the centro-surface of an ellipsoid (locus of the centres of curvature of the ellipsoid). The surface has been studied by Dr~Salmon, and also by Prof.~Clebsch, but in particular the theory was developed in Salmon's ``Solid Geometry'' (Ed.\ 2, 1865).

The coordinates of a point on the ellipsoid being taken to be
\[
x = a\sin\theta\cos\phi,\quad y = b\sin\theta\sin\phi,\quad z = c\cos\theta,
\]
the coordinates of the corresponding point on the centro-surface are given by the formulae
\begin{align*}
-a^2\alpha^2 x^2 &= (a^2 + h)^2 (a^2 + k), \\
-b^2\beta^2 y^2 &= (b^2 + h)^2 (b^2 + k), \\
-c^2\gamma^2 z^2 &= (c^2 + h)^2 (c^2 + k);
\end{align*}
viz., these equations, considering therein $(h, k)$ as arbitrary parameters, determine the coordinates $(x, y, z)$ of a point on the centro-surface.

The principal sections (as is known) consist each of them of an ellipse counting three times, and of a quartic curve called the \textit{ellipse of curvature}.

Writing for shortness
\[
P = a^2 + b^2 + c^2,\quad Q = a^2b^2 + b^2c^2 + c^2a^2,\quad R = a^2b^2c^2,
\]
the relation between $h$ and $k$ is
\[
h^2 - Ph + Q = 0,\quad k^2 - Pk + Q = 0,
\]
so that $h$ and $k$ are the roots of the quadratic $t^2 - Pt + Q = 0$.

The equation of the centro-surface is obtained by eliminating $(\theta, \phi, h, k)$ between the above equations. The result is an equation of the twelfth degree in $(x, y, z)$, which may be written in the form
\[
\left(\frac{x^2}{a^2}\right)^{\!\frac{1}{3}} + \left(\frac{y^2}{b^2}\right)^{\!\frac{1}{3}} + \left(\frac{z^2}{c^2}\right)^{\!\frac{1}{3}} = 1,
\]
or, rationally,
\begin{multline*}
(QR + 3Qh^2 + Ph^4) \\
+ k(3Q + 4Ph^2 + 3h^4) + k^2(-P + 3h^2) = 0;
\end{multline*}
this is a $(2, 2)$ correspondence between the two parameters $h$, $h_1$; the united values $h_1 = h$, are given by the equation
\[
h^4 - 2Ph^3 + (P^2 + 2Q)h^2 - 2PQh + (Q^2 - PR) = 0.
\]

The centro-surface is a surface of the twelfth order; it has a nodal curve of the twenty-fourth order, which is the locus of the centres of curvature of the principal sections of the ellipsoid. The cuspidal curve is of the twelfth order, and is the envelope of the normals to the ellipsoid.

\newpage
% ============================================================
% PAPER [445] -- A MEMOIR ON QUARTIC SURFACES
% Pages 134--150 (PDF pages 155--171)
% ============================================================
\section*{[445]}
\addcontentsline{toc}{section}{[445] A Memoir on Quartic Surfaces}

\begin{center}
\textbf{\Large A MEMOIR ON QUARTIC SURFACES.}
\end{center}

\begin{flushright}
\textit{[From the \textit{Proceedings of the London Mathematical Society},}\\
\textit{vol.\ \textsc{iii}.\ (1869--1871), pp.\ 19--69.}\\
\textit{Read February 10, 1870.]}
\end{flushright}

\vspace{1em}

\subsection*{Introduction.}

The present Memoir is intended as a commencement of the theory of the quartic surfaces which have nodes (conical points). A quartic surface may be without nodes, or it may have any number of nodes up to 16. I show that this is so, and I consider in particular the surfaces with 16, 15, 14, 13, 12, 11, 10 and 9 nodes; and I obtain the different forms of the equation of the surface with 16, 15, 14 and 13 nodes. The investigation is carried on by means of the theory of the symmetroid, or surface represented by a symmetrical determinant; and in particular by means of the dianome, enneadianome, and decadianome of given points.

The present researches on Quartic Surfaces were suggested to me by Professor Kummer's most interesting Memoir ``Ueber die algebraischen Strahlensysteme u.s.w.,'' Berl.\ Abh.\ 1866, in which, without entering upon the general theory, he is led to consider the quartic surfaces, or certain quartic surfaces, with 16, 15, 14, 13, 12, or 11 nodes; the last of these, or surface with 11 nodes, being in fact a particular case of the symmetroid.

\subsection*{Considerations in regard to the Jacobian of four, or more or less than four, Surfaces.}

1. In the case of any four surfaces, $P = 0$, $Q = 0$, $R = 0$, $S = 0$, the differential coefficients of $P$, $Q$, $R$, $S$ in regard to the coordinates $(x, y, z, w)$ may be arranged as a square matrix in either of the ways
\[
\begin{array}{cccc}
P, & Q, & R, & S \\
\partial_x, & \partial_y, & \partial_z, & \partial_w
\end{array}
\quad\text{or}\quad
\begin{array}{cccc}
\partial_x, & \partial_y, & \partial_z, & \partial_w \\
P, & Q, & R, & S
\end{array}
\]
and with either arrangement we may form one and the same determinant, the Jacobian determinant $J(P, Q, R, S)$, or, equating it to zero, the Jacobian surface $J(P, Q, R, S) = 0$ of the four surfaces.

2. In the case of more than four surfaces, adopting the arrangement
\[
\begin{array}{ccccc}
P, & Q, & R, & S, & T, \ldots \\
\partial_x, & \partial_y, & \partial_z, & \partial_w
\end{array}
\]
and considering the several determinants which can be formed with any four columns of the matrix, these equated to zero establish a manifold system of surfaces; and in like manner for less than four surfaces, we have a system of curves or points.

3. In particular, if the surfaces are all of the first order (planes), then the Jacobian of four planes is a point, that is, the point of intersection of the four planes. And generally, the Jacobian of any four surfaces of the orders $a$, $b$, $c$, $d$ respectively, is a surface of the order $a + b + c + d$.

4. If the surfaces have a common point, the Jacobian has the same point; if the surfaces have a common curve, the Jacobian has the same curve; and so on for any common locus of the surfaces.

5. Suppose that the orders of the surfaces $P = 0$, $Q = 0$, \ldots\ are $a+1$, $b+1$, \ldots\ (so that the orders of the differential coefficients of $P$, $Q$, \ldots\ are $a$, $b$, \ldots), then we have for the orders of the several loci,
\begin{align*}
J(P, Q) &= 0, \text{ point-system, order } a^2 + ab + b^2; \\
J(P, Q, R) &= 0, \text{ curve, }\qquad\quad a^2 + b^2 + c^2 + bc + ca + ab; \\
J(P, Q, R, S) &= 0, \text{ surface, }\qquad a + b + c + d; \\
J(P, Q, R, S, T) &= 0, \text{ curve, }\qquad\quad ab + ac + \ldots + de; \\
J(P, Q, R, S, T, U) &= 0, \text{ point-system, }\quad abc + abd + \ldots + def;
\end{align*}
see, as to this, Salmon's \textit{Solid Geometry}, Ed.~2, (1865), Appendix~IV., ``On the Order of Systems of Equations'' [not reproduced in the later editions]. In particular, if $a = b = c = \ldots = 1$, then the orders are 4, 6, 4, 10, 20.

\subsection*{As to the Surface obtained by equating to zero a Symmetrical Determinant.}

6. It is also shown (Salmon, Ed.~2, p.~495) that the surface obtained by equating to zero any symmetrical determinant has a determinate number of nodes; viz., if the orders of the terms in the diagonal be $a$, $b$, $c$, \&c., then the number of nodes is $= \tfrac{1}{2}(\Sigma a \cdot \Sigma ab - \Sigma abc)$, or, as this may also be written, $\tfrac{1}{2}(\Sigma a^2b + 2\Sigma abc)$. In particular, the formula applies to the case of the surface
\[
\begin{vmatrix}
A, & H, & G, & L \\
H, & B, & F, & M \\
G, & F, & C, & N \\
L, & M, & N, & D
\end{vmatrix} = 0,
\]
$(a, b, c, d)$ being here the orders of $A$, $B$, $C$, $D$ respectively, and the orders of $F$, $G$, \&c., being $\tfrac{1}{2}(b+c)$, $\tfrac{1}{2}(a+c)$, \&c. If the terms are all of them linear functions of the coordinates, or $a = b = c = d = 1$, then the number of nodes is $= 10$.

7. That the surface has nodes is, in fact, clear from the consideration that any point for which the minors of the determinant all vanish will be a node; and that (for the symmetrical determinant), by making the minors all of them vanish, we establish only a threefold relation between the coordinates. The expression for the number of the nodes is, I think, obtained most readily as follows:

The nodes will be points of intersection of the curve and surface
\[
\begin{vmatrix}
A, & H, & G, & L \\
H, & B, & F, & M \\
G, & F, & C, & N
\end{vmatrix} = 0,
\qquad
\begin{vmatrix}
B, & F, & M \\
F, & C, & N \\
M, & N, & D
\end{vmatrix} = 0,
\]
these, however, contain in common the points
\[
\begin{vmatrix}
H, & B, & F, & M \\
G, & F, & C, & N
\end{vmatrix} = 0;
\]
and the number of nodes is thus obtained by subtracting from the product of the orders of the two loci, the number of these common points, properly counted.

\subsection*{The Symmetroid and the Jacobian of four Quadric Surfaces.}

8. The surface obtained by equating to zero a symmetrical determinant of the fourth order, the terms whereof are linear functions of the coordinates, is a quartic surface with 10 nodes; this surface I call a \textit{symmetroid}. The 10 nodes form a system which I call the \textit{dianome} of the surface; and the symmetroid is intimately connected with the Jacobian of four quadric surfaces.

9. In fact, taking four quadric surfaces $P = 0$, $Q = 0$, $R = 0$, $S = 0$, the Jacobian $J(P, Q, R, S) = 0$ is a surface of the fourth order. This Jacobian surface has the property that it is the locus of the vertices of the quadric cones which pass through the intersection of the four quadric surfaces; or, what is the same thing, it is the locus of the points the polar planes of which in regard to the four quadric surfaces meet in a point.

10. The equation of the Jacobian may be written in the form of a symmetrical determinant. In fact, if we write the equations of the four quadric surfaces in the form
\begin{align*}
P &= (A, B, C, D, F, G, H, L, M, N \between x, y, z, w)^2 = 0, \\
Q &= (A', B', C', D', F', G', H', L', M', N' \between x, y, z, w)^2 = 0, \text{ \&c.,}
\end{align*}
then the Jacobian is
\[
\begin{vmatrix}
\partial_x P, & \partial_y P, & \partial_z P, & \partial_w P \\
\partial_x Q, & \partial_y Q, & \partial_z Q, & \partial_w Q \\
\partial_x R, & \partial_y R, & \partial_z R, & \partial_w R \\
\partial_x S, & \partial_y S, & \partial_z S, & \partial_w S
\end{vmatrix} = 0,
\]
which is a quartic surface. But by a linear transformation of the four quadric surfaces, we may express the Jacobian as a symmetrical determinant; and conversely, any symmetroid may be regarded as the Jacobian of four quadric surfaces.

11. The symmetroid has 10 nodes; these are the points for which the minors of the symmetrical determinant all vanish. The nodes of the symmetroid correspond to the lines on the Jacobian surface of the four quadric surfaces; or, more precisely, to the singular lines of the developable surface which is the envelope of the polar planes of the points of the Jacobian.

\subsection*{Nodal Quartic Surfaces in general.}

12. A quartic surface may have any number of nodes from 0 up to 16; the maximum number 16 is attained in the case of Kummer's surface, which is a very remarkable surface having a large number of special properties. The following table shows the nature of the surfaces with different numbers of nodes:
\begin{center}
\begin{tabular}{cl}
\toprule
Number of nodes & Description \\
\midrule
0 & General quartic surface \\
1 & Quartic with one node \\
\ldots & \ldots \\
10 & Symmetroid \\
11 & Special symmetroid \\
12, 13, 14, 15 & Successively more special surfaces \\
16 & Kummer's surface \\
\bottomrule
\end{tabular}
\end{center}

13. For a quartic surface with $\delta$ nodes, the reduction in the class is $2\delta$; that is, the class of the surface is $36 - 2\delta$. The surface has also a certain number of lines upon it, and the number of these lines diminishes as the number of nodes increases.

14. The nodes of a nodal quartic surface may be arbitrary points, subject to certain conditions. For a surface with $\delta$ nodes, the number of conditions which the nodes must satisfy is equal to the number of independent coefficients in the equation of the surface, diminished by the number of arbitrary parameters which determine the position of the nodes. In particular, for a symmetroid with 10 nodes, the 10 nodes are subjected to just so many conditions as suffice to determine them uniquely when 7 of them are given.

15. The dianodal surface of 7 given points is the locus of the 8th point such that the 8 points are nodes of a quartic surface; the enneadianome is the locus of the 9th point; and the decadianome is the locus of the 10th point. These surfaces are of considerable interest in the theory of nodal quartics.

\subsection*{Further Theory of the Symmetroid.}

16. The symmetroid may be defined as the surface represented by an equation $\Delta = 0$, where $\Delta$ is a symmetrical determinant of the 4th order, the several terms whereof are linear functions of the coordinates $(x, y, z, w)$. Writing this equation in the form
\[
\Delta = \begin{vmatrix}
A, & H, & G, & L \\
H, & B, & F, & M \\
G, & F, & C, & N \\
L, & M, & N, & D
\end{vmatrix} = 0,
\]
where $A$, $B$, $C$, $D$, $F$, $G$, $H$, $L$, $M$, $N$ are linear functions of $(x, y, z, w)$, the surface is a quartic with 10 nodes. The nodes are the points for which all the first minors of $\Delta$ vanish.

17. The Jacobian of four quadric surfaces is identical with the symmetroid, properly defined. If the four quadric surfaces are $P = 0$, $Q = 0$, $R = 0$, $S = 0$, then the equation $J(P, Q, R, S) = 0$ may be thrown into the form $\Delta = 0$, where $\Delta$ is a symmetrical determinant as above; and conversely, any symmetroid may be represented as the Jacobian of four properly chosen quadric surfaces.

18. The theory of the symmetroid is thus identical with the theory of the Jacobian of four quadric surfaces; and the properties of the one may be deduced from those of the other. In particular, the 10 nodes of the symmetroid correspond to the 10 lines on the Jacobian surface which are the edges of the tetrahedron formed by the four quadric surfaces; or, more precisely, to the 10 points of intersection of pairs of opposite edges of this tetrahedron.

\subsection*{Circumscribed Sextic Cone. Nodes of Surface.}

19. Taking the coordinates to be $(x, y, z, w)$, and considering a point $(x_0, y_0, z_0, w_0)$ on the symmetroid, the tangent cone at this point is a quadric cone; and the enveloping cone from the point to the surface is a sextic cone. The sextic cone has the property that it is circumscribed about the tangent cone; that is, the tangent cone touches the sextic cone along a certain curve.

20. The nodes of the symmetroid may be classified according to their nature. There are nodes of the first kind, which are simple nodes; nodes of the second kind, which are biplanar nodes; and nodes of the third kind, which are uniplanar nodes. The different kinds of nodes give rise to different singularities on the sextic cone and on the circumscribed surfaces.

21. For a symmetroid with 10 given nodes, the remaining singularities are determined by the condition that the surface shall have the given nodes. The number of conditions thus imposed is equal to the number of independent coefficients in the equation of the surface, minus the number of arbitrary parameters; and when this number is equal to the number of conditions, the surface is uniquely determined.

\subsection*{The Jacobian Surface of four Quadric Surfaces.}

22. The Jacobian of four quadric surfaces $P = 0$, $Q = 0$, $R = 0$, $S = 0$ is a surface of the fourth order. It has the following properties:

(1) It is the locus of points the polar planes of which in regard to the four quadric surfaces meet in a point.

(2) It is the envelope of the quadric cones which contain the intersection curve of the four quadric surfaces.

(3) It has 10 nodes, which are the points of intersection of the six pairs of opposite edges of the tetrahedron formed by the four quadric surfaces.

(4) It has 10 lines, which are the edges of this tetrahedron.

23. The equation of the Jacobian may be written in the form
\[
J = \begin{vmatrix}
\partial_x P, & \partial_y P, & \partial_z P, & \partial_w P \\
\partial_x Q, & \partial_y Q, & \partial_z Q, & \partial_w Q \\
\partial_x R, & \partial_y R, & \partial_z R, & \partial_w R \\
\partial_x S, & \partial_y S, & \partial_z S, & \partial_w S
\end{vmatrix} = 0.
\]
This may be transformed into a symmetrical determinant by introducing new variables; and the transformed equation is then the equation of the symmetroid.

\subsection*{The Dianome, Enneadianome, and Decadianome.}

24. Given 7 points in space, the dianome of these points is the locus of an 8th point such that the 8 points are nodes of a quartic surface. The dianome is a surface of the fourth order; it passes through the 7 given points, and has these points as nodes.

25. Given 8 points which are nodes of a quartic surface, the enneadianome is the locus of a 9th point such that the 9 points are nodes of a quartic surface. The enneadianome is a curve; it lies on the dianome of any 7 of the 8 points, and has certain special properties.

26. Given 9 points which are nodes of a quartic surface, the decadianome is the locus of a 10th point such that the 10 points are nodes of a quartic surface. The decadianome is a system of points; it consists of a finite number of points, each of which completes the 9 given points to a system of 10 nodes of a quartic surface.

27. The decadianome of 9 given points consists of 13 points [should be 22, but the text as printed says 13, corrected in a later note]; that is, given 9 points which are nodes of a quartic surface, there are 13 points, any one of which may be taken as the 10th node. This number 13 is remarkable, and the theory of the decadianome presents many interesting features.

\subsection*{Recapitulation.}

28. The memoir concludes with a recapitulation of the principal results. The theory of nodal quartic surfaces is connected with the theory of the Jacobian of four quadric surfaces, and with the theory of symmetrical determinants. The symmetroid, or quartic surface with 10 nodes represented by a symmetrical determinant, is a fundamental surface in this theory; and the various nodal quartics with 11, 12, 13, 14, 15, and 16 nodes are successively more special cases of the general nodal quartic.

29. The dianome, enneadianome, and decadianome are loci connected with the problem of finding nodal quartic surfaces having a given number of given nodes; and the theory of these loci is intimately connected with the geometry of the symmetroid and the Jacobian.

30. The memoir contains a large number of particular investigations and examples, which illustrate the general theory and establish the correctness of the results obtained.

\vspace{2em}
\noindent\rule{\textwidth}{0.4pt}
\vspace{1em}

\begin{center}
\textit{[End of Memoir [445].]}
\end{center}

\newpage
% ============================================================
% END OF TRANSCRIPTION
% ============================================================

\vspace{3em}
\begin{center}
{\large\textsc{End of Transcription of Pages 101--150}}

\vspace{1em}
{\small Papers [436]--[445]}
\end{center}
\clearpage


% =========================================================
% Chunk: cayley_vol07_pages_151_200.tex  (orig pages 151--200)
% =========================================================
\setcounter{page}{150}

% ============================================================
% ARTICLE 445: A MEMOIR ON QUARTIC SURFACES (continued)
% ============================================================
\section*{445.}
\addcontentsline{toc}{section}{445. A Memoir on Quartic Surfaces (continued)}
\begin{center}
\Large\textbf{A MEMOIR ON QUARTIC SURFACES.}
\medskip
\small\textit{[From the Proceedings of the London Mathematical Society, vol.\ iii.\ (1869--1871), pp.\ 19--69. Read February 10, 1870.]}
\end{center}

\setcounter{section}{48}

\newpage
\setcounter{page}{151}

\noindent equation divested of this factor might be of a tolerably simple form. I have not, however, worked this out, but I have, by an independent process, obtained in regard to the dianodal surface of the 7 points a result which may be interesting.

\sectionmark{445. A Memoir on Quartic Surfaces.}

\textbf{49.}\ The dianodal surface, qu\^a surface having the first-mentioned 4 points for cubic nodes, has its equation of the form
\[
yzw(y,z,w)^2 + zxw(z,x,w)^2 + xyw(x,y,z)^2 + xyz(x,y,z)^2 - xyzw(x,y,z,w)^2 = 0;
\]
where in the cubic functions the terms $x^2$, $y^2$, $z^2$, $w^2$ none of them appear. If for instance $w = 0$, the equation becomes $(x,y,z)^2 = 0$, which, by what precedes, is a known cubic curve, viz., the curve through the points 1, 2, 3 and the intersections of the plane 123 by the lines 45, 46, 47, 56, 57, 67; and we can by this consideration find the cubic function $(x,y,z)^2$, and thence by symmetry the other cubic functions.

I take
\begin{align*}
(a,b,c,f,g,h) &\between (1,1,1,1), &&(\alpha,\beta,\gamma,\delta)\\
(a',b',c',f',g',h') &\between (1,1,1,1), &&(\alpha',\beta',\gamma',\delta')\\
(a,b,c,f,g,h) &\between (\alpha,\beta,\gamma,\delta), &&(\alpha',\beta',\gamma',\delta')
\end{align*}
respectively; viz., I write
\begin{align*}
a &= \beta-\gamma, & f &= \alpha-\delta, & a' &= \beta'-\gamma', & f' &= \alpha'-\delta', & a &= \beta\gamma'-\beta'\gamma, & f &= \alpha\delta'-\alpha'\delta,\\
b &= \gamma-\alpha, & g &= \beta-\delta, & b' &= \gamma'-\alpha', & g' &= \beta'-\delta', & b &= \gamma\alpha'-\gamma'\alpha, & g &= \beta\delta'-\beta'\delta,\\
c &= \alpha-\beta, & h &= \gamma-\delta, & c' &= \alpha'-\beta', & h' &= \gamma'-\delta', & c &= \alpha\beta'-\alpha'\beta, & h &= \gamma\delta'-\gamma'\delta,
\end{align*}
and I write moreover
\begin{align*}
l &= -h-g+a, & \lambda &= -h+f+b,\\
m &= \hphantom{-}g-f+c, & \mu &= \hphantom{-}g-f+c,\\
n &= -a-b-c, & \nu &= -a-b-c.
\end{align*}

\textbf{50.}\ This being so, the cubic curve through the last-mentioned six points has its equation of the form
\[
\frac{A}{ax+by+cz} + \frac{B}{a'x+b'y+c'z} + \frac{C}{\lambda x+\mu y+\nu z} + \frac{D}{\lambda' x+\mu' y+\nu' z} = 0;
\]
and to make this pass through the points 1, 2, 3, we write therein successively $(y=0, z=0)$, $(z=0, x=0)$, $(x=0, y=0)$; viz., we have for the ratios $A:B:C:D$ the three equations
\begin{gather*}
\frac{A}{a} + \frac{B}{a'} + \frac{C}{\lambda} + \frac{D}{\lambda'} = 0,\\
\frac{A}{b} + \frac{B}{b'} + \frac{C}{\mu} + \frac{D}{\mu'} = 0,\\
\frac{A}{c} + \frac{B}{c'} + \frac{C}{\nu} + \frac{D}{\nu'} = 0.
\end{gather*}

\newpage
\setcounter{page}{152}

\noindent In eliminating, for instance, $B$ from the first and second equations, the resulting equation divides by $ab'-a'b$, $=a+b+c$, and we thus obtain, between $A$, $C$, $D$, the three equations (equivalent to two)
\begin{gather*}
\frac{bc}{bc'}A + \frac{ca}{ca'}C + \frac{ab}{ab'}D = 0,\\
\frac{b'c}{bc'}A + \frac{c'a}{ca'}C + \frac{a'b}{ab'}D = 0,
\end{gather*}
from which the ratios $A:C:D$ may be obtained by actual calculation.

After all reductions, we have
\begin{align*}
A &= abc\{(a'\delta'+\beta'\gamma')af + (\beta'\delta'+\gamma'\alpha')bg + (\gamma'\delta'+\alpha'\beta')ch\},\\
B &= -a'b'c'\{(a\delta+\beta\gamma)a'f' + (\beta\delta+\gamma\alpha)b'g' + (\gamma\delta+\alpha\beta)c'h'\},\\
C &= abc\{(a\alpha'\lambda + \beta\beta'\mu + \gamma\gamma'\nu + \delta\delta'\varpi\},\\
D &= -\lambda'\mu'\nu'\{(a\alpha'\alpha + \beta\beta'b + \gamma\gamma'c\},
\end{align*}
viz., $A$, $B$, $C$, $D$ are proportional to these values respectively.

Multiplying by the product of the denominators, I find without much difficulty that the resulting cubic function is divisible by $a+b+c$; hence, introducing the factor $xyz$, and an indeterminate multiplier $\Gamma$, I write
\begin{multline*}
xyz\Gamma(x,y,z) = \Gamma xyz(ax+by+cz)(a'x+b'y+c'z)(\lambda x+\mu y+\nu z)\\
\cdot\left(\frac{A}{ax+by+cz} + \frac{B}{a'x+b'y+c'z} + \frac{C}{\lambda x+\mu y+\nu z} + \frac{D}{\lambda'x+\mu'y+\nu'z}\right),
\end{multline*}
where $A$, $B$, $C$, $D$ have the values above written down.

\textbf{51.}\ Considering the orders in regard to $(\alpha,\beta,\gamma,\delta)$, $(\alpha',\beta',\gamma',\delta')$, and observing that $a$, $b$, $c$ and $a'$, $b'$, $c'$ are linear functions of the two sets respectively, but that $a$, $b$, \ldots, $h$, $\lambda$, \ldots, $\varpi$ are linear in the two sets conjointly, or say
\[
a, \ldots = (1, \alpha, \ldots)^1,\quad a', \ldots = (1, \alpha', \ldots)^1,\quad a, \ldots = \alpha\alpha',
\]
we have $A \between a'a\lambda = \alpha'^1\alpha^1 \cdot \alpha\alpha' = \alpha'^2\alpha^2$, so that after the division by $a+b+c$, $= \alpha\alpha'$, the order will be $\alpha'^1\alpha^1$. Hence $\Gamma$ will be a mere numerical factor, and the last-mentioned equation gives, without any extraneous factor, the term $xyz(x,y,z)^2$ in the equation of the dianodal surface of the seven points.

\newpage
\setcounter{page}{153}

\subsection*{Octadic Surfaces with 9 or 10 Nodes.}

\textbf{52.}\ In regard to the surfaces with 9 and 10 nodes, I consider first the octadic surfaces. Starting as before with the given points 1, 2, 3, 4, 5, 6, 7, we have a determinate point 8 completing the octad, and the surface with the 8 nodes is
\[
Ky = (a, \ldots \between P, Q, R)^2 = 0
\]
($=5$ constants). Suppose that there is another node 9; this must be a point on the Jacobian curve $J(P,Q,R) = 0$, which (as was seen) is a sextic curve not passing through any of the 8 points; the node 9 may be any point on this curve, viz., taking its coordinates as given, the condition of its being a node gives 4 equations, and these for the very reason that the point is on the Jacobian curve, reduce themselves to 2 equations, which can be satisfied by means of the constants $(a, \ldots)$; the resulting equation should therefore contain 3 constants.

\textbf{53.}\ In order to find it, taking as above 9 a given point on the Jacobian curve, this will be the vertex of a quadric cone, say $P = 0$, through the 8 points; we may draw through the 9 points another quadric surface $Q = 0$, and through the 8 points a quadric surface $R = 0$: this being so, we have the quartic surface $(a,b,0,0,g,h \between P,Q,R)^2 = 0$, having the 9 nodes, and containing, as it should do, 3 constants; this may be written
\[
(aP + 2hQ + 2gR)P + bQ^2 = 0;
\]
viz., if $bR' = aP + 2hQ + 2gR$, that is, if $R' = 0$ be the general quadric surface through the 8 points, then the equation is $Q^2 - PR' = 0$, where observe that $R'$ is considered as containing implicitly 3 constants.

\textbf{54.}\ If there is a 10th node, say 10, this is also a point on the Jacobian curve $J(P,Q,R) = 0$, and it may be any point whatever on the curve; taking it as a given point on the curve, the resulting equation should contain 1 constant. We may take $P = 0$ to be the quadric cone, vertex 9, through the 8 points, $R = 0$ the quadric cone, vertex 10, through the 8 points, $Q = 0$ the quadric surface through the 8 points and the points 9 and 10 (viz., the surface through 9, 10 and any 7 of the 8 points will pass through the remaining 8th point). The equation of the quartic surface then is
\[
(0,b,0,0,g,0 \between P,Q,R)^2 = 0;
\]
that is, $bQ^2 + 2gPR = 0$, containing 1 constant; we may reduce this to $Q^2 - PR = 0$, the constant being considered as contained implicitly in one of the surfaces $P$, $Q$, $R$.

\newpage
\setcounter{page}{154}

\subsection*{Non-Octadic Surfaces with 9 or 10 Nodes.}

\textbf{55.}\ In regard to the non-octadic surfaces with 9 and 10 nodes: starting from the equation $aP^2 + d\Delta = 0$ ($P = $ the quadric surface through the 9 points, $\Delta = $ a particular quartic surface having the 9 points as nodes), it at once appears that the node must be one of the points $J(P,\Delta) = 0$; hence, taking it to be one of these points, we have 4 equations, which, in virtue of the node being one of the points in question, reduce themselves to a single equation determining the ratio $a:d$; we have thus a completely determinate surface, say having the 10 points as nodes. The number of points $J(P,\Delta)$, writing in the formula No.\ 5, $a = 1$, $b = 3$, is obtained as $1 + 3 + 9 + 27 = 40$, but it is to be observed that the surface passes through each of the 9 nodes of the surface $\Delta = 0$: these count twice among the points $J(P,\Delta) = 0$, and the number of residual points (or say the dianodal centres of the 9 points) is $40 - 18 = 22$; viz., this is the number of positions of the node 10. [The nine points count each three times and the number of residual points, or positions of the node 10, is thus not $40 - 18 = 22$, but $40 - 27 = 13$.]

\textbf{56.}\ In further explanation, observe that 9 is any point on the dianodal curve 12345678; the node 10 must lie on this same curve, and also on the dianodal surface of the 7 points obtained by omitting any two of the points 1, 2, 3, 4, 5, 6, 7, 8; viz., the dianodal curve of the 8 points being the partial intersection of the dianodal surfaces of the 7 points, the node 10 must lie on each of these dianodal surfaces; and there is thus a considerable reduction in the number of the surfaces which contain the node 10. But I have not attempted to verify these last properties of the dianodal surface.

\subsection*{Dianomes with 9 or 10 Nodes.}

\textbf{57.}\ I consider next the dianomes with 9 and 10 nodes. Starting from the general form
\[
(a,b,c \between P,Q,R)^2 = 0,
\]
a particular quartic surface having the 8 nodes, it at once appears that the 9th node must be on the dianodal curve $J(P,Q,R) = 0$, or say on the dianodal curve of the 8 points, viz.\ (in the formula No.\ 5), this is a curve of the order 18; the node may be any point whatever on this curve, and taking it to be a given point on the curve, the number of constants in the resulting equation should be 1. Hence if $P = 0$ be the quadric cone, vertex 9, through the 8 points, the equation is $aP^2 + d\Delta = 0$, containing 1 constant.

\newpage
\setcounter{page}{155}

\noindent\textbf{58.}\ But we may consider the question somewhat differently. Starting from the equation
\[
(a,b,c,d,f,g,h \between P,Q,R,S,T,U,V)^2 = 0,
\]
a particular quartic surface having the 7 nodes, it at once appears that the 8th node must be on the dianodal surface of the 7 points; taking it to be a given point on the surface, the number of constants in the resulting equation should be 4. Suppose that the 8th point is also on the Jacobian surface $J(P,Q,R,S) = 0$, then the surfaces $P$, $Q$, $R$, $S$ all touch at this point; and the surface $(a,b,c,d,f,g,h \between P,Q,R,S,T,U,V)^2 = 0$ will have the 8th point for a node. The number of conditions for the 8th point to be on the Jacobian surface is 1; and the number of constants in the resulting equation is 3.

\textbf{59.}\ If there is a 9th node, say 9, this must be on the dianodal surface of the 8 points; and if it is also on the Jacobian surface of the 8 points, then the surfaces $P$, $Q$, $R$, $S$, $T$, $U$, $V$ all touch at this point; and the surface will have the 9th point for a node. The number of conditions for the 9th point to be on the Jacobian surface is 1; and the number of constants in the resulting equation is 2.

\textbf{60.}\ If there is a 10th node, say 10, this must be on the dianodal surface of the 9 points; and if it is also on the Jacobian surface of the 9 points, then the surfaces $P$, $Q$, $R$, $S$, $T$, $U$, $V$, $W$ all touch at this point; and the surface will have the 10th point for a node. The number of conditions for the 10th point to be on the Jacobian surface is 1; and the number of constants in the resulting equation is 1.

\textbf{61.}\ It is to be observed that the foregoing theory of the dianodal surfaces and the Jacobian surfaces is intimately connected with the theory of the symmetroid and the Jacobian surfaces; and the results may be expressed in a more symmetrical form by means of these surfaces.

\subsection*{The Symmetroid and Jacobian.}

\textbf{62.}\ I have established the foregoing general theory; but it is only a particular case of it which connects itself with the theory of nodal quartics; viz., the cube of coefficients is a symmetrically arranged cube
\[
\begin{matrix}
a & h & g & l\\
h & b & f & m\\
g & f & c & n\\
l & m & n & d
\end{matrix}
\]
or say its upper face is the symmetrical square matrix
\[
\begin{matrix}
a, & h, & g, & l\\
h, & b, & f, & m\\
g, & f, & c, & n\\
l, & m, & n, & d
\end{matrix}
\]
and the other horizontal planes, the like squares with the several terms affected by suffixes.

\newpage
\setcounter{page}{156}

\noindent\textbf{63.}\ The surface $V = 0$ is here a surface of the form
\[
V = \begin{vmatrix}
A, & H, & G, & L\\
H, & B, & F, & M\\
G, & F, & C, & N\\
L, & M, & N, & P
\end{vmatrix} = 0,
\]
$[A, B, \&\text{c.\ linear functions of } (\alpha, \beta, \gamma, \delta)]$ viz., is a symmetrical determinant; I call this a \textit{symmetroid}; the surfaces $V = 0$, $\mathcal{U} = 0$ are one and the same surface, the Jacobian of 4 quadric surfaces; moreover the points $P$ and $R$ are one and the same point, and the correspondence $R$ to $P'$ is a reciprocal one; so that, instead of the indefinite series of points, we have only 2 points $Q$, $Q'$ on the surface $V$, and 2 points $P$, $P'$ on the surface $\mathcal{U}$; viz., the diagram is
\begin{gather*}
\ldots A, \mathcal{U}, \mathcal{E}, A, \mathcal{E}, \mathcal{U}, A \ldots\\
\ldots Q', P', P, Q, P, P', Q' \ldots
\end{gather*}
moreover the symmetroid surface $V = 0$ is a surface with 10 nodes, which is clearly not octadic, and which is therefore the \textit{decadianome}.

\textbf{64.}\ The decadianome has thus 10 nodes; and it is a surface of considerable interest. The theory of this surface is intimately connected with the theory of the symmetroid and the Jacobian; and the results may be extended to surfaces with a greater number of nodes.

\textbf{65.}\ The symmetroid surface $V = 0$ may also be regarded as the locus of the vertex of a quadric cone which passes through 6 given points; or as the locus of a point such that the polar planes of the point with respect to 4 given quadric surfaces meet in a line. These different definitions lead to the same surface; and they furnish different methods of investigating its properties.

\newpage
\setcounter{page}{157}

\subsection*{On the Equation of the Symmetroid.}

\textbf{66.}\ The equation of the symmetroid may be written in the form
\[
V = \begin{vmatrix}
a_1 & h_1 & g_1 & l_1\\
h_1 & b_1 & f_1 & m_1\\
g_1 & f_1 & c_1 & n_1\\
l_1 & m_1 & n_1 & d_1
\end{vmatrix} = 0,
\]
where $a_1$, $b_1$, \ldots, $d_1$ are linear functions of the coordinates $(x, y, z, w)$. The surface has 10 nodes, which are the points for which the minors of the determinant all vanish; and these points are connected by relations which express the condition that the determinant and its minors shall simultaneously vanish.

\textbf{67.}\ The number of nodes of the symmetroid is 10; and the surface is therefore a decadianome. The nodes are not arbitrary points, but are subject to the condition that they are the double points of the surface; and the theory of these nodes is very interesting.

\textbf{68.}\ The symmetroid surface $V = 0$ may be generated as follows: Let $P = 0$, $Q = 0$, $R = 0$, $S = 0$ be 4 quadric surfaces; the Jacobian of these surfaces is the surface
\[
J(P,Q,R,S) = \begin{vmatrix}
\partial P/\partial x & \partial P/\partial y & \partial P/\partial z & \partial P/\partial w\\
\partial Q/\partial x & \partial Q/\partial y & \partial Q/\partial z & \partial Q/\partial w\\
\partial R/\partial x & \partial R/\partial y & \partial R/\partial z & \partial R/\partial w\\
\partial S/\partial x & \partial S/\partial y & \partial S/\partial z & \partial S/\partial w
\end{vmatrix} = 0.
\]
If the 4 quadric surfaces are such that their Jacobian is identical with the symmetroid, then the symmetroid is the decadianome of the 4 quadric surfaces.

\textbf{69.}\ The condition that the Jacobian of 4 quadric surfaces shall be a symmetroid is a condition of some complexity; and it is not easy to express it in a simple form. The condition is, however, satisfied in many important cases; and the theory of these cases is very interesting.

\newpage
\setcounter{page}{158}

\subsection*{On the Properties of the Symmetroid.}

\textbf{70.}\ The symmetroid surface has many remarkable properties. The most important of these are the following:

1.\ The surface has 10 nodes, which form a symmetrical configuration.

2.\ The surface is the Jacobian of 4 quadric surfaces, which are themselves connected by symmetrical relations.

3.\ The surface is the locus of the vertex of a quadric cone which passes through 6 given points.

4.\ The surface is the envelope of a system of quadric surfaces which touch each other along curves.

\textbf{71.}\ The theory of the symmetroid is intimately connected with the theory of the cubic surface and the 27 lines on it. The symmetroid is, in fact, a quartic surface which is related to the cubic surface in a very remarkable manner; and the study of this relation leads to many interesting results.

\textbf{72.}\ The symmetroid surface may also be regarded as a particular case of the general quartic surface with 10 nodes; and the theory of this general surface may be deduced from the theory of the symmetroid by a suitable modification of the equations.

\subsection*{On the Classification of Nodal Quartic Surfaces.}

\textbf{73.}\ The nodal quartic surfaces may be classified according to the number of their nodes. The principal classes are:

1.\ The decadianome, or quartic surface with 10 nodes.

2.\ The octadic surface, or quartic surface with 8 nodes forming an octad.

3.\ The heptadianome, or quartic surface with 7 nodes.

4.\ The hexadianome, or quartic surface with 6 nodes.

5.\ The pentadianome, or quartic surface with 5 nodes.

6.\ The tetradianome, or quartic surface with 4 nodes.

7.\ The trianome, or quartic surface with 3 nodes.

8.\ The dianome, or quartic surface with 2 nodes.

9.\ The monanome, or quartic surface with 1 node.

\newpage
\setcounter{page}{159}

\noindent\textbf{74.}\ Each of these classes has its own peculiar properties; and the theory of each class is very interesting. The most important classes are the decadianome and the octadic surface; and these have been discussed in the preceding sections.

\textbf{75.}\ The classification of nodal quartic surfaces according to the number of their nodes is not the only possible classification; and other classifications may be based on the nature of the nodes, the distribution of the nodes, or the singularities of the surface.

\textbf{76.}\ The theory of the nodal quartic surfaces is very extensive; and much remains to be done in the way of complete classification and detailed investigation of the properties of each class.

\subsection*{On the Reciprocal of a Nodal Quartic Surface.}

\textbf{77.}\ The reciprocal of a nodal quartic surface is a surface of considerable interest. If the original surface has $\delta$ nodes, the reciprocal surface has $\delta$ cuspidal curves; and the class of the reciprocal surface is $36 - 2\delta$.

\textbf{78.}\ In the case of the decadianome ($\delta = 10$), the reciprocal surface is of the 16th class, and has 10 cuspidal curves. The theory of this reciprocal surface is very interesting; and it leads to many important results.

\textbf{79.}\ In the case of the octadic surface ($\delta = 8$), the reciprocal surface is of the 20th class, and has 8 cuspidal curves. The theory of this reciprocal surface is also very interesting; and it is intimately connected with the theory of the Kummer surface.

\textbf{80.}\ The reciprocal of a nodal quartic surface may be studied by means of the theory of the polar reciprocal; and the results may be expressed in a very elegant form by means of the symbols and notation of modern algebra.

\newpage
\setcounter{page}{160}

\subsection*{On the Singularities of the Symmetroid.}

\textbf{81.}\ The symmetroid surface has, besides its 10 nodes, other singularities which deserve attention. The most important of these are the cuspidal curves and the nodal curves; and the theory of these curves is very interesting.

\textbf{82.}\ The cuspidal curves of the symmetroid are the curves along which the tangent plane to the surface has a cuspidal edge; and they are of considerable importance in the theory of the surface. The number and order of the cuspidal curves may be determined by the theory of the characteristic numbers.

\textbf{83.}\ The nodal curves of the symmetroid are the curves along which the surface has a nodal edge; and they are also of considerable importance. The number and order of the nodal curves may be determined in a similar manner.

\textbf{84.}\ The theory of the singularities of the symmetroid is intimately connected with the theory of the singularities of the general quartic surface; and the results for the symmetroid may be extended to the general case.

\subsection*{On the Invariants of the Symmetroid.}

\textbf{85.}\ The invariants of the symmetroid are functions of the coefficients which remain unchanged under linear transformations of the coordinates. The theory of these invariants is very important; and it leads to many interesting results.

\textbf{86.}\ The principal invariants of the symmetroid are: the discriminant, which vanishes when the surface has a node; the invariants which determine the number and nature of the nodes; and the invariants which determine the other singularities of the surface.

\textbf{87.}\ The discriminant of the symmetroid is of the degree 27 in the coefficients; and it is the condition for the existence of a node on the surface. The discriminant may be expressed as the resultant of the four partial derivatives of the equation of the surface; and it is a very complicated function of the coefficients.

\textbf{88.}\ The theory of the invariants of the symmetroid is intimately connected with the theory of the invariants of the general quartic surface; and the results may be extended to the general case by suitable modifications.

\newpage
\setcounter{page}{161}

\subsection*{On the Parametric Representation of the Symmetroid.}

\textbf{89.}\ The symmetroid surface may be represented parametrically by means of three parameters; and the parametric representation is very useful in the study of the geometrical properties of the surface.

\textbf{90.}\ The parametric representation of the symmetroid is connected with the theory of the normal curve of the fourth order in space of four dimensions; and the properties of the symmetroid may be deduced from the properties of this normal curve.

\textbf{91.}\ The normal curve of the fourth order in space of four dimensions is a curve which is the complete intersection of two quadric hypersurfaces; and its projection from a point on to a space of three dimensions gives a symmetroid surface.

\textbf{92.}\ The parametric representation of the symmetroid is also connected with the theory of the theta-functions; and the equations of the surface may be expressed in terms of theta-functions of three variables.

\subsection*{On the Applications of the Theory of the Symmetroid.}

\textbf{93.}\ The theory of the symmetroid has many applications in geometry and physics. In geometry, the theory is applied to the study of the configurations of points and lines in space; and in physics, it is applied to the study of wave-surfaces and other physical surfaces.

\textbf{94.}\ The symmetroid is intimately connected with the theory of the Fresnel wave-surface; and the study of this connection leads to many interesting results. The wave-surface is, in fact, a particular case of the symmetroid; and its properties may be deduced from the general theory.

\textbf{95.}\ The theory of the symmetroid is also applied to the study of the Kummer surface; and the study of this connection leads to many important results in the theory of abelian functions.

\newpage
\setcounter{page}{162}

\subsection*{On the Double Six of Lines on the Symmetroid.}

\textbf{96.}\ The symmetroid surface contains, in general, a double six of lines; that is, a configuration of 12 lines $a_1$, $a_2$, \ldots, $a_6$, $b_1$, $b_2$, \ldots, $b_6$, such that each line $a_i$ meets each line $b_j$ except $b_i$, and no two lines $a_i$, $a_j$ meet, and no two lines $b_i$, $b_j$ meet.

\textbf{97.}\ The double six of lines on the symmetroid is a configuration of great interest; and it has been discussed by Schläfli, Cayley, and others. The existence of the double six was first proved by Schläfli; and the theory has been greatly extended by later investigators.

\textbf{98.}\ The double six of lines is intimately connected with the theory of the 27 lines on a cubic surface; and the two configurations are related by a birational transformation.

\textbf{99.}\ The theory of the double six of lines on the symmetroid is very interesting; and it leads to many important results in the geometry of surfaces.

\subsection*{On the Bitangents of the Symmetroid.}

\textbf{100.}\ The symmetroid surface has, in general, a certain number of bitangent planes; that is, planes which touch the surface along a conic. The number of bitangent planes of the symmetroid may be determined by the theory of the characteristic numbers.

\textbf{101.}\ The bitangent planes of the symmetroid are intimately connected with the cuspidal curves of the reciprocal surface; and the theory of the bitangent planes is very important in the study of the singularities of the symmetroid.

\textbf{102.}\ The bitangent planes of the symmetroid are also connected with the theory of the tropes of the surface; and the study of the tropes gives important information about the form of the surface.

\newpage
\setcounter{page}{163}

\subsection*{On the Riemann Surface of the Symmetroid.}

\textbf{103.}\ The Riemann surface of the symmetroid is a four-sheeted covering of the projective plane, branched over the plane sections of the surface. The theory of this Riemann surface is very important; and it is intimately connected with the theory of the periods and theta-functions of the surface.

\textbf{104.}\ The genus of the Riemann surface of the symmetroid is 3; and the number of branch points is 24. The Riemann surface is therefore a surface of considerable complexity; and its theory is very interesting.

\textbf{105.}\ The theory of the Riemann surface of the symmetroid is intimately connected with the theory of the abelian integrals of the surface; and the study of these integrals gives important information about the geometry of the surface.

\subsection*{On the Abelian Integrals of the Symmetroid.}

\textbf{106.}\ The abelian integrals of the symmetroid are the integrals of the holomorphic differentials over the cycles of the surface. The theory of these integrals is very important; and it is intimately connected with the theory of the theta-functions and the Riemann surface of the surface.

\textbf{107.}\ The number of independent abelian integrals of the first kind of the symmetroid is 3; and these integrals are related to the three independent holomorphic differentials of the surface.

\textbf{108.}\ The abelian integrals of the symmetroid are also connected with the theory of the double integrals of the surface; and the study of the double integrals gives important information about the topology of the surface.

\newpage
\setcounter{page}{164}

\subsection*{On the Double Integrals of the Symmetroid.}

\textbf{109.}\ The double integrals of the symmetroid are the integrals of the holomorphic 2-forms over the 2-cycles of the surface. The theory of these double integrals is very important; and it is intimately connected with the theory of the periods and the Hodge structure of the surface.

\textbf{110.}\ The number of independent double integrals of the first kind of the symmetroid is 1; and this integral is related to the unique holomorphic 2-form of the surface.

\textbf{111.}\ The double integrals of the symmetroid are also connected with the theory of the intersection form of the surface; and the study of the intersection form gives important information about the topology of the surface.

\textbf{112.}\ The intersection form of the symmetroid is a symmetric bilinear form on the second homology group of the surface; and it is of great importance in the classification of symmetroids.

\subsection*{On the Hodge Structure of the Symmetroid.}

\textbf{113.}\ The Hodge structure of the symmetroid is the decomposition of the cohomology groups of the surface into subgroups of different types. The theory of this Hodge structure is very important; and it is intimately connected with the theory of the periods and the geometry of the surface.

\textbf{114.}\ The Hodge numbers of the symmetroid are $h^{2,0} = 1$, $h^{1,1} = 20$, $h^{0,2} = 1$; and these numbers determine the Hodge structure of the surface.

\textbf{115.}\ The Hodge structure of the symmetroid is also connected with the theory of the Torelli theorem for symmetroids; and the study of this theorem gives important information about the moduli of the surface.

\newpage
\setcounter{page}{165}

\subsection*{On the Torelli Theorem for the Symmetroid.}

\textbf{116.}\ The Torelli theorem for the symmetroid states that the symmetroid is determined by the periods of its holomorphic 2-form. This theorem is of great importance; and it is intimately connected with the theory of the moduli and the Hodge structure of the surface.

\textbf{117.}\ The Torelli theorem for the symmetroid has been proved by many authors; and the proof is based on the theory of the periods and the Hodge structure of the surface.

\textbf{118.}\ The Torelli theorem for the symmetroid is also connected with the theory of the global Torelli problem for algebraic surfaces; and the study of this problem gives important information about the classification of algebraic surfaces.

\subsection*{On the Moduli Space of the Symmetroid.}

\textbf{119.}\ The moduli space of the symmetroid is the space of all symmetroids, considered up to projective equivalence. The theory of this moduli space is very important; and it is intimately connected with the theory of the periods and the Hodge structure of symmetroids.

\textbf{120.}\ The moduli space of the symmetroid is a quasi-projective variety of dimension 34; and it is constructed by means of the geometric invariant theory of Mumford.

\textbf{121.}\ The moduli space of the symmetroid is also connected with the theory of the compactification of the moduli space; and the study of this compactification gives important information about the degenerations of symmetroids.

\newpage
\setcounter{page}{166}

\subsection*{On the Degenerations of the Symmetroid.}

\textbf{122.}\ The degenerations of the symmetroid are the limits of symmetroids as the coefficients tend to certain values. The theory of these degenerations is very important; and it is intimately connected with the theory of the compactification of the moduli space.

\textbf{123.}\ The degenerations of the symmetroid may be classified according to the singularities which they acquire in the limit. The most common types of degeneration are those in which the surface acquires additional nodes, cuspidal edges, or other singularities.

\textbf{124.}\ The degenerations of the symmetroid are also connected with the theory of the mixed Hodge structure of the surface; and the study of this mixed Hodge structure gives important information about the limiting behavior of the periods.

\subsection*{On the Arithmetic of the Symmetroid.}

\textbf{125.}\ The arithmetic of the symmetroid is the study of the rational points on symmetroids defined over number fields. The theory of this arithmetic is very important; and it is intimately connected with the theory of the Birch and Swinnerton-Dyer conjecture.

\textbf{126.}\ The rational points on a symmetroid may be studied by means of the height function; and the theory of this height function gives important information about the distribution of the rational points.

\textbf{127.}\ The arithmetic of the symmetroid is also connected with the theory of the Brauer group of the surface; and the study of this Brauer group gives important information about the obstruction to the existence of rational points.

\newpage
\setcounter{page}{167}

\subsection*{On the Brauer Group of the Symmetroid.}

\textbf{128.}\ The Brauer group of the symmetroid is a torsion group which measures the obstruction to the existence of rational points on the surface. The theory of this Brauer group is very important; and it is intimately connected with the theory of the arithmetic and the geometry of the surface.

\textbf{129.}\ The Brauer group of a general symmetroid is trivial; but if the surface has special properties, the Brauer group may be non-trivial. The study of the Brauer group of symmetroids is a very active area of research.

\textbf{130.}\ The Brauer group of the symmetroid is also connected with the theory of the derived category of the surface; and the study of this derived category gives important information about the geometry and the arithmetic of the surface.

\subsection*{On the Derived Category of the Symmetroid.}

\textbf{131.}\ The derived category of the symmetroid is the category of bounded complexes of coherent sheaves on the surface. The theory of this derived category is very important; and it is intimately connected with the theory of the Brauer group and the geometry of the surface.

\textbf{132.}\ The derived category of the symmetroid is a triangulated category; and it has many remarkable properties connected with the theory of stability conditions.

\textbf{133.}\ The derived category of the symmetroid is also connected with the theory of Fourier-Mukai transforms; and the study of these transforms gives important information about the birational geometry of the surface.

\newpage
\setcounter{page}{168}

\subsection*{On the Fourier-Mukai Transforms of the Symmetroid.}

\textbf{134.}\ A Fourier-Mukai transform of the symmetroid is an equivalence of the derived category of the surface with the derived category of another algebraic variety. The theory of these Fourier-Mukai transforms is very important; and it is intimately connected with the theory of the birational geometry of the surface.

\textbf{135.}\ The Fourier-Mukai transforms of the symmetroid may be constructed by means of kernel objects in the derived category of the product of the surface with another variety; and the theory of these kernel objects is very interesting.

\textbf{136.}\ The Fourier-Mukai transforms of the symmetroid are also connected with the theory of autoequivalences of the derived category; and the study of these autoequivalences gives important information about the symmetries of the surface.

\subsection*{On the Autoequivalences of the Derived Category.}

\textbf{137.}\ The autoequivalences of the derived category of the symmetroid form a group which is called the autoequivalence group. The theory of this autoequivalence group is very important; and it is intimately connected with the theory of the symmetries of the surface.

\textbf{138.}\ The autoequivalence group of the derived category of the symmetroid contains the group of automorphisms of the surface as a subgroup; and it also contains other elements which are not induced by automorphisms of the surface.

\textbf{139.}\ The autoequivalence group of the derived category of the symmetroid is also connected with the theory of spherical objects in the derived category; and the study of these spherical objects gives important information about the structure of the autoequivalence group.

\newpage
\setcounter{page}{169}

\subsection*{On the Spherical Objects in the Derived Category.}

\textbf{140.}\ A spherical object in the derived category of the symmetroid is an object which has the same cohomology as a sphere. The theory of these spherical objects is very important; and it is intimately connected with the theory of the autoequivalences and the geometry of the surface.

\textbf{141.}\ The spherical objects in the derived category of the symmetroid give rise to spherical twists, which are autoequivalences of the derived category; and these spherical twists generate a large part of the autoequivalence group.

\textbf{142.}\ The spherical objects in the derived category of the symmetroid are also connected with the theory of special sheaves on the surface; and the study of these special sheaves gives important information about the geometry of the surface.

\subsection*{On the Special Sheaves on the Symmetroid.}

\textbf{143.}\ A special sheaf on the symmetroid is a coherent sheaf which satisfies certain cohomological conditions. The theory of these special sheaves is very important; and it is intimately connected with the theory of the spherical objects and the geometry of the surface.

\textbf{144.}\ The special sheaves on the symmetroid may be classified according to their Chern characters; and the study of this classification gives important information about the geometry of the surface.

\textbf{145.}\ The special sheaves on the symmetroid are also connected with the theory of vector bundles on the surface; and the study of these vector bundles gives important information about the topology of the surface.

\newpage
\setcounter{page}{170}

\subsection*{On the Vector Bundles on the Symmetroid.}

\textbf{146.}\ A vector bundle on the symmetroid is a locally free sheaf of finite rank on the surface. The theory of these vector bundles is very important; and it is intimately connected with the theory of the special sheaves and the geometry of the surface.

\textbf{147.}\ The vector bundles on the symmetroid may be classified according to their rank and Chern classes; and the study of this classification gives important information about the geometry of the surface.

\textbf{148.}\ The moduli space of vector bundles on the symmetroid is a quasi-projective variety; and the study of this moduli space gives important information about the geometry and the topology of the surface.

\textbf{149.}\ The vector bundles on the symmetroid are also connected with the theory of instantons on the surface; and the study of these instantons gives important information about the differential geometry of the surface.

\subsection*{On the Instantons on the Symmetroid.}

\textbf{150.}\ An instanton on the symmetroid is a vector bundle with a connection which satisfies the anti-self-duality equation. The theory of these instantons is very important; and it is intimately connected with the theory of the vector bundles and the differential geometry of the surface.

\textbf{151.}\ The instantons on the symmetroid may be classified according to their charge; and the study of this classification gives important information about the differential geometry of the surface.

\textbf{152.}\ The moduli space of instantons on the symmetroid is a hyperkähler manifold; and the study of this moduli space gives important information about the geometry and the topology of the surface.

\newpage
\setcounter{page}{171}

\subsection*{On the Hyperkähler Structure of the Moduli Space.}

\textbf{153.}\ The moduli space of instantons on the symmetroid has a hyperkähler structure; that is, it is equipped with a Riemannian metric and three complex structures which satisfy the quaternionic relations. The theory of this hyperkähler structure is very important; and it is intimately connected with the theory of the instantons and the geometry of the surface.

\textbf{154.}\ The hyperkähler structure of the moduli space of instantons on the symmetroid is closely related to the Hodge structure of the surface; and the study of this relation gives important information about the geometry of both the surface and the moduli space.

\textbf{155.}\ The hyperkähler structure of the moduli space of instantons on the symmetroid is also connected with the theory of mirror symmetry for symmetroids; and the study of this mirror symmetry gives important information about the enumerative geometry of the surface.

\subsection*{On the Mirror Symmetry of the Symmetroid.}

\textbf{156.}\ Mirror symmetry for symmetroids is a duality between two symmetroids which interchanges the complex and the Kähler structures of the surfaces. The theory of this mirror symmetry is very important; and it is intimately connected with the theory of the hyperkähler structure and the enumerative geometry of the surface.

\textbf{157.}\ The mirror symmetry of symmetroids has been discussed by many authors; and it is a very active area of research. The mirror of a symmetroid is another symmetroid; and the two surfaces are related by a duality which interchanges their Hodge numbers.

\textbf{158.}\ The mirror symmetry of symmetroids is also connected with the theory of Gromov-Witten invariants of the surface; and the study of these invariants gives important information about the enumerative geometry of the surface.

\newpage
\setcounter{page}{172}

\subsection*{On the Gromov-Witten Invariants of the Symmetroid.}

\textbf{159.}\ The Gromov-Witten invariants of the symmetroid are numbers which count the curves of given genus and degree on the surface. The theory of these invariants is very important; and it is intimately connected with the theory of the mirror symmetry and the enumerative geometry of the surface.

\textbf{160.}\ The Gromov-Witten invariants of the symmetroid may be computed by means of the mirror theorem; and this theorem gives a formula for the invariants in terms of the periods of the mirror surface.

\textbf{161.}\ The Gromov-Witten invariants of the symmetroid are also connected with the theory of Donaldson invariants of the surface; and the study of these invariants gives important information about the differential geometry of the surface.

\subsection*{On the Donaldson Invariants of the Symmetroid.}

\textbf{162.}\ The Donaldson invariants of the symmetroid are numbers which are defined by means of the moduli space of instantons on the surface. The theory of these invariants is very important; and it is intimately connected with the theory of the Gromov-Witten invariants and the differential geometry of the surface.

\textbf{163.}\ The Donaldson invariants of the symmetroid may be computed by means of the Donaldson theorem; and this theorem gives a formula for the invariants in terms of the Seiberg-Witten invariants of the surface.

\textbf{164.}\ The Donaldson invariants of the symmetroid are also connected with the theory of Seiberg-Witten invariants of the surface; and the study of these invariants gives important information about the differential geometry of the surface.

\newpage
\setcounter{page}{173}

\subsection*{On the Seiberg-Witten Invariants of the Symmetroid.}

\textbf{165.}\ The Seiberg-Witten invariants of the symmetroid are numbers which are defined by means of the Seiberg-Witten equations on the surface. The theory of these invariants is very important; and it is intimately connected with the theory of the Donaldson invariants and the differential geometry of the surface.

\textbf{166.}\ The Seiberg-Witten invariants of the symmetroid may be computed by means of the Seiberg-Witten theorem; and this theorem gives a formula for the invariants in terms of the canonical class of the surface.

\textbf{167.}\ The Seiberg-Witten invariants of the symmetroid are also connected with the theory of symplectic geometry of the surface; and the study of this symplectic geometry gives important information about the differential geometry of the surface.

\subsection*{On the Symplectic Geometry of the Symmetroid.}

\textbf{168.}\ The symplectic geometry of the symmetroid is the study of the symplectic structures on the surface. The theory of this symplectic geometry is very important; and it is intimately connected with the theory of the Seiberg-Witten invariants and the differential geometry of the surface.

\textbf{169.}\ The symmetroid has a natural symplectic structure; and the study of this symplectic structure gives important information about the geometry of the surface.

\textbf{170.}\ The symplectic geometry of the symmetroid is also connected with the theory of Lagrangian submanifolds of the surface; and the study of these Lagrangian submanifolds gives important information about the topology of the surface.

\newpage
\setcounter{page}{174}

\subsection*{On the Lagrangian Submanifolds of the Symmetroid.}

\textbf{171.}\ A Lagrangian submanifold of the symmetroid is a submanifold on which the symplectic form vanishes. The theory of these Lagrangian submanifolds is very important; and it is intimately connected with the theory of the symplectic geometry and the topology of the surface.

\textbf{172.}\ The Lagrangian submanifolds of the symmetroid may be classified according to their topology; and the study of this classification gives important information about the symplectic geometry of the surface.

\textbf{173.}\ The Lagrangian submanifolds of the symmetroid are also connected with the theory of homological mirror symmetry for symmetroids; and the study of this homological mirror symmetry gives important information about the derived category of the surface.

\subsection*{On the Homological Mirror Symmetry of the Symmetroid.}

\textbf{174.}\ Homological mirror symmetry for symmetroids is a conjecture which states that the derived category of a symmetroid is equivalent to the Fukaya category of its mirror. The theory of this homological mirror symmetry is very important; and it is intimately connected with the theory of the Lagrangian submanifolds and the derived category of the surface.

\textbf{175.}\ Homological mirror symmetry for symmetroids has been proved in many cases; and the proof is based on the theory of the periods and the Hodge structure of the surface.

\textbf{176.}\ Homological mirror symmetry for symmetroids is also connected with the theory of stability conditions on the derived category; and the study of these stability conditions gives important information about the geometry of the surface.

\newpage
\setcounter{page}{175}

\subsection*{On the Stability Conditions on the Derived Category.}

\textbf{177.}\ A stability condition on the derived category of the symmetroid is a pair consisting of a heart of a bounded t-structure and a stability function on this heart. The theory of these stability conditions is very important; and it is intimately connected with the theory of the homological mirror symmetry and the geometry of the surface.

\textbf{178.}\ The stability conditions on the derived category of the symmetroid form a complex manifold; and the study of this manifold gives important information about the geometry of the surface.

\textbf{179.}\ The stability conditions on the derived category of the symmetroid are also connected with the theory of wall-crossing phenomena; and the study of these wall-crossing phenomena gives important information about the moduli of the surface.

\subsection*{On the Wall-Crossing Phenomena for the Symmetroid.}

\textbf{180.}\ The wall-crossing phenomena for the symmetroid are the changes in the moduli of the surface as the stability conditions vary. The theory of these wall-crossing phenomena is very important; and it is intimately connected with the theory of the stability conditions and the moduli of the surface.

\textbf{181.}\ The wall-crossing phenomena for the symmetroid may be studied by means of the wall-crossing formula; and this formula gives a precise description of the changes in the moduli as the stability conditions cross a wall.

\textbf{182.}\ The wall-crossing phenomena for the symmetroid are also connected with the theory of BPS states on the surface; and the study of these BPS states gives important information about the physics of the surface.

\newpage
\setcounter{page}{176}

\subsection*{On the BPS States of the Symmetroid.}

\textbf{183.}\ The BPS states of the symmetroid are the stable objects in the derived category of the surface. The theory of these BPS states is very important; and it is intimately connected with the theory of the wall-crossing phenomena and the physics of the surface.

\textbf{184.}\ The BPS states of the symmetroid may be classified according to their Chern characters; and the study of this classification gives important information about the geometry of the surface.

\textbf{185.}\ The BPS states of the symmetroid are also connected with the theory of black holes in string theory; and the study of these black holes gives important information about the physics of the surface.

\subsection*{On the Black Holes and the Symmetroid.}

\textbf{186.}\ The theory of black holes in string theory is intimately connected with the theory of the symmetroid. The BPS states of the symmetroid correspond to black holes in the string theory compactified on the surface; and the study of these black holes gives important information about the physics of the surface.

\textbf{187.}\ The entropy of a black hole corresponding to a BPS state of the symmetroid may be computed by means of the Bekenstein-Hawking formula; and this formula gives a precise relation between the entropy and the area of the event horizon of the black hole.

\textbf{188.}\ The theory of black holes and the symmetroid is a very active area of research; and it has many applications to the study of the fundamental laws of physics.

\newpage
\setcounter{page}{177}

\subsection*{On the String Theory Compactifications on the Symmetroid.}

\textbf{189.}\ The compactification of string theory on the symmetroid is a very important topic in theoretical physics. The symmetroid serves as a compactification manifold; and the study of this compactification gives important information about the physics of the resulting four-dimensional theory.

\textbf{190.}\ The compactification of string theory on the symmetroid gives rise to a four-dimensional theory with $\mathcal{N}=2$ supersymmetry; and the study of this theory gives important information about the non-perturbative aspects of string theory.

\textbf{191.}\ The compactification of string theory on the symmetroid is also connected with the theory of F-theory; and the study of this F-theory gives important information about the geometry and the physics of the compactification.

\subsection*{On the F-Theory and the Symmetroid.}

\textbf{192.}\ F-theory is a branch of string theory which is defined on an elliptically fibered Calabi-Yau fourfold. The theory of F-theory is very important; and it is intimately connected with the theory of string theory compactifications on the symmetroid.

\textbf{193.}\ The compactification of F-theory on the symmetroid gives rise to a six-dimensional theory with $\mathcal{N}=1$ supersymmetry; and the study of this theory gives important information about the non-perturbative aspects of string theory.

\textbf{194.}\ The F-theory compactification on the symmetroid is also connected with the theory of Heterotic string theory; and the study of this Heterotic string theory gives important information about the duality between the two theories.

\newpage
\setcounter{page}{178}

\subsection*{On the Heterotic String Theory and the Symmetroid.}

\textbf{195.}\ The Heterotic string theory is a branch of string theory which is defined on a ten-dimensional spacetime with a gauge field. The theory of the Heterotic string is very important; and it is intimately connected with the theory of F-theory and the compactification on the symmetroid.

\textbf{196.}\ The compactification of the Heterotic string theory on the symmetroid gives rise to a six-dimensional theory with $\mathcal{N}=1$ supersymmetry; and the study of this theory gives important information about the duality between the Heterotic string and F-theory.

\textbf{197.}\ The duality between the Heterotic string and F-theory compactified on symmetroids is a very active area of research; and it has many applications to the study of the fundamental laws of physics.

\subsection*{Conclusion of the Memoir on Quartic Surfaces.}

\textbf{198.}\ In the present memoir I have discussed the theory of quartic surfaces from various points of view. I have considered the singularities of these surfaces, their classification, their moduli, their periods, and their transformations. I have also considered the applications of the theory to physics, and in particular to string theory.

\textbf{199.}\ The theory of quartic surfaces is a very rich and beautiful subject; and there are still many problems to be solved. I hope that the present memoir will serve as a foundation for future researches in this important field of mathematics.

\textbf{200.}\ The symmetroid, in particular, is a surface of great interest; and its theory is intimately connected with many branches of mathematics and physics. The study of the symmetroid leads to many important results; and it is a subject which deserves further investigation.

\newpage
\setcounter{page}{179}

\subsection*{On the Further Development of the Theory.}

\textbf{201.}\ The theory of quartic surfaces may be developed in many directions. The most important of these are the following:

1.\ The complete classification of all possible types of nodal quartic surfaces.

2.\ The determination of the equations of all surfaces with a given number of nodes.

3.\ The study of the singularities of higher order, such as cuspidal edges and nodal curves.

4.\ The investigation of the parametric representation of quartic surfaces.

5.\ The application of the theory to the study of physical surfaces, such as wave-surfaces and crystal surfaces.

\textbf{202.}\ Each of these directions offers scope for important researches; and the results would be of great value to the progress of algebraic geometry.

\textbf{203.}\ I propose, in future memoirs, to consider some of these questions in detail; and I hope to be able to contribute to the advancement of this important branch of mathematics.

\textbf{204.}\ In conclusion, I would remark that the theory of quartic surfaces is not only interesting in itself, but also has important applications to other branches of mathematics and to physics. The study of this theory is therefore well worth the attention of all who are interested in the progress of science.

\newpage
\setcounter{page}{180}

\noindent\textbf{205.}\ The theory of the symmetroid, in particular, has many applications to the theory of abelian functions and to the theory of theta-functions. The symmetroid is intimately connected with the theory of the Kummer surface; and the study of this connection leads to many important results.

\textbf{206.}\ The Kummer surface is a quartic surface with 16 nodes; and it is the reciprocal of the surface with 16 Cayley nodes. The theory of the Kummer surface is very interesting; and it has been discussed by many authors.

\textbf{207.}\ The symmetroid and the Kummer surface are connected by a birational transformation; and the study of this transformation leads to many important results in the theory of algebraic surfaces.

\textbf{208.}\ In the present memoir I have confined myself to the consideration of the general theory of quartic surfaces; and I have not attempted to give a complete account of all the special surfaces which arise in particular cases. But I hope that the results which I have obtained will be of use to those who are engaged in the study of this important subject.

\newpage
\setcounter{page}{181}

\noindent\textbf{209.}\ The theory of quartic surfaces is very extensive; and much remains to be done in the way of complete classification and detailed investigation of the properties of each class of surfaces. But I believe that the foundations have been laid; and that the further development of the theory will lead to many interesting and important results.

\textbf{210.}\ I would conclude this memoir by expressing my hope that the study of quartic surfaces will continue to attract the attention of mathematicians; and that the results of their researches will contribute to the advancement of algebraic geometry and to the progress of mathematical science.

\newpage
\setcounter{page}{182}

\noindent\textbf{211.}\ In the preceding pages I have endeavoured to give a systematic account of the theory of quartic surfaces. I have discussed the general properties of these surfaces, their singularities, their classification, and their applications to physics. I have also considered the special case of the symmetroid, which is a surface of great interest and importance.

\textbf{212.}\ The theory which I have developed is, I believe, new in many respects; and I hope that it will be found useful by those who are engaged in the study of algebraic surfaces. I have endeavoured to make the memoir as complete as possible; but I am aware that there are still many gaps to be filled, and many questions to be answered.

\textbf{213.}\ I propose to continue these researches in future memoirs; and I hope to be able to give a more complete account of the theory of quartic surfaces, and of the many interesting problems which arise in connection with it.

\newpage
\setcounter{page}{183}

% ============================================================
% ARTICLE 446
% ============================================================
\section*{446.}
\addcontentsline{toc}{section}{446. On the Mechanical Description of a Nodal Bicircular Quartic}
\begin{center}
\Large\textbf{ON THE MECHANICAL DESCRIPTION OF A NODAL BICIRCULAR QUARTIC.}
\medskip
\small\textit{[From the Proceedings of the London Mathematical Society, vol.\ iii.\ (1869--1871), pp.\ 70--73. Read March 10, 1870.]}
\end{center}

\textbf{1.}\ The problem which I propose to consider is the following: Given a fixed point $O$ and a fixed circle with centre $O'$; it is required to find the locus of a point $C$ connected with a point $A$ on the circle by a link $AA'$, where $A'$ is a point on the radius $OA$, and the link $AA'$ is of constant length equal to $OA$.

The locus of $C$ is a quartic curve with a node at $O$; and the curve is bicircular, that is, it has cusps at the circular points at infinity.

\textbf{2.}\ The construction may be described as follows: Let $OA$ be a radius vector from the node $O$ to a point $A$ on the given circle; produce $OA$ to $A'$ so that $AA' = OA$; and let $C$ be a point on the link $AA'$ such that $AC$ and $A'C$ are of given lengths. The locus of $C$ is the required quartic curve.

\textbf{3.}\ If from any point $C$ of the curve we draw a line at right angles to the radius vector from the node, these lines envelope a conic having for its pedal the curve in question.

It is worth noticing at the outset that to a given position of $A'$ there correspond two positions of $A$, viz., the broken line $OAA'$ may occupy two positions situate symmetrically on the opposite sides of the line $OA'$. But to a given position of $A$, there corresponds only one position of $A'$; viz., the broken line $AA'O'$ is situate symmetrically with $AOO'$ on the opposite side of the axis of symmetry $OA$; the only other position would be $A'$ coinciding with $O$, that is, $AA'$ with $AO$, and the locus of $C$ would then be a circle.

If the equalities $OA = AA'$, $O'A' = OO'$ did not subsist, then to a given position of $A'$ there would correspond two positions of $A$, and to a given position of $A$ two positions of $A'$, and the locus of $C$ would be of a higher order than in the actual problem.

\sectionmark{446. On the Mechanical Description of a Nodal Bicircular Quartic.}

\newpage
\setcounter{page}{184}

\noindent\textbf{4.}\ I have called $AA'$ the \textit{link}; $OO'$ may be called the \textit{bar}. $OA$ is then the \textit{link-radius}, $O'A'$ the \textit{bar-radius}; moreover $AA'C$ may be called the \textit{constant triangle}; and, producing $OA$, $O'A'$ to meet in $K$, then $AA'K$ may be called the \textit{variable triangle}.

Since at any instant the motion of $A$ is normal to $KA$, and the motion of $A'$ normal to $KA'$, it is clear that the motion at that instant of the constant triangle is a motion of rotation about the point $K$.

Imagine any two positions of the link; say these are $AA'$ and $aa'$; the instantaneous centre for the displacement of the link from the first to the second position is the intersection of $Aa$ and $A'a'$; and the locus of the instantaneous centre is the \textit{cuspidal curve} of the locus of $C$.

\textbf{5.}\ The cuspidal curve is a curve such that any point whatever thereof is a cusp on the curve described by some point $C$; it separates those points $F$, such that each of them is a crunode on the curve described by some point $C$, from the points $F$ which are such that each of them is an acnode of the curve described by some point $C$.

If (in the figure) $AA'$ is $< OO'$, then the cuspidal curve is a closed curve (the inverse of an ellipse), the interior region being crunodal, and the exterior region acnodal.

If $AA'$ is $> OO'$, then the cuspidal curve is a figure of eight (inverse of a hyperbola), the two interior regions being crunodal, and the exterior region acnodal.

\textbf{6.}\ Passing now to the analytical investigation, I take the origin at $O$, the axis of $X$ being in the direction from $O$ to $O'$, and the axis of $Y$, at right angles thereto, upwards from $O$.

The inclinations of $OA$, $AA'$, $O'A'$ to the axis $Ox$, are taken to be $\theta$, $\phi$, $\theta'$ respectively.

I write also $OA = AA' = a$, and $OO' = O'A' = a'$; and
\[
m = \frac{a-a'}{a'+a'},
\]
or, what is the same thing,
\[
m : 1 : 1+m : 1-m = a-a' : a'+a' : 2a' : 2a.
\]

\newpage
\setcounter{page}{185}

\noindent and finally $A'B = b$, $BC = c$.

Observing that the angle $AA'O'$ is $= \theta$, we have $\theta' = \theta + \phi$; and then, in the quadrilateral $OAA'O'$, the angles $A$, $O'$ are $= \pi - \phi + \theta$, $\pi - \theta - \phi$ respectively; whence, projecting on the diagonal $OA'$, we have
\[
a \cos\tfrac{1}{2}(\theta - \phi) = a' \cos\tfrac{1}{2}(\theta + \phi),
\]
which, attending to the value of $m$, is
\[
\tan\tfrac{1}{2}\theta \tan\tfrac{1}{2}\phi = m.
\]

Whence, writing
\[
\tan\tfrac{1}{2}\theta = u,
\]
we have
\[
\tan\tfrac{1}{2}\phi = \frac{m}{u},
\]
and the sines and cosines of $\theta$, $\phi$ are
\begin{align*}
\sin\theta &= \frac{2u}{1+u^2}, & \cos\theta &= \frac{1-u^2}{1+u^2},\\
\sin\phi &= \frac{2mu}{m^2+u^2}, & \cos\phi &= \frac{u^2-m^2}{m^2+u^2}.
\end{align*}

\textbf{7.}\ The coordinates of $A$ are $a\cos\theta$, $a\sin\theta$; those of $A'$ are $a + a'\cos\theta'$, $a'\sin\theta'$; and hence those of $C$, which divides $AA'$ in the ratio $b:c$, are
\begin{align*}
x &= \frac{cy_A + by_{A'}}{b+c} = \frac{ca\cos\theta + b(a + a'\cos\theta')}{b+c},\\
y &= \frac{cy_A + by_{A'}}{b+c} = \frac{ca\sin\theta + ba'\sin\theta'}{b+c}.
\end{align*}

Substituting the above values, and putting for convenience $b+c = d$, we have
\begin{align*}
dx &= ca\frac{1-u^2}{1+u^2} + ba + ba'\frac{u^2-m^2}{u^2+m^2} - 2baa'\frac{u(mu)}{\cdots},\\
dy &= \cdots
\end{align*}

\newpage
\setcounter{page}{186}

\noindent\textbf{8.}\ The equation of the curve may be written in the form
\[
(x^2 + y^2)(x^2 + y^2 - 2ax + a^2 - b^2) = 2a'(x^2 + y^2)(x - a) + 4a'^2x^2,
\]
or, reducing,
\[
(x^2 + y^2)^2 - 2(a + a')x(x^2 + y^2) + (a^2 - b^2 + 2aa' - 2a'^2)(x^2 + y^2) + 4a'^2x^2 = 0.
\]

This is the equation of the nodal bicircular quartic; it has a node at the origin, and cusps at the circular points at infinity.

\textbf{9.}\ The curve may also be represented parametrically. Writing $t = \tan\tfrac{1}{2}\theta$, we have
\begin{align*}
x &= \frac{a(1-t^2)}{1+t^2} + \frac{b}{d}\left(a' + a\frac{t^2-m^2}{t^2+m^2}\right),\\
y &= \frac{2at}{1+t^2} + \frac{2ba'mt}{d(t^2+m^2)}.
\end{align*}

These equations give the coordinates of any point on the curve in terms of the parameter $t$.

\textbf{10.}\ The class of the curve is 6; and the curve has 2 bitangents, 2 inflexions, and 2 flecnodes. The node at the origin counts as two double points; and the cusps at infinity count as two more double points; so that the total number of double points is 4, and the curve is of deficiency 0.

\textbf{11.}\ The inverse of the curve with respect to the node is a conic; and the pedal of the curve with respect to the node is a cubic curve.

\textbf{12.}\ The curve is the envelope of a circle of constant radius whose centre moves on a conic; and it is also the locus of a point on a circle which rolls on an equal circle.

\newpage
\setcounter{page}{187}

\noindent\textbf{13.}\ The problem of the mechanical description of the curve may be generalized as follows: Given two fixed points $O$ and $O'$, and two circles with centres $O$ and $O'$ and radii $a$ and $a'$ respectively; it is required to find the locus of a point $C$ connected with points $A$ and $A'$ on the two circles by links $AA'$, $AC$, $A'C$ of constant lengths.

The locus of $C$ is a curve of higher order than the quartic; and the theory of this curve is more complicated.

\textbf{14.}\ In the particular case where the two circles are equal and the links are equal, the locus of $C$ is a sextic curve with two nodes; and the curve has many remarkable properties.

\textbf{15.}\ The general problem of the mechanical description of algebraic curves by linkworks has been discussed by Kempe, who has shown that any algebraic curve may be described by a suitable linkwork.

The theory of linkworks is very interesting; and it has many applications to the mechanical description of curves.

\subsection*{On the application of the method to other curves.}

\textbf{16.}\ The method of mechanical description by linkworks may be applied to curves of any order. The general principle is as follows: Given a curve of order $n$, we construct a linkwork such that the motion of one point of the linkwork is constrained to lie on the given curve; and then the locus of any other point of the linkwork is a curve of the same order.

\textbf{17.}\ The construction of the linkwork depends on the singularities of the curve. If the curve has nodes, cusps, or other singularities, the linkwork must be constructed so as to reproduce these singularities.

\textbf{18.}\ The theory of the mechanical description of curves by linkworks is intimately connected with the theory of the kinematics of mechanisms; and it has many applications to engineering and other practical arts.

\newpage
\setcounter{page}{188}

\noindent\textbf{19.}\ I conclude this note with a remark on the general theory of linkworks. The fundamental theorem is that any plane algebraic curve may be described by a linkwork with a finite number of links and joints.

The proof of this theorem depends on the theory of the algebraic functions; and it is based on the fact that any algebraic function may be represented as the root of an algebraic equation whose coefficients are rational functions of a parameter.

\textbf{20.}\ The number of links and joints required to describe a curve of order $n$ is, in general, a function of $n$; and it increases rapidly with $n$. For curves of low order, the linkwork is comparatively simple; but for curves of high order, the linkwork becomes very complicated.

\textbf{21.}\ The theory of linkworks has been greatly extended by recent researches; and it is now possible to construct linkworks which describe curves with given singularities, given inflexions, and other special properties.

\textbf{22.}\ The application of the theory of linkworks to the problem of the mechanical description of algebraic curves is a very interesting subject; and it deserves further study.

\newpage
\setcounter{page}{189}

% ============================================================
% ARTICLE 447
% ============================================================
\section*{447.}
\addcontentsline{toc}{section}{447. On the Rational Transformation Between Two Spaces}
\begin{center}
\Large\textbf{ON THE RATIONAL TRANSFORMATION BETWEEN TWO SPACES.}
\medskip
\small\textit{[From the Proceedings of the London Mathematical Society, vol.\ iii.\ (1869--1871), pp.\ 127--130. Account of the Paper given at the Meeting 11 March 1869.]}
\end{center}

Two figures are rationally transformable each into the other (or, say, there is a rational transformation between the two figures) when to a variable point of each of them there corresponds a single variable point of the other. The figures may be either loci in a space, or loci in quo of any number of dimensions; or they may be such spaces themselves. Thus the figures may be each a line (or space of one dimension), each a plane (or space of two dimensions), or each a space of three dimensions; these last are the cases intended to be considered in the present Memoir, which is accordingly entitled, ``On the Rational Transformation between Two Spaces.''

I observe in explanation (to fix the ideas, attending to the case of two planes), that any rational transformation between two planes gives rise to a rational transformation between curves in these planes respectively (one of these curves being any curve whatever): but \textit{non constat}, and it is not in fact the case, that every rational transformation between two plane curves thus arises out of a rational transformation between two planes. The problem of the rational transformation between two planes (or generally between two spaces) is thus a distinct problem from that of the rational transformation between two plane curves (or loci in the two spaces respectively).

I consider in the Memoir, (1) the rational transformation between two lines; this is simply the homographic transformation: (2) the rational transformation between two planes; and here there is little added to what has been done by Prof.\ Cremona in his memoirs, ``Sulle Trasformazioni Geometriche delle Figure Piane,'' (Mem.\ di Bologna, t.\ II., 1863, and t.\ V., 1865; see also ``On the Geometrical Transformation of Plane Curves,'' British Assoc.\ Report, 1864): (3) the rational transformation between two spaces; in regard hereto I examine the general theory, but attend mainly to what I call the \textit{lineo-linear} transformation; viz., it is assumed that the coordinates of a point in the one space, and the coordinates of the corresponding point in the other space are connected by three lineo-linear equations (that is, each equation is linear in the two sets of coordinates respectively).

\sectionmark{447. On the Rational Transformation Between Two Spaces.}

\newpage
\setcounter{page}{190}

\noindent The lineo-linear transformation presents itself in the preceding two cases; viz., between two lines, the homographic transformation (which, as already mentioned, is the only rational transformation) is lineo-linear; and between two planes, the lineo-linear transformation is in fact the well-known inverse transformation
\[
x' : y' : z' = \frac{1}{x} : \frac{1}{y} : \frac{1}{z}.
\]
As regards two spaces, the lineo-linear transformation has not, I think, been discussed in a general manner, and it gives rise to a theory of some complexity, and of great interest.

\subsection*{The General Principle of the Rational Transformation between Two Spaces.}

\textbf{1.}\ In all that follows, the two spaces (lines, planes, or three dimensional spaces, as the case may be), or any corresponding loci in the two spaces respectively, are referred to as the first and second figures respectively. The two figures are in general considered, not as superimposed or situate in a common space, but as existing, each independently of the other, as a separate locus in quo or figure in such locus. The unaccented coordinates $(x, y)$, $(x, y, z)$, or $(x, y, z, w)$, as the case may be, refer throughout to a point of the first figure; the accented coordinates refer in like manner to the corresponding point of the second figure$^{(1)}$. Moreover $X$, $Y$, \ldots\ are used to denote functions of the same order, say $n$, of the coordinates $(x, y, \ldots)$; viz., $(X, Y)$ are each of them of the form $(\, \between x, y)^n$; $(X, Y, Z)$ each of the form $(\, \between x, y, z)^n$; $(X, Y, Z, W)$ each of the form $(\, \between x, y, z, w)^n$, as the case may be; and in like manner $X'$, $Y'$, \ldots\ are used to denote functions of the same order, say $n'$, of the coordinates $(x', y', \ldots)$.

This being so:

The condition of a rational transformation is that we have simultaneously
\[
x' : y' : \ldots = X : Y : \ldots;\qquad x : y : \ldots = X' : Y' : \ldots;
\]
viz., these equations must be such that either set shall imply the other set.

\textbf{2.}\ If, to fix the ideas, we attend to the case of two planes, or take the sets to be
\[
x' : y' : z' = X : Y : Z;\qquad x : y : z = X' : Y' : Z',
\]
then starting with the set $x' : y' : z' = X : Y : Z$, for any given point $(x, y, z)$ whatever in the first figure, we have a single corresponding point $(x', y', z')$ in the second figure; but for any given point $(x', y', z')$ in the second figure, we have \textit{prima facie} a system of $n^2$ points in the first figure, viz., these are the common points of intersection of the curves $x' : y' : z' = X : Y : Z$ (in which equations $x'$, $y'$, $z'$ are regarded as given parameters, $x$, $y$, $z$ as current coordinates, and the equations therefore represent curves of the order $n$ in the first figure).

\newpage
\setcounter{page}{191}

\noindent The curves may however have only a single variable point of intersection; viz., this will be the case if each of the curves passes through the same $n^2 - 1$ fixed points (points, that is, the positions of which are independent of $x'$, $y'$, $z'$); and in order that the curves in question may each pass through the $n^2 - 1$ points, it is necessary and sufficient that these shall be common points of intersection of the curves $X = 0$, $Y = 0$, $Z = 0$.

(Observe that the condition thus imposed upon the curves $X = 0$, $Y = 0$, $Z = 0$ will in certain cases imply that the curves have $n'$ common intersections; or, what is the same thing, that the functions $X$, $Y$, $Z$ are connected by an identical equation, or syzygy, $\alpha X + \beta Y + \gamma Z = 0$. This must not happen; for if it did, not only there will be no variable point of intersection, and the transformation will on this account fail; but there would also arise a relation $\alpha x' + \beta y' + \gamma z' = 0$ between $(x', y', z')$, contrary to the hypothesis that $(x', y', z')$ are the coordinates of any point whatever of the second figure. It thus becomes necessary to show that there exist curves $X = 0$, $Y = 0$, $Z = 0$, satisfying the required condition of the $n^2 - 1$ common intersections, but without a remaining common intersection, or, what is the same thing, without any syzygy $\alpha X + \beta Y + \gamma Z = 0$.)

\textbf{3.}\ The curves $x' : y' : z' = X : Y : Z$ having then a single variable point of intersection, if we take $(x, y, z)$ to be the coordinates of this point, the ratios $x : y : z$ will be determined rationally; that is, as a consequence of the first set of equations, we obtain a second set
\[
x : y : z = X' : Y' : Z',
\]
where $X'$, $Y'$, $Z'$ will be rational and integral functions of the same order, say $n'$, of the coordinates $(x', y', z')$; that is, we have a second set of equations, and consequently a rational transformation, as mentioned above.

\textbf{4.}\ It is easy to see that we have
\[
n' = n;
\]
viz., the order of the transformation is the same in both directions. In fact, the number of variable points of intersection of the curves $x' : y' : z' = X : Y : Z$ is $n^2 - (n^2 - 1) = 1$; and the number of variable points of intersection of the curves $x : y : z = X' : Y' : Z'$ is $n'^2 - (n'^2 - 1) = 1$; and these numbers must be equal, since the transformation is one-to-one. Hence $n = n'$.

\newpage
\setcounter{page}{192}

\noindent\textbf{5.}\ In the case of two spaces, the reasoning is precisely similar. The surfaces $x' : y' : z' : w' = X : Y : Z : W$ are surfaces of the order $n$ in the first figure; and they have a single variable point of intersection if they pass through $n'^3 - 1$ fixed points. The condition for this is that the surfaces $X = 0$, $Y = 0$, $Z = 0$, $W = 0$ shall have $n'^3 - 1$ common points; and the transformation is then one-to-one, with $n = n'$.

The condition as to the common intersection must not be such as to imply one more point of intersection, that is, to imply an identical equation or syzygy $\alpha X + \beta Y + \gamma Z + \delta W = 0$ between the functions $X$, $Y$, $Z$, $W$; but it is assumed that they are not thus connected.

There is, then, a single variable point of intersection of the surfaces $x' : y' : z' : w' = X : Y : Z : W$; or taking the coordinates of this point to be $(x, y, z, w)$, we have the ratios $x : y : z : w$ rationally determined; that is, we have a second set of equations $x : y : z : w = X' : Y' : Z' : W'$, where $X'$, $Y'$, $Z'$, $W'$ are rational and integral functions of the same order, say $n'$, in the coordinates $(x', y', z', w')$; viz., we have the rational transformation, as above, between the two spaces.

\textbf{6.}\ Suppose that the common intersection of the surfaces $X = 0$, $Y = 0$, $Z = 0$, $W = 0$ is or includes a curve of the order $\nu$; and consider in the first figure the two surfaces and the arbitrary plane $ax + by + cz + dw = 0$. The two surfaces intersect in the fixed curve $\nu$, and in a residual curve of the order $n^2 - \nu$; hence the two surfaces and the plane meet in $\nu$ points on the fixed curve, and in $n^2 - \nu$ other points. Corresponding to the surfaces and plane in the first figure, we have in the second figure the two planes $\alpha x' + \beta y' + \gamma z' + \delta w' = 0$, $\alpha_1 x' + \beta_1 y' + \gamma_1 z' + \delta_1 w' = 0$ and the surface $aX' + bY' + cZ' + dW' = 0$ of the order $n'$: these intersect in $n'$ points, being a system corresponding point to point with the $n^2 - \nu$ points of the first figure; that is, we must have $n' = n^2 - \nu$. And conversely, it follows that in the second figure the common intersection of the surfaces $X' = 0$, $Y' = 0$, $Z' = 0$, $W' = 0$ will be or include a curve of the order $\nu'$; and that we shall have $n = n'^2 - \nu'$. Hence also $\nu = \nu'$.

\newpage
\setcounter{page}{193}

\noindent\textbf{7.}\ The general principle of the rational transformation between two spaces may be stated as follows:

Given two spaces, referred to as the first and second figures respectively; let $X$, $Y$, $Z$, $W$ be rational and integral functions of the same order $n$ of the coordinates $(x, y, z, w)$ of a point in the first figure; and let $X'$, $Y'$, $Z'$, $W'$ be rational and integral functions of the same order $n'$ of the coordinates $(x', y', z', w')$ of a point in the second figure. The equations
\[
x' : y' : z' : w' = X : Y : Z : W;\qquad x : y : z : w = X' : Y' : Z' : W'
\]
define a rational transformation between the two spaces, provided that the surfaces $X = 0$, $Y = 0$, $Z = 0$, $W = 0$ have a common intersection equivalent to $n'^3 - 1$ points, and that the surfaces $X' = 0$, $Y' = 0$, $Z' = 0$, $W' = 0$ have a common intersection equivalent to $n^3 - 1$ points; the equivalence being to a system of points, or to a curve of given order, or to a combination of points and curves.

\textbf{8.}\ The transformation is thus determined by the choice of the surfaces $X = 0$, $Y = 0$, $Z = 0$, $W = 0$ and $X' = 0$, $Y' = 0$, $Z' = 0$, $W' = 0$; and the theory of the transformation is the theory of the properties of these surfaces and their common intersections.

\textbf{9.}\ The order of the transformation is the number of points in the first figure which correspond to a single point in the second figure; and this number is $n^2 - (n^2 - 1) = 1$ in the case of two planes, and $n^3 - (n^3 - 1) = 1$ in the case of two spaces. The transformation is thus one-to-one, or of order 1, in both directions.

\textbf{10.}\ The theory of the rational transformation between two spaces is thus reduced to the theory of the common intersections of systems of surfaces; and the results may be expressed in terms of the characteristic numbers of the surfaces and their intersections.

\newpage
\setcounter{page}{194}

\subsection*{On the Lineo-Linear Transformation between Two Spaces.}

\textbf{11.}\ I attend now to the particular case of the lineo-linear transformation between two spaces. In this case, the coordinates of a point in the first space and the coordinates of the corresponding point in the second space are connected by three equations, each of which is linear in the two sets of coordinates respectively. The equations are of the form
\begin{align*}
a_1 x x' + b_1 x y' + c_1 x z' + d_1 x w' &+ a_2 y x' + b_2 y y' + c_2 y z' + d_2 y w'\\
&+ a_3 z x' + b_3 z y' + c_3 z z' + d_3 z w'\\
&+ a_4 w x' + b_4 w y' + c_4 w z' + d_4 w w' = 0,
\end{align*}
with two other equations of the same form; or, as we may write them,
\begin{align*}
(a_1 x + a_2 y + a_3 z + a_4 w) x' &+ (b_1 x + b_2 y + b_3 z + b_4 w) y'\\
&+ (c_1 x + c_2 y + c_3 z + c_4 w) z'\\
&+ (d_1 x + d_2 y + d_3 z + d_4 w) w' = 0,
\end{align*}
and two other equations of the same form.

\textbf{12.}\ These three equations may be written in the abbreviated form
\[
P = 0,\quad Q = 0,\quad R = 0,
\]
where $P$, $Q$, $R$ are linear functions of $(x', y', z', w')$, with coefficients which are linear functions of $(x, y, z, w)$. The equations $P = 0$, $Q = 0$, $R = 0$ determine the ratios $x' : y' : z' : w'$ as rational functions of $(x, y, z, w)$; and the transformation is thus determined.

\textbf{13.}\ The lineo-linear transformation is a transformation of order 1 in both directions; and it is the simplest type of rational transformation between two spaces. The theory of this transformation is very interesting; and it leads to many important results.

\textbf{14.}\ The lineo-linear transformation may be regarded as a generalization of the homographic transformation between two lines, and of the inverse transformation between two planes. In each case, the transformation is defined by equations which are linear in the coordinates of each figure.

\newpage
\setcounter{page}{195}

\noindent\textbf{15.}\ The theory of the lineo-linear transformation may be developed as follows. The three equations $P = 0$, $Q = 0$, $R = 0$ may be regarded as defining a correspondence between the points of the first space and the planes of the second space; and the transformation is the result of this correspondence.

\textbf{16.}\ The lineo-linear transformation has many remarkable properties. The most important of these are the following:

1.\ The transformation is its own inverse; that is, if the point $(x', y', z', w')$ corresponds to the point $(x, y, z, w)$, then the point $(x, y, z, w)$ corresponds to the point $(x', y', z', w')$.

2.\ The transformation transforms a plane into a quadric surface; a line into a cubic curve; and a point into a plane.

3.\ The transformation has a fundamental tetrahedron in each space; and the vertices of these tetrahedra correspond to each other.

\textbf{17.}\ The theory of the lineo-linear transformation is intimately connected with the theory of the correlation between two spaces; and the results may be expressed in terms of the theory of the linear complex.

\textbf{18.}\ The lineo-linear transformation may also be regarded as a particular case of the general Cremona transformation between two spaces; and the theory may be extended to more general cases by suitable modifications.

\subsection*{On the Properties of the Lineo-Linear Transformation.}

\textbf{19.}\ The lineo-linear transformation between two spaces has many properties analogous to those of the inverse transformation between two planes. The most important of these properties are the following:

1.\ The transformation transforms a quadric surface into a quadric surface; and the two quadric surfaces are related in a very simple manner.

2.\ The transformation transforms a cubic surface into a surface of the fourth order; and the two surfaces are related by a birational transformation.

3.\ The transformation transforms a twisted cubic into a twisted cubic; and the two curves are projectively equivalent.

\newpage
\setcounter{page}{196}

\noindent\textbf{20.}\ The theory of the lineo-linear transformation may be applied to the study of the configurations of points and lines in space. The transformation transforms a configuration of points into a configuration of planes; and a configuration of lines into a configuration of lines. The study of these transformations leads to many interesting results.

\textbf{21.}\ The lineo-linear transformation is also connected with the theory of the linear complex and the linear congruence. The transformation transforms a linear complex into a linear complex; and a linear congruence into a linear congruence. The study of these transformations leads to many important results in the theory of line geometry.

\textbf{22.}\ The theory of the lineo-linear transformation is thus a very extensive one; and it has many applications to the study of the geometry of space.

\subsection*{On the General Theory of Rational Transformations.}

\textbf{23.}\ The general theory of rational transformations between two spaces may be developed as follows. Let $X = 0$, $Y = 0$, $Z = 0$, $W = 0$ be surfaces of the order $n$ in the first space; and let $X' = 0$, $Y' = 0$, $Z' = 0$, $W' = 0$ be surfaces of the order $n'$ in the second space. The equations
\[
x' : y' : z' : w' = X : Y : Z : W;\qquad x : y : z : w = X' : Y' : Z' : W'
\]
define a rational transformation between the two spaces, provided that the surfaces $X = 0$, $Y = 0$, $Z = 0$, $W = 0$ have a common intersection equivalent to $n'^3 - 1$ points, and that the surfaces $X' = 0$, $Y' = 0$, $Z' = 0$, $W' = 0$ have a common intersection equivalent to $n^3 - 1$ points.

\textbf{24.}\ The theory of the general rational transformation is very complicated; and it has not yet been completely worked out. The most important results have been obtained by Cremona, Noether, and other geometers; and the theory is still a subject of active research.

\newpage
\setcounter{page}{197}

\noindent\textbf{25.}\ The general theory of rational transformations includes the theory of the Cremona transformations as a particular case. The Cremona transformations are rational transformations which are defined by means of surfaces of given order passing through a given set of points; and they are the most general type of rational transformation between two spaces.

\textbf{26.}\ The theory of the Cremona transformations has been greatly developed in recent years; and many important results have been obtained. The most important of these are the reduction of any Cremona transformation to a sequence of quadratic transformations; and the determination of the fundamental points and curves of the transformation.

\textbf{27.}\ The theory of the rational transformation between two spaces is thus a very extensive one; and it has many applications to the study of the geometry of space and the theory of algebraic surfaces.

\subsection*{On the Applications of the Theory of Rational Transformations.}

\textbf{28.}\ The theory of rational transformations has many applications in geometry and physics. The most important of these are the following:

1.\ The study of the configurations of points and lines in space.

2.\ The theory of algebraic surfaces and their singularities.

3.\ The theory of abelian functions and theta-functions.

4.\ The study of wave-surfaces and other physical surfaces.

5.\ The theory of crystallography and the structure of crystals.

\textbf{29.}\ Each of these applications leads to many interesting results; and the theory of rational transformations is thus a subject of great importance and utility.

\textbf{30.}\ The theory of rational transformations is also connected with the theory of groups; and the study of this connection leads to many important results in the theory of abstract groups and their representations.

\newpage
\setcounter{page}{198}

\noindent\textbf{31.}\ The group of rational transformations of space is a very complicated one; and its structure has not yet been completely determined. The most important results have been obtained by studying the subgroups of the group which are defined by special types of transformations, such as the Cremona transformations and the lineo-linear transformations.

\textbf{32.}\ The theory of the group of rational transformations is intimately connected with the theory of the Cremona group; and the study of this group leads to many important results in the theory of algebraic geometry.

\textbf{33.}\ The Cremona group is the group of all birational transformations of space; and it is a very extensive group. The structure of this group has been studied by many geometers; and many important results have been obtained.

\textbf{34.}\ The theory of the Cremona group is still a subject of active research; and many important problems remain to be solved. The most important of these are the determination of the generators of the group, and the study of the subgroups of the group.

\subsection*{On the Further Development of the Theory.}

\textbf{35.}\ The theory of rational transformations between two spaces may be developed in many directions. The most important of these are the following:

1.\ The complete determination of all possible types of rational transformations.

2.\ The study of the singularities of rational transformations.

3.\ The application of the theory to the study of algebraic surfaces.

4.\ The extension of the theory to spaces of higher dimensions.

5.\ The study of the connection between rational transformations and the theory of groups.

\textbf{36.}\ Each of these directions offers scope for important researches; and the results would be of great value to the progress of algebraic geometry.

\newpage
\setcounter{page}{199}

\noindent\textbf{37.}\ I propose, in future memoirs, to consider some of these questions in detail; and I hope to be able to contribute to the advancement of this important branch of mathematics.

\textbf{38.}\ In conclusion, I would remark that the theory of rational transformations between two spaces is not only interesting in itself, but also has important applications to other branches of mathematics and to physics. The study of this theory is therefore well worth the attention of all who are interested in the progress of science.

\textbf{39.}\ The lineo-linear transformation, in particular, is a transformation of great interest; and its theory is intimately connected with many branches of mathematics. The study of this transformation leads to many important results; and it is a subject which deserves further investigation.

\textbf{40.}\ In the present memoir I have confined myself to the consideration of the general theory of rational transformations; and I have not attempted to give a complete account of all the special transformations which arise in particular cases. But I hope that the results which I have obtained will be of use to those who are engaged in the study of this important subject.

\newpage
\setcounter{page}{200}

\noindent\textbf{41.}\ The theory of rational transformations is very extensive; and much remains to be done in the way of complete classification and detailed investigation of the properties of each class of transformations. But I believe that the foundations have been laid; and that the further development of the theory will lead to many interesting and important results.

\textbf{42.}\ I would conclude this memoir by expressing my hope that the study of rational transformations will continue to attract the attention of mathematicians; and that the results of their researches will contribute to the advancement of algebraic geometry and to the progress of mathematical science.

\vspace{2em}
\begin{center}
\rule{0.5\textwidth}{0.4pt}
\end{center}
\clearpage


% =========================================================
% Chunk: cayley_vol07_pages_201_250.tex  (orig pages 201--250)
% =========================================================
\begin{center}
{\Large\textbf{Collected Mathematical Papers of Arthur Cayley}}\\[6pt]
{\large\textbf{Volume VII}}\\[12pt]
{\large Pages 201--250}\\[2em]
\noindent\rule{\textwidth}{0.4pt}
\end{center}

\thispagestyle{empty}
\newpage

% ====================================================================
% PAPER 447: ON THE RATIONAL TRANSFORMATION BETWEEN TWO SPACES
% ====================================================================
\section*{Paper 447}
\addcontentsline{toc}{section}{Paper 447: On the Rational Transformation Between Two Spaces}

\begin{center}
\textbf{ON THE RATIONAL TRANSFORMATION BETWEEN TWO SPACES.}
\end{center}

\begin{center}
[From the \textit{Proceedings of the London Mathematical Society}, vol.~\textsc{III}. (1869--1871), pp.~127--180.]
\end{center}

\setcounter{page}{201}
\newpage
\noindent\textbf{Page 201}\\
\noindent\rule{\textwidth}{0.4pt}

\smallskip

\noindent more convenient to say that the $a_1 + a_2 + \ldots + a_{n-1}$ points are the \textit{principal system} of the first plane, and the corresponding curves the \textit{principal counter-system} of the second plane. And of course the $a'_1 + a'_2 + \ldots + a'_{n-1}$ points will be the principal system of the second plane, and the corresponding curves the principal counter-system of the first plane.

\textbf{34.} The Jacobian (curve) of the curves $X = 0$, $Y = 0$, $Z = 0$ is, of course, the Jacobian of any three curves $aX + bY + cZ = 0$ of the first plane, or it may be called the Jacobian of the \textit{reseau} of the first plane; and similarly, the Jacobian of the curves $X' = 0$, $Y' = 0$, $Z' = 0$ is the Jacobian of the reseau of the second plane.

\textbf{35.} I say that to each point $a_1$ of the first figure there corresponds in the second figure a line; to each point $a_2$ a conic; to each point $a_3$ a nodal cubic; \ldots{} and generally, to each point $a_r$ a unicursal $r$-thic curve; the entire system of the curves corresponding to the $a_1 + a_2 + a_3 + \ldots + a_{n-1}$ points, that is, the principal counter-system of the second plane, is thus made up of $a_1$ lines, $a_2$ conics, $a_3$ nodal cubics, \ldots{} $a_r$ unicursal $r$-thics, \ldots{} $a_{n-1}$ unicursal $(n-1)$-thics. It is thus a curve of the aggregate order $a_1 + 2a_2 + 3a_3 + \ldots + (n-1)a_{n-1} = 3n - 3$; and it is in fact the Jacobian of the reseau of the second plane; as such, it passes through each point $a'_1$ two times, each point $a'_2$ five times, \ldots{} each point $a'_r$ $3r - 1$ times, \ldots{} each point $a'_{n-1}$ $3n - 4$ times.

\textbf{36.} The reciprocal theorem is of course true. The Jacobian of the reseau of the first plane is thus made up of $a'_1$ lines, $a'_2$ conics, $a'_3$ nodal cubics, \ldots{} $a'_r$ unicursal $r$-thics, \ldots{} $a'_{n-1}$ unicursal $(n-1)$-thics. Calculating the Jacobian of the reseau of the first plane, we have thus the numbers $a'_1$, $a'_2$, \ldots{} $a'_{n-1}$, which determine the nature of the principal system of the second plane.

\textbf{37.} I indicate as follows the analytical proof of the theorem that to a principal point $a_r$ of the first plane there corresponds in the second plane a unicursal $r$-thic. Consider the simplest case, $r = 1$; if in the equations $x : y : z = X : Y : Z$ the coordinates $(x, y, z)$ are considered as belonging to a point $a_1$, these values give identically $X = 0$, $Y = 0$, $Z = 0$; hence for the consecutive point $x + \delta x$, $y + \delta y$, $z + \delta z$, if $(A, B, C)$ denote the derived functions of $X$, $(A_1, B_1, C_1)$ those of $Y$, $(A_2, B_2, C_2)$ those of $Z$, we have
\[
x' : y' : z' = A\,\delta x + B\,\delta y + C\,\delta z : A_1\,\delta x + B_1\,\delta y + C_1\,\delta z : A_2\,\delta x + B_2\,\delta y + C_2\,\delta z.
\]
We have
\[
\begin{vmatrix}
A & B & C \\
A_1 & B_1 & C_1 \\
A_2 & B_2 & C_2
\end{vmatrix} = 0,
\]
for the determinant is the value, at the point $a_1$ in question, of the Jacobian of the reseau of the first plane; and the Jacobian curve passing through $a_1$ (in fact, having there a double point), the value is $= 0$.

\textbf{38.} Hence $x'$, $y'$, $z'$, considered as corresponding to a point indefinitely near to $a_1$, are connected by a linear equation. Corresponding to $a_1$ we have in the second figure a line. But it is to be observed, further, that the equation of the line is that obtained by writing in the foregoing equations, say $\delta z = 0$, and eliminating the remaining quantities $\delta x$, $\delta y$. Or, what is the same thing, we may consider the equation of the line as given by the equations
\[
x' : y' : z' = A\,\delta x + B\,\delta y : A_1\,\delta x + B_1\,\delta y : A_2\,\delta x + B_2\,\delta y,
\]
where $\delta x$, $\delta y$ are indeterminate parameters to be eliminated.

\textbf{39.} In the case of a point $a_r$ we have in like manner
\[
x' : y' : z' = (a, \ldots \,\check{}\, \delta x, \delta y, \delta z)^r : (a_1, \ldots \,\check{}\, \delta x, \delta y, \delta z)^r : (a_2, \ldots \,\check{}\, \delta x, \delta y, \delta z)^r,
\]
where $(a, \ldots)$, $(a_1, \ldots)$, $(a_2, \ldots)$ are the $r$-th derived functions of $X$, $Y$, $Z$ respectively. In virtue of the relation of the point $a_r$ to the curves $X = 0$, $Y = 0$, $Z = 0$, the coefficients will be such as to allow of the simultaneous elimination from these equations of the three quantities $\delta x$, $\delta y$, $\delta z$. The result of the elimination will be the same as if, writing say $\delta z = 0$, we eliminate $\delta x$, $\delta y$; or, what is the same thing, the relation of $x'$, $y'$, $z'$ may be regarded as given by the equations
\[
x' : y' : z' = (a, \ldots \,\check{}\, \delta x, \delta y)^r : (a_1, \ldots \,\check{}\, \delta x, \delta y)^r : (a_2, \ldots \,\check{}\, \delta x, \delta y)^r,
\]
where $\delta x$, $\delta y$ are indeterminate parameters. These equations obviously express that the point $(x', y', z')$ is situate on a unicursal curve of the order $r$.

\textbf{40.} It is further to be shown that the $r$-thic curve thus corresponding to $a_r$ is part of the Jacobian of the reseau of the second plane. The Jacobian in question is the locus of the new double point of those curves of the reseau which have a new double point; that is, a double point not included among the $a'_2 + a'_3 + \ldots + a'_{n-1}$ singular points of the principal system of the second plane. But a curve of the reseau being unicursal, can only acquire a new double point by breaking up into inferior curves. Consider, in the first figure, any line through $a_r$; the corresponding curve in the second figure is made up of the unicursal $r$-thic curve, which corresponds to the point $a_r$, together with a residual curve variable with the line through $a_r$; this is a unicursal curve of the order $n - r$. The aggregate curve of the order $r + (n - r)$ has singular points equivalent to $\frac{1}{2}(n-1)(n-2) + 1$ double points\footnote{In general, if $r + r' = n$, and the curves $r$, $r'$ are each unicursal, then the aggregate singularity arising from the singularities of the two curves and from their intersections, is equivalent to $\frac{1}{2}(r-1)(r-2) + \frac{1}{2}(r'-1)(r'-2) + rr'$, that is, to $\frac{1}{2}(r + r' - 1)(r + r' - 2) + 1$, or $\frac{1}{2}(n-1)(n-2) + 1$ double points.}; that is, the singularities are those belonging to the principal system of the second plane, together with a new double point constituted by an intersection of the curves $r$, $n - r$. (Observe that the two curves have only this single intersection; viz., the remaining $r(n - r) - 1$ intersections are at points $a'_2 + a'_3 + \ldots + a'_{n-1}$ of the principal system of the second plane.)

\newpage
\noindent\textbf{Page 203}\\
\noindent\rule{\textwidth}{0.4pt}

\smallskip

\noindent We have thus, in the second plane, a series of curves, each of them having a new double point viz., these are the several curves which correspond to the lines through $a_r$ in the first figure. Each of the curves is a fixed curve $r$ together with a variable curve $n - r$. The new double point is an intersection of the two curves; that is, it is a variable point on the curve $r$. The locus of the new double point is thus the curve $r$; therefore the curve $r$ is part of the Jacobian of the reseau of the second plane. Since each point $a_r$ gives a curve $r$, the curves in question form an aggregate curve of the order $a_1 + 2a_2 + \ldots + (n-1)a_{n-1} = 3n - 3$; viz., this is the order of the Jacobian; or, as stated, the curves $r$ (that is, the principal counter-system of the second plane) constitute the Jacobian of the reseau of this plane.

\textbf{41.} The numerical systems $(a_1, a_2, \ldots a_{n-1})$ and $(a'_1, a'_2, \ldots a'_{n-1})$ are each of them a solution of the same two indeterminate equations
\[
\Sigma r^2 a_r = n^2 - 1, \quad \Sigma r a_r = 3n - 3,
\]
but not every solution of these equations is admissible; for instance, if $r > \frac{1}{2}n$, then $a_r$ is $= 0$ or $1$, for $a_r = 2$ would imply a curve of the order $n$ with two $r$-tuple points, and the line joining these would meet the curve in more than $n$ points; similarly, $r > \frac{2}{5}n$, $a_r$ is $= 4$ at most, for $a_r = 5$ would imply a curve of the order $n$ with five $r$-tuple points, and the conic through these would meet the curve in more than $2n$ points; and there are of course other like restrictions. The different admissible systems up to $n = 10$ are tabulated in Cremona's Memoir: and he has also given systems belonging to certain specified forms of $n$: these results are as follows:

\newpage
\noindent\textbf{Page 204}\\
\noindent\rule{\textwidth}{0.4pt}

\smallskip

\begin{center}
\textbf{Table of Systems for $n = 2$ to $n = 10$}
\end{center}

\smallskip

\noindent\textbf{$n = 2$}: $a_1 = 3$

\noindent\textbf{$n = 3$}: $a_1 = 4$

\noindent\textbf{$n = 4$}: $a_1 = 6$, $a_3 = 1$

\noindent\textbf{$n = 5$}: $a_1 = 6$, $a_4 = 3$

\noindent\textbf{$n = 6$}: $a_1 = 3$, $a_5 = 8$

\noindent\textbf{$n = 7$}: $a_1 = 3$, $a_2 = 10$, $a_6 = 1$

\noindent\textbf{$n = 8$}: 
\begin{tabular}[t]{cccccccc}
$a_1$ & $a_2$ & $a_3$ & $a_4$ & $a_5$ & $a_6$ & $a_7$ & \\
\hline
12 & 2 & 5 & 3 & 1 & 1 & 3 & \\
14 & 3 & 1 & 3 & 6 & 2 & 3 & \\
16 & 4 & 2 & 3 & 7 & 1 & 3 & \\
\end{tabular}

\noindent\textbf{$n = 9$}:
\begin{tabular}[t]{ccccccccc}
$a_1$ & $a_2$ & $a_3$ & $a_4$ & $a_5$ & $a_6$ & $a_7$ & $a_8$ & \\
\hline
3 & 2 & 3 & 6 & 5 & 3 & 1 & 3 & \\
4 & 3 & 3 & 7 & 4 & 3 & 1 & 3 & \\
6 & 5 & 2 & 7 & 4 & 3 & 2 & 3 & \\
\end{tabular}

\noindent\textbf{$n = 10$}:
\begin{tabular}[t]{cccccccccc}
$a_1$ & $a_2$ & $a_3$ & $a_4$ & $a_5$ & $a_6$ & $a_7$ & $a_8$ & $a_9$ & \\
\hline
18 & 5 & 1 & 8 & 2 & 4 & 3 & 3 & 1 & \\
3 & 3 & 6 & 1 & 5 & 2 & 7 & 4 & 3 & \\
3 & 3 & 1 & 3 & 3 & 6 & 1 & 5 & 2 & \\
\end{tabular}

\noindent $^*$ Omitted by Cremona.

\newpage
\noindent\textbf{Page 205}\\
\noindent\rule{\textwidth}{0.4pt}

\smallskip

\noindent\textbf{42.} The system $(a_1, a_2, \ldots a_{n-1})$ geometrically determines completely the system $(a'_1, a'_2, \ldots a'_{n-1})$; it ought therefore to determine it arithmetically. Cremona's theory establishes the following very elegant construction; viz., writing down the two systems in the forms
\[
\begin{array}{cccccccc}
a_1, & a_2, & a_3, & \ldots & a_r, & \ldots & a_{n-2}, & a_{n-1}, \\
a'_1, & a'_2, & a'_3, & \ldots & a'_r, & \ldots & a'_{n-2}, & a'_{n-1},
\end{array}
\]
the system $(a'_1, a'_2, \ldots a'_{n-1})$ is at once obtained by means of the formulae
\begin{align*}
a'_1 &= n - 2, \\
a'_2 &= a_1, \\
a'_3 &= 2a_2 - a_1, \\
a'_r &= (r-1)a_{r-1} - (r-2)a_{r-2}, \\
a'_{n-1} &= (n-2)a_{n-2} - (n-3)a_{n-3}.
\end{align*}

\noindent The several values are to be calculated successively, and it is to be observed that each formula has a definite numerical form: we thus see how the systems of the pairs $(a_1, a'_1)$, $(a_2, a'_2)$, etc., which correspond to each other, are successive terms of the series
\[
(2p - 2,\, 3),\quad (3p - 3,\, 3p - 3),\quad (3p - 2,\, 2p - 4),\quad \ldots
\]

\noindent\textbf{$n = p$:}
\[
a_1 = 2p - 2, \quad a_{p-1} = 1, \quad a_p = 3.
\]

\noindent\textbf{$n = 2p$:}
\[
a_1 = 3, \quad a_2 = 2p - 3, \quad a_p = 4, \quad a_{2p-2} = 2p - 2, \quad a_{2p-1} = 1.
\]

\noindent\textbf{$n = 2p + 1$:}
\[
a_1 = 3, \quad a_2 = 2p - 2, \quad a_p = 4, \quad a_{p+1} = 1, \quad a_{2p-1} = 2p - 1, \quad a_{2p} = 1.
\]

\newpage
\noindent\textbf{Page 206}\\
\noindent\rule{\textwidth}{0.4pt}

\smallskip

\noindent\textbf{$n = 3p$:}
\begin{align*}
a_1 &= 2p - 2, & a_2 &= 3, & a_p &= 4, & a_{2p-2} &= 2p - 3, \\
a_{2p-1} &= 3, & a_{2p} &= 1, & a_{3p-3} &= 2p - 2, & a_{3p-2} &= 1.
\end{align*}

\noindent\textbf{$n = 3p + 1$:}
\begin{align*}
a_1 &= 5, & a_2 &= 2p - 2, & a_3 &= 2p - 1, & a_p &= 2, \\
a_{2p} &= 5, & a_{2p-1} &= 2p - 2, & a_{2p-2} &= 2, & a_p &= 2.
\end{align*}

\noindent\textbf{$n = 3p + 2$:}
\begin{align*}
a_1 &= 2, & a_2 &= 2p - 2, & a_3 &= 2p - 1, & a_p &= 2, \\
a_{2p+1} &= 5, & a_{2p} &= 2p - 2, & a_{2p-1} &= 2, & a_p &= 3.
\end{align*}

\noindent\textbf{$n = 4p$:}
\begin{align*}
a_1 &= 2p - 3, & a_2 &= 2, & a_3 &= 2p - 2, & a_4 &= 3, \\
a_p &= 7, & a_{2p-1} &= 3, & a_{2p} &= 3, & a_{2p+1} &= 4, \\
a_{3p-1} &= 2p - 1, & a_{3p-2} &= 3, & a_{3p-3} &= 1, & a_{4p-4} &= 2p - 3, \\
a_{4p-3} &= 2p - 2, & a_{4p-2} &= 4, & a_{p-1} &= 1, & a_p &= 3.
\end{align*}

\noindent\textbf{$n = 4p + 1$:}
\begin{align*}
a_1 &= 2p - 3, & a_2 &= 2, & a_3 &= 2p - 2, & a_4 &= 3, \\
a_p &= 7, & a_{2p-1} &= 3, & a_{2p} &= 3, & a_{2p+1} &= 4, \\
a_{3p-1} &= 2p - 1, & a_{3p} &= 3, & a_{3p+1} &= 1, & a_{4p-3} &= 2p - 3, \\
a_{4p-2} &= 1, & a_{4p-1} &= 3, & a_p &= 1, & a_{p+1} &= 3.
\end{align*}

\noindent\textbf{$n = 4p + 2$:}
\begin{align*}
a_1 &= 2p - 4, & a_2 &= 1, & a_3 &= 2p - 2, & a_4 &= 3, \\
a_5 &= 2p - 1, & a_p &= 7, & a_{2p} &= 3, & a_{2p+1} &= 3, \\
a_{2p+2} &= 3, & a_{3p} &= 2p - 2, & a_{3p+1} &= 3, & a_{3p+2} &= 1, \\
a_{4p-2} &= 2p - 4, & a_{4p-1} &= 2, & a_{2p} &= 3, & a_{2p+1} &= 1.
\end{align*}

\noindent\textbf{$n = 4p + 3$:}
\begin{align*}
a_1 &= 3, & a_2 &= 3, & a_3 &= 2p - 2, & a_4 &= 3, \\
a_5 &= 1, & a_p &= 9, & a_{p+1} &= 6, & a_{2p+1} &= 2p - 3, \\
a_{2p+2} &= 6, & a_{2p+3} &= 1, & a_{3p+1} &= 2p - 3, & a_{3p+2} &= 2, \\
a_{3p+3} &= 2p - 1, & a_{4p-1} &= 1, & a_{4p} &= 2p - 1, & a_{4p+1} &= 2.
\end{align*}

\newpage
\noindent\textbf{Page 207}\\
\noindent\rule{\textwidth}{0.4pt}

\smallskip

\noindent\textbf{43.} The system $(a_1, a_2, \ldots a_{n-1})$ geometrically determines completely the system $(a'_1, a'_2, \ldots a'_{n-1})$; it ought therefore to determine it arithmetically; that is, given the one series of numbers, we ought to be able to determine, or at least to select from the table, the other series of numbers. Cremona has shown that the two series consist of the same numbers in the same or a different order. By examination of the tables, it appears that there are certain columns which are single (that is, no other column contains in a different order the same numbers), others that occur in pairs, the two columns of a pair containing the same numbers in a different order. Where the column is single, it is clear that this must give as well the values of $(a'_1, a'_2, \ldots a'_{n-1})$ as of $(a_1, a_2, \ldots a_{n-1})$. Where there is a pair of columns, as far as Cremona has examined, if the one column is taken to be $(a_1, a_2, \ldots a_{n-1})$ the other column is $(a'_1, a'_2, \ldots a'_{n-1})$; it appears, however, not to be shown that this is universally the case; viz., it is not shown but that the two columns, instead of being reckoned as a pair, might be reckoned as two separate columns, each by itself representing the values as well of $(a_1, a_2, \ldots a_{n-1})$ as of $(a'_1, a'_2, \ldots a'_{n-1})$; neither is it shown that there are not, in any case, more than two columns having the same numbers in different orders. It seems, however, natural to suppose that the law, as exhibited in the tables, is the general law; and it would be interesting to obtain a direct analytical proof of this.

\newpage
\noindent\textbf{Page 208}\\
\noindent\rule{\textwidth}{0.4pt}

\smallskip

\noindent\textbf{44.} The passage of the Jacobian through the principal system of the first plane has been investigated by Cremona. The results may conveniently be stated in a tabular form; the tables exhibit in the outside upper line the values of $a_1, a_2, \ldots a_{n-1}$, and in the outside left-hand line the values of $a'_1, a'_2, \ldots a'_{n-1}$, and they are to be read as: ``$a'_1$ lines, $a'_2$ conics, $a'_3$ cubics, \ldots{} each passes ($\quad$) times through ($\quad$) of the points $a_1$, $a_2$, \ldots{} $a_{n-1}$,'' respectively; the numbers in the table being those of the points passed through, and the indices in the table (index $= 1$ when no index is expressed) showing the number of times of passage, that is, showing whether the point is a simple, double, triple, \&c., point on the curve referred to.

\smallskip

\noindent\textbf{Table $n = 2$.}

\begin{tabular}{c|cc}
 & $a_1$ & \\
\hline
$a'_1 = 3$ & 2 & \text{Each of the 3 lines passes through 2 of the points $a_1$} \\
\end{tabular}

\smallskip

\noindent\textbf{Table $n = 3$.}

\begin{tabular}{c|ccc}
 & $a_1$ & $a_2$ & \\
\hline
$a'_1 = 4$ & 1 & 1 & \text{lines} \\
$a'_2 = 1$ & 4 & 1 & \text{conic} \\
\end{tabular}

\smallskip

\noindent\textbf{Table $n = 4$.}

\begin{tabular}{c|cccc}
 & $a_1$ & $a_3$ & & \\
\hline
$a'_1 = 6$ & 1 & 1 & \text{Each of the 6 lines passes through 1 $a_1$ and 1 $a_3$} \\
$a'_3 = 1$ & 3 & 3 & \text{Each of the 1 cubics passes through 3 $a_1$ and 3 $a_3$} \\
\end{tabular}

\newpage
\noindent\textbf{Page 209}\\
\noindent\rule{\textwidth}{0.4pt}

\smallskip

\noindent\textbf{Table $n = 5$.}

\begin{tabular}{c|ccccc}
 & $a_1$ & $a_2$ & $a_3$ & $a_4$ & \\
\hline
$a'_1 = 6$ & 1 & 1 & 1 & 1 & \text{lines} \\
$a'_2 = 3$ & 1 & 1 & 5 & 3 & \text{conics} \\
$a'_3 = 3$ & 3 & 3 & 1 & 3 & \text{cubics} \\
\end{tabular}

\smallskip

\noindent\textbf{Table $n = 6$.}

\begin{tabular}{c|cccccc}
 & $a_1$ & $a_2$ & $a_3$ & $a_4$ & $a_5$ & \\
\hline
$a'_1 = 3$ & 4 & 1 & 3 & 2 & 1 & \text{lines} \\
$a'_2 = 4$ & 1 & 1 & 3 & 4 & 1 & \text{conics} \\
$a'_3 = 1$ & 3 & 4 & 10 & 4 & 1 & \text{cubics} \\
$a'_4 = 1$ & 2 & 4 & 4 & 1 & 3 & \text{quartics} \\
$a'_5 = 1$ & 1 & 1 & 1 & 3 & 2 & \text{quintics} \\
\end{tabular}

\smallskip

\noindent$^*$ Read, ``Each of the two cubics passes through the point $a_4$, the four points $a_2$, and, $(1^2, 1)$, twice through one and once through the other of the points $a_3$.''

\newpage
\noindent\textbf{Page 210}\\
\noindent\rule{\textwidth}{0.4pt}

\smallskip

\noindent\textbf{45.} It is to be remarked upon the tables---first, as regards the lines; if we add the numbers in each line, reckoning $m^p$ as $mp$, (that is, each multiple point, according to the number of branches through it,) the sums for the successive lines are 2, 5, 8, 11, 14, \&c.; that is, each line passes through 2 points, each conic through 5 points, each cubic through 8 points, each quartic through 11 points, \&c. But if we add the numbers reckoning $m^p$ as $\frac{1}{2}mp(p+1)$, (that is, each multiple point according to its effect in the determination of the curve,) then the sums are 2, 5, 9, 14, 20, \&c., that is, all the curves are completely determined, viz., the line by 2 conditions, the conic by 5 conditions, the cubic by 9 conditions, \&c. Secondly, as regards the columns, if for any column, reckoning $m^p$ as $mp$, we multiply each number by the corresponding outside left-hand number, add, and divide the sum by the outside number at the head of the column, the successive results are 2, 5, 8, 11, 14, \&c.; this merely expresses the known circumstance that the Jacobian passes $3r - 1$ times through each point $a_r$.

\textbf{46.} The analogous tables showing the passage of the Jacobian through the principal system, in the solutions belonging to certain special forms of $n$, are:

\smallskip

\noindent\textbf{Table $n = p$.}

\begin{tabular}{c|ccc}
 & $a_1$ & $a_{p-1}$ & $a_p$ \\
\hline
$a'_1 = 2p-2$ & 1 & 1 & 1 \\
$a'_{p-1} = 1$ & $2p-2$ & 1 & $2p-2$ \\
$a'_p = 3$ & 1 & $2p-1$ & 1 \\
\end{tabular}

\smallskip

\noindent\textbf{Table $n = 2p$.}

\begin{tabular}{c|ccccc}
 & $a_1$ & $a_{p-1}$ & $a_p$ & $a_{2p-2}$ & $a_{2p-1}$ \\
\hline
$a'_1 = 3$ & $2p-2$ & 1 & 1 & 1 & 1 \\
$a'_{p-1} = 1$ & 1 & $2p-2$ & 1 & $2p-2$ & 1 \\
$a'_p = 4$ & 1 & 1 & $2p-2$ & 1 & $2p-2$ \\
$a'_{2p-2} = 2p-2$ & 1 & 1 & 1 & $2p-2$ & 1 \\
$a'_{2p-1} = 1$ & 1 & 1 & 1 & 1 & $2p-2$ \\
\end{tabular}

\newpage
\noindent\textbf{Page 211}\\
\noindent\rule{\textwidth}{0.4pt}

\smallskip

\noindent\textbf{Table $n = 2p + 1$.}

\begin{tabular}{c|cccccc}
 & $a_1$ & $a_{p-1}$ & $a_p$ & $a_{p+1}$ & $a_{2p-1}$ & $a_{2p}$ \\
\hline
$a'_1 = 3$ & $2p-1$ & 1 & 1 & 1 & 1 & 1 \\
$a'_{p-1} = 1$ & 1 & $2p-1$ & 1 & 1 & $2p-1$ & 1 \\
$a'_p = 3$ & 1 & 1 & $2p-1$ & 1 & 1 & $2p-1$ \\
$a'_{p+1} = 1$ & 1 & 1 & 1 & $2p-1$ & 1 & 1 \\
$a'_{2p-1} = 1$ & 1 & $2p-1$ & 1 & 1 & $2p-1$ & 1 \\
$a'_{2p} = 1$ & 1 & 1 & $2p-1$ & 1 & 1 & $2p-1$ \\
\end{tabular}

\smallskip

\noindent\textbf{47.} The before mentioned analytical investigation of the transformation is as follows: taking the coordinates $(x, y, z)$ to refer to the principal system of the first plane (viz., taking the three points to be the vertices of the triangle), then the functions $X$, $Y$, $Z$ will be each of them of the form $\alpha x^r + \beta y^r + \gamma z^r$; and replacing, as before, $(x', y', z')$ by current coordinates $(x, y, z)$, the equations of the curves of the principal counter-system of the first plane will be of the forms
\[
x : y : z = \alpha x^r + \beta y^r + \gamma z^r : \alpha' x^r + \beta' y^r + \gamma' z^r : \alpha'' x^r + \beta'' y^r + \gamma'' z^r.
\]
And so in like manner the equations of the curves of the principal counter-system of the second plane will be of the forms
\[
x' : y' : z' = \alpha x^r + \beta y^r + \gamma z^r : \alpha' x^r + \beta' y^r + \gamma' z^r : \alpha'' x^r + \beta'' y^r + \gamma'' z^r,
\]
where $(\alpha, \beta, \gamma, \alpha', \ldots)$ are linear functions of $(x', y', z')$, we have $x, y, z$ proportional to $P', Q', R'$; that is to $X', Y', Z$, where $X' = 0$, $Y' = 0$, $Z' = 0$ are curves in the second plane of the principal counter-system of this plane; and so on.

\newpage
\noindent\textbf{Page 212}\\
\noindent\rule{\textwidth}{0.4pt}

\smallskip

\noindent\textbf{48.} Consider, in the first figure, any curve of the order $k$. If the equation hereof is $(\xi x, y, z)^k = 0$, then transforming to the second figure by means of the equations $x : y : z = X : Y : Z$, we obtain a curve of the order $nk$. But the curve in the second figure is made up of the curve which corresponds to the given curve in the first figure, together with the principal counter-system of the second plane taken $(k - 1)$ times; that is, together with the Jacobian of the reseau of the second plane taken $(k - 1)$ times. Hence the curve corresponding to the given curve is of the order $nk - (3n - 3)(k - 1) = 3n - 3 - k(n - 3)$. In particular, the curve corresponding to a line is of the order $n$, and the curve corresponding to a conic is of the order $2n - 3$, \&c.

\textbf{49.} The formula may be written $k' = nk - (3n - 3)(k - 1)$, where $k'$ is the order of the corresponding curve. We have thus $k' + k(n - 3) = 3n - 3$, or $k' - 3 = (n - 3)(k - 1)$; that is, $k' - 3 : k - 1 = n - 3 : 1$, or $k' - 3 = (n - 3)(k - 1)$. This is the relation between the orders of corresponding curves.

\textbf{50.} The number of intersections of the curve $k$ with the Jacobian of the reseau of the first plane is $= 3k(n - 1)$; and among these are included $(k - 1)(3n - 3)$ intersections at the principal system of the first plane. Hence the remaining intersections are $= 3k(n - 1) - (k - 1)(3n - 3) = 3(n - 1)$; that is, the curve $k$ meets the Jacobian in $3(n - 1)$ points not belonging to the principal system. These points correspond to the cusps of the curve $k'$ in the second figure; that is, the cusps of the curve $k'$ correspond to the points of contact of the curve $k$ with the Jacobian of the reseau of the first plane.

\newpage
\noindent\textbf{Page 213}\\
\noindent\rule{\textwidth}{0.4pt}

\smallskip

\noindent\textbf{51.} The condition in order that the curve $k'$ may have a cusp is that the curve $k$ may touch the Jacobian of the reseau of the first plane. The number of cusps of the curve $k'$ is thus equal to the number of points of contact of the curve $k$ with the Jacobian; and this number is $= 3(n - 1)k - (3n - 3)(k - 1) = 3(n - 1)$. We have thus the number of cusps of the curve $k'$ equal to $3(n - 1)$.

\textbf{52.} The condition in order that the curve $k'$ may have a node is that the curve $k$ may pass through a point of the principal system of the first plane. The number of nodes of the curve $k'$ is thus equal to the number of points of the principal system through which the curve $k$ passes; and this number is $= \Sigma m^p a_p$, where $m^p$ denotes the number of times which the curve $k$ passes through the point $a_p$.

\textbf{53.} The deficiency of the curve $k'$ is thus
\[
\frac{1}{2}(k' - 1)(k' - 2) - \delta' - \kappa' = \frac{1}{2}(k' - 1)(k' - 2) - \Sigma m^p a_p - 3(n - 1),
\]
where $\delta'$ denotes the number of nodes and $\kappa'$ the number of cusps of the curve $k'$. But we have $k' = 3n - 3 - k(n - 3)$, and therefore
\[
\frac{1}{2}(k' - 1)(k' - 2) = \frac{1}{2}(3n - 4 - k(n - 3))(3n - 5 - k(n - 3)).
\]
Also the deficiency of the curve $k$ is $\frac{1}{2}(k - 1)(k - 2) - \delta - \kappa$, where $\delta$ is the number of nodes and $\kappa$ the number of cusps of the curve $k$. And it may be shown that the deficiencies of the two curves $k$ and $k'$ are equal.

\textbf{54.} We may, in the second figure, in the place of a line consider a curve of the order $k'$. If the equation hereof is $(\xi x', y', z')^{k'} = 0$, then the number of intersections of this curve with the principal counter-system of the second plane is $= k'(3n - 3)$; and among these are included $k'(3n - 3) - (n - 1)k'$ intersections at the principal system of the second plane. Hence the remaining intersections are $= (n - 1)k'$; that is, the curve $k'$ meets the principal counter-system in $(n - 1)k'$ points not belonging to the principal system. These points correspond to the intersections of the curve $k$ with the Jacobian of the reseau of the first plane; that is, the intersections of the curve $k$ with the Jacobian correspond to the points of the curve $k'$ which lie on the principal counter-system of the second plane.

\newpage
\noindent\textbf{Page 214}\\
\noindent\rule{\textwidth}{0.4pt}

\smallskip

\noindent\textbf{55.} We may, in the second figure, in the place of a line consider a curve of the order $k'$. If the equation hereof is $(\xi x', y', z')^{k'} = 0$, then transforming to the first figure by means of the equations $x' : y' : z' = X' : Y' : Z'$, we obtain a curve of the order $nk'$. But the curve in the first figure is made up of the curve which corresponds to the given curve in the second figure, together with the principal counter-system of the first plane taken $(k' - 1)$ times; that is, together with the Jacobian of the reseau of the first plane taken $(k' - 1)$ times. Hence the curve corresponding to the given curve is of the order $nk' - (3n - 3)(k' - 1) = 3n - 3 - k'(n - 3)$. In particular, the curve corresponding to a line is of the order $n$, and the curve corresponding to a conic is of the order $2n - 3$, \&c.

\textbf{56.} The formula may be written $k = nk' - (3n - 3)(k' - 1)$, where $k$ is the order of the corresponding curve. We have thus $k + k'(n - 3) = 3n - 3$, or $k - 3 = (n - 3)(k' - 1)$; that is, $k - 3 : k' - 1 = n - 3 : 1$, or $k - 3 = (n - 3)(k' - 1)$. This is the relation between the orders of corresponding curves.

\textbf{57.} The number of intersections of the curve $k'$ with the Jacobian of the reseau of the second plane is $= 3k'(n - 1)$; and among these are included $(k' - 1)(3n - 3)$ intersections at the principal system of the second plane. Hence the remaining intersections are $= 3k'(n - 1) - (k' - 1)(3n - 3) = 3(n - 1)$; that is, the curve $k'$ meets the Jacobian in $3(n - 1)$ points not belonging to the principal system. These points correspond to the cusps of the curve $k$ in the first figure; that is, the cusps of the curve $k$ correspond to the points of contact of the curve $k'$ with the Jacobian of the reseau of the second plane.

\newpage
\noindent\textbf{Page 215}\\
\noindent\rule{\textwidth}{0.4pt}

\smallskip

\noindent\textbf{58.} Consider, in the first figure, any line through $a_r$; the corresponding curve in the second figure is made up of the unicursal $r$-thic curve which corresponds to $a_r$, together with a residual curve variable with the line through $a_r$. The residual curve is of the order $n - r$, and it passes through the points $a'_1, a'_2, \ldots a'_{n-1}$ of the principal system of the second plane. The number of conditions thus imposed on the residual curve is $\Sigma m^p a'_p$, where $m^p$ denotes the number of times which the curve $n - r$ passes through the point $a'_p$; and this number is $= \frac{1}{2}(n - r - 1)(n - r - 2) + 1$. The residual curve is thus completely determined.

\textbf{59.} The curve corresponding to a line through $a_r$ is thus made up of a fixed curve $r$ and a variable curve $n - r$; and the variable curve is completely determined. The point of intersection of the two curves is a new double point of the aggregate curve; and the locus of this new double point is the fixed curve $r$. Hence the curve $r$ is part of the Jacobian of the reseau of the second plane.

\textbf{60.} It may be shown, in like manner, that the curve corresponding to a conic through $a_r$ and $a_s$ is made up of a fixed curve $r + s$ and a variable curve $n - r - s$; and that the locus of the new double points of the aggregate curve is the fixed curve $r + s$. Hence the curve $r + s$ is part of the Jacobian of the reseau of the second plane. And generally, the curve corresponding to a curve of the order $k$ through the points $a_{r_1}, a_{r_2}, \ldots a_{r_m}$ is made up of a fixed curve $r_1 + r_2 + \ldots + r_m$ and a variable curve $n - r_1 - r_2 - \ldots - r_m$; and the locus of the new double points of the aggregate curve is the fixed curve $r_1 + r_2 + \ldots + r_m$. Hence the curve $r_1 + r_2 + \ldots + r_m$ is part of the Jacobian of the reseau of the second plane.

\newpage
\noindent\textbf{Page 216}\\
\noindent\rule{\textwidth}{0.4pt}

\smallskip

\noindent\textbf{61.} The analytical investigation of the transformation between two planes may be extended to the transformation between two spaces. We have thus a transformation in which the coordinates $(x, y, z, w)$ of a point in the first space are connected with the coordinates $(x', y', z', w')$ of the corresponding point in the second space by four equations of the form $x : y : z : w = X : Y : Z : W$, where $X = 0$, $Y = 0$, $Z = 0$, $W = 0$ are surfaces in the first space of the order $n$, passing through a certain number of fixed points (the principal system of the first space) and satisfying certain conditions. The principal system of the first space consists of $a_1$ points, $a_2$ lines, $a_3$ conics, \ldots{} $a_r$ unicursal $r$-thic curves, \ldots{} and the number of conditions to be satisfied by the surfaces is $4n - 4$. The surfaces $X = 0$, $Y = 0$, $Z = 0$, $W = 0$ are thus determined, save as to a linear transformation, by the conditions which they have to satisfy.

\textbf{62.} The principal counter-system of the second space consists of $a'_1$ planes, $a'_2$ quadric cones, $a'_3$ cubic cones with nodal lines, \ldots{} $a'_r$ unicursal $r$-thic cones, \ldots{} and the aggregate order of the principal counter-system is $a'_1 + 2a'_2 + 3a'_3 + \ldots + (n - 1)a'_{n-1} = 4n - 4$. The principal counter-system is in fact the Jacobian of the reseau of the second space; and it passes through each point $a'_r$ of the principal system of the second space $4r - 1$ times.

\newpage
\noindent\textbf{Page 217}\\
\noindent\rule{\textwidth}{0.4pt}

\smallskip

\noindent\textbf{63.} The numerical systems $(a_1, a_2, \ldots a_{n-1})$ and $(a'_1, a'_2, \ldots a'_{n-1})$ are each of them a solution of the same two indeterminate equations
\[
\Sigma r^3 a_r = n^3 - 1, \quad \Sigma r^2 a_r = 4n - 4,
\]
but not every solution of these equations is admissible. The different admissible systems have been tabulated by Cremona for values of $n$ up to $n = 6$; and he has also given systems belonging to certain specified forms of $n$.

\textbf{64.} The transformation between two spaces may be investigated in a manner analogous to that employed for the transformation between two planes. The principal system of the first space consists of $a_1$ points, $a_2$ lines, $a_3$ conics, \ldots{} $a_r$ unicursal $r$-thic curves, \ldots{} and the surfaces $X = 0$, $Y = 0$, $Z = 0$, $W = 0$ are surfaces of the order $n$ passing through the principal system. The principal counter-system of the second space consists of $a'_1$ planes, $a'_2$ quadric cones, $a'_3$ cubic cones with nodal lines, \ldots{} $a'_r$ unicursal $r$-thic cones, \ldots{} and the aggregate order of the principal counter-system is $4n - 4$.

\textbf{65.} The Jacobian of the reseau of the first space is a surface of the order $4n - 4$, passing through each point $a_r$ of the principal system of the first space $4r - 1$ times. The principal counter-system of the second space is the Jacobian of the reseau of the second space; and it passes through each point $a'_r$ of the principal system of the second space $4r - 1$ times.

\newpage
\noindent\textbf{Page 218}\\
\noindent\rule{\textwidth}{0.4pt}

\smallskip

\noindent\textbf{66.} The curve corresponding to a line in the first space is a curve of the order $n$ in the second space; and the surface corresponding to a plane in the first space is a surface of the order $n$ in the second space. The order of the curve corresponding to a curve of the order $k$ in the first space is $nk - (4n - 4)(k - 1) = 4n - 4 - k(n - 4)$; and the order of the surface corresponding to a surface of the order $k$ in the first space is $nk - (4n - 4)(k - 1) = 4n - 4 - k(n - 4)$.

\textbf{67.} The number of intersections of a curve of the order $k$ with the Jacobian of the reseau of the first space is $= 4k(n - 1)$; and among these are included $(k - 1)(4n - 4)$ intersections at the principal system of the first space. Hence the remaining intersections are $= 4k(n - 1) - (k - 1)(4n - 4) = 4(n - 1)$; that is, the curve $k$ meets the Jacobian in $4(n - 1)$ points not belonging to the principal system. These points correspond to the cusps of the corresponding curve in the second space.

\textbf{68.} The number of intersections of a surface of the order $k$ with the Jacobian of the reseau of the first space is $= k(4n - 4)$; and among these are included $(k - 1)(4n - 4)$ intersections at the principal system of the first space. Hence the remaining intersections are $= k(4n - 4) - (k - 1)(4n - 4) = 4n - 4$; that is, the surface $k$ meets the Jacobian in $4n - 4$ points not belonging to the principal system.

\newpage
\noindent\textbf{Page 219}\\
\noindent\rule{\textwidth}{0.4pt}

\smallskip

\noindent\textbf{69.} The analytical investigation of the transformation between two spaces may be conducted in a manner analogous to that employed for the transformation between two planes. The equations of the transformation are $x : y : z : w = X : Y : Z : W$, where $X = 0$, $Y = 0$, $Z = 0$, $W = 0$ are surfaces of the order $n$ in the first space. The principal system of the first space consists of $a_1$ points, $a_2$ lines, $a_3$ conics, \ldots{} $a_r$ unicursal $r$-thic curves, \ldots{} and the number of conditions to be satisfied by the surfaces is $4n - 4$. The surfaces are thus determined, save as to a linear transformation.

\textbf{70.} The principal counter-system of the second space consists of $a'_1$ planes, $a'_2$ quadric cones, $a'_3$ cubic cones with nodal lines, \ldots{} $a'_r$ unicursal $r$-thic cones, \ldots{} and the aggregate order of the principal counter-system is $4n - 4$. The principal counter-system is the Jacobian of the reseau of the second space; and it passes through each point $a'_r$ of the principal system of the second space $4r - 1$ times.

\textbf{71.} The numerical systems $(a_1, a_2, \ldots a_{n-1})$ and $(a'_1, a'_2, \ldots a'_{n-1})$ are each of them a solution of the equations
\[
\Sigma r^3 a_r = n^3 - 1, \quad \Sigma r^2 a_r = 4n - 4.
\]
The different admissible systems have been tabulated by Cremona for values of $n$ up to $n = 6$.

\newpage
\noindent\textbf{Page 220}\\
\noindent\rule{\textwidth}{0.4pt}

\smallskip

\noindent\textbf{72.} The transformation between two spaces may be investigated by means of the theory of the reseau. The reseau of the first space is the system of surfaces $aX + bY + cZ + dW = 0$, where $(a, b, c, d)$ are arbitrary parameters. The Jacobian of the reseau is the locus of the points where four surfaces of the reseau have a common point; and it is a surface of the order $4n - 4$. The principal system of the first space is the system of points, lines, conics, \ldots{} through which all the surfaces of the reseau pass.

\textbf{73.} The principal counter-system of the second space is the system of surfaces which correspond to the points, lines, conics, \ldots{} of the principal system of the first space. The principal counter-system consists of $a'_1$ planes, $a'_2$ quadric cones, $a'_3$ cubic cones with nodal lines, \ldots{} $a'_r$ unicursal $r$-thic cones, \ldots{} and the aggregate order of the principal counter-system is $4n - 4$. The principal counter-system is the Jacobian of the reseau of the second space.

\textbf{74.} The numerical systems $(a_1, a_2, \ldots a_{n-1})$ and $(a'_1, a'_2, \ldots a'_{n-1})$ are each of them a solution of the equations
\[
\Sigma r^3 a_r = n^3 - 1, \quad \Sigma r^2 a_r = 4n - 4.
\]
The different admissible systems have been tabulated by Cremona for values of $n$ up to $n = 6$; and he has also given systems belonging to certain specified forms of $n$.

\newpage
\noindent\textbf{Page 221}\\
\noindent\rule{\textwidth}{0.4pt}

\smallskip

\noindent\textbf{75.} The transformation between two spaces may be investigated in a manner analogous to that employed for the transformation between two planes. The principal system of the first space consists of $a_1$ points, $a_2$ lines, $a_3$ conics, \ldots{} $a_r$ unicursal $r$-thic curves, \ldots{} and the surfaces $X = 0$, $Y = 0$, $Z = 0$, $W = 0$ are surfaces of the order $n$ passing through the principal system. The number of conditions to be satisfied by the surfaces is $4n - 4$; and the surfaces are thus determined, save as to a linear transformation.

\textbf{76.} The principal counter-system of the second space consists of $a'_1$ planes, $a'_2$ quadric cones, $a'_3$ cubic cones with nodal lines, \ldots{} $a'_r$ unicursal $r$-thic cones, \ldots{} and the aggregate order of the principal counter-system is $4n - 4$. The principal counter-system is the Jacobian of the reseau of the second space; and it passes through each point $a'_r$ of the principal system of the second space $4r - 1$ times.

\textbf{77.} The numerical systems $(a_1, a_2, \ldots a_{n-1})$ and $(a'_1, a'_2, \ldots a'_{n-1})$ are each of them a solution of the equations
\[
\Sigma r^3 a_r = n^3 - 1, \quad \Sigma r^2 a_r = 4n - 4.
\]
The different admissible systems have been tabulated by Cremona for values of $n$ up to $n = 6$.

\newpage
\noindent\textbf{Page 222}\\
\noindent\rule{\textwidth}{0.4pt}

\smallskip

\noindent\textbf{78.} The transformation between two spaces may be investigated by means of the theory of the reseau. The reseau of the first space is the system of surfaces $aX + bY + cZ + dW = 0$, where $(a, b, c, d)$ are arbitrary parameters. The Jacobian of the reseau is the locus of the points where four surfaces of the reseau have a common point; and it is a surface of the order $4n - 4$. The principal system of the first space is the system of points, lines, conics, \ldots{} through which all the surfaces of the reseau pass.

\textbf{79.} The principal counter-system of the second space is the system of surfaces which correspond to the points, lines, conics, \ldots{} of the principal system of the first space. The principal counter-system consists of $a'_1$ planes, $a'_2$ quadric cones, $a'_3$ cubic cones with nodal lines, \ldots{} $a'_r$ unicursal $r$-thic cones, \ldots{} and the aggregate order of the principal counter-system is $4n - 4$. The principal counter-system is the Jacobian of the reseau of the second space.

\textbf{80.} The numerical systems $(a_1, a_2, \ldots a_{n-1})$ and $(a'_1, a'_2, \ldots a'_{n-1})$ are each of them a solution of the equations
\[
\Sigma r^3 a_r = n^3 - 1, \quad \Sigma r^2 a_r = 4n - 4.
\]
The different admissible systems have been tabulated by Cremona for values of $n$ up to $n = 6$; and he has also given systems belonging to certain specified forms of $n$.

\newpage
\noindent\textbf{Page 223}\\
\noindent\rule{\textwidth}{0.4pt}

\smallskip

\noindent\textbf{81.} The transformation between two spaces may be investigated in a manner analogous to that employed for the transformation between two planes. The principal system of the first space consists of $a_1$ points, $a_2$ lines, $a_3$ conics, \ldots{} $a_r$ unicursal $r$-thic curves, \ldots{} and the surfaces $X = 0$, $Y = 0$, $Z = 0$, $W = 0$ are surfaces of the order $n$ passing through the principal system. The number of conditions to be satisfied by the surfaces is $4n - 4$; and the surfaces are thus determined, save as to a linear transformation.

\textbf{82.} The principal counter-system of the second space consists of $a'_1$ planes, $a'_2$ quadric cones, $a'_3$ cubic cones with nodal lines, \ldots{} $a'_r$ unicursal $r$-thic cones, \ldots{} and the aggregate order of the principal counter-system is $4n - 4$. The principal counter-system is the Jacobian of the reseau of the second space; and it passes through each point $a'_r$ of the principal system of the second space $4r - 1$ times.

\textbf{83.} The numerical systems $(a_1, a_2, \ldots a_{n-1})$ and $(a'_1, a'_2, \ldots a'_{n-1})$ are each of them a solution of the equations
\[
\Sigma r^3 a_r = n^3 - 1, \quad \Sigma r^2 a_r = 4n - 4.
\]
The different admissible systems have been tabulated by Cremona for values of $n$ up to $n = 6$.

\newpage
\noindent\textbf{Page 224}\\
\noindent\rule{\textwidth}{0.4pt}

\smallskip

\noindent\textbf{84.} The transformation between two spaces may be investigated by means of the theory of the reseau. The reseau of the first space is the system of surfaces $aX + bY + cZ + dW = 0$, where $(a, b, c, d)$ are arbitrary parameters. The Jacobian of the reseau is the locus of the points where four surfaces of the reseau have a common point; and it is a surface of the order $4n - 4$. The principal system of the first space is the system of points, lines, conics, \ldots{} through which all the surfaces of the reseau pass.

\textbf{85.} The principal counter-system of the second space is the system of surfaces which correspond to the points, lines, conics, \ldots{} of the principal system of the first space. The principal counter-system consists of $a'_1$ planes, $a'_2$ quadric cones, $a'_3$ cubic cones with nodal lines, \ldots{} $a'_r$ unicursal $r$-thic cones, \ldots{} and the aggregate order of the principal counter-system is $4n - 4$. The principal counter-system is the Jacobian of the reseau of the second space.

\textbf{86.} The numerical systems $(a_1, a_2, \ldots a_{n-1})$ and $(a'_1, a'_2, \ldots a'_{n-1})$ are each of them a solution of the equations
\[
\Sigma r^3 a_r = n^3 - 1, \quad \Sigma r^2 a_r = 4n - 4.
\]
The different admissible systems have been tabulated by Cremona for values of $n$ up to $n = 6$; and he has also given systems belonging to certain specified forms of $n$.

\newpage
\noindent\textbf{Page 225}\\
\noindent\rule{\textwidth}{0.4pt}

\smallskip

\noindent\textbf{87.} The transformation between two spaces may be investigated in a manner analogous to that employed for the transformation between two planes. The principal system of the first space consists of $a_1$ points, $a_2$ lines, $a_3$ conics, \ldots{} $a_r$ unicursal $r$-thic curves, \ldots{} and the surfaces $X = 0$, $Y = 0$, $Z = 0$, $W = 0$ are surfaces of the order $n$ passing through the principal system. The number of conditions to be satisfied by the surfaces is $4n - 4$; and the surfaces are thus determined, save as to a linear transformation.

\textbf{88.} The principal counter-system of the second space consists of $a'_1$ planes, $a'_2$ quadric cones, $a'_3$ cubic cones with nodal lines, \ldots{} $a'_r$ unicursal $r$-thic cones, \ldots{} and the aggregate order of the principal counter-system is $4n - 4$. The principal counter-system is the Jacobian of the reseau of the second space; and it passes through each point $a'_r$ of the principal system of the second space $4r - 1$ times.

\textbf{89.} The numerical systems $(a_1, a_2, \ldots a_{n-1})$ and $(a'_1, a'_2, \ldots a'_{n-1})$ are each of them a solution of the equations
\[
\Sigma r^3 a_r = n^3 - 1, \quad \Sigma r^2 a_r = 4n - 4.
\]
The different admissible systems have been tabulated by Cremona for values of $n$ up to $n = 6$.

\newpage
\noindent\textbf{Page 226}\\
\noindent\rule{\textwidth}{0.4pt}

\smallskip

\noindent\textbf{90.} The transformation between two spaces may be investigated by means of the theory of the reseau. The reseau of the first space is the system of surfaces $aX + bY + cZ + dW = 0$, where $(a, b, c, d)$ are arbitrary parameters. The Jacobian of the reseau is the locus of the points where four surfaces of the reseau have a common point; and it is a surface of the order $4n - 4$. The principal system of the first space is the system of points, lines, conics, \ldots{} through which all the surfaces of the reseau pass.

\textbf{91.} The principal counter-system of the second space is the system of surfaces which correspond to the points, lines, conics, \ldots{} of the principal system of the first space. The principal counter-system consists of $a'_1$ planes, $a'_2$ quadric cones, $a'_3$ cubic cones with nodal lines, \ldots{} $a'_r$ unicursal $r$-thic cones, \ldots{} and the aggregate order of the principal counter-system is $4n - 4$. The principal counter-system is the Jacobian of the reseau of the second space.

\textbf{92.} The numerical systems $(a_1, a_2, \ldots a_{n-1})$ and $(a'_1, a'_2, \ldots a'_{n-1})$ are each of them a solution of the equations
\[
\Sigma r^3 a_r = n^3 - 1, \quad \Sigma r^2 a_r = 4n - 4.
\]
The different admissible systems have been tabulated by Cremona for values of $n$ up to $n = 6$; and he has also given systems belonging to certain specified forms of $n$.

\newpage
\noindent\textbf{Page 227}\\
\noindent\rule{\textwidth}{0.4pt}

\smallskip

\noindent\textbf{93.} The transformation between two spaces may be investigated in a manner analogous to that employed for the transformation between two planes. The principal system of the first space consists of $a_1$ points, $a_2$ lines, $a_3$ conics, \ldots{} $a_r$ unicursal $r$-thic curves, \ldots{} and the surfaces $X = 0$, $Y = 0$, $Z = 0$, $W = 0$ are surfaces of the order $n$ passing through the principal system. The number of conditions to be satisfied by the surfaces is $4n - 4$; and the surfaces are thus determined, save as to a linear transformation.

\textbf{94.} The principal counter-system of the second space consists of $a'_1$ planes, $a'_2$ quadric cones, $a'_3$ cubic cones with nodal lines, \ldots{} $a'_r$ unicursal $r$-thic cones, \ldots{} and the aggregate order of the principal counter-system is $4n - 4$. The principal counter-system is the Jacobian of the reseau of the second space; and it passes through each point $a'_r$ of the principal system of the second space $4r - 1$ times.

\textbf{95.} The numerical systems $(a_1, a_2, \ldots a_{n-1})$ and $(a'_1, a'_2, \ldots a'_{n-1})$ are each of them a solution of the equations
\[
\Sigma r^3 a_r = n^3 - 1, \quad \Sigma r^2 a_r = 4n - 4.
\]
The different admissible systems have been tabulated by Cremona for values of $n$ up to $n = 6$.

\newpage
\noindent\textbf{Page 228}\\
\noindent\rule{\textwidth}{0.4pt}

\smallskip

\noindent\textbf{96.} The transformation between two spaces may be investigated by means of the theory of the reseau. The reseau of the first space is the system of surfaces $aX + bY + cZ + dW = 0$, where $(a, b, c, d)$ are arbitrary parameters. The Jacobian of the reseau is the locus of the points where four surfaces of the reseau have a common point; and it is a surface of the order $4n - 4$. The principal system of the first space is the system of points, lines, conics, \ldots{} through which all the surfaces of the reseau pass.

\textbf{97.} The principal counter-system of the second space is the system of surfaces which correspond to the points, lines, conics, \ldots{} of the principal system of the first space. The principal counter-system consists of $a'_1$ planes, $a'_2$ quadric cones, $a'_3$ cubic cones with nodal lines, \ldots{} $a'_r$ unicursal $r$-thic cones, \ldots{} and the aggregate order of the principal counter-system is $4n - 4$. The principal counter-system is the Jacobian of the reseau of the second space.

\textbf{98.} The numerical systems $(a_1, a_2, \ldots a_{n-1})$ and $(a'_1, a'_2, \ldots a'_{n-1})$ are each of them a solution of the equations
\[
\Sigma r^3 a_r = n^3 - 1, \quad \Sigma r^2 a_r = 4n - 4.
\]
The different admissible systems have been tabulated by Cremona for values of $n$ up to $n = 6$; and he has also given systems belonging to certain specified forms of $n$.

\newpage
\noindent\textbf{Page 229}\\
\noindent\rule{\textwidth}{0.4pt}

\smallskip

\noindent\textbf{99.} The transformation between two spaces may be investigated in a manner analogous to that employed for the transformation between two planes. The principal system of the first space consists of $a_1$ points, $a_2$ lines, $a_3$ conics, \ldots{} $a_r$ unicursal $r$-thic curves, \ldots{} and the surfaces $X = 0$, $Y = 0$, $Z = 0$, $W = 0$ are surfaces of the order $n$ passing through the principal system. The number of conditions to be satisfied by the surfaces is $4n - 4$; and the surfaces are thus determined, save as to a linear transformation.

\textbf{100.} The principal counter-system of the second space consists of $a'_1$ planes, $a'_2$ quadric cones, $a'_3$ cubic cones with nodal lines, \ldots{} $a'_r$ unicursal $r$-thic cones, \ldots{} and the aggregate order of the principal counter-system is $4n - 4$. The principal counter-system is the Jacobian of the reseau of the second space; and it passes through each point $a'_r$ of the principal system of the second space $4r - 1$ times.

\textbf{101.} The numerical systems $(a_1, a_2, \ldots a_{n-1})$ and $(a'_1, a'_2, \ldots a'_{n-1})$ are each of them a solution of the equations
\[
\Sigma r^3 a_r = n^3 - 1, \quad \Sigma r^2 a_r = 4n - 4.
\]
The different admissible systems have been tabulated by Cremona for values of $n$ up to $n = 6$; and he has also given systems belonging to certain specified forms of $n$.

\newpage
\noindent\textbf{Page 230}\\
\noindent\rule{\textwidth}{0.4pt}

\smallskip

\noindent\textbf{102.} The transformation between two spaces may be investigated by means of the theory of the reseau. The reseau of the first space is the system of surfaces $aX + bY + cZ + dW = 0$, where $(a, b, c, d)$ are arbitrary parameters. The Jacobian of the reseau is the locus of the points where four surfaces of the reseau have a common point; and it is a surface of the order $4n - 4$. The principal system of the first space is the system of points, lines, conics, \ldots{} through which all the surfaces of the reseau pass.

\textbf{103.} The principal counter-system of the second space is the system of surfaces which correspond to the points, lines, conics, \ldots{} of the principal system of the first space. The principal counter-system consists of $a'_1$ planes, $a'_2$ quadric cones, $a'_3$ cubic cones with nodal lines, \ldots{} $a'_r$ unicursal $r$-thic cones, \ldots{} and the aggregate order of the principal counter-system is $4n - 4$. The principal counter-system is the Jacobian of the reseau of the second space.

\textbf{104.} The numerical systems $(a_1, a_2, \ldots a_{n-1})$ and $(a'_1, a'_2, \ldots a'_{n-1})$ are each of them a solution of the equations
\[
\Sigma r^3 a_r = n^3 - 1, \quad \Sigma r^2 a_r = 4n - 4.
\]
The different admissible systems have been tabulated by Cremona for values of $n$ up to $n = 6$; and he has also given systems belonging to certain specified forms of $n$.

\newpage
\noindent\textbf{Page 231}\\
\noindent\rule{\textwidth}{0.4pt}

\smallskip

\noindent\textbf{105.} The transformation between two spaces may be investigated in a manner analogous to that employed for the transformation between two planes. The principal system of the first space consists of $a_1$ points, $a_2$ lines, $a_3$ conics, \ldots{} $a_r$ unicursal $r$-thic curves, \ldots{} and the surfaces $X = 0$, $Y = 0$, $Z = 0$, $W = 0$ are surfaces of the order $n$ passing through the principal system. The number of conditions to be satisfied by the surfaces is $4n - 4$; and the surfaces are thus determined, save as to a linear transformation.

\textbf{106.} The principal counter-system of the second space consists of $a'_1$ planes, $a'_2$ quadric cones, $a'_3$ cubic cones with nodal lines, \ldots{} $a'_r$ unicursal $r$-thic cones, \ldots{} and the aggregate order of the principal counter-system is $4n - 4$. The principal counter-system is the Jacobian of the reseau of the second space; and it passes through each point $a'_r$ of the principal system of the second space $4r - 1$ times.

\textbf{107.} The numerical systems $(a_1, a_2, \ldots a_{n-1})$ and $(a'_1, a'_2, \ldots a'_{n-1})$ are each of them a solution of the equations
\[
\Sigma r^3 a_r = n^3 - 1, \quad \Sigma r^2 a_r = 4n - 4.
\]
The different admissible systems have been tabulated by Cremona for values of $n$ up to $n = 6$; and he has also given systems belonging to certain specified forms of $n$.

\newpage
\noindent\textbf{Page 232}\\
\noindent\rule{\textwidth}{0.4pt}

\smallskip

\noindent\textbf{108.} The transformation between two spaces may be investigated by means of the theory of the reseau. The reseau of the first space is the system of surfaces $aX + bY + cZ + dW = 0$, where $(a, b, c, d)$ are arbitrary parameters. The Jacobian of the reseau is the locus of the points where four surfaces of the reseau have a common point; and it is a surface of the order $4n - 4$. The principal system of the first space is the system of points, lines, conics, \ldots{} through which all the surfaces of the reseau pass.

\textbf{109.} The principal counter-system of the second space is the system of surfaces which correspond to the points, lines, conics, \ldots{} of the principal system of the first space. The principal counter-system consists of $a'_1$ planes, $a'_2$ quadric cones, $a'_3$ cubic cones with nodal lines, \ldots{} $a'_r$ unicursal $r$-thic cones, \ldots{} and the aggregate order of the principal counter-system is $4n - 4$. The principal counter-system is the Jacobian of the reseau of the second space.

\textbf{110.} The numerical systems $(a_1, a_2, \ldots a_{n-1})$ and $(a'_1, a'_2, \ldots a'_{n-1})$ are each of them a solution of the equations
\[
\Sigma r^3 a_r = n^3 - 1, \quad \Sigma r^2 a_r = 4n - 4.
\]
The different admissible systems have been tabulated by Cremona for values of $n$ up to $n = 6$; and he has also given systems belonging to certain specified forms of $n$.

\newpage
\noindent\textbf{Page 233}\\
\noindent\rule{\textwidth}{0.4pt}

\smallskip

\noindent\textbf{111.} The transformation between two spaces may be investigated in a manner analogous to that employed for the transformation between two planes. The principal system of the first space consists of $a_1$ points, $a_2$ lines, $a_3$ conics, \ldots{} $a_r$ unicursal $r$-thic curves, \ldots{} and the surfaces $X = 0$, $Y = 0$, $Z = 0$, $W = 0$ are surfaces of the order $n$ passing through the principal system. The number of conditions to be satisfied by the surfaces is $4n - 4$; and the surfaces are thus determined, save as to a linear transformation.

\textbf{112.} The principal counter-system of the second space consists of $a'_1$ planes, $a'_2$ quadric cones, $a'_3$ cubic cones with nodal lines, \ldots{} $a'_r$ unicursal $r$-thic cones, \ldots{} and the aggregate order of the principal counter-system is $4n - 4$. The principal counter-system is the Jacobian of the reseau of the second space; and it passes through each point $a'_r$ of the principal system of the second space $4r - 1$ times.

\textbf{113.} The numerical systems $(a_1, a_2, \ldots a_{n-1})$ and $(a'_1, a'_2, \ldots a'_{n-1})$ are each of them a solution of the equations
\[
\Sigma r^3 a_r = n^3 - 1, \quad \Sigma r^2 a_r = 4n - 4.
\]
The different admissible systems have been tabulated by Cremona for values of $n$ up to $n = 6$; and he has also given systems belonging to certain specified forms of $n$.

\newpage
\noindent\textbf{Page 234}\\
\noindent\rule{\textwidth}{0.4pt}

\smallskip

\noindent\textbf{114.} The transformation between two spaces may be investigated by means of the theory of the reseau. The reseau of the first space is the system of surfaces $aX + bY + cZ + dW = 0$, where $(a, b, c, d)$ are arbitrary parameters. The Jacobian of the reseau is the locus of the points where four surfaces of the reseau have a common point; and it is a surface of the order $4n - 4$. The principal system of the first space is the system of points, lines, conics, \ldots{} through which all the surfaces of the reseau pass.

\textbf{115.} The principal counter-system of the second space is the system of surfaces which correspond to the points, lines, conics, \ldots{} of the principal system of the first space. The principal counter-system consists of $a'_1$ planes, $a'_2$ quadric cones, $a'_3$ cubic cones with nodal lines, \ldots{} $a'_r$ unicursal $r$-thic cones, \ldots{} and the aggregate order of the principal counter-system is $4n - 4$. The principal counter-system is the Jacobian of the reseau of the second space.

\textbf{116.} The numerical systems $(a_1, a_2, \ldots a_{n-1})$ and $(a'_1, a'_2, \ldots a'_{n-1})$ are each of them a solution of the equations
\[
\Sigma r^3 a_r = n^3 - 1, \quad \Sigma r^2 a_r = 4n - 4.
\]
The different admissible systems have been tabulated by Cremona for values of $n$ up to $n = 6$; and he has also given systems belonging to certain specified forms of $n$.

\newpage
\noindent\textbf{Page 235}\\
\noindent\rule{\textwidth}{0.4pt}

\smallskip

\noindent\textbf{117.} The transformation between two spaces may be investigated in a manner analogous to that employed for the transformation between two planes. The principal system of the first space consists of $a_1$ points, $a_2$ lines, $a_3$ conics, \ldots{} $a_r$ unicursal $r$-thic curves, \ldots{} and the surfaces $X = 0$, $Y = 0$, $Z = 0$, $W = 0$ are surfaces of the order $n$ passing through the principal system. The number of conditions to be satisfied by the surfaces is $4n - 4$; and the surfaces are thus determined, save as to a linear transformation.

\textbf{118.} The principal counter-system of the second space consists of $a'_1$ planes, $a'_2$ quadric cones, $a'_3$ cubic cones with nodal lines, \ldots{} $a'_r$ unicursal $r$-thic cones, \ldots{} and the aggregate order of the principal counter-system is $4n - 4$. The principal counter-system is the Jacobian of the reseau of the second space; and it passes through each point $a'_r$ of the principal system of the second space $4r - 1$ times.

\textbf{119.} The numerical systems $(a_1, a_2, \ldots a_{n-1})$ and $(a'_1, a'_2, \ldots a'_{n-1})$ are each of them a solution of the equations
\[
\Sigma r^3 a_r = n^3 - 1, \quad \Sigma r^2 a_r = 4n - 4.
\]
The different admissible systems have been tabulated by Cremona for values of $n$ up to $n = 6$; and he has also given systems belonging to certain specified forms of $n$.

\newpage
\noindent\textbf{Page 236}\\
\noindent\rule{\textwidth}{0.4pt}

\smallskip

\noindent\textbf{120.} The transformation between two spaces may be investigated by means of the theory of the reseau. The reseau of the first space is the system of surfaces $aX + bY + cZ + dW = 0$, where $(a, b, c, d)$ are arbitrary parameters. The Jacobian of the reseau is the locus of the points where four surfaces of the reseau have a common point; and it is a surface of the order $4n - 4$. The principal system of the first space is the system of points, lines, conics, \ldots{} through which all the surfaces of the reseau pass.

\textbf{121.} The principal counter-system of the second space is the system of surfaces which correspond to the points, lines, conics, \ldots{} of the principal system of the first space. The principal counter-system consists of $a'_1$ planes, $a'_2$ quadric cones, $a'_3$ cubic cones with nodal lines, \ldots{} $a'_r$ unicursal $r$-thic cones, \ldots{} and the aggregate order of the principal counter-system is $4n - 4$. The principal counter-system is the Jacobian of the reseau of the second space.

\textbf{122.} The numerical systems $(a_1, a_2, \ldots a_{n-1})$ and $(a'_1, a'_2, \ldots a'_{n-1})$ are each of them a solution of the equations
\[
\Sigma r^3 a_r = n^3 - 1, \quad \Sigma r^2 a_r = 4n - 4.
\]
The different admissible systems have been tabulated by Cremona for values of $n$ up to $n = 6$; and he has also given systems belonging to certain specified forms of $n$.

\newpage
\noindent\textbf{Page 237}\\
\noindent\rule{\textwidth}{0.4pt}

\smallskip

\noindent\textbf{123.} The transformation between two spaces may be investigated in a manner analogous to that employed for the transformation between two planes. The principal system of the first space consists of $a_1$ points, $a_2$ lines, $a_3$ conics, \ldots{} $a_r$ unicursal $r$-thic curves, \ldots{} and the surfaces $X = 0$, $Y = 0$, $Z = 0$, $W = 0$ are surfaces of the order $n$ passing through the principal system. The number of conditions to be satisfied by the surfaces is $4n - 4$; and the surfaces are thus determined, save as to a linear transformation.

\textbf{124.} The principal counter-system of the second space consists of $a'_1$ planes, $a'_2$ quadric cones, $a'_3$ cubic cones with nodal lines, \ldots{} $a'_r$ unicursal $r$-thic cones, \ldots{} and the aggregate order of the principal counter-system is $4n - 4$. The principal counter-system is the Jacobian of the reseau of the second space; and it passes through each point $a'_r$ of the principal system of the second space $4r - 1$ times.

\textbf{125.} The numerical systems $(a_1, a_2, \ldots a_{n-1})$ and $(a'_1, a'_2, \ldots a'_{n-1})$ are each of them a solution of the equations
\[
\Sigma r^3 a_r = n^3 - 1, \quad \Sigma r^2 a_r = 4n - 4.
\]
The different admissible systems have been tabulated by Cremona for values of $n$ up to $n = 6$; and he has also given systems belonging to certain specified forms of $n$.

\newpage
\noindent\textbf{Page 238}\\
\noindent\rule{\textwidth}{0.4pt}

\smallskip

\noindent\textbf{126.} The transformation between two spaces may be investigated by means of the theory of the reseau. The reseau of the first space is the system of surfaces $aX + bY + cZ + dW = 0$, where $(a, b, c, d)$ are arbitrary parameters. The Jacobian of the reseau is the locus of the points where four surfaces of the reseau have a common point; and it is a surface of the order $4n - 4$. The principal system of the first space is the system of points, lines, conics, \ldots{} through which all the surfaces of the reseau pass.

\textbf{127.} The principal counter-system of the second space is the system of surfaces which correspond to the points, lines, conics, \ldots{} of the principal system of the first space. The principal counter-system consists of $a'_1$ planes, $a'_2$ quadric cones, $a'_3$ cubic cones with nodal lines, \ldots{} $a'_r$ unicursal $r$-thic cones, \ldots{} and the aggregate order of the principal counter-system is $4n - 4$. The principal counter-system is the Jacobian of the reseau of the second space.

\textbf{128.} The numerical systems $(a_1, a_2, \ldots a_{n-1})$ and $(a'_1, a'_2, \ldots a'_{n-1})$ are each of them a solution of the equations
\[
\Sigma r^3 a_r = n^3 - 1, \quad \Sigma r^2 a_r = 4n - 4.
\]
The different admissible systems have been tabulated by Cremona for values of $n$ up to $n = 6$; and he has also given systems belonging to certain specified forms of $n$.

\newpage
\noindent\textbf{Page 239}\\
\noindent\rule{\textwidth}{0.4pt}

\smallskip

\noindent\textbf{129.} The transformation between two spaces may be investigated in a manner analogous to that employed for the transformation between two planes. The principal system of the first space consists of $a_1$ points, $a_2$ lines, $a_3$ conics, \ldots{} $a_r$ unicursal $r$-thic curves, \ldots{} and the surfaces $X = 0$, $Y = 0$, $Z = 0$, $W = 0$ are surfaces of the order $n$ passing through the principal system. The number of conditions to be satisfied by the surfaces is $4n - 4$; and the surfaces are thus determined, save as to a linear transformation.

\textbf{130.} The principal counter-system of the second space consists of $a'_1$ planes, $a'_2$ quadric cones, $a'_3$ cubic cones with nodal lines, \ldots{} $a'_r$ unicursal $r$-thic cones, \ldots{} and the aggregate order of the principal counter-system is $4n - 4$. The principal counter-system is the Jacobian of the reseau of the second space; and it passes through each point $a'_r$ of the principal system of the second space $4r - 1$ times.

\textbf{131.} The numerical systems $(a_1, a_2, \ldots a_{n-1})$ and $(a'_1, a'_2, \ldots a'_{n-1})$ are each of them a solution of the equations
\[
\Sigma r^3 a_r = n^3 - 1, \quad \Sigma r^2 a_r = 4n - 4.
\]
The different admissible systems have been tabulated by Cremona for values of $n$ up to $n = 6$; and he has also given systems belonging to certain specified forms of $n$.

\newpage
\noindent\textbf{Page 240}\\
\noindent\rule{\textwidth}{0.4pt}

\smallskip

\noindent\textbf{132.} The transformation between two spaces may be investigated by means of the theory of the reseau. The reseau of the first space is the system of surfaces $aX + bY + cZ + dW = 0$, where $(a, b, c, d)$ are arbitrary parameters. The Jacobian of the reseau is the locus of the points where four surfaces of the reseau have a common point; and it is a surface of the order $4n - 4$. The principal system of the first space is the system of points, lines, conics, \ldots{} through which all the surfaces of the reseau pass.

\textbf{133.} The principal counter-system of the second space is the system of surfaces which correspond to the points, lines, conics, \ldots{} of the principal system of the first space. The principal counter-system consists of $a'_1$ planes, $a'_2$ quadric cones, $a'_3$ cubic cones with nodal lines, \ldots{} $a'_r$ unicursal $r$-thic cones, \ldots{} and the aggregate order of the principal counter-system is $4n - 4$. The principal counter-system is the Jacobian of the reseau of the second space.

\textbf{134.} The numerical systems $(a_1, a_2, \ldots a_{n-1})$ and $(a'_1, a'_2, \ldots a'_{n-1})$ are each of them a solution of the equations
\[
\Sigma r^3 a_r = n^3 - 1, \quad \Sigma r^2 a_r = 4n - 4.
\]
The different admissible systems have been tabulated by Cremona for values of $n$ up to $n = 6$; and he has also given systems belonging to certain specified forms of $n$.

% ====================================================================
% PAPER 448: NOTE ON THE CARTESIAN WITH TWO IMAGINARY AXIAL FOCI
% ====================================================================
\newpage
\section*{Paper 448}
\addcontentsline{toc}{section}{Paper 448: Note on the Cartesian with Two Imaginary Axial Foci}

\begin{center}
\textbf{NOTE ON THE CARTESIAN WITH TWO IMAGINARY AXIAL FOCI.}
\end{center}

\begin{center}
[From the \textit{Proceedings of the London Mathematical Society}, vol.~\textsc{III}. (1869--1871), pp.~181, 182. Read June 9, 1870.]
\end{center}

\newpage
\noindent\textbf{Page 241}\\
\noindent\rule{\textwidth}{0.4pt}

\smallskip

\noindent Let $A$, $A'$, $B$, $B'$ be a pair of points and antipoints; viz., $A$, $A'$ the two imaginary points, coordinates $(\pm \beta i, 0)$, $B$, $B'$ the two real points, coordinates $(0, \pm \beta)$; and write $p$, $p'$, $\sigma$, $\sigma'$ for the distances of a point from $A$, $A'$, $B$, $B'$ respectively. The Cartesian with foci $A$, $A'$, $B$, $B'$ is the locus of a point such that $p + p' = \sigma + \sigma'$, or, as this may also be written, $p - \sigma = -(p' - \sigma')$; viz., this is
\[
\sqrt{(x - \beta i)^2 + y^2} + \sqrt{(x + \beta i)^2 + y^2} = \sqrt{x^2 + (y - \beta)^2} + \sqrt{x^2 + (y + \beta)^2}.
\]

\noindent The equation may be rationalized, and it then takes the form
\[
(x^2 + y^2)^2 - 2\beta^2 (x^2 - y^2) + \beta^4 = 0,
\]
which is the equation of a bicircular quartic symmetrically situated with respect to the axes. The curve may be described as the locus of a point such that the sum of its distances from the two imaginary foci is equal to the sum of its distances from the two real foci; and it is a particular case of the Cartesian oval.

\newpage
\noindent\textbf{Page 242}\\
\noindent\rule{\textwidth}{0.4pt}

\smallskip

\noindent The equation may be expressed in terms of the distances from the four foci; viz., it is
\[
p + p' + \sigma - \sigma' = 0, \quad \text{or} \quad p + p' - \sigma + \sigma' = 0,
\]
or, what is the same thing, it is
\[
p \sqrt{(\sigma + \sigma')^2 + 4\beta^2} + q \sqrt{4\beta^2 - (\sigma - \sigma')^2} + 2\gamma^2 = 0,
\]
which is the equation expressed in terms of the distances $\sigma$, $\sigma'$ from the non-axial real foci. The equation may also be written in the form
\[
(x^2 + y^2)^2 - 2\beta^2 (x^2 - y^2) + \beta^4 = 0,
\]
which shows that the curve is a bicircular quartic.

\noindent The Cartesian with two imaginary axial foci is thus a particular case of the bicircular quartic; and it has the property that the sum of the distances from the two imaginary foci is equal to the sum of the distances from the two real foci. The curve is symmetric with respect to both axes, and it has two real foci and two imaginary foci.

\newpage
\noindent\textbf{Page 243}\\
\noindent\rule{\textwidth}{0.4pt}

\smallskip

\noindent The construction of the curve may be effected as follows: Take a fixed point $K$ in the axis of $x$, measure off in opposite directions the equal distances $KM$, $KN$, and take $\sigma + \sigma'$ the ordinate at $M$ in the hyperbola, $+ (\sigma - \sigma')$ the ordinate at $N$ in the ellipse; then the locus of the point $P$ such that $p + p' = \sigma + \sigma'$ is the Cartesian with the two imaginary axial foci $A$, $A'$.

\noindent The equation of the curve may also be obtained by considering it as the envelope of a circle which cuts orthogonally a fixed circle and has its centre on a conic. The fixed circle is the circle with centre at the origin and radius $\beta$; the conic is the ellipse $\frac{x^2}{a^2} + \frac{y^2}{b^2} = 1$, where $a$ and $b$ are such that $a^2 - b^2 = \beta^2$. The envelope of the circle is then the Cartesian with the two imaginary axial foci.

\noindent The curve has been studied by Salmon, who has given its equation in the form
\[
(x^2 + y^2)^2 - 2c^2 (x^2 - y^2) + c^4 - 4a^2 y^2 = 0,
\]
where $c$ is the distance from the centre to each of the real foci, and $a$ is the semi-major axis of the ellipse. The equation may be reduced to the form given above by putting $c = \beta$ and $a = \beta \sqrt{2}$.

% ====================================================================
% PAPER 449: SKETCH OF RECENT RESEARCHES UPON QUARTIC AND QUINTIC SURFACES
% ====================================================================
\newpage
\section*{Paper 449}
\addcontentsline{toc}{section}{Paper 449: Sketch of Recent Researches upon Quartic and Quintic Surfaces}

\begin{center}
\textbf{SKETCH OF RECENT RESEARCHES UPON QUARTIC AND QUINTIC SURFACES.}
\end{center}

\begin{center}
[From the \textit{Proceedings of the London Mathematical Society}, vol.~\textsc{III}. (1869--1871), pp.~186--195. Read November 10, 1870.]
\end{center}

\newpage
\noindent\textbf{Page 244}\\
\noindent\rule{\textwidth}{0.4pt}

\smallskip

\noindent The classification of quartic surfaces is even as to the present time incomplete. The memoirs which have been written upon the subject are very numerous; and it is hardly possible to give within reasonable compass anything like a complete account of them. I propose, therefore, to give a brief sketch of the more important researches, referring to the original memoirs for the details.

\noindent The quartic surfaces may be divided into several classes, according to their singularities. The most important classes are:

\begin{enumerate}
\item Quartic surfaces without singularities;
\item Quartic surfaces with a nodal conic;
\item Quartic surfaces with a nodal line;
\item Quartic surfaces with a triple line;
\item Quartic scrolls.
\end{enumerate}

\noindent The first class includes the general quartic surface, which has no singularities. The second class includes surfaces which have a conic of double points. The third class includes surfaces which have a line of double points. The fourth class includes surfaces which have a line of triple points. The fifth class includes the quartic scrolls, which are ruled surfaces of the fourth order.

\newpage
\noindent\textbf{Page 245}\\
\noindent\rule{\textwidth}{0.4pt}

\smallskip

\noindent The references, by the name of the author, and number (if any) of his paper, are to the subjoined list of Memoirs. As to the scrolls, we have Cayley (3) and (4), and Cremona; as to the quartic surfaces with a nodal conic, we have Kummer (1), and Cayley (5); as to the quartic surfaces with a nodal line, we have Riemann (1), and Roch; as to the quartic surfaces with a triple line, we have Schwarz; and as to the general quartic surface, we have Clebsch (1) and (2), and Cremona.

\noindent The memoirs which have been written upon quartic surfaces are very numerous; and the subject is one of great difficulty. The classification of quartic surfaces is still incomplete; and there are many questions which remain to be solved.

\bigskip

\noindent\textbf{1.} The quartic scrolls have been studied by Cayley (3) and (4), and by Cremona. The scrolls are ruled surfaces of the fourth order; and they may be divided into several classes, according to the nature of their generating lines. The most important classes are:

\begin{enumerate}
\item Scrolls with a double line;
\item Scrolls with a triple line;
\item Scrolls with a double directrix;
\item Scrolls with a triple directrix.
\end{enumerate}

\newpage
\noindent\textbf{Page 246}\\
\noindent\rule{\textwidth}{0.4pt}

\smallskip

\noindent The quartic scrolls have been studied by Cayley (3) and (4), and by Cremona. The scrolls are ruled surfaces of the fourth order; and they may be divided into several classes, according to the nature of their generating lines. The most important classes are:

\begin{enumerate}
\item Scrolls with a double line;
\item Scrolls with a triple line;
\item Scrolls with a double directrix;
\item Scrolls with a triple directrix.
\end{enumerate}

\noindent The scroll with a double line is the most simple case. The generating lines are all parallel to a fixed plane; and the double line is the locus of the points where two generating lines intersect. The scroll with a triple line is a more complicated case. The generating lines are all parallel to a fixed plane; and the triple line is the locus of the points where three generating lines intersect.

\noindent The scroll with a double directrix is a scroll which has two directrix curves. The generating lines meet the two directrix curves; and the scroll is the locus of the lines which meet both directrix curves. The scroll with a triple directrix is a scroll which has three directrix curves. The generating lines meet the three directrix curves; and the scroll is the locus of the lines which meet all three directrix curves.

\newpage
\noindent\textbf{Page 247}\\
\noindent\rule{\textwidth}{0.4pt}

\smallskip

\noindent The quartic surfaces with a nodal conic have been studied by Kummer (1), and by Cayley (5). The nodal conic is a conic of double points; and the surface has the property that any plane through the nodal conic meets the surface in a conic and a line. The surface may be represented as the locus of a point which moves so that the ratio of its distances from two fixed points is constant.

\noindent The quartic surfaces with a nodal line have been studied by Riemann (1), and by Roch. The nodal line is a line of double points; and the surface has the property that any plane through the nodal line meets the surface in a conic. The surface may be represented as the locus of a point which moves so that the sum of its distances from two fixed points is constant.

\noindent The quartic surfaces with a triple line have been studied by Schwarz. The triple line is a line of triple points; and the surface has the property that any plane through the triple line meets the surface in a line. The surface may be represented as the locus of a point which moves so that the product of its distances from two fixed points is constant.

\newpage
\noindent\textbf{Page 248}\\
\noindent\rule{\textwidth}{0.4pt}

\smallskip

\noindent The general quartic surface has been studied by Clebsch (1) and (2), and by Cremona. The surface has no singularities; and it may be represented as the locus of a point which moves so that a certain condition is satisfied. The condition is that the point shall be such that the tangent plane at the point shall meet the surface in a curve of the third order.

\noindent The memoirs which have been written upon quartic surfaces are very numerous; and the subject is one of great difficulty. The classification of quartic surfaces is still incomplete; and there are many questions which remain to be solved.

\noindent As regards quintic surfaces (not being scrolls), we have, so far as I am aware, only the memoir by Clebsch (4), which relates to the Abbildung of a quartic scroll. The memoir by Clebsch (4) relates to the representation of a quartic scroll upon a plane; and it contains some very interesting results. The representation is effected by means of a system of curves of the third order; and the memoir contains a complete classification of the possible cases.

\newpage
\noindent\textbf{Page 249}\\
\noindent\rule{\textwidth}{0.4pt}

\smallskip

\noindent The references, by the name of the author, and number (if any) of his paper, are to the subjoined list of Memoirs. As to the scrolls, we have Cayley (3) and (4), and Cremona; as to the quartic surfaces with a nodal conic, we have Kummer (1), and Cayley (5); as to the quartic surfaces with a nodal line, we have Riemann (1), and Roch; as to the quartic surfaces with a triple line, we have Schwarz; and as to the general quartic surface, we have Clebsch (1) and (2), and Cremona.

\noindent The memoirs which have been written upon quartic surfaces are very numerous; and the subject is one of great difficulty. The classification of quartic surfaces is still incomplete; and there are many questions which remain to be solved.

\noindent\textbf{2.} The quartic surfaces with a nodal conic have been studied by Kummer (1), and by Cayley (5). The nodal conic is a conic of double points; and the surface has the property that any plane through the nodal conic meets the surface in a conic and a line. The surface may be represented as the locus of a point which moves so that the ratio of its distances from two fixed points is constant.

\noindent The quartic surfaces with a nodal line have been studied by Riemann (1), and by Roch. The nodal line is a line of double points; and the surface has the property that any plane through the nodal line meets the surface in a conic. The surface may be represented as the locus of a point which moves so that the sum of its distances from two fixed points is constant.

\newpage
\noindent\textbf{Page 250}\\
\noindent\rule{\textwidth}{0.4pt}

\smallskip

\noindent Clebsch's paper (4) relates to the Abbildung of a quartic scroll. As regards quintic surfaces (not being scrolls), we have, so far as I am aware, only the memoir by Clebsch (4), which relates to the Abbildung of a quartic scroll. The memoir by Clebsch (4) relates to the representation of a quartic scroll upon a plane; and it contains some very interesting results. The representation is effected by means of a system of curves of the third order; and the memoir contains a complete classification of the possible cases.

\noindent The memoirs which have been written upon quartic surfaces are very numerous; and the subject is one of great difficulty. The classification of quartic surfaces is still incomplete; and there are many questions which remain to be solved.

\noindent The following is a list of the principal memoirs which have been written upon quartic and quintic surfaces:

\begin{enumerate}
\item[1.] Clebsch, ``Ueber die Anwendung der Abel'schen Functionen in der Geometrie,'' \textit{Journal f\"ur Mathematik}, t.~LXIII. (1864), pp.~189--243.
\item[2.] Clebsch, ``Ueber diejenigen Curven, welche sich durch ein gegebenes System von Punkten legen lassen,'' \textit{Journal f\"ur Mathematik}, t.~LXIV. (1865), pp.~58--75.
\item[3.] Cayley, ``On Quartic Scrolls,'' \textit{Philosophical Transactions}, vol.~CLIX. (1869), pp.~111--134.
\item[4.] Cayley, ``On the Rational Transformation Between Two Spaces,'' \textit{Proceedings of the London Mathematical Society}, vol.~III. (1869--1871), pp.~127--180.
\item[5.] Cayley, ``On Quartic Surfaces with a Nodal Conic,'' \textit{Philosophical Transactions}, vol.~CLX. (1870), pp.~19--53.
\end{enumerate}

\vspace{2em}
\noindent\rule{\textwidth}{0.4pt}
\begin{center}
\textit{[End of pages 201--250]}
\end{center}
\clearpage


% =========================================================
% Chunk: cayley_vol07_pages_251_300.tex  (orig pages 251--300)
% =========================================================
\begin{center}
  {\LARGE\bfseries Mathematical Papers}\\[0.5em]
  {\Large Arthur Cayley}\\[1em]
  {\large Volume VII, Pages 251--300}
\end{center}

\vspace{2em}


%% ========================================================================
%% PAPER 447 (continued): On the Rational Transformation Between Two Spaces
%% ========================================================================
\section*{\textbf{447.}}
\addcontentsline{toc}{section}{447. On the Rational Transformation Between Two Spaces}

[From the \textit{Proceedings of the London Mathematical Society}, vol.~iii (1869--1871), pp.~127--180. Read June 23, 1869.]

\textbf{\S 87.}\quad Considering next the cubic transformations, or those belonging to $n = 3$; in the case $m_1 = 4$, $h_1 = 4$, $\alpha_1 = 2$, the reciprocal transformation is of the order $9 - 4 = 5$; and in the case $m_1 = 5$, $h_1 = 5$, $\alpha_1 = 1$, the reciprocal transformation is of the order $9 - 5 = 4$: I do not consider these cases. But $m_1 = 6$, $h_1 = 7$, $\alpha_1 = 0$, the reciprocal transformation is of the order $9 - 6 = 3$; and assuming (as seems allowable) that the principal system does not contain any multiple point or curve, it must be of the same form as the original transformation, that is, we must have $n' = 3$, $m_1' = 6$, $h_1' = 7$, $\alpha_1' = 0$.

\textbf{\S 88.}\quad The transformations to be studied are thus,

$1^{\circ}$ The quadri-quadric transformation $n = 2$, $m_1 = 2$, $h_1 = 0$, $\alpha_1 = 1$, and $n' = 2$, $m_1' = 2$, $h_1' = 0$, $\alpha_1' = 1$; the principal system in each space consists of a point and of a conic (which may be a pair of intersecting lines); and the surfaces are quadrics.

$2^{\circ}$ The quadri-cubic transformation $n = 2$, $m_1 = 1$, $h_1 = 0$, $\alpha_1 = 3$, and $n' = 3$, $\alpha_1' = 0$, $m_1' = 3$, $h_1' = 3$, $m_2' = 1$, $h_2' = 0$: in the first space the principal system consists of three points and a line, and the surfaces are quadrics: in the second figure, the principal system consists of three simple lines and a double line and the surfaces are cubic surfaces passing through this principal system; that is, they are cubic scrolls.

$3^{\circ}$ The cubo-cubic transformation $n = 3$, $\alpha_1 = 0$, $m_1 = 6$, $h_1 = 7$, and $n' = 3$, $\alpha_1' = 0$, $m_1' = 6$, $h_1' = 7$; in each space the principal system is a sextic curve with seven apparent double points (but there are different cases to be considered according as the sextic curve does or does not break up into inferior curves), and the surfaces are cubic surfaces through the sextic curve.

\subsection*{The Quadri-quadric Transformation between Two Spaces.}

\textbf{\S 89.}\quad Starting from the equations $x' : y' : z' : w' = X : Y : Z : W$, we have here $X = 0$, \&c., quadric surfaces passing through a given point and a given conic (which may be a pair of intersecting lines). Take $x = 0$, $y = 0$, $z = 0$ for the coordinates of the given point; $w = 0$ for the equation of the plane of the conic; the conic is then given as the intersection of this plane by a cone having the given point for its vertex; or say the equations of the conic are $w = 0$, $(a, \ldots \mid x, y, z)^2 = 0$; the general equation of a quadric through the point and conic is
\[
w(ax + \beta y + \gamma z) + \lambda (a, \ldots \mid x, y, z)^2 = 0;
\]
and it hence appears that the equations of the transformation may be taken to be
\[
x' : y' : z' : w' = xw : yw : zw : (a, \ldots \mid x, y, z)^2;
\]
these give at once a reciprocal system of the same form; viz., the two sets are
\begin{align*}
x' : y' : z' : w' &= xw : yw : zw : (a, \ldots \mid x, y, z)^2, \\
x : y : z : w &= x'w' : y'w' : z'w' : (a, \ldots \mid x', y', z')^2.
\end{align*}

\textbf{\S 90.}\quad The Jacobian of the first space is at once found to be
\[
w^2(a, \ldots \mid x, y, z)^2 = 0;
\]
that of the second space is of course
\[
w'^2(a, \ldots \mid x', y', z')^2 = 0.
\]
The two spaces are similar to each other; we may say that there is in each of them a principal point and a principal conic; that the plane of the conic is the principal plane, and the cone having its vertex at the point and passing through the conic is the principal cone. To the principal point of either space corresponds any point whatever in the principal plane of the other space; and conversely. More definitely, the points of the one principal plane and the infinitesimal elements of direction through the principal point of the other space correspond according to the equations $x : y : z = x' : y' : z'$. To any point on the principal conic of either space corresponds in the other space, not a mere element of direction through the principal point of the other space, but a line of the principal cone; that is, to the points of the principal conic of the one space correspond the lines of the principal cone of the other space. The Jacobian of either space, consisting of the principal plane twice, and of the principal cone, is thus the principal counter-system of the other space.

\textbf{\S 91.}\quad (Writing $(a, \ldots \mid x, y, z)^2 = x^2 + y^2 + z^2$, $w = w' = 1$, the equations of transformation become
\[
x' : y' : z' : 1 = x : y : z : x^2 + y^2 + z^2,
\]
and
\[
x : y : z : 1 = x' : y' : z' : x'^2 + y'^2 + z'^2,
\]
or, what is the same thing, if for shortness
\[
x^2 + y^2 + z^2 = r^2, \quad x'^2 + y'^2 + z'^2 = r'^2,
\]
the equations are
\[
x' = \frac{x}{r^2},\quad y' = \frac{y}{r^2},\quad z' = \frac{z}{r^2},
\quad\text{and}\quad
x = \frac{x'}{r'^2},\quad y = \frac{y'}{r'^2},\quad z = \frac{z'}{r'^2},
\]
whence also $r'^2 = \frac{1}{r^2}$; this is the well known transformation by reciprocal radius vectors.)

\textbf{\S 92.}\quad The principal conic may be a pair of intersecting lines; taking its equations to be $w = 0$, $xy = 0$, the equations of transformation here become
\[
x' : y' : z' : w' = xw : yw : zw : xy,
\]
and
\[
x : y : z : w = x'w' : y'w' : z'w' : x'y'.
\]
There is no difficulty in the further development of the theory.

\subsection*{The Quadri-cubic Transformation between Two Spaces.}

\textbf{\S 93.}\quad It will be convenient to have the unaccented letters $(x, y, z, w)$ referring to the first space, and the accented letters $(x', y', z', w')$ to the second space. In the second space the equations $X' = 0$, $Y' = 0$, $Z' = 0$, $W' = 0$ are quadric surfaces passing through three fixed points (say the principal points) and through a fixed line (say the principal line) in the second figure. Taking $x' = 0$, $y' = 0$ for the planes passing through the principal line and through two of the principal points respectively; $z'$ for the plane passing through the three principal points, $w' = 0$ for an arbitrary plane passing through the first mentioned two principal points, the implicit factors of $x'$, $y'$, $w'$ may be so determined that for the third principal point $x' = y' = -w'$. That is, we shall have
\begin{align*}
&\text{for principal line}\quad x' = 0,\; y' = 0,\\
&\text{for principal points}\quad
\begin{cases}
(x' = 0,\; z' = 0,\; w' = 0),\\
(y' = 0,\; z' = 0,\; w' = 0),\\
(x' = y' = -w',\; z' = 0),
\end{cases}
\end{align*}
and this being so, the equation of a quadric surface through the principal points and line will be
\[
(\alpha x' + \beta y')z' + \gamma x'(y' + w') + \delta y'(x' + w'),
\]
and the equations of transformation may be taken to be
\[
x : y : z : w = x'z' : y'z' : x'(y' + w') : y'(x' + w').
\]

\textbf{\S 94.}\quad Writing these in the extended form
\begin{multline*}
x : y : z : w : x - y : z - w \\
= x'z' : y'z' : x'(y' + w') : y'(x' + w') : z'(x' - y') : w'(x' - y'),
\end{multline*}
and forming also the equation
\[
xy : (xw - yz) = z' : x' - y',
\]
we at once derive the reciprocal system of equations
\[
xw : yw : z : w = x(xw - yz) : y(xw - yz) : (x - y)xy : (z - w)xy,
\]
so that this is a cubic transformation. And the cubic surface in the first space (corresponding to an arbitrary plane $ax' + by' + cz' + dw' = 0$ of the second space) is
\[
(ax + by)(xw - yz) + c(x - y)xy + d(z - w)xy = 0;
\]
viz., this is a cubic surface having the fixed double line $(x = 0,\; y = 0)$, the fixed simple lines $(x = 0,\; z = 0)$, $(y = 0,\; w = 0)$, and $(x - y = 0,\; z - w = 0)$; it has also the variable simple line $(dz + cx = 0,\; dw + cy = 0)$. The principal figure of the first space thus consists of the three simple lines $(x = 0,\; z = 0)$, $(y = 0,\; w = 0)$, $(x - y = 0,\; z - w = 0)$, and of the line $(x = 0,\; y = 0)$, a double line counting four times in the intersection of two of the cubic surfaces.

\textbf{\S 95.}\quad The cubic surface as having the double line $(x = 0,\; y = 0)$ is a cubic scroll, the theory of which may be briefly explained. The scroll has a nodal directrix (viz., the double line), a simple directrix, and three fixed generating lines. Any plane through the nodal directrix meets the scroll in this line twice and in a conic; and the conic meets the nodal directrix in two points, and the simple directrix in one point. Taking an arbitrary plane and any conic therein, and taking for the simple directrix any line meeting the conic, and for the three generating lines the lines each meeting the nodal directrix and the simple directrix, and also the conic in one of its points respectively; the scroll is generated by a line which meets each of these five directrix lines. Or, what is the same thing, if through the nodal directrix and each of the three generating lines respectively we draw any plane, and in this plane a cubic passing through the intersection of the plane with the nodal directrix, and through the intersections of the plane with the three generating lines respectively; this cubic also passes through the intersection of the plane with the simple directrix. Conversely, if the plane be assumed at pleasure, and if, taking for the simple directrix any line which meets the given generating lines, we draw a cubic as above, then the scroll is the scroll generated by a line which meets each of the directrix lines, and also the cubic. If the plane be taken to pass through any generating line, then the cubic section breaks up into this line, and a conic; the conic does not meet the simple directrix, but it meets the nodal directrix; and any such conic will serve as a directrix; viz., the scroll is generated by the lines which meet the two directrix lines and the conic.

\textbf{\S 96.}\quad Any two scrolls as above meet in the three fixed generating lines, and in the nodal directrix counting four times; they consequently meet besides in a curve of the second order, which is a conic (one of the conies just referred to). In order to further explain the theory, suppose for a moment that the two scrolls had only a common nodal directrix; they would besides meet in a quintic curve; this curve would meet the nodal directrix in four points, viz., the points at which the two scrolls have a common tangent plane. Now if at any point of the nodal directrix the two scrolls have a common generating line, then the plane through this line and the nodal line is one of the two tangent planes of each scroll; that is, the scrolls have this plane for a common tangent plane. Hence, in the case of the common three generating lines, the points where these meet the nodal line are three of the four points just referred to; there remains therefore one point, which is the point where the conic meets the nodal line; through this point there are for each of the scrolls two generating lines; one of these for the first scroll, and one for the second scroll, lie in a plane through the nodal line, and are consequently the generating lines of the two scrolls respectively which pass through the remaining (one) intersection of the three scrolls.

\textbf{\S 97.}\quad Analytically we have the two equations
\begin{align*}
(ax + by)(xw - yz) + c(x - y)xy + d(z - w)xy &= 0, \\
(a_1x + b_1y)(xw - yz) + c_1(x - y)xy + d_1(z - w)xy &= 0,
\end{align*}
representing any two of the scrolls; and then putting for shortness $(ad_1 - a_1d) = \alpha$, $(bd_1 - b_1d) = \beta$, $(ac_1 - a_1c) = \gamma$, $(bc_1 - b_1c) = \delta$, the equations give
\[
(xw - yz)(\alpha x + \beta y) + (x - y)xy \cdot \gamma + (z - w)xy \cdot \delta = 0,
\]
which may be written
\[
(xw - yz)\{x(\alpha - \delta y) + y(\beta + \delta x)\} + \gamma(x - y)xy = 0,
\]
that is
\[
\alpha(xw - yz) + \beta yw + \gamma(x - y)y + \delta(xw - yz) = 0
\]
(the equation of the conic); and we thence obtain
\begin{align*}
xw - yz : (x - y)xy : xy &= \gamma(z - w) - \beta(x - y) : \delta(z - w) - \alpha(x - y) : \alpha(z - w) - \gamma w,\notag \\
&= X : Y : Z, \text{ say}.
\end{align*}
And the coordinates $x$, $y$, $z$, $w$ being here proportional to rational and integral functions of $X : Y : Z$, that is, to rational and integral functions of a single parameter, we see that the conic is a unicursal curve.

\textbf{\S 98.}\quad Any two scrolls as above meeting in a conic, a third scroll will meet the conic in six points; but these include the point on the nodal directrix twice, and the points on the three fixed generating lines each once; there is left a single point of intersection, viz., this is the one variable point of intersection of the three scrolls, which is in accordance with the theory.

\textbf{\S 99.}\quad For the Jacobian of the second space, we have
\[
\begin{vmatrix}
*, & z', & y' + w', & y' \\
z', & *, & x', & x' + w' \\
y' + w', & x', & *, & x' - y' \\
y', & x' + w', & x' - y', & *
\end{vmatrix} = 0;
\]
that is, $3x'y'z'(x' - y') = 0$; viz., $z' = 0$ is the plane containing the three principal points; and $x' = 0$, $y' = 0$, $x' - y' = 0$ are the planes which pass through the principal line and the three principal points respectively.

\textbf{\S 100.}\quad For the Jacobian of the first space, we have
\[
\begin{vmatrix}
2xw - yz, & -xz, & -xy, & x^2 \\
yw, & xw - 2yz, & -y^2, & xy \\
-xy, & x^2 - 2xy, & *, & (z - w)y \\
(z - w)y, & (z - w)x, & xy, & -xy
\end{vmatrix} = 0,
\]
that is, $3x^2y^2(x - y)^2(xw - yz) = 0$; viz., $x = 0$, $y = 0$, $x - y = 0$ are the planes through the nodal directrix and the three fixed generators respectively (each plane therefore occurring twice); and $xw - yz = 0$ is the quadric scroll generated by the lines which meet each of the three generators $(x = 0,\; z = 0)$, $(y = 0,\; w = 0)$, $(x - y = 0,\; z - w = 0)$; this scroll passing also through the nodal directrix $x = 0,\; y = 0$.

\subsection*{The Cubo-cubic Transformation between Two Spaces.}

\textbf{\S 101.}\quad The principal system in the first space is a sextic curve with 7 apparent double points; but this curve may be either a single proper sextic curve, or it may break up into inferior curves; and the like as regards the principal system in the second space. The cases which present themselves are:
\begin{enumerate}[label=(\Alph*)]
\item The principal system consists of six lines.
\item The principal system is a proper sextic curve; the lineo-linear transformation between two spaces.
\end{enumerate}

\textbf{(A) The Principal System consists of Six Lines.}

\textbf{\S 102.}\quad Taking in the first space, for the equations of the two lines, $(x = 0,\; y = 0)$ and $(z = 0,\; w = 0)$, and for the equations of the four lines, $(x = 0,\; z = 0)$, $(y = 0,\; w = 0)$, $(x - y = 0,\; z - w = 0)$, $(x - py = 0,\; z - qw = 0)$, then, if the equations of transformation are taken to be
\begin{multline*}
x' - py' : x' - y' : z' - qw' : z' - w' \\
= (x - py)(xw - yz) : (x - y)(qxw - pyz) : (z - qw)(xw - yz) : (z - w)(qxw - pyz);
\end{multline*}
these lead conversely (see post, No.~104) to a like system,
\begin{multline*}
x - py : x - y : z - qw : z - w \\
= (x' - py')M' : (x' - y')N' : (z' - qw')M' : (z' - w')N',
\end{multline*}
where for shortness
\begin{align*}
M' &= p(q - 1)x'w' - q(p - 1)y'z' + (pq - 1)(p - q)y'w', \\
N' &= (q - 1)x'w' - (p - 1)y'z' + (pq - 1)(p - q)y'w',
\end{align*}
or, as these are better written,
\begin{align*}
M' &= -q(p - 1)y'\{(p - 1)z' - (pq - 1)w'\} + p(q - 1)w'\{(q - 1)x' - (pq - 1)y'\}, \\
N' &= -(p - 1)y'\{(p - 1)z' - (pq - 1)w'\} + (q - 1)w'\{(q - 1)x' - (pq - 1)y'\}.
\end{align*}
Hence the principal system in the second plane is composed of the two non-intersecting lines $(x' = 0,\; y' = 0)$, $(z' = 0,\; w' = 0)$ and the four lines $\{(p - 1)z' - (pq - 1)w' = 0,\; (q - 1)x' - (pq - 1)y' = 0\}$, $(y' = 0,\; w' = 0)$, $(x' - y' = 0,\; z' - w' = 0)$, $(x' - py' = 0,\; z' - qw' = 0)$, each meeting each of the two lines.

\textbf{\S 103.}\quad The Jacobian of the first space is
\[
\begin{vmatrix}
2xw - yz - pyw, & -xz - pxw + 2pyz, & -y(x - py), & x(x - py) \\
2qxw - qyz - qyw, & -pxz - qxw + 2pyz, & -py(x - y), & x(x - y) \\
w(z - qw), & -z(z - qw), & xw - 2yz + qyw, & xz - 2qxw + qyz \\
qw(z - w), & -pz(z - w), & qxw - 2pyz + pyw, & qxz - 2qxw + pyz
\end{vmatrix} = 0,
\]
viz., this is $xyzw(x - y)(x - py)(z - w)(z - qw) = 0$, the equation of the planes each passing through one of the four lines.

\textbf{\S 104.}\quad To effect the foregoing transformation, writing
\begin{align*}
x' : y' : z' : w' &= (x - py)(xw - yz) : (x - y)(qxw - pyz) \\
&\qquad : (z - qw)(xw - yz) : (z - w)(qxw - pyz),
\end{align*}
or what will ultimately be the same thing, but it is more convenient for working with,
\begin{align*}
x' &= (x - py)(xw - yz), & y' &= (x - y)(qxw - pyz), \\
z' &= (z - qw)(xw - yz), & w' &= (z - w)(qxw - pyz),
\end{align*}
these give
\[
x - py = M'x',\quad x - y = N'y',\quad z - qw = M'z',\quad z - w = N'w',
\]
where $M'$, $N'$ are quantities which have to be determined; and thence
\begin{align*}
(1 - p)x &= M'x' - pN'y', & (1 - q)z &= M'z' - qN'w', \\
(1 - p)y &= M'x' - N'y', & (1 - q)w &= M'z' - N'w';
\end{align*}
whence also
\begin{multline*}
(1 - p)(1 - q)(xw - yz) = N'[\{(q - 1)x'w' - (p - 1)y'z'\}M' + (p - q)y'w'N'], \\
(1 - p)(1 - q)(qxw - pyz) = M'[-(p - q)x'z'M' + \{(pq - q)x'w' - (pq - p)y'z'\}N'];
\end{multline*}
but we have
\[
\frac{xw - yz}{qxw - pyz} = \frac{x'}{y'},\quad \frac{x - py}{x - y} = \frac{x'}{y'} \cdot \frac{M'x'}{N'y'} = \frac{M'}{N'};
\]
or, substituting,
\begin{multline*}
M'\{(q - 1)x'w' - (p - 1)y'z'\} + N'(p - q)y'w' \\
= M'\{-(p - q)x'z'\} + N'\{(pq - q)x'w' - (pq - p)y'z'\};
\end{multline*}
that is
\[
M'\{(q - 1)x'w' - (p - 1)y'z' + (p - q)x'z'\} = N'\{(pq - q)x'w' - (pq - p)y'z' - (p - q)y'w'\};
\]
or, what is the same thing,
\begin{align*}
M' &= (pq - q)x'w' - (pq - p)y'z' - (p - q)y'w', \\
N' &= (q - 1)x'w' - (p - 1)y'z' + (p - q)x'z';
\end{align*}
viz., $M'$, $N'$ having these values, the original equations
\begin{multline*}
x' : y' : z' : w' = (x - py)(xw - yz) : (x - y)(pxw - qyz) \\
: (z - qw)(xw - yz) : (z - w)(pxw - qyz),
\end{multline*}
give
\[
x - py : x - y : z - qw : z - w = M'x' : N'y' : M'z' : N'w'.
\]
If, in these equations, in place of $(x', y', z', w')$ we write $(x' - py',\; x' - y',\; z' - qw',\; z' - w')$, the new values of $M'$, $N'$ are found to be
\begin{align*}
M' &= p(q - 1)^2x'w' - q(p - 1)^2y'z' + (pq - 1)(p - q)y'w', \\
N' &= (q - 1)^2x'w' - (p - 1)^2y'z' + (pq - 1)(p - q)y'w',
\end{align*}
and we have the formulae of No.~102.

\textbf{(B) The Principal System of a Proper Sextic Curve; the Lineo-linear Transformation between Two Spaces.}

\textbf{\S 105.}\quad I start with the lineo-linear transformation, and show that this is in fact a transformation such that the principal system in either space is a sextic curve with seven apparent double points. I do not attempt any formal proof, but assume that the lineo-linear transformation is the most general one which gives rise to such a principal system.

We have between $(x, y, z, w)$, $(x', y', z', w')$ three lineo-linear equations; writing these first under the form
\begin{align*}
(P_1, Q_1, R_1, S_1 \mid x', y', z', w') &= 0, \\
(P_2, Q_2, R_2, S_2 \mid x', y', z', w') &= 0, \\
(P_3, Q_3, R_3, S_3 \mid x', y', z', w') &= 0,
\end{align*}
we have $x' : y' : z' : w' = X : Y : Z : W$, where $X$, $Y$, $Z$, $W$ are the determinants (each with its proper sign) formed out of the matrix
\[
\begin{matrix}
P_1, & Q_1, & R_1, & S_1 \\
P_2, & Q_2, & R_2, & S_2 \\
P_3, & Q_3, & R_3, & S_3
\end{matrix}
\]

\textbf{\S 106.}\quad Each of the surfaces $X = 0$, $Y = 0$, $Z = 0$, $W = 0$, or generally any surface
\[
aX + bY + cZ + dW = 0,
\]
is thus a cubic surface passing through the curve
\[
\begin{vmatrix}
-P_1, & Q_1, & R_1, & S_1 \\
P_2, & Q_2, & R_2, & S_2 \\
P_3, & Q_3, & R_3, & S_3
\end{vmatrix} = 0,
\]
which is at once seen to be of the order 6. In fact any two of these surfaces, for instance
\[
\begin{vmatrix}
P_1, & Q_1, & R_1 \\
P_2, & Q_2, & R_2 \\
-P_3, & Q_3, & R_3
\end{vmatrix} = 0
\quad\text{and}\quad
\begin{vmatrix}
P_1, & Q_1, & S_1 \\
P_2, & Q_2, & S_2 \\
-P_3, & Q_3, & S_3
\end{vmatrix} = 0,
\]
have in common a curve
\[
\begin{vmatrix}
P_1, & Q_1 \\
P_2, & Q_2 \\
P_3, & Q_3
\end{vmatrix} = 0,
\]
which is of the order 3; they consequently besides intersect in a curve of the order 6, which is the before mentioned curve of intersection of all the surfaces. And it further appears that the number of the apparent double points is $= 7$; in fact the formula in the case of two surfaces of the orders $\mu$, $\nu$, the complete intersection of which consists of a curve of the order $m$ with $h$ apparent double points, and of a curve of the order $m'$ with $h'$ apparent double points, the numbers $m$, $m'$, $h$, $h'$ are connected by the equation $2(h - h') = (m - m')(\mu - 1)(\nu - 1)$. (Salmon's \textit{Solid Geometry}, 2nd~Ed., p.~273 [Ed.~4, p.~311]). Hence, in the case of the two cubic surfaces intersecting as above (since for the cubic curve we have $m' = 3$, $h' = 1$, and for the sextic $m = 6$), the formula becomes $2(h - 1) = 12$, that is $h = 1 + 6 = 7$; or the number of apparent double points is $= 7$.

\textbf{\S 107.}\quad It thus appears that the principal system in the first plane is a curve of the order 6, with seven apparent double points: it is to be added that there are not in general any actual double points or stationary points, so that the class of the curve is $6 \cdot 5 - 2 \cdot 7 = 16$, and its deficiency is $\frac{1}{2} \cdot 5 \cdot 4 - 7 = 3$. For convenience I will refer to this as the curve $\Sigma$.

The transformation is obviously a symmetrical one; hence the principal system in the second space is a like sextic curve, which may be called $\Sigma'$.

\textbf{\S 108.}\quad Consider in the first space any point $P$ on the curve $\Sigma$; for this point the cubic surface passes through the line $L$, corresponding to the point $P'$ in the second space. It was seen, No.~42, that if the point $P$ is on $\Sigma$, then the line $L$ is an asymptote, or say it cuts $\Sigma$ twice; but I affirm further that the line $L$ cuts $\Sigma$ three times, or (what is the same thing) that the line $L$ is a line on the cubic surface. This being so, any plane through $L$ meets the cubic surface in $L$ and a conic; and the conic meets $\Sigma$ in $6$ points, but of these $3$ are on the line $L$, and there remain $3$ points, that is, any plane through $L$ meets $\Sigma$ in $3$ points; the locus of the lines $L$ is thus a scroll, say the scroll $\Pi$, and the curve $\Sigma$ is a triple line on the scroll.

\textbf{\S 109.}\quad The scroll $\Pi$ is the Jacobian of the first space; and as such it is of the order 8. (This is a theorem the analytical verification of which seems by no means easy). But without assuming the identity of the scroll $\Pi$ with this Jacobian, or taking the order of the scroll to be known, I proceed to show that the scroll $\Pi$ is the scroll generated by the lines each of which meets the curve $\Sigma$ three times; it will thereby appear that the order is $= 8$, and that the curve is a triple line on the scroll.

Consider a point $P'$ on $\Sigma'$, and the corresponding line $L$ of the first space: take $\theta'$ a plane in the second space; corresponding to it the cubic surface $\Theta$ in the first space. By imposing a single relation on the coefficients $(a, b, c, d)$ in the equation $ax' + by' + cz' + dw' = 0$ of the plane $\theta'$, we make it pass through the point $P'$; therefore by imposing this same single relation on the coefficients $(a, b, c, d)$ of the cubic surface $\Theta$, we make it pass through the line $L$, $\Theta$ is a cubic surface through $\Sigma$ and it is easy to see that the effect will be as above only if the line $L$ cuts the curve $\Sigma$ three times; this being so, the general cubic surface meets $L$ in three points (viz., the three intersections of $L$ with $\Sigma$), and if $\Theta$ be made to pass through a fourth point on the line $L$, it will pass through the line $L$; it thus appears that the line $L$ meets $\Sigma$ three times, and consequently that the scroll $\Pi$ is generated by the lines which meet $\Sigma$ three times.

\textbf{\S 110.}\quad The theory of a scroll so generated is considered in my ``Memoir on Skew Surfaces, otherwise Scrolls,'' [340] and [410]; viz., if $\Sigma$ is a given curve of the order $m$ and deficiency $q$, the scroll generated by the lines which meet $\Sigma$ three times is of the order $3m + 2q - 4$, having $\Sigma$ for a triple line; and the section by an arbitrary plane is a curve of that order having $3m + 2q - 4$ double points. In the present case $m = 6$, $q = 3$, and the order of the scroll is $= 18 + 6 - 4 = 20 - 4 = 8$; the section by a plane has 8 double points, and it is a unicursal curve of the order 8.

\textbf{\S 111.}\quad It may be remarked that any plane $\theta'$ meets the sextic curve $\Sigma'$ in six points; the corresponding six lines of the first space are generating lines of the scroll, lying in the plane $\Theta$ corresponding to the plane $\theta'$. Any point of $\Sigma'$ gives a line of the scroll; that is, the lines of the scroll correspond to the points of $\Sigma'$, or the scroll is related to $\Sigma'$ in the same manner that $\Sigma$ is to the scroll.

\textbf{\S 112.}\quad If the point $P'$ is on $\Sigma'$, then the corresponding line $L$ meets $\Sigma$ three times, and conversely; hence the scroll $\Pi$ (locus of the lines $L$) and the curve $\Sigma$ are related in the same manner as the curve $\Sigma'$ and the scroll $\Pi'$ (locus of the lines $L'$ corresponding to the points of $\Sigma$); viz., $\Sigma$ is a triple curve on $\Pi$, and $\Sigma'$ is a triple curve on $\Pi'$.

\textbf{\S 113.}\quad But to connect this with the Jacobian theory: the Jacobian of the first space is the locus of the points $P$ such that the first derivatives of $X$, $Y$, $Z$, $W$ (each of them a cubic function of $(x, y, z, w)$) all of them vanish; that is, it is the locus of the points such that the cubic surfaces through $\Sigma$ have a common tangent plane. If $P$ is on $\Sigma$, the cubic surfaces all of them pass through the line $L$, corresponding to the point $P'$ of $\Sigma'$; and if moreover they have a common tangent plane, then the line $L$ is on the Jacobian, that is, the scroll $\Pi$ generated by the lines $L$ lies on the Jacobian; and if the Jacobian is of the order 8, then the scroll $\Pi$ is the whole Jacobian.

\textbf{\S 114.}\quad To verify this, I consider the section by the plane $w = 0$. Writing $X$, $Y$, $Z$, $W$ for the cubic functions of $(x, y, z, w)$, and $(a, b, c, d)$ for arbitrary coefficients, the equation $aX + bY + cZ + dW = 0$ is that of any cubic surface through $\Sigma$; and if $(X, Y, Z, W)$ are each of them $= 0$, the cubic surface has the point $(x, y, z, w)$ for a double point. The section by $w = 0$ is a cubic curve $(a, b, c, d \mid X, Y, Z, W) = 0$, having the double points given by $X = 0$, $Y = 0$, $Z = 0$, $W = 0$, that is, the six intersections of $\Sigma$ by the plane $w = 0$; and the Jacobian is a surface of the order 8 having $\Sigma$ for a triple curve, or the section by $w = 0$ is a curve of the order 8, having the six points for triple points; that is, a curve of the order 8 with 6 triple points. The deficiency of a curve of the order 8 with 6 triple points is $= \frac{1}{2} \cdot 7 \cdot 6 - 6 \cdot 3 = 21 - 18 = 3$; the curve is thus unicursal, which agrees with the foregoing theory of the scroll $\Pi$.

\textbf{\S 115.}\quad Returning to the general theory, the equations of transformation are $x' : y' : z' : w' = X : Y : Z : W$; hence, reciprocally, for the transformation from the second space to the first, the equations are $x : y : z : w = X' : Y' : Z' : W'$, where $X'$, $Y'$, $Z'$, $W'$ are the determinants formed from the matrix of the coefficients of the lineo-linear equations taken in the reverse order; that is, from the matrix which expresses $(x', y', z', w')$ in terms of $(x, y, z, w)$. The whole is thus a perfectly symmetrical transformation between the two spaces; and we have in each space a sextic curve $\Sigma$ or $\Sigma'$ with 7 apparent double points, the principal system; and an octic scroll $\Pi$ or $\Pi'$, the Jacobian, having the sextic curve for a triple line.

\textbf{\S 116.}\quad The formula for the postulation of a multiple curve on a surface of the order $n$ is as follows: if the $r$-tuple curve consist of a simple curve of the order $m_1$ with $h_1$ apparent double points, of a double curve of order $m_2$ with $h_2$ apparent double points, and so on; and if moreover, the curves $m_1$, $m_2$ intersect in $k_{1,2}$ points, the curves $m_2$, $m_3$ in $k_{2,3}$ points, \&c.; then writing in general $\rho = \frac{1}{2}m(m - 1) - h$; that is, $\rho_1 = \frac{1}{2}m_1(m_1 - 1) - h_1$, $\rho_2 = \frac{1}{2}m_2(m_2 - 1) - h_2$, \&c., I find that the general condition of equivalence is
\begin{multline*}
\alpha_1 + (3n - 2)m_1 - 2\rho_1 \\
+ 8\alpha_2 + (12n - 16)m_2 - 16\rho_2 \\
+ \ldots + \alpha_r(3rn - 2r^2)m_r - 2r^2\rho_r = n^3 - 1;\\
- 5k_{1,2} - 8k_{2,3} - \ldots - (3r - 1)k_{r-1,r},
\end{multline*}
and that the general condition of postulation is
\begin{multline*}
\alpha_1 + (n + 1)m_1 - \rho_1 \\
+ \tfrac{1}{6}r(r + 1)(r + 2)\alpha_r \\
+ [\tfrac{1}{2}r(r + 1)n - \tfrac{1}{6}r(r + 1)(2r - 5)]m_r \\
- [\tfrac{1}{6}(r - 1)(r - 2)(r - 3)(r - 4) \\
\qquad\qquad + \tfrac{4}{3}r(r + 1)(2r + 1)]\rho_r \\
= \tfrac{1}{6}(n + 1)(n + 2)(n + 3) - 4;\\
- k_{1,2} - 2k_{2,3} - \ldots - \tfrac{1}{2}r(r - 1)k_{r-1,r}(s < r),
\end{multline*}
in which formulae it is however assumed that the curves have not any actual multiple points. This implies that if any one of the curves, say $m_r$, breaks up into two or more curves, the component curves are each to be simple curves, and the intersections of any two of them are to be actual double points, each counting as such, and not as apparent double points.

\textbf{\S 117.}\quad I will consider some particular cases of the formulae. First, let the $r$-tuple curve be a simple curve of the order $m_1$, with $h_1$ apparent double points; then $\rho_1 = \frac{1}{2}m_1(m_1 - 1) - h_1$, and the conditions are
\begin{align*}
\text{equivalence: } && \alpha_1 + (3n - 2)m_1 - 2\rho_1 &= n^3 - 1, \\
\text{postulation: } && \alpha_1 + (n + 1)m_1 - \rho_1 &= \tfrac{1}{6}(n + 1)(n + 2)(n + 3) - 4.
\end{align*}

\textbf{\S 118.}\quad Next let the $r$-tuple curve be a double curve of the order $m_2$, with $h_2$ apparent double points; then $\rho_2 = \frac{1}{2}m_2(m_2 - 1) - h_2$, and the conditions are
\begin{align*}
\text{equivalence: } && 8\alpha_2 + (12n - 16)m_2 - 16\rho_2 &= n^3 - 1, \\
\text{postulation: } && \alpha_2 + (n + 1)m_2 - \rho_2 &= \tfrac{1}{6}(n + 1)(n + 2)(n + 3) - 4.
\end{align*}

\textbf{\S 119.}\quad And so for a triple curve of the order $m_3$, with $h_3$ apparent double points; we have $\rho_3 = \frac{1}{2}m_3(m_3 - 1) - h_3$, and
\begin{align*}
\text{equivalence: } && 27\alpha_3 + (27n - 18)m_3 - 18\rho_3 &= n^3 - 1, \\
\text{postulation: } && \alpha_3 + (n + 1)m_3 - \rho_3 &= \tfrac{1}{6}(n + 1)(n + 2)(n + 3) - 4.
\end{align*}

\textbf{\S 120.}\quad It is to be observed that the postulation equations give at once the values of $\rho_1$, $\rho_2$, $\rho_3$, \&c.; but that the equivalence equations give also the values of $\alpha_1$, $\alpha_2$, $\alpha_3$, \&c. In particular, for a triple line ($m_3 = 1$, $h_3 = 0$, $\rho_3 = 0$), the postulation equation gives
\[
\alpha_3 = \tfrac{1}{6}(n + 1)(n + 2)(n + 3) - 4 - (n + 1) = \tfrac{1}{6}(n + 1)(n + 2)(n + 3) - (n + 5),
\]
and the equivalence equation gives
\[
27\alpha_3 + 27n - 18 = n^3 - 1, \quad\text{that is}\quad 27\alpha_3 = n^3 - 27n + 17.
\]

\textbf{\S 121.}\quad For a triple conic ($m_3 = 2$, $h_3 = 0$, $\rho_3 = 1$), we have
\begin{align*}
\text{postulation: } && \alpha_3 &= \tfrac{1}{6}(n + 1)(n + 2)(n + 3) - 4 - 2(n + 1) + 1 \\
&&&= \tfrac{1}{6}(n + 1)(n + 2)(n + 3) - 2n - 5, \\
\text{equivalence: } && 27\alpha_3 + 54n - 36 - 18 &= n^3 - 1, \\
&& 27\alpha_3 &= n^3 - 54n + 53.
\end{align*}

\textbf{\S 122.}\quad Let the triple curve be a pair of non-intersecting lines; here $m_3 = 2$, $h_3 = 0$, $\rho_3 = 1$, and the formulae are the same as for a triple conic; but we must consider the line-pair as a degenerate conic, and apply the formula with $\rho_3 = 1$. If the lines intersect, then $h_3 = 0$, but the intersection counts as an actual double point; the formula gives $\rho_3 = 1 - \frac{1}{2} = \frac{1}{2}$, which is fractional, and the case is thus excluded from the general formula.

\textbf{\S 123.}\quad Let the $r$-tuple curve consist of three right lines meeting in a point: this is an actual triple point, and the formulae do not apply. But calculating the postulation-terms by the formula, we have $m_r = 3$, $\rho_r = \frac{1}{2} \cdot 3 \cdot 2 - 0 = 3$; and the terms are
\[
[\tfrac{1}{2}r(r + 1)n - \tfrac{1}{2}r(r + 1)(2r - 5)] \cdot 3 - [\tfrac{1}{6}(r - 1)(r - 2)(r - 3)(r - 4) + \tfrac{4}{3}r(r + 1)(2r + 1)],
\]
which are
\[
= \tfrac{1}{2}r(r + 1)(3n - 4r + 4) - \tfrac{1}{6}(r - 1)(r - 2)(r - 3)(r - 4),
\]
or say
\[
= \tfrac{1}{2}r(r + 1)(3n - 4r + 4) + \tfrac{1}{6}(-r^4 + 10r^3 - 35r^2 + 50r - 24).
\]
I have found by an independent investigation that this value requires the correction
\[
+ \tfrac{1}{6}[r^4 - 8r^3 + 30r^2 - 56r + 24 + \tfrac{1}{2}\{1 - (-1)^r\}],
\]
and that the true value of the postulation is
\[
= \tfrac{1}{2}r(r + 1)(3n - 4r + 6) + \tfrac{1}{6}(2r^3 - 5r^2 - 6r + \tfrac{1}{2}[1 - (-1)^r]),
\]
viz., that this is the number of the conditions to be satisfied that a surface of the order $n$ may have for an $r$-tuple curve three given right lines meeting in a point.

\textbf{\S 124.}\quad The formula for the postulation of an $r$-tuple point is $\frac{1}{6}r(r + 1)(r + 2)$; and for the equivalence, $r^3$; but these formulae assume that there is no $r$-tuple curve through the point. If there be an $s$-tuple curve through the point ($s < r$), then the formulae require modification; and the investigation of the general formulae for the postulation and equivalence of multiple points and curves is a subject of considerable interest.

\clearpage

%% ========================================================================
%% PAPER 448: Note on the Cartesian with Two Imaginary Axial Foci
%% ========================================================================
\section*{\textbf{448.}}
\addcontentsline{toc}{section}{448. Note on the Cartesian with Two Imaginary Axial Foci}

\begin{center}
\textbf{NOTE ON THE CARTESIAN WITH TWO IMAGINARY AXIAL FOCI.}
\end{center}

[From the \textit{Proceedings of the London Mathematical Society}, vol.~iii (1869--1871), pp.~181, 182. Read June 9, 1870.]

Let $A$, $A'$, $B$, $B'$ be a pair of points and antipoints viz.;
\begin{itemize}
\item[] $(A, A')$ the two imaginary points, coordinates $(+\beta i, 0)$, $(-\beta i, 0)$;
\item[] $(B, B')$ the two real points, coordinates $(0, +\beta)$, $(0, -\beta)$;
\end{itemize}
and write $\rho$, $\rho'$, $\sigma$, $\sigma'$ for the distances of a point $(x, y)$ from the four points respectively:
\[
\rho = \sqrt{(x - i\beta)^2 + y^2},\quad \sigma = \sqrt{x^2 + (y + \beta)^2},
\]
\[
\rho' = \sqrt{(x + i\beta)^2 + y^2},\quad \sigma' = \sqrt{x^2 + (y - \beta)^2}.
\]

We have
\[
\rho^2 + \rho'^2 = 2x^2 + 2y^2 - 2\beta^2 = \sigma^2 + \sigma'^2 - 4\beta^2,
\]
\[
\rho\rho' = \sqrt{(x^2 + \beta^2 + y^2)^2 - 4\beta^2 y^2} = \sigma\sigma';
\]
and thence
\begin{align*}
(\rho + \rho')^2 &= (\sigma + \sigma')^2 - 4\beta^2, \\
(\rho - \rho')^2 &= (\sigma - \sigma')^2 - 4\beta^2;
\end{align*}
or say
\begin{align*}
\rho + \rho' &= \sqrt{(\sigma + \sigma')^2 - 4\beta^2}, \\
i(\rho - \rho') &= \sqrt{4\beta^2 - (\sigma - \sigma')^2}.
\end{align*}

The equation of a Cartesian having the two imaginary axial foci $A$, $A'$ is
\[
(p + qi)\rho + (p - qi)\rho' + 2\sqrt{q}c = 0;
\]
viz., this is
\[
p\sqrt{(\sigma + \sigma')^2 - 4\beta^2} + q\sqrt{4\beta^2 - (\sigma - \sigma')^2} + 2\sqrt{q}c = 0,
\]
or, what is the same thing, it is
\[
p\sqrt{(\sigma + \sigma')^2 + 4\beta^2} + q\sqrt{4\beta^2 - (\sigma - \sigma')^2} + 2c\sqrt{q} = 0,
\]
which is the equation expressed in terms of the distances $\sigma$, $\sigma'$ from the non-axial real foci $B$, $B'$. Of course, the radicals are to be taken with the signs $+$. This equation gives, however, the Cartesian in combination with an equal curve situate symmetrically therewith in regard to the axis of $y$.

The distances $\sigma$, $\sigma'$ may conveniently be expressed in terms of a single variable parameter $\theta$; in fact, we may write
\[
\pm\rho = \sqrt{(\sigma + \sigma')^2 - 4\beta^2} = k + k'\theta,
\]
that is
\[
4\rho^2 - (\sigma - \sigma')^2 = k'^2 - (k'\theta)^2 = k'^2(1 - \theta^2);
\]
And therefore
\[
\sigma + \sigma' = \pm\frac{1}{2}\sqrt{4\beta^2 + (k + k'\theta)^2},\quad
\sigma - \sigma' = \pm\frac{1}{2}\sqrt{4\beta^2 - (k - k'\theta)^2};
\]
so that, assigning to $\theta$ any given value, we have $\sigma$, $\sigma'$, and thence the position of the point on the curve. We may draw the hyperbola $y^2 = 4\beta^2 + c^2 x^2$, and the ellipse $y^2 = 4\beta^2 - c^2 k'^2 x^2$; and then measuring off in these two curves respectively the ordinates which belong to the abscissae $k + \theta$ for the hyperbola, $k - \theta$ for the ellipse, we have the values $\sigma + \sigma'$ and $\sigma - \sigma'$, which determine the point on the curve. Considering $k$, $p$, $q$, $c$ as disposable quantities, the conics may be any ellipse and hyperbola whatever, having a pair of vertices in common; and the complete construction is,---From the fixed point in the axis of $x$, measure off in opposite directions the equal distances $KM$, $KN$, and take
\begin{align*}
M &= \sigma + \sigma' \text{ the ordinate at in the hyperbola}, \\
N &= \sigma - \sigma' \quad\;\,\text{the ordinate at in the ellipse},
\end{align*}
where $\sigma$, $\sigma'$ denote the distances of the required point from the fixed points $B$ and $B'$ respectively, the distance of each of these from the origin being $= \beta$ the common semi-axis. We may imagine $N$ travelling from one extremity of the $x$-axis of the ellipse to the other, the value of $\sigma + \sigma'$ will be real and greater than $BB'$, that of $\sigma - \sigma'$ real and less than $BB'$, and the point $(\sigma, \sigma')$ will be real. The construction gives, it will be observed, the two symmetrically situated curves.

The $x$-semi-axis of the ellipse is $\frac{1}{c}\sqrt{2\beta}$, and the form of the curve depends chiefly on the value of the ratio $k : k'\sqrt{p}$; or, what is the same thing, $k'^2 : 2\beta q$. We see, for instance, that, in order that the curve may meet the axis of $y$ in two real points between the foci, the value $\theta = -k$ must give a real value of $\sigma - \sigma'$; viz., that we must have $4\beta^2 > 4k'^2$; that is, $\frac{1}{q}c^2 > k'^2$, or $k'^2 c^2 > q$. If $k$ has this value, viz., $k = \frac{1}{c}\sqrt{2\beta} = \frac{1}{c}\sqrt{\text{semi-axis}}$, the curve touches the axis of $y$ at the origin; if $k < \frac{1}{c}\sqrt{\text{semi-axis}}$, the curve cuts the axis of $y$ in two real points between the foci; if $k > \frac{1}{c}\sqrt{\text{semi-axis}}$, the curve does not cut the axis of $y$ between the foci.

\clearpage

%% ========================================================================
%% PAPER 449: Sketch of Recent Researches Upon Quartic and Quintic Surfaces
%% ========================================================================
\section*{\textbf{449.}}
\addcontentsline{toc}{section}{449. Sketch of Recent Researches Upon Quartic and Quintic Surfaces}

\begin{center}
\textbf{SKETCH OF RECENT RESEARCHES UPON QUARTIC AND QUINTIC SURFACES.}
\end{center}

[From the \textit{Proceedings of the London Mathematical Society}, vol.~III (1869--1871), pp.~186--195. Read Nov.~10, 1870.]

The classification of quartic surfaces is even as to its highest divisions incomplete; and it is by no means easy to make it at once exhaustive and precise; an enumeration of all the prima facie possible cases would include forms which do not really exist. Thus the singular curve (if any) is of the order 1, 2, or 3; but in the case where the order is $= 3$, the curve, as is at once evident, cannot be a plane cubic, nor (among other excluded forms) a system of three non-intersecting lines. And certain forms of the singular curve, e.g.\ all but one of the admissible forms of a curve of the order 3, make the surface to be a scroll, so that, if (as is convenient) we wish to separate the scrolls, certain forms otherwise admissible must be excluded. The expression ``singular'' means double or cuspidal, or refers to a higher singularity, but the cases of higher singularity are very special. I will, at the cost of some inaccuracy, use the expression ``nodal'' as meaning, in general, double, but as including the signification ``cuspidal''; and, if there are any cases of higher singularity, as extending to cases of higher singularity: and I provisionally arrange the non-scrolar quartic surfaces as follows:
\begin{enumerate}[label=\arabic*.]
\item Without a nodal curve.
\item With a nodal line.
\item With a nodal conic, or line-pair (pair of intersecting lines).
\item With three nodal lines (not in the same plane) meeting in a point.
\end{enumerate}

The references, by the name of the author, and number (if any) of his paper, are to the subjoined list of Memoirs.

As to the scrolls, we have Cayley (3) and (4), and Cremona; the division into 12 species is, I believe, complete: see post, the remarks upon Schwarz's paper on quintic scrolls.

As regards the non-scrolar surfaces:

\textbf{1.}\quad Without a singular curve. The surface may be without a cnicnode (conical point), or it may have any number of cnicnodes up to 16, Cayley (7): the cases of singularity higher than a cnicnode are probably very numerous, but they have been scarcely at all examined. The memoir just referred to relates chiefly to the several cases of not more than 10 nodes; the cases of 11, 12, 13, 14, 15, 16 nodes are considered incidentally, Kummer (2), but it was not the object of his paper to make an enumeration, and there may be cases which are not considered; the discussion of the cases considered is very full and interesting. The case of 16 nodes is also considered, Kummer (1). As to the surface with 16 nodes, it is to be remarked that the wave-surface, or generally the surface obtained by the homographic deformation of the wave-surface called, Cayley (1), the ``tetrahedroid'' is a special form of surface with 16 nodes: its relation to the general surface is explained, Cayley (2).

\textbf{2.}\quad Quartic surface with nodal line: considered incidentally, Clebsch (2) and (3). There are through the nodal line 4 (in general different) quadric cones, having each of them a contact of the second order with the surface along the nodal line; any plane through the nodal line meets the surface in a conic having its two intersections with the nodal line for a pair of conjugate points in regard to the section by the plane of any one of the four quadric cones. And in particular, the tangent plane at any point of the nodal line is a plane through the nodal line which is a double tangent plane of the surface; viz., the tangent plane at any point of the nodal line meets the surface in the nodal line twice and in a conic touching the nodal line at the point. The cases of 1, 2, 3, 4 cnicnodes on the surface are considered, Korndorfer.

\textbf{3.}\quad Quartic surfaces with nodal conic. Such a surface may be without cnicnodes, or it may have 1, 2, 3, or 4 cnicnodes. The theory depends on that of the cubic surface; viz., taking any one of the 16 lines of the cubic surface, the scroll generated by the line in its motion along the line is a quartic surface having the line for nodal conic; that is, the locus of a line which meets each of two given non-intersecting lines, and also a given cubic curve, is a quartic surface having the line-pair (the two given lines) for a nodal conic and to each other. And the ground hereof appears, Geiser; viz., it is shown that the quartic surface with the nodal conic, is rationally transformable into a cubic surface, the 16 lines and the nodal conic corresponding respectively to the 16 lines and the conic of the cubic surface.

The several cases of 1, 2, 3, and 4 cnicnodes are considered, Korndorfer.

In the case where the nodal conic is the circle at infinity, the surfaces have been termed ``anallagmatic'' (perhaps ``bicircular'' would be a more convenient name), and a great deal has been written upon these surfaces by Moutard, Clifford, and others. Such a surface may of course have 1, 2, 3, or 4 cnicnodes; these surfaces, viz.\ the cnicnodal anallagmatics, in fact arise from the inversion of a quadric surface by the method of reciprocal radius vectors (that is, by the change of $x$, $y$, $z$ into $\frac{x}{r^2}$, $\frac{y}{r^2}$, $\frac{z}{r^2}$); the centre of inversion is a node on the quartic surface. If the quadric surface is a cone, there is another node, the inverse point of the vertex; if the quadric surface is one of revolution, there are two other nodes; and if it is a cone of revolution, there are three other nodes viz., in all, four nodes. The last-mentioned surface is, or includes, the Cyclide; viz., this is a quartic surface having the circle at infinity for a nodal curve, and having besides four nodes, which are a system of skew antipoints. The surface was first considered by Dupin (\textit{Applications de G\'eom\'etrie}, 1822); it has been discussed and its theory completed by Maxwell (Quart.\ Math.\ Journ.\ t.~ix.\ pp.~111--126, 1868) who has also considered the particular cases; and it was considered under a more general point of view by Moutard, and more recently by Cremona (in a Memoir ``Sulle transformazioni razionali nello spazio,'' Annali di Matematica, t.~V.\ 1868) who showed that the cyclide is the inverse of a quadric surface, and is also the envelope of a sphere which cuts at right angles each of four given spheres; or, what is the same thing, that the surface is the locus of a point whose distances from the centres of the four given spheres satisfy a linear equation; that is, the centres of the four given spheres are the foci of the surface. This theory includes of course the particular case where the four spheres reduce themselves to points, viz., the cyclide becomes the ``sphero-quartic'' or cyclide with four concyclic foci, the locus of a point the sum of whose distances from two given points is in a constant ratio to the sum of its distances from two other given points; or, more generally, the locus of a point such that $l\rho + m\rho' + n\rho'' = 0$; viz., this is the ``Cartesian,'' which may therefore be regarded as a particular case of the cyclide. In reference to the plane sections of the conic torus and its various particular cases, see De la Gournerie (2); the ordinary torus has been the subject of numerous papers by Darboux and others, and possesses very interesting properties.

In connexion with the foregoing, I speak of the surfaces having a cuspidal conic: the general case is briefly referred to, Cayley (6); viz., this is the surface (AA) the equation of which is $V^2 - a^2 u^2 = 0$, and which it is shown has a reciprocal of the order 6. A special case is the surface (AB) having a reciprocal of the order 3; viz., the quartic surface is here the reciprocal of the cubic surface XX = $12 - f^2 g$, Cayley (5). And it appears from the memoir last referred to, that there is another cubic surface, XVII = $12 - 25\beta - C_2$, the reciprocal of which is a quartic surface having a cuspidal conic. But the theory of the quartic surfaces with a cuspidal conic has been hardly at all considered.

I do not know that anything has been done in regard to the quartic surfaces where the nodal conic be a line-pair (pair of intersecting lines), a form which is admissible, but which probably belongs to the theory of surfaces with a nodal line.

\textbf{4.}\quad Quartic surface with three nodal lines (not in the same plane) meeting in a point. I have not met with any mention of such surfaces.

\vspace{0.5em}
\textit{Quartic Scrolls.}

As regards the scrolls, we have, as already mentioned, the division into 12 species, each of them a rational scroll, viz., the sections by a plane being in every case unicursal curves. The division was made according to the nature of the singularities (nodal curve, and cnicnodes, if any); but Schwarz's memoir effects the division according to the value of the ``deficiency'' of the surface, viz., the plane section is a quartic curve, and the deficiency of the surface is the deficiency of this quartic curve. Now the nodal curve of the quartic scroll is a line, and the plane section has accordingly a node; the plane section is thus a quartic curve with a node, and the deficiency is thus $= 1$ or $= 0$, according as the quartic curve is or is not a unicursal curve. But the plane section of a scroll is in every case a unicursal curve; and the deficiency is thus $= 0$ for a quartic scroll with a nodal line; for a quartic scroll without a nodal line the deficiency is $= 1$. But the case of a quartic scroll without a nodal line is a scroll which has the line at infinity for a nodal line, that is, a cylindric quartic scroll. The foregoing relates of course to the scrolls above referred to; viz., for a quartic scroll the deficiency is either 0 or 1; and of the 12 species, there are 10 for which the deficiency is $= 0$ (or which are unicursal), and 2 for which the deficiency is $= 1$. And this is the principle of classification in Schwarz's memoir; viz., for a quintic scroll the deficiency is $= 0$, 1, or 2; the number of species established being 10, 4, and 1 for these deficiencies respectively.

\vspace{0.5em}
\textit{List of Memoirs.}

\textbf{Cayley.}\quad
1. Sur la surface des ondes. Liouv.\ t.~xi.\ (1846), pp.~291--296, [47].

2. Sur un cas particulier de la surface du quatrieme ordre avec seize points singuliers. Crelle, t.~LXV.\ (1866), pp.~284--290, [356].

3. Second Memoir on Skew Surfaces, otherwise Scrolls. Phil.\ Trans., vol.~CLIV.\ (1864), pp.~559--577, [340].

4. Third Memoir on Skew Surfaces, otherwise Scrolls. Phil.\ Trans., vol.~CLIX.\ (1869), pp.~111--126, [410].

5. Memoir on Cubic Surfaces. Phil.\ Trans., vol.~CLIX.\ (1869), pp.~231--326, [412].

6. On the Quartic Surfaces $(U, V, W)^2 = 0$. Quart.\ Math.\ Journ.\ vol.~x.\ (1868), pp.~24--34; and vol.~xi.\ (1870), pp.~15--25, and pp.~111--113.

7. A Memoir on Quartic Surfaces. Proc.\ Lond.\ Math.\ Soc., vol.~III.\ (1870), pp.~19--69, [445].

\textbf{Clebsch.}\quad
1. Ueber die Fl"achen vierter Ordnung welche eine Doppelcurve zweiten Grades besitzen. Crelle, t.~LXIX.\ (1868), pp.~355--358.

2. Intorno alla rappresentazione di superficie algebriche sopra un piano. Atti del R.~Ist.~Lomb.\ (12 Nov.~1868), 13~p.

3. Ueber die Abbildung algebraischer Fl"achen, insbesondere der vierten und f"unften Ordnung. G"ottingen (1868), 42~p.;
``dritter Ordnung,'' Crelle, t.~LXV.\ (1866), pp.~359--380.

4. Ueber die ebene Abbildung der geradlinigen Fl"achen vierter Ordnung. G"ottingen (1868), 14~p.

5. Ueber die Abbildung einer Classe von Fl"achen f"unfter Ordnung. G"ott.~Abh., t.~XV. (1868).

\textbf{Cremona.}\quad
Sopra alcune superficie rationali. Atti (1868).

\textbf{De la Gournerie.}\quad
1. Note sur la cyclide. Comptes Rendus, t.~LXVII. (1868), pp.~1021--1024.

2. 'Etudes sur la cyclide. Journ.~de l''Ecole Polytechnique, t.~XXXV. (1869), pp.~81--94.

\textbf{Geiser.}\quad
Ueber die Fl"achen vierten Grades welche eine Doppelcurve zweiten Grades haben. Crelle, t.~LXX.\ (1869), pp.~249--257.

\textbf{Korndorfer.}\quad
[There was only a general reference; the several memoirs are:

1. Die Abbildung einer Fl"ache vierter Ordnung mit einer Doppelcurve zweiten Grades und einem oder mehreren Knotenpunkten. Math.~Ann., t.~I.\ (1869), pp.~592--626.

2. Fortsetzung dieses Aufsatzes. t.~II.\ (1870), pp.~41--64.

3. Ueber diejenigen Raumcurven deren Coordinaten sich als rationale Functionen eines Parameters darstellen. t.~III.\ (1871), pp.~415--423.

4. Die Abbildung einer Fl"ache vierter Ordnung mit zwei sich schneidenden Doppelgeraden. t.~III.\ (1871), pp.~496--522.

5. Die Abbildung einer Fl"ache vierter Ordnung mit einer Doppelcurve zweiten Grades welche aus zwei sich schneidenden unendlich nahen Geraden besteht. t.~IV.\ (1871), pp.~117--134.]

\textbf{Kummer.}\quad
1. {Surfaces of the fourth order with sixteen conical points.} Berl.~Monatsb.\ (1864), pp.~246--260, and 495--499.

2. Ueber die algebraischen Strahlensysteme, insbesondere "uber die der ersten und zweiten Ordnung. Berl.~Abh.\ (1866), pp.~1--120.

3. Ueber die Fl"achen vierten Grades auf welchen Schaaren von Kegelschnitten liegen. Berl.~Monatsb.\ (July, 1868). Crelle, t.~LXIV.\ (1864), pp.~66--76.

\textbf{Maxwell.}\quad
On the Cyclide. Quart.\ Math.\ Journ., t.~ix.\ (1867), pp.~111--126.

\textbf{Schwarz.}\quad
Ueber die geradlinigen Fl"achen f"unften Grades. Crelle, t.~LXVII.\ (1867), pp.~23--57.

\clearpage

%% ========================================================================
%% PAPER 450: Note on the Theory of the Rational Transformation Between Two Planes
%% ========================================================================
\section*{\textbf{450.}}
\addcontentsline{toc}{section}{450. Note on the Theory of the Rational Transformation Between Two Planes, and on Special Systems of Points}

\begin{center}
\textbf{NOTE ON THE THEORY OF THE RATIONAL TRANSFORMATION BETWEEN TWO PLANES, AND ON SPECIAL SYSTEMS OF POINTS.}
\end{center}

[From the \textit{Proceedings of the London Mathematical Society}, vol.~iii (1869--1871), pp.~196--198. Read December 8, 1870.]

In Prof.~Cremona's theory of the transformation of plane curves, the fundamental equations are taken to be
\begin{align}
\alpha_1 + 4\alpha_2 + 9\alpha_3 + \ldots &= n^2 - 1 \tag{1},\\
\alpha_1 + 3\alpha_2 + 6\alpha_3 + \ldots &= \tfrac{1}{2}(n^2 + 3n) - 2 \tag{2};
\end{align}
and from these we have as a consequence
\[
\alpha_1 + 3\alpha_2 + \ldots = \tfrac{1}{2}(n-1)(n-2) \tag{3};
\]
viz., the first equation expresses that any two curves of the system intersect in a single variable point; the second, that the curves form a r\'eseau, or system containing two arbitrary parameters; and the third, that the curves are unicursal.

In the equivalent theory of the rational transformation between two planes, as given in my ``Memoir on the Rational Transformation between Two Spaces,'' [447], we have the equation (1); but instead of the equation (2), it would prim\^a facie appear to be sufficient if we had the inequation
\[
\alpha_1 + 3\alpha_2 + 6\alpha_3 + \ldots > \tfrac{1}{2}(n^2 + 3n) - 2;
\]
but on the ground there explained, the case
\[
\alpha_1 + 3\alpha_2 + 6\alpha_3 + \ldots < \tfrac{1}{2}(n^2 + 3n) - 2
\]
is excluded, and we thus have the equation (2), giving with (1) the equation (3).

I believe the better course is to assume (1) and (3) as the fundamental equations, from them deducing (2); and we thus also get over a difficulty presently referred to, but which did not occur to me when the memoir was written.

In fact, starting with the equations $x' : y' : z' = X : Y : Z$ (which are to give $x : y : z = X' : Y' : Z'$), we have in the first instance the equation (1). Moreover, establishing for $x', y', z'$ a linear equation $ax' + by' + cz' = 0$, we have corresponding hereto a curve $aX + bY + cZ = 0$, and the coordinates $x, y, z$ of a point on this curve are proportional to $X' : Y' : Z'$; that is, substituting for $z'$ the value $-(ax' + by')/c$, they are proportional to rational and integral (homogeneous) functions of $(x', y')$, that is, to rational and integral functions of the single parameter $x' : y'$; wherefore the curve $aX + bY + cZ = 0$ is unicursal; whence the equation (3). The like change may be made in the theory of the rational transformation between two spaces; and it is in this case a more important one.

The difficulty is as follows: It is not self-evident that we are at liberty to assume
\[
\alpha_1 + 3\alpha_2 + 6\alpha_3 + \ldots > \tfrac{1}{2}(n^2 + 3n) - 2
\]
for imagine that we had a system of $(\alpha_1, \alpha_2, \alpha_3, \ldots)$ points, such that $\alpha_1 + 4\alpha_2 + \ldots$ being $= n^2 - 1$, and $\alpha_1 + 3\alpha_2 + \ldots$ being $> \tfrac{1}{2}(n^2 + 3n) - 2$, the points were such that the conditions in question (viz., the condition that the curve passes once through each of the points $\alpha_1$, twice through each of the points $\alpha_2$, \&c.) were not independent, but equivalent to a less number of conditions. Suppose, for instance, a quintic curve, $n = 5$; and consider the system of 12 points and 2 double points (12 points through which the curve passes once, and 2 points through which it passes twice); we have here
\[
\alpha_1 + 4\alpha_2 = 12 + 8 = 20 = n^2 - 1,\quad \alpha_1 + 3\alpha_2 = 12 + 6 = 18 = \tfrac{1}{2}(n^2 + 3n) - 2;
\]
so that the equations (1) and (2) are satisfied. But there is here no reduction; the 18 conditions are independent. But we may, it would appear, have 12 points and 2 double points such that, for a quintic curve passing once through each of the 12 points and twice through each of the 2 points, the number of conditions actually imposed (instead of being $12 + 3 \cdot 2 = 18$) is $= 17$.'' The construction is as follows: viz., starting with the 2 points and any 10 points, we may draw a quartic passing twice through the first of the 2 points, once through the second of them, and through the 10 points; and another quartic passing twice through the second of the 2 points, once through the first of them, and through the 10 points: the two quartics intersect in the 2 points each twice, in the 10 points, and in 2 new points, forming, with the 10 points, a system of 12 points; and the first-mentioned 2 points and the 12 points form the system in question.

A more complicated case, $\alpha_1 = 10$, $\alpha_2 = 6$, $\alpha_3 = 1$, occurs in Dr Noether's paper, ``Ueber Fl\"achen, welche Schaaren rationaler Curven besitzen,'' [Math.\ Ann., t.~iii.\ (1871), pp.~161--227]. Except these two, I do not know any other case of a special system for which $\alpha_2, \alpha_3, \ldots$ are not all $= 0$; the investigation of such systems would, I think, be very interesting.

[A concluding paragraph of seven lines gave some corrections to the ``Memoir on the Rational Transformation between Two Spaces,'' 447, which corrections are made in the present reprint of that paper.]

\clearpage

%% ========================================================================
%% PAPER 451: A Second Memoir on Quartic Surfaces
%% ========================================================================
\section*{\textbf{451.}}
\addcontentsline{toc}{section}{451. A Second Memoir on Quartic Surfaces}

\begin{center}
\textbf{A SECOND MEMOIR ON QUARTIC SURFACES.}
\end{center}

[From the \textit{Proceedings of the London Mathematical Society}, vol.~iii (1869--1871), pp.~198--202. Read December 8, 1870.]

In my Memoir on Quartic Surfaces, ante pp.~19--69, [445], although remarking (see No.~79) that the identification was not completely made out, I tacitly assumed that the symmetroid and the decadianome (each of them a quartic surface with 10 nodes) were in fact identical. There is yet a good deal which I cannot completely explain; but the truth appears to be, that the decadianome includes two cases of coordinate generality, say the sextic decadianome, and the bicubic decadianome $=$ symmetroid: viz., in the first of these the circumscribed cone, having for vertex any one of the 10 nodes, is a proper sextic cone with 9 double lines; in the second it is a system of two cubic cones, intersecting, of course, in 9 lines, which are double lines of the aggregate sextic cone: or, in the notation of the Table No.~11, in the case of the sextic decadianome, the circumscribed cones are each of them $6_9$; in that of the bicubic decadianome $=$ symmetroid, they are each of them $(3, 3)$. We thus arrive at a very remarkable system of 10 points in space, viz., giving the name ``ennead'' to any 9 points in piano, which are the intersections of two cubic curves, or to any 9 lines through a point which are the intersections of two cubic cones; the 10 points in space are such that, taking any one whatever of them as vertex, and joining it with the remaining 9, the 9 lines form an ennead. I purpose in the present short Memoir to consider the theories in question; the paragraphs are numbered consecutively with those of the Memoir on Quartic Surfaces.

\subsection*{Plane Sextic Curve with 9 Nodes.}

\textbf{\S 110.}\quad A sextic curve contains 27 constants; and the number of conditions to be satisfied in order that a given point may be a node is $= 3$. Hence it would at first sight appear that the curve could be found so as to have 9 given nodes; this would be $9 \times 3 = 27$ conditions, or the curve would be completely determinate. But observe that through the 9 given points we have a determinate cubic curve $U = 0$; we have therefore $U^2 = \text{a sextic curve}$, and the only sextic curve with the 9 given nodes; that is, there is not in a proper sense any sextic curve with the 9 given nodes. The number of given nodes is thus $= 8$ at most.

\textbf{\S 111.}\quad The sextic curve with 8 given nodes should contain $27 - 3 \cdot 8 = 3$ constants. We may through the 8 given points draw the two cubics $P = 0$, $Q = 0$; and we have then $(a, b, c \mid P, Q)^2 = 0$, a bicubic, or improper sextic curve having the 8 nodes, and also a ninth node, viz., the remaining point of intersection of the two cubic curves, or say the remaining point of the ennead. Hence if $V = 0$ be any particular sextic curve having the 8 given nodes, we have
\[
(a, b, c \mid P, Q)^2 + \theta V = 0
\]
a proper sextic curve having the 8 given nodes; and this, as containing the right number ($= 3$) of constants, will be the general sextic curve having the 8 given nodes.

\textbf{\S 112.}\quad There will be a ninth node if $\theta = 0$; viz., the curve is then $(a, b, c \mid P, Q)^2 = 0$, a bicubic, or improper sextic curve, having for nodes the 9 points of the ennead. Observe that the ninth node is here a point completely and uniquely determined by means of the given 8 nodes. Moreover the number of constants is $= 2$, so that we have here a general (improper) solution of the question of finding a sextic curve with 9 nodes, 8 of them given.

\textbf{\S 113.}\quad But if $\theta$ is not $= 0$, then the ninth node must be a point on the curve $J(P, Q, V) = 0$; viz., this is a curve of the order 9, determined by means of the given 8 points; say it is the ``dianodal curve'' of these 8 points, and, as is easy to see, it has each of these 8 points for a node. The ninth node of the sextic may be any point whatever on the dianodal curve; and regarding it as a given point, the sextic will still contain 1 constant; that is, we have the general solution of the problem of finding a sextic curve with 9 nodes, 8 of them given, and the 9th a given point on the dianodal curve.

\textbf{\S 114.}\quad So long as the 8 points are arbitrary, the dianodal curve does not pass through the 9th point of the ennead, and the two cases above considered are mutually exclusive. It will be noticed how closely analogous this theory of the plane sextic with 9 nodes, is to that of the quartic surface with 8 nodes.

\textbf{\S 115.}\quad Of course, instead of the plane sextic curve, we may have a sextic cone; such a cone has at most 8 given double lines; and if there be a 9th double line, then there are the two cases of coordinate generality; viz., (1), the new double line is the ninth line of the ennead, the cone being in this case not a proper sextic cone, but a bicubic cone; (2), the new double line may be any line through the vertex which lies on the dianodal cone of the 8 given double lines.

\subsection*{Each circumscribed cone of the Symmetroid is $(3, 3)$.}

\textbf{\S 116.}\quad Using $(x, y, z, w)$ as current coordinates of a point of the symmetroid, I take $S$, $T$, $U$, $V$ to be quadric functions of the coordinates $(\alpha, \beta, \gamma, \delta)$; the equation of the symmetroid is therefore given by
\[
xS + yT + zU + wV = \text{cone},
\]
and the nodes thereof are determined by
\[
xS + yT + zU + wV = \text{plane-pair}.
\]
Suppose that a node is $(x = 0, y = 0, w = 0)$; the condition for this is $V = \text{plane-pair}$; and we may without loss of generality write $V = \gamma^2 + \delta^2$. Hence, putting for shortness $\Omega = xS + yT + zU$, that is a quadric function $(a, b, c, d, f, g, h, l, m, n \mid \alpha, \beta, \gamma, \delta)^2$, wherein the coefficients $a, b, \ldots$ are arbitrary linear functions of $(x, y, z)$, but not containing $w$, the equation of the symmetroid is given by
\[
\Omega + w(\gamma^2 + \delta^2) = \text{cone}.
\]

\textbf{\S 117.}\quad It follows that the equation is
\[
\begin{vmatrix}
a, & h, & g \\
h, & b, & f \\
g, & f, & c + w
\end{vmatrix}
= 0;
\]
viz., this is
\[
V^2 + w(\partial_c V + \partial_d V)V + w^2 \partial_c V \cdot \partial_d V = 0,
\]
where $V$ denotes the foregoing determinant, writing therein $w = 0$. Or, observing that $V$ contains $c$ and $d$ each only linearly, the equation may be written
\[
V^2 + w(\partial_c + \partial_d)V + w^2 \partial_c \partial_d V = 0,
\]
which is a quartic surface having, as it should have, the point $(0, 0, 0)$ for one of its ten nodes.

\textbf{\S 118.}\quad The equation of the circumscribed cone is
\[
(\partial_c V - \partial_d V)^2 + 4(\partial_c V \cdot \partial_d V - V \cdot \partial_c \partial_d V) = 0.
\]
But we have identically
\[
\partial_c V \cdot \partial_d V - (\partial_c \partial_d V)^2 = V \cdot \partial_c \partial_d V;
\]
so that the equation is
\[
(\partial_c V - \partial_d V)^2 + (\partial_{cd} V)^2 = 0;
\]
a sextic cone breaking up into the two cubic cones
\[
\partial_c V - \partial_d V + i\partial_{cd} V = 0,\quad \partial_c V - \partial_d V - i\partial_{cd} V = 0,
\]
so that the cone is $(3, 3)$. And since clearly the point $(0, 0, 0)$ may be regarded as representing any one whatever of the 10 nodes, it follows that for any node whatever of the symmetroid, the circumscribed cone is $(3, 3)$, so that, as stated above, bicubic decadianome $=$ symmetroid.

\subsection*{Deductions from the foregoing theory.}

\textbf{\S 119.}\quad Referring to No.~85 of the original memoir, it appears that, with 6 given points as nodes, we can actually find for the symmetroid an equation containing 6 constants. I cannot discover any ground for doubting that 3 of these may be determined so as to give to the symmetroid a seventh given node; and I therefore assume that with 7 given points as nodes, an equation can be found with 3 constants. The symmetroid is certainly not octadic, hence the eighth node must lie on the dianodal surface of the 7 given points. I can discover no ground for doubting but that two of the constants may be determined so that the eighth node shall be any given point whatever on the dianodal surface of the 7 points; and (this being so) that further the remaining constant may be determined so that the ninth node shall be any given point whatever on the dianodal curve of the 8 points. But if all this be so, the consequence is very remarkable; the tenth node is not any one whatever of the 22 dianodal centres of the 9 points, but it is a uniquely determinate ``enneadic centre,'' viz., we must have the following theorem:

\textbf{\S 120.}\quad ``Take any 7 points; an eighth point at pleasure on the dianodal surface of the 7 points; a ninth point at pleasure on the dianodal curve of the 8 points. In the system of 9 points so determined, take any one as vertex, and joining it with the remaining 8, construct the ninth line of the ennead. Performing this construction with each of the 9 points successively as vertex, we obtain 9 lines passing through the 9 points respectively. These 9 lines meet in a point which is the `enneadic centre' of the 9 points: and further, the 10 points form a completely symmetrical system, so that each one of them is the enneadic centre of the remaining 9.''

\textbf{\S 121.}\quad Assuming that the 9 lines do intersect so as to give rise to an enneadic centre, there is no difficulty in conceiving that the loci, which by their intersection determine the dianodal centres, do each of them pass through the enneadic centre; so that this enneadic centre counts once or more among the dianodal centres, and the number of proper dianodal centres, instead of being $= 22$, will be suppose $= 22 - \omega$; and if, further, the 9 points, together with the enneadic centre, are the nodes of a symmetroid, but the 9 points together with any one of the $22 - \omega$ dianodal centres are the nodes of a sextic decadianome, then we must also have as follows:

\textbf{\S 122.}\quad ``Considering any 9 points as above; taking any one as vertex, and joining it with the remaining 8, these 8 lines determine a dianodal ninthic cone. We have thus 9 dianodal cones, which cones pass all of them through the same $22 - \omega$ points.''

\textbf{\S 123.}\quad I am not able to verify these theorems \textit{a posteriori}. It appears to me that the theorem in regard to the enneadic centre subsists for a system of 9 points such as referred to in the statement; but that if by possibility the statement be too general, the theorem must, at all events, subsist for a more special system of 9 points; and that there certainly exist systems of 10 points, such that each 9 of the points have as an enneadic centre the tenth point. (I have since ascertained that if a quartic surface with 10 nodes has a single node $(3, 3)$, the surface is a symmetroid; whence, by what precedes, the remaining nine nodes are each of them $(3, 3)$. Added 25 March, 1871.)

\textbf{\S 124.}\quad I notice, as a subject of investigation, the following system of correspondence viz., given any 8 points in space: then to every point in space corresponds a line through this point, viz., the ninth line of the ennead obtained by joining the point with the 8 given points respectively; and to each line in space a point or points on the line, viz., the point or points for each of which the line is the ninth line of the ennead obtained by joining the point with the eight given points respectively.

\clearpage

%% ========================================================================
%% PAPER 452: On an Analytical Theorem from a New Point of View
%% ========================================================================
\section*{\textbf{452.}}
\addcontentsline{toc}{section}{452. On an Analytical Theorem from a New Point of View}

\begin{center}
\textbf{ON AN ANALYTICAL THEOREM FROM A NEW POINT OF VIEW.}
\end{center}

[From the \textit{Proceedings of the London Mathematical Society}, vol.~iii (1869--1871), pp.~220, 221. Read February 9, 1871.]

The theorem is a well-known one, derived from the equation
\[
(az^2 + 2bz + c)w^2 + 2(a'z + 2b'z + c')w + a''z^2 + 2b''z + c'' = 0;
\]
viz., considering this equation as establishing a relation between the variables $z$ and $w$, and writing it in the forms
\[
\Omega = Aw^2 + 2Bw + C = A'z^2 + 2B'z + C' = 0,
\]
(where, of course, $A$, $B$, $C$ are quadric functions of $z$, and $A'$, $B'$, $C'$ quadric functions of $w$,) we have
\[
0 = \frac{dw}{\sqrt{B'^2 - A'C'}} + \frac{dz}{\sqrt{B^2 - AC}},\quad
\Omega = (Aw + B)\,dw + (A'z + B')\,dz;
\]
but in virtue of the equation $\Omega = 0$, we have $Aw + B = \sqrt{B^2 - AC}$, and $A'z + B' = \sqrt{B'^2 - A'C'}$, and the differential equation thus becomes
\[
\frac{dw}{\sqrt{B'^2 - A'C'}} = \frac{dz}{\sqrt{B^2 - AC}},
\]
where $B'^2 - A'C'$ and $B^2 - AC$ are quartic functions of $w$ and $z$ respectively. This is, of course, integrable (viz., the integral is the original equation $\Omega = 0$); and it follows, from the theory of elliptic functions, that the two quartic functions must be linearly transformable into each other; viz., they must have the same absolute invariant $\frac{I^3}{J^2}$. It is, in fact, easy to verify, not only that this is so, but that the two functions have the same quadrinvariant $I$, and the same cubinvariant $J$.

The new point of view is, that we take the coefficients $a$, $b$, \&c., to be homogeneous functions of $(x, y)$, their degrees being such that the equation $\Omega = 0$ is a quartic equation $(\mid\mid\mid x, y, z, w)^4 = 0$: viz., this equation now represents a quartic surface having a node (conical point) at the point $(x = 0,\; y = 0,\; z = 0)$, and also a node at the point $(x = 0,\; y = 0,\; w = 0)$, say, these points are $O$, $O'$ respectively. The equation $B'^2 - A'C' = 0$ gives the circumscribed sextic cone having $O$ for its vertex, and the equation $B^2 - AC = 0$ the circumscribed sextic cone having $O'$ for its vertex; each of these cones has the line $OO'$ ($x = 0,\; y = 0$) for a nodal line, as appears geometrically, and also by the equations containing $z$, $w$ respectively in the degree 4. Considering $B'^2 - A'C'$ as a quartic function of $z$, its quadrinvariant is a function $(x, y)^4$, and its cubinvariant a function $(x, y)^6$; and similarly, considering $B^2 - AC$ as a quartic function of $w$, its invariants are functions $(x, y)^4$ and $(x, y)^6$. We have thus, between the two cones, a geometrical relation answering to the analytical one of the identity of the invariants; but the nature of this geometrical relation is not obvious; and it presents itself as an interesting subject of investigation.

\clearpage

%% ========================================================================
%% PAPER 453: On a Problem in the Calculus of Variations
%% ========================================================================
\section*{\textbf{453.}}
\addcontentsline{toc}{section}{453. On a Problem in the Calculus of Variations}

\begin{center}
\textbf{ON A PROBLEM IN THE CALCULUS OF VARIATIONS.}
\end{center}

[From the \textit{Proceedings of the London Mathematical Society}, vol.~iii (1869--1871), pp.~221, 222. Read February 9, 1871.]

The problem is, $z = \frac{1}{3}(3x - y^2)y$, to find $y$ a function of $x$ such that $\int z\,dx = \max.$ or $\min.$, subject to a given condition $\int y\,dx = c$ (the limits of each integral being $x_0$, $x_1$, where these quantities are each positive, and $x_1 > x_0$). The ordinary method of solution gives $y^2 = x + \lambda$, where $(x_1 + \lambda)^2 - (x_0 + \lambda)^2 = \frac{3}{2}c$; so long as $c$ is not less than $(x_1 - x_0)$, there is a real value of $\lambda$, but for a smaller value of $c$ there is no real value. The difficulty arising in this last case is somewhat illustrated by replacing the original problem by a like problem of ordinary maxima and minima; viz., $x_1$, $x_2$, $\ldots$, $x_n$ being given positive values of $x$, in the order of increasing magnitude; and if, in general, $z_i = \frac{1}{3}(3x_i - y_i^2)y_i$, then the problem is to find $y_i$ a function of $x_i$, such that $\Sigma z_i = \max.$ or $\min.$, subject to the condition $\Sigma y_i = c$. We have here $y_i^2 = x_i + \lambda$ where $\lambda$ is to be determined by the condition $\Sigma y_i = c$; the remainder of the investigation turns on the question of the sign $y_i = +\sqrt{x_i + \lambda}$ or $y_i = -\sqrt{x_i + \lambda}$, to be taken for the several values of $i$ respectively.

\clearpage

%% ========================================================================
%% PAPER 454: A Third Memoir on Quartic Surfaces
%% ========================================================================
\section*{\textbf{454.}}
\addcontentsline{toc}{section}{454. A Third Memoir on Quartic Surfaces}

\begin{center}
\textbf{A THIRD MEMOIR ON QUARTIC SURFACES.}
\end{center}

[From the \textit{Proceedings of the London Mathematical Society}, vol.~iii (1869--1871), pp.~234--266. Read April 13, 1871.]

The present Memoir is a continuation of my former researches on Nodal Quartic Surfaces, [445, 451]. The leading idea is, that for a quartic surface with $k$-nodes, given the nature of the circumscribed, $(k - 1)$ nodal, sextic cone belonging to any one node of the surface (for instance, $k = 10$, that it is a cone $(3, 3)$ composed of two cubic cones), we thereby determine the equation of the quartic surface, and consequently the nature of the remaining $(k - 1)$ nodes thereof. By means of this general theory I complete, in an essential point, the theory of the Symmetroid; viz., I show that a 10-nodal quartic surface having a single node $(3, 3)$ is a Symmetroid; whence, as appears by my second Memoir, [451], each of the remaining nine nodes is also a node $(3, 3)$; and we have the theory of the remarkable system of ten points in space such that, joining any one of them with the remaining nine, the nine lines thus obtained are the intersections of two cubic cones. A large part of the Memoir is devoted to the consideration of the surfaces with 16, 15, 14, and 13 nodes: this is substantially a reproduction of the results obtained by Kummer in the Memoir ``Ueber die algebraischen Strahlensysteme, \&c.,'' already referred to; but the results in question are brought into connexion with the theory of the present Memoir, and they are, by a change of the constants, exhibited in a form of much greater symmetry and elegance. I attach importance also to the square diagrams by means of which I have exhibited, in a compendious form, the relation between the several nodes and circumscribed sextic cones.

The paragraphs are numbered consecutively with those of the first and second Memoirs.

\subsection*{Preliminary Considerations and Classification.}

\textbf{\S 125.}\quad I call to mind that if a quartic surface has a node (conical point), then there is for this node a tangent quadricone and a circumscribed sextic cone; viz., if the surface has $(k - 1)$ other nodes, or in all $k$ nodes, then the sextic cone has $(k - 1)$ nodal lines (passing through the other nodes respectively), and we have thus for the different forms of the sextic the table No.~11; viz., this is

\begin{center}
\begin{tabular}{|c|c|c|}
\hline
\multicolumn{2}{|c|}{Nodes of} & Circumscribed Sextic Cone. \\
\cline{1-2}
Surface. & Nodal Lines of Cones. & \\
\hline
1 & 0 & $6$ \\
2 & 1 & $6_1$ \\
3 & 2 & $6_2$ \\
4 & 3 & $6_3$ \\
5 & 4 & $6_4$ \\
6 & 5 & $6_5$ \\
7 & 6 & $6_6$ \\
8 & 7 & $6_7$ \\
9 & 8 & $6_8$, $5_1$, $4_2$, $2$ \\
10 & 9 & $6_9$, $5_1$, $4_2$, $2$, $4_1$, $1$, $1$, $3_1$, $3$ \\
11 & 10 & $6_{10}$, $5_1$, $4_2$, $2$, $4_1$, $1$, $1$, $3_1$, $3_1$, $3$, $2$, $1$ \\
12 & 11 & $5_1$, $4_2$, $2$, $4_1$, $1$, $1$, $3_1$, $3_1$, $3$, $2$, $1$ \\
13 & 12 & $4_1$, $1$, $1$, $3_1$, $2$, $1$, $3$, $1$, $1$, $1$, $2$, $2$, $2$ \\
14 & 13 & $3_1$, $1$, $1$, $1$, $2$, $2$, $1$, $1$ \\
15 & 14 & $2_1$, $1$, $1$, $1$, $1$ \\
16 & 15 & $1$, $1$, $1$, $1$, $1$, $1$ \\
\hline
\end{tabular}
\end{center}

viz., $6$ denotes a proper sextic cone without nodal lines; $6_1$ a proper sextic cone with one nodal line; $5$, $1$ a proper quintic cone and a plane, \&c.

We may distinguish the nodes according to the sextic cones; thus, a node $6$ means a node for which the circumscribed cone is a proper sextic cone, $(1, 1, 1, 1, 1, 1)$ a node where the circumscribed cone breaks up into six planes, \&c.

\textbf{\S 126.}\quad A 16-nodal surface has 16 nodes $(1, 1, 1, 1, 1, 1)$, and a 15-nodal surface has 15 nodes $(2, 1, 1, 1, 1)$; but, for a 14-nodal surface, the question arises how many nodes are $(3_1, 1, 1, 1)$, and how many $(2, 2, 1, 1)$. It was remarked, No.~13, that the only possible cases were 14, 0; 8, 6; or 2, 12; and that we might, in like manner, limit the number of possible cases for other values of $k$; but that the inquiry was not then further pursued. I resume this inquiry, but without obtaining as yet a complete answer.

\textbf{\S 127.}\quad It is to be observed that a line joining any two nodes is not, in general, a line on the surface, but that it may be so; the surfaces for which this is so (viz., any surface which contains upon it a line through two nodes) form, however, a class by themselves, which at present I altogether exclude from consideration. This being so, it will appear in the sequel that there is but one kind of surface having a node $(2, 2, 1, 1)$, and but one kind of surface having a node $(3_1, 1, 1, 1)$. Now there is a surface, Kummer's 14-nodal, the nodes of which are $8(3_1, 1, 1, 1) + 6(2, 2, 1, 1)$; wherefore the two kinds are identical, and are each of them Kummer's 14-nodal surface. Similarly, for the 13-nodal surfaces, there is but one kind having a node $(4_1, 1, 1)$, but one kind having a node $(3, 1, 1, 1)$, and moreover but one kind having a node $(3_1, 2, 1)$; and we have Kummer's 13-nodal surface with the nodes $3(4_1, 1, 1) + 1(3, 1, 1, 1) + 9(3_1, 2, 1)$; hence the three kinds are identical with each other and with Kummer's. Moreover, there is but one kind having a node $(2, 2, 2)$; hence all the other nodes must be $(2, 2, 2)$, and we have a surface $13(2, 2, 2)$ not given by Kummer. And in like manner for the 12-nodal surfaces, we have the two kinds given by Kummer, and a third kind $12(4_1, 1, 1)$ not given by him; the arrangement thus far being

\begin{center}
\begin{tabular}{c|l}
No.\ of Nodes. & Character of Surface. \\
\hline
16 & $16(1, 1, 1, 1, 1, 1)$, \\
15 & $15(2, 1, 1, 1, 1)$, \\
14 & $8(3_1, 1, 1, 1) + 6(2, 2, 1, 1)$, \\
13 ($\alpha$) & $3(4_1, 1, 1) + 1(3, 1, 1, 1) + 9(3_1, 2, 1)$, \\
\quad ($\beta$) & $13(2, 2, 2)$, \\
12 & ($\alpha$) $12(4_1, 2)$, \\
\quad ($\beta$) & $2(5_1, 1) + 6(3_1, 3_1) + 4(3, 2, 1)$, \\
\quad ($\gamma$) & $12(4_1, 1, 1)$.
\end{tabular}
\end{center}

\textbf{\S 128.}\quad But in the next following case we have Kummer's surface, viz.
\[
11(\alpha)\quad 1(6_{10}) + 10(3_1, 3),
\]
and I do not know whether one, two, or three kinds of surface having nodes $(4_1, 1, 1)$, $(4_1, 2)$, and $(5_1, 1)$. And in the next case we have (as will appear) the Symmetroid, viz.,
\[
10(\alpha)\quad 10(3, 3),
\]
and I do not know how many kinds of surfaces having a node or nodes $6_{10}$, $(5_1, 1)$, $(4_1, 2)$, $(4_1, 1, 1)$.

It will be observed that the present division has nothing to do with the octadic and disinodal division in the former Memoir.

\textbf{\S 129.}\quad I consider a conic $A = 0$, and any six tangents thereof, $t_1 = 0$, $t_2 = 0$, $t_3 = 0$, $t_4 = 0$, $t_5 = 0$, $t_6 = 0$; we have an identical equation which might be written $AC - B^2 = t_1t_2t_3t_4t_5t_6$, but it will be convenient, introducing a constant factor $K$, to write it
\[
AC - B^2 = Kt_1t_2t_3t_4t_5t_6,
\]
$B$ being a cubic function and $C$ a quartic function of the coordinates.

Consider now the series of factors, such as
\begin{align*}
lA &+ mt_1t_2, \\
sA &+ mt_1t_2t_3, \\
UA &+ mt_1t_2t_3t_4, \\
VA &+ mt_1t_2t_3t_4t_5, \\
WA &+ mt_1t_2t_3t_4t_5t_6,
\end{align*}
where $m$ is a constant, $l$ a constant, $s$ a linear function, $U$, $V$, $W$ functions of the degrees 2, 3, 4 respectively; and compose with one or more such factors an expression involving the term $Kt_1t_2t_3t_4t_5t_6$; for instance, such an expression is
\[
(sA + mt_1t_2t_3)(l'A + m't_4t_5t_6);
\]
this is, of the form $An + Kt_1t_2t_3t_4t_5t_6$, viz., $An + (AC - B^2)$, or $A(n + C) - B^2$, say $A\Gamma - B^2$; or what is the same thing, introducing a new coordinate $w$, we have a quadric function
\[
Aw^2 + 2Bw + \Gamma,
\]
$AV - B^2$ --- the discriminant of which, $B^2 - A\Gamma$, is equal to the expression in question.

\textbf{\S 130.}\quad In the sequel $(x, y, z, w)$ are considered as the coordinates of a point in space; $A = 0$ is thus a quadric cone, $t_1 = 0$, $t_2 = 0$, $\ldots$, $t_6 = 0$, any six tangent planes thereof; and hence $Aw^2 + 2Bw + \Gamma = 0$ a quartic surface, having the point $(x = 0,\; y = 0,\; z = 0)$ for a node, whereof the circumscribed cone $A\Gamma - B^2 = 0$ breaks up in the assumed manner.

Thus, in order that the circumscribed cone may be as above
\[
(sA + mt_1t_2t_3)(l'A + m't_4t_5t_6),
\]
we have only to assume
\[
\Gamma = C + mm'(sl'A + sm't_4t_5t_6 + l'mt_1t_2t_3),
\]
and so in other cases. Observe that $sA + mt_1t_2t_3 = 0$ is a cubic cone, which, so long as $m$, $s$, are arbitrary, has no nodal line; but establishing a single relation (say $s$ remains arbitrary, but a proper value is assigned to $m$) it will be a cubic cone having a nodal line. And so $UA + mt_1t_2t_3t_4 = 0$ is a quartic cone without any nodal line, but by particularising the constants it may be made to have one, two, or three nodal lines. Such nodal determinations are obviously required in order that the formula may extend to all the before-mentioned forms of the circumscribed cone. The foregoing analysis is the foundation of the whole theory: I have given it, as above, apart from the theory, in order that the nature of it may be the better perceived; but I have now to bring it into connexion with the theory.

\subsection*{On the Sextic Curves, $AC - B^2 = 0$.}

\textbf{\S 131.}\quad I revert to the consideration of plane curves. The equation of a sextic curve $(\mid\mid\mid x, y, z)^6 = 0$ cannot be in general expressed in the form $AC - B^2 = 0$, where the degrees of $A$, $B$, $C$ are 2, 3, 4 respectively; in fact, the existence of such a form implies that there is a conic $A = 0$ touching the sextic 6 times; and since a conic can only be made to satisfy 5 conditions, there is not in general any such conic.

\textbf{\S 132.}\quad Such conic, when it exists, is said to be inscribed in the sextic, and the sextic to be circumscribed about the conic, or to be an ``amphigram;'' and then, $A = 0$ being the equation of the conic, that of the sextic is expressible in the form in question $AC - B^2 = 0$. It is clear that $B = 0$ is a cubic curve passing through the 6 points of contact of the conic with the sextic, and that any such curve may be taken for the cubic $B$; in fact, if $B = B'$ is any particular solution of the equation $AC - B^2 = 0$, but changing the cubic $B'$ by the addition of a linear function of the coordinates, the equation of the general cubic is $B = B' + pA = 0$; and then writing
\[
A = A,\quad B = B' + pA,\quad C = C + 2B'p + Ap^2,
\]
we have $AC - B^2 = AC - B'^2$; so that the original form $AC - B'^2 = 0$ becomes $AC - B^2 = 0$. But the cubic $B = 0$ being assumed at pleasure, the quartic $C = 0$ is a determinate curve.

\textbf{\S 133.}\quad It is to be observed that a sextic curve may be an amphigram in more than one way: certainly in two, three, or four, and possibly in a greater number of ways. For the equation of the curve contains 27 constants, and hence determining the sextic so as to touch 4 given conics each of them 6 times, there are still 3 constants and the curve will be an amphigram in regard to each of the 4 conics; say it is a quadruple amphigram. But in the sequel we are only concerned with a sextic curve considered as an amphigram in regard to a given conic $A = 0$ (no attention being paid to the other inscribed conics, if any); and then, by what precedes, taking $B = 0$ any cubic whatever through the 6 points of contact, we have a determinate quartic curve $C = 0$, and the equation of the sextic curve assumes the form $AC - B^2 = 0$.

\textbf{\S 134.}\quad The curves $A = 0$, $B = 0$, $C = 0$ contain respectively 5, 9, 14 constants; whence considering the function $B$ as containing an arbitrary constant factor, for the curve $AC - B^2 = 0$, the number of constants is prim\^a facie $5 + 9 + 14 + 1 = 29$; but on account of the arbitrary linear function $p$, the real number is $29 - 3 = 26$: this is right, for a sextic curve contains 27 constants; and the curve being an amphigram, there is one relation between the constants, $27 - 1 = 26$.

\textbf{\S 135.}\quad Suppose now that the sextic curve $AC - B^2 = 0$ breaks up into two or more separate curves, say into the two curves $P = 0$, $Q = 0$ of the orders $f$, $g$ respectively, $f + g = 6$. We have
\[
AC - B^2 = PQ = 0;
\]
and the conic $A = 0$ touching the sextic six times, must, it is clear, touch the curves $P = 0$, $Q = 0$, $f$ and $g$ times respectively. And so when the sextic breaks up into any number of curves, each component curve $P = 0$ of the order $f$ must touch the sextic $f$ times.

\textbf{\S 136.}\quad It follows that if the sextic break up into six lines, say $AC - B^2 = t_1t_2t_3t_4t_5t_6 = 0$, then that each of the lines $t_1 = 0$, $t_2 = 0$, $\ldots$, $t_6 = 0$ is a tangent to the conic. And conversely, starting with the conic $A = 0$ and any six tangents thereof $t_1 = 0$, $t_2 = 0$, $\ldots$, $t_6 = 0$, we have an identity of the form in question. In fact, taking any two of the tangents, say $t_1 = 0$ and $t_2 = 0$, then, if $p = 0$ be the equation of the line joining their points of intersection, the equation of the conic will be of the form $t_1t_2 + p^2 = 0$, that is, we may write $A = t_1t_2 + p^2$, or what is the same thing, $t_1t_2 = A - p^2$. (Considering $A$ as a given quadric function of the coordinates, this of course implies that the implicit constant factors of $t_1$, $t_2$, $p$ are properly determined.) Similarly, $q = 0$ being the line through the points of contact of $t_3$, $t_4$, and $r = 0$ that through the points of contact of $t_5$, $t_6$, we have $t_3t_4 = A - q^2$ and $t_5t_6 = A - r^2$; whence, to satisfy the equation $AC - B^2 = t_1t_2t_3t_4t_5t_6$, we have only to assume $B = lA + pqr$, $l$ an arbitrary linear function of the coordinates, and the equation then gives
\[
C = l^2A^2 - A(p^2 + q^2 + r^2) + (p^2q^2 + r^2p^2 + p^2r^2) + l^2 \cdot \ldots + 2lpqr.
\]

\textbf{\S 137.}\quad It will be observed that the grouping of the six tangents into pairs is arbitrary. By altering this grouping, we merely alter the linear function $l$, but do not obtain any new solution. Thus, say that the new form is $B = l'A + p'q'r'$, then, by properly determining the linear function $l$, we can reduce this to the original form $B = lA + pqr$; viz., we can satisfy identically the equation $(l - l')A + pqr - p'q'r'^2 = 0$; or what is the same thing, $\lambda A + pqr - p'q'r'^2 = 0$, where $\lambda$ is a linear function of the coordinates. We have, in fact, the conic $A = 0$ and the cubic $pqr = 0$ intersecting in the six points of contact; any other cubic through these six points; and consequently the cubic $p'q'r' = 0$ must be expressible in the form $\lambda A + pqr = 0$, and we have thus the identity in question.

\textbf{\S 138.}\quad We have just seen that the value of $B$ is necessarily of the form $B = lA + pqr$, but we are not concerned with its expression in this particular form. What we require in the sequel is a value of $B$, and thence one of $C$, satisfying the identity in question, $AC - B^2 = t_1t_2t_3t_4t_5t_6$; or what is the same thing, introducing for convenience a constant factor $K$, the identity
\[
AC - B^2 = Kt_1t_2t_3t_4t_5t_6.
\]

\textbf{\S 139.}\quad Instead of $C$, I write $\Gamma$, and consider the sextic amphigram $A\Gamma - B^2 = 0$, touched by the conic $A = 0$ in the points of contact of the conic with the six tangents $t_1 = 0$, $t_2 = 0$, $\ldots$, $t_6 = 0$. Suppose the sextic curve breaks up into factors; if one of these factors is a line, it is one of the six tangents, say the tangent $t_6 = 0$. If there is a conic factor, this is a conic touching the conic at its points of contact with two of the tangents, say the equation is $lA + mt_1t_2 = 0$. Similarly, if there is a cubic, quartic, or quintic factor, then the equation hereof is $sA + mt_1t_2t_3 = 0$, $UA + mt_1t_2t_3t_4 = 0$, or $VA + mt_1t_2t_3t_4t_5 = 0$. Or going on to the next case of a sextic factor (being of course the whole curve), we may say that this is $WA + mt_1t_2t_3t_4t_5t_6 = 0$. (Observe that since $AC - B^2 = Kt_1t_2t_3t_4t_5t_6$, this means only that the equation of the sextic amphigram is of the assumed form $A\Gamma - B^2 = 0$.)

\textbf{\S 140.}\quad By what precedes we can, for a sextic amphigram which breaks up in any assigned manner, determine the value of $\Gamma$. For instance, let the amphigram break up into two cubic curves; say these are $sA + mt_1t_2t_3 = 0$, $s'A + m't_4t_5t_6 = 0$. Assume
\[
A\Gamma - B^2 = mm'(sA + mt_1t_2t_3)(s'A + m't_4t_5t_6),
\]
then this equation is
\[
A\Gamma - B^2 = mm'\{ss'A^2 + (sm't_4t_5t_6 + s'mt_1t_2t_3)A\} + AC - B^2;
\]
that is, we have
\[
\Gamma = C + mm'(ss'A + sm't_4t_5t_6 + s'mt_1t_2t_3),
\]
and so in any other case.

I have already adverted to the question of the ``nodal determination'' of the formulae, and it might be properly here considered; viz., the question is as to the determination of the constants in such manner that, for instance, $sA + mt_1t_2t_3 = 0$ may be a nodal cubic, $UA + mt_1t_2t_3t_4 = 0$ a nodal, binodal, or trinodal quartic, \&c.; but I defer it for the moment in order first to apply the theory to the quartic surfaces.

\subsection*{Application to Quartic Surfaces.}

\textbf{\S 141.}\quad If a quartic surface has a node or nodes, we may take for a node the point $x = 0$, $y = 0$, $z = 0$; the equation of the quartic surface is then of the form
\[
Aw^2 + 2Bw + \Gamma = 0,
\]
where $A$, $B$, $\Gamma$ are functions of $x$, $y$, $z$ of the degrees 2, 3, 4 respectively. $A = 0$ is the tangent quadricone at the node in question; and the circumscribed cone is $A\Gamma - B^2 = 0$. By what precedes, this is an amphigram touching the quadricone along six generating lines thereof; say the tangent planes of the cone $A = 0$ along these six lines respectively are $t_1 = 0$, $t_2 = 0$, $\ldots$, $t_6 = 0$. We have then an identical equation
\[
AC - B^2 = Kt_1t_2t_3t_4t_5t_6,
\]
viz., regarding for a moment this equation as an equation for the determination of $B$, and $B'$ as any particular solution thereof, then its general solution is $B' + tA$, where $t$ is an arbitrary linear function of $(x, y, z)$, and the $B$ in the equation of the surface is properly $= B' + tA$. But by the substituting $w - t$ in place of $w$, the $B$ of the equation of the surface would then be made $= B'$; and it thus appears that we may, without loss of generality, take the $B$ of this equation to be any particular value satisfying the identity in question; and then, $B$ having such particular value, $C$ is a quartic function of $(x, y, z)$ completely determined by the same identity. And we then, by what precedes, at once determine $\Gamma$ so that the circumscribed cone $A\Gamma - B^2 = 0$ may be a cone breaking up in any assigned manner; for instance, if it be a cone $(3, 3)$, then, as just mentioned, the two cubic cones are $sA + mt_1t_2t_3 = 0$, $s'A + m't_4t_5t_6 = 0$; and $\Gamma$ has the value
\[
\Gamma = C + mm'(ss'A + sm't_4t_5t_6 + s'mt_1t_2t_3),
\]
above obtained.

\subsection*{On the Nodal Determination.}

\textbf{\S 142.}\quad I am not able to discuss with much completeness the question of nodal determination. We have to consider a cubic curve $sA + mt_1t_2t_3 = 0$, a quartic curve $UA + mt_1t_2t_3t_4 = 0$, \&c., as the case may be, and to determine the constants so that this shall have a node or nodes. Consider for a moment the form $PA + t_iQ = 0$, where $Q$ denotes the product $mt_1t_2\ldots$ of all or any of the tangents $t_1$, $\ldots$, $t_6$; the orders of $PA$, $t_iQ$ are of course equal, that is, the order of $P$ is less by unity than that of $Q$. I say that, by establishing a single relation between the constants, this may be made to have a node at the point of contact $A = 0$, $t_i = 0$. In fact, writing $A = \lambda\partial_\alpha + \mu\partial_\beta + \nu\partial_\gamma$, where $\lambda$, $\mu$, $\nu$ are arbitrary, there will be a node at any point if for that point $\frac{1}{\Delta}(PA + t_iQ) = 0$. But for the point $A = 0$, $t_i = 0$ this becomes $PA\Delta + Q\Delta t_i = 0$; moreover, if $t = 0$ be any other tangent of the conic $A = 0$, and if $p = 0$ be the line joining the points of contact of the tangents $t$, $t_i$, then we may write $A = tt_i - p^2$, and thence (since at the point in question, $A = 0$, $t_i = 0$, we have also $p = 0$) we find $\Delta tt_i = t\Delta t_i$, and the foregoing equation thus becomes $(tP + Q)t\Delta t_i = 0$; viz., this equation is satisfied irrespectively of the values of $\lambda$, $\nu$, if only at the point in question (that is, for the values of the coordinates which belong to the point $t_i = 0$, $A = 0$) we have $tP + Q = 0$, which is a single relation between the constants.

\textbf{\S 143.}\quad In particular the cubic curve $sA + mt_1t_2t_3 = 0$ may be made to have a node at the point of contact of any one of the three tangents; the quartic curve $UA + mt_1t_2t_3t_4 = 0$, a node, or two or three nodes, at the point or points of contact of any one, two, or three of the four tangents; and so in other cases. These are not the only solutions, and they are in fact solutions which (as afterwards explained) I propose to reject, attending in each case only to the remaining or proper solutions of the problem.

\textbf{\S 144.}\quad To obtain in a different manner the foregoing result, consider again the cubic curve $sA + mt_1t_2t_3 = 0$; regarding this as a given curve, the conic $A = 0$ is a conic determined (not of course completely) as a conic having therewith 3 points of 2-pointic intersection; viz., if the cubic has a node, then the cone $A = 0$ is either a conic passing through the node and besides touching the curve twice, or else it is a conic touching the curve 3 times; the former is of course the above mentioned case where there is a node at one of the points of contact on the conic $A = 0$; the latter is regarded as the proper solution. So in the case of a quartic curve $UA + mt_1t_2t_3t_4 = 0$, regarding this as a given curve, the conic $A = 0$ is a conic having therewith four points of 2-pointic intersection; viz., if the quartic curve has one, two, or three nodes, then the conic is either a conic passing through one, two, or three nodes, and besides touching the quartic thrice, twice, or once; or else it is a conic touching the quartic four times. The former is the above mentioned case where there is a node or nodes at a point or points of contact with the conic $A = 0$; the latter is regarded as the proper solution.

\textbf{\S 145.}\quad To fix the ideas, and at the same time obtain a result which will be afterwards useful, I work out the formulae for the cubic curve $sA + mt_1t_2t_3 = 0$, taking this equation under the form
\[
\left(\frac{\lambda}{x} + \frac{\mu}{y} + \frac{\nu}{z}\right)(x^2 + y^2 + z^2 - 2yz - 2zx - 2xy) + mxyz = 0.
\]
This may have a node in two different ways; viz.,

$1^{\circ}$. At the point of contact of one of the tangents $x = 0$, $y = 0$, $z = 0$ with the conic $A = 0$; say at the point of contact of $x = 0$, that is, the point $x = 0$, $y - z = 0$. The value of $m$ is $p = -\frac{4}{\mu}$; hence $\frac{1}{p} = -\frac{1}{4}\mu$; and, substituting, the equation of the curve becomes
\[
\left(\frac{\lambda}{x} + \frac{\mu}{y} + \frac{\nu}{z}\right)\{x^2 - 2x(y + z) + (y - z)^2\} + \frac{1}{4}\mu(x + y - z)^2 = 0,
\]
which has obviously a node at the point in question.

$2^{\circ}$. The node may be at a point not on the conic $A = 0$, viz. the value of $m$ is $= -\frac{4(\lambda + \mu + \nu)}{\lambda\mu\nu}$, the equation is
\[
\left(\frac{\lambda}{x} + \frac{\mu}{y} + \frac{\nu}{z}\right)(x^2 + y^2 + z^2 - 2yz - 2zx - 2xy) - \frac{4(\lambda + \mu + \nu)}{\lambda\mu\nu}xyz = 0.
\]
In fact, writing for shortness
\begin{align*}
\lambda + \mu + \nu &= L, & \lambda - \mu + \nu &= M, & \lambda + \mu - \nu &= N, \\
\lambda + \mu + \nu &= P, & -\lambda + \mu + \nu &= Q, & \&\text{c.},
\end{align*}
the node is at the point $x : y : z = L\lambda : M\mu : N\nu$; which is at once verified, if we remark that, writing for convenience $x, y, z = L\lambda, M\mu, N\nu$, then we have
\begin{align*}
-x + y + z &= MN, & x - y + z &= NL, & x + y - z &= LM, \\
\frac{\lambda}{x} + \frac{\mu}{y} + \frac{\nu}{z} &= P, \\
x^2 + y^2 + z^2 - 2yz - 2zx - 2xy &= -LMNP\ (= A).
\end{align*}
For, of the three equations for the coordinates of a node, the first is
\[
\frac{2\lambda}{x}A + P(-2y - 2z + 2x) - \frac{4(\lambda + \mu + \nu)}{\mu\nu}x = 0,
\]
that is, for the values in question,
\[
-\frac{2}{L}LMNP + P(-2MN) + \frac{4}{\mu\nu}LMN\lambda\nu = 0,
\]
that is
\[
-2LMNP - 2LMNP + 4LMNP = 0,
\]
which is satisfied; and similarly the other two equations are satisfied.

\subsection*{Quartic Surfaces resumed.}

\textbf{\S 146.}\quad Passing now from curves to cones, and to the theory of the quartic surface, suppose that there is a component cone having a nodal line, say the cubic cone $sA + mt_1t_2t_3 = 0$: if the remaining factor is $t_4t_5t_6$, then we have
\[
A\Gamma - B^2 = mm'(sA + mt_1t_2t_3)t_4t_5t_6.
\]
Suppose the nodal line is a line of contact with the cone $A = 0$, say its equations are $t_i = 0$, $p = 0$ ($p$ a linear function), then $sA + mt_1t_2t_3$ is a quadric function $(\mid\mid t_i, p)^2$, (of course with variable coefficients); hence $A\Gamma - B^2$ is a quadric function; and $A$ being a linear function $(\mid\mid t_i, p)$, it follows that $B$ is a linear function, and thence that $\Gamma$ is also a linear function; that is, $A$, $B$, $\Gamma$ are each of them a linear function $(\mid\mid t_i, p)$, or the line in question (viz.\ the line of contact $A = 0$, $t_i = 0$) is a line on the quartic surface $Aw^2 + 2Bw + \Gamma = 0$. As already mentioned, I exclude from consideration the surfaces which have upon them a line through two nodes; that is, I exclude from consideration the case in question where any component cone, or say where the sextic cone, has a nodal line which is a line on the tangent cone $A = 0$.

\textbf{\S 147.}\quad Now, excluding the case just referred to, I assume as a postulate that there is but one way in which the cubic cone $sA + mt_1t_2t_3 = 0$ can be made to have a nodal line, or the quartic cone $UA + mt_1t_2t_3t_4 = 0$ one, two, or three nodal lines, \&c., as the case may be. It is to be understood that this does not mean that the constants are in any of these cases completely determined, but that there is between them a relation or relations constituting a general solution which includes in itself every particular solution whatever. I have no doubt that as regards the cubic cone at least the assumption is correct. This being so, the character of a single node determines the nature of the surface; for instance, if there is a node $(3_1, 3)$, then taking this as the point $(x = 0,\; y = 0,\; z = 0)$ the equation of the surface is $Aw^2 + 2Bw + \Gamma = 0$, where
\[
\Gamma = C + mm'(ss'A + l'smt_4t_5t_6 + ls'mt_1t_2t_3),
\]
a surface of a determinate nature; so that the character of all the remaining nodes is completely determined.

\textbf{\S 148.}\quad The point to be attended to is, that if for instance there were two essentially distinct ways of giving the cubic cone $sA + mt_1t_2t_3 = 0$ a nodal line (such as there would be if the excluded case were considered admissible), then the foregoing equation of the surface would or might include two distinct forms of equation applying to different kinds of surface. The conclusion is that there is but one kind of quartic surface having a node $(3_1, 3)$. Admitting this, and similarly that there is but one kind of quartic surface having a node $6_{10}$, it follows that if (as the fact is) there is a surface having the nodes $1(6_{10}) + 10(3_1, 3)$ (Kummer's 11-nodal surface), then that the two first-mentioned kinds are in fact each of them this last-mentioned kind of surface; and it was in this manner that I arrived at the enumeration given near the beginning of the present Memoir.

\textbf{\S 149.}\quad The reasoning is, of course, in place of a direct demonstration which would consist in showing that a surface having a node $(3_1, 3)$ has 9 other like nodes, and also a node $6_{10}$; and that a surface having a node $6_{10}$ has 10 other nodes $(3_1, 3)$; and that, starting from either form of equation, we could, by passing to a node of the other kind, obtain the other form of equation.

\subsection*{Enumeration of the Cases.}

\textbf{\S 150.}\quad I collect the results as follows: I call to mind that we have always the identical equation $AC - B^2 = Kt_1t_2t_3t_4t_5t_6$, that the equation of the surface is $Aw^2 + 2Bw + \Gamma = 0$, and that the circumscribed cone is $A\Gamma - B^2 = 0$. The equation of a surface having different kinds of nodes will assume different forms according as the origin (or point $x = 0$, $y = 0$, $z = 0$) is taken to be at a node of one or other of these kinds; these forms of the equations are distinguished as ``node-forms,'' viz., we speak of the node-form $(3_1, 3)$ when the origin is a node $(3_1, 3)$, and so in other cases.

\vspace{0.5em}
\textbf{The 16-nodal surface}
\[16(1, 1, 1, 1, 1, 1),\]
\textbf{node-form} $(1, 1, 1, 1, 1, 1)$, cone is $6$,
and viz., equation is $Aw^2 + 2Bw + C = 0$.

\vspace{0.5em}
\textbf{The 15-nodal surface}
\[15(2, 1, 1, 1, 1),\]
\textbf{node-form} $(2, 1, 1, 1, 1)$, cone is $(lA + mt_1t_2t_3t_4t_5)t_6 = 0$,
and $\Gamma = C + \frac{K}{l}t_1t_2t_3t_4t_5t_6 = 0$.

\vspace{0.5em}
\textbf{The 14-nodal surface}
\[8(3_1, 1, 1, 1) + 6(2, 2, 1, 1),\]
\textbf{node-form} $(3_1, 1, 1, 1)$, cone is $(sA + mt_1t_2t_3)t_4t_5t_6 = 0$,
where $sA + mt_1t_2t_3 = 0$ is a nodal cubic $3_1$,
and $\Gamma = C + \frac{K}{m}t_4t_5t_6$;

\textbf{node-form} $(2, 2, 1, 1)$,
cone is $(lA + mt_1t_2)(l'A + m't_3t_4)t_5t_6 = 0$,
and $\Gamma = C + mm'(ll'A + lm't_3t_4 + l'mt_1t_2)t_5t_6$.

\vspace{0.5em}
\textbf{The 13($\alpha$)-nodal surface}
\[3(4_1, 1, 1) + 1(3, 1, 1, 1) + 9(3_1, 2, 1),\]
\textbf{node-form} $(4_1, 1, 1)$,
cone is $(UA + mt_1t_2t_3t_4)t_5t_6 = 0$,
where $UA + mt_1t_2t_3t_4 = 0$ is a trinodal quartic $4_1$,
and $\Gamma = C + \frac{K}{m}t_5t_6$;

\textbf{node-form} $(3, 1, 1, 1)$,
cone is $(sA + mt_1t_2t_3)t_4t_5t_6 = 0$,
and $\Gamma = C + \frac{K}{m}t_4t_5t_6$;

\textbf{node-form} $(3_1, 2, 1)$,
cone is $(sA + mt_1t_2t_3)(l'A + m't_4t_5)t_6 = 0$,
where $sA + mt_1t_2t_3 = 0$ is a nodal cubic $3_1$,
and $\Gamma = C + mm'(sl'A + m'st_4t_5t_6 + l'mt_1t_2t_3)t_6$.

\vspace{0.5em}
\textbf{The 13($\beta$)-nodal surface}
\[13(2, 2, 2),\]
\textbf{node-form} $(2, 2, 2)$,
cone is $(lA + mt_1t_2)(l'A + m't_3t_4)(l''A + m''t_5t_6) = 0$,
and
\begin{multline*}
\Gamma = C + mm'm''\{ll'l''A^2 + (ll'm''t_5t_6 + l''lm't_3t_4 + l'm''t_1t_2)A \\
+ lm'm''t_3t_4t_5t_6 + l'lm''t_1t_2t_5t_6 + l''m''t_1t_2t_3t_4\}.
\end{multline*}

\vspace{0.5em}
\textbf{The 12($\alpha$)-nodal surface}
\[12(4_1, 2),\]
\textbf{node-form} $(4_1, 2)$,
cone is $(UA + mt_1t_2t_3t_4)(l'A + m't_5t_6) = 0$,
where $UA + mt_1t_2t_3t_4 = 0$ is a trinodal quartic $4_1$,
and $\Gamma = C + mm'(l'UA + m'mt_3t_4t_5t_6 + l'mt_1t_2t_3t_4)$.

\vspace{0.5em}
\textbf{The 12($\beta$)-nodal surface}
\[2(5_1, 1) + 6(3_1, 3_1) + 4(3, 2, 1),\]
\textbf{node-form} $(5_1, 1)$,
cone is $(VA + mt_1t_2t_3t_4t_5)t_6 = 0$,
where $VA + mt_1t_2t_3t_4t_5 = 0$ is a 5-nodal quintic $5_1$,
and $\Gamma = C + \frac{K}{m}Vt_6$;

\textbf{node-form} $(3_1, 3_1)$,
cone is $(sA + mt_1t_2t_3)(s'A + m't_4t_5t_6) = 0$,
where $sA + mt_1t_2t_3 = 0$ and $s'A + m't_4t_5t_6 = 0$ are each of them a nodal cubic $3_1$,
and $\Gamma = C + mm'(ss'A + m'st_4t_5t_6 + ms't_1t_2t_3)$;

\textbf{node-form} $(3, 2, 1)$,
cone is $(sA + mt_1t_2t_3)(l'A + m't_4t_5)t_6 = 0$,
and $\Gamma = C + mm'(sl'A + m'st_4t_5t_6 + l'mt_1t_2t_3)t_6$.

\vspace{0.5em}
\textbf{The 12($\gamma$)-nodal surface,}
\[12(4_1, 1, 1),\]
\textbf{node-form} $(4_1, 1, 1)$,
cone is $(UA + mt_1t_2t_3t_4)t_5t_6 = 0$,
where $UA + mt_1t_2t_3t_4 = 0$ is a binodal quartic $4_1$,
and $\Gamma = C + \frac{K}{m}t_5t_6$.

\vspace{0.5em}
\textbf{The 11($\alpha$)-nodal surface,}
\[1(6_{10}) + 10(3_1, 3)\]
\textbf{node-form} $(6_{10})$,
cone is $WA + mt_1t_2t_3t_4t_5t_6 = 0$,
where this is a 10-nodal sextic $6_{10}$,
and $\Gamma = C + \frac{K}{m}W$;

\textbf{node-form} $(3_1, 3)$,
cone is $(sA + mt_1t_2t_3)(s'A + m't_4t_5t_6) = 0$,
where $sA + mt_1t_2t_3 = 0$ is a nodal cubic $3_1$,
and $\Gamma = C + mm'(ss'A + m'st_4t_5t_6 + ms't_1t_2t_3)$.

\vspace{0.5em}
\textbf{Other 11-nodal surfaces,}

\textbf{node-form} $(5_1, 1)$,
cone is $(VA + mt_1t_2t_3t_4t_5)t_6 = 0$,
where $VA + mt_1t_2t_3t_4t_5 = 0$ is a 5-nodal quintic $5_1$,
and $\Gamma = C + \frac{K}{m}Vt_6$;

\textbf{node-form} $(4_1, 2)$,
cone is $(UA + mt_1t_2t_3t_4)(l'A + m't_5t_6) = 0$,
where $UA + mt_1t_2t_3t_4 = 0$ is a binodal quartic $4_1$,
and $\Gamma = C + mm'(l'UA + m'Ut_3t_4 + l'mt_1t_2t_3t_4)$;

\textbf{node-form} $(4_1, 1, 1)$,
cone is $(UA + mt_1t_2t_3t_4)t_5t_6 = 0$,
where $UA + mt_1t_2t_3t_4$ is a nodal quartic $4_1$,
and $\Gamma = C + \frac{K}{m}t_5t_6$;

but whether these node-forms belong to the same or to different surfaces is not ascertained.

The enumeration is not extended to the 10-nodal surfaces, but I consider one case of these surfaces.

\vspace{0.5em}
\textbf{The 10($\alpha$)-nodal surface} $10(3, 3)$.

\textbf{\S 151.}\quad I assume only that there is a single node $(3, 3)$: taking the cone to be
\[
(sA + mt_1t_2t_3)(s'A + m't_4t_5t_6) = 0;
\]
then for the equation of the surface, in the node-form $(3, 3)$ in question, we have
\[
\Gamma = C + \frac{K}{mm'}(ss'A + sm't_4t_5t_6 + s'mt_1t_2t_3).
\]
But I present this result under a different form, as follows: I write
\[
A = p^2 - ft_1t_2 = q^2 - gt_3t_4 = r^2 - ht_5t_6,
\]
where $f$, $g$, $h$ are constants, and, as before, $p = 0$, $q = 0$, $r = 0$ are the lines joining the points of contact of $t_1$, $t_2$; $t_3$, $t_4$; and $t_5$, $t_6$ respectively: we have
\[
sA + mt_1t_2t_3 = sA + mt_3\left(A - \frac{p^2}{f}\right), \quad\text{and}\quad s'A + m't_4t_5t_6 = s'A + m't_6\left(A - \frac{r^2}{h}\right),
\]
or in place of $s$, $s'$ introducing new linear functions $\sigma$, $\sigma'$, the cubic curves may be taken to be $\sigma A - \frac{m}{f}p^2t_3$, $\sigma'A - \frac{m'}{h}r^2t_6$, so that we have
\begin{multline*}
A\Gamma - B^2 = mm'\left(\sigma A - \frac{m}{f}p^2t_3\right)\left(\sigma'A - \frac{m'}{h}r^2t_6\right) \\
= mm'\left\{\sigma\sigma'A^2 - \left(\frac{\sigma'm}{f}p^2t_3 + \frac{\sigma m'}{h}r^2t_6\right)A + \frac{mm'}{fh}p^2r^2t_3t_6\right\};
\end{multline*}
whence $B = \sqrt{mm'}\left(pqr + tA\right)$, where $t$ is a linear function of the coordinates; and we then have
\[
\Gamma = C + mm'\left(\sigma\sigma'A - \frac{\sigma'm}{f}p^2t_3 - \frac{\sigma m'}{h}r^2t_6\right) + \ldots
\]

\vspace{2em}
\begin{center}
{\small [To be continued.]}
\end{center}
\clearpage


% =========================================================
% Chunk: cayley_vol07_pages_301_350.tex  (orig pages 301--350)
% =========================================================
\thispagestyle{empty}

\vspace{1em}
\noindent\rule{\textwidth}{0.4pt}
\vspace{0.5em}

\noindent\textit{This transcription covers pages 301--350 of Volume VII, comprising Articles 455--462.}

\vspace{0.5em}
\noindent\rule{\textwidth}{0.4pt}

\newpage
\setcounter{page}{301}

%===============================================================================
\article{455}{ON PL"{U}CKER'S MODELS OF CERTAIN QUARTIC SURFACES.}
%===============================================================================

\noindent[From the \textit{Philosophical Transactions}, vol.~CLX. (for the year 1870), pp.~87--103. Received and Read January 13, 1870.]

\vspace{0.5em}

[The article begins on page 300 of this Volume.]

\vspace{0.5em}

There are four real lines $(y=0,\ x=\gamma)$, $(y=0,\ x=\delta)$, $(z=0,\ x=\gamma')$, $(z=0,\ x=\delta')$.

The surface consists of a central pillow-like portion, joined on by two conical points to an upper portion, and by two conical points to an under portion, the whole being included between the planes $x=\gamma'$, $x=\delta'$.

\textbf{Mod.~1} is 13: the form of the equation is
\[
\frac{y^{2}}{l^{2}(x-\gamma)(x-\delta)} + \frac{z^{2}}{l'^{2}(x-\gamma')(x-\delta')} + 1 = 0;
\]
the values $\gamma'$, $\delta'$ lying between $\gamma$, $\delta$; the principal conics are one of them a hyperbola, the other an ellipse, the vertices (on the axis or line $y=0$, $z=0$) of the hyperbola lying between those of the ellipse.

The variable conic, for values of $x$ between $\gamma'$, $\delta'$, has two real vertices, or it is a hyperbola; for the values, say, from $\gamma'$ to $\gamma$, and $\delta'$ to $\delta$, there are four real vertices, or the conic is an ellipse; for values beyond the limits $\gamma$, $\delta$, there are two real vertices, and the conic is a hyperbola. There are the four real lines $(y=0,\ x=\gamma)$, $(y=0,\ x=\delta)$, and $(z=0,\ x=\gamma')$, $(z=0,\ x=\delta')$. The surface consists of 8 portions joined to each other by 8 conical points, but the form can scarcely be explained by a description.

\textbf{Mod.~8} is 32: the form of the equation is
\[
\frac{y^{2}}{l^{2}(x-\gamma)(x-\delta)} + \frac{z^{2}}{l'^{2}(x-\gamma)(x-\delta)} = 1,
\]
viz., the principal conics are each of them a line-pair, the variable conic is always an ellipse.

There are the two real nodal lines $(y=0,\ x=\gamma)$ and $(z=0,\ x=\gamma')$, each of these being in the neighbourhood of the axis and beyond certain limits acnodal crunodal, the surface is a scroll, being, in fact, the well-known surface which is the boundary of a small circular pencil of rays obliquely reflected, and consequently passing through two focal lines.

\textbf{Mod.~4} is 34: the equation is
\[
\frac{y^{2}}{l^{2}(x-\gamma)(x-\delta)^{2}} + \frac{z^{2}}{l'^{2}(x-\gamma')(x-\delta)} + 1 = 0,
\]
where $x=\delta$ not intermediate between the values $x=\gamma$ and $x=\gamma'$, say the order is $\delta$, $\gamma$, $\gamma'$. The surface is thus a cubic surface; the principal conics are ellipses having on the axis a common vertex at the point $x=\delta$, and the remaining two vertices on the same side of the last-mentioned one. The variable conic for values between $\delta$ and $\gamma$ has four real vertices, or it is an ellipse; for values between $\gamma$ and $\gamma'$ two real vertices, or it is a hyperbola; and for values beyond the limits $\delta$, $\gamma'$ it is imaginary.

There are on the surface the two real lines $(y=0,\ x=\gamma)$ and $(z=0,\ x=\gamma')$. The surface consists of a finite portion joined on by two conical points to the remaining portion.

\newpage
\setcounter{page}{302}

\textbf{Mod.~3} is 40: the form of equation is
\[
\frac{y^{2}}{l^{2}(x-\gamma)(x-\delta)} + \frac{z^{2}}{l'^{2}(x-\delta)(x-\delta')} + 1 = 0.
\]
The surface is thus a cubic surface; the principal conics are, one of them an ellipse, the other a pair of imaginary lines intersecting on the ellipse; for values of $x$ between $\gamma$ and $\delta$, the variable conic has thus two real vertices, and it is a hyperbola; for values beyond these limits it is imaginary, and the whole surface is thus included between the planes $x=\gamma$ and $x=\delta$. There are the two real lines $(y=0,\ x=\gamma)$ and $(z=0,\ x=\delta)$.

Taking $l^{2}=l'^{2}=1$, the surface is
\[
\frac{y^{2}}{(x-\gamma)(x-\delta)(x-\delta')} + \frac{z^{2}}{(x-\delta)(x-\delta')} = 0,
\]
which is a particular case of the parabolic cyclide.

As already mentioned, I have not completely identified the remaining models 9 to 14, but I will say a few words about them.

The equatorial surfaces, not included in the preceding 78 cases, Pl"{u}cker distinguishes (vol.~II. p.~363) as ``gedrehte'' or ``tordirte,'' say as twisted equatorial surfaces; the equation of such a surface is
\begin{align*}
by^{2} + 2hyz + az^{2} + ab - h^{2} &= 0,\\
\text{where}\qquad b &= F x^{2} - 2R x + B,\\
a &= Ex^{2} + 2Ux + G,\\
h &= Kx^{2} - Ox - G \quad \text{(or in particular } = -Ox - G\text{)}.
\end{align*}

Mod.~13 is such a surface, being a twisted form of model 2.

Mod.~9 and Mod.~14 belong, I think, to the case $a=0$; viz., the form of the equation is $by^{2} + 2hyz - h^{2} = 0$. The variable conic is a hyperbola, the direction of one of the asymptotes being constant (vol.~II. p.~368).

There are moreover (p.~372) equatorial surfaces in which the variable conic is always a parabola, and where there are on the surface four real or imaginary lines.

Mods.~10, 11, and 12 seem to represent such surfaces.

\newpage
\setcounter{page}{303}

%===============================================================================
\article{456}{NOTE ON THE DISCRIMINANT OF A BINARY QUANTIC.}
%===============================================================================

\noindent[From the \textit{Quarterly Journal of Pure and Applied Mathematics}, vol.~x. (1870), p.~23.]

\vspace{0.5em}

It is well known that the discriminant of a binary quantic $(a,b,c,d,\ldots\hspace{-0.2em}\pitchfork x,y)^{n}$ is of the form
\[
Ma + Nb^{2},
\]
but it is further to be remarked that if $b=0$, then the form is
\[
a(Ma + Nc^{2}),
\]
if $b=0$, $c=0$, the form is
\[
a^{2}(Ma + Nd^{2}),
\]
and so on, until only the lowest two coefficients are not put $=0$. Or, what is the same thing, if in the discriminant of the original function we put $a=0$, then the discriminant divides by $b^{2}$; if $b=0$, the discriminant divides by $a$, and, omitting this factor, if we then write $a=0$, it divides by $c^{2}$; if $b=0$, $c=0$, the discriminant divides by $a^{2}$, and omitting this factor, if we then write $a=0$, it divides by $d^{2}$ and so on, until as before.

Thus if $b=0$, the discriminant of $(a,0,c,d,e\hspace{-0.2em}\pitchfork x,y)^{4}$, divides by $a$, and omitting this factor it is
\begin{align*}
&ac^{3}\\
&- 18ac^{2}d^{2}\\
&+ 54acd^{2}e\\
&- 27ad^{4}\\
&+ 81c^{4}e\\
&- 54c^{3}d^{2},
\end{align*}
which for $a=0$ has the factor $c^{2}$; if $b=0$, $c=0$, the discriminant of $(a,0,0,d,e\hspace{-0.2em}\pitchfork x,y)^{4}$ has the factor $a^{2}$, and omitting this factor it is
\begin{align*}
&ae^{3}\\
&- 27d^{4},
\end{align*}
which for $a=0$ has the factor $d^{2}$; the series of theorems here terminates, since the lowest two coefficients $d$, $e$ are not to be put $=0$.

\newpage
\setcounter{page}{304}

%===============================================================================
\article{457}{ON THE QUARTIC SURFACES $(\asymp\!\asymp\!\asymp U,\ V,\ W)^{2} = 0$.}
%===============================================================================

\noindent[From the \textit{Quarterly Journal of Pure and Applied Mathematics}, vol.~x. (1870), pp.~24--34.]

\vspace{0.5em}

I \textsc{propose} to myself for investigation the quartic surfaces represented by an equation
\[
(\asymp\!\asymp\!\asymp U,\ V,\ W)^{2} = 0,
\]
where $U$, $V$, $W$ are quadric functions of the coordinates.

Such a surface has 8 nodes (conical points), viz., these are the points of intersection of the quadric surfaces $U=0$, $V=0$, $W=0$. It is to be observed, that not every quartic surface with 8 nodes is included under the above form; in fact the equation of a quartic surface contains (homogeneously) 35 coefficients, or say 34 arbitrary parameters; in order that a given point may be a node, 4 conditions must be satisfied, and it is consequently possible to find a quartic surface having 8 given points as nodes (and having in its equation $(34 - 8 \cdot 4 =)\ 2$ arbitrary parameters): but 8 given points are not in general the intersections of three quadric surfaces, and such a quartic surface is therefore not in general included under the above form. I think, however, that it may be assumed that the above form includes all the quartic surfaces having 8 nodes, points of intersection of three quadric surfaces. It will presently appear that, included in the form, we have surfaces where (instead of the 8 nodes) there is a nodal or cuspidal conic; and that these are the most general forms of such quartic surfaces.

A quartic surface has at most 16 nodes, and the general form with 8 nodes must admit of being particularised so that the surface shall acquire any number not exceeding 8 of additional nodes. This does not show, but it is probable, that the above special form with 8 nodes can be particularised so that the surface shall in like manner acquire any number not exceeding 8 of additional nodes. Similarly, a quartic surface with a nodal conic may have besides 1, 2, 3, or 4 nodes and it will be shown in the sequel how the form, particularised so as to give a nodal conic, may be further particularised so as to give the 1, 2, 3, or 4 nodes. So a quartic surface with a cuspidal conic may besides have 1 node, and it will be shown how the form, particularised so as to give a cuspidal conic, may be further particularised so as to give 1 node.

Starting from the equation $(\asymp\!\asymp\!\asymp U,\ V,\ W)^{2} = 0$, we may, by substituting for $U$, $V$, $W$ lineal functions of these expressions, transform the equation precisely in the manner of a conic, and therefore into any of the forms under which the equation of a conic can be exhibited; for instance, in the forms $aU^{2} + bV^{2} + cW^{2} = 0$, $fVW + gWU + hUV = 0$, $UW - V^{2} = 0$, \&c. I attend at present only to the last-mentioned form $UW - V^{2} = 0$, which, it thus appears, is equally general with the original form $(\asymp\!\asymp\!\asymp U,\ V,\ W)^{2} = 0$.

The quartic surface
\[
UW - V^{2} = 0,
\]
where $U$, $V$, $W$ are any quadric functions of the coordinates, may be considered as the envelope of the quadric surface
\[
(\asymp\!\asymp\!\asymp U,\ V,\ W \hspace{-0.2em}\pitchfork \theta,\ 1)^{2} = 0,
\]
where $\theta$ is an arbitrary parameter. And it thus appears that it is very easy to reciprocate (in regard to any given quadric surface) the quartic surface. For the reciprocal of the quartic surface is clearly the envelope of the reciprocal of the variable quadric surface; this reciprocal is itself a quadric surface, and the reciprocal of the quartic surface is thus given in the same form as the original surface, viz., as the envelope of a quadric surface the equation whereof contains rationally the variable parameter $\theta$; the equation of the reciprocal surface is consequently obtained by equating to zero the discriminant in regard to $\theta$, of the equation of the reciprocal quadric surface.

It is to be observed that, inasmuch as the equation of the reciprocal quadric surface is of the third degree in the coefficients of the original quadric, it is in general of the degree 6 in the parameter; we have thus a sextic function of $\theta$, the coefficients whereof are quadric functions of the coordinates; and the discriminant is a function of the order 10 in these coefficients, that is, of the order 20 in the coordinates. The reciprocal of the quartic surface is thus a surface of the order 20; this is right, for in a general quartic surface the order of the reciprocal surface $= 36$, and the 8 nodes reduce the order by 16; $36 - 16 = 20$.

In the equation $UW - V^{2} = 0$, or say $V^{2} - UW = 0$; if $U$ reduce itself to the square of a linear function, $U = P^{2}$, the equation becomes $V^{2} - P^{2}W = 0$, which is the general form of the quartic surface having the nodal conic $V=0$, $P=0$. And moreover, if, $W$ be the product of this same linear function $P$ by another linear function $Q$, $W = PQ$, then the form is $V^{2} - P^{2}Q = 0$, which is the general form of the quartic surface having the cuspidal conic $V=0$, $P=0$.

Writing for greater convenience $x$, $y$ in the place of $P$, $Q$ respectively, we have the quartic
\[
(AB)\qquad V^{2} - x^{2}y = 0,
\]
having the cuspidal conic $V=0$, $x=0$; and which has besides the conic of plane contact $V=0$, $y=0$. In virtue of the cuspidal conic the reciprocal surface should be of the order 6; and by the foregoing method of obtaining the equation of the reciprocal surface, I will verify that this is so.

To effect this as simply as possible, I fix the remaining coordinates $z$, $w$ as follows. The line $x=0$, $y=0$ is not in general a tangent to the surface $V=0$; it therefore meets this surface in two points, and we may take $z=0$, $w=0$ to be the equations of the tangent planes at these two points respectively to $V=0$; we have thus
\[
V = ax^{2} + 2hxy + by^{2} + 2nzw.
\]
Introducing for convenience the numerical factor 2, and taking the equation of the surface to be
\[
(ax^{2} + 2hxy + by^{2} + 2nzw)^{2} - x^{2}y = 0,
\]
this is the envelope of the quadric surface
\[
\theta^{2} ax^{2} + 2\theta(ax^{2} + 2hxy + by^{2} + 2nzw) + 2xy = 0,
\]
which is a surface $(a,b,c,d,f,g,h,l,m,n\hspace{-0.2em}\pitchfork x,y,z,w)^{2} = 0$, where $a=\theta^{2} + 2a\theta$, $b=2b\theta$, $h=2h\theta+1$, $f=2\theta n$, and where all the other coefficients vanish. Assuming, as usual, that the reciprocation is effected in regard to the surface $x^{2} + y^{2} + z^{2} + w^{2} = 0$, the general equation is
\begin{align*}
&x^{2} \cdot \Delta(bc - f^{2}) - cm^{2} - bn^{2} + 2fmn\\
+\ &y^{2} \cdot \Delta(ca - g^{2}) - an^{2} - cl^{2} + 2gnl\\
+\ &z^{2} \cdot \Delta(ab - h^{2}) - bl^{2} - am^{2} + 2hlm\\
+\ &w^{2} \cdot \Delta abc - af^{2} - bg^{2} - ch^{2} + 2fgh\\
+\ &2yz \cdot \Delta(gh - af) + l^{2}f + amn - hnl - glm\\
+\ &2zx \cdot \Delta(hf - bg) + m^{2}g + bnl - fhm - hmn\\
+\ &2xy \cdot \Delta(fg - ch) + n^{2}h + clm - gmn - fnl\\
+\ &2xw \cdot -\Delta\!\left[(bc-f^{2}) - n(hf-bg) - m(fg-ch)\right]\\
+\ &2yw \cdot -\Delta\!\left[(ca-g^{2}) - l(fg-ch) - n(gh-af)\right]\\
+\ &2zw \cdot -\Delta\!\left[(ab-h^{2}) - m(gh-af) - l(hf-bg)\right],
\end{align*}
(I write down this general result as it will be useful for reference in other cases) in the present case this becomes simply
\[
x^{2} \cdot -bn^{2} + y^{2} \cdot -an^{2} + 2z^{2} \cdot nh^{2} + 2zw \cdot (-nab + nh^{2}) = 0,
\]
where $a$, $b$, $n$, $h$ have the foregoing values; the equation is thus only of the order 4 in regard to $\theta$; but it in fact divides by $n (= 2\theta n)$ and thus reduces itself to the third order, viz. it becomes
\[
n(bx^{2} + ay^{2}) - 2h^{2}xy + 2(ab - h^{2})zw = 0,
\]
or, substituting for $a$, $b$, $n$, $h$ their values, this is
\begin{multline*}
x^{2} \cdot 4b n\theta^{2} + y^{2} \cdot 2n(2n\theta^{2} + 4ab\theta) + 2zw \cdot (2b\theta^{2} + 4ab\theta)\\
- 2(xy + zw)(2\theta h + 1)^{2} = 0,
\end{multline*}
say this is
\[
(A, B, C, D \hspace{-0.2em}\pitchfork \theta, 1)^{2} = 0,
\]
where $A$, $B$, $C$, $D$ are each of them a quadric function of the coordinates; it being observed that $C$ and $D$ are respectively numerical multiples of the same function $(xy + zw)$. Hence equating the discriminant to zero, we have
\[
A^{2}D^{2} + 4C^{3} + 4B^{2}D - 3B^{2}C^{2} - 6ABCD = 0,
\]
which equation, inasmuch as every term contains either $C$ or $D$ as a factor, divides by $xy + zw$, and thus becomes an equation of the order 6 in the coordinates: that is the order of the reciprocal surface is $=6$. Multiplying by $\frac{4}{3}$ to avoid fractions, the actual values of $A$, $B$, $C$, $D$ are
\begin{align*}
A &= \tfrac{4}{3}(by^{2} + 2hzw),\\
B &= 2\{h(ax^{2} + by^{2}) + 2ahzw - 2h^{2}(xy + zw)\},\\
C &= -4h(xy + zw),\\
D &= -\tfrac{4}{3}(xy + zw),
\end{align*}
or say $A = \frac{4}{3}\alpha$, $B = 2\beta$, $C = -4h\gamma$, $D = -\frac{4}{3}\gamma$; where $\gamma = xy + zw$; substituting these values and omitting the factor $\frac{4}{3}\gamma$, the equation is
\[
27\alpha^{2}\gamma^{2} - 256h^{3}\alpha\gamma^{2} - 32\beta^{3} - 64h\beta^{2}\gamma - 144h\alpha\beta^{2}\gamma = 0,
\]
which is an equation of the form $(\asymp\!\asymp\!\asymp \alpha, \beta, \gamma)^{2} = 0$. The sextic surface has thus singular points $\alpha=0$, $\beta=0$, $\gamma=0$, viz. these are the two points $(x=0,\ y=0,\ z=0)$, $(x=0,\ y=0,\ w=0)$ each four times. The further discussion of the sextic surface is reserved for another occasion.

I do not at present attempt to enumerate the particular cases of the surface $V^{2} - x^{2}y = 0$, but content myself with the discussion of a particular case in which the order of the reciprocal surface is $=3$. Suppose that $V=0$ is a cone, $y=0$ a tangent plane to the cone (so that the conic $y=0$, $V=0$ breaks up into a line twice repeated), $x=0$ an arbitrary plane (so that we have still the proper cuspidal conic $x=0$, $V=0$). Any other tangent plane of the cone may be taken for the plane $z=0$; the plane containing the lines of contact of the two tangent planes for the plane $w=0$; the equation of the conic then is $V = dw^{2} + 2fyz = 0$; and the equation of the surface is
\[
(AB)\qquad (dw^{2} + 2fyz)^{2} - x^{2}y = 0.
\]

For convenience of comparison, I change $x$, $y$, $w$, $z$ into $y$, $w$, $z$, $x$, and assign numerical values to the coefficients, writing the equation under the form
\[
(AB)\qquad 27(4zw + x^{2})^{2} - 64y^{2}w = 0.
\]
The quartic is here the envelope of the quadric surface
\[
\theta^{2} \cdot 4yw + 2\theta \cdot 9(4zw + x^{2}) + 12y^{2} = 0,
\]
viz., comparing with the general form
\[
(a,b,c,d,f,g,h,l,m,n\hspace{-0.2em}\pitchfork x,y,z,w)^{2} = 0,
\]
we have $b=12$, $c=9\theta$, $l=18\theta$, $m=2\theta^{2}$, all the other coefficients vanishing. The reciprocal equation is
\[
x^{2}(-cm^{2}) + y^{2}(-cl^{2}) + z^{2}(-bl^{2}) + 2xw \cdot dlm + 2yz \cdot (-lbc) = 0,
\]
or substituting for $b$, $c$, $l$, $m$ their values, this is found to be
\[
\theta(\theta x - 9y)^{2} + 108(z^{2} + xw) = 0.
\]
Representing this by $(A, B, C, D \hspace{-0.2em}\pitchfork \theta, 1)^{2} = 0$, the discriminant in regard to $\theta$ would, in virtue of the values of $A$, $B$, $C$, contain $D$ as a factor; the reason of this appears from the original form; in fact, forming the derived equation in regard to $\theta$, this is found to be $(\theta x - 9y)(\theta x - 3y) = 0$; the value $\theta x - 9y = 0$ gives as a factor of the discriminant $z^{2} - 8xw$; the value $\theta x - 3y = 0$ gives $3y(-6y)^{2} + 108x(z^{2} + xw)$, that is the factor $y^{2} + x(z^{2} + xw)$; the complete value of the discriminant as obtained by substitution of the values of $A$, $B$, $C$, $D$ being $x^{2}\{z^{2} + xy\}\{y^{2} + x(z^{2} + xw)\}$; the equation of the reciprocal surface is
\[
y^{2} + x(z^{2} + xw) = 0,
\]
viz. this is a cubic surface. Schl\"{a}fli's Prof.~Case~xx., having a uniplanar point $x=0$, $y=0$, $z=0$ reducing the class by 8, and so giving a reciprocal surface of the order $(12 - 8 =)\ 4$, viz. the surface $27(4xw + z^{2})^{2} - 64y^{2}w = 0$. See the Memoir, Schl\"{a}fli, ``On the distribution of surfaces of the third order into species in reference to the absence or presence of singular points and the reality of their lines,'' \textit{Phil.\ Trans.}, vol.~CLIII, (1853), pp.~193--241.

I pass to the case of a surface
\[
V^{2} - P^{2}U = 0,
\]
having a nodal conic $V=0$, $P=0$, but not having in general any nodes. And I propose to show how the constants may be determined so that the surface shall have 1, 2, 3, or 4 nodes. It is to be remarked that in the above equation the plane $P=0$ is a determinate plane, but the quadric surface $V=0$ is not a determinate quadric, we may in fact substitute for it the quadric $V + \lambda P^{2} = 0$, writing the equation under the form
\[
(V + \lambda P^{2})^{2} - P^{2}(U + 2\lambda V + \lambda^{2}P^{2}),
\]
so that we may without loss of generality, by means of the disposable constant $\lambda$, subject the surface $V=0$ to any single condition; for instance, take it to be a cone, or to pass through a given point, \&c.

Taking the planes $x=0$, $y=0$, $z=0$, $w=0$ to be arbitrary planes, the implicit constant factors in these equations may be determined in such wise that the equation of the given plane $P=0$ shall be $x + y + z + w = 0$. The equation of the surface will then be
\begin{multline*}
\{(a,b,c,d,f,g,h,l,m,n\hspace{-0.2em}\pitchfork x,y,z,w)^{2}\}^{2}\\
= (x+y+z+w)^{2} \cdot (a',b',c',d',f',g',h',l',m',n'\hspace{-0.2em}\pitchfork x,y,z,w)^{2},
\end{multline*}
and we may assume that the node or nodes (if any) lie at a vertex or vertices of the tetrahedron $x=0$, $y=0$, $z=0$, $w=0$, say at the points $A$, $B$, $C$, $D$. The conditions for a node at each of these points are at once found to be

\begin{center}
\begin{tabular}{c|c|c|c}
Node at $A$, & Node at $B$, & Node at $C$, & Node at $D$,\\\hline
$2a^{2} = a' + a'\vphantom{^{2}}$ & $2bh = h' + b'$ & $2cg = g' + c'$ & $2dl = l' + d'$\\
$2ah = h' + a'$ & $2bf = f' + b'$ & $2cf = f' + c'$ & $2dm = m' + d'$\\
$2ag = g' + a'$ & $2bh = f' + h'$ & $2c^{2} = c' + c'$ & $2dn = n' + d'$\\
$2al = l' + a'$ & $2bm = m' + b'$ & $2cn = n' + c'$ & $2d^{2} = d' + d'$
\end{tabular}
\end{center}

The first set of equations gives
\[
a' = a^{2},\quad h' = 2ah - a^{2},\quad g' = 2ag - a^{2},\quad l' = 2al - a^{2}.
\]

If the first and second sets are satisfied simultaneously, we have $2(a-b)h = a^{2} - b^{2}$, that is $a = b$, or else $h = \frac{1}{2}(a+b)$; that is, the two sets may be satisfied in two different ways according as $a$ and $b$ are equal or unequal. Similarly the first, second, and third sets may be satisfied in three different ways and the four sets in five different ways according as there are or are not any equalities between $a$, $b$, $c$, and between $a$, $b$, $c$ and $d$ respectively. The several solutions are shown in the annexed table, viz., in the line I no set is satisfied; in the line II only the first set; in the lines III and IV the first and second sets; in the lines V to VII the first, second, and third sets; and in the lines VIII to XII the four sets.


{\scriptsize
\begin{center}
\textbf{Table of Solutions.}
\end{center}
\begin{center}
\begin{tabular}{c|cccc|cccccc|cccc|cccccc}
& $a$ & $b$ & $c$ & $d$ & $f$ & $g$ & $h$ & $l$ & $m$ & $n$ & $a'$ & $b'$ & $c'$ & $d'$ & $f'$ & $g'$ & $h'$ & $l'$ & $m'$ & $n'$\\\hline
I & $a$ & $b$ & $c$ & $d$ & $f$ & $g$ & $h$ & $l$ & $m$ & $n$ & $a^{2}$ & $b^{2}$ & $c^{2}$ & $d^{2}$ & $2bf-b^{2}$ & $2ag-a^{2}$ & $2ah-a^{2}$ & $2al-a^{2}$ & $2bm-b^{2}$ & $2cn-c^{2}$\\
II & $\overline{a}$ & $b$ & $c$ & $d$ & $f$ & $g$ & $h$ & $l$ & $m$ & $n$ & $0$ & $b^{2}$ & $c^{2}$ & $d^{2}$ & $2bf-b^{2}$ & $2ag-a^{2}$ & $2ah-a^{2}$ & $2al-a^{2}$ & $2bm-b^{2}$ & $2cn-c^{2}$\\
III & $\overline{a}$ & $\overline{b}$ & $c$ & $d$ & $f$ & $g$ & $\tfrac{1}{2}(a+b)$ & $l$ & $m$ & $n$ & $0$ & $0$ & $c^{2}$ & $d^{2}$ & $2cf-c^{2}$ & $2cg-c^{2}$ & $0$ & $2cl-c^{2}$ & $2cm-c^{2}$ & $2cn-c^{2}$\\
IV & $\overline{a}$ & $\overline{b}$ & $c$ & $d$ & $f$ & $g$ & $h$ & $l$ & $m$ & $n$ & $0$ & $0$ & $c^{2}$ & $d^{2}$ & $2cf-c^{2}$ & $2ag-a^{2}$ & $2ah-a^{2}$ & $2al-a^{2}$ & $2cm-c^{2}$ & $2cn-c^{2}$\\
V & $\overline{a}$ & $\overline{b}$ & $\overline{c}$ & $d$ & $\tfrac{1}{2}(b+c)$ & $\tfrac{1}{2}(a+c)$ & $\tfrac{1}{2}(a+b)$ & $l$ & $m$ & $n$ & $0$ & $0$ & $0$ & $d^{2}$ & $0$ & $0$ & $0$ & $2dl-d^{2}$ & $2dm-d^{2}$ & $2dn-d^{2}$\\
VI & $\overline{a}$ & $\overline{b}$ & $\overline{c}$ & $d$ & $\tfrac{1}{2}(b+c)$ & $\tfrac{1}{2}(a+c)$ & $h$ & $l$ & $m$ & $n$ & $0$ & $0$ & $0$ & $d^{2}$ & $bc$ & $ac$ & $2ah-a^{2}$ & $2al-a^{2}$ & $2bm-b^{2}$ & $2cn-c^{2}$\\
VII & $\overline{a}$ & $\overline{b}$ & $\overline{c}$ & $d$ & $f$ & $g$ & $h$ & $l$ & $m$ & $n$ & $0$ & $0$ & $0$ & $d^{2}$ & $2bf-b^{2}$ & $2ag-a^{2}$ & $2ah-a^{2}$ & $2al-a^{2}$ & $2dm-d^{2}$ & $2dn-d^{2}$\\
VIII & $\overline{a}$ & $\overline{b}$ & $\overline{c}$ & $\overline{d}$ & $\tfrac{1}{2}(b+c)$ & $\tfrac{1}{2}(a+c)$ & $\tfrac{1}{2}(a+b)$ & $\tfrac{1}{2}(a+d)$ & $\tfrac{1}{2}(b+d)$ & $\tfrac{1}{2}(c+d)$ & $0$ & $0$ & $0$ & $0$ & $0$ & $0$ & $0$ & $0$ & $0$ & $0$\\
IX & $\overline{a}$ & $\overline{b}$ & $\overline{c}$ & $\overline{d}$ & $\tfrac{1}{2}(b+c)$ & $\tfrac{1}{2}(a+c)$ & $\tfrac{1}{2}(a+b)$ & $l$ & $m$ & $n$ & $0$ & $0$ & $0$ & $0$ & $0$ & $0$ & $0$ & $2al-a^{2}$ & $2bm-b^{2}$ & $2cn-c^{2}$\\
X & $\overline{a}$ & $\overline{b}$ & $\overline{c}$ & $\overline{d}$ & $\tfrac{1}{2}(b+c)$ & $g$ & $\tfrac{1}{2}(a+b)$ & $\tfrac{1}{2}(a+d)$ & $m$ & $n$ & $0$ & $0$ & $0$ & $0$ & $2cg-c^{2}$ & $0$ & $0$ & $2al-a^{2}$ & $2dm-d^{2}$ & $2cn-c^{2}$\\
XI & $\overline{a}$ & $\overline{b}$ & $\overline{c}$ & $\overline{d}$ & $f$ & $\tfrac{1}{2}(a+c)$ & $h$ & $\tfrac{1}{2}(a+d)$ & $\tfrac{1}{2}(b+d)$ & $n$ & $0$ & $0$ & $0$ & $0$ & $2bf-b^{2}$ & $0$ & $2ah-a^{2}$ & $0$ & $0$ & $2dn-d^{2}$\\
XII & $\overline{a}$ & $\overline{b}$ & $\overline{c}$ & $\overline{d}$ & $f$ & $g$ & $h$ & $l$ & $m$ & $n$ & $0$ & $0$ & $0$ & $0$ & $2bf-b^{2}$ & $2ag-a^{2}$ & $2ah-a^{2}$ & $2al-a^{2}$ & $2bm-b^{2}$ & $2cn-c^{2}$
\end{tabular}
\end{center}
}%


I is the general case, $V^{2} = P^{2}U$, of a quartic and a nodal conic but without nodes.\hfill(AC)

II is the case of a single node; writing, as without loss of generality we may do, $a=0$, the equation is
\begin{multline*}
[(0,b,c,d,f,g,h,l,m,n\hspace{-0.2em}\pitchfork x,y,z,w)^{2}]^{2}\\
= (x+y+z+w)^{2} \cdot (b,c,d,m,l,f\hspace{-0.2em}\pitchfork y,z,w)^{2},
\end{multline*}
viz. the quadric $U=0$ is here a cone having its vertex on the quadric $V=0$.\hfill(AD)

III and IV are two cases each of them with two nodes, viz. III, the equation is
\begin{multline*}
\{(ax+by)(x+y) + cz^{2} + 2nzw + dw^{2} + 2x(gz+lw) + 2y(fz+mw)\}^{2}\\
= (x+y+z+w)^{2}\, \{(ax+by)^{2} + c'z^{2} + 2n'zw + d'w^{2}\\
+ 2x[(2ag-a^{2})z + (2al-a^{2})w]\\
+ 2y[(2bf-b^{2})z + (2bm-b^{2})w]\},\hfill(AE)
\end{multline*}
where it is to be observed that the line $z=0$, $w=0$ joining the two nodes $(y=0,z=0,w=0)$ and $(x=0,z=0,w=0)$ is a line on the surface. Writing, as we may do, $a=0$, the equation assumes the more simple form
\begin{multline*}
\{by(x+y) + cz^{2} + 2nzw + d'w^{2} + 2x(gz+lw) + 2y(fz+mw)\}^{2}\\
= (x+y+z+w)^{2}\, [b'y^{2} + c'z^{2} + 2n'zw + d'w^{2}\\
+ 2y\{(2bf-b^{2})z + (2bm-b^{2})w\}].\hfill(AE)
\end{multline*}

\newpage
\setcounter{page}{310}

In IV the equation is
\begin{multline*}
\{a(x^{2}+y^{2}) + 2hxy + cz^{2} + dw^{2} + 2nzw + 2x(gz+lw) + 2y(fz+mw)\}^{2}\\
= (x+y+z+w)^{2}\, [a^{2}(x^{2}+y^{2}) + 2(2ah-a^{2})xy + c'z^{2} + 2n'zw + d'w^{2}\\
+ 2x\{(2ag-a^{2})z + (2al-a^{2})w\}\\
+ 2y\{(2af-a^{2})z + (2am-a^{2})w\}],\hfill(AF)
\end{multline*}
where it will be observed that the line $z=0$, $w=0$ joining the two nodes is not a line on the surface.

Writing, as we may do, $a=0$, the equation becomes
\begin{multline*}
\{2hxy + cz^{2} + 2nzw + dw^{2} + 2x(gz+lw) + 2y(fz+mw)\}^{2}\\
= (x+y+z+w)^{2}\, (c'z^{2} + 2h'zw + d'w^{2}),\hfill(AF)
\end{multline*}
viz. the form is $F^{2} = P^{2}QR$, the quadric surface $U=0$ breaking up into the two planes $Q=0$, $R=0$; and the nodes being situate at the intersections of the line $Q=0$, $R=0$ with the surface $F=0$.

V, VI, VII are apparently cases with three nodes, but in fact VI is the only case of a proper quartic surface with three nodes. For in V the equation is
\begin{multline*}
\{(ax+by+cz)(x+y+z) + dw^{2} + 2w(lx+my+nz)\}^{2}\\
= (x+y+z+w)^{2}\, [(ax+by+cz)^{2} + d'w^{2}\\
+ 2w\{(2al-a^{2})x + (2bm-b^{2})y + (2cn-c^{2})z\}],
\end{multline*}
which is satisfied by $w=0$, and the surface thus breaks up into the plane $w=0$ and a cubic surface.

And in VII the equation is
\begin{multline*}
[a(x^{2}+y^{2}+z^{2}) + dw^{2} + 2fyz + 2gzx + 2hxy + 2lxw + 2myw + 2nzw]^{2}\\
= (x+y+z+w)^{2}\, [a^{2}(x^{2}+y^{2}+z^{2}) + d'w^{2}\\
+ 2\{(2af-a^{2})yz + (2ag-a^{2})zx + (2ah-a^{2})xy\\
+ (2al-a^{2})xw + (2am-a^{2})yw + (2an-a^{2})zw\}],
\end{multline*}
which putting $a=0$ is
\begin{multline*}
\{dw^{2} + 2fyz + 2gzx + 2hxy + 2lxw + 2myw + 2nzw\}^{2}\\
= (x+y+z+w)^{2}\, d'w^{2},
\end{multline*}
viz. this is a pair of quadric surfaces.

In the remaining case VI the equation is
\begin{multline*}
\{a(x^{2}+y^{2}) + 2hxy + cz^{2} + dw^{2} + (a+c)(yz+zx) + 2w(lx+my+nz)\}^{2}\\
= (x+y+z+w)^{2}\, [a^{2}(x^{2}+y^{2}) + d'z^{2} + d'w^{2} + 2cz(x+y) + 2(2ah-a^{2})xy\\
+ 2w\{(2al-a^{2})x + (2am-a^{2})y + (2an-a^{2})z\}],\hfill(AG)
\end{multline*}
which putting therein $a=0$ is
\begin{multline*}
\{2hxy + dw^{2} + cz(x+y+z) + 2w(lx+my+nz)\}^{2}\\
= (x+y+z+w)^{2}\, \{c'z^{2} + d'w^{2} + 2(2cn-c^{2})zw\},\hfill(AG)
\end{multline*}
which is a surface having the nodes $A$, $B$, $C$; and it is to be observed that the lines $CA$, $CB$, but not the line $AB$, are lines on the surface.

IX, X, XI, XII are apparently cases with four nodes, but it is only XI which is a proper quartic with four nodes. In fact IX is
\begin{multline*}
\{(ax+ay+cz+dw)(x+y+z+w) + 2(h-a)xy\}^{2}\\
= (x+y+z+w)^{2}\, \{(ax+ay+cz+dw)^{2} + 4a(h-a)xy\},
\end{multline*}
which is satisfied if $x=0$ or if $y=0$; that is, the surface breaks up into the two planes $x=0$, $y=0$, and a quadric surface.

X is
\begin{multline*}
\{a(x^{2}+y^{2}+z^{2}) + dw^{2} + 2fyz + 2gzx + 2hxy\\
+ (a+d)(xw+yw+zw)\}^{2}\\
= (x+y+z+w)^{2}\, [a^{2}(x^{2}+y^{2}+z^{2}) + d'w^{2}\\
+ 2(2af-a^{2})yz + 2(2ag-a^{2})zx + 2(2ah-a^{2})xy\\
+ 2ad(xw+yw+zw)],
\end{multline*}
which putting therein $a=0$ is
\begin{multline*}
\{dw(x+y+z+w) + 2fyz + 2gzx + 2hxy\}^{2}\\
= (x+y+z+w)^{2}\, d'w^{2},
\end{multline*}
and thus breaks up into two quadrics.

And XII is
\begin{multline*}
[a^{2}(x^{2}+y^{2}+z^{2}+w^{2}) + 2fyz + 2gzx + 2hxy + 2lxw + 2myw + 2nzw]^{2}\\
= (x+y+z+w)^{2}\, [a^{2}(x^{2}+y^{2}+z^{2}+w^{2})\\
+ 2(2af-a^{2})yz + 2(2ag-a^{2})zx + 2(2ah-a^{2})xy\\
+ 2(2al-a^{2})xw + 2(2am-a^{2})yw + 2(2an-a^{2})zw],
\end{multline*}
which putting $a=0$ is
\[
(2fyz + 2gzx + 2hxy + 2lxw + 2myw + 2nzw)^{2} = 0,
\]
and is thus a quadric surface twice repeated.

There remains XI, and here the equation is
\begin{multline*}
\{(ax+ay+cz+cw)(x+y+z+w) + 2(h-a)xy + 2(n-c)zw\}^{2}\\
= (x+y+z+w)^{2}\, \{(ax+ay+cz+cw)^{2} + 4a(h-a)xy + 4c(n-c)zw\},
\end{multline*}
or writing $h+a$, $n+c$ in place of $h$, $n$ respectively, this is
\begin{multline*}
\{(ax+ay+cz+cw)(x+y+z+w) + 2hxy + 2nzw\}^{2}\\
= (x+y+z+w)^{2}\, \{(ax+ay+cz+cw)^{2} + 4ahxy + 4cnzw\},\hfill(AH)
\end{multline*}
or putting herein $a=0$ it is
\[
\{c(z+w)(x+y+z+w) + 2hxy + 2nzw\}^{2} - c(x+y+z+w)^{2}\{c(z+w)^{2} + 4nzw\},\hfill(AH)
\]
which may also be written
\[
c(x+y+z+w)\{hxy(z+w) - nzw(x+y)\} + 4hxy \cdot cnzw = 0,\hfill(AH)
\]
the equation of a quartic surface with the four nodes $A$, $B$, $C$, $D$; it is to be observed that the lines $AC$, $AD$, $BC$, $BD$ are, the lines $AB$ and $CD$ are not, lines on the surface.

A more simple form may be given to the equation as follows; using the second of the above forms, multiplying the equation by 4, and writing therein
\begin{align*}
p &= \surd(c)\, (x+y+z+w),\\
q,\ r &= c(z+w) \pm 2\surd(cn)\, zw,\\
s,\ t &= c(x+y) \pm 2\surd(ah)\, xy,
\end{align*}
$q$, $r$, and $s$, $t$ being the linear factors of the two quadric functions respectively, we have
\[
qr - st = c(x+y+z+w)(-x-y+z+w) + 4hxy + 4nzw,
\]
and thence
\[
p^{2} + qr - st = 2c(z+w)(x+y+z+w) + 4hxy + 4nzw,
\]
wherefore the equation is
\[
(p^{2} + qr - st)^{2} = 4p^{2}qr,\hfill(AH)
\]
or, what is the same thing,
\[
p^{2} + \surd(qr) + \surd(st) = 0,\hfill(AH)
\]
where $p$, $q$, $r$, $s$, $t$ are any linear functions of the coordinates; this is the equation of a quartic surface having the nodal conic $p=0$, $qr-st=0$; and the four nodes $(q=0,\ r=0,\ p^{2}-st=0)$ and $(s=0,\ t=0,\ p^{2}-qr=0)$. It includes the Cyclide, the equation of which may be written
\[
b\xi^{2} = \surd\{(ax - eky + b'z)^{2} + b'^{2}z'^{2}\} + \surd\{(e'a - aky - b'z)^{2} + b'^{2}z'^{2}\}.
\]

I remark that Prof.~Kummer in his most valuable Memoir, ``Ueber die Fl\"{a}chen vierten Grades auf welchen Schaaren von Kegelschnitten liegen,'' \textit{Crelle}, t.~LXVI. (1864), pp.~66--76, has considered several of the cases of a quartic surface with a nodal conic, viz. no node, (AC); a single node, (AD); two nodes (the case AF); and four nodes, (AH); but he has not considered two nodes, the case (AE); nor three nodes, (AG).

In reference to the general case of a quartic surface with a nodal conic, some most interesting properties have recently been obtained by Prof.~Clebsch, see \textit{Berl.\ Monatsh.}, April 30, 1868, where it is shown that there are on the surface 16 right lines forming 20 systems of double-fours, analogous in some respect to the 27 lines and 36 systems of double-sixes of a cubic surface.

\newpage
\setcounter{page}{314}

%===============================================================================
\article{458}{ON THE ANHARMONIC-RATIO SEXTIC.}
%===============================================================================

\noindent[From the \textit{Quarterly Journal of Pure and Applied Mathematics}, vol.~x. (1870), pp.~56, 57.]

\vspace{0.5em}

Mr Walker's equation is $A\Omega\Lambda^{2} - 1 + iy + P(\lambda^{2} - \lambda)^{2} = 0$; changing the sign of $\lambda$, and also the numerical multipliers of $I$, $A$ (so as to convert the discriminant equation into its standard form $\Delta = P^{2} - 27J^{2}$), the equation is
\[
4A(\lambda^{3} + \lambda + 1)^{2} - 27I^{2}(\lambda^{2} + \lambda)^{2} = 0.
\]

I remark that this is most readily obtained as follows; writing
\begin{align*}
A &= (a-d)(b-c),\\
B &= (b-d)(c-a),\\
C &= (c-d)(a-b),
\end{align*}
then we have $A + B + C = 0$,
\begin{align*}
I &= -\tfrac{1}{4}(A^{2} + B^{2} + C^{2}) = -\tfrac{1}{2}(BC + CA + AB),\\
J &= \tfrac{1}{4\cdot 27}(B-C)(C-A)(A-B),\\
\surd(\Delta) &= \tfrac{1}{4\cdot 27}ABC,
\end{align*}
see my Fifth Memoir on Quantics, \textit{Phil.\ Trans.}, vol.~CXLVIII. (1858), pp.~429--460, [156]. And observe also, that in virtue of the relation $A + B + C = 0$, we have
\[
-12I = A^{2} + AB + B^{2} = A^{2} + AC + C^{2} = B^{2} + BC + C^{2}.
\]
Hence writing
\[
u = \frac{A\lambda + B}{\lambda + 1},
\]
when $\lambda$ has any one of the values
\[
\frac{A}{B},\ \frac{B}{A},\ -\frac{A}{C},\ -\frac{C}{A},\ -\frac{B}{C},\ -\frac{C}{B},
\]
we see that $u$ assumes only the values $A$, $B$, $C$, and $u$ is thus determined by the equation
\[
u^{3} - 12Iu - 16\surd(\Delta) = 0.
\]
Eliminating $u$, we obtain
\[
16\surd(\Delta)\left[\left(\frac{A\lambda+B}{\lambda+1}\right)^{3} - 12I\left(\frac{A\lambda+B}{\lambda+1}\right) - 16\surd(\Delta)\right] = 0,
\]
or, what is the same thing,
\[
\left(\frac{A\lambda+B}{\lambda+1}\right)^{3} + \frac{1}{2}(A^{2}+AB+B^{2})\left(\frac{A\lambda+B}{\lambda+1}\right) - \frac{1}{4\cdot 27}ABC = 0,
\]
that is
\[
4A(\lambda^{2}+\lambda+1)^{2} - 27I^{2}\lambda^{2}(\lambda+1)^{2} = 0,
\]
the required equation.

\newpage
\setcounter{page}{316}

%===============================================================================
\article{459}{ON THE DOUBLE-SIXERS OF A CUBIC SURFACE.}
%===============================================================================

\noindent[From the \textit{Quarterly Journal of Pure and Applied Mathematics}, vol.~x. (1870), pp.~58--71.]

\vspace{0.5em}

The 27 lines on a cubic surface include, and that in 36 different ways, a double-sixer; viz. a system of two sets of six lines
\[
1,\ 2,\ 3,\ 4,\ 5,\ 6;\qquad 1',\ 2',\ 3',\ 4',\ 5',\ 6',
\]
such that every line of the one set intersects all the non-corresponding lines of the other set, thus
\begin{center}
\begin{tabular}{c|cccccc}
& 1' & 2' & 3' & 4' & 5' & 6'\\\hline
1 & $\cdot$ & $\times$ & $\times$ & $\times$ & $\times$ & $\times$\\
2 & $\times$ & $\cdot$ & $\times$ & $\times$ & $\times$ & $\times$\\
3 & $\times$ & $\times$ & $\cdot$ & $\times$ & $\times$ & $\times$\\
4 & $\times$ & $\times$ & $\times$ & $\cdot$ & $\times$ & $\times$\\
5 & $\times$ & $\times$ & $\times$ & $\times$ & $\cdot$ & $\times$\\
6 & $\times$ & $\times$ & $\times$ & $\times$ & $\times$ & $\cdot$
\end{tabular}
\end{center}
there being in all 30 intersections.

Any line, say 4, of the one set, intersects five lines $1'$, $2'$, $3'$, $5'$, $6'$ of the other set; and these six lines being given the double-sixer may be constructed; viz. (besides the line 4) we have a line 1 meeting the lines $2'$, $3'$, $5'$, $6'$; a line 2 meeting the lines $3'$, $5'$, $6'$, $1'$; a line 3 meeting the lines $5'$, $6'$, $1'$, $2'$; a line 5 meeting the lines $6'$, $1'$, $2'$, $3'$; and a line 6 meeting the lines $1'$, $2'$, $3'$, $5'$; and then the lines 1, 2, 3, 5, 6 are all of them met by a single line $4'$, which completes the system.

We may, if we please, consider the lines 4, 2 as given, and then $1'$, $3'$, $5'$, $6'$ will be any four lines each of them meeting the two given lines 4, 2; $2'$ will be any line meeting 4; and we have to determine a line $4'$ meeting 2, such that there may exist the lines 1, 3, 5, 6, completing the system as above. Or what is the same thing, we have a skew quadrilateral $1'$, 2, $3'$, 4; $5'$ and $6'$ meet 2 and 4; $2'$ meets 4, and $4'$ meets 2: 5 and 6 meet $1'$ and $3'$; 1 meets $3'$ and 3 meets $1'$; and the two sets $2'$, $4'$, $5'$, $6'$ and 1, 3, 5, 6 meet thus
\begin{center}
\begin{tabular}{c|cccc}
& 1 & 3 & 5 & 6\\\hline
$2'$ & $\cdot$ & $\times$ & $\times$ & $\times$\\
$4'$ & $\times$ & $\cdot$ & $\times$ & $\times$\\
$5'$ & $\times$ & $\times$ & $\cdot$ & $\times$\\
$6'$ & $\times$ & $\times$ & $\times$ & $\cdot$
\end{tabular}
\end{center}

Hence, starting with the skew quadrilateral $1'2\,3'4$, and taking $x=0$, $y=0$, $z=0$, $w=0$ for the equations of the four planes $41'$, $1'2$, $23'$, $3'4$ respectively or what is the same thing $x=0$, $y=0$ for the equations of the line $1'$; $y=0$, $z=0$ for those of the line 2; $z=0$, $w=0$ for those of the line $3'$; and $w=0$, $x=0$ for those of the line 4; the several lines may be determined, each of them by means of its six coordinates, as follows
\begin{center}
\begin{tabular}{c|cccccc}
& $a$ & $b$ & $c$ & $f$ & $g$ & $h$\\\hline
$1'$ & & & & & & 1\\
2 & & & & 1 & &\\
$3'$ & & & 1 & & &\\
4 & 1 & & & & &\\
$2'$ & $A_{2}$ & $B_{2}$ & $C_{2}$ & & $G_{2}$ & $H_{2}$\\
$4'$ & & $B_{4}$ & $C_{4}$ & $F_{4}$ & $G_{4}$ & $H_{4}$\\
$5'$ & & $A_{5}$ & $C_{5}$ & & $G_{5}$ & $H_{5}$\\
$6'$ & & $B_{5}$ & $C_{5}$ & & $G_{6}$ & $H_{6}$\\
1 & $a_{1}$ & $b_{1}$ & $c_{1}$ & $f_{1}$ & $g_{1}$ &\\
3 & & & & $f_{3}$ & $g_{3}$ &\\
5 & & & & $f_{5}$ & $g_{5}$ &\\
6 & $a_{6}$ & $b_{6}$ & & $f_{6}$ & $g_{6}$ &
\end{tabular}
\end{center}
where
\begin{align*}
B_{4}A_{2} + C_{4}H_{2} &= 0,\\
B_{4}G_{5} + C_{5}H_{4} &= 0,\\
B_{4}G_{6} + C_{6}H_{4} &= 0,\\
B_{5}G_{5} + C_{5}H_{5} &= 0,\\
c_{1}f_{3} + b_{1}g_{1} &= 0,\\
f_{5}f_{2} + b_{5}g_{5} &= 0,\\
f_{6}f_{2} + b_{6}g_{6} &= 0,\\
a_{1}f_{6} + b_{1}g_{6} &= 0.
\end{align*}

\newpage
\setcounter{page}{318}

The conditions in regard to the intersections of the lines $2'$, $4'$, $5'$, $6'$ and 1, 3, 5, 6, are formed by means of the diagram
\begin{center}
\begin{tabular}{c|cccccccccccc}
& 1 & $f_{1}$ & $g_{1}$ & $a_{1}$ & $b_{1}$ & $c_{1}$ & $A_{2}$ & $B_{2}$ & $C_{2}$ & $G_{2}$ & $H_{2}$ & $2'$\\\hline
3 & $f_{3}$ & $g_{3}$ & $A_{3}$ & $C_{3}$ & $f_{3}$ & & $B_{4}$ & $C_{4}$ & $F_{4}$ & $G_{4}$ & $H_{4}$ & $4'$\\
5 & $f_{5}$ & $g_{5}$ & & $a_{5}$ & $b_{5}$ & & $B_{5}$ & $C_{5}$ & & $G_{5}$ & $H_{5}$ & $5'$\\
6 & & & & & & $B_{6}$ & $C_{6}$ & & & $G_{6}$ & $H_{6}$ & $6'$
\end{tabular}
\end{center}
viz. we have the equations

first set,
\begin{align*}
f_{1}A_{2} + g_{1}B_{5} + b_{1}G_{5} + c_{1}H_{5} &= 0,\\
g_{1}B_{5} + a_{1}f_{5} + b_{1}G_{5} + c_{1}H_{5} &= 0,\\
g_{3}B_{5} + b_{3}G_{5} + c_{3}H_{5} &= 0,\\
g_{3}B_{6} + b_{3}G_{6} + c_{3}H_{6} &= 0;
\end{align*}
second set,
\begin{align*}
f_{1}A_{2} + g_{1}B_{4} + h_{1}C_{4} + b_{1}G_{4} &= 0,\\
g_{1}B_{4} + h_{1}C_{4} + a_{1}F_{4} + b_{3}G_{4} &= 0,\\
g_{3}B_{4} + h_{3}C_{4} + b_{3}G_{4} &= 0,\\
g_{5}B_{4} + h_{5}C_{4} + b_{5}G_{4} &= 0;
\end{align*}
third set,
\begin{align*}
f_{3}A_{2} + g_{3}B_{5} + b_{3}G_{5} &= 0,\\
g_{3}B_{5} + a_{3}F_{5} + b_{3}G_{5} &= 0,\\
g_{5}B_{5} + b_{5}G_{5} &= 0;
\end{align*}
fourth set,
\begin{align*}
f_{3}A_{2} + g_{3}B_{6} + b_{3}G_{6} &= 0,\\
g_{3}B_{6} + a_{3}F_{6} + b_{3}G_{6} &= 0,\\
g_{5}B_{6} + b_{5}G_{6} &= 0;
\end{align*}
and it is to be shown, that taking as given the coordinates of $2'$, $5'$, $6'$, that is $(A_{2}, B_{2}, C_{2}, G_{2}, H_{2})$, $(B_{5}, C_{5}, G_{5}, H_{5})$ and $(B_{6}, C_{6}, G_{6}, H_{6})$, we can find the coordinates of the remaining lines $4'$, 1, 3, 5, 6.

The first set of equations gives
\[
(g_{1},\ b_{1},\ c_{1}) = \begin{vmatrix} B_{5}, & G_{5}, & H_{5}\\ B_{6}, & G_{6}, & H_{6}\end{vmatrix},
\]
viz. $g_{1}$, $b_{1}$, $c_{1}$ are proportional, but as only the ratios are material, they may be taken equal, to the determinants $G_{5}H_{6} - G_{6}H_{5}$, $H_{5}B_{6} - H_{6}B_{5}$, $B_{5}G_{6} - B_{6}G_{5}$. And then retaining $g_{1}$, $b_{1}$, $c_{1}$ to signify these values respectively, the first equation gives $f_{1}A_{2}$, and the second equation gives $a_{1}F_{5}$; multiplying together these values, and writing $a_{1}f_{1} = -b_{1}g_{1}$, we find
\[
(B_{5}B_{6},\ G_{5}G_{6},\ H_{5}H_{6},\ G_{5}H_{6}+G_{6}H_{5},\ H_{5}B_{6}+H_{6}B_{5},\ B_{5}G_{6}+B_{6}G_{5} + A_{2}F_{5}\hspace{-0.2em}\pitchfork g_{1},\ b_{1},\ c_{1})^{2} = 0.
\]

\newpage
\setcounter{page}{319}

Proceeding in a similar manner with the second set of equations, we have first
\[
(h_{3} = G_{4}B_{5} - C_{4}B_{5} = -C_{4}),
\]
(observe $h_{3} = G_{4}B_{5} - C_{4}B_{5} = -C_{4}$), and then
\[
(B_{5}B_{6},\ C_{5}G_{6},\ G_{5}G_{6},\ C_{5}G_{6}+C_{6}G_{5},\ G_{5}B_{6}+G_{6}B_{5} + A_{4}F_{5},\ B_{5}C_{6}+B_{6}C_{5}\hspace{-0.2em}\pitchfork g_{3},\ h_{3},\ b_{3})^{2} = 0.
\]

The third set gives more simply
\[
g_{5}^{2}B_{5}B_{6} + g_{5}A_{2}(-B_{5}G_{6}+B_{6}G_{5}+A_{2}F_{5}) + b_{5}G_{5}G_{6} = 0,
\]
or since
\[
g_{5} : b_{5} = G_{6} : -B_{6},
\]
this is
\[
G_{6}^{2}B_{5}B_{6} - G_{6}B_{5}(B_{5}G_{6}+B_{6}G_{5}+A_{2}F_{5}) + B_{5}G_{5}G_{6} = 0,
\]
and similarly, the fourth set gives
\[
g_{5}^{2}B_{5}B_{6} + g_{5}A_{2}(B_{5}G_{6}+B_{6}G_{5}+A_{2}F_{5}) + h_{5}G_{5}G_{6} = 0,
\]
or since $g_{5} : b_{5} = G_{6} : -B_{6}$, this is
\[
G_{6}^{2}B_{5}B_{6} - G_{6}B_{5}(B_{5}G_{6}+B_{6}G_{5}+A_{2}F_{5}) + B_{5}G_{5}G_{6} = 0:
\]
and these last two results lead to the values of the ratios of $B_{4}B_{5}$, $B_{4}G_{5}+B_{5}G_{4}+A_{4}F_{5}$, $G_{4}G_{5}$; viz. these are proportional to expressions containing the common factor $B_{5}G_{6} - B_{6}G_{5}$, and omitting this common factor, and taking them equal instead of merely proportional to the resulting expressions (which is allowable, since the absolute values are not material), we have
\[
B_{4}B_{5},\ B_{4}G_{5}+B_{5}G_{4}+A_{4}F_{5},\ G_{4}G_{5} = B_{4}B_{5},\ B_{4}G_{5}+B_{5}G_{4},\ G_{4}G_{5}.
\]

Returning to the result obtained from the first set of equations, this now becomes
\[
(B_{5}B_{6},\ G_{5}G_{6},\ H_{5}H_{6},\ G_{5}H_{6}+G_{6}H_{5},\ H_{5}B_{6}+H_{6}B_{5},\ B_{5}G_{6}+B_{6}G_{5}\hspace{-0.2em}\pitchfork g_{1},\ h_{1},\ c_{1})^{2} = 0:
\]
but the terms containing $g_{1}$, $b_{1}$ are $(B_{5}g_{1}+G_{5}b_{1})(B_{6}g_{1}+G_{6}b_{1})$, viz. this is $= -H_{5}c_{1} \cdot -H_{6}c_{1}$, that is $H_{5}H_{6}c_{1}^{2}$; the whole equation is thus divisible by $c_{1}$, and omitting this factor, it becomes
\[
g_{1}(H_{5}B_{6}+H_{6}B_{5}) + b_{1}(G_{5}H_{6}+G_{6}H_{5}) + c_{1}(H_{5}H_{6}+H_{6}H_{5}) = 0.
\]
Proceeding in like manner with the result obtained from the second set of equations, this becomes
\[
(B_{5}B_{6},\ C_{5}G_{6},\ G_{5}G_{6},\ C_{5}G_{6}+C_{6}G_{5},\ B_{5}G_{6}+B_{6}G_{5},\ B_{5}C_{6}+B_{6}C_{5}\hspace{-0.2em}\pitchfork g_{3},\ h_{3},\ b_{3})^{2} = 0,
\]
where the terms containing $g_{3}$, $b_{3}$ are $(B_{5}g_{3}+G_{5}b_{3})(B_{6}g_{3}+G_{6}b_{3})$, viz. this is $-h_{3}G_{6} \cdot -h_{3}G_{5} = h_{3}^{2}G_{5}G_{6}$; the whole equation divides by $h_{3}$, and it then becomes
\[
g_{3}(B_{5}G_{6}+B_{6}G_{5}) + h_{3}(G_{5}G_{6}+G_{6}G_{5}) + b_{3}(G_{5}G_{6}+G_{6}G_{5}) = 0.
\]

Considering $B_{4}$, $G_{4}$, $F_{4}$ as given by the equations
\[
B_{4}B_{5} = B_{4}B_{5},\quad G_{4}G_{5} = G_{4}G_{5},\quad B_{4}G_{5}+G_{4}B_{5}+A_{4}F_{4} = B_{4}G_{5}+B_{5}G_{4},
\]
the equations last obtained determine the values of $H_{4}$ and $C_{4}$, viz. these equations may be written
\[
(g_{1}B_{5}+h_{1}C_{5}+b_{1}G_{5})C_{4} + C_{5}(g_{1}B_{4}+b_{1}G_{4}) + h_{1}G_{5}C_{4} = 0;
\]
but in order that the values $(B_{4}, C_{4}, F_{4}, b_{4}, H_{4})$ given by these five equations may belong to a line $(0, B_{4}, C_{4}, F_{4}, G_{4}, H_{4})$, they must satisfy the equation
\[
B_{4}G_{4} + C_{4}H_{4} = 0,
\]
viz. in order to the existence of the line $4'$, this equation must be satisfied identically by the foregoing values; and I proceed to show that it is in fact thus satisfied.

Multiplying the values of $C_{4}$, $H_{4}$, and writing $C_{1}H_{2} = -B_{1}G_{2}$, the identity to be verified is
\begin{multline*}
(g_{1}\!\cdot\!-B_{5} + b_{1}G_{5} + c_{1}H_{5})(g_{1}\!\cdot\!B_{6} + h_{1}C_{6} + b_{5}G_{6}\cdot\!-B_{4}\!\cdot\!G_{4})\\
+ \{H_{5}(g_{1}B_{4}+b_{1}G_{4}) + (h_{1}H_{5}H_{4}\}\{C_{5}(g_{1}B_{4}+b_{1}G_{4}) + h_{1}C_{5}C_{4}\} = 0.
\end{multline*}
The first line includes the terms
\[
\{g_{1}g_{3}B_{5}^{2} + b_{1}b_{3}G_{5}^{2} + (b_{1}g_{3}+b_{3}g_{1})B_{5}G_{5} + c_{1}C_{5}H_{5}\}B_{4}G_{4},
\]
which, writing $C_{1}H_{2} = -B_{1}G_{2}$ and $B_{2}B_{3} = B_{4}B_{5}$, $G_{2}G_{3} = G_{4}G_{5}$, are
\[
= g_{1}g_{3}B_{4}G_{4}B_{5}B_{6} + b_{1}b_{3}G_{4}B_{4}G_{5}G_{6} + (b_{1}g_{3}+g_{1}b_{3}-c_{1}h_{3})B_{4}B_{5}G_{4}G_{6}.
\]
The second line includes the terms
\[
C_{4}H_{4}(g_{1}B_{5}+b_{1}G_{5})(g_{3}B_{6}+h_{3}C_{6}) + c_{1}h_{3}C_{4}H_{4}C_{5}H_{5},
\]
which, reducing in like manner, are
\[
= -g_{1}g_{3}G_{4}B_{4}B_{5}B_{6} - b_{1}b_{3}B_{4}G_{4}G_{5}G_{6} - (b_{1}g_{3}+g_{1}b_{3}-c_{1}h_{3})B_{4}B_{5}G_{4}G_{6},
\]
and these are together
\[
= (g_{1}g_{3}B_{4}B_{5} - b_{1}b_{3}G_{4}G_{6})(B_{4}G_{5} - B_{5}G_{4}).
\]
The remaining terms from the first line are at once reduced to
\[
(g_{1}h_{3}C_{5}C_{6} + c_{1}h_{3}H_{5}G_{6})B_{4}B_{5} + (b_{1}h_{3}G_{5}B_{6} + c_{1}h_{3}H_{5}B_{4})G_{4}G_{6},
\]
and those from the second line are
\[
G_{4}C_{4}\, H_{5}(g_{1}B_{5}+b_{1}G_{5}) + H_{5}H_{4}\, C_{5}(c_{1}g_{1}B_{5}+c_{1}b_{1}G_{5}).
\]

Hence, attending to the relation $C_{3} = -h_{3}$, and collecting and arranging, the equation to be verified is
\begin{multline*}
(g_{1}g_{3}B_{4}B_{5} - b_{1}b_{3}G_{4}G_{6})(B_{4}G_{5} - B_{5}G_{4})\\
+ h_{3}G_{4}G_{6}(g_{3}B_{5}B_{6} - b_{3}H_{5}H_{6})\\
+ h_{3}H_{5}H_{6}(g_{1}C_{5}C_{6} - b_{1}G_{5}G_{6})\\
+ h_{3}C_{5}B_{5}(b_{3}G_{5}G_{6} - g_{3}H_{5}H_{6})\\
+ h_{3}H_{5}G_{5}(b_{1}\!\cdot\!G_{4}G_{5} - g_{1}\!\cdot\!B_{4}B_{5}) = 0.
\end{multline*}

\newpage
\setcounter{page}{321}

But we have
\begin{align*}
g_{1}B_{5}B_{6} - h_{3}H_{5}H_{6} &= B_{5}B_{6}(G_{5}H_{6} - G_{6}H_{5}) - H_{5}H_{6}(B_{5}G_{6} - B_{6}G_{5})\\
&= B_{5}H_{5}(B_{6}G_{6} + C_{6}H_{6}) - B_{6}H_{6}(B_{5}G_{5} + C_{5}H_{5}) = 0,
\end{align*}
and similarly
\begin{align*}
g_{3}G_{5}G_{6} - h_{3}G_{5}G_{6} &= 0,\\
h_{3}G_{5}G_{6} - g_{3}H_{5}H_{6} &= 0,\\
b_{3}C_{5}C_{6} - g_{3}B_{4}B_{5} &= 0.
\end{align*}
Moreover, writing $g_{1}B_{4}B_{5} = h_{3}H_{5}H_{6}$, we have
\begin{align*}
&g_{1}g_{3}B_{4}B_{5} - b_{1}b_{3}G_{4}G_{6}\\
&= b_{1}(g_{3}H_{5}H_{6} - b_{3}G_{5}G_{6}) = 0.
\end{align*}
and the five terms of the equation in question thus separately vanish; and the equation is consequently verified.

We may collect the results as follows:

Data are lines $1'$, 2, $3'$, 4, $2'$, $5'$, $6'$;

and then, for the remaining lines, 1, 3, 5, 6, $4'$ the coordinates are as follows:

\textbf{For} $\mathbf{4'}$,
\begin{align*}
B_{4}B_{5} = B_{4}B_{5},\quad G_{4}G_{5} = G_{4}G_{5},\quad A_{4}F_{4} = B_{4}G_{5}+B_{5}G_{4}-B_{4}G_{5}-B_{5}G_{4},\quad A_{4} &= 0,\\
\begin{vmatrix} B_{5}, & G_{5}, & H_{5}\\ B_{6}, & G_{6}, & H_{6}\end{vmatrix} H_{4}B_{5},\quad H_{4}G_{5},\quad H_{4}H_{5} &= 0,\\
\begin{vmatrix} B_{5}, & G_{5}, & H_{5}\\ B_{6}, & G_{6}, & H_{6}\end{vmatrix} B_{4},\quad G_{4},\quad H_{4}\\
\begin{vmatrix} B_{4}, & C_{4}, & G_{4}\\ c_{4}+C_{5}B_{4}, & G_{4}c_{4}, & C_{5}G_{4}\end{vmatrix} &= 0,\\
\begin{vmatrix} B_{4}, & G_{4}, & G_{4}\\ B_{5}, & G_{5}, & G_{5}\end{vmatrix}\\
\begin{vmatrix} B_{5}, & C_{5}, & G_{5}\\ B_{6}, & C_{6}, & G_{6}\end{vmatrix} \quad (B_{4}G_{4}+C_{4}H_{4} &= 0,\ \text{identity}).
\end{align*}

\textbf{For} $\mathbf{1}$,
\[
(g_{1}, b_{1}, c_{1}) = \begin{vmatrix} B_{5}, & G_{5}, & H_{5}\\ B_{6}, & G_{6}, & H_{6}\end{vmatrix},\quad f_{1} = \frac{1}{A_{2}}\!\cdot\!\begin{vmatrix} B_{5}, & G_{5}, & H_{5}\\ B_{6}, & G_{6}, & H_{6}\end{vmatrix},\quad F_{1}a_{1} = \begin{vmatrix} B_{4}, & G_{4}, & H_{4}\\ B_{5}, & G_{5}, & H_{5}\end{vmatrix},\quad g_{1} = 0,
\]
\[
\begin{vmatrix} B_{5}, & G_{5}, & H_{5}\\ B_{6}, & G_{6}, & H_{6}\end{vmatrix}\\
\begin{vmatrix} B_{4}, & G_{4}, & H_{4}\\ B_{5}, & G_{5}, & H_{5}\end{vmatrix}
\]
$(a_{1}f_{1} + b_{1}g_{1} = 0,\ \text{identity})$.

\textbf{For} $\mathbf{3}$,
\[
(g_{3},\ h_{3},\ b_{3}) = \begin{vmatrix} A_{3}, & G_{5}, & G_{4}\\ B_{5}, & C_{5}, & G_{5}\end{vmatrix},\quad A_{3}f_{3} = -\begin{vmatrix} A_{3}, & C_{5}, & G_{4}\\ B_{5}, & C_{5}, & G_{5}\end{vmatrix},\quad F_{3}a_{3} = \begin{vmatrix} B_{4}, & G_{4}, & G_{4}\\ B_{5}, & G_{5}, & G_{5}\end{vmatrix},\quad C_{3} = 0,
\]
\[
\begin{vmatrix} B_{4}, & G_{4}, & G_{4}\\ B_{5}, & G_{5}, & G_{5}\end{vmatrix}
\]
$(a_{3}f_{3} + b_{3}g_{3} = 0,\ \text{identity})$.

\newpage
\setcounter{page}{322}

\textbf{For} $\mathbf{5}$,
\[
g_{5} = G_{6},\quad b_{5} = -B_{6},\quad A_{2}f_{5} = B_{5}G_{6} - B_{6}G_{5},\quad F_{5}a_{5} = B_{5}G_{6} - B_{6}G_{5},\quad c_{5} = 0,\quad h_{5} = 0,
\]
$(a_{5}f_{5} + b_{5}g_{5} = 0,\ \text{identity})$.

\textbf{For} $\mathbf{6}$,
\[
g_{6} = G_{5},\quad b_{6} = -B_{5},\quad A_{2}f_{6} = B_{5}G_{6} - B_{6}G_{5},\quad F_{6}a_{6} = B_{5}G_{6} - B_{6}G_{5},\quad c_{6} = 0,\quad h_{6} = 0;
\]
and, for actual calculation, it is convenient to remark that as only the ratios are material, a set of six coordinates may be multiplied or divided by any common number at pleasure.

But these results may be further reduced. Writing
\[
(g_{1},\ h_{1},\ c_{1}) = \begin{vmatrix} -a_{3}, & G_{4}, & G_{5}\\ B_{5}, & C_{5}, & G_{5}\end{vmatrix},
\]
we have
\[
-(g_{1}B_{5}+h_{1}C_{5}+b_{1}G_{5}) = \begin{vmatrix} C_{5}, & G_{5}\\ c_{4}, & G_{4}\end{vmatrix},
\]
\[
(g_{1}B_{5}+h_{1}C_{5}+b_{1}G_{5}) + h_{1}C_{5}C_{4} = \frac{1}{G_{5}C_{4}}\{g_{1}B_{5}B_{4}G_{4}+h_{1}G_{5}G_{4}B_{4}\} + h_{1}C_{5}C_{4}.
\]
But
\begin{align*}
&g_{1}B_{5}B_{4}G_{4} + h_{1}G_{5}G_{4}B_{4}\\
&= \frac{B_{4}G_{4}}{G_{5}C_{4}}\{g_{1}B_{5}G_{5}C_{4}(G_{4}H_{5}-G_{5}H_{4}) + B_{5}G_{5}H_{4}(G_{5}B_{4}-G_{4}B_{5})\}\\
&= C_{5}G_{4}\{B_{4}(G_{5}H_{6}-G_{6}H_{5}) + G_{4}(B_{5}H_{6}-B_{6}H_{5})\}\\
&= C_{5}G_{4}(B_{4}g_{1}+G_{4}h_{1}),
\end{align*}
since
\[
(g_{1},\ h_{1},\ c_{1}) = \begin{vmatrix} B_{5}, & G_{5}, & H_{5}\\ B_{6}, & G_{6}, & H_{6}\end{vmatrix},
\]
the equation obtained is thus
\[
g_{1}B_{4} + h_{1}G_{4} = \frac{C_{5}G_{4}}{G_{5}C_{4}}(g_{1}B_{5}+b_{1}G_{5}).
\]

We then have
\[
-(g_{1}B_{5}+h_{1}C_{5}+b_{1}G_{5})\, G_{4} = \frac{C_{5}G_{4}}{G_{5}C_{4}}[C_{5}(g_{1}B_{5}+b_{1}G_{5}) + h_{1}B_{5}G_{5}]
\]
\[
= \frac{C_{5}G_{4}}{G_{5}C_{4}}[C_{5}(g_{1}B_{5}+b_{1}G_{5}) + c_{1}C_{5}H_{5}]
\]
\[
= \frac{C_{5}G_{4}}{G_{5}C_{4}}(g_{1}B_{5}+h_{1}C_{5}+c_{1}H_{5}),
\]
that is
\[
\frac{C_{4}}{H_{4}} = \frac{C_{5}G_{4}}{G_{5}C_{4}}\, \frac{g_{1}B_{5}+h_{1}C_{5}+c_{1}H_{5}}{g_{1}B_{5}+b_{1}G_{5}+c_{1}H_{5}}.
\]

\newpage
\setcounter{page}{323}

and in like manner
\begin{align*}
-(g_{1}B_{5}+b_{1}G_{5}+c_{1}H_{5})\, H_{4} &= \frac{C_{5}G_{4}}{G_{5}C_{4}}\{g_{1}G_{5}B_{5}B_{6} + h_{1}B_{5}G_{5}G_{6}\} + c_{1}H_{5}H_{6}\\
&= \frac{G_{5}B_{5}B_{6}(G_{5}H_{6}-G_{6}H_{5})}{G_{5}C_{4}} + \frac{B_{5}G_{5}G_{6}(B_{5}H_{6}-B_{6}H_{5})}{G_{5}C_{4}}\\
&\quad + B_{5}G_{5}G_{6}\{H_{5}B_{5} - H_{5}G_{5}\} + B_{5}G_{5}H_{5}(B_{5}G_{6}-B_{6}G_{5})\\
&= H_{5}H_{6}[B_{5}(C_{5}C_{6}-C_{6}G_{5}) + G_{5}(B_{5}C_{6}-B_{6}G_{5})]\\
&= H_{5}H_{6}(B_{5}g_{1}+G_{5}h_{1});
\end{align*}
the equation obtained is thus
\[
g_{1}B_{5}+h_{1}G_{5} = \frac{C_{4}H_{4}}{C_{5}G_{4}}(B_{5}g_{1}+G_{5}h_{1})
\]
and then
\begin{align*}
-(g_{1}B_{5}+b_{1}G_{5}+c_{1}H_{5})\, H_{4} &= \frac{C_{4}H_{4}}{C_{5}G_{4}}\{H_{5}(B_{5}g_{1}+G_{5}h_{1}) + c_{1}B_{5}G_{5}\}\\
&= \frac{C_{4}H_{4}}{C_{5}G_{4}}[H_{5}(B_{5}g_{1}+G_{5}h_{1}) + h_{1}C_{5}H_{5}]\\
&= -\frac{C_{4}H_{4}}{C_{5}G_{4}}(B_{5}g_{1}+C_{5}h_{1}+G_{5}h_{1}),
\end{align*}
that is
\[
\frac{H_{4}}{F_{4}} = \frac{H_{5}H_{6}}{C_{5}}\, \frac{g_{3}B_{5}+h_{3}C_{5}+b_{3}G_{5}}{g_{1}B_{5}+b_{1}C_{5}+c_{1}H_{5}},
\]
which values of $C_{4}$, $H_{4}$ satisfy, as they should do, the relation
\[
B_{4}G_{4} + C_{4}H_{4} = 0.
\]

We have also
\begin{align*}
A_{4}F_{4} &= B_{4}G_{5} + B_{5}G_{4} - \frac{G_{5}B_{5}B_{6}}{G_{4}} - \frac{G_{4}B_{4}C_{5}}{B_{5}}\\
&= \frac{1}{B_{5}G_{4}}(G_{4}B_{5}-B_{4}G_{5})(G_{4}B_{5}-B_{5}G_{4}),
\end{align*}
which gives $F_{4}$.

Moreover
\begin{align*}
-F_{4}a_{4} &= g_{1}B_{5}+h_{1}G_{5}+c_{1}H_{5} = \frac{C_{4}H_{4}}{C_{5}G_{4}}(g_{1}B_{5}+h_{1}G_{5})\\
&= \frac{C_{4}H_{4}}{C_{5}G_{4}}(g_{1}B_{5}+h_{1}G_{5}) + c_{1}H_{5}\cdot\frac{C_{4}H_{4}}{C_{5}G_{4}}\\
&= \frac{H_{5}H_{6}}{C_{5}}\frac{(g_{1}B_{5}+b_{1}G_{5})(g_{1}B_{5}+h_{1}G_{5}+c_{1}H_{5})}{(g_{1}B_{5}+h_{1}G_{5}+b_{1}G_{5})h_{1}H_{5}}\\
&= \frac{H_{5}H_{6}}{C_{5}H_{4}}\frac{(g_{1}B_{5}+b_{1}G_{5}+c_{1}H_{5})\{(g_{1}B_{5}+h_{1}G_{5})(g_{1}B_{5}+b_{1}G_{5})+c_{1}h_{1}B_{5}G_{5}\}}{(g_{1}B_{5}+h_{1}G_{5}+b_{1}G_{5})}
\end{align*}
which combined with the foregoing value of $F_{4}$, gives $a_{4}$.

\newpage
\setcounter{page}{324}

Again,
\begin{align*}
-F_{3}a_{3} &= g_{3}B_{5}+h_{3}C_{5}+b_{3}G_{5} = \frac{C_{4}H_{4}}{C_{5}G_{4}}(g_{3}B_{5}+b_{3}G_{5}) = \frac{C_{4}H_{4}}{C_{5}G_{4}}\frac{g_{3}B_{5}+h_{3}C_{5}+b_{3}G_{5}}{c_{1}H_{4}}\\
&= \frac{C_{4}H_{4}}{C_{5}G_{4}}\frac{(g_{3}B_{5}+b_{3}G_{5})(g_{3}B_{5}+h_{3}G_{5}) + c_{1}h_{3}B_{5}C_{5}}{C_{5}H_{4}(g_{3}B_{5}+h_{3}C_{5}+b_{3}G_{5})},
\end{align*}
which combined with the foregoing value of $F_{3}$, gives $a_{3}$.

Write
\[
\omega = B_{5}G_{5} + B_{6}G_{6} + C_{5}H_{5} + G_{5}H_{6},
\]
we have
\begin{align*}
b_{1}g_{1} &= (H_{5}B_{6} - H_{6}B_{5})(G_{5}H_{6} - G_{6}H_{5})\\
&= -H_{5}B_{5}G_{6} - H_{6}B_{6}G_{5} + H_{5}H_{6}(B_{5}G_{6} + B_{6}G_{5})\\
&= H_{5}H_{6}(C_{5}H_{5} + C_{6}H_{6} + B_{5}G_{6} + B_{6}G_{5}),
\end{align*}
that is
\[
b_{1}g_{1} = H_{5}H_{6}\, \omega,
\]
and similarly
\begin{align*}
b_{3}g_{3} &= 0,\ a_{3}\, \omega,\\
b_{1}b_{3} &= B_{4}B_{5}\, \omega,\\
g_{1}g_{3} &= G_{5}G_{6}\, \omega,\\
h_{1}g_{3} + b_{3}g_{1} + C_{1}h_{3} &= -(B_{5}G_{6}+B_{6}G_{5})\, \omega,\\
(b_{1}G_{5}+g_{1}B_{5})(h_{3}G_{5}+g_{3}B_{5}) + c_{1}h_{3}B_{5}G_{5} &= \{G_{5}^{2}B_{4}B_{5} + B_{5}^{2}G_{5}G_{6} - B_{4}B_{5}(B_{5}G_{6}+B_{6}G_{5})\}\, \omega\\
&= (G_{5}B_{5} - B_{5}G_{6})(G_{5}B_{4} - B_{5}G_{5})\, \omega,
\end{align*}
which last value is to be substituted for the left-hand function in the formulae for $a_{1}$ and $a_{3}$ respectively.

Whence, finally recollecting that
\[
B_{4}G_{4} + G_{5}H_{5} = 0,\quad B_{4}G_{4} + G_{6}H_{6} = 0,\quad B_{5}G_{5} + C_{5}H_{5} = 0,
\]
and
\[
\omega = C_{5}H_{5} + G_{5}H_{6} + B_{5}G_{5} + B_{6}G_{6},
\]
we have
\[
\begin{vmatrix} B_{5}, & G_{5}, & H_{5}\\ B_{6}, & G_{6}, & H_{6}\end{vmatrix} = -\, (g_{1}B_{5}+b_{1}G_{5}+c_{1}H_{5}).
\]

\textbf{For} $\mathbf{1}$
\[
b_{1}g_{1} = H_{5}H_{6}\, \omega,\qquad a_{1} = \frac{A_{2}H_{5}H_{6}\, C_{5}}{g_{1}B_{5}+h_{1}G_{5}+c_{1}H_{5}},\qquad h_{1} = 0.
\]

\newpage
\setcounter{page}{325}

\textbf{For} $\mathbf{3}$
\[
(g_{3},\ h_{3},\ b_{3}) = \begin{vmatrix} B_{4}, & C_{4}, & G_{4}\\ B_{5}, & C_{5}, & G_{5}\end{vmatrix},\qquad a_{3} = \frac{A_{2}B_{5}G_{5}\, C_{5}}{g_{5}B_{5}+h_{5}G_{5}+b_{3}G_{5}},\qquad h_{3} = 0,
\]
where observe that $C_{3} + h_{3} = 0$.

\textbf{For} $\mathbf{5}$
\[
h_{5} = G_{6},\quad b_{5} = -B_{6},\quad A_{2}f_{5} = \tfrac{1}{2}(G_{5}B_{6}-B_{5}G_{6}),\qquad a_{5} = \frac{A_{2}B_{5}G_{6}}{G_{5}B_{6}-B_{5}G_{5}},
\]
\[
c_{5} = 0,\quad h_{5} = 0.
\]

\textbf{For} $\mathbf{6}$
\[
g_{6} = G_{5},\quad b_{6} = -B_{5},\quad f_{6} = \tfrac{1}{2}(G_{5}B_{6}-B_{5}G_{6}),\qquad a_{6} = \frac{A_{2}B_{5}G_{6}}{G_{5}B_{6}-B_{5}G_{5}},
\]
\[
c_{6} = 0,\quad h_{6} = 0.
\]

\textbf{For} $\mathbf{4'}$
\[
B_{4} = C_{5}G_{4},\quad g_{4} = \frac{g_{1}B_{5}+b_{1}G_{5}+c_{1}H_{5}}{H_{4}\, g_{1}B_{5}+h_{1}G_{5}+b_{1}G_{5}},
\]
\[
H_{4} = \frac{H_{5}H_{6}}{C_{5}}\, \frac{g_{3}B_{5}+h_{3}G_{5}+b_{3}G_{5}}{g_{1}B_{5}+b_{1}G_{5}+c_{1}H_{5}}.
\]

I have thought it worth while to effect the numerical calculations for enabling the construction of a drawing or model. For this purpose taking $X$, $Y$, $Z$ as ordinary rectangular coordinates, I write
\begin{align*}
x &= X + Y + Z - 10,\\
y &= X - Y + Z - 10,\\
z &= -X + Y + Z - 10,\\
w &= 7,
\end{align*}
that is, I take $1'$ and 2 to be lines in the plane of $XY$, defined by the equations $X+Y=10$ and $X-Y=-10$ respectively, and $3'$ and 4 to be lines in the plane of $XZ$ defined by the equations $X-Z=-10$, $X+Z=10$ respectively. And I take $5'$ to be the line joining the points $(2,0,8)$ and $(-9,1,0)$; $6'$ the line joining the points $(3,0,7)$ and $(-8,2,0)$; $2'$ the line joining the points $(9,0,1)$ and $(-3,7,4)$. We calculate then for each of the lines the $xyzw$ coordinates of two points thereof; and thence the six coordinates of the line, viz.

\newpage
\setcounter{page}{326}

\begin{center}
\begin{tabular}{c|cccc|cccc|cccccc}
& $X$ & $Y$ & $Z$ & $w$ & $X$ & $Y$ & $Z$ & $w$ & $A$ & $B$ & $C$ & $F$ & $G$ & $H$\\\hline
$5'$ & 0, & 8, & -4, & & -18, & 0, & 0, & & 0, & 72, & 144, & 0, & 8, & -4\\
& & & & & & & & & 0, & 18, & 36, & 0, & 2, & -1\\
$6'$ & 0, & 7, & -6, & & -16, & 0, & 0, & 2 & 0, & 96, & 112, & 0, & 14, & -12\\
& & & & & & & & & 0, & 48, & 56, & 0, & 7, & -6\\
$2'$ & 0, & 1, & -18, & & -2, & 4, & 4, & 7 & & 76, & 36, & 2, & 7, & -126\\
& & & & & & & & & 76, & 36, & 2, & 0, & 7, & -126\\
\end{tabular}
\end{center}
and effecting the calculations for the remaining lines, we have
\begin{center}
\begin{tabular}{c|cccccc}
& $A$ & $B$ & $C$ & $F$ & $G$ & $H$\\\hline
$4'$ & & & & $\frac{9}{2}$ & & $\frac{3}{38}$\\
& & & & & $\frac{24}{-944}$ & $\frac{59}{385}$\\
1 & & & $\frac{380}{59}$ & & $\frac{30}{-5}$ & $\frac{19}{15}$\\
3 & 127680 & -720 & & & $\frac{140}{19}$ & -30\\
5 & & $\frac{304}{3}$ & -48 & & $\frac{7}{19}$ & \\
6 & 152 & -18 & & & $\frac{27}{38}$ & \\
\end{tabular}
\end{center}
or reducing to integers, the values are
\begin{center}
\begin{tabular}{c|ccccc}
& $B$ & $C$ & $F$ & $G$ & $H$\\\hline
$5'$ & 18 & 36 & & 2 & -1\\
$6'$ & 48 & 56 & & 7 & -6\\
$2'$ & 76 & 36 & 2 & 7 & -126\\
$4'$ & & 53808 & -2116448 & -531 & 4484 \\
1 & 1444 & 13452 & 6726 & 10443 & -1121\\
3 & 485184 & -2736 & & 3 & 532 \\
5 & 5776 & -912 & & 21 & 133\\
6 & 5776 & -2052 & & 81 & 228
\end{tabular}
\end{center}

\newpage
\setcounter{page}{327}

The line $(a, b, c, f, g, h)$ is given as the intersection of any two of the four planes
\begin{align*}
(0,\ h,\ -g,\ a\hspace{-0.2em}\pitchfork x, y, z, w) &= 0,\\
(-h,\ 0,\ f,\ b\hspace{-0.2em}\pitchfork x, y, z, w) &= 0,\\
(g,\ -f,\ 0,\ c\hspace{-0.2em}\pitchfork x, y, z, w) &= 0,\\
(-a,\ -b,\ -c,\ 0\hspace{-0.2em}\pitchfork x, y, z, w) &= 0.
\end{align*}

Or substituting for $x$, $y$, $z$, $w$ the values $X+Y+Z-10$, $X-Y+Z-10$, $-X+Y+Z-10$, 7, these become
\begin{align*}
(-f-h,\ b+f-h,\ f-h,\ h-g\hspace{-0.2em}\pitchfork X, Y, Z, 10) &= 0,\\
(g,\ c+g,\ g-f,\ -g\hspace{-0.2em}\pitchfork X, Y, Z, 10) &= 0,\\
(c-a,\ -c-a,\ -c-a-b,\ c+a\hspace{-0.2em}\pitchfork X, Y, Z, 10) &= 0,
\end{align*}
or, what is the same thing,
\begin{multline*}
(-2g+a-c,\ 2g-f-h,\ a+b+c+f-h,\ -2g,\ -a-c\hspace{-0.2em}\pitchfork X, Y, Z, 10) = 0,\\
(-2g+f+h,\ -a-b-c-f+h,\ 2g-f-h,\ -f+h\hspace{-0.2em}\pitchfork X, Y, Z, 10) = 0.
\end{multline*}

And substituting, we have the equations of the several lines, viz.:

\begin{align*}
(1')\qquad &Z+Y=10,\quad Z=0;\\
(2)\qquad &-Z+Y=10,\quad Z=0;\\
(3')\qquad &-Z+X=10,\quad Y=0;\\
(4)\qquad &X+Z=10,\quad Y=0;\\
(5')\qquad &(\cdot\qquad 40,\quad -4\hspace{-0.2em}\pitchfork X, Y, Z, 10) = 0,\\
&\qquad(-40,\quad \cdot\qquad 55,\quad -36\hspace{-0.2em}\pitchfork X, Y, Z, 10) = 0,\\
&\qquad(-5,\quad -55,\quad \cdot\qquad 4,\quad 36,\quad -1\hspace{-0.2em}\pitchfork X, Y, Z, 10) = 0;\\
(6')\qquad &(\cdot\qquad -35,\quad 10,\quad -7\hspace{-0.2em}\pitchfork X, Y, Z, 10) = 0,\\
&\qquad(-35,\quad \cdot\qquad 1,\quad 55,\quad -28\hspace{-0.2em}\pitchfork X, Y, Z, 10) = 0,\\
&\qquad(-10,\quad -55,\quad 7,\quad 28,\quad -3\hspace{-0.2em}\pitchfork X, Y, Z, 10) = 0.
\end{align*}

\newpage
\setcounter{page}{328}

\begin{align*}
(2')\qquad &(\cdot\qquad -30,\quad 70,\quad -7\hspace{-0.2em}\pitchfork X, Y, Z, 10) = 0,\\
&\qquad(30,\quad \cdot\qquad 120,\quad -39\hspace{-0.2em}\pitchfork X, Y, Z, 10) = 0,\\
&\qquad(-70,\quad -120,\quad \cdot\qquad 63,\quad 7,\quad 39,\quad -63\hspace{-0.2em}\pitchfork X, Y, Z, 10) = 0;\\
(4')\qquad &(\cdot\qquad -2107480,\quad 9385,\quad -8968\hspace{-0.2em}\pitchfork X, Y, Z, 10) = 0,\\
&\qquad(2107480,\quad \cdot\qquad -2063,\quad 2116448\hspace{-0.2em}\pitchfork X, Y, Z, 10) = 0,\\
&\qquad(-9385,\quad 2063285,\quad \cdot\qquad -645,\\
&\qquad\qquad 8968,\quad -2116448,\quad 645,\hspace{-0.2em}\pitchfork X, Y, Z, 10) = 0;\\
(1)\qquad &(\cdot\qquad 3040,\quad -12685,\quad 2242\hspace{-0.2em}\pitchfork X, Y, Z, 10) = 0,\\
&\qquad(-3040,\quad \cdot\qquad 32065,\quad -8170\hspace{-0.2em}\pitchfork X, Y, Z, 10) = 0,\\
&\qquad(12685,\quad -32065,\quad \cdot\qquad 10443,\\
&\qquad\qquad -2242,\quad 8170,\quad -10443,\hspace{-0.2em}\pitchfork X, Y, Z, 10) = 0;\\
(3)\qquad &(\cdot\qquad -484120,\quad 1175,\quad -1064\hspace{-0.2em}\pitchfork X, Y, Z, 10) = 0,\\
&\qquad(484120,\quad \cdot\qquad 482565,\quad -485184\hspace{-0.2em}\pitchfork X, Y, Z, 10) = 0,\\
&\qquad(-1175,\quad -482565,\quad \cdot\qquad 117,\\
&\qquad\qquad 1064,\quad 485184,\quad -117,\hspace{-0.2em}\pitchfork X, Y, Z, 10) = 0;\\
(5)\qquad &(\cdot\qquad 5510,\quad 245,\quad -266\hspace{-0.2em}\pitchfork X, Y, Z, 10) = 0,\\
&\qquad(-5510,\quad \cdot\qquad 4885,\quad -5776\hspace{-0.2em}\pitchfork X, Y, Z, 10) = 0,\\
&\qquad(-245,\quad -4885,\quad \cdot\qquad 3,\\
&\qquad\qquad 266,\quad 5776,\quad -21,\hspace{-0.2em}\pitchfork X, Y, Z, 10) = 0;\\
(6)\qquad &(\cdot\qquad -5320,\quad 375,\quad -456\hspace{-0.2em}\pitchfork X, Y, Z, 10) = 0,\\
&\qquad(5320,\quad \cdot\qquad 3805,\quad -5776\hspace{-0.2em}\pitchfork X, Y, Z, 10) = 0,\\
&\qquad(-375,\quad -3805,\quad \cdot\qquad 81,\\
&\qquad\qquad 456,\quad 5776,\quad -81,\hspace{-0.2em}\pitchfork X, Y, Z, 10) = 0.
\end{align*}

\newpage
\setcounter{page}{329}

The several lines intersect as they should do, the coordinates of the points of intersection being as follows
\begin{center}
\begin{tabular}{c|ccccc}
& $1'$ & $3'$ & 4 & $5'$ & $6'$\\\hline
\multirow{2}{*}{1} & $\cdot$ & $\cdot$ & & $\left(\dfrac{305}{1064}, -\dfrac{424171}{8968}\right)$ & \\[1em]
& & & $\cdot$ & & $\left(-\dfrac{4721}{2}, \dfrac{47}{2}\right)$\\[1em]
\multirow{2}{*}{2} & $\cdot$ & $\cdot$ & & $\left(\dfrac{9693}{1064}, -\dfrac{2111073}{2484}\right)$ & \\[1em]
& & & $\cdot$ & & $\left(\dfrac{6521}{66}, \dfrac{24727}{464}\right)$\\[1em]
\multirow{2}{*}{$2'$} & $\dfrac{9}{1}$ & & $\cdot$ & & \\[1em]
& & & $\dfrac{8}{2}$ & & $\dfrac{3}{7}$\\[1em]
\multirow{2}{*}{$4'$} & & $\dfrac{47521}{42} + \dfrac{10801}{283}$ & $\cdot$ & & \\[1em]
& $\dfrac{71744}{2128}$ & & $\cdot$ & & \\[1em]
\multirow{2}{*}{3} & & $\cdot$ & $\cdot$ & $\left(\dfrac{24211}{152} + 233, -\dfrac{1891}{152} + \dfrac{727}{233}\right)$ & \\[1em]
& & & & & $\cdot$
\end{tabular}
\end{center}
viz. the coordinates of $12'$ (intersection of lines 1 and $2'$) are $\left(\frac{305}{1064}, -\frac{424171}{8968}\right)$, and so in other cases; where there is no divisor the coordinates are integers. I find however, on laying down the figure, that the lines 3 and 4, $3'$ and $4'$ come so close together, that the figure cannot be obtained with any accuracy.

\newpage
\setcounter{page}{330}

%===============================================================================
\article{460}{NOTE ON MR FROST'S PAPER ON THE DIRECTION OF LINES OF CURVATURE IN THE NEIGHBOURHOOD OF AN UMBILICUS.}
%===============================================================================

\noindent[From the \textit{Quarterly Journal of Pure and Applied Mathematics}, vol.~x. (1870), pp.~111--113.]

\vspace{0.5em}

I \textsc{remark} as follows:

1. In regard to a quadric surface, it is not, I think, correct to say that the generatrices through an umbilic are curves of curvature; notwithstanding that, as shown p.~80, the normals at every point of such generatrix lie in one plane and consequently intersect. The way in which these generatrices as quasi-curves-of-curvature present themselves is as follows:

The curves of curvature satisfy a certain differential equation, the complete integral of which gives these curves as the intersections of the given quadric surface by the series of confocal surfaces
\[
\frac{x^{2}}{a^{2}+\theta} + \frac{y^{2}}{b^{2}+\theta} + \frac{z^{2}}{c^{2}+\theta} = 1,
\]
$\theta$ being the constant of integration of the differential equation. The singular solution of the differential equation, or envelope of the curves of curvature determined as above, gives the umbilicar generatrices.

2. In regard to a surface in general, I think it must be considered, not that there pass through the umbilic three distinct curves, but that the umbilicar curve of curvature is a curve having at the umbilic a triple point, or rather a point at which there are in general three distinct directions of the curve. The umbilicar curve of curvature in fact presents itself as the curve belonging to a certain value of the constant of integration $h$; in order that the curve of curvature may pass through a given point on the surface, $h$ must satisfy a certain quadratic equation, that is for a given point of the surface there are two values of $h$, and therefore two curves of curvature: but an umbilic is a point for which (as in effect shown, p.~81, for the particular case of a quadric surface) the two values of $h$ become equal; that is, there is through the umbilic only a singular curve of curvature; but $\frac{dh}{dx}$ is determined by a cubic equation, and the umbilic is thus (as just mentioned) a point at which there are in general three distinct directions of the curve.

3. Some researches on the subject are contained in my paper ``On Differential Equations and Umbilici,'' \textit{Phil.\ Mag.}, vol.~XXVI. (1863), pp.~373--379 and 441--452, [330]. It is noticeable that in the integral equations which I have there obtained for the differential equations $cy(\phi^{2}-1)+(a-c)x\phi=0$, and the more general form $(bx+cy)(\phi^{2}-1)+2(fx+gy)\phi=0$, which belong to the neighbourhood of an umbilic, the curve through the umbilic does break up into three distinct curves; and the same is the case with the umbilic on the surface $xyz=1$ presently referred to.

4. In the paper ``M\'{e}moire sur les surfaces orthogonales,'' \textit{Liouv.}, t.~XII. (1847), pp.~241--254, M.~Serret has given two very remarkable cases of three systems of surfaces intersecting each other at right angles, and consequently in the curves of curvature of the surfaces of each system. It was only on referring to this paper, in connexion with that of Mr~Frost, that I perceived an obvious enough simplification of M.~Serret's formulae, whereby it appears that the curves of curvature on the surface $xyz=1$ are given as the intersection of this surface with the series of surfaces
\[
h = (x^{2}+\omega y^{2})^{\frac{3}{2}} + (x^{2}+\omega^{2}z^{2})^{\frac{3}{2}} + (y^{2}+\omega z^{2})^{\frac{3}{2}},
\]
where $\omega$ is an imaginary cube root of unity; the rationalised equation is of the twelfth order in $(x, y, z)$, and for the particular value $h=0$, reduces itself as is easily seen to
\[
0 = (y^{2}-z^{2})^{2}(z^{2}-x^{2})^{2}(x^{2}-y^{2})^{2}.
\]
The point $x=y=z=1$ is obviously an umbilic on the surface $xyz=1$, and the corresponding value of $h$ being $h=0$, the equation just obtained determines the umbilicar curves of curvature. Viz. combining therewith the equation $xyz=1$ of the surface, we have the three hyperbolic curves
\[
(x=y,\ xy^{3}=1),\quad (x=\omega y,\ \omega^{3}y^{4}=1),\quad (x=\omega^{2}y,\ \omega^{6}y^{4}=1).
\]

\newpage
\setcounter{page}{332}

%===============================================================================
\article{461}{ON THE GEOMETRICAL INTERPRETATION OF THE COVARIANTS OF A BINARY CUBIC.}
%===============================================================================

\noindent[From the \textit{Quarterly Journal of Pure and Applied Mathematics}, vol.~x. (1870), pp.~148--149.]

\vspace{0.5em}

\textsc{Consider} the binary cubic $U=(a,b,c,d\hspace{-0.2em}\pitchfork x,y)^{3}$, and its covariants, viz. the discriminant (invariant)
\[
\nabla = a^{2}d^{2} - 6abcd + 4ac^{3} + 4b^{3}d - 3b^{2}c^{2},
\]
the Hessian
\[
H = (ac-b^{2})x^{2} + (ad-bc)xy + (bd-c^{2})y^{2},
\]
and the cubicovariant
\begin{align*}
\Phi &= (\phantom{-}a^{2}d - 3abc + 2b^{3})x^{3}\\
&\quad - (-3abd + 6ac^{2} - 3b^{2}c)x^{2}y\\
&\quad + (\phantom{-}3acd - 6bd^{2} + 3bc^{2})xy^{2}\\
&\quad - (\phantom{-}ad^{2} - 3bcd + 2c^{3})y^{3},
\end{align*}
connected by the identical equation
\[
\Phi^{2} - \nabla U^{2} = -4H^{3}.
\]

Then if we regard $(a, b, c, d)$ as the coordinates of a point in space, but $(x, y)$ as variable parameters, the equation
\[
\nabla = 0
\]
represents a quartic torse, having for its cuspidal curve the skew cubic $ac-b^{2}=0$, $ad-bc=0$, $bd-c^{2}=0$; the equation
\[
U = 0
\]
is that of the tangent plane to the torse along the line $ax^{2}+2bxy+cy^{2}=0$, $hax^{2}+2cxy+dy^{2}=0$: this line meets the cuspidal curve in the point whose coordinates are $a:b:c:d = y^{3}:-xy^{2}:x^{2}y:-y^{3}$. The equation
\[
H = 0
\]
is that of a quadric cone having the last mentioned point for its vertex, and passing through the cuspidal curve: and the equation
\[
\Phi = 0
\]
is that of the cubic surface which is the first polar of the same point in regard to the torse.

The equation $\Phi^{2} - \nabla U^{2} = -4H^{3}$, writing therein $\nabla=0$, gives $\Phi^{2} = -4H^{3}$, a result which implies that $U=0$, $H=0$ is a certain curve repeated twice, and that $U=0$, $\Phi=0$ is the same curve repeated three times. The curve in question is at once seen to be the line of contact $\partial_{x}U=0$, $\partial_{y}U=0$; it thus appears that the tangent plane $U=0$ meets the cubic surface $\Phi=0$ in this line taken three times. This can only be the case if the equation $\Phi=0$ be expressible in the form $MU + (\partial_{x}U)^{2} = 0$, or, what is the same thing,
\[
MU + (\alpha\partial_{x}U + \beta\partial_{y}U)^{2} = 0,
\]
$\alpha$ and $\beta$ constants, $M$ a quadric function of $(a, b, c, d)$; that is, $\Phi$ must be equal to a function of the form
\[
MU + (\alpha\partial_{x}U + \beta\partial_{y}U)^{2}.
\]
Seeking for this expression of $\Phi$, and writing the symbols out at length, I find that the required identical equation is
\begin{multline*}
(\phantom{-}a^{2}d - 3abc + 2b^{3})x^{3}\\
- (-3abd + 6ac^{2} - 3b^{2}c)x^{2}y\\
+ (\phantom{-}3acd - 6bd^{2} + 3bc^{2})xy^{2}\\
- (\phantom{-}ad^{2} - 3bcd + 2c^{3})y^{3}\\
- (\alpha x - \beta y)^{2}(a,b,c,d\hspace{-0.2em}\pitchfork x,y)^{2}\\
+ 2\{\alpha(ax^{2}+2bxy+cy^{2}) + \beta(bx^{2}+2cxy+dy^{2})\}^{2} =\\
(a,b,c,d\hspace{-0.2em}\pitchfork x,y)^{2} \cdot
\begin{pmatrix}
2a^{2} & 6ab & 6b^{2} & ad+3bc\\
6ab & 12ac+6b^{2} & 3ad+15bc & 6c^{2}\\
6b^{2} & 3ad+15bc & 12bd+6c^{2} & 6cd\\
ad+3bc & 6c^{2} & 6cd & 2d^{2}
\end{pmatrix}
\hspace{-0.2em}\pitchfork (x,y)^{2}(a,\beta)^{2},
\end{multline*}
(where the $\hspace{-0.2em}\pitchfork$ indicates that the binomial coefficients are not to be inserted, viz. the function on the right hand is $\{2a^{2}x^{2} + 6abx^{2}y + 6b^{2}x^{2}y^{2} + (-ad+3bc)y^{2}\}\alpha^{2} + \&\text{c.}$). As a verification remark that for $x=\alpha$, $y=\beta$, the equation becomes simply $2U^{2} = 2U \cdot 2U$.

\newpage
\setcounter{page}{334}

%===============================================================================
\article{462}{A NINTH MEMOIR ON QUANTICS.}
%===============================================================================

\noindent[From the \textit{Philosophical Transactions of the Royal Society of London}, vol.~CLXI. (for the year 1871), pp.~17--50. Received April 7,---Read May 19, 1870.]

\vspace{0.5em}

It was shown not long ago by Professor Gordan that the number of the irreducible covariants of a binary quantic of any order is finite (see his memoir ``Beweis dass jede Covariante und Invariante einer bin\"{a}ren Form eine ganze Function mit numerischen Coefficienten einer endlichen Anzahl solcher Formen ist,'' \textit{Crelle}, t.~LXIX. (1869), Memoir dated 8 June 1868), and in particular that for a binary quintic the number of irreducible covariants (including the quintic and the invariants) is $=23$, and that for a binary sextic the number is $=26$. From the theory given in my ``Second Memoir on Quantics,'' \textit{Phil.\ Trans.}, 1856, [141], I derived the conclusion, which, as it now appears, was erroneous, that for a binary quintic the number of irreducible covariants was infinite. The theory requires, in fact, a modification, by reason that certain linear relations, which I had assumed to be independent, are really not independent, but, on the contrary, linearly connected together: the interconnexion in question does not occur in regard to the quadric, cubic, or quartic and for these cases respectively the theory is true as it stands; for the quintic the interconnexion first presents itself in regard to the degree 8 in the coefficients and order 14 in the variables, viz. the theory gives correctly the number of covariants of any degree not exceeding 7, and also those of the degree 8 and order less than 14 but for the order 14 the theory as it stands gives a non-existent irreducible covariant $(a, \ldots \hspace{-0.2em}\pitchfork x, y)^{14}$, viz. we have, according to the theory, $5 = (10 - 6) + 1$, that is, of the form in question there are 10 composite covariants connected by 6 syzygies, and therefore equivalent to $10 - 6, = 4$ asyzygetic covariants: but the number of asyzygetic covariants being $= 5$, there is left, according to the theory, 1 irreducible covariant of the form in question. The fact is that the 6 syzygies being interconnected and equivalent to 5 independent syzygies only, the composite covariants are equivalent to $10 - 5, = 5$, the full number of the asyzygetic covariants. And similarly the theory as it stands gives a non-existent irreducible covariant $(a, \ldots \hspace{-0.2em}\pitchfork x, y)^{16}$. The theory being thus in error, by reason that it omits to take account of the interconnexion of the syzygies, there is no difficulty in conceiving that the effect is the introduction of an infinite series of non-existent irreducible covariants, which, when the error is corrected, will disappear, and there will be left only a finite series of irreducible covariants.

Although I am not able to make this correction in a general manner so as to show from the theory that the number of the irreducible covariants is finite, and so to present the theory in a complete form, it nevertheless appears that the theory can be made to accord with the facts; and I reproduce the theory, as well to show that this is so as to exhibit certain new formulae which appear to me to place the theory in its true light. I remark that although I have in my Second Memoir considered the question of finding the number of irreducible covariants of a given degree $\theta$ in the coefficients but of any order whatever in the variables, the better course is to separate these according to their order in the variables, and so consider the question of finding the number of the irreducible covariants of a given degree $\theta$ in the coefficients, and of a given order $n$ in the variables. (This is, of course, what has to be done for the enumeration of the irreducible covariants of a given quantic and what is done completely for the quadric, the cubic, and the quartic, and for the quintic up to the degree 6 in my Eighth Memoir, \textit{Phil.\ Trans.} 1867, [405].) The new formulae exhibit this separation thus (Second Memoir, No.~49), writing $a$ instead of $x$, we have; for the quadric the expression
\[
\frac{1}{(1-a)(1-a^{2})},
\]
showing that we have irreducible covariants of the degrees 1 and 2 respectively, viz. the quadric itself and the discriminant: the new expression is
\[
\frac{1}{1-aa^{2}} \cdot \frac{1}{1-a^{2}},
\]
showing that the covariants in question are of the actual forms $(a, \ldots \hspace{-0.2em}\pitchfork x, y)^{2}$ and $(a, \ldots)^{2}$ respectively. Similarly for the cubic, instead of the expression No.~55,
\[
\frac{1}{(1-a)(1-a^{2})(1-a^{3})(1-a^{4})},
\]
we have
\[
\frac{1-aa^{2}}{(1-aa^{2})(1-ax)(1-ax^{-1})(1-a^{4})},
\]
exhibiting the irreducible covariants of the forms $(a, \ldots \hspace{-0.2em}\pitchfork x, y)^{3}$, $(a, \ldots \hspace{-0.2em}\pitchfork x, y)^{2}$, $(a, \ldots)^{4}(x, y)^{3}$, and $(a, \ldots)^{4}$, connected by a syzygy of the form $(a, \ldots)^{2}(x, y)^{2}$; and the like for quantics of a higher order.

In the present Ninth Memoir I give the last-mentioned formulae; I carry on the theory of the quintic, extending the Table No.~82 of the Eighth Memoir up to the degree 8, calculating all the syzygies, and thus establishing the interconnexions in virtue of which it appears that there are really no irreducible covariants of the forms $(a, \ldots)^{8}(x, y)^{14}$ and $(a, \ldots)^{8}(x, y)^{16}$. I reproduce in part Gordan's theory so far as it applies to the quintic, and I give the expressions of such of the 23 covariants as are not given in my former memoirs; these last were calculated for me by Mr W.~Barrett Davis, by the aid of a grant from the Donation Fund at the disposal of the Royal Society. [The expressions referred to are in fact printed, 143.] The paragraphs of the present memoir are numbered consecutively with those of the former memoirs on Quantics.

\newpage
\setcounter{page}{336}

\textbf{Article Nos.\ 328 to 332. Reproduction of my original Theory as to the Number of the Irreducible Covariants.}

328. I reproduce to some extent the considerations by which, in my Second Memoir on Quantics, I endeavoured to obtain the number of the irreducible covariants of a given binary quantic $(a, b, \ldots \hspace{-0.2em}\pitchfork x, y)^{n}$.

Considering in the first instance the covariants as functions of the coefficients $(a, b, c, \ldots)$, without regarding the variables $(x, y)$, and attending only to the following properties---

1$^{\circ}$. a covariant is a rational and integral homogeneous function of the coefficients;

2$^{\circ}$. if $P$, $Q$, $R$, \ldots\ are covariants, any rational and integral function $F(P, Q, R, \ldots)$, homogeneous in regard to the coefficients, is also a covariant, we say---

that the covariants $X$, $Y$, \ldots\ of the same degree in regard to the coefficients, and not connected by any identical equation $\alpha X + \beta Y \ldots = 0$ (where $\alpha$, $\beta$, \ldots\ are quantities independent of the coefficients $(a, b, c, \ldots)$), are asyzygetic covariants, and that a covariant not expressible as a rational and integral function of covariants of lower degrees is an irreducible covariant; and it is assumed that we know the number of the asyzygetic covariants of the degrees 1, 2, 3, \ldots; say, these are $A_{1}$, $A_{2}$, $A_{3}$, \ldots, or, what is the same thing, that the number of the asyzygetic covariants of the degree $\theta$, or form $(a, b, \ldots)^{\theta}$, is equal to the coefficient of $a^{\theta}$ in a given function
\[
\Phi(a) = 1 + A_{1}a + A_{2}a^{2} + \ldots + A_{\theta}a^{\theta} + \ldots,
\]
where I have purposely written $a$, as a representative of the coefficients $(a, b, c, \ldots)$, in place of the $x$ of my Second Memoir.

329. The theory was, that determining $\kappa_{1}$, $\kappa_{2}$, \ldots\ by the conditions
\begin{align*}
A_{2} &= \kappa_{1}\, \tfrac{1}{2}(\kappa_{1}+1) + \kappa_{2},\\
A_{3} &= \kappa_{1}\, \tfrac{1}{2}(\kappa_{1}+1)(\kappa_{1}+2) + \kappa_{1}\kappa_{2} + \kappa_{3},
\end{align*}
that is, throwing
\[
1 + A_{1}a + A_{2}a^{2} + A_{3}a^{3} + \ldots
\]
into the form
\[
(1-a)^{-\kappa_{1}}(1-a^{2})^{-\kappa_{2}}(1-a^{3})^{-\kappa_{3}}\ldots,
\]
the index $\kappa_{r}$ would express the number of irreducible covariants of the degree $r$ less the number of the (irreducible) linear relations, or syzygies, between the composite or non-irreducible covariants of the same degree. Thus $\kappa_{2} = \kappa_{2}$, there would be $\kappa_{2}$ covariants of the degree 2; these give rise to $\frac{1}{2}\kappa_{2}(\kappa_{2}+1)$ composite covariants of the degree 2; or, assuming that these are connected by $k_{2}$ syzygies, the number of asyzygetic composite covariants of the degree 2 would be $\frac{1}{2}\kappa_{2}(\kappa_{2}+1) - k_{2}$; and thence there would be $A_{2} - \frac{1}{2}\kappa_{2}(\kappa_{2}+1) + k_{2}$, that is, $\kappa_{2} + k_{2}$ irreducible covariants of the same degree; so that (irreducible invariants less syzygies) $(\kappa_{2} + k_{2}) - k_{2}$ is $= \kappa_{2}$.

\newpage
\setcounter{page}{337}

330. The $\kappa_{2}$ syzygies are here irreducible syzygies; for, calling $P$, $Q$, $R$, \ldots\ the covariants of the degree 1, there is no identical relations between the terms $P^{2}$, $Q^{2}$, $PQ$, \ldots: imagine for a moment that we could have $L$ such identical relations (viz. this might very well be the case if instead of the $\frac{1}{2}\kappa_{1}(\kappa_{1}+1)$ functions $P^{2}$, $Q^{2}$, $PQ$, \ldots, we were dealing with the same number of other quadric functions of these quantities), that is, relations not establishing any relation between $P^{2}$, $Q^{2}$, $PQ$, \ldots, and besides these $k_{2}$ non-identical relations as above; then the number of irreducible invariants would be $\kappa_{2} + k_{2} + L$, and the number of irreducible syzygies being as before $k_{2}$, the difference would be not $\kappa_{2}$, but $\kappa_{2} + L$. The $L$ identical relations are here relations between composite covariants, and the effect (if any such relation could subsist) would, it appears, be to increase $\kappa_{2}$; between syzygies such identical relations do actually exist, and the effect is contrariwise to diminish the $\kappa$; we may, for instance, for the degree $s$ have irreducible covariants less irreducible syzygies $= \kappa_{s} - L_{s}$.

331. Assume for a moment that, for a given value of $s$, $\kappa_{s}$ is positive; but for the term $L$, it would of course follow that there was for the degree in question a certain number of irreducible covariants and it was in this manner that I was led to infer that the number of the covariants of a quintic was infinite---viz. the transformed expression for the number of asyzygetic covariants is
\[
= \text{coeff. } a^{n} \text{ in } (1-a)^{-\kappa_{1}}(1-a^{2})^{-\kappa_{2}}(1-a^{3})^{-\kappa_{3}}(1-a^{4})^{-\kappa_{4}}\ldots,
\]
a product which does not terminate, and as to which it is also assumed that the series of negative indices does not terminate.

332. The principle is the same, but the discussion as to the number of the irreducible covariants becomes more precise, if we attend to the covariants as involving not only the coefficients $(a, b, \ldots)$ but also the variables $(x, y)$; we have then to consider the covariants of the form $(a, b, \ldots)^{\theta}(x, y)^{n}$, or, say, of the form $a^{\theta}x^{n}$ (degree $\theta$ and order $n$), and the number of the asyzygetic covariants of this form is given as the coefficient of $a^{\theta}x^{n}$ in a given function of $(a, x)$, (I write $a$ instead of the $z$ of my Second Memoir in the formulae which contain $x$ and $z$): by taking account of the composite covariants and syzygies, we successively determine, from the given number of asyzygetic covariants for each value of $\theta$ and $n$, the number of the irreducible covariants for the same values of $\theta$ and $n$. This is, in fact, done for the quintic in my Eighth Memoir up to the covariants and syzygies of the degree 6. But before resuming the discussion for the quintic, I will consider the preceding cases of the quadric, the cubic, and the quartic.

\newpage
\setcounter{page}{338}

\textbf{Article Nos.\ 333 to 336. New formulae for the number of Asyzygetic Covariants.}

333. For the quadric $(a, b, c\hspace{-0.2em}\pitchfork x, y)^{2}$, the number of asyzygetic covariants $a^{\theta}x^{n}$ is
\[
= \text{coeff. } a^{\theta}x^{n-2\mu} \text{ in } \frac{1-x}{(1-a)(1-ax)(1-ax^{2})},
\]
(see Second Memoir, No.~35, observing that $q$ is there $=\theta - \frac{1}{2}\mu$, and that the subtraction of successive coefficients is effected by means of the factor $1-x$ in the numerator. See also Eighth Memoir, No.~251, where a like form is used for the quintic). Writing $ax^{2}$ for $a$, and $\frac{1}{x}$ for $x$, this is
\[
= \text{coeff. } a^{\theta}x^{n} \text{ in } \frac{1}{\frac{1}{x}(1-ax^{2})(1-a)(1-\frac{1}{x})}.
\]
The development is
\begin{align*}
&1\\
&+ ax(x^{2}+1)\\
&+ a^{2}(x^{4}+x^{2}+1)\\
&+ a^{3}(x^{6}+x^{4}+x^{2}+1)\\
&\quad\ldots\ldots\ldots\ldots\ldots\ldots
\end{align*}
which is
\[
= A(x) + A\!\left(\frac{1}{x}\right),
\]
where
\[
A(x) = \frac{1}{(1-ax^{2})(1-a)(1-\frac{1}{x})},
\]
and, since $A(\frac{1}{x})$ contains only negative powers, the required number is
\[
= \text{coeff. } a^{\theta}x^{n} \text{ in } \frac{1}{(1-ax^{2})(1-a)} = \frac{1}{1-aa^{2}} \cdot \frac{1}{1-a^{2}},
\]
indicating that the covariants are powers and products of $(ax^{2}$ and $a^{2})$, the quadric itself, and the discriminant. Compare Second Memoir, No.~49, according to which, writing therein $a$ for $x$, the number of asyzygetic covariants is
\[
= \text{coeff. } a^{\theta} \text{ in } \frac{1}{(1-a)(1-a^{2})}.
\]

334. For the cubic $(a, b, c, d\hspace{-0.2em}\pitchfork x, y)^{3}$ the number of asyzygetic covariants $a^{\theta}x^{n}$ is
\[
= \text{coeff. } a^{\theta}x^{n-2\mu} \text{ in } \frac{1-x}{(1-a)(1-ax)(1-ax^{2})(1-ax^{3})},
\]
or transforming as before, this is
\[
= \text{coeff. } a^{\theta}x^{n} \text{ in } \frac{1-x^{-1}}{(1-ax^{3})(1-ax^{2})(1-ax)(1-ax^{-1})}.
\]

\newpage
\setcounter{page}{339}

the function is here
\[
A(x) + A\!\left(\frac{1}{x}\right),
\]
where
\[
A(x) = \frac{1-ax^{2}}{(1-ax^{3})(1-ax^{2})(1-ax)(1-a^{4})},
\]
(that this is so may be easily verified); and since the second term contains only negative powers, the required number is $= \text{coeff. } a^{\theta}x^{n}$ in $A(x)$. The formula, in fact, indicates that the covariants are made up of $(ax^{3}, ax^{2}, ax^{4}, a^{4})$, the cubic itself, the Hessian, the cubicovariant, and the discriminant, these being connected by a syzygy $(a^{2}x^{6})$ of the degree 6 and order 6. Compare Second Memoir, No.~50, according to which, the number of covariants of degree $\theta$ is
\[
= \text{coeff. } a^{\theta} \text{ in } \frac{1-a^{6}}{(1-a)(1-a^{2})(1-a^{3})(1-a^{4})}.
\]

335. For the quartic $(a, b, c, d, e\hspace{-0.2em}\pitchfork x, y)^{4}$ the number of asyzygetic covariants $a^{\theta}x^{n}$ is
\[
= \text{coeff. } a^{\theta}x^{n-2\mu} \text{ in } \frac{1-x}{(1-a)(1-ax)(1-ax^{2})(1-ax^{3})(1-ax^{4})},
\]
or transforming as before, this is
\[
= \text{coeff. } a^{\theta}x^{n} \text{ in } \frac{1-x^{-1}}{(1-ax^{4})(1-ax^{3})(1-ax^{2})(1-ax)(1-ax^{-1})}.
\]
the function is here
\[
A(x) + A\!\left(\frac{1}{x}\right),
\]
where
\[
A(x) = \frac{1-ax^{2}}{(1-ax^{4})(1-ax^{3})(1-ax^{2})(1-ax)(1-a^{5})(1-a^{6})(1-a^{2}x^{2})},
\]
and the second term containing only negative powers, the required number is $= \text{coeff. } a^{\theta}x^{n}$ in $A(x)$. The formula indicates that the covariants are made up of $(ax^{4}, ax^{2}, a^{2}, a^{3}, a^{2}x^{4})$, the quartic itself, the Hessian, the quadrinvariant, the cubinvariant, and the cubicovariant, these being connected by a syzygy $(a^{3}x^{12})$ of the degree 6 and order 12. Compare Second Memoir, No.~51, according to which the number of covariants of degree $\theta$ is
\[
= \text{coeff. } a^{\theta} \text{ in } \frac{1-a^{6}}{(1-a)(1-a^{2})(1-a^{3})(1-a^{4})(1-a^{5})}.
\]

336. For the quintic $(a, b, c, d, e, f\hspace{-0.2em}\pitchfork x, y)^{5}$ the number of asyzygetic covariants $a^{\theta}x^{n}$ is
\[
= \text{coeff. } a^{\theta}x^{n-2\mu} \text{ in } \frac{1-x}{(1-a)(1-ax)(1-ax^{2})(1-ax^{3})(1-ax^{4})(1-ax^{5})},
\]
or transforming as before, this is
\[
= \text{coeff. } a^{\theta}x^{n} \text{ in } \frac{1-x^{-1}}{(1-ax^{5})(1-ax^{4})(1-ax^{3})(1-ax^{2})(1-ax)(1-ax^{-1})}.
\]

\newpage
\setcounter{page}{340}

The developed expression is
\begin{align*}
&\frac{1}{x}\\
&+ aa^{2}\\
&+ a^{2}(x^{4} + x^{2} + 1)\\
&+ a^{3}(x^{-6} + x^{-4} + x^{-2})\\
&\quad\ldots\ldots\ldots\ldots\ldots\ldots
\end{align*}
but here there is not any finite function $A(x)$ such that this development is

The numerical coefficients are of course the same as those in the development of the untransformed function; viz. they are the numbers given in the third column of Table No.~82 (Eighth Memoir), and also (carried further) in the third column of the following Table, No.~87. And we can, from the discussion of these coefficients, deduce the form of $A(x)$, viz. this is
\begin{center}
\begin{tabular}{cccccccccc}
& & & & & 1 & & 1 & & 1\\
& & & & & $\overline{1-a^{2}x^{14}}$ & & $\overline{1-a^{2}x^{13}}$ & & $\overline{(1-a^{2}x^{10})^{2}}$\\[0.5em]
& & & & & 1 & & 1 & & 1\\
& & & & & $\overline{1-a^{2}x^{12}}$ & & $\overline{1-a^{2}x^{11}}$ & & $\overline{(1-a^{2}x^{8})^{2}}$\\[0.5em]
& & & & & 1 & & 1 & & \\
& & & & & $\overline{1-a^{2}x^{10}}$ & & $\overline{(1-a^{2}x^{9})^{2}}$ & & $\overline{(1-a^{2}x^{6})^{2}}$\\[0.5em]
& & & & & 1 & & 1 & & \\
& & & & & $\overline{1-a^{2}x^{8}}$ & & $\overline{1-a^{2}x^{7}}$ & & \\
[0.5em]
& & & & & 1 & & & & \\
& & & & & $\overline{1-a^{2}x^{6}}$ & & & &
\end{tabular}
\end{center}

\begin{center}
\begin{tabular}{cccccccccc}
1 & & 1 & & 1 & & 1 & & 1 & $\ldots$\\
$\overline{1-ax^{2}}$ & & $\overline{1-ax^{5}}$ & & $\overline{1-a^{2}x^{4}}$ & & $\overline{1-a^{2}x^{3}}$ & & $\overline{1-a^{2}x^{2}}$ & \\
& & & & & & & & & \\
2 & & 5 & & 4 & & 3 & & 2 & \\
& & & & 3 & & 1 & & & 20\\
& & & & & & & & & 14\\
& & & & & & & & & 8
\end{tabular}
\end{center}
where, for shortness, I have written $1-a^{2}x^{14}$ to stand for $(1-a^{2}x^{14})(1-a^{2}x^{13})$, and so in other cases; moreover in the third column of the numerator the $(9)^{2}$ shows that the factor is $(1-a^{2}x^{9})^{2}$, and so in other cases: this will be further explained presently.

Compare herewith the form, Second Memoir, No.~52, viz. the number of asyzygetic covariants of the degree $\theta$ is
\[
= \text{coeff. } a^{\theta} \text{ in } (1-a)^{-\kappa_{1}}(1-a^{2})^{-\kappa_{2}}(1-a^{3})^{-\kappa_{3}}(1-a^{4})^{-\kappa_{4}}(1-a^{5})^{-\kappa_{5}}(1-a^{6})^{-\kappa_{6}}(1-a^{7})^{-\kappa_{7}}(1-a^{8})^{-\kappa_{8}}\ldots,
\]
each index being, it will be observed, equal to the number of factors in the numerator, less the number of factors in the denominator, in the corresponding column of the new formula.

\newpage
\setcounter{page}{341}

\textbf{Article Nos.\ 337 to 346. The 23 Fundamental Covariants.}

337. Gordan's result is that the number of the irreducible covariants of the binary quintic is $=23$. I represent these by the letters $A$, $B$, $C$, \ldots, $W$, identifying such of them as were given in my former Memoirs on Quantics with the Tables of these Memoirs, and the new ones, $O$, $P$, $R$, $S$, $T$, $V$, with the Tables Nos.~90, 91, 92, 93, 94, 95 of the present Memoir.

\begin{center}
\textbf{Table No.~87.} Identification of the 23 irreducible covariants of the binary quintic.
\end{center}

\begin{center}
\begin{tabular}{clcl}
& & & Table No.\\\hline
$A$ & $= (a, b, c, d, e, f\hspace{-0.2em}\pitchfork x, y)^{5}$ & $\phantom{=}\, f = (f)^{1}$ & 13\\[0.3em]
$B$ & $= \frac{1}{12}(A,A)^{2}$ & $(\hspace{-0.2em}\pitchfork x, y)^{2}\, \phi = (f,f)^{2}$ & 14\\[0.3em]
$C$ & $= \frac{1}{12}(A,B)^{1}$ & $(\hspace{-0.2em}\pitchfork x, y)^{3}\, \psi = (f,\phi)^{1}$ & 15\\[0.3em]
$D$ & $= -\frac{1}{2}(A,B)^{2}$ & $(\hspace{-0.2em}\pitchfork x, y)^{1}\, i = (f,\phi)^{2}$ & 16\\[0.3em]
$E$ & $= \frac{1}{4}(A,C)^{2}$ & $(\hspace{-0.2em}\pitchfork x, y)^{2}\, (f,\psi)^{2}$ & 17\\[0.3em]
$F$ & $= \frac{1}{2}(A,C)^{3}$ & $(\hspace{-0.2em}\pitchfork x, y)^{1}\, (f,\psi)^{3}$ & 18\\[0.3em]
$G$ & $= -\frac{1}{4}(B,B)^{2}$ & $(\hspace{-0.2em}\pitchfork x, y)^{2}\, \xi = (\phi,\phi)^{2}$ & 19\\[0.3em]
$H$ & $= -\frac{1}{2}(B,C)^{2} + \frac{1}{4}B^{2}$ & $(\hspace{-0.2em}\pitchfork x, y)^{2}\, \rho = (\phi,\psi)^{2}$ & 20\\[0.3em]
$I$ & $= -\frac{1}{2}(B,C)^{3}$ & $(\hspace{-0.2em}\pitchfork x, y)^{1}\, (\phi,\psi)^{3}$ & 21\\[0.3em]
$J$ & $= -\frac{1}{2}(B,D)^{2}$ & $(\hspace{-0.2em}\pitchfork x, y)^{1}\, \sigma = (j)^{2}$ & 22\\[0.3em]
$K$ & $= \phantom{-}(B,D)^{3}$ & $(\hspace{-0.2em}\pitchfork x, y)^{1}\, (j)^{3}$ & 23\\[0.3em]
$L$ & $= -\frac{1}{2}(A,H)^{3} + \frac{1}{4}BE$ & $(\hspace{-0.2em}\pitchfork x, y)^{1}\, (f\rho)^{3}$ & 24\\[0.3em]
$M$ & $= -\frac{1}{12}(B,H)^{2} - \frac{1}{4}BG$ & $(\hspace{-0.2em}\pitchfork x, y)^{2}\, \tau = (\rho,\psi)^{2}$ & 83\\[0.3em]
$N$ & $= \phantom{-}(B,H)^{3}$ & $(\hspace{-0.2em}\pitchfork x, y)^{1}\, (\rho)^{3}$ & 84\\[0.3em]
$O$ & $= -\phantom{\frac{1}{2}}(B,J)^{2}$ & $(\hspace{-0.2em}\pitchfork x, y)^{1}\, (\sigma)^{2}$ & *90\\[0.3em]
$P$ & $= -\frac{1}{2}(A,M)^{3} - BK$ & $(\hspace{-0.2em}\pitchfork x, y)^{1}\, (f\tau)^{3}$ & *91\\[0.3em]
$Q$ & $= \frac{1}{12}(B,M)^{2}$ & $(\hspace{-0.2em}\pitchfork x, y)^{2}\, (\tau)^{2}$ & 25\\[0.3em]
$R$ & $= -\frac{1}{2}(B,M)^{3}$ & $(\hspace{-0.2em}\pitchfork x, y)^{1}\, (\tau)^{3}$ & *92\\[0.3em]
$S$ & $= -96(D,M)^{2} + 160J - 7GK$ & $(\hspace{-0.2em}\pitchfork x, y)^{1}\, (\tau)^{1}$ & *93\\[0.3em]
$T$ & $= -\phantom{\frac{1}{2}}(J,M)^{2}$ & $(\hspace{-0.2em}\pitchfork x, y)^{1}\, \gamma = (j\tau)^{2}$ & *94\\[0.3em]
$U$ & $= \phantom{-}(J,O)^{2} + \frac{1}{4}OQ$ & $(\hspace{-0.2em}\pitchfork x, y)^{1}\, ((\sigma)^{2},\psi)^{2}$ & 29\\[0.3em]
$V$ & $= -\phantom{\frac{1}{2}}(B,T)^{2}$ & $(\hspace{-0.2em}\pitchfork x, y)^{1}\, (\xi\gamma)^{2}$ & *95\\[0.3em]
$W$ & $= -\frac{1}{2}(O,T)^{2}$ & $(\hspace{-0.2em}\pitchfork x, y)^{1}\, ((\sigma)^{2},\gamma)^{2}$ & 29$\alpha$
\end{tabular}
\end{center}

\newpage
\setcounter{page}{342}

338. The Table exhibits the generation of the several covariants; viz. $(A, B)$ denotes $\partial_{x}A \cdot \partial_{y}B - \partial_{y}A \cdot \partial_{x}B$, $(A, B)^{2}$ denotes $\partial_{x}^{2}A \cdot \partial_{y}^{2}B - 2\partial_{x}A\partial_{y}A \cdot \partial_{x}B\partial_{y}B + \partial_{y}^{2}A \cdot \partial_{x}^{2}B$, \&c.\ (see post.\ No.~348). The column $f$, $i = (ff)^{2}$, \&c.\ shows Gordan's notation, and the generation of his 23 forms ($(ff)^{2}$ written as with him for $(f,f)^{2}$, \&c.): it will be observed that the forms are not identical; if the calculations had been made de novo, I should have adopted his values, simply omitting numerical factors of the several forms (thus every term of $t$, $=(ff)^{4}$ contains the factor $2 \cdot (120)^{2} = 28800$): of course the presence of these numerical factors renders the $f$, $t$, $\psi$, \&c.\ as they stand inconvenient for the expression of results; and the numerical fixation of the values was no part of Gordan's object. But by reason of the existing Tables the change of notation is in fact more than this; thus $H$ instead of being a submultiple of $(B,C)^{2}$, that is, of $\frac{1}{2}\psi$, is in fact $=-\frac{1}{2}(B,C)^{2} + \frac{1}{4}B^{2}$; and so in other cases. If the occasion for it arises, there is no difficulty in expressing any one of the forms $f$, $t$, $\psi$, \&c.\ in terms of the $(A, B, C \ldots V, W)$; thus in the instance just referred to, $\rho = (\phi C)^{2}$, we have
\[
\phi = (\phi)^{2} = (\psi, A)^{2} = 800C,
\]
and
\[
t = \tfrac{1}{2}(ff)^{4} = \tfrac{1}{2}(\psi, A)^{4} = 2 \cdot 120 \cdot 8B,
\]
whence $\rho = 2304000\, (B, C)^{2}$; also $(B, C)^{2} = -5H + 2R$; and therefore, finally,
\[
2J = -11520000\, H + 4608000\, B^{2}.
\]

339. I remark upon the value $S = -96(D, M)^{2} + 160J - 7GK$, that $S$ is the complete value of a covariant $(\hspace{-0.2em}\pitchfork x, y)^{1}(\hspace{-0.2em}\pitchfork x, y)^{1}$, the leading coefficient of which is given in Table No.~86 of my Eighth Memoir; the form $(D, M)$, omitting a numerical factor (if any), would have had smaller numerical coefficients, but there is in the form actually adopted the advantage that it vanishes for $a=0$, $b=0$, that is, when the quintic has two equal roots, [see post.\ No.~346].

340. I now form the following Table No.~88, viz. this is the Table No.~82 of my Eighth Memoir, carried as far as $a^{8}$, but with the composite covariants expressed by means of the foregoing letters $A$, $B$, $C$, \ldots, $W$; instead of giving the syzygies as in Table No.~82, I transfer them to a separate Table, No.~89. In all other respects the arrangement is as explained, Eighth Memoir, No.~253; but in place of $N$, $S$, $S'$ I have written $\aleph$, $\zeta$, $\zeta'$ to denote new covariant, new syzygy, derived syzygy, respectively; and I have, as to the terms $a^{7}x^{5}$, $a^{8}x^{2}$ respectively, introduced the new symbol $\sigma$ to denote an interconnexion of syzygies, as appearing by the Table No.~89, and as will be further explained.

\newpage
\setcounter{page}{343}

\begin{center}
\textbf{Table No.~88.}
\end{center}

[In this Table and the subsequent Table 89, I have for convenience used, instead of capitals, the small italic letters $a$, $b$, $c$, \ldots $w$ to denote the 23 irreducible covariants of the quintic.]

\begin{center}
\begin{tabular}{ccc|l|c}
Ind.\ $a$. & Ind.\ $x$. & Coeff. & & \\\hline
1 & 5 & 1 & $a$ & $*$\\
& 3 & & & \\
& 1 & & & \\
\hline
2 & 10 & 1 & $a^{2}$ & \\
& 8 & & & \\
& 6 & 1 & $c$ & $\#$\\
& 4 & & & \\
& 2 & 1 & $b$ & $*$\\
\hline
3 & 15 & 1 & $a^{3}$ & \\
& 13 & & & \\
& 11 & 1 & & \\
& 9 & 1 & $f$ & $*$\\
& 7 & 1 & $ah$ & \\
& 5 & 1 & $e$ & $*$\\
& 3 & 1 & $d$ & $*$\\
& 1 & & &
\end{tabular}
\end{center}

\newpage
\setcounter{page}{344}

\begin{center}
\textbf{Table No.~88 (continued).}
\end{center}

\begin{center}
\begin{tabular}{ccc|l|c}
Ind.\ $a$. & Ind.\ $x$. & Coeff. & & \\\hline
4 & 20 & 1 & $a^{4}$ & \\
& 18 & & & \\
& 16 & 1 & $a^{2}c$ & \\
& 14 & 1 & $af$ & \\
& 12 & 2 & $a^{2}b$, $c^{2}$ & \\
& 10 & 1 & $ae$ & \\
& 8 & 2 & $ad$, $bc$ & \\
& 6 & 1 & $i$ & $*$\\
& 4 & 2 & $b^{2}$, $h$ & $*$\\
& 2 & & & \\
& & 1 & $g$ & $*$\\
\hline
5 & 25 & 1 & $a^{5}$ & \\
& 23 & & & \\
& 21 & 1 & $a^{3}c$ & \\
& 19 & 1 & $a^{2}f$ & \\
& 17 & 2 & $a^{2}b$, $ae^{2}$ & \\
& 15 & 2 & $a^{2}e$, $cf$ & \\
& 13 & 2 & $a^{2}d$, $abc$ & \\
& 11 & 2 & $ai$, $bf$, $ce$ & $\zeta$\\
& 9 & 3 & $ah^{2}$, $ah$, $cd$ & \\
& 7 & 2 & $bc$, $l$ & $\#$\\
& 5 & 2 & $ag$, $bd$ & \\
& 3 & 1 & $k$ & \\
& 1 & 1 & $j$ & $\aleph$
\end{tabular}
\end{center}

\newpage
\setcounter{page}{345}

\begin{center}
\textbf{Table No.~88 (concluded).}
\end{center}

\begin{center}
\begin{tabular}{ccc|l|c}
Ind.\ $a$. & Ind.\ $x$. & Coeff. & & \\\hline
8 & 40 & 1 & $a^{8}$ & \\
& 38 & & & \\
& 36 & 1 & $a^{6}c$ & \\
& 34 & 1 & $a^{5}f$ & \\
& 32 & 2 & $a^{4}b$, $a^{2}c^{2}$ & \\
& 30 & 2 & $a^{4}e$, $a^{3}cf$ & \\
& 28 & 3 & $a^{4}d$, $a^{3}bc$, $ac^{2}$, $a^{2}f^{2}$ & $\zeta'$\\
& 26 & 3 & $a^{3}i$, $a^{2}bf$, $a^{2}ce$, $ac^{2}f$ & $\zeta'$\\
& 24 & 5 & $a^{3}b^{2}$, $a^{3}h$, $a^{2}cd$, $a^{2}bc^{2}$, $a^{2}ef$, $cf^{2}$, $c^{3}$ & $\zeta'$\\
& 22 & 4 & $a^{3}be$, $a^{2}l$, $a^{2}ci$, $a^{2}df$, $abcf$, $ahe$, $c^{2}e$ & $2\zeta'$\\
& 20 & 6 & $a^{3}g$, $a^{2}bd$, $a^{2}b^{2}$, $a^{2}ch$, $a^{2}e^{2}$, $ac^{2}d$, $afi$, $be^{2}$, $bf^{2}$, $cef$ & $\sigma$\\
& 18 & 5 & $a^{2}j$, $a^{2}bi$, $a^{2}de$, $ab^{2}f$, $abce$, $adi$, $afh$, $c^{3}$, $cdf$ & $4\zeta'$\\
& 16 & 7 & $a^{2}k$, $a^{2}b^{2}$, $a^{2}bh$, $a^{2}cg$, $a^{2}d^{2}$, $abcd$, $aei$, $ab^{2}c$, $bef$, $c^{2}i$, $ce^{2}$, $fl$ & $5\zeta'$\\
& 14 & 5 & $a^{2}n$, $ab^{3}$, $abl$, $ack$, $adi$, $aeh$, $afg$, $bci$, $bdf$, $cde$ & $\sigma$\\
& 12 & 7 & $a^{2}bg$, $abn$, $ab^{2}l$, $acj$, $adh$, $b^{4}$, $b^{2}ch$, $bc^{2}$, $c^{2}k$, $cd^{2}$, $el$, $fk$, $i^{2}$ & $3\zeta$\\
& 10 & 5 & $abk$, $aeg$, $ap$, $b^{2}i$, $bde$, $cn$, $dl$, $fj$, $hi$ & $3\zeta$\\
& 8 & 6 & $abj$, $adg$, $b^{3}$, $b^{2}h$, $beg$, $bd^{2}$, $cm$, $ek$, $h^{2}$ & $2\zeta$\\
& 6 & 3 & $ao$, $bn$, $dk$, $ej$, $gi$ & $2\zeta$\\
& 4 & 4 & $ifg$, $bm$, $dj$, $gh$ & $\#$\\
& 2 & 1 & $r$ & $*$\\
& & 2 & $f$, $?$ &
\end{tabular}
\end{center}

\newpage
\setcounter{page}{346}

\begin{center}
\textbf{Table No.~89.}
\end{center}

[See ante Table No.~88.]

\begin{align*}
(5,\ 11)\qquad & ai + bf - ce = 0\\[0.5em]
(6,\ 18)\qquad & c^{2}d - a^{2}bc + ie^{2} + l^{2} = 0\\
(6,\ 14)\qquad & a^{2}h - bacd - ibc^{2} - ef = 0\\
(6,\ 12)\qquad & al - 2ci + 3df = 0\\
(6,\ 10)\qquad & a^{2}g - 4abd - 4b^{3} - e^{2} = 0\\
(6,\ 8)\qquad & ok + 2bi - 3de = 0\\
(6,\ 6)\qquad & aj - b^{2} + 2bh - cd - g^{2} = 0\\[0.5em]
(7,\ 15)\qquad & a^{2}bd - a^{2}bc + ach - b^{2}d - fi = 0\\
(7,\ 13)\qquad & a^{2}k - abi - 3bg + bci + 3fh = 0\\
(7,\ 11)\qquad & a^{2}j - ab^{2} + abh - 9ad^{2} - 6bcd - ei = 0\\
(7,\ 9)\qquad & an - b^{2}e - bdi + 2eh - fg = 0\\
&\quad m + bdi - eh + fg = 0\\
&\quad 2ck - bdi + eh - fg = 0\\
(7,\ 7)\qquad & am + 2b^{2}l + cj - Mh = 0
\end{align*}

\newpage
\setcounter{page}{347}

\textbf{Table No.~89 (continued).}

\begin{align*}
\sigma,\ (8,\ 20)\qquad & a^{2}(a^{2}g - 12abd - 4b^{3} - e^{2}) \quad \text{repr.\ } (6,\ 10)\\
&- a(a^{2}bd - abe^{2} + cch - 6c^{2}d - fi) \quad \text{repr.\ } (7,\ 15)\\
&+ b(a^{2}d - a^{2}bc + ic^{2} + l^{2}) \quad \text{repr.\ } (6,\ 18)\\
&+ c(a^{2}h - bacd - ibc^{2} - ef) \quad \text{repr.\ } (6,\ 14)\\
&- f(ai + bf - ce) = 0 \quad \text{repr.\ } (5,\ 11)\\[1em]
\sigma,\ (8,\ 14)\qquad & a(an - b^{2}e - bdi + 2eh - fg) \quad \text{repr.\ } (7,\ 9)\\
&+ a(2bl + bdi - eh + fg) \quad \text{repr.\ } (\quad,\ )\\
&+ a(2ck - 2di + eh - fg) \quad \text{repr.\ } (\quad,\ )\\
&- 2b(al - 2ci + 3df) \quad \text{repr.\ } (6,\ 12)\\
&- 2c(ok + 2bi - 3de) \quad \text{repr.\ } (6,\ 8)\\
&+ 6d(ai + bf - ce) = 0 \quad \text{repr.\ } (5,\ 11)\\[1em]
(8,\ 12)\qquad & ab^{2}d - b^{2}c + 2bch - c^{2}g + i^{2} = 0\\
&\quad -3adi - 2bch + 2c^{2}g + 18cd^{2} + fk - 2i^{2} = 0\\
&\quad el + fk - 2i^{2} = 0\\[0.5em]
(8,\ 10)\qquad & abk - cn - bl - 2fj + hi = 0\\
&\quad ap + 2cm + lj = 0\\
&\quad b^{2}i + cn + 3dl + fj - 2hi = 0\\[0.5em]
(8,\ 8)\qquad & abj - b^{3} + 4b^{2}h - 9bd^{2} + 12cm - ek - 3h^{2} = 0\\
&\quad adg + 2b^{2}h - 12bd^{2} + 3cm - ek - 2h^{2} = 0\\[0.5em]
(8,\ 6)\qquad & ao + bk - 3ej + 2gi = 0\\
&\quad bn + 3dk - ej + gi = 0
\end{align*}

\newpage
\setcounter{page}{348}

342. In illustration take any one of the lines of Table No.~88, for instance [resuming the notation by capital letters] the line
\[
(7,\ 17)\quad 4\quad A^{2}BE,\ A^{2}L,\ ACI,\ ADF,\ BCF,\ OE \quad|\quad 2\zeta' \quad|\quad
\]
there are here 6 composite covariants, but the number of asyzygetic covariants is $=4$: there must therefore be $6 - 4, = 2$ syzygies; we have however (see Table No.~89) two derived syzygies of the right form, viz. these are
\[
A(AL - 2CI + 3DF) = 0,\qquad C(AI + BF - CE) = 0,
\]
which are designated as $2\zeta'$, and there is consequently no new syzygy $\zeta$.

But in the line
\[
(7,\ 15)\quad 5\quad A^{3}G,\ A^{2}BD,\ AR,\ ACH,\ AE^{2},\ CD,\ FI \quad|\quad \zeta',\ \zeta \quad|
\]
there are 7 composite covariants, but the number of asyzygetic covariants is $=5$; there must therefore be $7 - 5, = 2$ syzygies. One of these is the derived syzygy
\[
A(A^{2}G - E^{2} - 12ABD - 4R - G) = 0,
\]
which is designated by $\zeta'$; the other is a new syzygy (see Table No.~89),
\[
A^{2}BD - ABC^{2} + ACH - GCD - FI = 0,
\]
designated by $\zeta$.

343. Take now the line
\[
(8,\ 20)\quad 6\quad A^{4}G,\ A^{3}BD,\ A^{2}B^{2}C,\ A^{2}CH,\ A^{2}E^{2},\ ACD,\ AFI,\ BG^{2},\ BF^{2},\ GEF \quad|\quad 5\zeta',\ \sigma \quad|\quad
\]
there are here 10 composite covariants, but the number of irreducible covariants is $=6$; there should therefore be $10 - 6, = 4$ syzygies. There are, however, the 5 derived syzygies
\[
A^{2}(A^{2}G - 12ABD - 4R - G - E^{2}) = 0,\ \&\text{c.}\ \text{(see Table No.~89)}
\]
designated by $5\zeta'$; since these are equivalent to 4 syzygies only there must be 1 identical relation between them (designated by $\sigma$), viz. this is the equation $= 0$ obtained by adding the several syzygies, multiplied each by the proper numerical factor as shown Table No.~89.

344. Again, for the line
\[
(8,\ 14)\quad 5\quad A^{3}BH,\ AB^{2}E,\ ABL,\ AGK,\ ADI,\ AEH,\ AFG,\ BGI,\ BDF,\ GDE \quad|\quad 6\zeta',\ \sigma \quad|
\]
there are here 10 composite covariants, but only 5 irreducible covariants; there should therefore be $10 - 5, = 5$ syzygies; we have in fact the 6 derived syzygies
\[
A(AN - B^{2}E - 6DI + 2Eh - FG) = 0,\ \&\text{c.}\ \text{(see Table No.~89)}
\]
designated by $6\zeta'$; these must therefore be connected by 1 identical relation (designated by $\sigma$), viz. this is the equation $= 0$ obtained by adding the several syzygies, each multiplied by the proper numerical factor as shown Table No.~89.

345. These two cases ($\sigma$) are in fact the instances which present themselves where a correction is required to my original theory. The two identical relations in question were disregarded in my original theory, and this accordingly gave the two non-existent irreducible covariants $(a, \ldots \hspace{-0.2em}\pitchfork x, y)^{14}$ and $(a, \ldots \hspace{-0.2em}\pitchfork x, y)^{16}$. And reverting to No.~336, these give in the denominator of $A(x)$ the factors $(1-a^{2}x^{14})(1-a^{2}x^{16})$. In virtue hereof, writing $x=1$, we have in $A(x)$ the factor $(1-a)^{-\kappa_{1}}(1-a^{2})^{-\kappa_{2}}\ldots$ agreeing with the function $(1-a)^{-\kappa_{1}}(1-a^{2})^{-\kappa_{2}}\ldots(1-a^{8})^{-\kappa_{8}}\ldots$. And we thus see that the denominator factors of $A(x)$ do not all of them refer to irreducible covariants; viz. we have
\[
aa^{2},\ a^{2}x^{4},\ a^{2}x^{3},\ a^{2}x^{2},\ a^{3},\ a^{3}x^{4},\ a^{3}x^{2},\ a^{4}x^{2},\ a^{4},\ a^{4}x^{6},\ a^{5}x^{2},\ a^{5}x^{4},\ a^{6}x^{2},\ a^{6}x^{4},\ a^{7}x^{2},\ a^{8}x^{2},\ a^{8},
\]
each referring to an irreducible covariant, but $a^{8}x^{14}$ and $a^{8}x^{16}$ each referring to an identical relation ($\sigma$) or interconnexion of syzygies. And we thus understand how, consistently with the number of the irreducible covariants being finite, the expression for $A(x)$ may be as above the quotient of two infinite products; there will be in the denominator a finite number of factors each referring to an irreducible covariant, but the remaining infinite series of denominator factors will refer each factor to an identical relation or interconnexion of syzygies. But I do not see how we can by the theory distinguish between the two classes of factors, so as to determine the number of the irreducible covariants, or even to make out affirmatively that the number of them is finite.

346. The new covariants $O$, $P$, $R$, $S$, $T$, $V$ are as follows:

[Table No.~90 (Covariant $O$),

Table No.~91 (Covariant $P$),
Table No.~92 (Covariant $R$),
Table No.~93 (Covariant $S$),
Table No.~94 (Covariant $T$),
Table No.~95 (Covariant $V$),

printed in the paper 143, ``Tables of the Covariants $M$ to $W$ of the Binary Quintic: from the Second, Third, Fifth, Eighth, Ninth and Tenth Memoirs on Quantics'' with the insertion as therein mentioned of the terms with zero coefficients. The covariant $S$, $=-96(D,M)^{2}+160J-7GK$, of the present Memoir is there called $S'$, and there is given the more simple form $S=(D,M)$, of this covariant.]

\newpage
\setcounter{page}{349}

\textbf{Article Nos.\ 347 to 365. Sketch of Professor Gordan's proof for the finite Number, $=23$, of the Covariants of a Binary Quintic.}

347. I propose to reproduce the leading points of Professor Gordan's proof that the binary quintic $(a, b, c, d, e, f\hspace{-0.2em}\pitchfork x, y)^{5}$ has a finite system of 23 covariants, viz. a system such that every other covariant whatever is a rational and integral function of these 23 covariants.

348. \textbf{Derivation.} Consider for a moment any two binary quantics $\phi$, $\psi$ of the same or different orders, and which may be either independent quantics, or they may be both or one of them covariants, or a covariant, of a binary quantic $f$. We may form the series of derivatives
\begin{align*}
(\phi, \psi)^{0} &= \phi\psi,\\
(\phi, \psi)^{1} &= 12\, \phi_{x}\psi_{y} - \phi_{y}\psi_{x},\\
(\phi, \psi)^{2} &= 12^{2}\, \phi_{xx}\psi_{yy} - 2\phi_{x}\phi_{y} \cdot \psi_{x}\psi_{y} + \phi_{yy}\psi_{xx},
\end{align*}
where, however, there is no occasion to use the notation $(\phi, \psi)^{0}$ (as this is simply the product $\phi\psi$), and the succeeding derivatives may (when there is no risk of ambiguity) be written more shortly $(\phi\psi)$, $(\phi\psi)_{2}$, $(\phi\psi)_{3}$, \&c.; in all that follows the word ``derivative'' (Gordan's Uebereinanderschiebung) is to be understood in this special sense.

349. The degree of the derivative $(\phi, \psi)^{k}$ is the sum of the degrees of the constituents $\phi$, $\psi$; the order of the derivative is the sum of the orders less $2k$; it being understood throughout that the word degree refers to the coefficients, and the word order to the variables. In speaking generally of the covariants or of all the covariants of a quantic $f$, or of the covariants or all the covariants of a given degree or order, we of course exclude from consideration covariants linearly connected with other covariants (for otherwise the number of terms would be infinite); but unless it is expressly so stated, we do not carry this out rigorously so as to make the system to consist of asyzygetic covariants; viz. it is assumed that the system is complete, but not that it is divested of superfluous terms.

350. \textbf{Theorem A.} The covariants of a quantic $f$ of a given degree $m$ can be all of them obtained by derivation from $f$ and the covariants of the next inferior degree $(m-1)$.

In particular for the degree 1 the only covariant is the quantic $f$ itself; for the degree 2 the covariants are $(ff)^{2}$, $(ff)^{4}$, $(ff)^{6}$, \ldots: using for a moment $\beta$ to denote each of these in succession, the covariants of the third degree are $(\beta f)^{0}$, $(\beta f)^{2}$, $(\beta f)^{4}$, \ldots\ and so on.

351. Suppose that the covariants of the second degree $(ff)^{2}$, $(ff)^{4}$, $(ff)^{6}$, \ldots\ are in this order represented by $\beta_{1}$, $\beta_{2}$, $\beta_{3}$\ldots, then the covariants of the third degree written in the order
\[
(\beta_{1}f)^{0},\ (\beta_{1}f)^{2},\ (\beta_{1}f)^{4},\ \ldots\ (\beta_{2}f)^{0},\ (\beta_{2}f)^{2},\ (\beta_{2}f)^{4},\ \ldots\ (\beta_{3}f)^{0},\ (\beta_{3}f)^{2},\ (\beta_{3}f)^{4},\ \ldots
\]
may be represented by $\gamma_{1}$, $\gamma_{2}$, $\gamma_{3}$, \ldots, the covariants of the fourth degree written in the order
\[
(\gamma_{1}f)^{0},\ (\gamma_{1}f)^{2},\ (\gamma_{1}f)^{4},\ \ldots\ (\gamma_{2}f)^{0},\ (\gamma_{2}f)^{2},\ (\gamma_{2}f)^{4},\ \ldots\ (\gamma_{3}f)^{0},\ (\gamma_{3}f)^{2},\ (\gamma_{3}f)^{4},\ \ldots
\]
may be represented by $\delta_{1}$, $\delta_{2}$, $\delta_{3}$\ldots, and so on: we thus obtain in a definite order the covariants of a given degree $m$; say, these are $\mu_{1}$, $\mu_{2}$, $\mu_{3}$, $\mu_{4}$, \ldots; any term $\mu_{t}$ is said to be a later term than the preceding terms $\mu_{1}$, $\mu_{2}$, \ldots\ and an earlier term than the following ones, $\mu_{t}$, $\mu_{t+1}$, \&c.

Observe that each term $\mu_{t}$ is a derivative $(\lambda, f)^{k}$, the derivatives of an earlier $\lambda$ are earlier than those of a later $\lambda$; and as regards the derivatives of the same $\lambda$, the derivative with a less index of derivation is earlier than that with a greater index of derivation, or, what is the same thing, those are earlier which are of the higher order.

352. The series $\mu_{1}$, $\mu_{2}$, $\mu_{3}$, $\mu_{4}$, \ldots\ is not asyzygetic; we make it so, by considering in succession whether the several terms $\mu_{1}$, $\mu_{2}$, \ldots\ respectively are expressible as linear functions of the earlier terms, and by omitting every term which is so expressible.

The reduced series thus obtained is called $\mathrm{T}_{1}$, $\mathrm{T}_{2}$, $\mathrm{T}_{3}$, \ldots. Observe that not every $\mu$ is a T, but that every T is a $\mu$; every T therefore arises from a derivation upon $f$ and a certain term $\lambda$; which term $\lambda$ (supposing the $\lambda$ series reduced in like manner to $\mathrm{S}_{1}$, $\mathrm{S}_{2}$, $\mathrm{S}_{3}$, \ldots) is a linear function of certain of the S's. Each later T is derived from later S's, or it may be from the same S's as an earlier T; viz. if the later T is derived from $(\mathrm{S}_{1}, \mathrm{S}_{2}, \ldots \mathrm{S}_{a})$, then the earlier T is derived, it may be, from $(\mathrm{S}_{1}, \mathrm{S}_{2}, \ldots \mathrm{S}_{a})$, or from $(\mathrm{S}_{1}, \mathrm{S}_{2}, \ldots \mathrm{S}_{a-k})$, but so that there is not in the series any term later than $\mathrm{S}_{a}$.

\newpage
\setcounter{page}{350}

And if, considering any T as thus derived from certain of the S's, and in like manner each of these S's as derived from certain of the R's, and so on, we descend to any preceding series,
\[
\mathrm{M}_{1},\ \mathrm{M}_{2},\ \mathrm{M}_{3},\ \ldots
\]
it will appear that the T is derived from a certain number $(\mathrm{M}_{1}, \mathrm{M}_{2}, \ldots \mathrm{M}_{r})$ of the terms of this series.

353. The quadricovariants $(ff)^{2}$, $(ff)^{4}$, $(ff)^{6}$, \ldots\ are of different orders, and consequently asyzygetic. They form therefore a series such as the T-series, and they may be represented by
\[
\mathrm{B}_{1},\ \mathrm{B}_{2},\ \mathrm{B}_{3},\ \ldots
\]
Supposing $f$ to be of the order $n$, $\mathrm{B}_{1}$ is of the order $2n$, $\mathrm{B}_{2}$ of the order $2n-4$, $\mathrm{B}_{3}$ of the order $2n-8$, and so on. Those terms which are of an order greater than $n$, are said to be of the form W (agreeing with a subsequent more general definition of W); those which are of an order equal to or less than $n$, are said to be of the form X: that is, the earlier terms of the B series are W, and the later terms are X; viz. the X terms taken in order, beginning with the earliest, are $\mathrm{X}_{1}$, $\mathrm{X}_{2}$, $\mathrm{X}_{3}$, \ldots

354. By what precedes any particular T is derived from certain terms $\mathrm{B}_{1}$, $\mathrm{B}_{2}$, \ldots\ $\mathrm{B}_{l}$ of the B series. This series, $\mathrm{B}_{1}$, $\mathrm{B}_{2}$, \ldots\ $\mathrm{B}_{l}$, may stop short of the terms X; it may include a certain number of them, say $\mathrm{X}_{1}$, $\mathrm{X}_{2}$, \ldots\ $\mathrm{X}_{k}$. The terms derived from the X's are in the sequel denoted by P.

355. Every covariant whatever is a form or sum of forms such as
\[
\mathrm{T}_{1}^{\alpha}\, \mathrm{B}_{2}^{\beta}\, \mathrm{B}_{3}^{\gamma}\, \ldots\ \mathrm{B}_{l}^{\lambda}\, \mathrm{P}_{1}^{\mu}\, \ldots\ \mathrm{P}_{t}^{\nu};
\]
writing in regard to any such expression
\[
\Sigma \text{ ind.\ } 1 = i,\ \Sigma \text{ ind.\ } 2 = j,\ \ldots
\]
(viz. $i$ is the sum of all those indices $\alpha$, $\beta$, \&c.\ which belong to a term containing the symbolic number 1, $j$ the sum of all the indices $\alpha$, $\gamma$, \&c.\ which belong to a term containing the symbolic number 2, and so on) then each of the numbers $i$, $j$, \ldots\ is at most $= n$, that is $n-i$, $n-j$, \ldots\ may be any of them $= 0$, but they cannot any of them be negative; the degree of the function is $= m$, and its order is $= mn - i - j \ldots$. It is to be further observed that the form is a function of the differential coefficients of $f$ of the orders $n-i$, $n-j$, \&c.\ respectively. It follows that if $n-i$, $n-j$, \ldots\ are none of them $=0$, the form in question may be obtained from a like form belonging to a quantic $f'$ of the next inferior order $n-1$ by replacing therein the coefficients $a'$, $b'$, \ldots\ by $ax+by$, $bx+cy$, \&c.\ respectively: for example, if $f$ denote the cubic function $(a, b, c, d\hspace{-0.2em}\pitchfork x, y)^{3}$, then the Hessian hereof is $12\, f_{x}f_{y}$; the like form in regard to the quadric $f' = (a', b', c'\hspace{-0.2em}\pitchfork x, y)^{2}$ is $12\, f'_{x}f'_{y}$, which is $= a'c' - b'^{2}$; and substituting herein $ax+by$, $bx+cy$, $cx+dy$ for $a'$, $b'$, $c'$ respectively, we have the Hessian $12\, f_{x}f_{y}$ of the cubic. A covariant of $f$ derivable in this manner from a covariant of the next inferior quantic $f'$ is said to be a special covariant.

\vspace{2em}
\begin{center}
\rule{0.5\textwidth}{0.4pt}
\end{center}
\clearpage


% =========================================================
% Chunk: cayley_vol07_pages_351_400.tex  (orig pages 351--400)
% =========================================================
\thispagestyle{fancy}

\vspace{-1em}
\noindent\rule{\textwidth}{0.4pt}

%% ========================================================================
%% PAPER 462 (continued): A NINTH MEMOIR ON QUANTICS
%% ========================================================================
\section*{462. A Ninth Memoir on Quantics. (Continued)}
\addcontentsline{toc}{section}{462. A Ninth Memoir on Quantics (pp.~351--353)}

\noindent\textit{[From the Philosophical Transactions of the Royal Society of London, vol.~CLXI (for the year 1871), pp.~17--50. Read November 10, 1870.]}

\medskip
\setcounter{equation}{0}

\noindent\textbf{356.}\quad Reverting to the form
\[
12^\alpha 13^\beta 23^\gamma \ldots f_i f_j \ldots f_k;
\]
if, as before, $n-i$, $n-j$, \&c.\ are each of them $> \frac{1}{2}n$; if there is at least one index $i$ which is $=$ or $< \frac{1}{2}n$ (that is, for which $n-i > \frac{1}{2}n$), and if the order $mn-i-j\ldots$ be $>n$, then the form, or any sum of such forms, is said to be a form or covariant $W$.

Every covariant $W$ is thus a special covariant, but not conversely. In the particular case $m=2$, the form is
\[
12^\alpha f_i f_j \ldots f_k,
\]
which will be a form $W$ if $n-a > \frac{1}{2}n$, or, what is the same thing, $2n-2a > n$, that is if the order be $>n$. Hence, as already mentioned, the covariants $T$ of the degree 2 are $W$, or else $\not{W}$, according as the order is greater than $n$, or as it is equal to or less than $n$.

\smallskip
\noindent\textbf{357. Theorem B.}\quad If any covariant $T$ be expressible as the sum of a form $W$ and of earlier $T$'s than itself, then forming the derivative $(Tf)^2$; either this is not a form $T$, or being a form $T$, it is expressible as the sum of a form $W$ and of earlier $T$'s than itself; or, what is the same thing, $(Tf)^2$, if it be a form $T$, is (like the original $T$) the sum of a form $W$ and of earlier $T$'s than itself. Hence also every form $T$ is the sum of a form $W$, and of forms derived from the functions $X_1$, $X_2$, \ldots\ or, what is the same thing, every covariant whatever is of the form $W + P_X$.

\smallskip
\noindent\textbf{358.}\quad The proof that for a form $f$ of the order $n$ the number of covariants is finite, depends on the assumption that the number is finite for a form $f'$ of the next inferior order $n-1$: this being so, the number of the special covariants of $f$ will be finite; say these are $A_1$, $A_2$, $A_3$, \ldots\ ($f$ is itself one of the series, but we may separate it, and speak of the form $f$ and its special covariants): the forms $W$ are functions of the special covariants, and hence every covariant whatever of $f$ is of the form $F(A) + P_X$; but it requires still a long investigation to pass from this to the theorem of the existence of a finite number of forms $V$ such that every covariant whatever is $F(V)$. I pass this over,

\newpage
\noindent and consider only the covariants of a quantic of the order $n$; these are all of them of the form $W + P_X$.

\smallskip
\noindent\textbf{359.}\quad If in any covariant $W$ the variables be interchanged in any manner whatever, the result is still a covariant $W$; and if from the same or different covariants $W$ we form any linear function with numerical coefficients, the result is still a covariant $W$; but if we form from a covariant $W$ a symbolical product involving any new symbolical coefficients, the result is no longer a covariant $W$, but is (as already mentioned) a sum of a covariant $W$ and of earlier covariants $T$.

\smallskip
\noindent\textbf{360.}\quad It may be remarked that in the case of a single variable, or say of a binary quantic, the only possible index is $i=1$; this is $< \frac{1}{2}n$ except in the case $n=1$; and the covariants $W$ are consequently nonexistent except for the case in question $n=1$; but for $n=1$, they are in fact the only covariants; viz.\ $f=(ax+by)$ is the only covariant, and this is a $W$.

\medskip
\noindent\textbf{The Binary Quantics.}

\smallskip
\noindent\textbf{361.}\quad For the binary quantics the foregoing theorem may be stated as follows:\[6pt]
\textbf{Theorem.}\ \textit{Every covariant of a binary quantic of the order $n$ is expressible as a rational and integral function of the special covariants.} And observing that every rational and integral function of the special covariants may, in virtue of the linear relations which connect them, be expressed as a rational and integral function of a certain finite number of these special covariants, the theorem may be further stated in the form:\[6pt]
\textbf{Theorem.}\ \textit{Every covariant of a binary quantic of the order $n$ is expressible as a rational and integral function of a certain finite number of covariants.}

\smallskip
\noindent\textbf{362.}\quad For the binary quantic of the order $n$, writing as usual
\[
(a, b, \ldots \mathbin{\raisebox{-2pt}{$\overset{\displaystyle\curvearrowright}{\scriptstyle n}$}} \!\! l\,x, y)^n = (ax+by)^n = \ldots,
\]
the number of asympetic covariants is as already mentioned equal to the coefficient of $x^\alpha a^\rho$ in
\[
\frac{1}{1-ax} \cdot \frac{1}{(1-a^2x^2)(1-ax^2)} \cdots \frac{1}{(1-a^nx^n)},
\]
or, what is the same thing, equal to the coefficient of $a^\rho x^\sigma$ in
\[
\frac{1}{1-ax} \cdot \frac{1}{(1-a^2x^2)(1-ax^2)} \cdots \frac{1}{(1-a^nx^n)},
\]
where $\rho$ is the degree and $\sigma$ the order.

\smallskip
\noindent\textbf{363.}\quad For $U$ equal to any one of the asympetic covariants, the degree $\rho$ and the order $\sigma$ satisfy the condition $\rho n - 2\sigma = 0$ or is a positive integer; or, what is the same thing, $\rho n$ and $2\sigma$ are equal or differ by a multiple of $n$, or, again, $\rho$ and $2\sigma$ are equal or differ by a multiple of 2; or, finally, $\rho$ and $\sigma$ are both of them even, or one of them is even and the other odd. In particular, if $\rho$ is even, then $\sigma$ is even; and if $\rho$ is odd, then $\sigma$ is odd. Hence, also, the order of a covariant of an odd degree is odd; and the order of a covariant of an even degree is even.

\smallskip
\noindent\textbf{364.}\quad For the binary quantic of the order $n$, the complete asympetic system of covariants includes the form itself, and it is convenient to speak of the form as a covariant; but in regard to the covariants of a binary quantic, the term is frequently restricted to mean a covariant other than the form itself; and it is in this restricted sense that I shall in general employ the term.

\smallskip
\noindent\textbf{365.}\quad For the several values of $n$, the asympetic covariants (including the form itself) are shown in the following table:\[6pt]
\begin{center}
\begin{tabular}{c|l}
\toprule
\textbf{Order $n$} & \textbf{Asympetic Covariants} \\
\midrule
$n=2$ & $f$ \\
$n=3$ & $f$, $H$, $\Phi$ \\
$n=4$ & $f$, $H$, $i$, $T$ \\
$n=5$ & $f$, $H$, $i$, $T$, $A$, $B$, $C$ \\
$n=6$ & $f$, $H$, $i$, $T$, $A$, $B$, $C$, $D$, $E$ \\
\bottomrule
\end{tabular}
\end{center}

\newpage
\noindent\textbf{366.}\quad For the binary quantic of the order $n$, the number of asympetic covariants is equal to the number of partitions of any number into parts not exceeding $n$; or, what is the same thing, it is equal to the coefficient of $x^\alpha$ in
\[
\frac{1}{(1-x)(1-x^2)\cdots(1-x^n)}.
\]
For $n=2$, this is
\[
\frac{1}{(1-x)(1-x^2)} = 1 + x + 2x^2 + 2x^3 + 3x^4 + \ldots,
\]
where the coefficient of $x^\alpha$ gives the number of asympetic covariants of the order $\alpha$.

\medskip
\noindent\textbf{367.}\quad The generating function for the number of asympetic covariants of the binary quantic of the order $n$ is
\[
\frac{1}{1-ax} \cdot \frac{1}{(1-a^2x^2)(1-ax^2)} \cdots \frac{1}{(1-a^nx^n)},
\]
and the coefficient of $a^\rho x^\sigma$ in the development of this function gives the number of asympetic covariants of the degree $\rho$ and order $\sigma$.

\smallskip
\noindent\textbf{368.}\quad It is to be observed that the asympetic covariants are not in general independent; there exist syzygies, that is, linear relations with numerical coefficients between the products of the asympetic covariants. For example, for the quartic, there is a syzygy connecting $f^3$, $f^2H$, $fH^2$, $H^3$, $i^2$, $iT$, $T^2$; and for the sextic, there are syzygies connecting the products of the asympetic covariants in a more complicated manner.

\smallskip
\noindent\textbf{369.}\quad The question of the independence of the asympetic coviants, or, more generally, the question of the syzygies which exist between the products of the asympetic covariants, is one of great difficulty; and it has been completely solved only for the lower values of $n$. For $n=2$, there are of course no syzygies; for $n=3$, there is one syzygy of degree 12 and order 12; for $n=4$, there are syzygies of degrees 6 and 8; and so on.

\medskip
\noindent\textbf{370.}\quad The theory of the binary quantics may be extended to ternary and higher quantics; but the theory becomes much more complicated, and many of the results which have been established for binary quantics are still wanting for the higher quantics.

\newpage
%% ========================================================================
%% PAPER 463: NOTE ON A DIFFERENTIAL EQUATION
%% ========================================================================
\section*{463. Note on a Differential Equation.}
\addcontentsline{toc}{section}{463. Note on a Differential Equation (pp.~354--356)}

\noindent\textit{[From the Quarterly Journal of Pure and Applied Mathematics, vol.~VIII. (1867), pp.~120--123.]}

\medskip

\noindent\textbf{1.}\quad The differential equation
\[
\frac{d^ny}{dx^n} = x^m y,
\]
where $m$ is a positive integer, has been considered by various authors; and the solution has been expressed in terms of definite integrals, or by means of series. The equation is of considerable importance in Analysis; and it presents itself in various physical problems.

\smallskip
\noindent\textbf{2.}\quad If we assume
\[
y = \int e^{xt} \phi(t) \, dt,
\]
then substituting in the differential equation, we have
\[
\int t^n e^{xt} \phi(t) \, dt = x^m \int e^{xt} \phi(t) \, dt.
\]
Integrating by parts the right-hand side, we obtain
\[
\int t^n e^{xt} \phi(t) \, dt = \int e^{xt} \left(-\frac{d}{dt}\right)^m \phi(t) \, dt;
\]
and this will be satisfied if
\[
t^n \phi(t) = \left(-\frac{d}{dt}\right)^m \phi(t),
\]
or, what is the same thing,
\[
\frac{d^m \phi}{dt^m} = (-1)^m t^n \phi.
\]

\smallskip
\noindent\textbf{3.}\quad The equation for $\phi$ is of the same form as the original equation for $y$, but with the indices $m$ and $n$ interchanged; and the solution of the one equation may be made to depend on the solution of the other. In particular, if $m=1$, the equation for $\phi$ becomes
\[
\frac{d\phi}{dt} = -t^n \phi,
\]
whence
\[
\phi = C e^{-\frac{t^{n+1}}{n+1}}.
\]

\smallskip
\noindent\textbf{4.}\quad Substituting this value of $\phi$, we have
\[
y = C \int e^{xt - \frac{t^{n+1}}{n+1}} \, dt,
\]
where the integral is taken along any path in the plane of the complex variable $t$, which renders the expression finite. The limits of integration may be so chosen that the integrated terms vanish; and the solution thus obtained will be the complete solution, containing $n$ arbitrary constants.

\newpage
\noindent\textbf{5.}\quad For example, if $n=2$, we have
\[
y = C \int e^{xt - \frac{t^3}{3}} \, dt.
\]
The integral may be evaluated by taking the path of integration to be a contour enclosing the origin, or by taking it from $-\infty$ to $\infty$ along the real axis; and the two independent solutions thus obtained will be the two independent integrals of the differential equation $\dfrac{d^2y}{dx^2} = xy$.

\smallskip
\noindent\textbf{6.}\quad The differential equation $\dfrac{d^2y}{dx^2} = xy$ is a particular case of the equation of the parabolic cylinder; and its solutions are connected with the Airy integral and the Bessel functions of order $\pm\frac{1}{3}$. The solutions may be expressed in the form
\[
y = x^{\frac{1}{2}} \left\{ A J_{\frac{1}{3}}\left(\frac{2}{3} x^{\frac{3}{2}}\right) + B J_{-\frac{1}{3}}\left(\frac{2}{3} x^{\frac{3}{2}}\right) \right\},
\]
where $A$ and $B$ are arbitrary constants.

\smallskip
\noindent\textbf{7.}\quad The general equation $\dfrac{d^ny}{dx^n} = x^m y$ may be treated in a similar manner; and the solution may be expressed by means of definite integrals, or by series involving the gamma function. The theory presents many interesting analytical developments; but the complete discussion of the equation would occupy too much space for the present note.

\newpage
%% ========================================================================
%% PAPER 464: NOTE ON PLANA'S LUNAR THEORY
%% ========================================================================
\section*{464. Note on Plana's Lunar Theory.}
\addcontentsline{toc}{section}{464. Note on Plana's Lunar Theory (pp.~357--360)}

\noindent\textit{[From the Monthly Notices of the Royal Astronomical Society, vol.~XXV. (1864--1865), pp.~225--229.]}

\medskip

\noindent The lunar theory of Plana is founded upon a development of the disturbing function in terms of the elements of the lunar orbit; and the expressions for the coordinates of the Moon are obtained by a process of successive approximation. The theory has been carried to a very high degree of accuracy; and the results have been compared with observation with satisfactory agreement.

\smallskip

\noindent In the development of the disturbing function, Plana introduces certain quantities which are functions of the eccentricities and inclinations of the lunar and solar orbits; and these quantities are expanded in series proceeding according to powers of the eccentricities and inclinations. The coefficients of the several terms in these series are expressed by means of the Besselian functions; and the general terms of the series may be written down by means of a known formula.

\smallskip

\noindent The lunar coordinates are expressed in terms of the mean longitude, the mean anomaly, and the mean elongation; and the expressions involve the time explicitly, through the mean motions of the Moon and the Sun. The theory takes account of the secular variations of the lunar elements; and the long-period inequalities are included in the development.

\smallskip

\noindent In the application of the theory to the calculation of the lunar coordinates, it is necessary to substitute numerical values for the elements of the orbit; and the accuracy of the results depends upon the accuracy with which these elements are known. The comparison of the theory with observation affords a means of determining the elements with greater precision; and the corrections thus obtained are applied to the elements in the subsequent calculations.

\newpage
\noindent The disturbing function in Plana's theory is developed in the following manner. Let $r$, $r'$ be the radii vectores of the Moon and Sun respectively; let $v$, $v'$ be their true longitudes; and let $R$ be the distance between them. Then the disturbing function is
\[
\Omega = \frac{m'}{\sqrt{r^2 + r'^2 - 2rr' \cos(v-v')}} - \frac{m'r}{r'^2} \cos(v-v'),
\]
where $m'$ is the mass of the Sun. The first term represents the direct attraction of the Sun upon the Moon; and the second term represents the effect of the Sun's attraction upon the Earth.

\smallskip

\noindent The function $\Omega$ is expanded in a series of cosines of multiples of the argument $v-v'$; and the coefficients of these cosines are functions of the ratio $r/r'$. Since this ratio is small (its mean value being about $\frac{1}{400}$), the coefficients decrease rapidly; and the series is therefore convergent.

\smallskip

\noindent The expansion of $\Omega$ is effected by means of the formula
\[
\frac{1}{\sqrt{1 - 2\alpha \cos\theta + \alpha^2}} = \frac{1}{2} \sum_{n=-\infty}^{\infty} P_n(\cos\theta) \alpha^{|n|},
\]
where $P_n$ is the Legendre's coefficient of order $n$; or, what is the same thing,
\[
\frac{1}{\sqrt{1 - 2\alpha \cos\theta + \alpha^2}} = \sum_{n=0}^{\infty} P_n(\cos\theta) \alpha^n,
\]
for $\alpha < 1$.

\smallskip

\noindent The development of the disturbing function in Plana's theory is one of the most elaborate analytical works which have ever been executed; and it is remarkable for the elegance and completeness of the results. The theory has been verified by comparison with observation; and it forms the foundation of the lunar tables which are employed in the nautical almanacs.

\newpage
%% ========================================================================
%% PAPER 465: NOTE ON THE LUNAR THEORY
%% ========================================================================
\section*{465. Note on the Lunar Theory.}
\addcontentsline{toc}{section}{465. Note on the Lunar Theory (pp.~361--366)}

\noindent\textit{[From the Monthly Notices of the Royal Astronomical Society, vol.~XXV. (1864--1865), pp.~229--237.]}

\medskip

\noindent The lunar theory, as developed by Laplace, Damoiseau, Plana, Hansen, and other investigators, is founded upon the general principles of the Newtonian theory of gravitation; and the method of solution consists in developing the coordinates of the Moon in series proceeding according to powers of the disturbing force. The disturbing force being small (its effect upon the lunar motion being of the order of the ratio of the Sun's parallax to the Moon's parallax, or about $\frac{1}{400}$), the series converge rapidly; and a comparatively small number of terms suffices to give the coordinates with great accuracy.

\smallskip

\noindent The method of the variation of arbitrary constants, which is employed in the lunar theory, consists in assuming for the coordinates expressions of the same form as in the undisturbed elliptic motion, but with the elements of the orbit regarded as variable quantities. These variable elements are determined by differential equations, which express the condition that the coordinates shall satisfy the equations of motion; and the solution of these differential equations gives the variations of the elements produced by the disturbing force.

\smallskip

\noindent The lunar theory presents many interesting and difficult problems; among which may be mentioned the determination of the secular acceleration of the Moon's mean motion, the theory of the long-period inequalities, and the effect of the figure of the Earth upon the lunar motion. These problems have been treated by the great geometers of the last century and the present; and the results have been verified by comparison with the observations of the Moon, which extend over a period of more than two thousand years.

\smallskip

\noindent The following formulae relate to the expressions for the longitude, latitude, and parallax of the Moon in terms of the elements of the orbit. Let $l$ be the mean longitude, $g$ the mean anomaly, $e$ the eccentricity, $\gamma$ the tangent of the inclination, $c$ the mean distance of the node from the mean equinox, and $m$ the ratio of the mean motion of the Sun to that of the Moon. Then the longitude of the Moon may be expressed in the form
\[
v = l + P_1 \sin g + P_2 \sin 2g + P_3 \sin 3g + \ldots,
\]
where $P_1$, $P_2$, $P_3$, \ldots\ are functions of $e$, $\gamma$, $m$, and the other elements.

\newpage
\noindent\textbf{1.}\quad The expressions for the longitude and latitude of the Moon may be written in the following form. Let $\lambda$ be the longitude, $\beta$ the latitude, and $r$ the radius vector; then
\begin{align*}
\lambda &= l + A_1 \sin g + A_2 \sin 2g + \ldots + B_1 \sin(2l-2l') + B_2 \sin(4l-4l') + \ldots \\
&\quad + C_1 \sin(g + 2l - 2l') + C_2 \sin(g - 2l + 2l') + \ldots,\\
\beta &= \gamma \sin(l - \Omega) + D_1 \sin(3l - 2l' - \Omega) + D_2 \sin(l - 2l' + \Omega) + \ldots,
\end{align*}
where $l$ is the mean longitude of the Moon, $l'$ that of the Sun, $g$ the mean anomaly, $\Omega$ the longitude of the node, and $A_1$, $A_2$, \ldots, $B_1$, $B_2$, \ldots\ are coefficients depending on the elements of the orbit.

\smallskip

\noindent\textbf{2.}\quad The coefficients $A_1$, $A_2$, \ldots\ are functions of the eccentricity $e$; and they may be expanded in series proceeding according to powers of $e$. The general term of these series may be expressed by means of the Besselian functions; and the coefficients may be calculated to any required degree of accuracy.

\smallskip

\noindent\textbf{3.}\quad The coefficients $B_1$, $B_2$, \ldots\ relate to the variation and the parallactic inequality; and they are functions of the ratio $m$ of the mean motions. These coefficients have been calculated to a very high degree of accuracy; and the values have been verified by comparison with observation.

\smallskip

\noindent\textbf{4.}\quad The coefficients $C_1$, $C_2$, \ldots\ relate to the evection; and they are functions both of $e$ and of $m$. The evection is the largest of the lunar inequalities; and its theory has been treated with great care by all the investigators of the lunar theory.

\smallskip

\noindent\textbf{5.}\quad The coefficients $D_1$, $D_2$, \ldots\ relate to the reduction to the ecliptic; and they are functions of the inclination and of the ratio $m$. The reduction to the ecliptic is necessary because the latitude of the Moon is measured from the ecliptic, while the elements of the orbit are referred to the mean equator.

\newpage
\noindent\textbf{6.}\quad The expressions for the longitude and latitude involve also terms depending on the solar eccentricity $e'$, and on the action of the planets. These terms are smaller than those depending on the lunar eccentricity; but they are sensible, and they must be taken into account in the construction of accurate lunar tables.

\smallskip

\noindent\textbf{7.}\quad The differential equations which determine the variations of the elements of the lunar orbit are of the form
\begin{align*}
\frac{de}{dt} &= -\frac{a\sqrt{1-e^2}}{\mu e} \frac{\partial R}{\partial \varpi} + \ldots,\\
\frac{d\varpi}{dt} &= \frac{a\sqrt{1-e^2}}{\mu e} \frac{\partial R}{\partial e} + \ldots,
\end{align*}
where $R$ is the disturbing function, $a$ the semi-axis major, $\mu$ the sum of the masses of the Earth and Moon, and $\varpi$ the longitude of the perigee. Similar equations hold for the other elements of the orbit.

\smallskip

\noindent\textbf{8.}\quad The integration of these differential equations gives the secular and periodic variations of the elements. The secular variations are those which increase continually with the time; and the periodic variations are those which oscillate between fixed limits. The secular variations of the eccentricity and inclination are very small; but the secular variation of the mean motion is sensible, and it gives rise to the well-known acceleration of the Moon's mean motion.

\smallskip

\noindent\textbf{9.}\quad The theory of the acceleration of the Moon's mean motion has been the subject of much discussion; and the question whether the observed acceleration can be entirely accounted for by the gravitational action of the Sun has been investigated by many eminent mathematicians. It appears from the most recent researches that the gravitational theory accounts for the greater part of the observed acceleration; but that there remains a small residual which may be due to the tidal friction, or to other causes not yet fully understood.

\newpage
%% ========================================================================
%% PAPER 466: SECOND NOTE ON THE LUNAR THEORY
%% ========================================================================
\section*{466. Second Note on the Lunar Theory.}
\addcontentsline{toc}{section}{466. Second Note on the Lunar Theory (pp.~367--370)}

\noindent\textit{[From the Monthly Notices of the Royal Astronomical Society, vol.~XXV. (1864--1865), pp.~237--245.]}

\medskip

\noindent\textbf{1.}\quad In my former note on the lunar theory, I gave the general expressions for the longitude, latitude, and parallax of the Moon in terms of the elements of the orbit. I now propose to consider more particularly the expressions for the longitude, and to investigate the coefficients of the several inequalities up to the seventh order inclusive.

\smallskip

\noindent\textbf{2.}\quad The longitude of the Moon may be expressed in the form
\[
\lambda = l + \delta\lambda,
\]
where $l$ is the mean longitude, and $\delta\lambda$ is the sum of the several inequalities. The principal inequalities are the following:\[6pt]
\begin{tabular}{ll}
\quad 1.\ The Equation of the centre & $2e \sin g + \frac{5}{4}e^2 \sin 2g + \ldots$ \\
\quad 2.\ The Evection & $-\frac{15}{8}m e \sin(g - 2l + 2l') + \ldots$ \\
\quad 3.\ The Variation & $\frac{15}{4}m \sin(2l - 2l') + \ldots$ \\
\quad 4.\ The Parallactic inequality & $-\frac{15}{8}m^2 \sin(2l - 2l') + \ldots$ \\
\quad 5.\ The Annual equation & $-e' \sin l' + \ldots$ \\
\end{tabular}

\smallskip

\noindent\textbf{3.}\quad The coefficients of these inequalities are functions of the ratio $m$ of the mean motions; and they may be expanded in series proceeding according to powers of $m$. The calculation of these coefficients is one of the most laborious parts of the lunar theory; and it has been carried to a very high degree of accuracy by the various investigators.

\smallskip

\noindent\textbf{4.}\quad The equation of the centre is the inequality depending on the eccentricity of the lunar orbit; and its coefficient is the largest of all the inequalities. The expression for the equation of the centre may be written in the form
\[
\delta\lambda_1 = 2e \sin g + \frac{5}{4}e^2 \sin 2g + \frac{13}{12}e^3 \sin 3g + \ldots
\]
where $g$ is the mean anomaly. The coefficients of this series are known functions of the eccentricity; and they may be calculated to any required degree of accuracy.

\newpage
\noindent\textbf{5.}\quad The evection is the inequality depending on the combined action of the eccentricity and the solar disturbance; and it is the second in magnitude among the lunar inequalities. The expression for the evection may be written in the form
\[
\delta\lambda_2 = -\frac{15}{8}m e \sin(g - 2l + 2l') + \frac{75}{32}m^2 e \sin(g - 4l + 4l') + \ldots
\]
The coefficient of the principal term is $-\frac{15}{8}me$; and the succeeding terms are of higher orders in $m$ and $e$.

\smallskip

\noindent\textbf{6.}\quad The variation is the inequality depending on the solar disturbance alone; and it is the third in magnitude among the lunar inequalities. The expression for the variation may be written in the form
\[
\delta\lambda_3 = \frac{15}{4}m \sin(2l - 2l') + \frac{105}{16}m^2 \sin(4l - 4l') + \ldots
\]
The coefficient of the principal term is $\frac{15}{4}m$; and the succeeding terms are of higher orders in $m$.

\smallskip

\noindent\textbf{7.}\quad The parallactic inequality is the inequality depending on the solar parallax; and it is smaller than the variation. The expression for the parallactic inequality may be written in the form
\[
\delta\lambda_4 = -\frac{15}{8}m^2 \sin(2l - 2l') - \frac{105}{32}m^3 \sin(4l - 4l') + \ldots
\]
The coefficient of the principal term is $-\frac{15}{8}m^2$; and the succeeding terms are of higher orders in $m$.

\smallskip

\noindent\textbf{8.}\quad The annual equation is the inequality depending on the eccentricity of the Earth's orbit; and it is smaller than the parallactic inequality. The expression for the annual equation may be written in the form
\[
\delta\lambda_5 = -e' \sin l' - \frac{1}{2}e'^2 \sin 2l' + \ldots
\]
where $e'$ is the solar eccentricity, and $l'$ is the mean anomaly of the Sun.

\newpage
%% ========================================================================
%% PAPER 467: EXPRESSIONS FOR PLANA'S e, γ IN TERMS OF THE ELLIPTIC e, γ
%% ========================================================================
\section*{467. Expressions for Plana's $e$, $\gamma$ in terms of the Elliptic $e$, $\gamma$.}
\addcontentsline{toc}{section}{467. Expressions for Plana's $e$, $\gamma$ in terms of the Elliptic $e$, $\gamma$ (pp.~371--374)}

\noindent\textit{[From the Monthly Notices of the Royal Astronomical Society, vol.~XXV. (1864--1865), pp.~245--252.]}

\medskip

\noindent\textbf{1.}\quad In Plana's lunar theory, the quantities $e$ and $\gamma$ are not the eccentricity and tangent of the inclination of the actual lunar orbit, but are certain ideal elements, which are connected with the true elements by relations of a somewhat complicated character. The object of the present note is to establish these relations, and to express Plana's $e$ and $\gamma$ in terms of the elliptic (that is, the true) $e$ and $\gamma$.

\smallskip

\noindent\textbf{2.}\quad The relations between Plana's elements and the true elements may be obtained by comparing the expressions for the coordinates in the two theories. In the elliptic theory, the coordinates are expressed in terms of the true eccentricity $e$ and inclination $\gamma$; while in Plana's theory, they are expressed in terms of Plana's $e$ and $\gamma$. By equating the two expressions for the same coordinate, we obtain relations between the two sets of elements.

\smallskip

\noindent\textbf{3.}\quad The comparison gives the following equations:\begin{align*}
e_P &= e - \frac{1}{8}m^2 e + \frac{1}{4}m^2 e' \cos(\varpi - \varpi') + \ldots,\\
\gamma_P &= \gamma - \frac{1}{8}m^2 \gamma + \frac{3}{8}m^2 \gamma' \cos(\Omega - \Omega') + \ldots,
\end{align*}
where $e_P$, $\gamma_P$ are Plana's elements, and $e$, $\gamma$ are the elliptic elements; $e'$, $\gamma'$ are the solar eccentricity and the inclination of the solar orbit; $\varpi$, $\varpi'$ are the longitudes of the perigees; and $\Omega$, $\Omega'$ are the longitudes of the nodes.

\smallskip

\noindent\textbf{4.}\quad These equations show that Plana's $e$ differs from the true eccentricity by quantities of the order $m^2$; and that Plana's $\gamma$ differs from the true inclination by quantities of the same order. The differences are therefore small; but they are sensible, and they must be taken into account in the comparison of the theory with observation.

\newpage
\noindent\textbf{5.}\quad The relations between Plana's elements and the true elements may also be obtained by a direct transformation of the differential equations. If we write the equations of motion in terms of the true elements, and then transform to Plana's elements by means of the above formulae, we obtain the equations of motion in the form given by Plana.

\smallskip

\noindent\textbf{6.}\quad The inverse relations, expressing the true elements in terms of Plana's, are
\begin{align*}
e &= e_P + \frac{1}{8}m^2 e_P - \frac{1}{4}m^2 e' \cos(\varpi - \varpi') + \ldots,\\
\gamma &= \gamma_P + \frac{1}{8}m^2 \gamma_P - \frac{3}{8}m^2 \gamma' \cos(\Omega - \Omega') + \ldots
\end{align*}
These equations are obtained by solving the former equations with respect to $e$ and $\gamma$; and they are of the same degree of accuracy as the former.

\smallskip

\noindent\textbf{7.}\quad The corrections to Plana's elements, depending on the solar eccentricity and the inclination of the solar orbit, are of considerable importance in the lunar theory. They give rise to inequalities in the lunar motion, which are sensible in the comparison of the theory with observation; and they must be taken into account in the construction of accurate lunar tables.

\smallskip

\noindent\textbf{8.}\quad The foregoing results may be extended to the other elements of the lunar orbit; and similar relations may be established between the other elements of Plana's theory and the corresponding true elements. The complete investigation would occupy too much space for the present note; but the method is the same as that which has been employed for the eccentricity and the inclination.

\newpage
%% ========================================================================
%% PAPER 468: ADDITION TO SECOND NOTE ON THE LUNAR THEORY
%% ========================================================================
\section*{468. Addition to Second Note on the Lunar Theory.}
\addcontentsline{toc}{section}{468. Addition to Second Note on the Lunar Theory (pp.~375--376)}

\noindent\textit{[From the Monthly Notices of the Royal Astronomical Society, vol.~XXV. (1864--1865), pp.~252--254.]}

\medskip

\noindent\textbf{1.}\quad In my Second Note on the Lunar Theory, I gave the expressions for the principal inequalities of the Moon's longitude, up to terms of the seventh order inclusive. I now wish to add some remarks on the method of calculating the higher inequalities, and on the degree of accuracy which may be obtained by the theory.

\smallskip

\noindent\textbf{2.}\quad The method of successive approximation, which is employed in the lunar theory, consists in assuming for the coordinates expressions of the form
\[
x = x_0 + x_1 + x_2 + \ldots,
\]
where $x_0$ is the value of the coordinate in the undisturbed elliptic motion, and $x_1$, $x_2$, \ldots\ are the corrections depending on the disturbing force to the first, second, \ldots\ orders. The quantities $x_1$, $x_2$, \ldots\ are determined by differential equations, which are obtained by substituting the assumed expressions in the equations of motion and equating the terms of each order separately.

\smallskip

\noindent\textbf{3.}\quad The calculation of the higher inequalities is attended with considerable labour, on account of the great number of terms which must be taken into account. The number of inequalities of a given order increases rapidly with the order; and the calculation of the coefficients of these inequalities requires the solution of a large number of differential equations. The labour of the calculation has been greatly diminished by the introduction of mechanical methods; but it is still very considerable.

\smallskip

\noindent\textbf{4.}\quad The degree of accuracy which may be obtained by the lunar theory is very high. The coefficients of the principal inequalities are known with great precision; and the errors of the theory are chiefly due to the uncertainty of the elements of the orbit, and to the neglect of the higher inequalities. By carrying the calculation to a sufficient order, the errors of the theory may be made smaller than the errors of observation; and the theory may thus be used to determine the elements of the orbit with the greatest possible accuracy.

\newpage
%% ========================================================================
%% PAPER 469: ON AN EXPRESSION FOR THE ANGULAR DISTANCE OF TWO PLANETS
%% ========================================================================
\section*{469. On an Expression for the Angular Distance of Two Planets.}
\addcontentsline{toc}{section}{469. On an Expression for the Angular Distance of Two Planets (pp.~377--379)}

\noindent\textit{[From the Monthly Notices of the Royal Astronomical Society, vol.~XXV. (1864--1865), pp.~254--258.]}

\medskip

\noindent\textbf{1.}\quad The angular distance of two planets, as seen from the Earth, is an important element in the theory of planetary perturbations; and it is also of use in the determination of the orbits of the planets from observations. The expression for the angular distance may be obtained from the spherical triangle formed by the two planets and the Earth; and it may be developed in a series proceeding according to powers of the ratios of the distances to the radius of the sphere of observation.

\smallskip

\noindent\textbf{2.}\quad Let $r$, $r'$ be the distances of the two planets from the Earth; let $\lambda$, $\lambda'$ be their longitudes; and let $\beta$, $\beta'$ be their latitudes. Then the angular distance $z$ is given by the formula
\[
\cos z = \sin\beta \sin\beta' + \cos\beta \cos\beta' \cos(\lambda - \lambda').
\]
This formula is exact; but it is not convenient for calculation, on account of the number of trigonometrical functions which it involves.

\smallskip

\noindent\textbf{3.}\quad For the purpose of calculation, it is more convenient to have an expression for $\cos z$ in terms of the rectangular coordinates of the planets. If $x$, $y$, $z$ and $x'$, $y'$, $z'$ are the coordinates of the two planets referred to the ecliptic, then
\[
rr' \cos z = xx' + yy' + zz'.
\]
This expression may be developed in a series by substituting for the coordinates their values in terms of the elements of the orbits.

\smallskip

\noindent\textbf{4.}\quad The development of the expression for $\cos z$ gives
\begin{align*}
\cos z &= \cos(\lambda - \lambda') \cos\beta \cos\beta' + \sin\beta \sin\beta' \\
&= \left(1 - \frac{1}{2}\beta^2 + \ldots\right)\left(1 - \frac{1}{2}\beta'^2 + \ldots\right) \cos(\lambda - \lambda') + \beta\beta' + \ldots
\end{align*}
where $\beta$, $\beta'$ are expressed in parts of the radius. This development is convenient when the latitudes are small, as is usually the case with the planets.

\newpage
\noindent\textbf{5.}\quad The expression for the angular distance may also be obtained by means of the theory of spherical coordinates. If $\theta$, $\theta'$ are the angular distances of the planets from a fixed point on the sphere, and $\phi$, $\phi'$ are their position-angles, then
\[
\cos z = \cos\theta \cos\theta' + \sin\theta \sin\theta' \cos(\phi - \phi').
\]
This formula is sometimes useful in the theory of the mutual perturbations of the planets; and it may be employed in the calculation of the disturbing function.

\smallskip

\noindent\textbf{6.}\quad In the theory of the mutual perturbations of the planets, the expression for the angular distance enters into the disturbing function; and it is necessary to develop this expression in a series of cosines of multiples of the mean longitudes. The development is effected by means of the formula
\[
\frac{1}{\Delta} = \frac{1}{\sqrt{r^2 + r'^2 - 2rr' \cos z}},
\]
where $\Delta$ is the distance between the two planets. This expression is expanded in a series of Legendre's coefficients; and the coefficients are then developed in series of cosines of multiples of the mean longitudes.

\smallskip

\noindent\textbf{7.}\quad The calculation of the coefficients of the series for $\dfrac{1}{\Delta}$ is one of the most laborious parts of the theory of planetary perturbations. The coefficients are functions of the ratios of the distances and of the eccentricities and inclinations of the orbits; and they must be calculated to a high degree of accuracy in order to obtain satisfactory results. The labour of the calculation has been greatly diminished by the introduction of mechanical quadratures and other numerical methods.

\newpage
%% ========================================================================
%% PAPER 470: NOTE ON THE ATTRACTION OF ELLIPSOIDS
%% ========================================================================
\section*{470. Note on the Attraction of Ellipsoids.}
\addcontentsline{toc}{section}{470. Note on the Attraction of Ellipsoids (pp.~380--383)}

\noindent\textit{[From the Monthly Notices of the Royal Astronomical Society, vol.~XXV. (1864--1865), pp.~258--264.]}

\medskip

\noindent\textbf{1.}\quad The attraction of an ellipsoid upon an external point is a problem which has occupied the attention of mathematicians since the time of Newton. The problem was first completely solved by Ivory, who showed that the attraction of an ellipsoid upon an external point is the same as the attraction of a certain sphere upon the same point; and this result has been extended by Chasles and others to the case of heterogeneous ellipsoids.

\smallskip

\noindent\textbf{2.}\quad The method of Ivory consists in comparing the attraction of the given ellipsoid with that of a confocal ellipsoid passing through the attracted point. If $a$, $b$, $c$ are the semi-axes of the given ellipsoid, and $a'$, $b'$, $c'$ are those of the confocal ellipsoid, then
\[
\frac{x^2}{a^2} + \frac{y^2}{b^2} + \frac{z^2}{c^2} = 1, \qquad
\frac{x^2}{a'^2} + \frac{y^2}{b'^2} + \frac{z^2}{c'^2} = 1,
\]
are the equations of the two ellipsoids; and the condition that they are confocal gives
\[
a'^2 - a^2 = b'^2 - b^2 = c'^2 - c^2 = \theta,
\]
say, where $\theta$ is determined by the condition that the second ellipsoid passes through the attracted point.

\smallskip

\noindent\textbf{3.}\quad The potentials of the two ellipsoids at the attracted point are
\[
V = \pi\rho abc \int_0^\infty \frac{ds}{\sqrt{(a^2+s)(b^2+s)(c^2+s)}},
\]
\[
V' = \pi\rho a'b'c' \int_0^\infty \frac{ds}{\sqrt{(a'^2+s)(b'^2+s)(c'^2+s)}};
\]
and Ivory's theorem asserts that the attractions of the two ellipsoids at the attracted point are in the ratio of the masses of the ellipsoids. That is, if $X$, $Y$, $Z$ are the components of the attraction of the first ellipsoid, and $X'$, $Y'$, $Z'$ are those of the second, then
\[
\frac{X}{X'} = \frac{Y}{Y'} = \frac{Z}{Z'} = \frac{abc}{a'b'c'}.
\]

\smallskip

\noindent\textbf{4.}\quad This theorem may be proved by direct integration. The potential of a homogeneous ellipsoid at an external point $(x,y,z)$ is
\[
V = \pi\rho abc \int_\lambda^\infty \left(1 - \frac{x^2}{a^2+s} - \frac{y^2}{b^2+s} - \frac{z^2}{c^2+s}\right) \frac{ds}{\sqrt{(a^2+s)(b^2+s)(c^2+s)}},
\]
where $\lambda$ is the greatest root of the equation
\[
\frac{x^2}{a^2+\lambda} + \frac{y^2}{b^2+\lambda} + \frac{z^2}{c^2+\lambda} = 1.
\]

\newpage
\noindent\textbf{5.}\quad The components of the attraction are obtained by differentiating the potential with respect to $x$, $y$, $z$; and we find
\begin{align*}
X &= -\frac{\partial V}{\partial x} = 2\pi\rho abc \cdot x \int_\lambda^\infty \frac{ds}{(a^2+s)\sqrt{(a^2+s)(b^2+s)(c^2+s)}},\\
Y &= -\frac{\partial V}{\partial y} = 2\pi\rho abc \cdot y \int_\lambda^\infty \frac{ds}{(b^2+s)\sqrt{(a^2+s)(b^2+s)(c^2+s)}},\\
Z &= -\frac{\partial V}{\partial z} = 2\pi\rho abc \cdot z \int_\lambda^\infty \frac{ds}{(c^2+s)\sqrt{(a^2+s)(b^2+s)(c^2+s)}}.
\end{align*}
These expressions give the attraction of a homogeneous ellipsoid upon an external point in terms of definite integrals.

\smallskip

\noindent\textbf{6.}\quad For a point on the surface of the ellipsoid, we have $\lambda=0$; and the expressions become
\begin{align*}
X &= 2\pi\rho abc \cdot x \int_0^\infty \frac{ds}{(a^2+s)\sqrt{(a^2+s)(b^2+s)(c^2+s)}},\\
Y &= 2\pi\rho abc \cdot y \int_0^\infty \frac{ds}{(b^2+s)\sqrt{(a^2+s)(b^2+s)(c^2+s)}},\\
Z &= 2\pi\rho abc \cdot z \int_0^\infty \frac{ds}{(c^2+s)\sqrt{(a^2+s)(b^2+s)(c^2+s)}}.
\end{align*}
These integrals may be evaluated in terms of elliptic functions; and the result is the well-known formula for the attraction of an ellipsoid upon a point on its surface.

\smallskip

\noindent\textbf{7.}\quad The attraction of a heterogeneous ellipsoid, in which the density is a function of the distance from the centre, may be obtained by integrating the attraction of a series of concentric homogeneous ellipsoids. If the density is $\rho = f(\sqrt{x^2+y^2+z^2})$, then the potential is
\[
V = \int_0^1 \pi f(ar) \, r^2 dr \int_0^\infty \frac{ds}{\sqrt{(a^2r^2+s)(b^2r^2+s)(c^2r^2+s)}},
\]
and the components of the attraction are obtained by differentiating this expression.

\newpage
%% ========================================================================
%% PAPER 471: ON THE DETERMINATION OF A PLANET'S ORBIT
%% ========================================================================
\section*{471. On the Determination of a Planet's Orbit.}
\addcontentsline{toc}{section}{471. On the Determination of a Planet's Orbit (pp.~384--386)}

\noindent\textit{[From the Monthly Notices of the Royal Astronomical Society, vol.~XXV. (1864--1865), pp.~264--268.]}

\medskip

\noindent\textbf{1.}\quad The determination of the orbit of a planet from three observations is one of the fundamental problems of theoretical astronomy; and it has been treated by many eminent mathematicians, from the time of Gauss to the present day. The problem consists in finding the elements of the orbit, that is, the semi-axis major, the eccentricity, the inclination, the longitude of the node, the longitude of the perigee, and the mean motion, from three observed positions of the planet.

\smallskip

\noindent\textbf{2.}\quad The method of Gauss, which is the most celebrated of all the methods for the solution of this problem, consists in expressing the ratios of the triangles formed by the three observed positions of the planet and the Sun in terms of the elements of the orbit; and then determining the elements by means of these ratios. The method is elegant and complete; and it has been the foundation of all subsequent investigations on the subject.

\smallskip

\noindent\textbf{3.}\quad The fundamental equations of the method of Gauss are the following. Let $r_1$, $r_2$, $r_3$ be the radii vectores of the planet at the three epochs; let $2f_1$, $2f_2$, $2f_3$ be the areas of the triangles formed by the radii vectores; and let $\tau_1$, $\tau_2$, $\tau_3$ be the intervals of time between the observations. Then
\[
\frac{f_1}{\tau_1} = \frac{\sqrt{\mu p}}{2}, \qquad
\frac{f_2}{\tau_2} = \frac{\sqrt{\mu p}}{2}, \qquad
\frac{f_3}{\tau_3} = \frac{\sqrt{\mu p}}{2},
\]
where $\mu$ is the sum of the masses of the Sun and planet, and $p$ is the semi-parameter of the orbit.

\smallskip

\noindent\textbf{4.}\quad The ratios of the areas of the triangles are known from the observations; and the ratios of the intervals of time are also known. The equations therefore give the value of $p$; and this value being substituted in the equations of the orbit, the elements may be determined.

\smallskip

\noindent\textbf{5.}\quad The method of Gauss has been modified and improved by many subsequent investigators; and various methods have been proposed for the direct solution of the equations, without the use of successive approximations. The most important of these methods are those of Encke, Leverrier, and others; and they have been employed with great success in the calculation of the orbits of the planets and comets.

\newpage
%% ========================================================================
%% PAPER 472: NOTE ON LAMBERT'S THEOREM FOR ELLIPTIC MOTION
%% ========================================================================
\section*{472. Note on Lambert's Theorem for Elliptic Motion.}
\addcontentsline{toc}{section}{472. Note on Lambert's Theorem for Elliptic Motion (pp.~387--389)}

\noindent\textit{[From the Monthly Notices of the Royal Astronomical Society, vol.~XXV. (1864--1865), pp.~268--272.]}

\medskip

\noindent\textbf{1.}\quad Lambert's theorem is one of the most remarkable results in the theory of elliptic motion; and it has been the subject of many investigations by eminent mathematicians. The theorem asserts that the time of describing an arc of an ellipse depends only on the sum of the radii vectores of the extremities of the arc, and on the chord of the arc; or, in other words, that elliptic arcs which have the same chord and the same sum of radii vectores are described in the same time.

\smallskip

\noindent\textbf{2.}\quad The analytical expression of Lambert's theorem is as follows. Let $r$, $r'$ be the radii vectores of the extremities of an arc of an ellipse; let $c$ be the chord of the arc; and let $t$ be the time of describing the arc. Then
\[
\sqrt{\mu}\, t = a^{\frac{3}{2}} \left\{ (\epsilon - \sin\epsilon) - (\delta - \sin\delta) \right\},
\]
where $\mu$ is the sum of the masses, $a$ is the semi-axis major, and $\epsilon$, $\delta$ are angles determined by the equations
\[
\sin\frac{\epsilon}{2} = \frac{1}{2}\sqrt{\frac{r+r'+c}{a}}, \qquad
\sin\frac{\delta}{2} = \frac{1}{2}\sqrt{\frac{r+r'-c}{a}}.
\]

\smallskip

\noindent\textbf{3.}\quad The proof of Lambert's theorem may be obtained by a direct integration of the differential equations of elliptic motion. The time of describing an arc is expressed by the integral
\[
t = \int \frac{r^2 \, dv}{\sqrt{\mu p}},
\]
where $v$ is the true anomaly and $p$ is the semi-parameter. By transforming this integral by means of the relations between the true anomaly, the eccentric anomaly, and the mean anomaly, the result may be brought to the form given by Lambert's theorem.

\smallskip

\noindent\textbf{4.}\quad Another proof of Lambert's theorem may be obtained by means of the theory of elliptic functions. The expression for the time involves the elliptic integral of the second kind; and this integral may be transformed by means of the addition-formulae of the elliptic functions. The result of this transformation is Lambert's theorem; and the proof is in many respects simpler and more elegant than the direct proof.

\newpage
\noindent\textbf{5.}\quad Lambert's theorem may be extended to hyperbolic motion. In this case, the time of describing an arc is expressed by means of hyperbolic functions; and the formula is
\[
\sqrt{\mu}\, t = (-a)^{\frac{3}{2}} \left\{ (\sinh\epsilon - \epsilon) - (\sinh\delta - \delta) \right\},
\]
where $\epsilon$, $\delta$ are determined by equations similar to those for the elliptic case, but with hyperbolic functions in place of circular functions.

\smallskip

\noindent\textbf{6.}\quad The theorem may also be extended to parabolic motion. In this case, the semi-axis major is infinite; and the formula becomes
\[
t = \frac{1}{6\sqrt{\mu}} \left\{ (r+r'+c)^{\frac{3}{2}} \pm (r+r'-c)^{\frac{3}{2}} \right\},
\]
which is Barker's theorem. This theorem is a particular case of Lambert's theorem; and it may be obtained from it by making the eccentricity equal to unity.

\smallskip

\noindent\textbf{7.}\quad Lambert's theorem has many important applications in the theory of planetary motion. It is used in the determination of the orbits of planets and comets from observations; and it is also of use in the calculation of the perturbations of the planets. The theorem is one of the most elegant results in the theory of elliptic motion; and it is remarkable for its generality and simplicity.

\newpage
%% ========================================================================
%% PAPER 473: ON THE CONSTRUCTION OF THE UMBRAL OR PENUMBRAL CURVE &C.
%% ========================================================================
\section*{473. On the Construction of the Umbral or Penumbral Curve \&c.}
\addcontentsline{toc}{section}{473. On the Construction of the Umbral or Penumbral Curve (pp.~390--391)}

\noindent\textit{[From the Monthly Notices of the Royal Astronomical Society, vol.~XXV. (1864--1865), pp.~272--274.]}

\medskip

\noindent\textbf{1.}\quad The umbral and penumbral curves are the boundaries of the regions on the surface of the Earth where the Sun is totally or partially eclipsed by the Moon. The construction of these curves is a problem of considerable importance in the theory of eclipses; and it has been treated by many eminent astronomers.

\smallskip

\noindent\textbf{2.}\quad The umbral curve is the locus of points on the surface of the Earth where the Sun is totally eclipsed; and the penumbral curve is the locus of points where the Sun is partially eclipsed. These curves are determined by the condition that the apparent distances of the centres of the Sun and Moon, as seen from a point on the Earth's surface, are equal to the sum or difference of their apparent semidiameters.

\smallskip

\noindent\textbf{3.}\quad The condition for the umbral curve is
\[
\Delta = s - s',
\]
where $\Delta$ is the apparent distance of the centres of the Sun and Moon, $s$ is the apparent semidiameter of the Sun, and $s'$ is the apparent semidiameter of the Moon. The condition for the penumbral curve is
\[
\Delta = s + s'.
\]

\smallskip

\noindent\textbf{4.}\quad The construction of these curves on the Earth's surface may be effected by the following method. Let $S$ and $M$ be the centres of the Sun and Moon; let $O$ be the centre of the Earth; and let $P$ be a point on the surface of the Earth. The apparent distance of the centres of the Sun and Moon, as seen from $P$, is determined by the spherical triangle $SPM$; and the condition for the umbral or penumbral curve gives an equation between the coordinates of $P$.

\smallskip

\noindent\textbf{5.}\quad The equation of the umbral or penumbral curve may be written in the form
\[
\cos SP \cos MP + \sin SP \sin MP \cos SPM = \cos(s \mp s').
\]
This equation determines the locus of the point $P$ on the surface of the Earth; and it may be solved by expressing the coordinates of $P$ in terms of the latitude and longitude.

\newpage
%% ========================================================================
%% PAPER 474: ON THE GEOMETRICAL THEORY OF SOLAR ECLIPSES
%% ========================================================================
\section*{474. On the Geometrical Theory of Solar Eclipses.}
\addcontentsline{toc}{section}{474. On the Geometrical Theory of Solar Eclipses (pp.~392--396)}

\noindent\textit{[From the Monthly Notices of the Royal Astronomical Society, vol.~XXV. (1864--1865), pp.~274--282.]}

\medskip

\noindent The geometrical theory of solar eclipses is founded upon the consideration of the cones of shadow cast by the Moon. These cones are determined by the apparent semidiameters of the Sun and Moon, and by their relative positions; and their intersections with the surface of the Earth give the regions of total and partial eclipse.

\smallskip

\noindent The cone of total shadow, or the umbra, is the cone whose vertex is between the Moon and the Earth, and whose generators touch the Sun and Moon externally. The cone of partial shadow, or the penumbra, is the cone whose vertex is on the side of the Moon remote from the Earth, and whose generators touch the Sun and Moon internally on one side and externally on the other.

\smallskip

\noindent The intersections of these cones with the surface of the Earth are curves of a complicated character; and their determination requires the solution of a system of equations involving the coordinates of the Sun, Moon, and Earth. The following investigation is founded upon a method which was first employed by Bessel, and which has been adopted by most subsequent writers on the subject.

\smallskip

\noindent Let $\xi$, $\eta$, $\zeta$ be the coordinates of the Moon relative to the centre of the Earth; let $R$ be the distance of the Moon from the Earth; and let $s$, $s'$ be the apparent semidiameters of the Sun and Moon. Then the equations of the cones of shadow are
\begin{align*}
\xi^2 + \eta^2 &= (\zeta \tan s' - R \sin s)^2 \quad \text{(umbra)},\\
\xi^2 + \eta^2 &= (\zeta \tan s' + R \sin s)^2 \quad \text{(penumbra)}.
\end{align*}

\smallskip

\noindent The intersections of these cones with the surface of the Earth are obtained by substituting for $\xi$, $\eta$, $\zeta$ their values in terms of the latitude and longitude of the point on the Earth's surface. The resulting equations are of a complicated character; but they may be solved by numerical methods.

\newpage
\noindent\textbf{1.}\quad The fundamental equations of the geometrical theory of eclipses are the following. Let $x$, $y$, $z$ be the coordinates of the Moon relative to the centre of the Earth; let $l$ be the ratio of the apparent semidiameter of the Moon to the apparent semidiameter of the Sun; and let $\sin f$, $\sin f'$ be the sines of the angles which the generators of the umbral and penumbral cones make with the line joining the centres of the Sun and Moon. Then the equations of the umbral and penumbral surfaces are
\begin{align*}
(x-x_0)^2 + (y-y_0)^2 &= (z-z_0)^2 \tan^2 f \quad \text{(umbra)},\\
(x-x_0)^2 + (y-y_0)^2 &= (z-z_0)^2 \tan^2 f' \quad \text{(penumbra)},
\end{align*}
where $x_0$, $y_0$, $z_0$ are the coordinates of the vertex of the cone.

\smallskip

\noindent\textbf{2.}\quad The coordinates of the vertex of the umbral cone are determined by the condition that the distances of the vertex from the centres of the Sun and Moon are in the ratio of their apparent semidiameters. If $D$ is the distance between the centres of the Sun and Moon, then
\[
\frac{z_0}{R'} = \frac{s'}{s-s'}, \qquad z_0 = \frac{D s'}{s-s'},
\]
where $R'$ is the distance of the Sun from the Earth. The coordinates $x_0$, $y_0$ are equal to the coordinates of the Moon multiplied by the ratio $\dfrac{z_0}{R}$.

\smallskip

\noindent\textbf{3.}\quad The intersections of the umbral and penumbral surfaces with the surface of the Earth are the curves of total and partial eclipse. These curves may be constructed by the following method. For a given time, calculate the coordinates of the Moon and the values of $s$, $s'$; then calculate the coordinates of the vertex of the cone; and finally, determine the intersection of the cone with the sphere of the Earth.

\smallskip

\noindent\textbf{4.}\quad The equation of the intersection of the cone with the sphere of the Earth may be written in the form
\[
\xi^2 + \eta^2 + \zeta^2 = R_\oplus^2, \qquad
\xi^2 + \eta^2 = (\zeta - z_0)^2 \tan^2 f,
\]
where $R_\oplus$ is the radius of the Earth, and $\xi$, $\eta$, $\zeta$ are the coordinates of a point on the curve. Eliminating $\xi^2 + \eta^2$, we obtain
\[
R_\oplus^2 - \zeta^2 = (\zeta - z_0)^2 \tan^2 f,
\]
which is a quadratic equation for $\zeta$.

\newpage
\noindent\textbf{5.}\quad The solution of this equation gives two values for $\zeta$; and each value corresponds to a point on the curve. The values of $\xi$ and $\eta$ are then determined by the equation of the cone; and the latitude and longitude of the point are obtained by transforming the rectangular coordinates to geographical coordinates.

\smallskip

\noindent\textbf{6.}\quad The curve of total eclipse is thus determined for each instant of time; and the path of the eclipse on the Earth's surface is the locus of these curves. The limits of the path are determined by the condition that the quadratic equation for $\zeta$ has equal roots; and this condition gives the equations of the limiting curves.

\smallskip

\noindent\textbf{7.}\quad The duration of the total eclipse at any point on the path is determined by the time during which the point remains within the umbral cone. This time may be calculated from the motion of the Moon and the rotation of the Earth; and it is a maximum at the centre of the path, and decreases to zero at the limits.

\smallskip

\noindent\textbf{8.}\quad The foregoing theory may be extended to the case of annular eclipses. In this case, the vertex of the umbral cone is between the Earth and the Moon; and the eclipse is annular when the apparent diameter of the Moon is less than that of the Sun. The condition for an annular eclipse is that the umbral cone intersects the surface of the Earth beyond the vertex; and the theory of annular eclipses is in all respects similar to that of total eclipses.

\smallskip

\noindent\textbf{9.}\quad The penumbral cone gives rise to partial eclipses; and the region of partial eclipse is much more extensive than the region of total eclipse. The limits of the partial eclipse are determined by the intersection of the penumbral cone with the surface of the Earth; and they may be calculated by the same method as the limits of the total eclipse.

\newpage
%% ========================================================================
%% PAPER 475: ON A PROPERTY OF THE STEREOGRAPHIC PROJECTION
%% ========================================================================
\section*{475. On a Property of the Stereographic Projection.}
\addcontentsline{toc}{section}{475. On a Property of the Stereographic Projection (pp.~397--399)}

\noindent\textit{[From the Monthly Notices of the Royal Astronomical Society, vol.~XXV. (1864--1865), pp.~282--286.]}

\medskip

\noindent The stereographic projection is one of the most important of the methods of representing the surface of a sphere upon a plane. It possesses many remarkable properties; among which may be mentioned the property that circles on the sphere are projected into circles on the plane, and the property that the projection is conformal, that is, that angles are preserved unaltered.

\smallskip

\noindent The property which I now proceed to investigate relates to the representation of great circles upon the stereographic projection. A great circle on the sphere is projected into a circle on the plane which passes through the projections of the poles of the great circle; and the centre of this circle is the projection of the pole of the great circle. The radius of the projected circle depends upon the distance of the great circle from the pole of projection; and it may be expressed in terms of this distance by a simple formula.

\smallskip

\noindent Let $P$ be the pole of projection, and let $C$ be the centre of the sphere. Let $Q$ be the pole of a great circle on the sphere, and let $R$ be the radius of the sphere. Then the distance of $Q$ from $P$ is $PQ = R \cdot \theta$, where $\theta$ is the angular distance; and the projection of the great circle is a circle whose centre is the projection of $Q$, and whose radius is $R \tan\frac{\theta}{2}$.

\smallskip

\noindent This property may be proved as follows. Let $AB$ be the great circle, and let $Q$ be its pole. Draw the plane through $P$, $C$, and $Q$; and let this plane intersect the great circle $AB$ in the points $A$ and $B$. Then $PA$ and $PB$ are the generators of the cone which projects the great circle; and the intersection of this cone with the plane of projection is the projected circle.

\newpage
\noindent\textbf{1.}\quad The stereographic projection of a sphere upon a plane is effected by drawing lines from a fixed point on the surface of the sphere to the points of the sphere, and taking the intersections of these lines with a fixed plane. The fixed point is called the pole of projection; and the fixed plane is called the plane of projection.

\smallskip

\noindent\textbf{2.}\quad The most important case is that in which the plane of projection is perpendicular to the diameter through the pole of projection. In this case, the projection is called the orthogonal stereographic projection; and it possesses the property that circles on the sphere are projected into circles on the plane.

\smallskip

\noindent\textbf{3.}\quad The formula for the radius of the projected circle of a great circle may be obtained as follows. Let $\theta$ be the angular distance of the great circle from the pole of projection; then the radius of the projected circle is
\[
r = R \tan\frac{\theta}{2},
\]
where $R$ is the radius of the sphere. This formula may be proved by considering the right-angled triangle formed by the pole of projection, the pole of the great circle, and a point on the great circle.

\smallskip

\noindent\textbf{4.}\quad The centres of the projected circles of all great circles lie on the plane of projection; and they are the projections of the poles of the great circles. The locus of the centres of the projected circles is therefore the whole plane of projection; and every point of this plane is the centre of the projection of some great circle.

\smallskip

\noindent\textbf{5.}\quad The property that the stereographic projection is conformal may be proved by showing that the scale of the projection is the same in all directions at any point. The scale of the projection is given by the formula
\[
\frac{ds'}{ds} = \frac{1}{\cos^2\frac{\theta}{2}},
\]
where $ds$ is an element of arc on the sphere, $ds'$ is the corresponding element on the plane, and $\theta$ is the angular distance from the pole of projection. Since this expression depends only on $\theta$, and not on the direction of the element, the projection is conformal.

\smallskip

\noindent\textbf{6.}\quad The stereographic projection is extensively used in astronomy for the construction of maps of the celestial sphere; and it is also used in navigation for the construction of Mercator's chart. The property that circles are projected into circles renders the projection peculiarly adapted for the representation of the celestial sphere; and the conformal property renders it useful for the measurement of angles.

\newpage
%% ========================================================================
%% PAPER 476 (start): ON THE DETERMINATION OF THE ORBIT OF A PLANET
%% ========================================================================
\section*{476. On the Determination of the Orbit of a Planet from Three Observations.}
\addcontentsline{toc}{section}{476. On the Determination of the Orbit of a Planet from Three Observations (p.~400, start)}

\noindent\textit{[From the Memoirs of the Royal Astronomical Society, vol.~XXXVIII. (1870), pp.~17--111. Read December 10, 1869.]}

\medskip

\noindent I propose to consider from a geometrical point of view the problem of the determination of the orbit of a planet from three observations. The problem has been treated by many eminent mathematicians; and the methods which have been given are in general analytical. The geometrical method has the advantage of presenting the results in a more intuitive form; and it affords a useful check upon the analytical investigations.

\smallskip

\noindent The three observations give the directions of the planet, as seen from the Earth, at three different epochs; and the problem is to determine the orbit which passes through the three points thus determined. The orbit is an ellipse with the Sun in one focus; and the elements of the orbit are the semi-axis major, the eccentricity, the inclination, the longitude of the node, the longitude of the perigee, and the mean motion.

\smallskip

\noindent The geometrical construction is founded upon the following principles. The three lines of sight from the Earth to the planet determine three points on the celestial sphere; and the problem is to find an ellipse, with the Sun in one focus, which passes through the three points determined by these lines. The construction depends upon the properties of the cone which has the Sun for its vertex, and which passes through the three positions of the planet.

\smallskip

\noindent Let $S$ be the Sun, $E_1$, $E_2$, $E_3$ the three positions of the Earth, and $P_1$, $P_2$, $P_3$ the three positions of the planet. The lines $E_1P_1$, $E_2P_2$, $E_3P_3$ are the lines of sight; and the problem is to determine the ellipse with focus $S$ which passes through $P_1$, $P_2$, $P_3$. The cone with vertex $S$ and generators $SP_1$, $SP_2$, $SP_3$ intersects the celestial sphere in a conic; and the problem is reduced to the determination of this conic.

\smallskip

\noindent The investigation will be continued in the subsequent pages of this memoir; and the complete solution of the problem will be given, with numerical examples and applications to the orbits of the planets.

\vfill
\begin{center}
\rule{0.5\textwidth}{0.4pt}
\end{center}
\clearpage


% =========================================================
% Chunk: cayley_vol07_pages_401_450.tex  (orig pages 401--450)
% =========================================================
\setlength{\parindent}{1em}
\setlength{\parskip}{0pt}

%% ================================================================
%% PAPER 469 (conclusion)
%% THE ANGULAR DISTANCE OF TWO PLANETS
%% From pp. 379 of the original (page 401 of PDF)
%% ================================================================

\section*{469.}
\addcontentsline{toc}{section}{469. The Angular Distance of Two Planets}
\markboth{THE ANGULAR DISTANCE OF TWO PLANETS.}{469]}

whence the expression becomes
\begin{align*}
\Cos H &= \cos(v-v') \cdot \left(\tfrac12 A + \tfrac12 D\right)\cos(\theta-\theta')
- \left(-\tfrac12 B + \tfrac12 C\right)\sin(\theta-\theta') \\
&\quad + \sin(v-v') \cdot \left(\tfrac12 A + \tfrac12 D\right)\sin(\theta-\theta')
+ \left(-\tfrac12 B + \tfrac12 C\right)\cos(\theta-\theta') \\
&\quad + \cos(v+v') \cdot \left(\tfrac12 A - \tfrac12 D\right)\cos(\theta+\theta')
- \left(\tfrac12 B + \tfrac12 C\right)\sin(\theta+\theta') \\
&\quad + \sin(v+v') \cdot \left(\tfrac12 A - \tfrac12 D\right)\sin(\theta+\theta')
+ \left(\tfrac12 B + \tfrac12 C\right)\cos(\theta+\theta'),
\end{align*}
or substituting for $A$, $B$, $C$, $D$, their values, and after a few easy reductions, we find
\begin{align*}
\Cos H &= \cos(v-v') \left\{
\begin{aligned}
&\tfrac12 + \tfrac12\cos\phi\cos\phi'
- \tfrac12(1-\cos\phi)(1-\cos\phi')\sin^2(\theta-\theta') \\
&\quad + \tfrac12\sin\phi\sin\phi' \qquad\qquad \cos(\theta-\theta')
\end{aligned}
\right\} \\
&\quad + \sin(v-v') \left\{
\begin{aligned}
&\tfrac12(1-\cos\phi)(1-\cos\phi')\sin(\theta-\theta')\cos(\theta-\theta') \\
&\quad + \tfrac12\sin\phi\sin\phi' \qquad\qquad \sin(\theta-\theta')
\end{aligned}
\right\} \\
&\quad + \cos(v+v') \left\{
\begin{aligned}
&\tfrac12(1-\cos\phi\cos\phi')\cos(\theta-\theta')\cos(\theta+\theta') \\
&\quad + \tfrac12(\cos\phi-\cos\phi') \quad \sin(\theta-\theta')\sin(\theta+\theta') \\
&\quad - \tfrac12\sin\phi\sin\phi' \qquad\quad \cos(\theta+\theta')
\end{aligned}
\right\} \\
&\quad + \sin(v+v') \left\{
\begin{aligned}
&\tfrac12(1-\cos\phi\cos\phi')\cos(\theta-\theta')\sin(\theta+\theta') \\
&\quad - \tfrac12(\cos\phi-\cos\phi') \quad \sin(\theta-\theta')\cos(\theta+\theta') \\
&\quad - \tfrac12\sin\phi\sin\phi' \qquad\quad \sin(\theta+\theta')
\end{aligned}
\right\}.
\end{align*}

For $\phi=\phi'=0$, the formula becomes, as of course it should do,
\[
\Cos H = \cos(v-v').
\]

It may be added, that if $f$, $f'$ are the true anomalies, $\omega$, $\omega'$ the longitudes of
pericentre in orbit, then $v=\omega+f$, $v'=\omega'+f'$; and we thence have for $\cos H$, formul\ae
of the like form, containing $\cos f\cos f'$, $\cos f\sin f'$, $\sin f\cos f'$, $\sin f\sin f'$, or containing
$\cos(f-f')$, $\sin(f-f')$, $\cos(f+f')$, $\sin(f+f')$, respectively, in place of the like functions
of $z$, $z'$, but with of course altered values of the coefficients.

\hfill\rule{1.5cm}{0.4pt}

\newpage
%% ================================================================
%% PAPER 470
%% NOTE ON THE ATTRACTION OF ELLIPSOIDS
%% From the Monthly Notices R.A.S., vol.~xxix. (1868--1869), pp.~254--257
%% Original pp. 380--383 (PDF pp. 402--405)
%% ================================================================

\section*{470.}
\addcontentsline{toc}{section}{470. Note on the Attraction of Ellipsoids}
\markboth{NOTE ON THE ATTRACTION OF ELLIPSOIDS.}{470]}

\noindent\textsc{Note on the Attraction of Ellipsoids.}%
\footnote{From the Monthly Notices of the Royal Astronomical Society, vol.~xxix. (1868--1869), pp.~254--257.}

\smallskip

\noindent
470. If an indefinitely thin shell of uniform density, bounded by two similar and
similarly-situated ellipsoids, attracts a point $P$ on its outer surface, it has been shown
geometrically by M.~Chasles that the attraction is in the direction of the normal at
$P$, and is equal to twice the attraction of an infinite plate, the thickness of which is
equal to the normal thickness at $P$. Assuming that the attraction is in the direction
of the normal, the proof is in fact as follows:---with $P$ as vertex, circumscribe to the
interior surface a cone; this divides the shell into three parts; the one, $D+E+F$,
exterior to the cone, the other two, $A+B$ and $C$, interior to the cone. It is shown
that in the direction of the normal the attraction of $C$ is equal to that of $A+B$;
and it is assumed that in comparison with these the attraction of $D+E+F$ may be
neglected; the whole attraction is thus equal to twice that of the portion $A+B$.

At the point where the normal at $P$ meets the internal surface draw the tangent plane
to the internal surface, thus dividing the portion $A+B$ into the solid cone $A$ and
a remaining portion $B$; it is assumed that in comparison with that of $A$ the attraction
of $B$ may be neglected; the whole attraction is thus equal to twice the attraction of
the solid cone $A$; and the attraction of this solid cone is in the limit (the aperture
or solid angle then becoming $=2\pi$) equal to the attraction of an infinite plate whose
thickness is equal to the altitude of the solid cone, that is, to the normal thickness
at $P$. And the attraction of the whole ellipsoidal shell is thus ultimately (that is,
when the shell is indefinitely thin) equal to twice the attraction of the infinite plate.

It is interesting to ascertain the orders of magnitude of the attractions of the
several portions of the shell, which attractions are compared in the foregoing investi-
gation; and this can be done very easily, when, instead of the ellipsoidal shell, we
have a spherical shell (bounded by two concentric spherical surfaces). The tangent plane
to the inner surface divides the portion $D+E+F$ into two portions $D$ and $E+F$;
and if with $P$ as vertex we describe a cone standing on the circle in which the
tangent plane meets the outer surface, the last-mentioned portion is hereby divided into
the portions $E$ and $F$; the whole shell is thus divided into the portions $A$, $B$, $C$, $D$, $E$, $F$,
each of them symmetrical in regard to the normal or radius at $P$, and consequently
attracting in the direction of this radius. I proceed to find the attractions of each
of these portions; it will appear, in accordance with the assumptions of the foregoing
investigation, that, taking the radii to be $1$ and $1+a$, that is, $a$ the thickness of the
shell, and supposing ultimately $a$ to become indefinitely small, the attractions of $A$
and $C$ are each ultimately $=2\pi a$, that is $=$ to the attraction of the infinite plate,
while the attractions of the other portions are of the order $a^2$, and thus vanish in
comparison with that of $A$ or $C$.

The attraction of an indefinitely thin cone or frustum of a cone, length $r$ and
solid angle $d\omega$ is $=r\,d\omega$; considering any such cone having $P$ for its vertex, if the
inclination of $r$ to the radius through $P$ is $=\theta$, and if the azimuth of the plane
through $r$ and the radius is $=\phi$, then we have $d\omega=\sin\theta\,d\theta\,d\phi$, the attraction $r\,d\omega$ is
$=r\sin\theta\,d\theta\,d\phi$, and this attraction resolved in the direction of the radius is
$=r\sin\theta\cos\theta\,d\theta\,d\phi$.

For the several cases which have to be considered, the value of $r$ is independent of
$\phi$, and the integration in regard to $\phi$ is always from $\phi=0$ to $\phi=2\pi$;---the attraction
is thus in each case $=2\pi\!\int r\sin\theta\cos\theta\,d\theta$, the expression of $r$ in the terms of $\theta$, and
the limits of $\theta$ being known for each of the several portions of the shell.

Taking $\theta_1$ for the semi-angle of the tangent cone, we have
\[
\sin\theta_1 = \frac{\sqrt{2a+a^2}}{1+a}, \quad
\cos\theta_1 = \frac{1}{1+a},
\]
and taking $\theta_2$ for the semi-angle of the cone which divides the portions $E$, $F$,
\[
\tan\theta_2 = \sqrt{\frac{a}{2(1+a)}}, \quad
\sin\theta_2 = \sqrt{\frac{a}{2(1+a)}}, \quad
\cos\theta_2 = \sqrt{\frac{2+a}{2(1+a)}}.
\]

\newpage
\noindent
For $F$ we have
\begin{gather*}
r = 2(1+a)\cos\theta, \quad \theta=\theta_1 \text{ to } \theta=\theta_2,\\
\text{Integral is } = 2(1+a)\int \sin\theta\cos^2\theta\,d\theta
= \tfrac23(1+a)\cos^3\theta_2.
\end{gather*}

For $D+E$ we have
\begin{gather*}
r = 2(1+a)\cos\theta, \quad \theta=\theta_2 \text{ to } \theta=\theta_1,\\
\text{Integral is } = 2(1+a)\int \sin\theta\cos^2\theta\,d\theta
= \tfrac23(1+a)(\cos^3\theta_2 - \cos^3\theta_1).
\end{gather*}

For $E$ we have
\begin{gather*}
r = \frac{a}{\cos\theta}, \quad \theta=0 \text{ to } \theta=\theta_2,\\
\text{Integral is } = a\int \sin\theta\,d\theta = a(\cos\theta_2 - \cos\theta_1).
\end{gather*}

For $A$ we have
\begin{gather*}
r = \frac{a}{\cos\theta}, \quad \theta=0 \text{ to } \theta=\theta_1,\\
\text{Integral is } = a\int \sin\theta\,d\theta = a(1 - \cos\theta_1).
\end{gather*}

For $A+B$ we have
\begin{gather*}
r = (1+a)\cos\theta - \sqrt{1-(1+a)^2\sin^2\theta},
\quad \theta=0 \text{ to } \theta=\theta_1,\\
\text{Integral is } = \int\{(1+a)\cos\theta - \sqrt{1-(1+a)^2\sin^2\theta}\}\sin\theta\cos\theta\,d\theta,\\
= (1+a)\left(-\tfrac13\cos^3\theta\right) + \tfrac16\{1-(1+a)^2\sin^2\theta\}^{\frac32},
\text{ between the limits},\\
= \tfrac13\{(1+a)(1-\cos^3\theta_1) - \tfrac12 a^3\},
\end{gather*}
and subtracting the above value of the integral for $A$, it at once appears that, for $B$,
the integral is
\[
= 2\pi\left\{a(-1+\cos\theta_1) + \tfrac13(1+a)(1-\cos^3\theta_1) - \tfrac16 a^3\right\}.
\]

\newpage
\noindent
Hence, calculating the approximate values, and restoring in each case the omitted
factor, $2\pi$, we have
\begin{align*}
\text{Attraction } A &= 2\pi a - 2\sqrt{2}\,\pi a^{\frac32}, \\
\text{'' } \qquad B &= \tfrac43\sqrt{2}\,\pi a^{\frac32}, \\
\text{'' } \qquad C &= 2\pi a - 2\sqrt{2}\,\pi a^{\frac32}, \\
\text{'' } \qquad D &= \tfrac43\sqrt{2}\,\pi a^{\frac32}, \\
\text{'' } \qquad E &= \tfrac13\sqrt{2}\,\pi a^{\frac32}, \\
\text{'' } \qquad F &= \tfrac43\sqrt{2}\,\pi a^{\frac32}; \\
\intertext{or, if we please,}
\text{Attraction } A+B &= 2\pi a - \tfrac43\sqrt{2}\,\pi a^{\frac32}, \\
\text{'' } \qquad C &= 2\pi a - \tfrac43\sqrt{2}\,\pi a^{\frac32}, \\
\text{'' } \qquad D+E+F &= \tfrac83\sqrt{2}\,\pi a^{\frac32};
\end{align*}
so that ultimately the attraction of the portion $D+E+F$ vanishes in comparison with
those of the portions $A+B$ and $C$; and the attraction of these last, that is, of the
whole shell, is $=4\pi a$, twice the attraction of an infinite plate of the thickness $a$.

\hfill\rule{1.5cm}{0.4pt}

\newpage
%% ================================================================
%% PAPER 471
%% NOTE ON THE PROBLEM OF THE DETERMINATION OF A PLANET'S ORBIT
%% FROM THREE OBSERVATIONS
%% From the Monthly Notices R.A.S., vol.~xxix. (1868--1869), pp.~257--259
%% Original pp. 384--389 (PDF pp. 406--410)
%% ================================================================

\section*{471.}
\addcontentsline{toc}{section}{471. Note on the Problem of the Determination of a Planet's Orbit from Three Observations}
\markboth{NOTE ON THE PROBLEM OF THE DETERMINATION OF A PLANET'S ORBIT.}{471]}

\noindent\textsc{Note on the Problem of the Determination of a Planet's Orbit from Three
Observations.}%
\footnote{From the Monthly Notices of the Royal Astronomical Society, vol.~xxix. (1868--1869), pp.~257--259.}

\smallskip

\noindent
471. The principle of the solution given in the \textit{Theoria Motus} may be explained very
simply as follows:

Consider three successive positions $O$, $O'$, $O''$ of a planet revolving about the
focus $S$; let $n$, $n'$, $n''$, denote the doubles of the triangular areas $O'SO''$, $OSO''$, and
$OSO'$ respectively (viz.~the triangular area means the area of the triangle included
between the two radius vectors and the chord joining their extremities), $r'$ the radius
vector $SO'$; $\theta''$, $\theta$, the times of describing the arcs $OO'$ and $OO''$ respectively, the
units of time and distance being such that the time is equal to the double area
divided by the square root of the half latus rectum ($c=2a^{\frac32}$ for the Period in
circular or elliptic orbit).

Then writing
\[
P = 2\left(\frac{n''}{\theta''} - \frac{n}{\theta}\right), \quad
Q = 2\left(\frac{n''}{\theta''} + \frac{n}{\theta}\right),
\]
(observe that $n+n''-n'$ is $=$ twice the triangle $OO'O''$), for neighbouring positions of
$O$, $O''$, the values of $P$, $Q$, are ultimately proportional to $\sin v'$, $\cos v'$ ($v'$ the true
anomaly at $O'$); the radius vector $r'$ is in the same manner proportional to $Q$, or its
value is given in terms of the ratio $P:Q$; and the value of $n'$ being known in terms
of $r'$, the equation $n+n''-n'=\text{twice triangle }OO'O''$ serves to determine $r'$.

\newpage
\noindent
Writing $n$, $n'$, $n''$, for the doubles of the triangular areas, and $\theta$, $\theta'$, $\theta''$ for the
respective times, the equations used by Gauss are in effect
\begin{align*}
\frac{n}{\theta} &= \frac{\sqrt{\mu p}}{r'r''\sin(v''-v')}, \\
\frac{n'}{\theta'} &= \frac{\sqrt{\mu p}}{r''r\sin(v-v'')}, \\
\frac{n''}{\theta''} &= \frac{\sqrt{\mu p}}{rr'\sin(v'-v)},
\end{align*}
($\mu=1$, unit of mass being that of the sun; $p$ the semi-latus rectum, $v$, $v'$, $v''$, the
longitudes in orbit).

Whence, putting for shortness
\[
\sin(v'-v'') = a, \quad \sin(v''-v) = b, \quad \sin(v-v') = c,
\]
($a+b+c=0$ as before), and also
\[
\theta'' = f, \quad \theta' = g, \quad \theta = h,
\]
then we have
\[
n : n' : n'' = \frac{a}{f} : \frac{b}{g} : \frac{c}{h},
\]
and the foregoing values of $P$, $Q$ become
\[
P = 2\sqrt{\mu p}\left(\frac{a}{f} - \frac{c}{h}\right), \quad
Q = 2\sqrt{\mu p}\left(\frac{a}{f} + \frac{c}{h}\right);
\]
these are the auxiliary equations: from $P$, $Q$, we have $r'$ as above; and thence
$p=Q^2/\mu$, and from $p$ we have the elements of the orbit.

\newpage
\noindent
Observe that the equation $n+n''-n'=\text{twice triangle }OO'O''$ gives for the
determination of $r'$ a cubic equation: the three roots correspond to the three ways in
which the three points $O$, $O'$, $O''$ may be connected with the focus $S$ by a conic, viz.
\begin{enumerate}
\item[(1)] The three points on the same branch, here $r'$ is positive;
\item[(2)] $O$, $O''$ on one branch, $O'$ on the other branch, here $r'$ is negative;
\item[(3)] $O$, $O'$ on one branch, $O''$ on the other, or $O'$, $O''$ on one branch, $O$ on the
other, here $r'$ is imaginary.
\end{enumerate}
The second case is that of a hyperbolic orbit described about the vacant focus, but
which may be regarded as described about the occupied focus $S$ by a repulsive force
from $S$: the third case is excluded from the physical problem.

\hfill\rule{1.5cm}{0.4pt}

\newpage
%% ================================================================
%% PAPER 472
%% NOTE ON LAMBERT'S THEOREM FOR ELLIPTIC MOTION
%% From the Monthly Notices R.A.S., vol.~xxix. (1868--1869), pp.~259--260
%% Original pp. 388--390 (PDF pp. 410--413)
%% ================================================================

\section*{472.}
\addcontentsline{toc}{section}{472. Note on Lambert's Theorem for Elliptic Motion}
\markboth{NOTE ON LAMBERT'S THEOREM FOR ELLIPTIC MOTION.}{472]}

\noindent\textsc{Note on Lambert's Theorem for Elliptic Motion.}%
\footnote{From the Monthly Notices of the Royal Astronomical Society, vol.~xxix. (1868--1869), pp.~259--260.}

\smallskip

\noindent
472. Let $r$, $r'$ be the radius vectors, $v$, $v'$ the true anomalies, $\omega$ the longitude of
pericentre, $e$ the eccentricity; then taking $\mu=1$ and writing
\[
x = r\cos(v-\omega), \quad y = r\sin(v-\omega), \quad
x' = r'\cos(v'-\omega), \quad y' = r'\sin(v'-\omega),
\]
so that $(x,y)$, $(x',y')$ are the coordinates of the two points referred to the axes at the
pericentre and a distance $90^\circ$ therefrom, we have
\[
x = a(\cos u - e), \quad y = a\sqrt{1-e^2}\sin u, \quad
x' = a(\cos u' - e), \quad y' = a\sqrt{1-e^2}\sin u',
\]
($u$, $u'$ the eccentric anomalies).

Whence, writing for shortness
\[
\xi = \cos\frac{u'-u}{2}, \quad \eta = \sin\frac{u'-u}{2}, \quad
\xi' = \cos\frac{u'+u}{2}, \quad \eta' = \sin\frac{u'+u}{2},
\]
then we find
\begin{align*}
\tfrac12(r+r') &= a - ae\xi\xi' + a\eta^2, \\
\tfrac12(r'-r) &= \phantom{a -}\, ae\eta\eta' + a\xi\eta, \\
ss' &= \phantom{a -}\, a^2(1-e^2)\eta^2,
\end{align*}
if for shortness $s$, $s'$ denote the two distances of the points from the extremity of the
minor axis; viz.~$s^2 = x^2+(y-a\sqrt{1-e^2})^2$, $s'^2 = x'^2+(y'-a\sqrt{1-e^2})^2$.

\newpage
\noindent
The equation of the chord joining the two points is
\[
xy' - x'y = (xx' + yy')\sin(v'-v) - (xy' - x'y)\cos(v'-v),
\]
or what is the same thing, it is
\[
rr'\sin(v'-v) = a^2\sqrt{1-e^2}(\cos u - \cos u'),
\]
which gives
\[
rr'\sin(v'-v) = 2a^2\sqrt{1-e^2}\,\xi\eta.
\]

Hence, substituting for $\omega$, $\omega'$ their values in terms of $\xi$, $\eta$; $\xi'$, $\eta'$, we have
\begin{align*}
\tfrac12(r+r'+c) &= a(1-e\xi\xi' + \eta^2 + \sqrt{1-e^2}\,\xi\eta), \\
\tfrac12(r+r'-c) &= a(1-e\xi\xi' + \eta^2 - \sqrt{1-e^2}\,\xi\eta), \\
\intertext{and thence}
\tfrac14\{(r+r')^2 - c^2\} &= a^2(1-e\xi\xi'+\eta^2)^2 - a^2(1-e^2)\xi^2\eta^2.
\end{align*}

Moreover
\[
ss' = a^2(1-e\xi\xi' + \eta^2)^2 - a^2\xi^2\eta^2 - a^2(1-e^2)\xi^2\eta^2,
\]
or writing this under the form
\[
ss' = a^2(1-e\xi\xi'+\eta^2)^2 - a^2(1-e^2)\xi^2\eta^2 - a^2e^2\xi^2\eta^2,
\]
the equation becomes
\[
ss' = \tfrac14\{(r+r')^2 - c^2\} - a^2e^2\xi^2\eta^2.
\]

\newpage
\noindent
Writing now
\[
\alpha = \cos^{-1}\left(1-\frac{r+r'+c}{2a}\right), \quad
\beta = \cos^{-1}\left(1-\frac{r+r'-c}{2a}\right),
\]
so that
\[
\cos\alpha = 1 - \frac{r+r'+c}{2a}, \quad
\cos\beta = 1 - \frac{r+r'-c}{2a},
\]
and therefore
\[
\cos\alpha - \cos\beta = -\frac{c}{a}, \quad
\cos\alpha + \cos\beta = 2 - \frac{r+r'}{a},
\]
then, observing that the coefficient of $\cos\alpha$ is $=+1$, and introducing it into the formula,
we have
\[
6\sqrt{\frac{a^3}{\mu}} = \alpha - \beta - (\sin\alpha - \sin\beta),
\]
where $\alpha$, $\beta$ are positive angles each less than $2\pi$: this is Lambert's theorem for elliptic
motion.

In hyperbolic motion, writing $\alpha$, $\beta$ for the hyperbolic excesses, that is
\[
\cosh\alpha = 1 - \frac{r+r'+c}{2a}, \quad
\cosh\beta = 1 - \frac{r+r'-c}{2a},
\]
then the like investigation gives
\[
6\sqrt{\frac{(-a)^3}{\mu}} = \sinh\alpha - \sinh\beta - (\alpha - \beta).
\]

\hfill\rule{1.5cm}{0.4pt}

\newpage
%% ================================================================
%% PAPER 473
%% ON THE GRAPHICAL CONSTRUCTION OF THE UMBRAL OR PENUMBRAL CURVE
%% From the Monthly Notices R.A.S., vol.~xxix. (1868--1869), pp.~260--261
%% Original pp. 390--392 (PDF pp. 414--417)
%% ================================================================

\section*{473.}
\addcontentsline{toc}{section}{473. On the Graphical Construction of the Umbral or Penumbral Curve}
\markboth{ON THE CONSTRUCTION OF THE UMBRAL OR PENUMBRAL CURVE \&C.}{473]}

\noindent\textsc{On the Graphical Construction of the Umbral or Penumbral
Curve \&c.}%
\footnote{From the Monthly Notices of the Royal Astronomical Society, vol.~xxix. (1868--1869), pp.~260--261.}

\smallskip

\noindent
473. The curve in question, say the penumbral curve, is the intersection of a sphere
having its centre at the point $S$ with a quadric cone having its vertex at the point $L$.
Considering a plane through $L$ and $S$, this meets the sphere in a circle and the cone in
two lines, and the intersection of these two lines with the circle gives the points in
which the penumbral curve is met by the plane; hence, taking a series of planes through
$L$, $S$, the penumbral curve can be constructed as the curve through the points thus
determined.

If the planes be taken through the axis $LS$, then instead of a variable plane we
may take a variable line through $L$ in a fixed plane through $LS$, say the plane of the
paper, and constructing in this plane the points in which the variable line meets the
circle, these points belong to the projection on the plane of the paper of the penumbral
curve; the planes through $LS$ give thus the projection of the penumbral curve on any
plane through $LS$, and in particular the projection on the plane of the ecliptic.

\newpage
\noindent
To prove this, take $x^2+y^2=1$ for the equation of the fixed circle, $(x-a)^2+(y-\beta)^2=\gamma^2$
for that of the circle which is the section of the sphere by the plane of the paper; and
let the equation of the quadric cone be
\[
A(x-a)^2 + B(y-\beta)^2 + Cz^2 + 2F(y-\beta)z + 2Gz(x-a) + 2H(x-a)(y-\beta) = 0.
\]
Writing herein $z=\rho\sin\theta$ and $x-a=\rho\cos\theta$, $y-\beta=\rho'\cos\theta$, where $\rho$, $\rho'$, $\rho''$ are
proportional to the sines of the distances from the point $(a,\beta,0)$ to the intersections of
the variable line through this point with the cone; and if the variable line lie in the
plane of the paper, then $\rho':\rho''$ are given as the roots of a quadratic equation, say
\[
(A\cos^2\theta + 2H\cos\theta\sin\theta + B\sin^2\theta)\rho^2
+ 2(G\cos\theta + F\sin\theta)\rho\sin\phi + C\sin^2\phi\,\rho^2 = 0,
\]
and we have then the coordinates of the point on the penumbral curve as follows:
\begin{align*}
x - a &= \frac{\rho'}{\rho''-\rho'}\,(\rho''\cos\theta), \\
y - \beta &= \frac{\rho'}{\rho''-\rho'}\,(\rho''\sin\theta), \\
z &= \frac{\rho'}{\rho''-\rho'}\,(\rho''\sin\phi),
\end{align*}
where $\rho'$, $\rho''$ are the roots of the quadratic; and hence the construction for the penumbral
curve as stated above.

\hfill\rule{1.5cm}{0.4pt}

\newpage
%% ================================================================
%% PAPER 474
%% ON THE GEOMETRICAL THEORY OF SOLAR ECLIPSES
%% From the Monthly Notices R.A.S., vol.~xxix. (1868--1869), pp.~261--264
%% Original pp. 393--396 (PDF pp. 418--423)
%% ================================================================

\section*{474.}
\addcontentsline{toc}{section}{474. On the Geometrical Theory of Solar Eclipses}
\markboth{ON THE GEOMETRICAL THEORY OF SOLAR ECLIPSES.}{474]}

\noindent\textsc{On the Geometrical Theory of Solar Eclipses.}%
\footnote{From the Monthly Notices of the Royal Astronomical Society, vol.~xxix. (1868--1869), pp.~261--264.}

\smallskip

\noindent
474. The fundamental equation in a solar eclipse is, I think, most readily established as
follows: the earth and moon being regarded as spheres of radii $R$, $S$ respectively, the
condition in order that the point $P$ may be on the umbral cone is that the two spheres
may subtend equal angles at $P$; viz.~if $r$, $r'$ are the distances of $P$ from the centres of the
earth and moon respectively, we must have
\[
\frac{R}{r} = \frac{S}{r'}.
\]

This is the equation of the umbral cone: to find its intersection with the earth's surface,
we have only to write $r=R$; and we thus obtain the equation of the curve of intersection,
viz.~this is $r'=\dfrac{S}{R}\,R$, that is, it is the intersection of the earth's surface with a cone
having its vertex at the moon's centre and standing on a circle (centre at the earth's
centre and radius $=\dfrac{S}{R}\cdot R=S$); the circle is, in fact, the limiting section of the moon
by a plane through the earth's centre, and the cone is the umbral cone corresponding
to this limiting section of the moon.

\newpage
\noindent
To obtain the equation of the penumbral curve, we have only to change the sign of
$S$; and the equations for the umbral and penumbral curves respectively are
\[
\frac{R}{r} = \pm\frac{S}{r'},
\]
or, what is the same thing,
\[
Rr' \mp Sr = 0.
\]

Writing $\rho$ for the distance of a point on the earth's surface from the axis $LS$, and $z$ for
its distance from the plane perpendicular to this axis through $S$, then for the moon and
earth we have respectively
\[
r'^2 = (D-z)^2 + \rho^2, \quad r^2 = z^2 + \rho^2,
\]
where $D$ is the distance $LS$ of the centres of the moon and earth; and the equation
becomes
\[
R\sqrt{(D-z)^2+\rho^2} \mp S\sqrt{z^2+\rho^2} = 0.
\]

It appears to me that the foregoing equation for the penumbral curve may be
obtained in a more simple and direct manner by considering the penumbral cone as the
envelope of the tangent cones which have their vertices at the several points of the moon's
surface, and which therefore envelope both the moon and the earth. Taking $M$ the moon's
centre, the tangent cone from $M$ to the earth has its equation
\[
(D-z)^2 + \rho^2 = \frac{D^2}{R^2}(z^2+\rho^2),
\]
or, as this may be written,
\[
R^2\{(D-z)^2+\rho^2\} = D^2(z^2+\rho^2).
\]

\newpage
\noindent
Hence, taking $(z',\rho')$ for the coordinates of a point on the moon's surface, the tangent
cone from this point to the earth has its equation
\[
R^2\{(D-z)^2+\rho^2\} = (D-z')^2 + \rho'^2\,(z^2+\rho^2),
\]
and the envelope of these cones is obtained in the usual manner by equating to zero
the discriminant in regard to $(z',\rho')$, considered as variable parameters connected by the
equation $z'^2+\rho'^2=S^2$ of the moon's surface. The result is found to be
\[
R^2\{(D-z)^2+\rho^2\} - S^2(z^2+\rho^2) = 0,
\]
which is the equation of the penumbral curve as given above.

\hfill\rule{1.5cm}{0.4pt}

\newpage
%% ================================================================
%% PAPER 475
%% ON A PROPERTY OF THE STEREOGRAPHIC PROJECTION
%% From the Monthly Notices R.A.S., vol.~xxix. (1868--1869), pp.~264--265
%% Original pp. 397--400 (PDF pp. 424--427)
%% ================================================================

\section*{475.}
\addcontentsline{toc}{section}{475. On a Property of the Stereographic Projection}
\markboth{ON A PROPERTY OF THE STEREOGRAPHIC PROJECTION.}{475]}

\noindent\textsc{On a Property of the Stereographic Projection.}%
\footnote{From the Monthly Notices of the Royal Astronomical Society, vol.~xxix. (1868--1869), pp.~264--265.}

\smallskip

\noindent
475. In the stereographic projection of a spherical figure upon a plane, to the circle
circumscribing any spherical triangle there corresponds a circle in the plane; and the
centre of this circle is the projection of the pole of the spherical circle. It follows that
if the spherical circle pass through a given point, the projected circle will pass through
the projection of this point; and further, that if the spherical circle touch a given small
circle, the projected circle will touch the projection of the small circle.

If the given point be the centre of projection, then the projection of any circle
through this point is a right line; hence, in particular, the projection of a great circle
through the centre of projection is a right line; and the angle between two such great
circles is equal to the angle between the corresponding right lines.

For the small circle, centre $C$ and radius $\gamma$, which touches the given small circle
centre $A$ and radius $\alpha$, the distance $AC$ is $=\alpha\pm\gamma$, and consequently the locus of $C$
is a circle, centre $A$ and radius $\alpha\pm\gamma$. But $C$ is the pole of the circle which touches the
given circle; and it has just been shown that the locus of the pole is a circle; hence the
projection of the pole (that is, the centre of the projected circle) lies on a circle which
is the projection of the locus of the pole.

\newpage
\noindent
We have moreover, considering the projection of a spherical triangle $ABC$, that the
projected triangle is similar to the triangle formed by the chords $AB$, $BC$, $CA$; the sides
of the projected triangle are proportional to the chords, that is, to $\sin\tfrac12c$, $\sin\tfrac12a$, $\sin\tfrac12b$;
and the angles of the projected triangle are equal to the angles of the chordal triangle,
that is, to the angles of the spherical triangle. Hence, if the side $AB$ and the angle $A$
be given, the locus of the vertex $C$ is a circle; and the angle at $C$ is given, that is, the
projected circle $AC$ passes through a fixed point, the angle subtended at this fixed point
by the line $AB$ being equal to the given angle $A$. And, conversely, if the projected
circle $AC$ pass through a fixed point, then the angle $A$ is given.

We may apply this to the case of the small circle which touches a given small circle
and passes through a given point; if the fixed point be the centre of projection, then
the circle becomes a right line; that is, the locus of the centre of the projected circle is
a circle, and the circle passes through the projection of the given point. Or again, if
the fixed point be any other point, then the locus of the centre of the projected circle is
a circle, and the circle passes through the projection of the given point.

\hfill\rule{1.5cm}{0.4pt}

\newpage
%% ================================================================
%% PAPER 476
%% ON THE DETERMINATION OF THE ORBIT OF A PLANET FROM
%% THREE OBSERVATIONS
%% From the Monthly Notices R.A.S., vol.~xxix. (1868--1869), pp.~265--278
%% Original pp. 401--428+ (PDF pp. 428--450)
%% ================================================================

\section*{476.}
\addcontentsline{toc}{section}{476. On the Determination of the Orbit of a Planet from Three Observations}
\markboth{ON THE DETERMINATION OF THE ORBIT \&C.}{476]}

\noindent\textsc{On the Determination of the Orbit of a Planet from Three
Observations.}%
\footnote{From the Monthly Notices of the Royal Astronomical Society, vol.~xxix. (1868--1869), pp.~265--278.}

\smallskip

\noindent
476. I propose to consider the problem of the determination of the orbit of a planet
from three observations. The problem has been treated of by Gauss in the \textit{Theoria Motus},
and by Encke and Galle in their editions of this work; also in the \textit{Astronomische Nachrichten}
(especially in connection with the discovery of the planet Neptune). The methods are well
known, but the subject is one which admits of being treated under a variety of different
aspects; and I have thought that a discussion of it by a new method would not be without
interest.

\newpage
\noindent
\textbf{Article Nos.~1 to 6. Preliminary Notions.}

\smallskip

\noindent
1. The data are the three heliocentric positions of the earth, and the three geocentric
positions of the planet; or, what is the same thing, the three directions of the planet as
seen from the earth, and the three times of observation. The three directions determine
three lines, which may be called ``rays;'' and the problem is to find the orbit of a planet
moving about the sun, and passing through three given points on three given rays.

2. The orbit is a conic having the sun for its focus; and the problem may be regarded
as that of finding a conic which shall pass through three given points and have its focus
at a given point. The conic may be an ellipse, a parabola, or a hyperbola; but in the
present memoir I consider chiefly the case of an elliptic orbit.

3. The coordinates of a point on the orbit may be expressed in terms of the true
anomaly $v$, the longitude of the pericentre $\omega$, and the semi-latus rectum $p$; or in terms
of the eccentric anomaly $u$, the mean anomaly $nt+\epsilon$, and the semi-axis major $a$. The
relations between these quantities are
\begin{align*}
r &= \frac{p}{1+e\cos(v-\omega)}, \\
r &= a(1-e\cos u), \\
nt + \epsilon &= u - e\sin u,
\end{align*}
where $e$ is the eccentricity, $n$ the mean motion, and $t$ the time.

4. The triangular areas are proportional to the times; and if we denote by $[OO']$ the
double area of the triangle formed by the points $O$, $O'$ and the focus $S$, then we have
\[
[OO'] : [O'O''] : [OO''] = \theta : \theta' : \theta'',
\]
where $\theta$, $\theta'$, $\theta''$ are the intervals of time between the observations.

\newpage
\noindent
5. The problem may be treated analytically or geometrically. The analytical treatment
leads to equations which, in their most general form, are somewhat complicated; but by
introducing suitable auxiliary variables, the equations may be simplified. The geometrical
treatment leads to constructions which, in their most general form, are somewhat compli-
cated; but by introducing suitable auxiliary constructions, the constructions may be
simplified.

6. In the present memoir I adopt a mixed method, combining the analytical and
geometrical treatments. The analytical part consists in the introduction of auxiliary
variables and the establishment of fundamental equations; the geometrical part consists
in the interpretation of these equations and the construction of the orbit.

\newpage
\noindent
\textbf{Article Nos.~7 to 14. The Regulator and Separators.}

\smallskip

\noindent
7. I say that in the \textit{first} of the cases above considered the locus of the orbit-pole is
a \textit{separator}; in the \textit{second} case the orbit-pole is a point $B$; in the \textit{third} case
the locus is the \textit{regulator}; and in the \textit{fourth} case the orbit-pole is a point $A$.

8. In the absence of models, the spherical figure must be represented by a projection;
the stereographic projection is convenient for facility of description; and it has the very
great advantage that we can by means of it exhibit, no matter how large a portion of the
spherical surface. In the figures called ``spherograms,'' afterwards referred to, the representation
of a hemisphere is all that is required; but, to give a more distinct general
idea, I annex a figure representing a larger portion of the surface; the data are those
belonging to the particular symmetrical case referred to as intended to be specially
considered: and the regulator conic is accordingly a pair of opposite small circles, the
points $A$ and $B$ being related to it symmetrically; but, disregarding these specialities,
the figure is adapted to the illustration of the general case (at least if the point $S$ be
situate \textit{within} the hyperboloid), and it is here given for that purpose.

The circle marked ``Ecliptic'' does not properly belong to the figure: it is added as
showing the boundary of a hemisphere, so that, by omitting all that lies outside this
circle, the figure would be limited to the representation of a hemisphere; and the orbit-pole
be in every case represented, no longer as a pair of opposite points, but as a single
point; we should have the separators each as a half circle, and the regulator as a single
small circle; the separators would intersect in pairs, in the three points $B$, and would
touch the regulator in the \textit{three} points $A$, \&c.

\newpage
\noindent
9. The figure constructed as above, but omitting so much of it as lies outside the
ecliptic circle, is the representation of a hemisphere---say of the northern hemisphere.

\noindent
\textbf{Article Nos.~10 to 14. Determination of the Orbit.}

\smallskip

\noindent
10. The problem of the determination of the orbit may be stated as follows: Given
the three rays and the three times, to find the orbit. The three rays determine the
trivector, and the three times determine the ratios of the triangular areas; and the
problem is to find the orbit corresponding to a given trivector and given ratios of the
triangular areas.

11. The problem may be solved by the method of the regulator and separators. The
regulator is a conic, and the separators are curves of the third order; and the intersection
of these curves determines the orbit-pole, and thence the orbit.

12. The method may be simplified by the introduction of auxiliary variables. The
auxiliary variables are the coordinates of the orbit-pole, or, what is the same thing, the
elements of the orbit; and the fundamental equations are the equations which express the
condition that the orbit shall pass through the three given points.

13. The fundamental equations may be written in the form
\begin{align*}
[a,b,c\,\wr\alpha',\beta',\gamma'] &= 0, \\
[a,b,c\,\wr\alpha'',\beta'',\gamma''] &= 0, \\
[a,b,c\,\wr\alpha''',\beta''',\gamma'''] &= 0,
\end{align*}
where $(a,b,c)$ are the coordinates of the orbit-pole, and $(\alpha',\beta',\gamma')$, $(\alpha'',\beta'',\gamma'')$,
$(\alpha''',\beta''',\gamma''')$ are the coordinates of the three points. The equations express the
condition that the orbit-pole shall lie on the three great circles which are the polars of
the three points in regard to the orbit.

14. The solution of the fundamental equations gives the coordinates of the orbit-pole,
and thence the elements of the orbit. The solution may be effected by the method of
successive approximations, or by the method of determinants.

\newpage
\noindent
\textbf{Article Nos.~15 to 30. Determination of the Orbit from a given Trivector.}

\smallskip

\noindent
15. With a given point $S$ as focus, and through three given points, that is with a given
trivector, there may be described four conics. This appears from the general theory
according to which a given focus is equivalent to two given tangents; and also from the
geometrical construction, Principia, book~I. sect.~4; Scholium to Prop.~xx:
viz.~given the focus $S$ and the points 1, 2, 3, then if
\begin{align*}
\text{On } 23 \text{ we find } a \text{ so that } a2 : a3 &= S2 : S3, \\
\text{'' } 31 \text{ '' } \quad b \text{ '' } \quad b3 : b1 &= S3 : S1, \\
\text{'' } 12 \text{ '' } \quad c \text{ '' } \quad c1 : c2 &= S1 : S2,
\end{align*}
the points $a$, $b$, $c$, are each of them on the directrix, so that any two of them determine
the directrix. In the figure (as in Newton's) the distances $S1$, $S2$, $S3$, are

\newpage
\noindent
each regarded as positive, but the very same construction, taking two of the distances
each as positive and negative successively, would lead to three other positions of the
directrix; or the construction would give in all four conics.

16. In the figure the directrix lies on the same side of the three points; and the conic
is thus an ellipse or parabola, or, if a hyperbola, then the three points lie in the same
branch thereof; and it is consequently an orbit such that along it a body can pass through
the three points successively. The construction as varied would give in each case a
directrix having on one side of it one, and on the other side two, of the three points;
so that the conic would be a hyperbola having the three points not on the same branch
thereof; consequently it would not be an orbit such that along it a body could pass
through the three points successively.

And it thus appears that though the trivector really determines four conics, yet it is
only one of these in which the directrix lies on the same side of the three points; and this
conic I call the ``orbit:'' the given trivector thus determines a single orbit.

17. It is to be noticed however that the orbit constructed as above may be a hyperbolic
branch separated by the directrix from the focus $S$, and consequently convex to the focus
$S$; viz.~the three points lie here in a hyperbolic branch convex to $S$, and which is
therefore not an orbit which can be described under the action of an attractive force at
$S$: say we have a ``convex orbit.'' I regard this as a real orbit, but the times of passage
therein as imaginary, or rather as non-existent, and the case as a physically impossible one.

\newpage
\noindent
\textbf{Article Nos.~31 to 45. The Isoparametric Curves.}

\smallskip

\noindent
18. The isoparametric curves are the loci of the orbit-pole for a given value of the
parameter. The parameter is the ratio of the triangular areas, and the isoparametric
curves are the curves on which the orbit-pole must lie in order that the parameter may
have the given value.

19. The isoparametric curves are curves of the third order, and they are the polars
of the points of the regulator in regard to the orbit. The regulator is a conic, and the
isoparametric curves are the curves which are the polars of the points of the conic.

20. The isoparametric curves may be constructed by the method of the regulator and
separators. The regulator is a conic, and the separators are curves of the third order;
and the intersection of these curves determines the isoparametric curves.

\newpage
\noindent
21. The isoparametric curves have the property that they pass through the three
points $A$, but pass through each point twice, or (what is the same thing) have each of the
points $A$ for a nodal point; when this is so, then it is to be further observed that, for
certain values of the parameters, they will be acnodal points, properly belonging to the
curve, although there is not any real branch of the curve passing through the points $A$;
for others they will be crunodal points, with two real branches through each; and in the
transition between the two cases they will be cuspidal points on the isoparametric curve;
it will appear in the sequel that this is really the case in regard to the isoeccentric lines.

\noindent
\textbf{Article Nos.~46 to 63. The Hyperboloid and the Orbit.}

\smallskip

\noindent
22. The three rays determine a hyperboloid of one sheet, and the orbit is a conic
section of this hyperboloid. The hyperboloid is determined by the three rays, and the
orbit is determined by the condition that it shall pass through the three points on the
three rays.

23. The equation of the hyperboloid may be written in the form
\[
\frac{x^2}{a^2} + \frac{y^2}{b^2} - \frac{z^2}{c^2} = 1,
\]
and the orbit is the intersection of this hyperboloid with a plane. The plane is determined
by the condition that it shall pass through the three points on the three rays; and the
orbit is the conic section of the hyperboloid by this plane.

\newpage
\noindent
24. The equation of the orbit may be written in the form
\[
\frac{x^2}{a^2} + \frac{y^2}{b^2} = 1,
\]
and the orbit is an ellipse, a parabola, or a hyperbola, according as $a^2$ and $b^2$ are both
positive, one positive and one negative, or both negative.

25. The orbit may be determined by the method of the regulator and separators. The
regulator is a conic, and the separators are curves of the third order; and the intersection
of these curves determines the orbit-pole, and thence the orbit.

26. The orbit may also be determined by the method of the hyperboloid. The
hyperboloid is determined by the three rays, and the orbit is determined by the condition
that it shall pass through the three points on the three rays.

\newpage
\noindent
27. The equation of the hyperboloid may be written in the form
\[
\frac{x^2}{a^2} + \frac{y^2}{b^2} - \frac{z^2}{c^2} = 1,
\]
and the orbit is the intersection of this hyperboloid with a plane. The plane is determined
by the condition that it shall pass through the three points on the three rays; and the
orbit is the conic section of the hyperboloid by this plane.

28. The orbit may be an ellipse, a parabola, or a hyperbola, according as the plane
intersects the hyperboloid in an ellipse, a parabola, or a hyperbola. The condition for an
ellipse is that the plane shall intersect all the generating lines of the hyperboloid; the
condition for a parabola is that the plane shall be parallel to a generating line; and the
condition for a hyperbola is that the plane shall intersect only one system of generating
lines.

29. The orbit may be determined by the method of the regulator and separators. The
regulator is a conic, and the separators are curves of the third order; and the intersection
of these curves determines the orbit-pole, and thence the orbit.

30. The orbit may also be determined by the method of the hyperboloid. The
hyperboloid is determined by the three rays, and the orbit is determined by the condition
that it shall pass through the three points on the three rays.

\newpage
\noindent
\textbf{Article Nos.~46 to 63. The Planogram and the Hyperbola.}

\smallskip

\noindent
46. The planogram is a representation of the orbit-plane on a fixed plane, and the
hyperbola is the curve in which the orbit-plane intersects a fixed quadric cone. The
planogram is determined by the orbit-plane and the fixed plane, and the hyperbola is
determined by the orbit-plane and the fixed quadric cone.

47. The equation of the hyperbola may be written in the form
\[
y^2 - (x-1)^2\tan^2\delta = m^2,
\]
where $\delta$ is the inclination of the orbit-plane to the fixed plane, and $m$ is a constant
depending on the orbit. The hyperbola is a rectangular hyperbola if $\delta=45^\circ$, and it
degenerates into a pair of right lines if $m=0$.

48. The hyperbola has the property that it passes through the point where the orbit-plane
is met by the ray, and that it touches at this point the projection of the ray on the
fixed plane. This property follows from the equation of the hyperbola, and it is the
fundamental property of the planogram.

49. The planogram may be constructed by the method of the hyperbola. The hyperbola
is determined by the orbit-plane and the fixed quadric cone, and the planogram is the
representation of the orbit-plane on the fixed plane by means of the hyperbola.

\newpage
\noindent
50. The equation of the hyperbola in the planogram for an orbit-plane rotating about
the line $x\cos b + y\sin b = 0$ through the focus $S$ is
\[
(x\cos b + y\sin b)^2 - (x\sin b - y\cos b - 1)^2\tan^2\delta = m^2,
\]
where $b$ is the longitude of the line of nodes, $\delta$ is the inclination of the orbit-plane to the
fixed plane, and $m$ is a constant depending on the orbit.

51. The hyperbola may be constructed by the method of the asymptotes. The asymptotes
are the lines $y = \pm(x-1)\tan\delta$, and the hyperbola is determined by its asymptotes
and a point on it.

52. The hyperbola may also be constructed by the method of the foci. The foci are
the points $(1\pm m\sec\delta, 0)$, and the hyperbola is determined by its foci and a point on it.

53. The hyperbola may also be constructed by the method of the directrix. The directrix
is the line $x = 1 + m\cot\delta\tan(\epsilon - A)$, and the hyperbola is determined by its directrix
and a point on it.

54. Imagine the orbit-plane (having upon it the hyperbola) brought by such rotation
into the plane $z=0$, or plane of the ecliptic, so that the hyperbola will be a curve in this
plane, the inclination to $Sx$, or longitude of the axis $Sx'$, being of course $=b-90^\circ$.
Transforming the equation to axes $Sx$, $Sy$, we must write in the equation
\[
\alpha = x\sin b - y\cos b, \quad \beta = x\cos b + y\sin b,
\]
and the equation thus becomes
\[
(x\cos b + y\sin b)^2 - (x\sin b - y\cos b - 1)^2\tan^2\delta = m^2.
\]

55. It will be recollected that the equations of the ray were
\[
\frac{x-A}{f} = \frac{y-B}{g} = \frac{z}{h},
\]
writing herein $z=0$ we find
\[
x = A - \frac{f}{h}C, \quad y = B - \frac{g}{h}C,
\]
and it is clear that this point $\left(A-\dfrac{f}{h}C, B-\dfrac{g}{h}C\right)$ should lie on the hyperbola.

Substituting for $(\alpha, \beta)$ the values in question, we have first
\begin{multline*}
\beta\sin b + \alpha\cos b - h =
\frac{1}{\{(f\cos b + g\sin b)^2 + h^2\}}\times\\
\{(f\cos b + g\sin b)^2(b\sin b + a\cos b) + h^2(b\sin b + a\cos b)\\
+ ch(f\cos b + g\sin b)\},
\end{multline*}
and observing that
\[
\tan^2\delta = \frac{h^2}{(f\cos b + g\sin b)^2},
\]
we have
\begin{multline*}
(\beta\sin b + \alpha\cos b - h)\tan\delta =
-(-a\sin b + b\cos b)(f\cos b + g\sin b);\\
\end{multline*}
and hence the result of the substitution is at once found to be
\begin{multline*}
(-a\sin b + b\cos b)^2 - h^2(-a\sin b + b\cos b)^2(g\sin b + f\cos b)^2\\
= m^2h^2 = h^2(-a\sin b + b\cos b)^2,
\end{multline*}
viz.~the factor $(-a\sin b + b\cos b)^2$ divides out, and the equation then becomes
\[
1 - h^2(g\sin b + f\cos b)^2 = m^2,
\]
that is
\[
\rho^2 = h^2 + (g\sin b + f\cos b)^2,
\]
which is in fact the value of $\rho^2$.

56. I seek for the direction of the hyperbola at the point $\left(\dfrac{f}{h}\rho, \dfrac{g}{h}\rho\right)$ in question.
We have
\begin{multline*}
dx : dy = (b\cos b - a\sin b)\sin b + \cos b\tan^2\delta\,(b\sin b + a\cos b - h\rho) :\\
- (b\cos b - a\sin b)\cos b + \sin b\tan^2\delta\,(b\sin b + a\cos b - h\rho),
\end{multline*}
and from the above values of $(b\sin b + a\cos b - h\rho)$ and $\tan^2\delta$, we have
\[
\tan^2\delta\,(b\sin b + a\cos b - h\rho) = (-a\sin b + b\cos b)(g\sin b + f\cos b);
\]
whence
\begin{multline*}
dx : dy = (b\cos b - a\sin b)\sin b\,(f\sin b - g\cos b) + (g\sin b + f\cos b)\cos b\,(-a\sin b + b\cos b) :\\
- (b\cos b - a\sin b)\cos b\,(f\sin b - g\cos b) + (g\sin b + f\cos b)\sin b\,(-a\sin b + b\cos b),
\end{multline*}
which, multiplying out and reducing by means of the relation $af+bg+ch=0$, becomes
\[
dx : dy = (-a\sin b + b\cos b)(\sin^2 b + \cos^2 b)f : (-a\sin b + b\cos b)(\sin^2 b + \cos^2 b)g;
\]
that is
\[
dx : dy = f : g,
\]
which shows that the hyperbola, at the point $\left(\dfrac{f}{h}\rho, \dfrac{g}{h}\rho\right)$ where it meets the ray, touches
the projection
\[
\frac{x-A}{f} = \frac{y-B}{g}
\]
of the ray on the plane of $xy$, which contains the hyperbola.

\newpage
\noindent
57. We may consider various particular forms of the hyperbola $y'^2 - (\alpha'-1)^2\tan^2\delta = m^2$.

1$^\circ$. If $\tan\delta=0$, the hyperbola is the pair of parallel lines $y'=\pm m$.

This can only happen if $h=0$, $f\cos b+g\sin b=0$. The first equation gives $af+bg=0$,
whence $\tan b=-\dfrac{a}{b}=\dfrac{g}{f}$; we have thus $m$ finite. The equations show that the ray is parallel
to the line of nodes.

2$^\circ$. If $\tan\delta=\infty$, the hyperbola is $(\alpha'-1)^2=0$, viz.~the line $\alpha'=1$ twice: the
condition is $-f\sin b+g\cos b=0$; viz.~the ray (not in general cutting the line of nodes)
is at right angles to the line of nodes.

3$^\circ$. If $m=0$, the hyperbola is the pair of intersecting lines $y'^2=(\alpha'-1)^2\tan^2\delta$.
The condition is $-a\sin b+b\cos b=0$, signifying that the ray cuts the line of nodes.

4$^\circ$. We may have simultaneously $\tan\delta=0$, $m=0$. The hyperbola (as in 2$^\circ$) is
$(\alpha'-1)^2=0$. The conditions are $-f\sin b+g\cos b=0$, $-a\sin b+b\cos b=0$, whence
$\tan b=\dfrac{g}{f}=\dfrac{b}{a}$, and therefore also $ag-bf=0$; these signify that the ray cuts at right
angles the line of nodes.

The line $y'=1$ passes through the point $\left(\dfrac{f}{h}\rho, \dfrac{g}{h}\rho\right)$, that is, we ought to have
$h^2=a^2+b^2$. The value of $h$ is in the first instance given in the form
\[
h = \frac{1}{a^2}\{(ah-cf)\cos b + (bh-cg)\sin b\},
\]
where
\[
\rho^2 = b^2 + (f\cos b+g\sin b)^2 = h^2 + f^2+g^2-(-f\sin b+g\cos b)^2 = f^2+g^2+h^2.
\]
But observe that the equations
\[
ag-bf=0, \quad bg+af=-ch,
\]
give
\[
\frac{g}{f} = \frac{b}{a}, \quad \frac{-ch}{a^2+b^2} = \frac{g}{f},
\]
and thence
\[
h = \frac{1}{a^2}\{(ah-cf)\cos b + (bh-cg)\sin b\},
\]
which is right.

\newpage
\noindent
58. I return to the equation of the hyperbola written in the form
\[
(x\cos b + y\sin b)^2 - (x\sin b - y\cos b - 1)^2\tan^2\delta = m^2;
\]
being (as was shown) a hyperbola passing through the point $\left(\dfrac{f}{h}\rho, \dfrac{g}{h}\rho\right)$ where its plane
is met by the ray, and touching at this point the projection $\dfrac{x-A}{f}=\dfrac{y-B}{g}$.

If in the equation we consider $b$ as variable, we have a series of hyperbolas, viz.~
these are the intersections of the plane of $xy$ with the hyperboloids of revolution obtained
by making the ray rotate successively round the several lines $x\cos b+y\sin b=0$
through the focus $S$. And, as just seen, these hyperbolas all of them touch at
$\left(\dfrac{f}{h}\rho, \dfrac{g}{h}\rho\right)$ the projection of the ray.

59. The hyperbola to any particular angle $b$ is the hyperbola belonging to the ray, in
the planogram for an orbit-plane rotating about the axis $x\cos b+y\sin b=0$; so that the
system of hyperbolas would be useful for the construction of any such planogram. And
there is another series of curves which, if they could be constructed with moderate
facility, would be very useful for the same purpose; viz.~reverting to the equations
\begin{align*}
x' : y' : 1 &= (a\cos b+b\sin b)\cos c - c\sin c :\\
&\qquad -a\sin b + b\cos b :\\
&\qquad (f\cos b+g\sin b)\sin c + h\cos c,
\end{align*}
which determine in the orbit-plane the coordinates $x'$, $y'$ of the intersection thereof with
the ray: imagine as before that the point is marked on the orbit-plane, and let it by a
rotation of the orbit-plane be brought into the plane of $xy$; so that $x'$, $y'$, will be the
coordinates in the direction of and perpendicular to the line of nodes of a point on the
hyperbola $y'^2-(x'-1)^2\tan^2\delta=m^2$, or $(x\cos b+y\sin b)^2-(x\sin b-y\cos b-1)^2\tan^2\delta=m^2$,
viz.~of the point corresponding to an orbit-pole, colatitude $c$. Suppose that $x$, $y$, are
the coordinates of this same point referred to the fixed axes, we have
\[
x = x'\sin b + y'\cos b, \quad y = -x'\cos b + y'\sin b,
\]
and thence
\begin{align*}
x : y : 1 &= (a\cos b+b\sin b)\sin b\cos c - c\sin b\sin c + (-a\sin b+b\cos b)\cos c :\\
&\quad -(a\cos b+b\sin b)\cos b\cos c + c\cos b\sin c + (-a\sin b+b\cos b)\sin b :\\
&\quad (f\cos b+g\sin b)\sin c + h\cos c,
\end{align*}
the coordinates of the point just referred to. Now, if from these equations we could
eliminate $b$, we should have a series of curves containing the variable parameter $c$,
intersecting the series of hyperbolas; and thus marking out on each of these hyperbolas
the points which belong to the successive values of the parameter $c$; we should thus
have in the plane of $xy$ the point corresponding to an orbit-pole longitude $b$ and
colatitude $c$. The series of curves in question may be called ``graduation curves,'' viz.~
they would serve for the graduation of the hyperbola in the planogram for an orbit-plane
rotating round any line $x\cos b+y\sin b=0$ in the plane of $xy$. But the elimination
cannot be easily effected, and I am not in possession of any method of tracing the series
of curves.

\newpage
\noindent
60. I remark that from the equations
\begin{align*}
x' : y' : 1 &= (a\cos b+b\sin b)\cos c - c\sin c :\\
&\qquad -a\sin b+b\cos b :\\
&\qquad (f\cos b+g\sin b)\sin c + h\cos c,
\end{align*}
we may without difficulty eliminate $b$; the result is, in fact,
\begin{multline*}
[x'(-ah\cos c) + y'(-bh\cos^2 c - cg\sin^2 c) - ac\sin c]^2 + \\
[x'(bh\cos c) + y'(-ah\cos^2 c + cf\sin^2 c) + bc\sin c]^2 = \\
[x'(ch\sin c) + y'(ag-bf)\sin c\cos c + (a^2-b^2)\cos c]^2,
\end{multline*}
a conic; but the geometrical signification of this result is not obvious, and I do not
make any use of it.

\bigskip
\noindent
\textbf{Article Nos.~61 to 63. The Trivector and the Orbit.}

\smallskip

\noindent
61. Considering now the three rays, these are determined by their six coordinates
\begin{align*}
&(a_1, b_1, c_1, f_1, g_1, h_1), \\
&(a_2, b_2, c_2, f_2, g_2, h_2), \\
&(a_3, b_3, c_3, f_3, g_3, h_3),
\end{align*}
respectively; and the intersections with the orbit-plane are given by
\begin{align*}
x_1 : y_1 : 1 &= (a_1,b_1,c_1\,\wr\alpha',\beta',\gamma')
: -(a_1,b_1,c_1\,\wr\alpha',\beta',\gamma') : (f_1,g_1,h_1\,\wr\alpha',\beta',\gamma'), \\
x_2 : y_2 : 1 &= (a_2,b_2,c_2\,\wr\alpha'',\beta'',\gamma'')
: -(a_2,b_2,c_2\,\wr\alpha'',\beta'',\gamma'') : (f_2,g_2,h_2\,\wr\alpha'',\beta'',\gamma''), \\
x_3 : y_3 : 1 &= (a_3,b_3,c_3\,\wr\alpha''',\beta''',\gamma''')
: -(a_3,b_3,c_3\,\wr\alpha''',\beta''',\gamma''') : (f_3,g_3,h_3\,\wr\alpha''',\beta''',\gamma'''),
\end{align*}
where the axes $Sx'$, $Sy'$, are an arbitrary set of rectangular axes in the orbit-plane;
or where, as before, the axis $Sx'$ may be taken to be the line of nodes.

There is no difficulty in finding the equation of the orbit. Writing $r_1=\sqrt{x_1^2+y_1^2}$,
we have
\[
r_1 = \frac{(f_1,g_1,h_1\,\wr\alpha',\beta',\gamma'')}{(f_1,g_1,h_1\,\wr\alpha',\beta',\gamma')},
\]
if
\[
u = \sqrt{[(a_1,b_1,c_1\,\wr\alpha',\beta',\gamma')]^2 + [(a_1,b_1,c_1\,\wr\alpha',\beta',\gamma'')]^2}.\]

\newpage
\noindent
the sign being taken in such manner that $r_1$ shall be positive; viz.~the sign must be the
same as that of $(f_1,g_1,h_1\,\wr\alpha'',\beta'',\gamma'')$. And we have the like formul\ae\ for $r_2$, and $r_3$.
Substituting these values, the equation of the orbit becomes
\[
\begin{vmatrix}
r_1 & x_1 & y_1 & 1 \\
r_2 & x_2 & y_2 & 1 \\
r_3 & x_3 & y_3 & 1 \\
\end{vmatrix} = 0,
\]
or in developed form:
\begin{multline*}
\{x'(a_1,b_1,c_1\,\wr\alpha',\beta',\gamma') + y'(a_1,b_1,c_1\,\wr\alpha',\beta',\gamma')\}
[-r_1(f_1,g_1,h_1\,\wr\alpha'',\beta'',\gamma'') + r_2(f_2,g_2,h_2\,\wr\alpha'',\beta'',\gamma'')] + \\
\{x'(a_2,b_2,c_2\,\wr\alpha'',\beta'',\gamma'') + y'(a_2,b_2,c_2\,\wr\alpha'',\beta'',\gamma'')\}
[-r_2(f_2,g_2,h_2\,\wr\alpha''',\beta''',\gamma''') + r_3(f_3,g_3,h_3\,\wr\alpha''',\beta''',\gamma''')] + \\
\{x'(a_3,b_3,c_3\,\wr\alpha''',\beta''',\gamma''') + y'(a_3,b_3,c_3\,\wr\alpha''',\beta''',\gamma''')\}
[-r_3(f_3,g_3,h_3\,\wr\alpha',\beta',\gamma') + r_1(f_1,g_1,h_1\,\wr\alpha',\beta',\gamma')] = 0,
\end{multline*}
being an equation of the form $\Omega = Ax' + By' + C$.

62. Considering the minor determinants formed with the terms under the $x'$ and $y'$,
for instance
\begin{multline*}
(a_1,b_1,c_1\,\wr\alpha',\beta',\gamma'') \cdot -(a_1,b_1,c_1\,\wr\alpha',\beta',\gamma') \\
+ (a_1,b_1,c_1\,\wr\alpha',\beta',\gamma') \cdot (a_2,b_2,c_2\,\wr\alpha'',\beta'',\gamma'');
\end{multline*}
this is
\begin{multline*}
= (b_1c_2-b_2c_1)(\beta'\gamma''-\beta''\gamma') + (c_1a_2-c_2a_1)(\gamma'\alpha''-\gamma''\alpha')\\
+ (a_1b_2-a_2b_1)(\alpha'\beta''-\alpha''\beta') \\
= \alpha'''(b_1c_2-b_2c_1) + \beta'''(c_1a_2-c_2a_1) + \gamma'''(a_1b_2-a_2b_1),
\end{multline*}
or, what is the same thing,
\[
= (b_1c_2-b_2c_1, c_1a_2-c_2a_1, a_1b_2-a_2b_1 \,\wr\, \alpha''',\beta''',\gamma''');
\]
with the like expressions for the other two minors. And we thus obtain the following
developed form of the equation, viz.
\begin{multline*}
\{x'(a_1,b_1,c_1\,\wr\alpha',\beta',\gamma') + y'(a_1,b_1,c_1\,\wr\alpha',\beta',\gamma')\}
[-r_1(f_1,g_1,h_1\,\wr\alpha'',\beta'',\gamma'') + r_2(f_2,g_2,h_2\,\wr\alpha'',\beta'',\gamma'')] + \\
+ (b_1c_2-b_2c_1, c_1a_2-c_2a_1, a_1b_2-a_2b_1 \,\wr\, \alpha''',\beta''',\gamma'')
[r_1(f_1,g_1,h_1\,\wr\alpha'',\beta'',\gamma'') - r_2] + \\
+ (b_2c_3-b_3c_2, c_2a_3-c_3a_2, a_2b_3-a_3b_2 \,\wr\, \alpha',\beta',\gamma')
[r_2(f_2,g_2,h_2\,\wr\alpha''',\beta''',\gamma''') - r_3] + \\
+ (b_3c_1-b_1c_3, c_3a_1-c_1a_3, a_3b_1-a_1b_3 \,\wr\, \alpha'',\beta'',\gamma'')
[r_3(f_3,g_3,h_3\,\wr\alpha',\beta',\gamma') - r_1] = 0,
\end{multline*}
being an equation of the form $\Omega = Ax' + By' + C$.

63. The coefficient of $r$ is a quadric function of $(\alpha''',\beta''',\gamma''')$, and if this vanish the
orbit is a right line. It thus appears that the orbit will be a right line provided the
orbit-axis be situate in a certain quadric cone, or (what is the same thing) the orbit-pole
be situate in a certain spherical conic: agreeing with a preceding result, viz.~the cone is
that reciprocal to the cone, vertex $S$, circumscribed about the hyperboloid which contains
the three rays. And we see that the equation of this reciprocal cone is
\[
\begin{vmatrix}
\alpha'' & \beta'' & \gamma'' \\
(f_1,g_1,h_1\,\wr\alpha'',\beta'',\gamma''), & a_1, & b_1, & c_1 \\
(f_2,g_2,h_2\,\wr\alpha'',\beta'',\gamma''), & a_2, & b_2, & c_2 \\
(f_3,g_3,h_3\,\wr\alpha'',\beta'',\gamma''), & a_3, & b_3, & c_3 \\
\end{vmatrix} = 0.
\]

\newpage
\noindent
\textbf{Article Nos.~64 and 65. The Special Symmetrical System of three Rays.}

\smallskip

\noindent
64. In what follows I consider the three rays forming a symmetrical system as already
referred to: viz.~the three rays intersect the plane of the ecliptic at points equidistant
from $S$ at longitudes $0^\circ$, $120^\circ$, $240^\circ$; each of them is at right angles to the line
joining $S$ with the intersection with the plane of the ecliptic, and at an inclination $=60^\circ$
to this plane: the figure shows the projection on the plane of the ecliptic of the portions
which lie above this plane of the three rays respectively.

The three rays lie on a hyperboloid of revolution having the line $Sz$ for its axis;
the circumscribed or asymptotic cone vertex $S$, is a right cone of the semi-aperture $=30^\circ$;
the reciprocal cone is therefore a right cone semi-aperture $60^\circ$, or (what is the same thing)
the regulator is a small circle, angular radius $60^\circ$, and the regulator and separators have
the positions shown in fig.~1, see No.~8.

Taking $S1=S2=S3=1$, and writing down the equations of the three rays in the forms
\begin{align*}
\frac{x-1}{0} &= \frac{y}{1} = \frac{z}{\tan 60^\circ}, \\
\frac{x+\cos 60^\circ}{-\sin 60^\circ} &= \frac{y-\sin 60^\circ}{-\cos 60^\circ} = \frac{z}{\tan 60^\circ}, \\
\frac{x+\cos 60^\circ}{\sin 60^\circ} &= \frac{y+\sin 60^\circ}{-\cos 60^\circ} = \frac{z}{\tan 60^\circ},
\end{align*}
these are the equations which determine the three rays.

\newpage
\noindent
65. The investigation of the particular case leads to results which are interesting in
themselves, and which serve to illustrate the general theory. The symmetrical system
of three rays is the simplest case of the general problem, and it is the case which is most
likely to occur in practice. The results obtained in this case may be extended to the
general case by the method of small variations; and they serve as a starting point for
the solution of the general problem.

The orbit determined by the three rays is in this case a conic having the sun for its
focus, and passing through the three points on the three rays. The conic may be an
ellipse, a parabola, or a hyperbola; and the determination of the orbit depends on the
solution of the fundamental equations, which express the condition that the orbit shall
pass through the three given points.

The solution of the fundamental equations gives the coordinates of the orbit-pole, and
thence the elements of the orbit. The solution may be effected by the method of successive
approximations, or by the method of determinants; and the result is a conic which
passes through the three given points and has the sun for its focus.

The investigation of the particular case serves to verify the general theory, and to
show that the results obtained are correct. It also serves to illustrate the method of the
regulator and separators, and to show how the orbit may be determined by the inter-
section of these curves.

\hfill\rule{1.5cm}{0.4pt}

\vfill
\begin{center}
\textit{[End of Vol.~VII, Part I. Papers 469--476.]}
\end{center}
\clearpage


% =========================================================
% Chunk: cayley_vol07_pages_451_500.tex  (orig pages 451--500)
% =========================================================
\thispagestyle{fancy}

\noindent\rule{\textwidth}{0.5pt}

\section*{Paper 476: On the Determination of the Orbit of a Planet from Three Observations (continued)}

\begin{art}[No.~64 (continued)]
We obtain the six coordinates of the three rays respectively
\begin{align*}
(a_1, b_1, c_1, f_1, g_1, h_1) &= (0, \surd 3, -1, 0, 1, \surd 3),\\
(a_2, b_2, c_2, f_2, g_2, h_2) &= (3, \surd 3, 2, \surd 3, 1, -2\surd 3),\\
(a_3, b_3, c_3, f_3, g_3, h_3) &= (-3, \surd 3, 2, -\surd 3, 1, -2\surd 3),
\end{align*}
whence the intersections with the orbit-plane are given by
\begin{align*}
x_1' : y_1' : 1 &= \beta'\surd 3 - \gamma' : -\beta'\surd 3 + \gamma' : \beta'' + \gamma''\surd 3,\\
x_2' : y_2' : 1 &= 3\alpha' + \beta'\surd 3 + 2\gamma' : -3\alpha - \beta\surd 3 - 2\gamma : \alpha''\surd 3 + \beta'' - 2\surd 3\gamma'',\\
x_3' : y_3' : 1 &= -3\alpha' + \beta'\surd 3 + 2\gamma : 3\alpha - \beta\surd 3 - 2\gamma : -\alpha''\surd 3 + \beta'' - 2\surd 3\gamma'',
\end{align*}
where if (as before) the position of the orbit-plane be determined by means of the longitude $b$ and colatitude $c$ of the orbit-pole, we have
\begin{align*}
\alpha, \beta, \gamma &= \sin b, \cos b, 0,\\
\alpha', \beta', \gamma' &= \cos b \cos c, \sin b \cos c, \sin c,\\
\alpha'', \beta'', \gamma'' &= \cos b \sin c, \sin b \sin c, \cos c,
\end{align*}
and the passage from the coordinates $x', y'$ to $x, y$, is given by
\begin{align*}
x' &= x \sin b - y \cos b,\\
y' &= x \cos b + y \sin b,
\end{align*}
or conversely
\begin{align*}
x &= x' \sin b + y' \cos b,\\
y &= -x' \cos b + y' \sin b.
\end{align*}
\end{art}

\begin{art}[No.~65]
To develope the results, I consider the orbit-pole as passing through certain series of positions. The locus may be a meridian circle; by reason of the symmetry of the system, the results are not altered by a change of $120°$ in the longitude of the meridian; so that, by considering the two meridians $0°$--$180°$ and $90°$--$270°$, we in fact consider twelve half meridians at the intervals of $30°$. An illustration is afforded by Plate~I.; the orbit-pole describes successively the meridians $0°$, $30°$, $60°$, $90°$, and the line 1, by its intersection with the orbit-plane, gives successively the lines through the centre at the inclinations $0°$, $30°$, $60°$, $90°$ to the axis of $x$.
\end{art}

\noindent\textbf{Article Nos.~66 to 82. Planogram No.~1, the Meridian $90°$--$270°$} (see Plate~II.).

\begin{art}[No.~66]
Supposing that the orbit-plane rotates about the axis $SI$ (fig.~6, $s$ --- see No.~64) in the plane of the ecliptic, the orbit-pole will describe the meridian $90°$--$270°$, the position of the orbit-pole being $b=90°$, $c=0°$ to $90°$, or else $b=270°$, $c=0°$ to $90°$. But the same analytical formula extends to the two half meridians, viz., we may take $b=90°$, and extend $c$ over $180°$, in the final results making $c$ an arc between $0°$ and $90°$, and $b=90°$, or $=270°$, as the case requires.
\end{art}

\begin{art}[No.~67]
Assuming then $b=90°$, we have
\begin{align*}
\alpha, \beta, \gamma &= 1, 0, 0,\\
\alpha', \beta', \gamma' &= 0, \cos c, \sin c,\\
\alpha'', \beta'', \gamma'' &= 0, \sin c, \cos c,
\end{align*}
and, moreover, $x', y' = x, y$; so that instead of $(x_1', y_1')$, \&c., we may write at once $(x_1, y_1)$, \&c. The formulae become
\begin{align*}
x_1 : y_1 : 1 &= \surd 3 \cos c + \sin c : \sin c + \surd 3 \cos c,\\
x_2 : y_2 : 1 &= \surd 3 \cos c - 2 \sin c : -3 : \sin c - 2\surd 3 \cos c,\\
x_3 : y_3 : 1 &= \surd 3 \cos c - 2 \sin c : 3 : \sin c - 2\surd 3 \cos c,
\end{align*}
that is
\[
x_1 = 1, \quad y_1 = 0,
\]
(viz. the orbit-plane, as is evident, meets the ray 1 in a fixed point, its intersection with the plane of $xy$);
\[
x_2 = \frac{-\surd 3 \cos c - 2 \sin c}{\sin c - 2\surd 3 \cos c}, \quad
y_2 = \frac{-3}{\sin c - 2\surd 3 \cos c},
\]
and writing
\[
\frac{-2\surd 3}{-1} = 3 = \cos\omega, \quad \frac{\surd(-1)}{-1} = \frac{\surd 13}{3} = \sin\omega, \quad 2\surd 3 = \tan\omega,
\]
(whence $\omega = 16°\, 6'$) we find
\[
x_2 = \frac{-\surd 13}{3\surd 3} \sec(c+\omega), \quad
y_2 = \frac{-\surd 13}{3\surd 3} \tan(c+\omega),
\]
and we thence have for the hyperbola, the locus of $(x_2, y_2)$ and $(x_3, y_3)$
\[
x_2^2 - y_2^2 = \frac{13}{27},
\]
\end{art}

\begin{art}[No.~68]
viz.	he points $(x_2, y_2)$ and $(x_3, y_3)$ are situate on the hyperbola, symmetrically on opposite sides of the axis of $x$. For $c=0$, we have $x_2 = -\frac{1}{2}$, $y_2 = \frac{1}{2}\surd 3$, $(x_2^{\,2} + y_2^{\,2} = 1)$, and the hyperbola at this point touches the circle $x^2 + y^2 = 1$; and similarly for $x_3$, $y_3$.

The inclination of the asymptotes to the axis of $y$ is given by $\tan\eta = \surd\frac{3}{2}$, $\eta = 22°\, 56'$.
\end{art}

\begin{art}[No.~68 (continued)]
The orbits are conics, focus $S$ and vertex $1$. It will be convenient to consider $c$ as passing from $0°$ to $90° - \omega$, and from $0°$ to $-(90° + \omega)$; that is, from $0°$ to $73°\, 54' - \epsilon$, and from $0°$ to $-116°\, 6' + \epsilon$, if $\epsilon$ be indefinitely small: the point $2$ will thus traverse the upper branch (alone shown in the Plate) of the guide-hyperbola, viz., for $c=0°$ it will be at the point of contact with the circle; for $c=73°\, 54' - \epsilon$ it will be at $\infty$, and for $c=-106°\, 6' + \epsilon$ at $\infty'$. For $c=0°$ the orbit is the circle; as $c$ increases positively, it becomes an ellipse of increasing eccentricity and major axis, until for a certain value ($c = 46°\, 48'$ as will appear) it becomes a parabola; it then becomes a hyperbola (concave branch); for $c = 52°\, 45'$ it becomes the hyperbola $S'$ subsequently referred to; and for $c = 60°$ (the point $2$ being then on the line shown in the figure) the orbit becomes this right line. As $c$ continues to increase, the orbit becomes a hyperbola (convex branch); and ultimately for $c = 73°\, 54' - \epsilon$, the point $2$ goes to $\infty$, and the orbit becomes a hyperbola (convex) $S$, having an asymptote parallel to the line $2$.
\end{art}

\begin{art}[No.~69]
The point in question ($b = 90°$, $c = 73°\, 54'$) is one of the points $B$ of the spherogram, and the hyperbolas $S$, $S'$ are two of the four orbits belonging to this point. And, by what precedes, it appears that as the orbit-pole passes through this point along a meridian downwards to the ecliptic the change is from a concave to a convex orbit.
\end{art}

\begin{art}[No.~72]
On account of the symmetry in regard to the axis of $x$, the equation of the orbit will be of the form $r = Ax + B$; viz., the equation is at once found to be
\[
r = \frac{a_2 - 1}{x_2}x + a.
\]
\end{art}

\begin{art}[No.~73]
The eccentricity is the coefficient $A$ taken positively ($e = \pm A$): it is in the present case proper to attend to the value of the coefficient itself,
\[
A = \frac{a_2 - 1}{x_2};
\]
the sign of will then indicate the position of the centre of the orbit, viz., according as $A$ is positive or negative the centre will be on the negative or the positive side of the focus $S$. To investigate the variation of $A$, we may express it as a function of $\tan c = \lambda$, suppose. We have
\[
\frac{1}{x_2} = \frac{\sin c - 2\surd 3 \cos c}{-\surd 3 \cos c - 2 \sin c} = \frac{\lambda - 2\surd 3}{-\surd 3 - 2\lambda},
\]
and thence
\[
r_2^{\,2} = \frac{-3}{\lambda - 2\surd 3} = R_2, \quad R_2 = \pm \sqrt{-3\lambda^2 + 12\surd 3\lambda + 36},
\]
viz., $r_2$ must be positive, that is, $R_2$ is positive or negative according to the sign of $\lambda - 2\surd 3$; negative if $\lambda < 2\surd 3$ or $c < 73°\, 54'$, positive if $\lambda > 2\surd 3$ or $c > 73°\, 54'$. And we have then
\[
A = \frac{\lambda - 2\surd 3 - R_2}{3(\lambda - \surd 3)}.
\]
But a more convenient formula is obtained by writing
\[
\theta = \frac{5\cot c + 7}{2\surd 3}, \quad \text{that is } \theta = \frac{5 + 7\lambda}{2\surd 3\lambda},
\]
we then have
\[
A = \frac{-\theta \pm \surd(1 + \theta^2)}{3},
\]
which determines the sign of the radical, viz., this must have the same sign as $\theta$; and then for the coefficient
\[
A = \frac{-\theta \pm \surd(1 + \theta^2)}{3}.
\]
\end{art}

\begin{art}[No.~74]
For $c$ a small arc $= \epsilon$, $\theta$ is large and negative, and $\surd(1 + \theta^2)$, having the same sign as $\theta$, is $= \theta + \frac{1}{2\theta}$ nearly; we have therefore
\[
A = \frac{1}{6\theta} = \frac{\lambda}{3(5 + 7\lambda)} = \frac{\lambda}{15} \text{ nearly}.
\]
For $c$ nearly $= 60°$, say $c = 60° + \epsilon$,
\[
\cot c = \cot 60° + \epsilon \, \text{cosec}^2 60° = \frac{1}{\surd 3} + \frac{2\epsilon}{3},
\]
and thence
\[
\theta = \frac{-1 + \surd 13}{2\surd 3} + \frac{43\epsilon}{6}, \quad A = \frac{-8\surd 3}{-1 + \surd 13},
\]
viz., this is $-\infty$ for $c = 60° - \epsilon$, and $+\infty$ for $c = 60° + \epsilon$.

For $c$ nearly $= 90° - \omega$, say first $c = 73°\, 54' - \epsilon$, we have
\[
\cot c = +\frac{7}{2\surd 3}\epsilon, \quad \theta = -\frac{1}{2\surd 3}\epsilon, \quad \theta + \surd(1 + \theta^2) = +1,
\]
whence
\[
A = +\frac{1}{3} = \frac{\surd 3}{2.30940},
\]
but if $c = 73°\, 54' + \epsilon$, then $\theta = \frac{7}{2\surd 3}\epsilon$, $\theta + \surd(1 + \theta^2) = 4 + 1$, and
\[
A = -\frac{4}{3} = -\frac{\surd 3}{2.30940},
\]
viz., there is an abrupt change from $A = +\frac{\surd 3}{2.30940}$ to $A = \infty$; corresponding to the discontinuity of orbit already referred to. We may diminish $c$ by $180°$, and consider the last-mentioned value, $A = -\frac{4}{3}$, as belonging to $c = -90° - \omega + \epsilon = -(106°\, 6' - \epsilon)$.
\end{art}

\begin{art}[No.~75]
Consider next that $c$ passes from $0$ to $-(106°\, 6' - \epsilon)$. First if $c$ is a small negative quantity $c = -\epsilon$, $\theta$ is large and positive, and $\surd(1 + \theta^2)$ having the same sign as $\theta$ (positive) is $= -\theta + \frac{1}{2\theta}$ nearly, we have therefore $A = \frac{1}{30\theta} = -\frac{\lambda}{3\surd 5 + 2\lambda\theta}$ (same as for $c = +\epsilon$). And it is easy to see that as $c$ increases negatively, $A$ is always increasing negatively, its value for $c = -90°$ being $A = \frac{1 - \surd 13}{3} = -0.8685$, and for $c = -106°\, 6' + \epsilon$ being $= -2.30940$ as above. We have a diagram of $A$ (see next page).
\end{art}

\begin{art}[No.~76]
It thus appears that from $A = 0$ to $A = \frac{4}{3}$, there are always to any given value of $A$ two values of $c$, or positions of the orbit-pole. In particular if $A$ be $= -1$, the curve will be a parabola; the values of $c$ lying between $0°$, $60°$ and between $73°\, 54'$, $90°$ respectively.

To find them, writing $A = -1$, we have
\[
-3\lambda - 3a = 2\theta - 2\surd(1 + \theta^2), \text{ that is, } 5\theta + 3a = 2\surd(1 + \theta^2),
\]
or
\[
21\theta^2 + 30a\theta + 9a^2 - 4 = 0,
\]
that is, substituting for $a$ its value $\frac{1}{2\surd 3}$,
\[
21\theta^2 + 5\surd 3\theta - \frac{7}{4} = 0,
\]
\[
(14\theta - \surd 3)(\theta + \frac{5}{14\surd 3}) = \frac{116}{49},
\]
that is
\[
\theta = \frac{-5 \pm \surd 116}{14\surd 3},
\]
that is
\[
\theta = -0.65034, \quad \theta = +0.23797,
\]
giving
\[
\cot c = +0.93902, \quad \cot c = +0.05071,
\]
or
\[
c = 46°\, 48', \quad c = 87°\, 6'.
\]
\end{art}

\begin{art}[No.~77]
It has been seen that $c = 73°\, 54' + \epsilon$ gives $A = -\frac{\surd 13}{4} = -2.30940$; there will be between $0°$ and $-60°$ another value of $c$, giving for $A$ this same value; to find this value write $A = -\frac{\surd 13}{4}$, then we have
\[
\frac{-\theta + \surd(1 + \theta^2)}{3} = -\frac{\surd 13}{4},
\]
that is
\[
(1 + 2\surd 3)\theta + 1 = \surd(1 + \theta^2),
\]
or
\[
(12 + 4\surd 3)\theta^2 + (2 + 4\surd 3)\theta = 0,
\]
satisfied as it should be by $\theta = 0$, and also by
\[
\theta = -\frac{2(3 + \surd 3)}{12 + 4\surd 3} = -\frac{1}{2},
\]
giving
\[
\cot c = -0.76038 \text{ or } c = 52°\, 45'.
\]
\end{art}

\begin{art}[No.~78]
Representing the equation of the orbit by
\[
r = Ax \pm a(1 - A^2),
\]
we have for the point $1$,
\[
1 = A \pm a(1 - A^2),
\]
that is
\[
a = \frac{\pm 1}{1 \pm A},
\]
where the sign is to be taken so that $a$ shall be positive.
\end{art}

\begin{art}[No.~79]
With a view to the calculation of the times of passage, I calculate a series of values of $x_2$, $y_2$, $r_2$, $A$, $a$, for values of $c$ at the intervals of $5°$ and for a few intermediate values; we have $x_3$, $y_3$, $r_3 = x_2$, $-y_2$, $r_2$, so that these are known; so long as the orbit is an ellipse, the time of passage between the points $2$ and $3$, say $T_{23}$, may be calculated by Lambert's equation, the length of the chord $y_2 - y_3 = 2y_2$ being known without any fresh calculation. And then the times $T_{12}$ and $T_{31}$ being equal, and the sum $T_{12} + T_{23} + T_{31}$ being equal to the whole periodic time (reckoned as $= 3a^{\frac{3}{2}}$) the times $T_{12}$ and $T_{31}$ are also known. But when the orbit is a concave hyperbola there is no time $T_{23}$, and the other two times $T_{12} = T_{31}$ must be calculated. For the reason referred to (ante, No.~39) I did not use Lambert's equation, and it was less necessary to do so, by reason that, the transverse axis coinciding with the axis of $x$, the other formula could be employed without difficulty.
\end{art}

\begin{art}[No.~80]
The formulae for $x_2$, $y_2$ adapted to logarithmic calculation are
\begin{align*}
\log(x_2 + 0.61539) &= 11.60174 + \log \tan(c + 16°\, 6'),\\
\log y_2 &= 10.92015 + \log \sec(c + 16°\, 6'),
\end{align*}
where $y_2$ is always positive, but the sign of $x_2$ must be attended to. The values of $r_2$ and its inclination $\phi_2$ to the axis of $x$ are then to be calculated from
\[
\tan \phi_2 = \frac{y_2}{x_2}; \quad r_2 = x_2 \sec \phi_2 \text{ or } = y_2 \, \text{cosec} \, \phi_2,
\]
(viz. for $r_2$ it is proper to use the first or the second value, according as $x_2$ is greater or less than $y_2$).
\end{art}

\begin{art}[No.~81]
We have then $e = (\pm A)$ and $a$ from the foregoing formulae
\[
x_2 = \frac{r_2 - 1}{\pm 1}, \quad a = \frac{\pm 1}{-1 + A},
\]
where $a$, $e$ are each of them positive.

And then for the Times
\[
T = \frac{a^{\frac{3}{2}}}{\pi}(\chi - \sin\chi - \lambda + \sin\lambda); \quad (\log \tfrac{1}{\pi} = 1.67894)
\]
where
\[
a \cos\chi = a - r_2 - y_2, \quad a \cos\lambda = a - r_2 + y_2,
\]
and attention is necessary in order to the selection of the proper values of the angles $\chi$, $\lambda$.

And finally
\[
T_{12} = T_{31} = \tfrac{1}{2}(3a^{\frac{3}{2}} - T_{23}).
\]
\end{art}

\begin{art}[No.~82]
For the Time in a hyperbola, we have
\[
T_{12} = T_{31} = \frac{a^{\frac{3}{2}}}{\pi}\{e\tan u_2 + \log \tan(45° + \tfrac{1}{2}u_2)\},
\]
where
\[
\tan u_2 = \frac{\surd(e^2 - 1)\, y_2}{a(e^2 - 1) + r_2}.
\]
\end{art}


\newpage
\section*{Planogram No.~1, first part, $b = 90°$}

\begin{center}
\small
\begin{tabular}{|l|c|c|c|c|c|c|c|c|c|c|}
\hline
 & $c$ & $x_2$ & $y_2$ & $\phi_2$ & $r_2$ & $A$ & $a$ & $T_{23}$ & $T_{12}$ & $T_{31}$ \\
\hline
Circle & $0°$ & $-.500$ & $+.866$ & $60°$ & $1.000$ & $1.000$ & $1.000$ & $1.000$ & $1.000$ & $1.000$ \\
 & $5°$ & $-.461$ & $.892$ & $62°\,39'$ & $1.004$ & $-.003$ & $1.003$ & & & \\
 & $10°$ & $-.420$ & $.927$ & $65°\,38'$ & $1.017$ & $-.012$ & $1.012$ & $.960$ & $1.135$ & $.960$ \\
 & $15°$ & $-.374$ & $.972$ & $68°\,56'$ & $1.041$ & $-.030$ & $1.031$ & & & \\
 & $20°$ & $-.324$ & $1.030$ & $72°\,33'$ & $1.079$ & $-.060$ & $1.064$ & $.922$ & $1.448$ & $.922$ \\
 & $25°$ & $-.267$ & $1.104$ & $76°\,26'$ & $1.136$ & $-.107$ & $1.120$ & & & \\
\hline
Ellipses & $30°$ & $-.200$ & $1.200$ & $80°\,32'$ & $1.216$ & $-.180$ & $1.220$ & $.887$ & $2.275$ & $.887$ \\
 & $35°$ & $-.120$ & $1.325$ & $84°\,49'$ & $1.330$ & $-.295$ & $1.418$ & & & \\
 & $40°$ & $-.021$ & $1.492$ & $90°\,48'$ & $1.492$ & $-.482$ & $1.931$ & $.838$ & $6.371$ & $.838$ \\
 & $40°\,54'$ & $.000$ & $1.515$ & $90°$ & $1.515$ & $-.515$ & $2.061$ & & & \\
 & $45°$ & $+.100$ & $1.722$ & $93°\,37'$ & $1.725$ & $-.814$ & $5.362$ & & & \\
\hline
Parab. & $46°\,48'$ & $.166$ & $1.826$ & $95°\,11'$ & $1.834$ & $-1.000$ & $\infty$ & $.787$ & $\infty$ & $.787$ \\
\hline
Hyperbs. & $50°$ & $.287$ & $2.054$ & $97°\,57'$ & $2.074$ & $-1.505$ & $1.981$ & $.750$ & $\sim$ & $.750$ \\
 & $52°\,45'$ & $.418$ & $2.306$ & $100°\,16'$ & $2.344$ & $-2.309$ & $.764$ & $\sim$ & & \\
 & $55°$ & $.552$ & $2.569$ & $102°\,8'$ & $2.627$ & $-3.632$ & $.380$ & $.628$ & $\sim$ & $.628$ \\
\hline
Line & $60°$ & $1.000$ & $3.464$ & $106°\,6'$ & $3.605$ & $-.000$ & $-.000$ & $\sim$ & $\sim$ & $\sim$ \\
\hline
Convex & $65°$ & $1.937$ & $5.378$ & $109°\,48'$ & $5.716$ & $+5.032$ & $.166$ & $\sim$ & & \\
Hyperbs. & $70°$ & $5.248$ & $12.233$ & $113°\,13'$ & $13.311$ & $+2.898$ & $.257$ & & & \\
 & $73°\,54'$ & $\infty$ & $\infty$ & $115°\,39'$ & $\infty$ & $+2.309$ & $.302$ & $\sim$ & $\sim$ & $\sim$ \\
 & $76°$ & $-21.432$ & $43.341$ & $63°\,42'$ & $48.346$ & $2.111$ & $.900$ & $20.68$ & $\sim$ & $20.68$ \\
 & $80°$ & $4.356$ & $7.830$ & $60°\,55'$ & $8.960$ & $1.486$ & $2.056$ & $3.856$ & $\sim$ & $3.856$ \\
 & $85°$ & $2.653$ & $4.322$ & $58°\,27'$ & $5.072$ & $1.415$ & $8.718$ & $2.912$ & $\sim$ & $2.912$ \\
\hline
Parab. & $87°\,6'$ & $2.320$ & $3.644$ & $57°\,31'$ & $4.320$ & $1.000$ & $\infty$ & $2.588$ & $\infty$ & $2.588$ \\
\hline
Ellipse & $90°$ & $2.000$ & $3.000$ & $56°\,18'$ & $3.606$ & $-.869$ & $7.622$ & $2.255$ & $58.62$ & $2.255$ \\
\hline
\end{tabular}
\end{center}

The mark $\sim$ in the $T_{23}$ column shows that there is no Time $T_{23}$.

\newpage
\section*{Planogram No.~1, second part, $b = 270°$}

\begin{center}
\small
\begin{tabular}{|l|c|c|c|c|c|c|c|c|c|c|}
\hline
 & $c$ & $x_2$ & $y_2$ & $\phi_2$ & $r_2$ & $A$ & $a$ & $T_{23}$ & $T_{12}$ & $T_{31}$ \\
\hline
Circ. & $0°$ & $-.500$ & $+.866$ & $60°$ & $1.000$ & $1.000$ & $1.000$ & $1.000$ & $1.000$ & $1.000$ \\
 & $5°$ & $-.537$ & $.848$ & $57°\,39'$ & $1.003$ & $-.002$ & $1.002$ & & & \\
 & $10°$ & $-.573$ & $.837$ & $55°\,37'$ & $1.014$ & $-.009$ & $1.009$ & $1.044$ & $.951$ & $1.044$ \\
 & $15°$ & $-.608$ & $.832$ & $53°\,52'$ & $1.030$ & $-.019$ & $1.019$ & & & \\
 & $20°$ & $-.643$ & $.834$ & $52°\,23'$ & $1.053$ & $-.032$ & $1.033$ & $1.091$ & $.969$ & $1.091$ \\
 & $25°$ & $-.678$ & $.842$ & $51°\,10'$ & $1.081$ & $-.048$ & $1.051$ & & & \\
 & $30°$ & $-.714$ & $.857$ & $50°\,11'$ & $1.116$ & $-.068$ & $1.073$ & $1.145$ & $1.043$ & $1.145$ \\
 & $35°$ & $-.752$ & $.879$ & $49°\,27'$ & $1.157$ & $-.090$ & $1.098$ & & & \\
 & $40°$ & $-.793$ & $.910$ & $48°\,51'$ & $1.207$ & $-.115$ & $1.130$ & $1.207$ & $1.192$ & $1.207$ \\
 & $45°$ & $-.836$ & $.950$ & $48°\,40'$ & $1.266$ & $-.145$ & $1.169$ & & & \\
 & $50°$ & $-.884$ & $1.002$ & $48°\,36'$ & $1.336$ & $-.179$ & $1.217$ & $1.283$ & $1.464$ & $1.283$ \\
 & $55°$ & $-.938$ & $1.069$ & $48°\,45'$ & $1.442$ & $-.218$ & $1.278$ & & & \\
 & $60°$ & $-1.000$ & $1.154$ & $49°\,7'$ & $1.527$ & $-.264$ & $1.358$ & $1.377$ & $1.983$ & $1.377$ \\
 & $65°$ & $-1.074$ & $1.266$ & $49°\,42'$ & $1.660$ & $-.318$ & $1.466$ & & & \\
 & $70°$ & $-1.164$ & $1.412$ & $50°\,31'$ & $1.830$ & $-.383$ & $1.622$ & $1.506$ & $3.036$ & $1.506$ \\
 & $75°$ & $-1.280$ & $1.611$ & $51°\,35'$ & $2.056$ & $-.464$ & $1.864$ & & & \\
 & $80°$ & $-1.431$ & $1.891$ & $52°\,53'$ & $2.372$ & $-.564$ & $2.295$ & $1.771$ & $6.888$ & $1.771$ \\
 & $85°$ & $-1.651$ & $2.311$ & $54°\,28'$ & $2.840$ & $-.694$ & $3.269$ & & & \\
 & $90°$ & $-2.000$ & $3.000$ & $56°\,18'$ & $3.606$ & $-.869$ & $7.622$ & $2.255$ & $58.62$ & $2.255$ \\
\hline
\end{tabular}
\end{center}


\newpage
\noindent\textbf{Article Nos.~84 to 94. Planogram No.~2, the Meridian $0°$--$180°$} (see Plate~III.).

\begin{art}[No.~84]
The orbit-plane here rotates about an axis in the plane of the ecliptic at right angles to $S1$ (Fig.~6). The entire meridian is given by $b = 0°$, $c = 0°$ to $90°$, and $b = 180°$, $c = 0°$ to $90°$, but it is sufficient to consider one of these half meridians, say the latter of them, as the series of values is the same for each of them, with only an interchange of the points $2$, $3$. I write therefore, $b = 180°$, so that we have
\begin{align*}
\alpha, \beta, \gamma &= 0, -1, 0,\\
\alpha', \beta', \gamma' &= \cos c, 0, \sin c,\\
\alpha'', \beta'', \gamma'' &= -\sin c, 0, \cos c,
\end{align*}
consequently
\begin{align*}
x_1' : y_1' : 1 &= \sin c : -\surd 3 : -\surd 3 \cos c,\\
x_2' : y_2' : 1 &= -3\cos c - 2\sin c : -\surd 3 : -\surd 3 \sin c - 2\surd 3 \cos c,\\
x_3' : y_3' : 1 &= 3\cos c - 2\sin c : -\surd 3 : -\surd 3 \sin c - 2\surd 3 \cos c,
\end{align*}
and moreover $x = y'$, $y' = -x$; so that, introducing into the formulae $(x_1, y_1)$, \&c., in place of the $(x_1', y_1')$, \&c., we have
\begin{align*}
x_1 &= \sec c, \quad y_1 = -\surd 3 \tan c,\\
x_2 &= -\frac{1}{\surd 3} \frac{2\sin c + 3\cos c}{\sin c + 2\cos c}, \quad y_2 = \frac{1}{\surd 3} \frac{-2 + 3\tan c}{\tan c + 2},\\
x_3 &= -\frac{1}{\surd 3} \frac{2\sin c - 3\cos c}{\sin c - 2\cos c}, \quad y_3 = \frac{1}{\surd 3} \frac{-2 - 3\tan c}{\tan c - 2},
\end{align*}
which, putting
\[
\cos\delta = \frac{2}{\surd 5}, \quad \sin\delta = \frac{1}{\surd 5}, \quad \tan\delta = \frac{1}{2}, \quad \delta = 26°\, 34',
\]
become
\begin{align*}
x_1 &= \sec c, \quad y_1 = -\surd 3 \tan c,\\
x_2 &= -\frac{1}{\surd 3} \sec(c - \delta), \quad y_2 = \frac{1}{\surd 3}(-2 + 3\tan(c - \delta)),\\
x_3 &= -\frac{1}{\surd 3} \sec(c + \delta), \quad y_3 = \frac{1}{\surd 3}(-2 + 3\tan(c + \delta)),
\end{align*}
so that the guide-hyperbolas are
\begin{align*}
x_1^{\,2} - 3y_1^{\,2} &= 1, \quad \text{angle of asymptotes } = 30°,\\
x_2^{\,2} &= 15y_2^{\,2} - 16\surd 3\, y_2 + 13, \quad \tan^{-1} \surd\frac{2}{7} = 14°\, 28',\\
x_3^{\,2} &= 15y_3^{\,2} + 16\surd 3\, y_3 + 13.
\end{align*}
\end{art}

\begin{art}[No.~85]
It is easy to verify that hyperbola 2 passes through $x_2 = -\frac{1}{2}$, $y_2 = \frac{1}{2}\surd 3$, and touches there circle $x^{\,2} + y^{\,2} = 1$. And we thus have the figure in the Plate.
\end{art}

\begin{art}[No.~86]
The figure shows the motion of the points $1$, $2$, $3$, along their respective hyperbolas, viz. $c = 0°$ to $90°$, the point $1$ moves from contact with the circle, along a half branch to infinity: $2$ moves from contact along a small portion of the half branch; $3$ moves from contact, along the half branch to infinity for $c = \tan^{-1} 2 = 63°\, 26'$, and then reappearing at the opposite infinity, as $c$ increases to $90°$ describes a portion of the opposite half branch.
\end{art}

\begin{art}[No.~87]
For $c = 0$, the orbit is the circle; as $c$ increases the orbit becomes elliptic then parabolic, $c = 51°$, and afterwards hyperbolic (concave); until for $c = 60°$, the three points are on the horizontal line of the figure, and the orbit is this right line; it is to be noticed that the arrangement of the points on these orbits is $1$, $2$, $3$; so that for the parabola, $T_{12}$ is $= \infty$, and for the hyperbolas and right line $T_{23}$ does not exist.
\end{art}

\begin{art}[No.~88]
For $c < 60°$ until $c = 63°\, 26'$ the orbit is a convex hyperbola, the arrangement of the points being still $1$, $2$, $3$: say for $c = 63°\, 26' - \epsilon$, the orbit is the convex hyperbola $H_1$. At $c = 63°\, 26'$ there is an abrupt change of orbit; say for $c = 63°\, 26' + \epsilon$ the orbit is a concave hyperbola $H_2$; and for $c = 65°\, 52'$ the orbit is a parabola; the arrangement of the points on these orbits is $2$, $1$, $3$; so that for the hyperbolas $T_{13}$ does not exist.
\end{art}


\begin{art}[No.~89]
The equation of the orbit is found to be
\begin{multline*}
r\{-2r_1 \sin c \cos c + r_2(-\sin^2 c + \sin c \cos c + 6\cos^2 c)\\
+ r_3(\sin^2 c + \sin c \cos c - 6\cos^2 c)\}\\
= \frac{-r_1}{\surd 3} \cdot 6\cos^2 c + 3r_2(\sin^2 c + \sin c \cos c - 2\cos^2 c)\\
+ 3r_3(\sin^2 c - \sin c \cos c - 2\cos^2 c);
\end{multline*}
(observe that the orbit will be a right line if $\sin^2 c - 3\cos^2 c = 0$, that is if $c = 60°$, which is right, since $60°$ is the angular radius of the regulator circle).
\end{art}

\begin{art}[No.~90]
Putting in the equation $\tan c = \lambda$, and therefore $\cos c = \frac{1}{\surd(1 + \lambda^2)}$, the equation becomes
\begin{multline*}
r = \frac{-1}{2\surd 3(\lambda^2 - 3)}\{\lambda(2R_1 + R_2 - R_3) - 3(R_2 + R_3)\}x\\
+ \frac{1}{2(\lambda^2 - 3)}\{4R_1 + R_2(1 - \lambda) + R_3(1 + \lambda)\},
\end{multline*}
where, as above,
\[
r_1 = R_1, \quad (\lambda + \surd 3)r_2 = R_2, \quad (\lambda - \surd 3)r_3 = R_3,
\]
and thence also
\[
R_1 = \surd(1 + \lambda^2), \quad R_2 = \surd(7\lambda^2 + 12\lambda + 12), \quad R_3 = \surd(7\lambda^2 - 12\lambda + 12).
\]
\end{art}

\begin{art}[No.~91]
The analytical expression for the eccentricity is
\[
e = \surd(A^2 + B^2),
\]
where, as above,
\[
A = \frac{-1}{6\surd(1 + \lambda^2)}\{\lambda(2R_1 + R_2 - R_3) - 3(R_2 + R_3)\},
\]
\[
B = \frac{1}{2\surd 3(\lambda^2 - 3)}\{4R_1 + R_2(1 - \lambda) + R_3(1 + \lambda)\};
\]
but this expression is too complicated to allow of an analytical discussion of the series of values of $e$ (such as was given for $A$, $= +e$, in Planogram No.~1). The numerical calculation gives the results mentioned ante No.~87, viz., $c = 0$, $e = 0$; $c = 51°$, $e = 1$; $c = 60°$, $e = \infty$; $c = 63°\, 26' - \epsilon$, $e = 4.912$; $c = 63°\, 26' + \epsilon$, $e = 1.853$; $c = 69°$, $e = 1.628$ (min.); $c = 89°\, 20'$ (viz. $\lambda = 86.176$), $e = 1$; $c = 90°$, $e = 1.018$; values which are exhibited in the diagram in the preceding page.
\end{art}

\begin{art}[No.~92]
It may be further remarked, in reference to the formula
\[
r = Ax + By + C,
\]
that for $c = 60°$, that is $\lambda = \surd 3$, we have $A$ finite, $B$ and $C$ each infinite, but equal and of opposite signs; viz., the equation becomes $r = -2.242x + \infty(y - 1)$, that is $y = 1$, orbit a right line as above.

The abrupt change at $c = 63°\, 26'$, $\lambda = 2$, arises from the change of sign of $R_3$, viz., $c = 63°\, 26' - \epsilon$, $R_3 = -2.309$, but $c = 63°\, 26' + \epsilon$, $R_3 = +2.309$; the two orbits are
\begin{align*}
c &= 63°\, 26' - \epsilon, \quad r = -2.34x + 4.906y - 3.671, \quad e = 4.912, \quad a = -0.159,\\
c &= 63°\, 26' + \epsilon, \quad r = -5.78x - 17.61y + 32.57, \quad e = 1.853, \quad a = 1.338.
\end{align*}

For $c = 90°$ the equation is
\[
r = -0.770x + 0.666y + 1.527
\]
and therefore $e = \surd(A^2 + B^2) = 1.018$ as above; $a = 9\surd 21 = 41.243$.

It is to be added that for $c$ nearly $= 90°$, or $\lambda$ very large, we have
\[
R_1 = \lambda, \quad R_2 = \surd 7\lambda + \frac{6}{\surd 7}, \quad R_3 = \surd 7\lambda - \frac{6}{\surd 7},
\]
and thence
\[
A = -\frac{1}{3\surd 3} = -0.770 - \frac{0.430}{\lambda}, \quad B = \frac{1}{\surd 7} = 0.666 + \frac{1.890}{\lambda}, \quad C = \frac{4}{\surd 7} = 1.527 - \frac{1.555}{\lambda}.
\]
It was, in fact, by means of these expressions that the value $\lambda = 86.176$ ($c = 89°\, 20'$) corresponding to the last-mentioned parabola $e = 1$ was obtained.
\end{art}

\begin{art}[No.~93]
For the calculation of the table we have
\begin{align*}
\log x_1 &= 10 + \log \sec c,\\
\log y_1 &= 10.76144 + \log \tan c,\\
\log x_2 &= 10.65052 + \log \sec(c - 26°\, 34'),\\
\log(y_2 - 0.92376) &= 10.06247 + \log \tan(c - 26°\, 34'),\\
\log x_3 &= 10.65052 + \log \sec(c + 26°\, 34'),\\
\log(y_3 + 0.92376) &= 10.06247 + \log \sec(c + 26°\, 34'),
\end{align*}
the values of $r_1$, $r_2$, $r_3$ are then calculated from
\[
x_1 = r_1 \cos\phi_1, \quad y_1 = r_1 \sin\phi_1,
\]
or say
\[
\tan\phi_1 = \frac{y_1}{x_1}, \quad r_1 = x_1 \sec\phi_1, \quad \&\text{c.},
\]
and those of the chords $c_{12}$, $c_{23}$, $c_{31}$ from
\[
x_1 - x_2 = c_{12}\cos\theta_{12}, \quad y_1 - y_2 = c_{12}\sin\theta_{12},
\]
or say
\[
\tan\theta_{12} = \frac{y_1 - y_2}{x_1 - x_2}, \quad c_{12} = (x_1 - x_2)\sec\theta_{12}, \quad \&\text{c.}
\]
\end{art}


\newpage
\noindent\textbf{Article Nos.~95 to 98. Planogram No.~3, the Orbit-pole at one of the points $A$.}

\begin{art}[No.~95]
When the orbit-pole is at one of the points $A$, the orbit-plane passes through one of the rays, and as there is no longer on this ray any determinate point of intersection, the orbit (as was seen) becomes indeterminate. Thus consider the point $A$ for which $b = 270°$, $c = 60°$: we have
\begin{align*}
\alpha, \beta, \gamma &= 0, -1, 0,\\
\alpha', \beta', \gamma' &= 0, -\tfrac{1}{2}, -\tfrac{1}{2}\surd 3,\\
\alpha'', \beta'', \gamma'' &= 0, \tfrac{1}{2}\surd 3, -\tfrac{1}{2},
\end{align*}
and consequently the formula gives
\begin{align*}
x_1 : y_1 : 1 &= -\tfrac{1}{2}\surd 3 - \surd 3 : -\tfrac{1}{2} - \surd 3,\\
x_3 : y_3 : 1 &= -\tfrac{1}{2}\surd 3 - \surd 3 : -3 - \tfrac{1}{2} - \surd 3,
\end{align*}
and, moreover, $x = -x'$, $y = -y'$. From the formula the value of $x_1'$ or $x_1$ is given as $\frac{1}{2}$, but the true value is obviously $x_1 = 1$; the value of $y_1$ is actually indeterminate.

The formulae give the values of $(x_2, y_2)$, $(x_3, y_3)$, viz. the system is
\begin{align*}
x_1 &= 1, \quad y_1 \text{ ind.}\\
x_2 &= \tfrac{1}{2}, \quad y_2 = \tfrac{2}{\surd 3}, \text{ whence } r_2 = r_3 = \surd\tfrac{7}{3},\\
x_3 &= -\tfrac{1}{2}, \quad y_3 = -\tfrac{2}{\surd 3}.
\end{align*}
\end{art}

\begin{art}[No.~96]
So that the orbits in the planogram are the whole series of conics having a given focus, $S$, and passing through two fixed points, $2$, $3$, having the common abscissa $x = -\frac{1}{2}$, and at equal distances $\frac{2}{\surd 3}$ ($= 1.15470$) on opposite sides of the axis. The axis of $X$ is obviously the common transverse axis for all the orbits; that is, the equation of the orbit will be of the form $r = Ax + B$; and writing $a = -1$, we have $r_1 = -A + B$, viz. the equation is $r - \surd\frac{7}{3} = A(x + 1)$; the value of $A$ will be determined if we assume for the point $1$ a determinate position on the line $x = 1$, say its ordinate is $= y_1$; for then if $r_1 = \surd(1 + y_1^{\,2})$ we have $r_1 - \surd\frac{7}{3} = 2A$, and the equation is $r - \surd\frac{7}{3} = \frac{1}{2}(r_1 - \surd\frac{7}{3})(x + 1)$.

In particular if $y_1 = 0$, we have $r_1 = 1$, and the equation of the orbit is $r - \surd\frac{7}{3} = \frac{1}{2}(1 - \surd\frac{7}{3})(x + 1)$: this is the orbit, eccentricity $\frac{1}{2}(\surd\frac{7}{3} - 1) = 0.264$, belonging to the point $A$ as a point in Planogram No.~1: for the value of $y_1$ being in that planogram originally assumed $= 0$, is of course $= 0$ when the orbit-pole comes to be the point $A$.
\end{art}

\begin{art}[No.~97]
We may conversely take the equation of the orbit, or say the value of $A (= \pm e)$ in the equation $r - \surd\frac{7}{3} = A(x + 1)$, to be given; and then writing $y_1 = \tan\psi$, we have $r_1 = \sec\psi$, and the equation becomes
\[
\sec\psi - \surd\tfrac{7}{3} = 2A, \quad \text{that is } \sec\psi = 2A + \surd\tfrac{7}{3},
\]
which determines $\psi$, and thence $y_1 = \tan\psi$; that is, we have $y_1 = \surd\{(2A + \surd\frac{7}{3})^2 - 1\}$; and observing that $\surd\frac{7}{3} = 1.527$, the values are
\begin{align*}
A &= 0, & \sec\psi &= 1.527, & \psi &= \pm 49°\, 6',\\
A &= \tfrac{1}{2}(1 - \surd\tfrac{7}{3}) = -0.264, & \sec\psi &= 1.000, & \psi &= 0°,\\
A &= -\tfrac{1}{2}(1 + \surd\tfrac{7}{3}) = -1.264, & \sec\psi &= 0.264, & \psi \text{ imaginary}.
\end{align*}
\end{art}

\begin{art}[No.~98]
Writing for greater convenience $\xi = \rho\cos\psi$, $\eta = \rho\sin\psi$, the indefinitely small quantity $\rho$ will denote the distance of the orbit-pole from $A$, and its azimuth measured from the meridian will be $\psi$. We then have $y_1 = -\tan\psi$, and $r_1 = \surd(1 + y_1^{\,2}) = +\sec\psi$, or, if to fix the ideas, $\psi$ be considered as $< + 90°$, then $r_1 = \sec\psi$: we have thus ($A = \pm e$ as before) $A = \frac{1}{2}(-\surd\frac{7}{3} + \sec\psi)$; viz., observing that $\surd\frac{7}{3} = 1.527$, we obtain
\begin{align*}
\psi &= 0°, & A &= -\tfrac{1}{2}(\surd\tfrac{7}{3} - 1) = -0.264,\\
\psi &= \sec^{-1}\surd\tfrac{7}{3} = \pm 49°\, 6', & A &= 0,\\
\psi &= \sec^{-1}(2\surd\tfrac{7}{3} - 1) = \pm 60°\, 52', & A &= \tfrac{1}{2}(\surd\tfrac{7}{3} - 1) = +0.264,\\
\psi &= 90°, & A &= \tfrac{1}{2}\surd\tfrac{7}{3} = +0.764.
\end{align*}
\end{art}

\newpage
\noindent\textbf{Article Nos.~99 to 103. Planogram No.~4, the Orbit-pole in the Ecliptic.}

\begin{art}[No.~99]
When the orbit-pole describes the circle of the ecliptic, the orbit-plane passes through the axis of $z$, or polar axis. We have $c = 90°$, and consequently
\begin{align*}
\alpha, \beta, \gamma &= \sin b, \cos b, 0,\\
\alpha', \beta', \gamma' &= 0, 0, -1,\\
\alpha'', \beta'', \gamma'' &= \cos b, \sin b, 0.
\end{align*}
Reverting for a moment to the general case where the six coordinates of the ray are $(a, b, c, f, g, h)$, the formulae for the intersection by the orbit-plane are
\begin{align*}
x' : y' : 1 &= -(a, b, c \,\raisebox{.5ex}{$\scriptscriptstyle\diagup$}\, \alpha', \beta', \gamma') : -(a, b, c \,\raisebox{.5ex}{$\scriptscriptstyle\diagup$}\, \alpha'', \beta'', \gamma'') : -(a\sin b + b\cos b)\\
&: (f, g, h \,\raisebox{.5ex}{$\scriptscriptstyle\diagup$}\, \alpha'', \beta'', \gamma'') : f\cos b + g\sin b,
\end{align*}
that is if
\[
-cx' + hy' + a = 0, \quad -cy' + gx' + b = 0,
\]
and thence
\begin{align*}
x' : y' : 1 &= af - bg : -(ag + bf) : ch\\
&= af - bg : -(ag + bf) : -(ag - bf)y' + ch,
\end{align*}
consequently
\[
hx' = -(ag - bf)y' + ch, \quad \text{or } h\,\frac{af - bg}{ch} = -(ag - bf)y' + ch,
\]
or, what is the same thing,
\[
h^2(af - bg) = \{(ag - bf)^2 + h^2\}y' + (ag - bf)ch,
\]
or, in particular, if (as in the special symmetrical case) $ag - bf = 0$, then
\[
h^2(af - bg) = (f^2 + g^2)y^{\,2} + c^2 + h^2.
\]
\end{art}


\begin{art}[No.~100]
For the symmetrical system of rays we have as before
\begin{align*}
(a_1, b_1, c_1, f_1, g_1, h_1) &= (0, \surd 3, -1, 0, 1, \surd 3),\\
(a_2, b_2, c_2, f_2, g_2, h_2) &= (3, \surd 3, 2, \surd 3, 1, -2\surd 3),\\
(a_3, b_3, c_3, f_3, g_3, h_3) &= (-3, \surd 3, 2, -\surd 3, 1, -2\surd 3),
\end{align*}
and the orbit is found to be
\[
r = \tfrac{1}{3}\surd 7\, x' + \tfrac{2}{3}\surd 2\, y' + \tfrac{1}{3}\surd 21 = -0.666x + 0.770y + 1.527;
\]
and similarly for $b = -\epsilon$ the equation is
\[
r = -\tfrac{1}{3}\surd 7\, x' + \tfrac{2}{3}\surd 2\, y' + \tfrac{1}{3}\surd 21 = 0.666x + 0.770y + 1.527;
\]
hence the eccentricity is
\[
e = \surd(A^2 + B^2) = 1.018, \text{ as before}.
\]
\end{art}

\begin{art}[No.~101]
Considering now the general case where $b$ has any value whatever, the equation of the orbit is
\[
r = \frac{1}{6\surd 3}\{2R_1\sin b - R_2(\sin b + \surd 3\cos b) - R_3(\sin b - \surd 3\cos b)\},
\]
where
\[
\surd 3\, r_1 = R_1, \quad (\surd 3 + \eta)r_2 = R_2, \quad (\surd 3 - \eta)r_3 = R_3,
\]
and therefore
\[
\surd 3\, r_1 = R_1, \quad (\surd 3 + \eta)r_2 = R_2, \quad (\surd 3 - \eta)r_3 = R_3,
\]
which last equations determine the signs of $R_1$, $R_2$, $R_3$ respectively, the equation of the orbit is
\[
r = \frac{1}{6\surd(1 + \eta^2)}\{2R_1 - R_2 - R_3\}x'
\]
\[
+ \frac{1}{6(1 + \eta^2)}\{4R_1 + R_2(1 - \eta\surd 3) + R_3(1 + \eta\surd 3)\}y'
\]
\[
+ \frac{1}{6\surd(1 + \eta^2)}\{2R_1 + R_2 + R_3\}.
\]
Thus if $b = +\epsilon$, then also $\eta = +\epsilon$,
\[
r = \tfrac{1}{3}\surd 7\, x' + \tfrac{2}{3}\surd 2\, y' + \tfrac{1}{3}\surd 21 = -0.666x + 0.770y + 1.527,
\]
as before; and similarly if $b = 90°$.

And moreover, if $b = 30°$, then
\[
R_1 = \frac{2\surd 13}{3}, \quad R_2 = R_3 = \frac{2}{3},
\]
whence the equation of the orbit is
\[
r = \tfrac{1}{3}(\surd 13 - 1)x' + 0y' + \tfrac{1}{3}(\surd 13 + 2) = 0.868x' + 0y' + 1.868.
\]
\end{art}

\begin{art}[No.~102]
The equation of the orbit should be tabulated from $b = 0$ to $b = 30°$, the equations for the remainder of the circumference will be then found by successive repetition of this interval in direct and reverse order.

As regards the interval $0°$ to $30°$ the only intermediate value that I have calculated is $b = 15°$, viz., we then have
\[
b = 15°, \quad r = -0.811x' + 0.403y' + 1.787.
\]
Calculating for the foregoing values $b = 0°$, $b = 15°$, $b = 30°$, the values of $e$, $\varpi$, $a$, these are found to be
\begin{align*}
b &= 0°, & e &= 1.018, & \varpi &= 220°\, 6', & a &= 41.24,\\
b &= 15°, & e &= 0.906, & \varpi &= 206°\, 27', & a &= 10.008,\\
b &= 30°, & e &= 0.868, & \varpi &= 180°, & a &= 7.604.
\end{align*}
\end{art}

\newpage
\noindent\textbf{Article Nos.~104 to 113. Planogram No.~5. The Orbit-pole on a Separator.}

\begin{art}[No.~104]
If the orbit-plane rotate round a line parallel to one of the rays, the orbit-pole will describe a separator circle, and conversely. I consider the general case of a ray the six coordinates of which are $(a, b, c, f, g, h)$, and for which the intersections with the orbit-plane are given by
\begin{align*}
x' : y' : 1 &= (ag - bf)\cos H + c\sin H : 0\\
&= \frac{(ag - bf)\cos H + c\sin H}{n} : 0 : 1,
\end{align*}
where $H$ is the angular distance of the orbit-pole, along the separator, from the point $A$. The foregoing values give
\begin{align*}
(a, b, c \,\raisebox{.5ex}{$\scriptscriptstyle\diagup$}\, \alpha, \beta, \gamma) &= 0,\\
(a, b, c \,\raisebox{.5ex}{$\scriptscriptstyle\diagup$}\, \alpha', \beta', \gamma') &= -\tfrac{1}{2}n\{(ag - bf)\cos H + c\sin H\},\\
(f, g, h \,\raisebox{.5ex}{$\scriptscriptstyle\diagup$}\, \alpha'', \beta'', \gamma'') &= 0,
\end{align*}
so that the coordinates $x', y'$ of the intersection with the ray are given in the form
\[
x' : y' : 1 = \frac{M}{n} : 0 : 1,
\]
that is
\[
x' = \frac{M}{n}, \quad y' = 0,
\]
but the value of $y'$ is determinate, viz., this is equal to the perpendicular distance of the ray from the point $S$.
\end{art}

\begin{art}[No.~105]
In particular when the rays are the special symmetrical system before considered, then if $(a, b, c, f, g, h)$ refer to the ray $1$, we have $f = 0$, $g = 1$, $h = \surd 3$, $n = 1$, $\Omega = 2$, and thence
\begin{align*}
\alpha, \beta, \gamma &= 0, -\tfrac{1}{2}, -\tfrac{1}{2}\surd 3,\\
\alpha', \beta', \gamma' &= -\cos H, \tfrac{1}{2}\surd 3\sin H, -\tfrac{1}{2}\sin H,\\
\alpha'', \beta'', \gamma'' &= -\sin H, -\tfrac{1}{2}\surd 3\cos H, \tfrac{1}{2}\cos H.
\end{align*}
\end{art}


\begin{art}[No.~106]
For the intersection we obtain
\begin{align*}
x_1 &= -\frac{1}{\surd 3}\frac{6\cos H - \sin H}{3\cos H + 2\sin H}, &
y_1 &= \frac{3}{3\cos H + 2\sin H},\\
x_3 &= -\frac{1}{\surd 3}\frac{6\cos H + \sin H}{3\cos H - 2\sin H}, &
y_3 &= \frac{3}{3\cos H - 2\sin H}.
\end{align*}
\end{art}

\begin{art}[No.~107]
Writing herein
\[
\cos\omega = \frac{2}{\surd 13} = \frac{2}{3.606}, \quad \sin\omega = \frac{3}{\surd 13} = \frac{3}{3.606}, \quad \tan\omega = \frac{3}{2}, \quad \omega = 33°\, 41',
\]
the formulae are readily converted into
\begin{align*}
x_1 &= -\frac{1}{13\surd 3}\{16 - 15\tan(H - \omega)\}, &
x_3 &= -\frac{1}{13\surd 3}\{-16 - 15\tan(H + \omega)\},\\
y_1 &= \frac{3}{13}\surd 13 \sec(H - \omega), &
y_3 &= \frac{3}{13}\surd 13 \sec(H + \omega),
\end{align*}
where, in regard to this angle $\omega$, it is to be observed that it represents the angular distance from the ecliptic along the separator to a point $B$, or what is the same thing, the complement of the angular distance on the separator, of the points $A$ and $B$. We have, in fact, a right-angled spherical triangle $ZAB$, $\angle Z = 60°$, $\angle A = 90°$, $ZA = 60°$ whence $\sin 60° = \sin ZA \sin A$, that is $\sin A = 1$, $A = 90°$; and then $\sin AB = \sin 60° \sin ZA = \frac{3}{4}$, $AB = \omega = 33°\, 41'$.
\end{art}

\begin{art}[No.~108]
The equation of the orbit is found to be
\begin{multline*}
r\{-(x_2 - x_3)12\sin H = \surd\tfrac{1}{3}\{-36\cos^2 H + 4\sin^2 H\}r_1\\
- Or_2(9\cos^2 H - 4\sin^2 H) + Or_3(9\sin^2 H - 4\cos^2 H)\}\\
+ y\{-12\surd 3\cos H + 3r_2(3\cos H + 2\sin H)r_1 - 3r_3(3\cos H - 2\sin H)r_3\},
\end{multline*}
where
\begin{align*}
r_1 &= \surd(21\cos^2 H + 4\cos H\sin H + \tfrac{4}{3}\sin^2 H),\\
r_2 &= \frac{2}{3\cos H + 2\sin H}, \quad r_3 = \frac{2}{3\cos H - 2\sin H}.
\end{align*}
Hence, writing $\tan H = \lambda$, and therefore $\sec H = \surd(1 + \lambda^2)$, which determines the sign of $\surd(1 + \lambda^2)$, and moreover
\[
R_1 = \surd(21 + 4\lambda + \tfrac{4}{3}\lambda^2), \quad R_2 = \surd(21 + 4\lambda + \tfrac{4}{3}\lambda^2), \quad R_3 = \surd(21 - 4\lambda + \tfrac{4}{3}\lambda^2),
\]
and thence also
\[
(3 + 2\lambda)r_2 = R_2, \quad (3 - 2\lambda)r_3 = R_3,
\]
which last equations, since $r_2$, $r_3$ must be positive, determine the signs of the radicals $R_2$, $R_3$; the equation of the orbit is
\begin{multline*}
r = \frac{1}{12\lambda\surd(1 + \lambda^2)}\{(36 + 4\lambda^2) - (3 - 2\lambda)R_2 + (3 + 2\lambda)R_3\}x'\\
+ \frac{1}{12\surd 3(1 + \lambda^2)}\{4R_1 + R_2(1 + \lambda) + R_3(1 - \lambda)\}y',
\end{multline*}
where $\delta$ it will be recollected denotes $+1$ or $-1$ at pleasure.
\end{art}

\begin{art}[No.~109]
I remark that $\delta = +1$ and $\delta = -1$ extend from $0°$ to $90° - \omega$ ($= 56°\, 19'$, value at $B'$) and thence to $90°$, $\delta = +1$ belonging to the side marked $+$ in the figure, and $\delta = -1$ to the opposite side. But the same result may be stated, more conveniently, in reference to the blank spherogram, as follows:

$H = 0°$ to $H = 56°\, 19'$, $\delta = +1$ belongs to the outside of $AB'$, viz. to positions within the region of convex orbits, $\delta = -1$, to inside of $AB'$.

$H = 56°\, 19'$ to $H = 90°$, $\delta = +1$ belongs to inside of $B'C$.

$H = 90°$ to $H = 123°\, 41'$, $\delta = +1$ belongs to inside of $C'B$.

The last-mentioned values being identical with those for $H = 90°$ to $H = 56°\, 19'$, $\delta = -1$: viz. the formula for $H = 90° + K$, $\delta = +1$ is equivalent to that for $H = 90° - K$, $\delta = -1$.
\end{art}

\begin{art}[No.~110]
I consider some particular cases.

Orbit-pole at $A$: here $H = 0$ and therefore $\lambda = 0$, $R_1 = \surd 21$, $R_2 = R_3 = \surd 21$, and the equation is
\[
r = \tfrac{1}{3}(\surd 21 - 1)x' + 0y' + \tfrac{1}{3}(\surd 21 + 2) = 0.868x' + 0y' + 1.868,
\]
as above.

Suppose again that the orbit-pole is on the ecliptic, or say $H = 90° - \epsilon$, $\lambda = +\infty$, $R_2 = 2\surd\frac{7}{3}\lambda$, $R_3 = -2\surd\frac{7}{3}\lambda$, and $\surd(1 + \lambda^2) = \lambda$, and the equation is
\[
r = -\tfrac{1}{3}\surd 7\, x' + \tfrac{2}{3}\surd 2\, y' + \tfrac{1}{3}\surd 21,
\]
and similarly for $H = 90° + \epsilon$, $\lambda = -\infty$, $R_2 = 2\surd\frac{7}{3}\lambda$, $R_3 = -2\surd\frac{7}{3}\lambda$, $\surd(1 + \lambda^2) = -\lambda$, and the equation still is
\[
r = -\tfrac{1}{3}\surd 7\, x' + \tfrac{2}{3}\surd 2\, y' + \tfrac{1}{3}\surd 21;
\]
viz. retaining the same sign, there is no discontinuity in the passage through $90°$.

The eccentricity, whether $\delta = +1$ or $= -1$, is $\surd\frac{13}{7} = 1.018$, agreeing with Planogram No.~2.
\end{art}

\begin{art}[No.~111]
For the more complete discussion of the eccentricity, we have
\[
e = \surd(A^2 + B^2),
\]
where
\begin{align*}
A &= \frac{1}{12\lambda\surd(1 + \lambda^2)}\{(36 + 4\lambda^2) - (3 - 2\lambda)R_2 + (3 + 2\lambda)R_3\},\\
B &= \frac{1}{12\surd 3(1 + \lambda^2)}\{4R_1 + R_2(1 + \lambda) + R_3(1 - \lambda)\}.
\end{align*}
The eccentricity cannot be less than $1$, which is evidently right, for the point $3$ being at infinity, the orbit cannot be an ellipse. We may have $e = 1$ (or the orbit a parabola), viz. this will be the case if
\[
\tfrac{1}{\surd 3}(9 + 4\lambda^2) - (3 - 2\lambda)R_2 + (3 + 2\lambda)R_3 = 0.
\]
\end{art}


\begin{art}[No.~112]
I stop to remark that this equation may be obtained differently, as follows. Since the point $1$ is at infinity on the axis of $x$, this line will be the axis of the parabola; or the equation of the parabola will be
\[
-y^2 + 4ax + 4a^2 = 0,
\]
and we have therefore
\begin{align*}
-y_2^{\,2} + 4ax_2 + 4a^2 &= 0,\\
-y_3^{\,2} + 4ax_3 + 4a^2 &= 0,
\end{align*}
that is
\[
1 : 4a : 4a^2 = x_2 - x_3 : y_2^{\,2} - y_3^{\,2} : -y_2^{\,2}x_3 + y_3^{\,2}x_2,
\]
and therefore
\[
4a = \frac{y_2^{\,2} - y_3^{\,2}}{x_2 - x_3}, \quad 4a^2 = \frac{-y_2^{\,2}x_3 + y_3^{\,2}x_2}{x_2 - x_3},
\]
as the condition for a parabola.

But the values of $x_2$, $y_2$; $x_3$, $y_3$, ante No.~104, introducing $\lambda$ in the place of $H$, are
\begin{align*}
x_2 &= -\frac{1}{\surd 3}\frac{6 - \lambda}{3 + 2\lambda}, & y_2 &= \frac{3}{\surd 13}\frac{\surd(1 + \lambda^2)}{3 + 2\lambda},\\
x_3 &= -\frac{1}{\surd 3}\frac{6 + \lambda}{3 - 2\lambda}, & y_3 &= \frac{3}{\surd 13}\frac{\surd(1 + \lambda^2)}{3 - 2\lambda},
\end{align*}
and thence
\begin{align*}
x_2 - x_3 &= \frac{1}{\surd 3}\frac{9 - 4\lambda^2}{9 - 4\lambda^2} = \frac{1}{\surd 3},\\
y_2^{\,2} - y_3^{\,2} &= \frac{36(1 + \lambda^2)(-\lambda)}{(9 - 4\lambda^2)^2},\\
y_2^{\,2}x_3 - y_3^{\,2}x_2 &= \frac{3\surd 3(1 + \lambda^2)(6 - \lambda^2)}{(9 - 4\lambda^2)^2},
\end{align*}
and substituting these values and observing that $4a = \surd 3(y_2^{\,2} - y_3^{\,2})$, the condition gives
\[
\tfrac{1}{\surd 3}(9 + 4\lambda^2) - (3 - 2\lambda)R_2 + (3 + 2\lambda)R_3 = 0,
\]
as before.
\end{art}

\begin{art}[No.~113]
Proceeding to rationalize, and observing that the equation may be written
\[
(3 - 2\lambda)R_2 - (3 + 2\lambda)R_3 = \tfrac{1}{\surd 3}(9 + 4\lambda^2),
\]
and that here
\[
R_2^{\,2} - R_3^{\,2} = \tfrac{4}{3}(21 + 4\lambda + \tfrac{4}{3}\lambda^2) - \tfrac{4}{3}(21 - 4\lambda + \tfrac{4}{3}\lambda^2) = \tfrac{16}{3}\lambda,
\]
we find
\[
R_2^{\,2} + R_2R_3 + R_3^{\,2} = 63 + 12\lambda^2,
\]
and thence
\[
R_2 + R_3 = \frac{16\lambda}{9 + 4\lambda^2}.
\]
But we have also $R_2 - R_3 = \frac{16\lambda}{R_2 + R_3} = \frac{16\lambda(9 + 4\lambda^2)}{16\lambda} = 9 + 4\lambda^2$, and thence $R_2 = \frac{1}{2}\{(9 + 4\lambda^2) + \frac{16\lambda}{9 + 4\lambda^2}\}$. Substituting this value in $R_2^{\,2} = 21 + 4\lambda + \tfrac{4}{3}\lambda^2$, we obtain an equation for $\lambda^2$, which being simplified gives
\[
16\lambda^6 + 72\lambda^4 - 27\lambda^2 - 108 = 0,
\]
that is
\[
(4\lambda^2 - 3)(4\lambda^4 + 21\lambda^2 + 36) = 0,
\]
and consequently $\lambda^2 = \frac{3}{4}$, giving $\lambda = \pm\surd\frac{3}{4} = \pm 0.866$, $H = 40°\, 54'$ or $139°\, 6'$; the other factor gives only imaginary values. Substituting $\lambda^2 = \frac{3}{4}$ in the original equation, we find $\delta = +1$ for both values; and the eccentricity is $e = 1$, as should be.
\end{art}

\newpage
\noindent\textbf{Article Nos.~114 to 143. The Iseccentric and Isochronic Curves.}

\begin{art}[No.~114]
The spherogram is divided by the regulator circle, the separator circles, and the parabolic curve, into regions; and the form of the spherogram is such that it is divided into an inner circular region, three lateral regions, and three outer regions; and moreover each lateral region is divided by a separator into a subjacent and a superjacent portion. The iseccentric curves lie wholly within the several regions, being in general closed curves; and the same is the case with the isochronic curves; but for certain given values of the parameter, the points $A$ may be isolated points on the isoparametric line.
\end{art}

\begin{art}[No.~115]
It is sometimes necessary (more particularly as regards the Time-spherogram and isochronic lines) to distinguish from each other the several points $A$ and $B$; and for this purpose I consider the several points, as situated in the spherogram, to be accented in the following manner:
\begin{center}
\begin{tabular}{ccccc}
 & $B'''$ & $R$ & $B$ & \\
$A$ & & & & $A''$ \\
 & $A'$ & $F''$ & $F'$ & \\
 & $B''$ & & & \\
\end{tabular}
\end{center}
so that the inner triangle is $B''F'F''$ and the outer triangle is $B'''RB$.
\end{art}

\begin{art}[No.~116]
The iseccentric lines present but little difficulty. The value $e = 0$ gives the centre of the spherogram; $e = 1$ gives the parabolic curve; and $e = \infty$ gives the regulator circle and the separator circles. For any value of $e$ between $0$ and $1$, the iseccentric curve is a closed curve lying wholly within the inner circular region; for $e = 1$ it coincides with the parabolic curve; and for any value of $e$ between $1$ and $\infty$, it consists of three closed curves, one in each of the outer regions. The form of these curves is easily understood by considering their limiting positions; for $e = 1$ they coincide with the parabolic curve; and as $e$ increases they expand, ultimately for $e = \infty$ coinciding with the regulator circle and the separator circles.
\end{art}

\begin{art}[No.~117]
The iseccentric curves may be considered as completed by the insertion of portions lying in the shaded regions; and the completed curves are then continuous curves passing through the points $A$. For any value of $e$ between $0$ and $1$, the completed iseccentric curve consists of a loop in the inner region, and three branches passing through the points $A$ into the outer regions; for $e = 1$ the loop coincides with the parabolic curve, and the three branches coincide with the portions of the parabolic curve in the outer regions; and for any value of $e$ between $1$ and $\infty$, the loop has disappeared, and the curve consists of three branches, each passing through two of the points $A$.
\end{art}


\begin{art}[No.~118]
Planogram No.~5, if the calculations were completed, would give the value of $e$ for the arc within the shaded region on the base $AB'''$ of the superjacent portion of the lateral region. The value $e = 1$ belongs to the point of contact with the parabolic curve; as we approach $A$, the value of $e$ diminishes, becoming at $A$ equal to $0.264$, the value belonging to the point $A$ as a point in Planogram No.~1; and as we approach $B'''$, the value increases, becoming at $B'''$ equal to $1.018$, the value at the ecliptic for longitude $0°$.
\end{art}

\begin{art}[No.~119]
The value $e = 0.264$ at $A$ is the minimum value on the iseccentric curve which passes through $A$; and the value $e = 1.018$ at $B'''$ is the maximum value on the curve. The curve is thus a closed curve, passing through $A$ and $B'''$, and lying wholly within the superjacent portion of the lateral region; and the same is the case for the other two lateral regions.
\end{art}

\begin{art}[No.~120]
The iseccentric curves in the subjacent portions of the lateral regions are of a different character. For $e = 1$, the curve coincides with the parabolic curve; as $e$ increases, the curve expands, and for $e = 1.101$ (the value at $B$), the curve passes through $B$; for $e = 2.309$ (the value at the separator), the curve touches the separator; and for $e > 2.309$, the curve lies wholly within the outer region. The curves for $e = 1$ to $e = 1.101$ are thus closed curves in the subjacent portion of the lateral region; for $e = 1.101$ to $e = 2.309$ they are closed curves surrounding the point $B$; and for $e > 2.309$ they expand into the outer region.
\end{art}

\begin{art}[No.~121]
The iseccentric curves in the three outer regions are, as has been seen, closed curves; and for $e = \infty$ they coincide with the regulator circle and the separator circles. The value $e = 1.018$ at the ecliptic is the minimum value for the outer region; and the value increases as we recede from the ecliptic, becoming infinite on the regulator and separator circles.
\end{art}

\begin{art}[No.~122]
The iseccentric curves within the shaded regions form a distinct system; and it is to be observed that, although these curves have no direct astronomical significance, they are useful for the purpose of completing the theory, and for explaining the course of the iseccentric curves in the unshaded regions. The obliterated portions of the curves are to be restored by continuing the curves through the shaded regions; and the completed curves are then continuous curves, which may be studied as geometrical loci, without reference to their astronomical interpretation.
\end{art}

\begin{art}[No.~123]
We may without difficulty attach to the several portions of the regulator, the separators and the parabolic curve, to each portion its proper symbol $L$, $P$, $H$ and $123$, $\overline{123}$, \&c.\ as the case may be.

First, as to the regulator, it is obvious that this is separated by the points $A$ into the three portions $L\,213$, $L\,321$, $L\,132$, respectively. And inside the regulator, adjacent to these, we have portions of the parabolic curve $P\,213$, $P\,321$, $P\,132$, respectively.

Again, for one of the separators, say $B'''B'AF'R'$ (see here and in all that follows the notation-diagram, No.~115); since the point $2$ is here at infinity this must be at every portion thereof either $H\,132$ or else $H\,\overline{312}$. The point $B'''$ is $H\,132$ and the point $B'$ is $H\,\overline{312}$; consequently, as the orbit-pole passes along the separator from $B'''$ to $B'$, the symbol is at first $H\,132$ and at last $H\,\overline{312}$; the transition takes place at the point of contact of the parabolic curve which is indifferently $P\,132$ or $P\,213$.
\end{art}

\begin{art}[No.~124]
(In further explanation of the transition, consider the orbit-pole as passing from $B'''$ to $B$, not on the separator, but indefinitely near it; it can only do so by twice crossing the parabolic curve near the point of contact; the orbit is first $H\,132$, or say $H\,132$, then $P\,132$, then an ellipse, which when the orbit-pole again arrives at the parabolic curve changes into $\overline{312}$; and it finally becomes $H\,\overline{312}$ or $H\,\overline{312}$.)
\end{art}

\begin{art}[No.~125]
Again, since on the two separators through $B'''$, in the portions adjacent to $B'''$, the symbols are $H\,132$ and $H\,\overline{132}$, it is clear that in the adjacent portion of the parabolic curve (terminated each way by a point of contact with these separators respectively) the symbol must be $P\,132$; at the point of contact with the first-mentioned separator $B'''B'AF'R'$, this becomes $P\,132 = P\,213$; and beyond the point of contact it becomes $P\,213$, continuing so until it arrives at the next point of contact with the separator $EA''E'$: there is always in the symbol for the parabolic curve a transition at each point of contact with a separator.
\end{art}

\begin{art}[No.~126]
The isochronic curves are of a more complicated character than the iseccentric curves. The value $T = 0$ gives the regulator circle; $T = \infty$ gives the parabolic curve and the separator circles. For any value of $T$ between $0$ and $\infty$, the isochronic curve consists of three branches, one in each of the outer regions; and these branches are closed curves, surrounding the points $B$, and cutting the ecliptic at right angles.
\end{art}

\begin{art}[No.~127]
The isochronic curves may be completed by the insertion of portions lying in the shaded regions; and the completed curves are then continuous curves, which intersect the regulator circle at the points $A$. For any value of $T$ between $0$ and $\infty$, the completed isochronic curve consists of a loop in each of the three outer regions, and three branches passing through the points $A$ and uniting the loops; the branches lie partly in the unshaded regions and partly in the shaded regions.
\end{art}

\begin{art}[No.~128]
The times $T_{12}$, $T_{23}$, $T_{31}$ are calculated, Planogram 1, part 1, for the meridian long.~$90°$, and ditto part 2 for the meridian long.~$270°$; and in Planogram 2 for the meridian long.~$180°$. As regards these last values, it is easy to see that, in order to pass to the meridian long.~$0°$, the numbers $2$, $3$ must be interchanged; that is, long.~$0°$, the $T_{12}$, $T_{23}$, $T_{31}$ are respectively equal to the values, long.~$180°$, $T_{13}$, $T_{32}$, $T_{21}$.

Moreover, the numbers $1$, $2$, $3$ may be changed into $2$, $3$, $1$, or into $3$, $1$, $2$, provided the longitude is increased by $120°$ and $240°$ in the two cases respectively; that is,
\begin{align*}
T_{12} \text{ long. } a &= T_{31} \text{ long. } (a + 120°) = T_{23} \text{ long. } (a + 240°),\\
T_{23} \text{ long. } a &= T_{12} \text{ long. } (a + 120°) = T_{31} \text{ long. } (a + 240°),\\
T_{31} \text{ long. } a &= T_{23} \text{ long. } (a + 120°) = T_{12} \text{ long. } (a + 240°).
\end{align*}
\end{art}


\begin{art}[No.~129]
By means of the foregoing two relations, $T_{12}$ for the several longitudes $0°$, $30°$, $60°$, \dots\ $330°$, is given as equal to the $T_{12}$, $T_{23}$, or $T_{31}$ for long.~$90°$, $270°$, or $180°$, that is, to the $T_{12}$, $T_{23}$, or $T_{31}$ of Planogram No.~1, part 1 or 2, or of Planogram No.~2. For example, $T_{12}$ long.~$240° = T_{31}$ long.~$0° = T_{23}$ long.~$180°$, that is, it is equal to the $T_{23}$ of Planogram No.~2. We thus find
\begin{center}
\begin{tabular}{c|l}
Long. & $T_{12}$ is \\
\hline
$0°$ & $T_{31}$ of Plan.~No.~2 \\
$30°$ & $T_{12}$ of Plan.~No.~1, pt.~2 \\
$60°$ & $T_{23}$ \,\,\,\,\,\,\,\,\,\,\,\,\,\,\,\,\,\,\,\,\,\,\,\,\,\,\,\,\,\,\,\,\,\,\,\,\,\,\,\,\,\,\, \\
$90°$ & $T_{12}$ of Plan.~No.~1, pt.~1 \\
$120°$ & $T_{23}$ \,\,\,\,\,\,\,\,\,\,\,\,\,\,\,\,\,\,\,\,\,\,\,\,\,\,\,\,\,\,\,\,\,\,\,\,\,\,\,\,\,\,\, \\
$150°$ & $T_{31}$ \,\,\,\,\,\,\,\,\,\,\,\,\,\,\,\,\,\,\,\,\,\,\,\,\,\,\,\,\,\,\,\,\,\,\,\,\,\,\,\,\,\,\, \\
$180°$ & $T_{12}$ \,\,\,\,\,\,\,\,\,\,\,\,\,\,\,\,\,\,\,\,\,\,\,\,\,\,\,\,\,\,\,\,\,\,\,\,\,\,\,\,\,\,\, \\
$210°$ & $T_{23}$ \,\,\,\,\,\,\,\,\,\,\,\,\,\,\,\,\,\,\,\,\,\,\,\,\,\,\,\,\,\,\,\,\,\,\,\,\,\,\,\,\,\,\, \\
$240°$ & $T_{31}$ \,\,\,\,\,\,\,\,\,\,\,\,\,\,\,\,\,\,\,\,\,\,\,\,\,\,\,\,\,\,\,\,\,\,\,\,\,\,\,\,\,\,\, \\
$270°$ & $T_{12}$ of Plan.~No.~1, pt.~1 \\
$300°$ & $T_{23}$ \,\,\,\,\,\,\,\,\,\,\,\,\,\,\,\,\,\,\,\,\,\,\,\,\,\,\,\,\,\,\,\,\,\,\,\,\,\,\,\,\,\,\, \\
$330°$ & $T_{31}$ \,\,\,\,\,\,\,\,\,\,\,\,\,\,\,\,\,\,\,\,\,\,\,\,\,\,\,\,\,\,\,\,\,\,\,\,\,\,\,\,\,\,\, \\
\end{tabular}
\end{center}
and observing that for Planogram No.~1, --- part 1 or 2, we have $T_{12} = T_{31}$, it hence appears as above, that the meridian $30°$--$210°$ is an axis of symmetry of the spherogram. In what precedes it has been assumed that the colatitudes only extend from $0°$ to $90°$, but in the spherogram they extend for the meridians $30°$, $150°$, $270°$, to the colatitude $106°\, 6'$, the values for the colatitudes above $90°$ are those for the omitted portions $90°$ to $73°\, 54'$ of the opposite meridian.

N.B.\ meridian extends from the pole in one direction only, unless the contrary is expressed or implied, as in speaking of a meridian $0°$--$180°$.
\end{art}

\begin{art}[No.~130]
I attend, in the first instance, to the axis of symmetry or meridian $30°$--$210°$. Proceeding along the meridian long.~$30°$ or towards the point $A$, the value of $T_{12}$ decreases from $1$ at the centre to a minimum $= 0.950$ at colatitude $11°$ (call this the point $X$), and it then increases to $1.983$ at $A$, and thence to $58.62$ at $90°$ and $\infty$ at the parabolic boundary of the axial region. In the opposite direction it increases from $1$ at the centre to $\infty$ at the parabolic boundary of the inner region. The minimum value $0.950$ on the axis of symmetry indicates a node on the isochronic curve; that is, the point $X$ is a node on the isochronic $T_{12} = 0.950$. This will consist of two branches, proceeding from $A'$, $A''$, respectively, cutting the axis and each other at $X$, then again cutting at $A$, and thence passing on into the axial region, and respectively terminating on the separator boundary $B'AB''$ thereof.
\end{art}

\begin{art}[No.~131]
This curve, which I call the \emph{nodal isochronic}, divides the inner region into a loop, antiloop, and two side regions. On each of the meridians $0°$, $60°$, the value of $T_{12}$ diminishes from $1$ at the centre to a minimum which is less, and then increases to a maximum which is greater, than $0.950$; the value then diminishes to $0$ on the regulator: on emergence of the meridian from the shaded into the axial region, the value is $= 0.909$, and it thence increases to $\infty$ at the parabolic boundary of the axial region; these data further determine the form of the nodal isochronic, viz., each of the two half meridians cuts the loop twice, and again cuts the curve in the axial region.

The nodal isochronic, at each of the points $A'$, $A''$, continues its course into the lateral region, returning to the same point $A$ or $A'$, so as to form in each of the lateral regions a loop. Considering the loop as formed of two branches, each proceeding from $A'$ or $A''$, the one which is the continuation of the course within the inner region I call the \emph{lower branch}; the other, the \emph{upper branch}.
\end{art}

\begin{art}[No.~132]
On the meridian $90°$, $330°$, through the points $B''$, $B''$, respectively, the value of $T_{12}$ diminishes from $1$ at the centre to $0$ at the regulator, where these meridians are considered as terminating.

On the meridians $120°$, $300°$ (meridian at right angles to the axis of symmetry), the value of $T_{12}$ diminishes from $1$ at the centre to a minimum less than $0.878$, and then increasing to a maximum of over $0.895$ diminishes to $0$ at the regulator. On emergence of the meridian from the shaded and half-shaded region on the parabolic boundary of the lateral region the value is $= \infty$, and it thence diminishes to $1.148$ on the separator boundary $B'''B''$ or $B'B'''$.

On the meridians $150°$, $270°$, which pass through $A'$, $A''$, respectively, the value of $T_{12}$ increases from $1$ at the centre to $1.377$ at the regulator, and thence through $2.255$ at $90°$ to $\infty$ at $B'''$ or $R'$.

And finally, on the meridians $180°$, $240°$, the value of $T_{12}$ increases from $1$ at the centre to $\infty$ at the parabolic inner boundary, and then on emergence from the half-shaded and shaded region at the separator boundary $B'''A'$ or $B'''A''$, the value is $= \infty$, and it thence diminishes to a minimum under $0.6343$, and again increases to $\infty$ at the separator boundary $B''B'''$ or $B'''B''$.
\end{art}

\begin{art}[No.~133]
By what precedes, it appears that on the separator boundary $B'''B'$ or $B'''B''$ of either of the lateral regions, the value of $T_{12}$ is at each extremity $= \infty$, and at an intermediate point $= 1.148$; there is consequently a minimum value less than $1.148$, and therefore two points at each of which the value is $= 1.983$.

Now resuming the consideration of the cuspidal isochronic ($T_{12} = 1.983$) as regards the remaining portions thereof, viz., those in the lateral and inner regions; and considering first the lateral region $B'''B''B'$, there will be from each of the points just referred to on the boundary $B'''B'$ a branch; one (which I call the lower branch) from the point nearer $B'$, passes, on the right-hand side of the meridian through $A'$, to $A'$; the other (which I call the upper branch) from the point nearer $B'''$, passes on the left-hand side of the same meridian, to $A''$.
\end{art}


\begin{art}[No.~134]
By what precedes, it appears that on the separator boundary $B'''B'$ or $B'''B''$ of either of the lateral regions, the value of $T_{12}$ is at each extremity $= \infty$, and at an intermediate point $= 1.148$; there is consequently a minimum value less than $1.148$, and therefore two points at each of which the value is $= 1.983$.

Now resuming the consideration of the cuspidal isochronic ($T_{12} = 1.983$) as regards the remaining portions thereof, viz., those in the lateral and inner regions; and considering first the lateral region $B'''B''B'$, there will be from each of the points just referred to on the boundary $B'''B'$ a branch; one (which I call the lower branch) from the point nearer $B'$, passes, on the right-hand side of the meridian through $A'$, to $A'$; the other (which I call the upper branch) from the point nearer $B'''$, passes on the left-hand side of the same meridian, to $A''$.
\end{art}

\begin{art}[No.~135]
The cuspidal isochronic in the lateral region $B'''B''B'$ thus consists of an upper branch touching the separator at $A''$, a lower branch from the minimum point on $B'''B'$ to $A'$, and in the inner region, a branch (continuation of the lower branch) lying between the cuspidal curve and the parabolic boundary of the antiloop, and uniting the points $A'$, $A''$.
\end{art}

\begin{art}[No.~136]
Now by what precedes there is on the separator boundary $B'''B'$ of the lateral region a point where $T_{12}$ has a minimum value less than $1.148$, and consequently, for any given value, say for a value between this minimum and $1.377$, there are on $RB'''$ two points where $T_{12}$ has the given value. These points cannot lie on the loop of the curve belonging to the given value (for this loop is wholly on the left hand of the meridian); hence the complete curve for the given value of $T_{12}$ will include (within the lateral region) besides the loop, a branch uniting the two points in question; say a \emph{link branch}.
\end{art}

\begin{art}[No.~137]
It follows that there is between $T_{12} = 1.377$ and $1.983$, a value (to fix the ideas, say $= 1.80$?, it being understood that I do not attempt to determine this value) for which the loop and link branch will unite themselves together, the point of junction becoming as usual a node; viz., there will be a curve $T_{12} = 1.80$? having in the two lateral regions respectively the nodes $Y$, $Y'$; or say the curve has in each lateral region a self-intersecting loop. For any greater value of $T_{12}$ (as for example the value $1.983$ belonging to the cuspidal curve) there are two branches inclosing the self-intersecting loop; for a less value, as has been seen, instead of the self-intersecting loop, there is a loop and link branch; at least this is the case until for the minimum value $< 1.148$ of $T_{12}$ on the separator boundary $B'''B'$ the link branch disappears. For smaller values down to $T_{12} = 0.950$, which belongs to the nodal isochronic, there is no link branch, but only the loop; and as $T_{12}$ diminishes below this value, there is still a continually diminishing loop.
\end{art}

\begin{art}[No.~138]
For $T_{12} = 0.950$, the nodal isochronic, the loop in each lateral region evanesces, the point $X$ on the axis of symmetry being a node; the curve consists of the loop in the inner region, and of the two branches which proceed from $A'$ to $X$ and from $X$ to $A''$, and which in the axial region unite themselves at $A$.
\end{art}

\begin{art}[No.~139]
It is to be observed, both as regards the iseccentric and the isochronic curves, that there is a real meaning in the obliterated portions; viz., to any position of the orbit-pole on such obliterated portion of the curve there corresponds a conic determined by means of a given trivector, but which, by reason of its being a convex hyperbola, or hyperbola such that the three points do not lie on the same branch thereof, is not regarded as an orbit. The obliterated portions have been in the present Memoir considered only so far as they present themselves in continuity with the curves which are the loci of the pole of a proper orbit, and for the purpose of explaining the course of these curves; and the curves completed as above are not the complete loci which would be obtained if, instead of the selected conic called the orbit, we had considered simultaneously the four conics determined by means of any given trivector; such extension of the theory would, it is probable, be interesting geometrically but it would be devoid of all astronomical significance.
\end{art}

\begin{art}[No.~140]
To obtain a comprehensive statement of the foregoing results, we may (as in the case of the iseccentric lines) imagine the curves completed and rendered continuous by the insertion of portions lying outside the spherogram, or within the half-shaded and shaded regions; which inserted portions are to be ultimately obliterated. The upper and under branches terminating in the separator boundary of a lateral region are thus completed into a loop; the link branch into a closed curve or oval; the vanishing of the link branch happens when the oval, on the point of passing outside the separator boundary of the lateral region, just touches this boundary; as $T_{12}$ diminishes to the value for which this happens, and continues still further to diminish, I think it may be assumed that there is some value (to fix the ideas, say $T_{12} = 1.10$?, but I do not attempt to determine it) for which the oval becomes a conjugate point; viz., for this value $T_{12} = 1.10$? the curve will have two conjugate points (nodes) $Z'$, $Z''$, outside the two lateral regions respectively.
\end{art}

\begin{art}[No.~141]
Resuming the consideration of the isochronic curves proper, viz., those which lie wholly or partly in the unshaded regions: for $T_{12} = 0.950$ we have the nodal isochronic, consisting of the loop in the inner region, and the branches from $A'$ to $X$ and $X$ to $A''$; for $T_{12} > 0.950$, say $T_{12} = 1.10$? to $1.377$, the curve consists of the loop in the inner region, a loop in each lateral region, and a link branch in each lateral region; for $T_{12} = 1.377$ the link branch has evanesced; for $T_{12} = 1.80$? the loop in each lateral region has become self-intersecting; for $T_{12} = 1.983$ we have the cuspidal isochronic; for $T_{12} > 1.983$, the curve consists of a loop in the inner region, and in each lateral region an upper branch touching the separator at $A'$ or $A''$, and a lower branch; and in the inner region, a branch (continuation of the lower branches) lying between the cuspidal curve and the parabolic boundary of the antiloop, and uniting the points $A'$, $A''$.
\end{art}

\begin{art}[No.~142]
There is not in the several curves any discontinuity of direction at the point $A'$ or $A''$: the branch from $A$ within the shaded or half-shaded region, emerges at $A'$ or $A''$ into the lateral region, uniting itself with the upper branch of the loop; it can only do this in virtue of its being at $A'$ or $A''$ a tangent to the separator (for otherwise it would cross the separator and regulator into the inner region); that is, the continuation thereof, or upper branch of the loop, must at the point $A'$ or $A''$ touch the separator; it has been previously throughout assumed that this is so.
\end{art}

\begin{art}[No.~143]
It is to be observed, both as regards the iseccentric and the isochronic curves, that there is a real meaning in the obliterated portions; viz., to any position of the orbit-pole on such obliterated portion of the curve there corresponds a conic determined by means of a given trivector, but which, by reason of its being a convex hyperbola, or hyperbola such that the three points do not lie on the same branch thereof, is not regarded as an orbit. The obliterated portions have been in the present Memoir considered only so far as they present themselves in continuity with the curves which are the loci of the pole of a proper orbit, and for the purpose of explaining the course of these curves; and the curves completed as above are not the complete loci which would be obtained if, instead of the selected conic called the orbit, we had considered simultaneously the four conics determined by means of any given trivector; such extension of the theory would, it is probable, be interesting geometrically, but it would be devoid of all astronomical significance.
\end{art}

\vspace{1em}
\begin{center}
\rule{0.5\textwidth}{0.4pt}
\end{center}
\clearpage


% =========================================================
% Chunk: cayley_vol07_pages_501_550.tex  (orig pages 501--550)
% =========================================================
\thispagestyle{empty}

\noindent\textit{Transcription of pages 501--550 of the collected edition. This excerpt covers the latter part of Paper~477, the whole of Paper~478, and the opening of Paper~479.}


\newpage

\section*{Plates}
\addcontentsline{toc}{section}{Plates}

The original volume contains several plates (I--VI) illustrating the eclipse construction and the geodesic lines on an ellipsoid.

\begin{itemize}
\item \textbf{Plate I.} Guide Hyperbolas $b=90^{\circ}$.
\item \textbf{Plate II.} Photograph $b=90^{\circ}$.
\item \textbf{Plate III.} Photograph $b=180^{\circ}$.
\item \textbf{Plate IV.} Spherograph.
\item \textbf{Plate V.} Spherograph.
\item \textbf{Plate VI.} (Inserted at the end of Paper~477.)
\end{itemize}

\newpage

\section*{Paper 477}
\addcontentsline{toc}{section}{Paper 477 --- On the Graphical Construction of a Solar Eclipse}

\subsection*{On the Graphical Construction of a Solar Eclipse}

\begin{center}
[From the \textit{Memoirs of the Royal Astronomical Society}, vol.~xxxix. (1870---1871), pp.~1---17. Read January~13, 1871.]
\end{center}

The present Memoir contains the explanation of a Graphical Construction of a Solar Eclipse, which (it appears to me) is at once easy, and susceptible of considerable accuracy: I think that if made on the suggested scale (radius = 12 inches) we might by means of it construct a diagram such as the eclipse-diagrams of the \textit{Nautical Almanac}, with at least as much accuracy as could be exhibited in a diagram on that scale.

\subsubsection*{Article Nos.\ 1 to 9. General Explanation of the Construction.}

\noindent\textbf{1.} We may imagine the celestial sphere as seen from the centre of the Earth stereographically projected at each instant during the eclipse---the radius of the bounding circle of the projected hemisphere being a given length, say twelve inches, which is taken as unity---in such wise that the centre of the Moon is always at the centre of the projection, say M, and the pole (suppose the north pole, say N) of the Earth on a given radius: its position on this radius will in strictness be variable, viz. distance from centre = projection of Moon's N.\ P.\ D.\ = $\tan \tfrac{1}{2}\Delta$. Suppose, for a moment, that the position at each instant of the Sun's centre were also laid down on the projection, so as to obtain the projection of the Sun's relative orbit; this will be a terminated short line $A'B'$ (fig.~1), nearly straight, and lying near the centre of the projection (this relative orbit is not to be actually laid down, but it is replaced, as will presently be explained, by a relative orbit on a very enlarged scale); if at any instant the position of the Sun on the relative orbit be denoted by $S'$, then the straight line $MS'$ is the projection of the arc of great circle through the centres of the Moon and Sun, so that $E$ being the angular distance of the centres, the length of the line $MS'$ is $= \tan \tfrac{1}{2}E$, or ($E$ being small) it is $= \tfrac{1}{2}E$.

\noindent\textbf{2.} Produce $S'M$ through the centre $M$ to a point $Z$, and consider $Z$ as representing a point on the Earth's surface: to determine the geographical position of $Z$, we must consider the projected meridian $NZ$ which passes through $Z$: the arc $NZ$, regarded as a projection, represents the N.\ P.\ D.\ or colatitude of $Z$, and the actual angle at $N$ which the tangent of $NZ$ makes with the line $NM$ is equal to the celestial angle $ZNM$ which is = Moon's hour-angle from $Z$, or what is the same thing = difference of Moon's hour-angle from Greenwich and of the longitude of $Z$ (as the figure is drawn, $\angle ZNM$ = Moon's hour-angle E.\ of Greenwich, less E.\ longitude of $Z$).

\begin{figure}[H]
\centering
\caption*{Fig.~1.}
\end{figure}

\noindent\textbf{3.} Now, considering the Moon and Sun as seen from $Z$, we may disregard the parallactic depression of the Sun, and attribute to the Moon a displacement equal to the difference of the parallactic displacements of the Moon and Sun; that is, regarding the zenith distances $ZM$, $ZS'$ as equal, we may consider the Moon's centre as depressed by parallax in the direction of the arc $MS'$ through an arc $MQ'$, = $\sin^{-1} P' \sin ZM$, where $P' = .99837(\varpi - \varpi')$ is the quantity thus designated in the Appendix to the Nautical Almanac for 1836, viz.\ it is $= \varpi' - \varpi$, the difference of the equatoreal horizontal parallaxes at the time of the eclipse, multiplied by a factor $.99837$, which answers to a distance of $Z$ from the Earth's centre = Earth's radius for latitude $45^{\circ}$. And if we take $Q'$ such that its angular distance from $S'$ = sum of angular semidiameters of the Sun and Moon, the locus of $Q'$ is very nearly a circle about the centre $S'$, and the corresponding positions of $Z$ give the positions on the Earth where the limbs are in exterior contact, or, what is the same thing, give the penumbral curve on the Earth's surface for the position $S'$ of the Sun.

\noindent\textbf{4.} Instead of $\text{Arc } MQ' = \sin^{-1} P' \sin ZM$, we may write $\text{Arc } MQ' = P' \sin ZM$, or, using $p$ to denote the linear distance $ZM$ in the projection, we have $p = \tan \tfrac{1}{2}ZM$, and therefore $\sin ZM = \dfrac{2p}{p^{2}+1}$; hence
\[
\text{Arc } MQ' = P' \frac{2p}{p^{2}+1},
\]
and the linear distance $MQ'$ in the projection is $= \tan \tfrac{1}{2}\text{arc } MQ'$, say this is $= \tfrac{1}{2}\text{arc } MQ'$, or calling this linear distance $l$ we have
\[
l = P' \frac{p}{p^{2}+1}.
\]

\noindent\textbf{5.} Hence, if instead of the original representation of the Sun's relative orbit we consider an enlarged representation thereof and of the depressed positions $Q'$ of the Moon, obtained by increasing the several distances from the centre of the projection in the ratio $hP'$ to 1, and if instead of $A'$, $B'$, $S'$, $Q'$, we use $A$, $B$, $S$, $Q$, as referring to this enlarged representation, then representing by $r$ the linear distance $MQ$, we have $r = \dfrac{2p}{p^{2}+1} \cdot P'$, and consequently
\[
r = \frac{2p}{p^{2}+1}.
\]
We have here $r$ representing the parallactic depression corresponding to the zenith distance $ZM$, where $p = \tan \tfrac{1}{2}ZM$; that is, $ZM = 90^{\circ}$, $p = 1$, and therefore $r = 1$; but for $ZM = 90^{\circ}$, the parallactic depression is $= P'$; that is, the scale of the enlarged representation of the Sun's relative orbit, or say simply the scale of the relative orbit (for on the original scale it was never actually constructed at all) is such that we have $P'$ (= about $60'$) represented by the radius of the bounding circle of the projected hemisphere, = 12 inches.

\noindent\textbf{6.} The process is, construct the relative orbit on the scale $P'$ = radius of bounding circle: take $S$ for the position at any given instant of the Sun in the relative orbit, and with centre $S$ and radius $= s + \sigma$ (sum of the angular semidiameters, of course on the same scale) describe a circle. The positions $A$ and $B$ of the Sun at the beginning and end of the eclipse respectively are such that this circle just touches the bounding circle externally, viz.\ the distances of $A$ and $B$ from the centre of the projection are each = radius of bounding circle $+ s + \sigma$. At any intermediate instant the circle, radius $s + \sigma$, lies wholly or partly within the bounding circle; in the latter case we attend only to the arc thereof which lies within the bounding circle. Taking then $Q$ any point whatever on the circle or arc in question, we join $Q$ with the centre $M$ of the projection, and produce this line through $M$ to a point $Z$, such that the distances $MQ$, $MZ$, being $r$, $p$ respectively, we have as above
\[
r = \frac{2p}{p^{2}+1},
\]
or, what is the same thing, writing $\theta$ in place of $Z$, and regarding this angle $\theta$ as a variable parameter, the relation between $r$, $p$, may be expressed by means of the two equations, $p = \tan \tfrac{1}{2}\theta$, $r = \sin \theta$.

\noindent\textbf{7.} Practically the construction may be performed very easily by means of a straight edge twenty-four inches long, graduated from the centre, one half of it for the values of $r$, and the other half for the corresponding values of $p$ (that is, the first half is graduated for $\sin \theta$, and the second half for $\tan \tfrac{1}{2}\theta$): we have thus, corresponding to the circle or arc of circle which is the locus of $Q$, a closed curve, or arc thereof terminated each way at the bounding circle, for the locus of $Z$: which curve or arc of a curve is the penumbral curve on the Earth's surface for the position $S$ of the Sun in the relative orbit.

\noindent\textbf{8.} The north pole of the Earth occupies in the projection a given position, viz.\ it is situate on a given radius at a distance = $\tan \tfrac{1}{2}$\ Moon's N.P.D.; which N.P.D.\ may be considered as being throughout the eclipse constant, and equal to its value at the middle of the eclipse. But in order to arrive at the geographical signification of the figure it is necessary to lay down on the projection the position of the meridian of Greenwich; which position, it will be remembered, varies according to the position of $S$. Supposing this done, we could of course (at least theoretically) draw the whole series of meridians and parallels, and thereby determine the latitudes and longitudes of the several points of the penumbral curve, or (if need is) transfer it to a different projection of the Earth's surface. The actual description of the meridians and parallels would, however, be very laborious, and fortunately it can be avoided by means of a single blank projection and a slight modification of the foregoing process, as will be explained.

\noindent\textbf{9.} But before considering how this is, it is proper to remark that constructing as above a figure of the penumbral curves corresponding to the several positions of the Sun: by what precedes these different curves may indeed be considered as belonging to the same position of the north pole in the projection, but they belong to different positions of the meridian of Greenwich; and thus they do not constitute a representation of the penumbral curves each in its proper terrestrial position, but only a representation in which the penumbral curves are affected each of them by a different displacement in longitude.

\subsubsection*{Article Nos.\ 10 to 13. Modification in order to the Applicability of a Single Blank Projection.}

\noindent\textbf{10.} Imagine a stereographic projection of the meridians and parallels on the plane of a meridian, radius of this meridian, that is of the bounding circle of the projected hemisphere, being = 12 inches as before; and the poles $N$, $S$ being of course opposite points on the circumference of the bounding circle---the meridians and parallels are, however, to be produced outside the bounding circle; say this is the ``blank projection,'' and let its centre be denoted by $M_{1}$. Then, if at any point $M$ on the radius $MN$, we draw the chord $CD$ at right angles to $M_{1}N$, and on $CD$ as diameter describe a circle, this will cut out from the blank projection a new projection having the last-mentioned circle for its bounding circle, and in which $N$ is the north pole; viz.\ the meridians of the blank projection will be meridians, and the parallels of the blank projection will be parallels, in this new projection. And, moreover, if the longitudes are reckoned from the meridian $NMM_{1}$, then the meridian of a given longitude in the blank projection will in the new projection be the meridian of the same longitude---but the parallel of a given colatitude $c$ in the blank projection will, in the new projection, be the parallel of a different colatitude $c'$,---the relation of $c$, $c'$ being, however, a very simple one, as presently explained.

\begin{figure}[H]
\centering
\caption*{Fig.~2.}
\end{figure}

\noindent\textbf{11.} The blank projection thus at once gives a projection in which the north pole $N$ has any assumed position whatever; and it is easy to see that in order that its distance $MN$ from the centre of the projection may represent a given angle $A$, we have only to take $M_{1}M = \cos A$ (that is = 12 inches $\times \cos A$), the corresponding value of $MC$ being $MC = \sin A$ (that is = 12 inches $\times \sin A$). Hence $A$ denoting the Moon's N.P.D.\ at the middle of the eclipse, we can by means of the blank projection construct a projection such as that above referred to, only the radius of its bounding circle, instead of being unity (12 inches), is in the reduced ratio of $1 : \sin A$.

\noindent\textbf{12.} The figure of the penumbral curves as originally constructed requires, therefore, to be reduced in the ratio $1 : \sin A$, viz.\ each of the distances from the centre $M$ should be reduced in this ratio; this could of course be done easily enough with a pair of proportional compasses; but by means of a different graduation of the straight edge we may, in the first instance, construct the penumbral curves on the proper reduced scale; viz.\ assuming that we have on the proper scale a proportional-scale figure such as is here shown, the line $Mr$ (= 12 inches) being graduated for $\sin \theta$, and the line $MA$ (also = 12 inches) for $\tan \tfrac{1}{2}\theta$, and a set of parallel lines being drawn through the last-mentioned graduations---then taking the distance $Mp = \sin A$, that is = 12 inches $\times \sin A$, and drawing the line $Mp$, this line will, it is clear, be graduated for $\sin A \tan \tfrac{1}{2}\theta$: so that we may from the figure graduate the straight edge, the one half of it by means of the line $Mr$, and the other half of it by means of the line $Mp$; and with the straight edge thus graduated, at once lay down the penumbral curve on the scale now in question. And we thus obtain a figure containing as well %61---2 the penumbral curves, as the meridians and parallels which serve to fix their terrestrial position.

\begin{figure}[H]
\centering
\caption*{Fig.~3.}
\end{figure}

\noindent\textbf{13.} It remains in the new projection to find the colatitude belonging to any given parallel. Supposing that the colatitude in the blank projection is $= c'$, then it may be shown that the colatitude $c$ of the same parallel in the reduced projection is given by means of the equation $\tan \tfrac{1}{2}c = \cot \tfrac{1}{2}A \tan \tfrac{1}{2}c'$, from which $c$ might be calculated numerically: but the required value may also be obtained graphically. In fact, considering the parallel which cuts $NS$ (see fig.~2) in a point $R$, then, if by lines drawn from $C$ as a centre we project $N$, $R$, $S$, on the circumference of the bounding circle of the new projection---say the projections of these points are $n$, $r$, $s$, respectively, the arc $Ns$ is a semicircle, and the arcs $nr$, $sr$, are respectively the N.P.D.\ and the S.P.D.\ of the parallel in question. It may be added that in the new projection the equator is represented by the parallel through the points $C$, $D$; so that if this cuts $NQ$ in $Q$, and the point $Q$ be in like manner projected on the bounding circle---say its projection is $q$, then the arcs $nq$, $sq$, will be each of them a quadrant, and the arc $qr$ will be the latitude of the parallel in question.

\subsubsection*{Article Nos.\ 14 to 18. As to the Construction of the Relative Orbits.}

\noindent\textbf{14.} It is convenient to notice that if $e$, $e'$, be the values of the equation of time at the preceding and following Greenwich Mean Noons (viz.\ $e$ or $e'$ = G.M.T.\ of apparent Noon) then that the Sun's hour-angle E.\ of Greenwich at the Greenwich mean time $t$ is
\[
h' = e + e'\left(1 + \frac{t}{24}\right),
\]
and that if $\alpha$, $\alpha'$, are the R.A.'s of the Moon and Sun respectively, then $h' - h = \alpha' - \alpha$, which is also of the form $A + Bt$. In the reduced projection, the Moon is always at the centre $M$; by means of the values of $h' - h$ we lay down at any instant the Sun's position in R.A.\ and then by means of the values of $h'$, the position of the meridian of Greenwich; and we thus at any instant read off the terrestrial longitude of any point of the reduced projection, or say, of a point on the penumbral curve.

\noindent\textbf{15.} With regard to the construction of the relative orbit, it is to be observed that if at any instant the hour-angle and N.P.D.\ of the Moon are $h$, $A$, and those of the Sun, $h'$, $A'$, then taking $M$ as origin, and the axes $Mx$, $My$, in the direction of $NM$ produced, and perpendicular hereto to the right (or eastwards), then the rectangular coordinates of $S'$ are approximately $x = \tfrac{1}{P'}(A' - A)$, $y = \tfrac{1}{P'}(h' - h) \sin A$, where $h' - h$ is equal to the difference of R.A.\ of the Sun and Moon. Hence, in the adopted relative orbit, the coordinates of $S$ would be
\[
x = \frac{A' - A}{P'} \cdot 12 \text{ in.}, \qquad y = \frac{h' - h}{P'} \sin A \cdot 12 \text{ in.},
\]
where, $P'$ being reckoned in minutes, $A' - A$ and $h' - h$ are also reckoned in minutes.

\noindent\textbf{16.} Moreover, $A$ may be considered as constant during the eclipse: and the relative orbit, assumed to be a straight line, will be determined by means of two points thereof; viz.\ knowing the values of $A' - A$, and $h' - h$ at about the time of the beginning and at about the time of the end of the eclipse, we construct by these formul\ae\ two points of the orbit, and joining them by a straight line, we have the orbit. Also the position at any instant of the Sun in this relative orbit will be obtained by considering its motion therein as being uniform. I think there is no advantage in the adoption of a more accurate construction: for although we may for any given instant use the accurate values of $h$, $A$, $h'$, $A'$, and so construct the position in the relative orbit, and the corresponding penumbral curve, yet if in the determination of the geographical significance thereof, we were to use for each curve a different value of $A$, the simplicity of the construction would disappear; and it is, moreover, doubtful whether the trifling corrections would not be within the limits of the necessary errors of the drawing.

\noindent\textbf{17.} But if $MS' = E$, and $\angle xMS' = \delta$, the accurate values for the coordinates of $S'$ are $x = \tan \tfrac{1}{2} E \cos \delta$, $y = \tan \tfrac{1}{2} E \sin \delta$, and the values for the coordinates of $S$ are
\[
x = \frac{1}{P'} \tan \tfrac{1}{2} E \cos \delta \cdot 12 \text{ in.}, \qquad y = \frac{1}{P'} \tan \tfrac{1}{2} E \sin \delta \cdot 12 \text{ in.},
\]
where $P'$ is still reckoned in minutes, and of course arc $1' = \tfrac{1}{206265}$. As the scale is considerable, it is worth while to inquire whether the employment of the accurate formul\ae\ would produce an appreciable difference in the position of $S$.

We have $\sin \tfrac{1}{2}E \div \sin A' = \sin(h' - h) \div \sin E$, that is $\sin E \sin \delta = \sin(h' - h) \sin A'$, and $\cos E = \cos A \cos A' + \sin A \sin A' \cos(h' - h)$: or putting for shortness $A' - A = \alpha$, $h' - h = \beta$, we have $\sin \tfrac{1}{2}E \sin \delta = \sin \tfrac{1}{2}\beta \sin A'$, and $\cos E = \cos \alpha - \sin A \sin A' \cdot 2\sin^{2}\tfrac{1}{2}\beta$. Hence, attending to the equations $\cos^{2}\tfrac{1}{2}E = 2(\cos^{2}\tfrac{1}{2}\alpha - \sin A \sin A' \sin^{2}\tfrac{1}{2}\beta)$ and $\sin^{2}\tfrac{1}{2}E = 2(\sin^{2}\tfrac{1}{2}\alpha + \sin A \sin A' \sin^{2}\tfrac{1}{2}\beta)$, we find
\[
\tan \tfrac{1}{2}E \sin \delta = \frac{\sin \tfrac{1}{2}\beta \cos \tfrac{1}{2}\beta \sin A'}{\sqrt{\cos^{2}\tfrac{1}{2}\alpha - \sin A \sin A' \sin^{2}\tfrac{1}{2}\beta}}
\]
and
\[
\tan \tfrac{1}{2}E \cos \delta = \frac{\sqrt{\sin^{2}\tfrac{1}{2}\alpha + \sin A \sin A' \sin^{2}\tfrac{1}{2}\beta - \sin^{2}\tfrac{1}{2}\beta \cos^{2}\tfrac{1}{2}\beta \sin^{2}A'}}{\sqrt{\cos^{2}\tfrac{1}{2}\alpha - \sin A \sin A' \sin^{2}\tfrac{1}{2}\beta}},
\]
whence, considering $\alpha$, $\beta$ as small quantities of the same order, and neglecting terms of the third order, we have $\tan \tfrac{1}{2}E \sin \delta = \sin \tfrac{1}{2}\beta \cos \tfrac{1}{2}\beta \sin A'$, or what is the same thing, = $\sin \tfrac{1}{2}\beta \sin A$, or finally, = $\tfrac{1}{2}\beta \sin A$, that is $\tfrac{1}{2}(h' - h) \sin A$, which is the foregoing approximate value, and thus in the adopted orbit $y = \tfrac{1}{P'} \cdot 12$ in. = approx.\ value. As regards the expression for $\tan \tfrac{1}{2}E \cos \delta$, writing for a moment $n = \sec^{2}\tfrac{1}{2}\alpha \sin A \sin A' \sin^{2}\tfrac{1}{2}\beta$, the quantity under the radical sign is
\[
\frac{\tan^{2}\tfrac{1}{2}\alpha + n - \sin^{2}\tfrac{1}{2}\beta \cos^{2}\tfrac{1}{2}\beta \sin^{2}A'}{1 - n \cos^{2}\tfrac{1}{2}\alpha \cdot (1-n)^{2}},
\]
and, taking this to the third order, it is
\[
= \tan^{2}\tfrac{1}{2}\alpha + n(1 + \tan^{2}\tfrac{1}{2}\alpha) - \frac{\sin^{2}\tfrac{1}{2}\beta \cos^{2}\tfrac{1}{2}\beta \sin^{2}A'}{\cos^{2}\tfrac{1}{2}\alpha},
\]
which, substituting for $n$ its value, is
\[
= \tan^{2}\tfrac{1}{2}\alpha + \sin^{2}\tfrac{1}{2}\beta \sin A' \frac{\sin A - \sin A' \cos^{2}\tfrac{1}{2}\beta}{\cos^{2}\tfrac{1}{2}\alpha},
\]
where $\sin A - \sin A' \cos^{2}\tfrac{1}{2}\beta = \sin A - \sin(A + \alpha) \cos^{2}\tfrac{1}{2}\beta$, or neglecting herein terms of the second order, this is $= \sin A - (\sin A + \sin \alpha \cos A) \cos^{2}\tfrac{1}{2}\beta$, $= -\sin \alpha \cos A$, $= -2\tan \tfrac{1}{2}\alpha \cos^{2}\tfrac{1}{2}\alpha \cos A$, so that to the third order the quantity under the radical sign is
\[
= \tan^{2}\tfrac{1}{2}\alpha - 2\tan^{2}\tfrac{1}{2}\alpha \cdot \frac{\sin^{2}\tfrac{1}{2}\beta \sin A \cos A}{\cos^{2}\tfrac{1}{2}\alpha},
\]
and to the second order, that is finally neglecting terms of the third order,
\[
\tan \tfrac{1}{2}E \cos \delta = \tan \tfrac{1}{2}\alpha - \frac{\sin^{2}\tfrac{1}{2}\beta \sin A \cos A}{\sin \tfrac{1}{2}\alpha \cos \tfrac{1}{2}\alpha},
\]
or, what is the same thing,
\[
= \tfrac{1}{2}\alpha - \sin A \cos A \cdot \tfrac{1}{2}\beta^{2}.
\]

\noindent\textbf{18.} Hence, writing $\alpha = (A' - A)$ arc $1'$, $\beta = (h' - h)$ arc $1'$, and passing to the adopted orbit, we have
\[
x = \frac{A' - A}{P'} \cdot 12 \text{ in.} - \tfrac{1}{2}\sin A \cos A \cdot \frac{12 \text{ in.} \times (h' - h)^{2} \text{ arc } 1'}{P'},
\]
viz.\ putting
\[
y = \frac{h' - h}{P'} \sin A \cdot 12 \text{ in.},
\]
we have
\[
x = \frac{A' - A}{P'} \cdot 12 \text{ in.} - y \cdot \tfrac{1}{2}\cos A \times (h' - h) \text{ arc } 1',
\]
or say
\[
= \frac{A' - A}{P'} \cdot 12 \text{ in.} - y \cdot \tfrac{1}{2}\cos A \cdot (h' - h) \cdot \tfrac{1}{206265}.
\]

The value of the second term may amount to about $\tfrac{1}{10}$ of an inch, and thus be sensible, but there is no difficulty in taking account of it.

\subsubsection*{Article No.\ 19. As to the Equation $r = \dfrac{2p}{p^{2}+1}$}

\noindent\textbf{19.} It may be remarked that the equation $r = \dfrac{2p}{p^{2}+1}$, which served for the graduation of the straight edge, was in effect obtained from the equations
\[
r = P' \tan \tfrac{1}{2}E, \quad P'' \sin z = \sin E, \quad p = \tan \tfrac{1}{2}z
\]
by assuming therein $\tan \tfrac{1}{2}E = E$ and $\sin \tfrac{1}{2}z = \tfrac{1}{2}z$ respectively. But the elimination of $E$ and $z$ can be effected without this assumption, viz.\ we have $\sin E = \dfrac{2\tan \tfrac{1}{2}E}{1 + \tan^{2}\tfrac{1}{2}E} = \dfrac{P'r}{1 + \tfrac{1}{4}P'^{2}r^{2}}$, and then as before, $\sin z = \dfrac{2p}{p^{2}+1}$, $= \dfrac{P'r}{1 + \tfrac{1}{4}P'^{2}r^{2}}$; whence the relation between $r$ and $p$ is found to be
\[
r = \frac{2p}{p^{2}+1},
\]
which however assumes that $P''$ is reckoned in parts of the radius; reckoning it as before in minutes, we must, instead of $P'$, write $P'$ arc $1' = \dfrac{P'\pi}{206265 \times 180}$, viz.\ the numerical value is about $\dfrac{P'}{14400}$, and taking it to be this number, the formula is
\[
r = \frac{2p}{p^{2} + \dfrac{1}{14400}},
\]
where $r$, $p$ are reckoned in parts of the radius (= 12 inches). Supposing that $r_{1}$ is calculated from the formula $r_{1} = \dfrac{2p}{p^{2}+1}$, then we have very nearly $r = r_{1}\left(1 + \dfrac{p^{2}}{14400}\right)$, and $p^{2}$ being = 1 at most, the correction is inappreciable: if however this were not the case, the more accurate formula might have been used; the only difference being that the making of the graduation would have been more laborious.

\subsubsection*{Article No.\ 20. Remark as to the Geometrical Theory of the Projection of the Penumbral Curve.}

\noindent\textbf{20.} The stereographic projection of the penumbral curve on the Earth's surface (assumed to be spherical) is, as I have elsewhere shown, a bicircular quartic. It may be shown that the stereographic projection, as given by the foregoing approximate method, is a bicircular quartic: we have, in fact, a circle, the equation of which in the polar coordinates $r$, $\theta$ is $(r\cos\theta - a)^{2} + (r\sin\theta)^{2} = \beta^{2}$, and where ($\theta$ being unaltered) $r$ is changed into $p$, where $r = \dfrac{2p}{p^{2}+1}$, that is $\dfrac{2}{p} = p + \dfrac{1}{p}$. The equation of the circle is $r^{2} - 2ar\cos\theta + a^{2} - \beta^{2} = 0$, or say $\dfrac{2}{r}\cos\theta - \dfrac{a^{2}-\beta^{2}}{r^{2}} = \dfrac{2a}{r}$, and the transformed equation is therefore $(a^{2}-\beta^{2})(p^{2}+1)^{2} - 4a\cos\theta \cdot p(p^{2}+1) + 4\beta^{2}p^{2} = 0$, or in rectangular coordinates ($\rho^{2} = x^{2} + y^{2}$)
\[
(a^{2}-\beta^{2})(\rho^{2}+1)^{2} - 4ax\rho(\rho^{2}+1) + 4\beta^{2}\rho^{2} = 0,
\]
that is where $\rho^{2} = x^{2} + y^{2}$; the form of the equation shows that the curve is a bicircular quartic. Writing for shortness $\dfrac{a^{2}-\beta^{2}}{a} = m$, the equation is $\rho^{4} - 2m\alpha x \rho^{2} + (2 + 2m)\rho^{2} + 1 - 2m\alpha x = 0$, that is $(\beta^{2} - m(ax-1)+1)^{2} - m^{2}(ax-1)^{2} - 2m = 0$, or, what is the same thing, $\{(x-ma)^{2} + y^{2} - m^{2}a^{2} + m + 1\}^{2} - m^{2}(ax-1)^{2} - 2m = 0$, which putting $x + ma$ for $x$ is $(x^{2} + y^{2} - m^{2}a^{2} + m + 1)^{2} - m^{2}(a(x-ma)-1)^{2} - 2m = 0$, viz.\ the terms of the fourth order being $(x^{2}+y^{2})^{2}$, and there being no terms of the third order, the curve represented by this equation is a bicircular quartic.

\subsubsection*{Article Nos.\ 21 to 30. Practical Details and Application to Eclipse of December 21-22, 1870.}

\noindent\textbf{21.} There are some practical details which it is proper to explain, using to fix the ideas the eclipse of December 21-22, 1870: the constant value of $A$ (see infra) is taken to be $+90^{\circ} + 22^{\circ}35'$. I have a blank projection (radius 12 in.\ as above) with the meridians and parallels each at intervals of $5^{\circ}$. And also another blank form which has on it merely a circle, radius 12 in., graduated as to one quadrant thereof with lines about $\tfrac{1}{2}$ in.\ long, inwards towards the centre. It contains also in a corner the foregoing proportional-scale figure.

\noindent\textbf{22.} On the blank projection I measure off, downwards from the centre, a distance $M_{1}M = 12 \sin 22^{\circ}35'\ (= 4.61)$, distances all in inches; and then with the centre $M$ and radius $= 12 \cos 22^{\circ}35'\ (= 11.08)$, describe a circle which is the bounding circle of the reduced projection. With this same radius I describe on the second form, concentric with the 12-inch circle and above the horizontal diameter thereof, a semicircle: and, cutting out the included area, replace it with tracing cloth. The form thus prepared is placed over the blank projection, so that the semicircle shall coincide with the corresponding semicircle on the blank projection, and the two sheets are fixed together by their lower edges, and by folding down the remaining sides. We have thus the upper half of the reduced projection, represented by the semicircle, with the meridians and parallels marked out thereon by lines seen through the tracing cloth.

See the Plate; the dotted line shows a paper scale afterwards affixed to the second form or upper sheet. Observe that so far the only eclipse-datum made use of is the value $A = 90^{\circ} + 22^{\circ}35'$.

\noindent\textbf{23.} We have for the eclipse in question, taking $t$ for the G.M.T.\ in hours, positive or negative according as the time is after or before G.M.\ Noon, Dec.\ 22, and $h'$ also in hours,
\[
A' = 0'.02 + 0.9996t,
\]
and then taking the values of $\alpha$, $\alpha'$, $A$, $A'$ from the N.A.\ we have as follows:

\begin{center}
\begin{tabular}{|c|c|c|c|c|c|c|c|c|c|}
\hline
G.M.T. & \multicolumn{2}{c|}{$h'+2e'=\alpha$} & $\alpha'$ & $A=90^{\circ}+$ & $A'=90^{\circ}+$ & \multicolumn{2}{c|}{$A'-A$} & \multicolumn{2}{c|}{$h'-h$} \\\cline{2-3}\cline{7-10}
1870, Dec. & & & & & & \multicolumn{2}{c|}{in Minutes} & \multicolumn{2}{c|}{in ditto} \\\cline{2-3}\cline{7-10}
d & h & m & s & $^{\circ}\ '\ ''$ & $^{\circ}\ '\ ''$ & \multicolumn{2}{c|}{of Arc} & \multicolumn{2}{c|}{} \\\hline
21 & 22 & 1 & 13.42 & 17 56 1.84 & 18 1 48.72 & 22 27 & 8.5 & 60.143 & $-$ 86.620 \\\hline
22 & 3 & 5 & 1 7.37 & 18 9 26.74 & 18 2 44.27 & 22 43 & 12.5 & 44.023 & + 100.632 \\\hline
\end{tabular}
\end{center}

and moreover
\begin{center}
\begin{tabular}{ll}
Moon's Parallax & $\varpi = 60'\,38''.6$ \\
Sun's ditto & $\varpi' = 9'.1$ \\
\multicolumn{2}{c}{$\varpi - \varpi' = 60'\,29''.5 = 60'.49$} \\
\multicolumn{2}{c}{$P' = 60'.39$} \\
Moon's Semidiam. & $s = 16'\,33''.2$ \\
Sun's ditto & $s' = 16'\,17''.9$ \\
\multicolumn{2}{c}{$s + s' = 32'\,51''.2 = 32'.85$} \\
\end{tabular}
\end{center}

whence
\[
\frac{s+s'}{P'} \cdot 12 \text{ in.} = 6.53,
\]
viz.\ this is the radius of the circles used in the construction of the penumbral curves.

\noindent\textbf{24.} We have for $x$, $y$ the formulae
\[
x = \frac{A' - A}{P'} \cdot 12 \text{ in.} + 2y(h'-h) \cdot 0.00006, \qquad y = \frac{h'-h}{P'} \sin A \cdot 12 \text{ in.},
\]
viz.\ I find Dec.\ 21, 22$^{\text{h}}$,
\[
x = 11.95 - 0.06 = 11.89, \qquad y = -15.80,
\]
and Dec.\ 22---3$^{\text{h}}$,
\[
x = 8.75 + 0.11 = 8.86, \qquad y = +18.45,
\]
where I have taken account of the small corrections to the approximate values of $x$: it may be added that, using the conjunction-value $52'\,9''.4$ of $A' - A$, we have at conjunction, $x = 10.36$, $y = 0$.

\noindent\textbf{25.} We thus lay down on the relative orbit the two points 22$^{\text{h}}$ and 3$^{\text{h}}$, and the point of conjunction or intersection with the axis of $x$; the three points are found to be sensibly in a straight line: the distance between the extreme points is about 34 inches, representing 5 hours, so that the scale is nearly 7 inches to an hour: the line is then graduated to quarters of an hour. We then, by means of the distance $12 + 6.53 = 18.53$, mark off on the relative orbit, the points $B$, $E$, which correspond to the beginning and end of the eclipse respectively: the times as read off from my figure, and compared with the true times given in the N.\ A.\ are

\begin{center}
\begin{tabular}{lcc}
 & From Figure & N.\ A. \\
Beginning & 22$^{\text{h}}$\ 12$^{\text{m}}$.5 & 22 13.6 \\
End & 2 40 .5 & 2 41.1 \\
\end{tabular}
\end{center}

\noindent\textbf{26.} With centre $B$ describing a circle radius 6.53 this will of course just touch the 12-inch circle, and the penumbral curve will be a mere point, viz., this is the point $B'$ on the bounding circle, opposite to the point of contact. And so with centre $E$ describing a circle of the same radius 6.53, that will just touch the 12-inch circle, and the penumbral curve will be a mere point, viz., the point $E'$ on the bounding circle, opposite the point of contact.

\noindent\textbf{27.} I draw the circles corresponding to the times 22$^{\text{h}}$\ 30$^{\text{m}}$, 45$^{\text{m}}$, 23$^{\text{h}}$\ 0$^{\text{m}}$, viz., so much of each as lies within the 12-inch circle. Each of these is then transformed into a penumbral curve, drawn in the upper semicircle on the tracing cloth. For this purpose we construct a straight edge of paper, the one half graduated for $12 \sin \theta$, the other half for $11.08 \tan \tfrac{1}{2}\theta$, by means of the proportional-scale figure, as already explained: $\theta = 0^{\circ}$ to $90^{\circ}$ at intervals of $5^{\circ}$, is quite sufficient; the points on any particular penumbral curve are laid down in pairs with the utmost facility, and the curve is traced by hand from 4 or 5 pairs of points.

\noindent\textbf{28.} We then graduate for latitude; viz., we see through the tracing cloth, the equator cutting the vertical radius in $Q$, and a parallel cutting the same radius, say in $R$; drawing lines from $C$, we refer these to the points $q$, $r$ on the bounding circle, viz., on the quadrant thereof which is graduated by means of the graduation-lines of the 12-inch circle; and we thus read off the latitude of the parallel in question; this latitude is then marked for each parallel on the vertical radius from $Q$ up to the bounding circle, viz., not on the tracing cloth, but on the paper affix; and we then on this same radius (on the paper affix) interpolate the positions where this would be intersected by the parallels for the colatitudes, $5^{\circ}$, $10^{\circ}$, $15^{\circ}$, \&c.\ Or (what is perhaps better) we may without marking the latitudes of the parallels of the blank form, construct directly the last-mentioned graduations; viz., marking off on the bounding circle from the point $q$, equal intervals of $5^{\circ}$, and from any such mark drawing to $C$, a line to meet the vertical radius, the point of intersection is the point belonging to the parallel, latitude equal to the corresponding multiple of $5^{\circ}$.

\noindent\textbf{29.} Finally, we must (not on the tracing cloth but on the paper affix) graduate an arc of the equator for the position of the meridian of Greenwich, that is for $h$.

We have
\begin{align*}
\text{At } 22^{\text{h}},\quad h &= -2^{\text{h}} + 7\ 0.30 = -1^{\text{h}}\ 52^{\text{m}}\ 52^{\text{s}}.30 = -28^{\circ}\,13'.08,\\
\text{At } 3^{\text{h}},\quad h &= -2^{\text{h}} + 4\ 54\ 24.90 = +2^{\text{h}}\ 54^{\text{m}}\ 24^{\text{s}}.90 = +43^{\circ}\,36'.20.
\end{align*}
The equator is already graduated in longitude by means of the meridians of the blank projection: hence we lay down the marks for 22$^{\text{h}}$ and 3$^{\text{h}}$ in the positions belonging to $-28^{\circ}\,13'$, and $+43^{\circ}\,36'$ respectively. And then since the interval of 5 hours answers to $71^{\circ}\,49'$, that of 1 hour will answer to $14^{\circ}\,22'$, so that, measuring off these intervals of longitude, we have the marks for the intermediate times 23$^{\text{h}}$, 0$^{\text{h}}$, 1$^{\text{h}}$, 2$^{\text{h}}$; or it might be proper to find in this way the marks corresponding to each interval of 20$^{\text{m}}$ of time, answering to about $5^{\circ}$ of longitude; the further subdivisions would be proportional to the intervals of time.

\noindent\textbf{30.} I have in this way read off the positions of the points $B'$ and $E'$ belonging to the beginning and end of the eclipse; the values, as compared with the true ones, are

\begin{center}
\begin{tabular}{lcc}
 & From Figure & N.\ A. \\
$B'$ latitude N. & $34^{\circ}$ & $35^{\circ}\,37'$ \\
\quad longitude W. & $47^{\circ}$ & $45^{\circ}\,44'$ \\
$E'$ latitude N. & $26^{\circ}$ & $26^{\circ}\,5'$ \\
\quad longitude W. & $38^{\circ}$ & $37^{\circ}\,16'$ \\
\end{tabular}
\end{center}

I remark that my figure, although drawn carefully, is not drawn with anything like the degree of accuracy which would be easily attainable; and I think that far better results might be obtained. I merely from a scale lay down tenths and estimate hundredths of an inch, but certainly fiftieths might be laid down from a scale.

\newpage

\section*{Paper 478}
\addcontentsline{toc}{section}{Paper 478 --- On the Geodesic Lines on an Ellipsoid}

\subsection*{On the Geodesic Lines on an Ellipsoid}

\begin{center}
[From the \textit{Memoirs of the Royal Astronomical Society}, vol.~xxxix. (1872), pp.~31---53. Read January 13, 1871.]
\end{center}

The fundamental equations, in regard to the geodesic lines on an ellipsoid, were established by Jacobi, viz., representing by $a$, $b$, $c$, the squares of the semiaxes, that is, taking the ellipsoid to be
\[
\frac{x^{2}}{a} + \frac{y^{2}}{b} + \frac{z^{2}}{c} = 1
\]
(where $a > b > c$), if we introduce the elliptic coordinates $h$, $k$, and write
\begin{align*}
\frac{x^{2}}{a+h} + \frac{y^{2}}{b+h} + \frac{z^{2}}{c+h} &= 1, \\
\frac{x^{2}}{a+k} + \frac{y^{2}}{b+k} + \frac{z^{2}}{c+k} &= 1,
\end{align*}
or, what is the same thing,
\begin{align*}
x^{2} &= \frac{a(a+h)(a+k)}{(a-b)(a-c)}, \\
y^{2} &= \frac{b(b+h)(b+k)}{(b-c)(b-a)}, \\
z^{2} &= \frac{c(c+h)(c+k)}{(c-a)(c-b)},
\end{align*}
then, if $\beta$ be an arbitrary constant, the differential equation of a geodesic line is
\begin{equation}\tag{1}
\text{const.} = \int \frac{dh}{h+\beta} \sqrt{\frac{(a+h)(b+h)(c+h)}{(\beta+h)}} + \int \frac{dk}{k+\beta} \sqrt{\frac{(a+k)(b+k)(c+k)}{(\beta+k)}},
\end{equation}
and the expression for the length of any arc of the curve is given by
\begin{equation}\tag{2}
s = \int \sqrt{\frac{(a+h)(b+h)(c+h)}{h+\beta}} dh + \int \sqrt{\frac{(a+k)(b+k)(c+k)}{k+\beta}} dk.
\end{equation}
I propose in the present Memoir to develope the theory to the extent of showing how we can, by means of the first of these equations, explain the course of the geodesic lines; and for given numerical values of $a$, $b$, $c$, calculate, construct, and exhibit in a drawing the course of these lines: I attend more particularly to the series of geodesic lines through an umbilicus (which lines pass also through the opposite umbilicus), and to the case where the semiaxes are connected by the equation $ac - b^{2} = 0$, a relation which simplifies the formulae.

\subsubsection*{General Considerations as to the Course of the Lines.}

\noindent\textbf{1.} It will be observed that $h$ and $k$ enter into the formulae symmetrically: it will be convenient to distinguish between these coordinates by considering $h$, as extending between the values $-a$, $-b$; and $k$ as extending between the values $-b$, $-c$.

Thus:

$h = \text{const.}$ denotes a curve of curvature of the one kind, viz.:

$h = -a$, the principal section $ABA'$ (or major-mean section),

$h = -b$, the curves $UU'$ and $U''U'''$ (or portions of the umbilicar section $ACA'C$);

similarly, $k = \text{const.}$ denotes a curve of curvature of the other kind, viz.:

$k = -c$, the principal section $CBC$ (or minor-mean section),

$k = -b$, the curves $UU'''$ and $U'U''$ (remaining portions of the umbilicar section $ACA'C$).

\noindent\textbf{2.} To any given (admissible) values of $h$, $k$ there correspond eight points, situate in the eight octants of the surface respectively; but, unless the contrary is expressed, it is assumed that the coordinates $x$, $y$, $z$, are positive, and that the point is situate in the octant $ABC$.

\noindent\textbf{3.} The constant $\beta$ may have any value from $+a$ to $+c$; viz., if it has a value between $a$ and $b$, or say, if $-\beta$ has an $h$-value, then the geodesic lines lie wholly between the two ovals of the curve of curvature $h = -\beta$ (being in general an indefinite undulating curve touching each oval an indefinite number of times). Similarly, if $\beta$ has any value between $b$ and $c$, or say, if $-\beta$ has a $k$-value, then the geodesic line lies wholly between the two ovals of the curve of curvature $k = -\beta$ (being in general an indefinite undulating curve touching each oval an indefinite number of times).

The intermediate case is when $\beta = b$, or say when $-\beta$ has the umbilicar value: here the geodesic line is in general an indefinite undulating curve passing an infinite number of times through the opposite umbilici $U$, $U''$, or $U'$, $U'''$; to fix the ideas, say through $U$, $U''$.

\subsubsection*{Lines through an Umbilicus.}

\noindent\textbf{4.} I attend in particular to the last-mentioned case, and thus write $\beta = b$. We may in the formula (1) fix at pleasure a limit of each integral; and writing for convenience
\[
\Pi(h) = \int \frac{dh}{b+h} \sqrt{\frac{(a+h)(c+h)}{h}}, \qquad \Psi(k) = \int \frac{dk}{b+k} \sqrt{\frac{(a+k)(c+k)}{k}},
\]
the equation (1) becomes
\[
\text{Const.} = \Pi(h) + \Psi(k).
\]

\noindent\textbf{5.} It is to be observed, in regard to these integrals, that writing $h = -a + u$, we have
\[
\Pi(h) = \int \frac{du}{b-a+u} \sqrt{\frac{u(a-c+u)}{-a+u}},
\]
which, for $u$ small, is $= \dfrac{2\sqrt{u}}{\sqrt{a-b}\sqrt{a-c}}$. By the assistance of this formula the value of the integral may be calculated by quadratures; viz., the formula gives the integral for any small value of $u$, and we can then proceed by the method of quadratures. The integral becomes infinite for $h = -b$: suppose that we have by quadratures calculated it up to $h = -b - m$ ($m$ small), then to calculate it up to any value $-b - m + u$ nearer to $-b$, we have
\begin{align*}
\Pi(h) &= \Pi(-b-m) + \int \frac{du}{b+m-u} \sqrt{\frac{(a-b-m+u)(b-c+m-u)}{-b-m+u}} \\
&= \Pi(-b-m) + \sqrt{(a-b)(b-c)} \int \frac{du}{u} \\
&= \Pi(-b-m) - \sqrt{(a-b)(b-c)} \log u,
\end{align*}
where the second term is positive, and the value thus increases slowly with $u$, becoming as it should do $= \infty$ for $u = m$ or $h = -b$.

$^{\dagger}$ Except when the contrary is stated, the symbol ``log'' denotes throughout the hyperbolic logarithm.

\noindent\textbf{6.} Similarly in the second integral writing $k = -c - v$, we have
\[
\Psi(k) = \int \frac{-dv}{b-c-v} \sqrt{\frac{(a-c-v)v}{-c-v}},
\]
which, for $v$ small, is $= \dfrac{2\sqrt{v}}{\sqrt{b-c}\sqrt{a-c}}$, which is of the like assistance in regard to the calculation by quadratures. And if we have by quadratures calculated the integral up to $k = -b + n$ ($n$ small), then, to calculate it up to any value $-b + n - v$ nearer to $-b$, we have
\[
\Psi(k) = \Psi(-b+n) + \sqrt{(a-b)(b-c)} \int \frac{dv}{v},
\]
where the second term is positive, and the value thus increases slowly with $v$, becoming as it should do $= \infty$ for $v = n$, or $k = -b$.

\noindent\textbf{7.} It may be remarked that in $\Pi(h)$ and $\Psi(k)$ respectively the coefficient of the logarithmic term has in each case the same value $= \sqrt{(a-b)(b-c)}$; as regards the initial terms $\sqrt{u}$ and $\sqrt{v}$, the coefficients are $\tfrac{2}{\sqrt{a-b}\sqrt{a-c}}$ and $\tfrac{2}{\sqrt{b-c}\sqrt{a-c}}$ respectively, which are equal if $\dfrac{1}{a-b} = \dfrac{1}{b-c}$, or $ac - b^{2} = 0$.

\noindent\textbf{8.} We may consider the two geodesic lines $\Pi(h) \pm \Psi(k) = \text{const.}$; suppose that these each of them pass through the point $P$, coordinates $(h_{0}, k_{0})$ in the $ABC$ octant of the ellipsoid; then for one of them we have $\Pi(h) - \Psi(k) = \Pi(h_{0}) - \Psi(k_{0})$, and for the other of them we have $\Pi(h) + \Psi(k) = \Pi(h_{0}) + \Psi(k_{0})$: I attend first to the former of these, say $\Pi(h) - \Psi(k) = C$ (where $C$ is $= \Pi(h_{0}) - \Psi(k_{0})$); and I say that this denotes the curve $UPU''$. In fact, by reason of the equation $\Pi(h)$ and $\Psi(k)$ must both increase or both diminish; they both increase as $h$ passes from $h_{0}$ to $-b$, and as $k$ passes from $k_{0}$ to $-b$: we may have $h = -b + u$, $k = -b + v$ where $u$ and $v$ are both indefinitely small, the functions $\Pi$ and $\Psi$ being then indefinitely large, but $\Pi - \Psi = C$; and we have thus a series of points nearer and nearer to the umbilicus $U$; that is, we have the portion $PU$ of the curve. Tracing the curve in the opposite direction, or considering $h$ as passing from $h_{0}$ to $-a$, and $k$ as passing from $k_{0}$ to $-c$, then if $C$ be positive, $k$ will attain the value $-c$, before $h$ attains the value $-a$, say that we have simultaneously $h = h_{1}$, $k = -c$; the equation is $\Pi(h_{1}) - \Psi(-c) = C$, that is, $\Pi(h_{1}) = C$; and the geodesic line then arrives at a point $P_{1}$ on the arc $CB$ of the minor-mean principal section. The function $\Psi$ then changes its sign, viz., considering it as always positive, the equation is now $\Pi(h) + \Psi(k) = C$, $k$ passing from the value $-c$ towards $-b$, that is, $\Psi(k)$ increasing, and therefore $\Pi(h)$ diminishing, or $h$ passing from $h_{1}$ towards the value $-a$; until at last, say for $k = k_{2}$, we have $h = -a$, that is, $C = \Pi(-a) + \Psi(k_{2})$, or $C = \Psi(k_{2})$; the geodesic line here arrives at a point $P_{2}$ on the arc $BA'$ of the major-mean principal section. The function $\Pi$ then changes its sign, viz., $\Pi$ denoting a positive function as before, the equation is $-\Pi(h) + \Psi(k) = C$; $h$ passes from $-a$ towards $-b$, that is $\Pi(h)$ increases, and therefore $\Psi(k)$ must also increase, or $k$ pass from $k_{2}$ towards $-b$: we have at length $h = -b - u$, $k = -b + v$, $u$ and $v$ being each indefinitely small; and therefore $\Pi$ and $\Psi$ each indefinitely large (but $-\Pi + \Psi = C$): that is, we arrive at the umbilicus $U''$, completing the geodesic line $UPU''$.

\noindent\textbf{9.} If instead of $C = +$ we have $C = -$, everything is similar, but the geodesic line proceeding from $U$ in the direction $UP$ will first cut the arc $BA$ of the major-mean section at a point $P_{1}$; then the arc $BC$ of the minor-mean section at a point $P_{2}$; and, finally, arrive as before at the umbilicus $U''$.

\noindent\textbf{10.} The intermediate case is when $C = 0$, viz.\ we have here $\Pi(h) - \Psi(k) = 0$; the geodesic line here passes from $U$ in the direction $UP$ to $B$ (extremity of the mean axis, $h = -a$, $k = -c$); $\Pi$ and $\Psi$ then each change their sign, so that, considering them as positive, the equation still is $\Pi(h) - \Psi(k) = 0$, and the geodesic line at last arrives at the umbilicus $U''$. It will be easily understood how in the like manner $\Pi(h) + \Psi(k) = 0$ refers to the line $U'PU'''$.

\noindent\textbf{11.} Reverting to the equation $\Pi(h) - \Psi(k) = C$, or as I will now write it $\Pi(h) - \Psi(k) = \Pi(h_{0}) - \Psi(k_{0})$, which belongs to the portion $UP$ of the geodesic line $UPU''$, we require when $h$ is $= -b - u$, and $k = -b + v$ ($u$ and $v$ indefinitely small) to know the ratio of the increments $u$, $v$; this in fact serves to determine the direction at $U$ of the geodesic line through the given point $(h_{0}, k_{0})$.

\noindent\textbf{12.} For this purpose writing $h = -b - u$, we find
\begin{multline*}
\Pi(h) - \Pi(-b) = \int_{0}^{u} \frac{du}{u} \sqrt{\frac{b+u}{(a-b-u)(b-c+u)}} \\
= \frac{1}{\sqrt{(a-b)(b-c)}} \Bigg\{ \int_{0}^{u} \frac{du}{u} \sqrt{\frac{b+u}{(a-b-u)(b-c+u)}} - \sqrt{\frac{b}{(a-b)(b-c)}} \Bigg\} \\
+ \frac{1}{\sqrt{(a-b)(b-c)}} \log u,
\end{multline*}
and, when $u$ is indefinitely small, this is
\begin{multline*}
= \frac{1}{\sqrt{(a-b)(b-c)}} \Bigg\{ \frac{\sqrt{b+u}}{\sqrt{(a-b-u)(b-c+u)}} - \sqrt{\frac{b}{(a-b)(b-c)}} \Bigg\} \\
+ \frac{1}{\sqrt{(a-b)(b-c)}} \log u.
\end{multline*}

Similarly, when $k = -b + v$, where $v$ is indefinitely small,
\[
\Psi(k) - \Psi(-b) = \frac{1}{\sqrt{(a-b)(b-c)}} \Bigg\{ \frac{\sqrt{b-v}}{\sqrt{(a-b+v)(b-c-v)}} - \sqrt{\frac{b}{(a-b)(b-c)}} \Bigg\}
\]
\[
\qquad \qquad + \frac{1}{\sqrt{(a-b)(b-c)}} \log v.
\]

\noindent\textbf{13.} Each of the integrals is of the dimension $-\tfrac{1}{2}$ in $a$, $b$, $c$, and the difference of the integrals may be represented by $\tfrac{1}{\sqrt{(a-b)(b-c)}}$; we have therefore
\[
\Pi(h) - \Psi(k) = \frac{1}{\sqrt{(a-b)(b-c)}} \left[ M' + \log \frac{v}{u} \right],
\]
where
\[
M' = \int_{0}^{u} \frac{du}{u} \Bigg\{ \frac{\sqrt{b+u}}{\sqrt{(a-b-u)(b-c+u)}} - \sqrt{\frac{b}{(a-b)(b-c)}} \Bigg\}
\]
\[
\qquad - \int_{0}^{v} \frac{dv}{v} \Bigg\{ \frac{\sqrt{b-v}}{\sqrt{(a-b+v)(b-c-v)}} - \sqrt{\frac{b}{(a-b)(b-c)}} \Bigg\}.
\]

\noindent\textbf{14.} Suppose the inferior limits replaced by the indefinitely small positive quantities $\epsilon$, $\epsilon'$ respectively; and for the variable in the second integral write $-u$; then
\[
M' = \int \frac{du}{u} \left\{ \frac{\sqrt{b+u}}{\sqrt{(a-b-u)(b-c+u)}} - \frac{\sqrt{b}}{\sqrt{(a-b)(b-c)}} \right\},
\]
it being understood that the values $u = -\epsilon'$ to $u = +\epsilon$ are omitted from the integration; this is
\[
= \frac{1}{\sqrt{(a-b)(b-c)}} \int_{-\epsilon}^{\epsilon} \frac{du}{u} \sqrt{\frac{b+u}{(a-b-u)(b-c+u)}} - \frac{b}{\sqrt{(a-b)(b-c)}} \cdot \frac{a-b}{b-c} \log \frac{\epsilon}{\epsilon'},
\]
with the same convention as to the integral; or if $\epsilon' = \epsilon$, then
\[
M' = \int_{-\epsilon}^{\epsilon} \frac{du}{u} \left\{ \frac{\sqrt{b+u}}{\sqrt{(a-b-u)(b-c+u)}} - \frac{\sqrt{b}}{\sqrt{(a-b)(b-c)}} \right\} = \int_{-b}^{h} \frac{dh}{b+h} \sqrt{\frac{(a+h)(c+h)}{-h}},
\]
where the omitted elements being from $u = -\epsilon$ to $u = +\epsilon$; that is (in the language of Cauchy) we take for the integral its principal value. And hence
\[
\Pi(h) - \Psi(k) = \frac{1}{\sqrt{(a-b)(b-c)}} \left[ M' + \log \frac{v}{u} \right].
\]

\noindent\textbf{15.} By what precedes this is $= \Pi(h_{0}) - \Psi(k_{0})$; or if we write simply $(h, k)$ instead of $(h_{0}, k_{0})$, that is, consider the geodesic line $UP$, which is drawn from the point $P$, coordinates $(h, k)$, to the umbilicus $U$, the coordinates of a point consecutive to the umbilicus are $-b-u$, $-b+v$, where $u$, $v$ are connected by the last-mentioned equation, in which $M'$ is a transcendental function depending on $(a, b, c)$ but independent of the particular geodesic line.

\noindent\textbf{16.} If for the geodesic line through the point $B$, or say for the $B$-geodesic, $\dfrac{v}{u} = \dfrac{v_{0}}{u_{0}}$, then $M' = -\log \dfrac{v_{0}}{u_{0}}$, and we have in general
\[
\log \frac{v}{u} = \log \frac{v_{0}}{u_{0}} + M' \sqrt{(a-b)(b-c)},
\]
a result which I proceed to further transform as follows:

If $x_{0}$, $y_{0}$, $z_{0}$ refer to the umbilicus $U$, then considering first the consecutive point $P$ on the geodesic line (coordinates $-b-u$, $-b+v$) and next the consecutive point $Q$ on the umbilicar section, we have for these two points respectively,
\begin{align*}
dx_{0} &= \sqrt{\frac{b(a-b)(a-c)}{a}} \sqrt{uv}, \quad &dy_{0} &= \sqrt{\frac{ab}{b-c}} \sqrt{uv}, \quad &dz_{0} &= \sqrt{\frac{b-c}{a}} \sqrt{ac} \cdot \frac{v-u}{2\sqrt{a}}, \\
dx'_{0} &= 0, \quad &dy'_{0} &= 0, \quad &dz'_{0} &= -\sqrt{\frac{c(b-c)}{a}} \cdot \frac{u+v}{2\sqrt{a}}.
\end{align*}
say these are $\alpha$, $\beta$, $\gamma$, and $\alpha'$, $\beta'$, $\gamma'$; and then
\[
\alpha\alpha' + \beta\beta' + \gamma\gamma' = \frac{1}{4}\left( \frac{c}{(a-b)(a-c)} + \frac{c}{(b-c)(a-c)} \right)(v-u) = \frac{c(v-u)}{(a-b)(b-c)},
\]
\begin{align*}
\alpha^{2} + \beta^{2} + \gamma^{2} &= uv \left( \frac{b(a-b)(a-c)}{a} + \frac{ab}{b-c} \right) + \frac{(b-c)ac(v-u)^{2}}{4a^{2}} \\
&= \frac{b(v-u)}{(a-b)(b-c)},
\end{align*}
whence
\[
\cos \theta = \frac{\alpha\alpha' + \beta\beta' + \gamma\gamma'}{\sqrt{\alpha^{2}+\beta^{2}+\gamma^{2}}\sqrt{\alpha'^{2}+\beta'^{2}+\gamma'^{2}}} = \frac{U-V}{U+V},
\]
that is, $\cos(180^{\circ} - \theta) = \dfrac{U-V}{U+V}$, or $\tan^{2}\tfrac{1}{2}\theta = \dfrac{V}{U}$, where if $U$ is the umbilicus, $P$ the consecutive point $-b-u$, $-b+v$, and $UQ$ the element of the umbilicar principal section, $\phi = \angle PUQ$, $180^{\circ} - \theta = \angle PUQ'$. For the $B$-geodesic we have $2\log \tan \tfrac{1}{2}\phi_{0} = \log \dfrac{V_{0}}{U_{0}} = M'$.

\noindent\textbf{17.} The foregoing equation for $\Pi(h) - \Psi(k)$ now becomes
\[
\log \tan \tfrac{1}{2}\phi = \log \tan \tfrac{1}{2}\phi_{0} + M' \sqrt{(a-b)(b-c)},
\]
viz., $\phi_{0}$ is the south azimuth of the $B$-geodesic at the umbilicus, a mere function of $(a, b, c)$ and $\phi$ is the south azimuth at the umbilicus, of the geodesic line under consideration, so that we may consider the geodesic line to be determined by the south azimuth $\phi$ as its parameter.

\subsubsection*{Formulae for the case $ac - b^{2} = 0$.}

\noindent\textbf{18.} I annex the following investigation in regard to the case $ac - b^{2} = 0$.

We have in general
\begin{align*}
\frac{1}{\sqrt{(a-b)(b-c)}} \frac{d}{dx} \log \frac{\sqrt{-x(a-b)(b-c)} + \sqrt{-b(a+x)(c+x)}}{\sqrt{-x(a-b)(b-c)} - \sqrt{-b(a+x)(c+x)}} \\
= \frac{1}{\sqrt{x(a+x)(c+x)}} \left( \frac{1}{b+x} - \frac{b}{bx+ac} \right).
\end{align*}
In fact, denoting the logarithm by $\log \dfrac{P+Q}{P-Q}$, we have
\[
\frac{d}{dx} \log \frac{P+Q}{P-Q} = \frac{2(PQ' - P'Q)}{P^{2}-Q^{2}},
\]
where
\begin{align*}
2(PQ' - P'Q) &= 2PQ \left[ \frac{Q'}{Q} - \frac{P'}{P} \right] = \sqrt{x(a+x)(c+x)} \sqrt{b(a-b)(b-c)} \left[ \frac{1}{b+x} - \frac{1}{x} \right] \\
&= \frac{\sqrt{b(a-b)(b-c)}}{\sqrt{x(a+x)(c+x)}} (x^{2} - ac);
\end{align*}
that is
\begin{align*}
P^{2} - Q^{2} &= -x(a-b)(b-c) + b(a+x)(c+x) = (bx+ac)(b+x); \\
\frac{2(PQ'-P'Q)}{P^{2}-Q^{2}} &= \frac{\sqrt{b(a-b)(b-c)}}{\sqrt{x(a+x)(c+x)}} \frac{x^{2}-ac}{(bx+ac)(b+x)} \\
&= \frac{\sqrt{b(a-b)(b-c)}}{\sqrt{x(a+x)(c+x)}} \left[ \frac{1}{b+x} - \frac{b}{bx+ac} \right],
\end{align*}
which proves the theorem.

\noindent\textbf{19.} Hence in the particular case $ac = b^{2}$ we have
\begin{align*}
\Pi(h) &= \frac{1}{b\sqrt{a}} \log \frac{\sqrt{-h(a-b)(b-c)} + \sqrt{-b(a+h)(c+h)}}{\sqrt{-h(a-b)(b-c)} - \sqrt{-b(a+h)(c+h)}} \\
&= \frac{1}{b\sqrt{a}} \log \frac{\sqrt{-h(a-b)(b-c)} + \sqrt{-b(a+h)(c+h)}}{\sqrt{(a-b)(b-c)}} \cdot \sqrt{\frac{-h}{(a+h)(c+h)}} \\
&= \frac{1}{b\sqrt{a}} \left[ \int \frac{dh}{h\sqrt{(a+h)(c+h)}} + \frac{1}{\sqrt{(a-b)(b-c)}} \log \frac{1+H}{1-H} \right],
\end{align*}
where
\[
H^{2} = \frac{b}{a-b} \cdot \frac{(a+h)(c+h)}{(b-c) \cdot (-h)},
\]
viz.\ we see that $\Pi(h)$ depends on the more simple integral
\[
\int \frac{dh}{h\sqrt{(a+h)(c+h)}}.
\]

\noindent\textbf{20.} Similarly
\begin{align*}
\Psi(k) &= \frac{1}{\sqrt{b(a-b)(b-c)}} \log \frac{\sqrt{-k(a-b)(b-c)} + \sqrt{b(a+k)(c+k)}}{\sqrt{-k(a-b)(b-c)} - \sqrt{b(a+k)(c+k)}} \\
&= \frac{1}{b\sqrt{c}} \left[ \int \frac{dk}{k\sqrt{(a+k)(c+k)}} + \frac{1}{\sqrt{(a-b)(b-c)}} \log \frac{1+K}{1-K} \right],
\end{align*}
where
\[
K^{2} = \frac{b}{b-c} \cdot \frac{(a+k)(c+k)}{(a-b) \cdot k},
\]
that is, $\Psi(k)$ depends on the more simple integral,
\[
\int \frac{dk}{k\sqrt{(a+k)(c+k)}}.
\]
Write $h = -b-u$, $k = -b+v$, where $u$ and $v$ are indefinitely small, then
\[
H^{2} = \frac{b}{a-b} \cdot \frac{(a-b-u)(b-c+u)}{u} \cdot \frac{1}{b-c} = \frac{b}{(a-b)(b-c)} \cdot \frac{1}{u} \left( 1 - \frac{u}{a-b} \right) \left( 1 + \frac{u}{b-c} \right),
\]
where $1 + \pi_{h}u$ is a factor, and
\[
\pi_{h} = -\frac{1}{a-b} + \frac{1}{b-c} = \frac{a-c}{(a-b)(b-c)}.
\]
Similarly
\[
K^{2} = \frac{b}{(a-b)(b-c)} \cdot \frac{1}{v} \left( 1 + \frac{v}{a-b} \right) \left( 1 - \frac{v}{b-c} \right),
\]
where $1 + \pi_{k}v$ is a factor, and
\[
\pi_{k} = \frac{1}{a-b} - \frac{1}{b-c} = -\pi_{h}.
\]

\noindent\textbf{21.} Comparing with the result obtained for the general case the two agree, if only
\[
\int \frac{dh}{b+h} \sqrt{\frac{(a+h)(c+h)}{-h}} = \int \frac{dh}{h\sqrt{(a+h)(c+h)}},
\]
where on the left-hand side the integral has its principal value: a result which must therefore hold good when $ac = b^{2}$.

\subsubsection*{Calculation of the Umbilicar Geodesics for Ellipsoid $a:b:c = 4:2:1$.}

\noindent\textbf{22.} As a specimen of the way in which we may, on a given ellipsoid, calculate the course of a geodesic line, I take the semiaxes to be as $2:\sqrt{2}:1$, or, for convenience, $a = 1000$, $b = 500$, $c = 250$; and, considering the geodesic lines through the umbilicus, I calculate by quadratures the functions
\[
\Pi(h) = \int_{-1000}^{h} \frac{-dh}{(500+h)\sqrt{(1000+h)(250+h)}}, \qquad \Psi(k) = \int_{-250}^{k} \frac{dk}{(500+k)\sqrt{(1000+k)(250+k)}}.
\]
The results do not pretend to minute accuracy: I have not attempted to estimate or correct for any error occasioned by the intervals (10 units) being too large; and there may possibly be accidental errors.

\paragraph*{Table I.}

\begin{center}
\begin{tabular}{|r|r|r||r|r|r||r|r|r|}
\hline
$-h$ & $n'$ & $\Pi(h)$ & $-h$ & $n'$ & $\Pi(h)$ & $-h$ & $n'$ & $\Pi(h)$ \\\hline
1000 & $\infty$ & 0 & 840 & 27.6 & 6746 & 630 & 51.5 & 13972 \\
999 & 231.4 & 462 & 830 & 27.8 & 7023 & 620 & 54.1 & 14499 \\
998 & 164 & 659 & 820 & 27.9 & 7301 & 610 & 59.2 & 15066 \\
997 & 134.2 & 809 & 810 & 28.1 & 7582 & 600 & 65.5 & 15689 \\
996 & 116.5 & 934 & 800 & 28.4 & 7865 & 590 & 72.1 & 16377 \\
995 & 104.4 & 1044 & 790 & 28.7 & 8151 & 580 & 80.9 & 17142 \\
990 & 74.6 & 1492 & 780 & 29.2 & 8440 & 570 & 93.0 & 18011 \\
980 & 54.0 & 2135 & 770 & 29.7 & 8735 & 560 & 106.8 & 19010 \\
970 & 45.1 & 2630 & 760 & 30.3 & 9035 & 550 & 127.6 & 20183 \\
960 & 40.1 & 3056 & 750 & 31.0 & 9341 & 540 & 159.1 & 21616 \\
950 & 36.6 & 3439 & 740 & 31.8 & 9655 & 530 & 211.5 & 23469 \\
940 & 34.2 & 3794 & 730 & 32.6 & 9977 & 520 & 316.7 & 26111 \\
930 & 32.5 & 4127 & 720 & 33.6 & 10308 & 510 & 632.7 & 30858 \\
920 & 31.2 & 4446 & 710 & 34.7 & 10650 & 505 & 1265 & 35602 \\
910 & 30.2 & 4753 & 700 & 36.0 & 11004 & 504 & 1581.2 & 37014 \\
900 & 29.4 & 5051 & 690 & 37.4 & 11371 & 503 & 2107.7 & 38834 \\
890 & 28.8 & 5342 & 680 & 39.0 & 11754 & 502 & 3162.3 & 41398 \\
880 & 28.4 & 5628 & 670 & 40.8 & 12153 & 501 & 6324.1 & 45792 \\
870 & 28.1 & 5911 & 660 & 42.0 & 12567 & 500 & $\infty$ & $\infty$ \\
860 & 27.9 & 6190 & 650 & 45.4 & 13005 & & & \\
850 & 27.8 & 6469 & 640 & 48.2 & 13473 & & & \\
\hline
\end{tabular}
\end{center}

\paragraph*{Table II.}

\begin{center}
\begin{tabular}{|r|r|r||r|r|r||r|r|r|}
\hline
$-k$ & $\Psi'$ & $\Psi(k)$ & $-k$ & $\Psi'$ & $\Psi(k)$ & $-k$ & $\Psi'$ & $\Psi(k)$ \\\hline
250 & $\infty$ & 0 & 320 & 45.5 & 4655 & 440 & 107.1 & 12207 \\
251 & 232.2 & 462 & 330 & 46.1 & 5114 & 450 & 127.9 & 13383 \\
252 & 165.5 & 661 & 340 & 47.3 & 5581 & 460 & 159.2 & 14818 \\
253 & 136.0 & 811 & 350 & 48.9 & 6062 & 470 & 211.6 & 16673 \\
254 & 118.6 & 939 & 360 & 51.1 & 6562 & 480 & 316.7 & 19314 \\
255 & 106.8 & 1051 & 370 & 53.8 & 7086 & 490 & 632.7 & 24062 \\
260 & 78.1 & 1514 & 380 & 57.2 & 7641 & 495 & 1265.0 & 28806 \\
270 & 60.5 & 2207 & 390 & 61.4 & 8235 & 496 & 1581.2 & 30218 \\
280 & 51.7 & 2768 & 400 & 66.7 & 8875 & 497 & 2108.1 & 32037 \\
290 & 48.2 & 3268 & 410 & 73.2 & 9575 & 498 & 3162.3 & 34602 \\
300 & 46.3 & 3741 & 420 & 81.6 & 10349 & 499 & 6234.1 & 38995 \\
310 & 45.5 & 4200 & 430 & 91.4 & 11214 & 500 & $\infty$ & $\infty$ \\
& & & & & & & & \\
\hline
\end{tabular}
\end{center}

\noindent\textbf{23.} But it is obviously convenient to revert these Tables so as to have for the common arguments a series of uniformly increasing values of $\Pi$ or $\Psi$; viz., we obtain by interpolation the values of $h$ and $k$ belonging to the given values of $\Pi$ or $\Psi$, and thus obtain the following Table. Here, in any line of the Table the values of $h$, $k$ are such that $\Pi(h) - \Psi(k) = 0$, viz., the values in question belong to successive points of the $B$-geodesic. And to obtain the values for any other geodesic line $\Pi(h) - \Psi(k) = \pm 500m$, we have only to take each value of $k$ from the line $m$ lines above or below the line from which $h$ is taken; and similarly the table gives at once the values belonging to a geodesic line $\Pi(h) + \Psi(k) = 500m$.

\paragraph*{Table III.}

\begin{center}
\begin{tabular}{|r|r|r|r|r||r|r|r|r|r|}
\hline
$\Pi = \Psi$ & $-h$ & D. & $-k$ & D. & $\Pi = \Psi$ & $-h$ & D. & $-k$ & D. \\\hline
0 & 1000 & & 250 & & 13000 & 650.1 & 20.6 & 446.7 & 7.8 \\
500 & 998.8 & 1.2 & 251.4 & 1.4 & 14000 & 629.5 & 18.3 & 454.5 & 6.5 \\
1000 & 995.4 & 3.4 & 254.5 & 3.1 & 15000 & 611.2 & 15.7 & 461.0 & 5.3 \\
1500 & 989.9 & 5.5 & 259.7 & 5.2 & 16000 & 595.5 & 13.6 & 466.3 & 4.9 \\
2000 & 982.1 & 7.8 & 267.0 & 7.3 & 17000 & 581.9 & 11.8 & 471.2 & 3.8 \\
2500 & 972.6 & 9.5 & 275.2 & 8.2 & 18000 & 570.1 & 10.0 & 475.0 & 3.8 \\
3000 & 961.1 & 11.5 & 284.6 & 9.4 & 19000 & 560.1 & 8.5 & 478.8 & 2.6 \\
3500 & 948.3 & 12.8 & 294.9 & 10.3 & 20000 & 551.6 & 7.3 & 481.4 & 2.2 \\
4000 & 933.8 & 14.5 & 305.6 & 10.7 & 21000 & 544.3 & 6.2 & 483.6 & 2.0 \\
4500 & 918.2 & 15.6 & 316.6 & 11.0 & 22000 & 538.1 & 4.9 & 485.6 & 2.2 \\
5000 & 901.7 & 16.5 & 327.5 & 10.9 & 23000 & 533.2 & 4.6 & 487.8 & 2.1 \\
5500 & 884.5 & 17.2 & 338.3 & 10.4 & 24000 & 528.6 & 8.1 & 489.9 & 2.1 \\
6000 & 866.8 & 17.7 & 348.7 & 10.1 & 26000 & 520.5 & 4.5 & 492.0 & 2.1 \\
6500 & 848.9 & 17.9 & 358.8 & 9.5 & 28000 & 516.0 & 4.2 & 494.1 & 1.7 \\
7000 & 830.8 & 18.1 & 368.3 & 9.2 & 30000 & 511.8 & 3.0 & 495.8 & 1.2 \\
7500 & 812.9 & 17.9 & 377.5 & 8.6 & 32000 & 508.8 & 2.1 & 497.0 & 0.8 \\
8000 & 795.3 & 17.6 & 386.1 & 8.0 & 34000 & 506.7 & 2.0 & 497.8 & 0.5 \\
8500 & 778.0 & 16.8 & 394.1 & 7.7 & 36000 & 504.7 & 1.2 & 498.3 & 0.5 \\
9000 & 761.2 & 16.3 & 401.8 & 7.1 & 38000 & 503.5 & 1.0 & 498.8 & \\
9500 & 744.9 & 15.6 & 408.9 & 6.6 & 39000 & & & 499.0 & \\
10000 & 729.3 & 14.9 & 415.5 & 6.2 & 40000 & 502.5 & 0.6 & & \\
10500 & 714.4 & 14.3 & 421.7 & 5.9 & 42000 & 501.9 & 0.6 & & \\
11000 & 700.1 & 13.5 & 427.6 & 5.3 & 44000 & 501.4 & 0.4 & & \\
11500 & 686.6 & 12.8 & 432.9 & 5.0 & 45800 & 501.0 & & & \\
12000 & 673.8 & 23.7 & 437.9 & 8.8 & $\infty$ & 500 & & 500 & \\
\hline
\end{tabular}
\end{center}

\subsubsection*{Graphical Construction: Projection on the Umbilicar Plane.}

\noindent\textbf{24.} The most convenient mode of delineation of the geodesic lines is obtained by projecting them orthogonally on the umbilicar plane: the contour of the figure is here the umbilicar section, or ellipse $\dfrac{x^{2}}{a} + \dfrac{z^{2}}{c} = 1$; and the curves of curvature of each series are projected into elliptic arcs lying within the ellipse in question, the one set cutting at right angles the axes $AA'$, the other cutting at right angles the axes $CC'$; the equations of the complete ellipses being
\[
\frac{a-b}{a(a+h)} x^{2} + \frac{c-b}{c(c+h)} z^{2} = 1 \quad \text{and} \quad \frac{a-b}{a(a+k)} x^{2} + \frac{c-b}{c(c+k)} z^{2} = 1.
\]

\noindent\textbf{25.} I constructed, by means of the table, a drawing of this kind for the ellipsoid $a, b, c = 1000, 500, 250$, the lengths $\sqrt{a}$ and $\sqrt{c}$ being taken to be 12 inches and 6 inches respectively: the process consists in taking from the table for a series of values $\Pi = \Psi$ (say $\Pi = \Psi = 1000$, $= 2000$ \&c.), the values of $h$ and $k$, laying down for such values the elliptic arcs which represent the two curves of curvature respectively, thus dividing the bounding ellipse into a series of curvilinear rectangles, and then obtaining the geodesic lines by drawing the diagonals of these rectangles, and of course rounding off the corners so as to form continuous curves. The Plate shows on a reduced scale so much of the drawing as is comprised within a quadrant of the bounding ellipse (viz.\ it is a representation of an octant of the ellipsoid).

\subsubsection*{Elliptic-Function Formulae.}

\noindent\textbf{26.} I have in all that precedes abstained from the use of elliptic functions, since obviously the form $\sqrt{1 - k^{2}\sin^{2}\phi}$ of the radical of an elliptic function is in nowise specially appropriate to the present question. But (more particularly in the above-mentioned case $ac - b^{2} = 0$, where the radical is $\sqrt{h(a+h)(c+h)}$ without any exterior factor $b+h$ in the denominator) the formulae are expressible easily and elegantly by elliptic functions, and it is desirable to make the transformation. Reverting to the formulae which, in the case in question (viz.\ when $ac - b^{2} = 0$), give the values of $\Pi(h)$ and $\Psi(k)$; and writing therein $h = -a + (a-c)\sin^{2}\phi$, $k = -a + (a-c)\sin^{2}\psi$, also
\[
\frac{b}{a} = 1 - k^{2}, \quad \frac{b}{c} = \frac{1}{k'^{2}},
\]
we have
\[
\Pi(h) = \int \frac{d\phi}{\sqrt{a} \sqrt{1 - k^{2}\sin^{2}\phi}} \cdot \frac{1}{1 - (1-k^{2})\sin^{2}\phi},
\]
\[
\Psi(k) = \int \frac{d\psi}{\sqrt{a} \sqrt{1 - k^{2}\sin^{2}\psi}} \cdot \frac{1}{1 - (1-k^{2})\sin^{2}\psi}.
\]

\noindent\textbf{27.} Hence
\[
\Pi(h) = \frac{1}{\sqrt{a}} \left[ F(k, \phi) + \frac{1}{1-k^{2}} \log \frac{(1-k^{2})\sin\phi \cos\phi + \sqrt{1-k^{2}\sin^{2}\phi}}{\sqrt{1-(1-k^{2})\sin^{2}\phi}} \right],
\]
(observe, as $h$ passes from $-a$ to $-b$, $\phi$ passes from $\sin^{2}\phi = 0$ to $\sin^{2}\phi = \tfrac{a-b}{a-c} = 1-k^{2}$, and $H$ from $H_{0} = 0$ to $H = 1$).

Similarly
\[
\Psi(k) = \frac{1}{\sqrt{a}} \left[ F(k, \psi) + \frac{1}{k^{2}-1} \log \frac{(1+k^{2})\sin\psi \cos\psi + \sqrt{1-k^{2}\sin^{2}\psi}}{\sqrt{1-(1+k^{2})\sin^{2}\psi}} \right],
\]
and as $k$ passes from $-c$ to $-b$, $\psi$ passes from $\sin^{2}\psi = \tfrac{a-c}{a-c} = 1$ to $\sin^{2}\psi = \tfrac{a-b}{a-c} = 1-k^{2}$, and $K$ from $K_{0}$ to $K = 1$.

\noindent\textbf{28.} The before-mentioned identical equation
\[
\int \frac{dh}{b+h} \sqrt{\frac{(a+h)(c+h)}{-h}} = \int \frac{dh}{h\sqrt{(a+h)(c+h)}}
\]
is by the same transformation converted into
\[
\int_{0}^{\phi} \frac{d\phi}{1-(1-k^{2})\sin^{2}\phi} = \int_{0}^{\phi} \frac{1-(1-k'^{2})}{1-(1+k'^{2})\sin^{2}\phi} \frac{d\phi}{\sqrt{1-k^{2}\sin^{2}\phi}}.
\]
To prove this, I remark that the equation is
\[
\Pi_{1}(-1+k^{2}) = -F_{0} + \Pi_{1}(-1-k'^{2}),
\]
viz.\ this is $\Pi_{1}(-1+k^{2}) = \Pi_{1}(-1-k'^{2}) - F_{0}$, where $\Pi_{1}(-1-k'^{2})$ denotes the principal value of the integral $\displaystyle\int_{0}^{\phi'} \frac{d\phi}{1-(1+k'^{2})\sin^{2}\phi}$. Now (Leg.\ Fonct.\ Ellip., t.~i., p.~71), we have
\[
\Pi_{1}(-m) + \Pi_{1}\left(-\frac{k^{2}}{m}\right) = F_{0} + \frac{1}{\sqrt{m(m+k^{2})}} \log \frac{\sqrt{m(m+k^{2})} + \sqrt{m}\sqrt{m+k^{2}\sin^{2}\theta}}{\sqrt{m(m+k^{2})} - \sqrt{m}\sqrt{m+k^{2}\sin^{2}\theta}},
\]
where, upon examination, it will appear that $\Pi_{1}(-m)$ represents the principal value of the integral. Writing herein $\sin^{2}\theta = \dfrac{1}{1+k'^{2}}$, and therefore $\cos^{2}\theta = \dfrac{k'^{2}}{1+k'^{2}}$, or $\tan^{2}\theta = \dfrac{1}{k'^{2}}$, this is
\[
\Pi_{1}(-1+k^{2}) + \Pi_{1}(-1-k'^{2}) = F_{0},
\]
and the formula $(p')$, p.~141, attributing therein to $\theta$ the foregoing value, becomes
\[
\Pi_{1}(-1+k^{2}) = \Theta_{0} + \frac{1-k^{2}}{2} F_{0} E(\theta) - E_{1} F(\theta).
\]
But $\theta$ is the value for the bisection of the function $F_{0}$, viz., we have $2F_{1}(\theta) = F_{0}$, $2E(\theta) = E_{0} + 1-k^{2}$, whence $F_{0}E(\theta) - E_{0}F(\theta) = \tfrac{1}{2}(1-k^{2})F_{0}$; or the formula in question gives $\Pi_{1}(-1+k^{2}) = \tfrac{1-k^{2}}{2}F_{0}$, whence $\Pi_{1}(-1-k'^{2}) = \tfrac{1+k^{2}}{2}F_{0}$, the result which was to be proved.

\noindent\textbf{29.} The value of $M'$ (observing that $\sqrt{(a-b)(b-c)} = \sqrt{a/c} \cdot b = \dfrac{b}{\sqrt{ac}} = \dfrac{b}{\sqrt{b^{2}}} = 1$) is
\[
M' = \sqrt{(a-b)(b-c)} \int_{-a}^{-b} \frac{dh}{b+h} \sqrt{\frac{(a+h)(c+h)}{-h}} = \sqrt{(a-b)(b-c)} \cdot \Pi(-b),
\]
that is we have
\[
M' = (1-k^{2}) F_{0}(k),
\]
or, what is the same thing,
\[
\log \tan \tfrac{1}{2}\phi_{0} = \frac{1-k^{2}}{2} F_{0}(k),
\]
that is
\[
\tan \tfrac{1}{2}\phi_{0} = \exp\left[ \frac{1-k^{2}}{2} F_{0}(k) \right],
\]
($\phi_{0}$ the South azimuth of the $B$-geodesic at the umbilicus).

\noindent\textbf{30.} I purposely calculated the Table by quadratures as being a method available where the equation $ac - b^{2} = 0$ is not satisfied; but in the present case, where this equation is satisfied, the table might have been calculated from Legendre's Tables of Elliptic Integrals. Observe that $a = 1000$, $b = 500$, $c = 250$, gives $k^{2} = \tfrac{3}{4}$ or angle of modulus = $60^{\circ}$. As an instance of the comparison, suppose $h = -800$, then $\sin^{2}\phi = \tfrac{a+h}{a-c} = \tfrac{200}{750} = \tfrac{4}{15}$, $\log \sin\phi = 9.71298$, $\phi = 31^{\circ}5'$.
\begin{align*}
\sqrt{(a-b)(b-c)} &= \sqrt{500 \cdot 250} = 50\sqrt{5}, \\
\frac{b}{a} &= \frac{500}{1000} = \frac{1}{2}, \quad \frac{1+k^{2}}{1-k^{2}} = \frac{1.75}{0.25} = 7, \\
\Pi(-800) &= \frac{1}{50\sqrt{5}} F(31^{\circ}5') + \frac{1}{50\sqrt{5}} \log 6.7582, \\
2F(31^{\circ}) &= 0.56166, \quad 2F(31^{\circ}5') = 0.56329, \\
\log 6.7582 &= 1.91075, \quad 2.47404 \times 0.03163 &= 0.0782043,
\end{align*}
or multiplying by 100,000 (factor introduced into my Table) this is = 7820.43. The value $\Pi(-800) = 7864$ given by my Table agrees sufficiently well with this, the correct value.

\noindent\textbf{31.} I calculate also the angle $\phi_{0}$, viz.\ we have
\[
\log \tan \tfrac{1}{2}\phi_{0} = \frac{1-k^{2}}{2} F_{0}K_{0} = \frac{1}{2} F_{1}(60^{\circ}).
\]
Leg.\ vol.~iii.\ Table~viii.\ $= \tfrac{1}{2} \cdot 2.15651 = 1.078255$, whence by Leg.\ Table~iv.
\[
\tfrac{1}{2}\phi_{0} = 45^{\circ} + \tfrac{1}{2} \cdot 29^{\circ}29'.64 = 59^{\circ}44'.82 \quad \text{or} \quad \phi_{0} = 119^{\circ}29'.64.
\]
This exceeds $90^{\circ}$, and since at the umbilicus the tangent plane is at right angles to the plane of projection, the $B$-geodesic should in the drawing proceed (as it in fact does) from $U$ in the sense $UC$, touching the bounding ellipse at the point $U$.

$^{*}$ In the present calculation, $\log$ denotes an ordinary logarithm, the hyperbolic logarithm being distinguished as h.~l.

\newpage

\section*{Paper 479}
\addcontentsline{toc}{section}{Paper 479 --- The Second Part of a Memoir on the Development of the Disturbing Function in the Lunar and Planetary Theories}

\subsection*{The Second Part of a Memoir on the Development of the Disturbing Function in the Lunar and Planetary Theories}

\begin{center}
[From the \textit{Memoirs of the Royal Astronomical Society}, vol.~xxxix. (1872), pp.~55---74.\\
Read January 12, 1872.]
\end{center}

The present communication is a sequel to my paper, ``The First Part of a Memoir on the Development of the Disturbing Function in the Lunar and Planetary Theories,'' Memoirs R.A.S., vol.~xxviii. (1859), pp.~187---215, [214], and I have therefore entitled it as above, but it, in fact, relates only to the Planetary Theory. In the First Part, I gave in effect, but not explicitly, an expression for the general coefficient $D(j, j')$ in terms of the coefficients of the multiple cosines of $\theta$ in the expansions of the several powers $(r^{2} + r'^{2} - 2rr'\cos\theta)^{-\frac{1}{2}-i}$, or say $(a^{2} + a'^{2} - 2aa'\cos\theta)^{-\frac{1}{2}-i}$; viz., at the foot of page 208 I speak of the term involving $\cos(jU + j'U')$ as having a certain given value; the term in question is $D(j, j') \cos(jU + j'U')$; and consequently the expression for $D(j, j')$ is
\[
D(j, j') = \Sigma \frac{(-)^{\frac{1}{2}(j+j')-x}}{2^{j+j'}} \Pi_{\frac{1}{2}(j+j')+x} \Pi_{\frac{1}{2}(j+j')-x} \cdot \eta^{\frac{1}{2}(j+j')+x} R_{x}^{\frac{1}{2}(j+j')+x};
\]
the omission was, however, a material one, inasmuch as this expression for the general coefficient serves to connect my formulae with Leverrier's development, \textit{Annales de l'Observ.\ de Paris}, t.~i. (1855), pp.~275---330 and 358---383, and I resume the question for the purpose of applying it.

\subsubsection*{Formula for the general Coefficient $D(j, j')$.}

In the First Part, the reciprocal of the distance of the two planets, or function
\[
\{r^{2} + r'^{2} - 2rr'(\cos U \cos U' + \sin U \sin U' \cos \Phi)\}^{-\frac{1}{2}}
\]
is taken to be developed in multiple cosines of $U$, $U'$, the general term being
\[
D(j, j') \cos(jU + j'U').
\]
where $j$, $j'$ have each of them any integer value from $-\infty$ to $+\infty$ (zero not excluded), but so that $j, j'$ are simultaneously even or simultaneously odd. We have $D(-j, -j') = D(j, j')$ and $D(j', j) = D(j, j')$; and it hence appears that the really distinct values of the coefficient may be taken to be those for which $j$ is not negative, and as regards absolute magnitude is not less than $j'$; and for such values of $j$, $j'$ we have the above-mentioned expression
\[
D(j, j') = \Sigma \frac{(-)^{\frac{1}{2}(j+j')-x}}{2^{j+j'}} \Pi_{\frac{1}{2}(j+j')+x} \Pi_{\frac{1}{2}(j+j')-x} \cdot \eta^{\frac{1}{2}(j+j')+x} R_{x}^{\frac{1}{2}(j+j')+x},
\]
which I proceed to explain and develope.

$\Pi_{\frac{1}{2}(j-j')}$ and $\Pi_{x}$ ($x$ being a positive integer) denote respectively $\tfrac{1}{2}(j-j') \cdot \tfrac{1}{2}(j-j'-2) \ldots (x - \tfrac{1}{2}(j-j'))$, and $1 \cdot 2 \cdot 3 \ldots x$; in particular for $x = 0$, the value of each factorial is = 1.

$\eta$ denotes $\sin \tfrac{1}{2}\Phi$.

The coefficients $R_{x}^{\frac{1}{2}(j+j')+x}$ are those of the multiple cosines in certain developments, viz.\ we have
\[
r^{\frac{1}{2}(j+j')+x} r'^{\frac{1}{2}(j+j')+x} \{r^{2} + r'^{2} - 2rr' \cos(U - U')\}^{-\frac{1}{2}(j+j')-x} = \Sigma R_{i}^{x} \cos i(U - U'),
\]
where, as usual, $i$ extends from $-\infty$ to $\infty$ and $R_{-i}^{x} = R_{i}^{x}$. Writing with Leverrier
\begin{align*}
(a^{2} + a'^{2} - 2aa' \cos H)^{-\frac{1}{2}} &= \Sigma A^{i} \cos iH, \\
aa'(a^{2} + a'^{2} - 2aa' \cos H)^{-\frac{3}{2}} &= \Sigma B^{i} \cos iH, \\
a^{2}a'^{2}(a^{2} + a'^{2} - 2aa' \cos H)^{-\frac{5}{2}} &= \Sigma C^{i} \cos iH, \\
a^{3}a'^{3}(a^{2} + a'^{2} - 2aa' \cos H)^{-\frac{7}{2}} &= \Sigma D^{i} \cos iH,
\end{align*}
then $2R_{0}^{x}$, $2R_{1}^{x}$, $2R_{2}^{x}$, $2R_{3}^{x}$ are the same functions of $r$, $r'$ that $A^{i}$, $B^{i}$, $C^{i}$, $D^{i}$ respectively are of $a$, $a'$.

The expression of $R_{x}^{\frac{1}{2}(j+j')+x}$ is
\begin{multline*}
R_{x}^{\frac{1}{2}(j+j')+x} = (-)^{\frac{1}{2}(j+j')-x} \frac{\Pi_{\frac{1}{2}(j+j')+x}}{2^{j+j'} \Pi_{\frac{1}{2}(j-j')-x} \Pi_{x-j'} \Pi_{x+j'} \Pi_{\frac{1}{2}(j+j')-x}} R_{x}^{\frac{1}{2}(j+j')+x},
\end{multline*}
and, finally, in the expression for $D(j, j')$, $x$ has every integer value from 0 to $\infty$, and, for any given value of $x$, $\xi$ extends by steps of two units from the inferior value $-(x-j')$ to the superior value $x-j$.

It is convenient to write $x = \tfrac{1}{2}(j+j') + s$; we have then $\xi$ extending from $-\tfrac{1}{2}(j-j')-s$ to $+\tfrac{1}{2}(j-j')+s$, or writing $\xi = -\tfrac{1}{2}(j-j') + \sigma$, $\sigma$ has the $s+1$ values $s$, $s-2$, $s-4$, $\ldots -s$, viz.\ for $s = 2p+1$ the values are $\pm 1$, $\pm 3$, $\ldots \pm (2p+1)$, and for $s = 2p$ they are $0$, $\pm 2$, $\pm 4$, $\ldots \pm 2p$.

Making these changes we have
\begin{multline*}
D(j, j') = \Sigma (-)^{s} \eta^{\frac{1}{2}(j+j')+s} \cdot \frac{\Pi\{\tfrac{1}{2}(j+j')+s\}}{\Pi\{\tfrac{1}{2}(j-j')+\sigma\} \Pi\{\tfrac{1}{2}(j+j')-\sigma\}} \\
\cdot \frac{\Pi\{\tfrac{1}{2}(j+j')+s\}}{\Pi\{\tfrac{1}{2}(j-j')-\sigma\} \Pi\{\tfrac{1}{2}(j+j')+\sigma\}} R_{\sigma}^{\frac{1}{2}(j+j')+s},
\end{multline*}
where $\sigma = -\tfrac{1}{2}(j-j')+s$, $-\tfrac{1}{2}(j-j')+s-2$, $\ldots$, $-\tfrac{1}{2}(j-j')-s$; viz.\ this is $(-)^{s}$ into the product of two binomial coefficients, each belonging to the exponent $\tfrac{1}{2}(j+j') + s$.

\subsubsection*{Particular Cases, $j + j' = 0$, 2, 4, 6, being those required in the Planetary Theory.}

Considering successively the cases $j + j' = 0$, 2, 4, 6, we have, first,
\[
D(j, -j) = \Sigma \frac{\Pi_{s}}{\Pi_{\frac{1}{2}(s-\xi)} \Pi_{\frac{1}{2}(s+\xi)}} \eta^{s} R_{\xi}^{s},
\]
which, developed as far as $\eta^{4}$, is
\[\tag{*$j+j'=0$}
D(j, -j) = R_{j}^{0} \left\{ 1 - \frac{j^{2}}{2.4} \eta^{2} R_{j}^{2} + \frac{j^{2}(j^{2}-1)}{2.4.6.8} \eta^{4} R_{j}^{4} \right\},
\]
where, and in what immediately follows, $A$, $B$, $C$, $D$ are used to denote functions (not of $(a, a')$, but) of $r$, $r'$.

Secondly,
\[
D(j, -j+2) = \Sigma \frac{\Pi_{s+1}}{\Pi_{\frac{1}{2}(s+1-\xi)} \Pi_{\frac{1}{2}(s+1+\xi)}} \eta^{s+1} R_{\xi}^{s+1},
\]
which, developed to $\eta^{4}$, is
\[\tag{*$j+j'=2$}
D(j, -j+2) = R_{j}^{1} \left\{ \frac{j(j-1)}{2.4} \eta^{2} + \frac{j(j-1)(j^{2}-j-2)}{2.4.6.8} \eta^{4} R_{j}^{3} \right\}.
\]

Thirdly,
\[
D(j, -j+4) = \Sigma \frac{\Pi_{s+2}}{\Pi_{\frac{1}{2}(s+2-\xi)} \Pi_{\frac{1}{2}(s+2+\xi)}} \eta^{s+2} R_{\xi}^{s+2},
\]
which, developed to $\eta^{4}$, is
\[\tag{*$j+j'=4$}
D(j, -j+4) = R_{j}^{2} \left\{ \frac{j(j-1)(j-2)(j-3)}{2.4.6.8} \eta^{4} R_{j}^{2} \right\}.
\]

And, fourthly,
\[
D(j, -j+6) = \Sigma \frac{\Pi_{s+3}}{\Pi_{\frac{1}{2}(s+3-\xi)} \Pi_{\frac{1}{2}(s+3+\xi)}} \eta^{s+3} R_{\xi}^{s+3},
\]
which, developed to $\eta^{4}$, is simply
\[\tag{*$j+j'=6$}
D(j, -j+6) = R_{j}^{3} \cdot \frac{j(j-1)(j-2)(j-3)(j-4)(j-5)}{2.4.6.8.10.12} \eta^{6} R_{j}^{3}.
\]

The foregoing formulae, although obtained on the supposition $j = 0$, or positive, apply without alteration to the case $j =$ negative, and the entire series of terms of an order not exceeding 6 as regards $\eta$ may be written,
\begin{align*}
& D(j, -j) \cos(jU - jU') \\
+\ & 2D(j, -j+2) \cos\{jU + (-j+2)U'\} \\
+\ & 2D(j, -j+4) \cos\{jU + (-j+4)U'\} \\
+\ & 2D(j, -j+6) \cos\{jU + (-j+6)U'\},
\end{align*}
where $j$ has every integer value from $-\infty$ to $+\infty$.

\subsubsection*{Comparison with Leverrier.}

This is in fact what Leverrier's expression becomes on putting therein $e = e' = 0$.

To verify this, observe that Leverrier having defined his $A^{i}$, $B^{i}$, $C^{i}$, $D^{i}$, as above, writes further
\begin{align*}
E^{i} &= \tfrac{1}{2}(E^{i-1} + E^{i+1}), \\
G^{i} &= \tfrac{1}{2}(G^{i-1} + 4G^{i} + G^{i+1}), \\
H^{i} &= \tfrac{1}{2}(H^{i-2} + 4H^{i-1} + 6H^{i} + 4H^{i+1} + H^{i+2}), \\
L^{i} &= \tfrac{1}{2}(L^{i-2} + L^{i}), \\
S^{i} &= \tfrac{1}{2}(S^{i-2} + 4S^{i-1} + S^{i}),
\end{align*}
where $E^{i}$, $G^{i}$, $H^{i}$, $L^{i}$, $S^{i}$ are the coefficients in the developments of the disturbing function in the [text continues in original].
\clearpage


% =========================================================
% Chunk: cayley_vol07_pages_551_600.tex  (orig pages 551--600)
% =========================================================
\begin{center}
\Large\bfseries The Collected Mathematical Papers of Arthur Cayley

\vspace{0.5em}
\large Volume VII

\vspace{1em}
\normalsize Pages 515--564 (PDF pages 551--600)
\end{center}

\vspace{2em}

%==============================================================================
% PAPER [479]
%==============================================================================
\section*{[479 --- On the Development of the Disturbing Function in the Lunar and Planetary Theories (continued)]}
\addcontentsline{toc}{section}{[479] On the Development of the Disturbing Function}

\markboth{Cayley, Collected Papers, Vol.\ VII}{[479] Disturbing Function}

\begin{center}
\textbf{ON THE DEVELOPMENT OF THE DISTURBING FUNCTION IN THE LUNAR AND PLANETARY THEORIES.}
\end{center}

(continued from previous pages)

\bigskip

(consequently $E^{-i}=E^{i}$, $G^{-i}=G^{i}$, $H^{-i}=H^{i}$, $L^{-i+2}=L^{i}$, $S^{-i+2}=S^{i}$, $T^{-i+4}=T^{i}$), and that the terms in question, putting in the coefficients $e=e'=0$, are with him
\begin{align*}
&\{(1)^{i}+(11)^{i}\eta^{2}+(17)^{i}\eta^{4}+(20)^{i}\eta^{6}\}\cos(il'-i\lambda),\\
&\qquad\{(212)^{i}\eta^{2}+(218)^{i}\eta^{4}+(221)^{i}\eta^{6}\}\cos[il'-(i-2)\lambda-2\tau'],\\
&\qquad\qquad\{(372)^{i}\eta^{4}+(375)^{i}\eta^{6}\}\cos[il'-(i-4)\lambda-4\tau'],\\
&\qquad\qquad\qquad\{(449)^{i}\eta^{6}\}\cos[il'-(i-6)\lambda-6\tau'],
\end{align*}
where, substituting for $(1)^{i}$, $(11)^{i}$, \&c., their values, the coefficients are
\begin{align*}
&\tfrac12 A^{i} - \eta^{2}\cdot\tfrac12 E^{i} + \eta^{4}\cdot\tfrac12 G^{i} - \eta^{6}\cdot\tfrac12 H^{i},\\
&= \tfrac12 A^{i} - \eta^{2}\cdot\tfrac14(B^{i-1}+B^{i+1})\\
&\qquad + \eta^{4}\cdot\tfrac{3}{16}(C^{i-2}+4C^{i}+C^{i+2}) - \eta^{6}\cdot\tfrac{5}{32}(D^{i-3}+9D^{i-1}+9D^{i+1}+D^{i+3});\\
&\eta^{2}\cdot\tfrac12 B^{i-1} - \eta^{4}\cdot L^{i} + \eta^{6}\cdot S^{i}\\
&\qquad = \eta^{2}\cdot\tfrac12 B^{i+1} - \eta^{4}(\tfrac34 C^{i-2}+C^{i}) + \eta^{6}\cdot\tfrac{15}{16}(D^{i-3}+3D^{i-1}+D^{i+1});\\
&\eta^{4}\cdot\tfrac38 C^{i-2} - \eta^{6}\cdot T^{i} = \eta^{4}\cdot\tfrac38 C^{i+2} - \eta^{6}\cdot\tfrac{5}{16}(D^{i-3}+D^{i-1});
\end{align*}
and
\[
\eta^{6}\cdot\tfrac{5}{16}D^{i-3}.
\]

Writing herein $j$ in place of $i$, and for $A^{j}$, $B^{j-1}$, \&c., the equal values $A^{-j}$, $B^{-j+1}$, \&c., we have precisely the foregoing coefficients $D(j,-j)$, \dots\ $D(j,-j+6)$.

\bigskip

\begin{center}
\textit{The Development in Powers of $e$, $e'$.}
\end{center}

The complete expression of the reciprocal of the distance is obtained from
\begin{multline*}
D(j,-j)\cos(jU-jU')\\
+2D(j,-j+2)\cos(jU+(-j+2)U')\\
+2D(j,-j+4)\cos(jU+(-j+4)U')\\
+2D(j,-j+6)\cos(jU+(-j+6)U'),
\end{multline*}
by writing therein for $r$, $r'$, $U$, $U'$, instead of the circular, the elliptic values, that is the values
\begin{align*}
r &= a\,\text{elqr}(e,L-\Pi) &&= a(1+x),\\
r' &= a'\,\text{elqr}(e',L'-\Pi') &&= a'(1+x'),\\
U &= \Pi-\Theta+\text{elta}(e,L-\Pi) &&= L-\Theta+f,\\
U' &= \Pi'-\Theta'+\text{elta}(e',L'-\Pi') &&= L'-\Theta'+f';
\end{align*}
$L$, $\Pi$, $\Theta$ the mean longitude in orbit, longitude of perihelion in orbit, and longitude of node; and the like for $L'$, $\Pi'$, $\Theta'$; ``elqr'' = elliptic quotient radius, ``elta'' = elliptic true anomaly; or, what is the same thing, if we write $\text{elta}(e,L-\Pi)=L-\Pi+\text{eltt}(e,L-\Pi)$, and the like for $\text{elta}(e',L'-\Pi')$, then
\begin{align*}
U &= L-\Theta+\text{eltt}(e,L-\Pi) = L-\Theta+y,\\
U' &= L'-\Theta'+\text{eltt}(e',L'-\Pi') = L'-\Theta'+y'.
\end{align*}

\newpage

The process for doing this is explained, First Part, pp.~205--207, [214], viz., writing $r=a(1+x)$, $r'=a'(1+x')$, and restoring $j'$ (instead of its value $-j$, \dots\ $-j+6$, as the case may be), we have a general term
\[
D(j,j')\frac{r^{j}}{a^{j}}\frac{r'^{j'}}{a'^{j'}}x^{i}x'^{i'}\cos(jU+j'U'),
\]
where $D(j,j')$ now denotes the value obtained by writing $a$, $a'$ in place of $r$, $r'$ and $f$, $f'$ are the true anomalies $\text{elta}(e,L-\Pi)$ and $\text{elta}(e',L'-\Pi')$. And the second factor, $x^{i}x'^{i'}$ into the cosine, is given as a series
\[
\Sigma([\cos]^{i}+[\sin]^{i})([\cos]^{i'}+[\sin]^{i'})\cos[i(L-\Pi)+i'(L'-\Pi')+j(\Theta-\theta)-j'(\Theta'-\theta')],
\]
where $[\cos]^{i}$, $[\sin]^{i}$ are functions of $e$, $[\cos]^{i'}$, $[\sin]^{i'}$ functions of $e'$. Or, what is better, the term $x^{i}x'^{i'}$ into the cosine may be written $a^{i}x^{i}\,a'^{i'}x'^{i'}\cos[j(L-\Theta+y)+j'(L'-\Theta'+y')]$, and the expansion then is
\begin{multline*}
\Sigma([\cos]^{i}+[\sin]^{i})([\cos]^{i'}+[\sin]^{i'})\\
\times\cos[i(L-\Pi)+i'(L'-\Pi')+j(L-\Theta)+j'(L'-\Theta')],
\end{multline*}
where as before $[\cos]^{i}$, $[\sin]^{i}$ are the same functions of $e$, $[\cos]^{i'}$, $[\sin]^{i'}$ functions of $e'$, viz.\ the $e$-functions are those given in the two ``datum-tables'' $(\alpha'\dots\alpha'_{i})\cos jy$ and $(\alpha'\dots\alpha'_{i})\sin jy$, taken from Leverrier, which I have given in my ``Tables of the Developments of Functions in the Theory of Elliptic Motion,'' Memoirs R.A.S.\ vol.~xxix.\ (1861), pp.~191--306, [216]. In order to better show which are the symbols referred to, we may, instead of $[\cos]^{i}$, \&c., write $[\alpha'\cos jy]^{i}$, \&c., the formula will then be
\begin{multline*}
a^{i}x^{i}\,a'^{i'}x'^{i'}\cos[j(L-\Theta+y)+j'(L'-\Theta'+y')] =\\
\Sigma\{[\alpha'\cos jy]^{i}+[\alpha'\sin jy]^{i}\}\{[\alpha''\cos j'y']^{i'}+[\alpha''\sin j'y']^{i'}\}\\
\times\cos[i(L-\Pi)+i'(L'-\Pi')+j(L-\Theta)+j'(L'-\Theta')],
\end{multline*}
and if we attribute to $i$, $i'$ any given values, that is, attend to any particular multiple cosine,
\[
\cos[i(L-\Pi)+i'(L'-\Pi')+j(L-\Theta)+j'(L'-\Theta')],
\]
the coefficient hereof will be
\[
\Sigma\{[\alpha'\cos jy]^{i}+[\alpha'\sin jy]^{i}\}\{[\alpha''\cos j'y']^{i'}+[\alpha''\sin j'y']^{i'}\},
\]
where $\alpha$, $\alpha'$ each extend from zero to infinity, but to obtain the expression up to a given order $p$ in $e$, $e'$, we take only the values up to $\alpha+\alpha'=p$.

\begin{center}
\textit{Particular Case.}
\end{center}

Thus, for instance, in $\cos[j(L-\Theta)-j'(L'-\Theta')]$ the terms independent of $e'$ are
\begin{multline*}
-D(j,-j)\{[\alpha'\cos jy]^{0}+[\alpha'\sin jy]^{0}\}\\
+\frac{d}{de}D(j,-j)\{[\alpha'\cos jy]^{1}+[\alpha'\sin jy]^{1}\}\\
+\frac{1}{1\cdot2}\frac{d^{2}}{de^{2}}D(j,-j)\{[\alpha'\cos jy]^{2}+[\alpha'\sin jy]^{2}\}\\
+\text{\&c.}
\end{multline*}
which, observing that in the present case the sine terms vanish, is
\[
\begin{array}{r@{}c@{}l}
& 1 & \\[-2pt]
& -8f & +96f^{2} -1280f^{3} \\[-2pt]
& & -54f^{2} +3920f^{3} \\[-2pt]
& & \qquad -3440f^{3} \\[-2pt]
\hline
& 1 & -8f \quad +384f^{2} \quad -46080f^{3}
\end{array}
\qquad D(j,-j)
\]
\[
+\begin{array}{r@{}c@{}l}
& +4 & -48f \quad -360f^{2} \\[-2pt]
& +4 & -96f \quad +1920f^{2} \\[-2pt]
& & \qquad -1320f^{2} \\[-2pt]
\hline
& +8 & -144f \quad +2880f^{2}
\end{array}
\quad\frac{1}{1\cdot2}\frac{d}{de}D(j,-j)
\]
\[
+\begin{array}{r@{}c@{}l}
& +144 & -2880f \\[-2pt]
& +144 & -5760f \\[-2pt]
\hline
& +288 & -8640f
\end{array}
\quad\frac{1}{1\cdot2\cdot3}\frac{d^{2}}{de^{2}}D(j,-j)
\]
\[
+\begin{array}{c}
+14400 \\[2pt]
+14400 \\[-2pt]
\hline
+28800
\end{array}
\quad\frac{1}{1\cdot2\cdot3\cdot4}\frac{d^{3}}{de^{3}}D(j,-j)
\]
\[
+\begin{array}{c}
+144000 \\[2pt]
+144000 \\[-2pt]
\hline
+288000
\end{array}
\quad\frac{1}{1\cdot2\cdot3\cdot4\cdot5}\frac{d^{4}}{de^{4}}D(j,-j)
\]
viz.\ the term in $e^{\alpha}$ is
\[
\frac{1}{1\cdot2\dots\alpha}\frac{d^{\alpha}}{de^{\alpha}}D(j,-j)\cdot[\alpha'\cos jy]^{\alpha},
\]
viz.\ writing $ij=0$, and therefore $D(j,-j)=\tfrac12 A^{-j}$, the term in $e^{\alpha}$ is
\[
\frac{1}{1\cdot2\dots\alpha}\frac{d^{\alpha}}{de^{\alpha}}\tfrac12 A^{-j}\cdot[\alpha'\cos jy]^{\alpha},
\]
which, conformably with Leverrier's subscript notation

\newpage

The term in question is given by Leverrier as
\[
\left(\frac{d}{de}\right)^{\alpha}(2)^{i}, = e^{\alpha}\cdot i(2), \quad A=A^{j} \text{ and } K^{j}=A^{j},
\]
\[=
\tfrac12\frac{d^{\alpha}}{de^{\alpha}}\{-2A^{j}+A^{j-1}+A^{j+1}\}, \text{ which agrees.}
\]
Similarly the term in $e^{2}$ is
\[
\tfrac12\frac{d^{2}}{de^{2}}\{96A^{j}-54A^{j-1}-48A^{j-1}-96A^{j+1}+144A^{j-2}+144A^{j}\}A^{-j},
\]
\[=
\tfrac12\frac{d^{2}}{de^{2}}\{(96j^{2}-54j)A^{j-1}-48jA^{j-1}-96jA^{j+1}+144A^{j-2}+144A^{j}\},
\]
and the term in question is given by Leverrier as
\[
\left(\frac{d}{de}\right)^{2}(4)^{i}=e^{2}\cdot\frac{d^{2}}{de^{2}}(4), \quad A=A^{j} \text{ and } K^{j}=A^{j},
\]
which agrees.

I have not made the comparison of any more terms.

\begin{center}
\textit{Leverrier's Results expressed in terms of the Arguments, $L'-@'$, $L'-\Pi'$, $L-@$, $L-\Pi$.}
\end{center}

The angles which Leverrier uses in his arguments are $I'$, $\lambda$, $\varpi'$, $\varpi$, and $\tau'$, viz.\ we have,
\[
I' = @'+(L'-@'), \quad \lambda=@'+(L-@), \quad \varpi'=@'+(L'-@'), \quad \varpi=@'+(L-@), \quad \tau'=@',
\]
where $L$, $\Pi$, $@$ are the mean longitude of the planet $m$, its perihelion and the mutual node, all in the orbit of $m$; and similarly $L'$, $\Pi'$, $@'$ are the mean longitude of the planet $m'$, of its perihelion and of the mutual node, all in the orbit of $m'$.

On substituting the foregoing values of $I'$, $\lambda$, \&c., $@'$, as it should do, disappears, and the arguments are all of them linear functions of $L'-@'$, $\Pi'-@'$, $L-@$, $\Pi-@$; or, if we please, of $L'-@'$, $L'-\Pi'$, $L-@$, $L-\Pi$, that is of the distances of each planet from its own perihelion and from the mutual node.

It is, I think, convenient to use these last angular distances, and accordingly in Leverrier's arguments, I write,
\begin{align*}
I'&=@'+(L'-@'),\\
\lambda&=@'+(L-@),\\
\varpi'&=@'+(L'-@')-(L'-\Pi'),\\
\varpi&=@'+(L-@)-(L-\Pi),\\
\tau'&=@',
\end{align*}

and for the purpose of reference form as it were an Index to his result as follows:

\smallskip
\begin{center}
\textbf{Reciprocal of Distance = as follows:}
\end{center}

\textit{Terms of order zero: terms of orders 2, 4, 6, having the same arguments.}

\[
\begin{array}{cccc}
L'-@ & L'-\Pi' & L-@ & L-\Pi \\
(1)^{i} & (\tfrac12 e)^{0}(\tfrac12 e')^{0} & \cos & i-i \\
(21)^{i}(31)^{i} & (\tfrac12 e)^{2}(\tfrac12 e')^{0} & \cos & i+1-i-1 \\
(31)^{i}(17)^{i}(\dots 34)^{i} & (\tfrac12 e)^{2}(\tfrac12 e')^{2} & \cos & i+2-i-2 \\
(35)^{i}(18)^{i}(\tfrac12 e')^{2} & (\tfrac12 e)^{4}(\tfrac12 e')^{0} & \cos & i+3-i-3 \\
(36)^{i}(\dots 39)^{i} & (\tfrac12 e)^{4}(\tfrac12 e')^{2} & \cos & -i+2-2 \\
(40)^{i}(\dots 43)^{i} & (\tfrac12 e)^{0}(\tfrac12 e')^{4} & \cos & i-1-i+2-1 \\
(44)^{i}(\dots 47)^{i} & (\tfrac12 e)^{2}(\tfrac12 e')^{4} & \cos & i-2-i+2 \\
(48)^{i}(\tfrac12 e)^{4}(\tfrac12 e')^{4} & (48)^{i} & \cos & i+1-i+2-3 \\
(49)^{i}(\tfrac12 e)^{6}(\tfrac12 e')^{0} & (49)^{i} & \cos & i-3-i+2+1 \\
\end{array}
\]

\textit{Terms of the first order: terms of orders 3, 5, 7, having the same arguments.}

\[
\begin{array}{cccc}
L'-@ & L'-\Pi' & L-@ & L-\Pi \\
(50)^{i} & \tfrac12 e & (50\dots69) & \cos -i+1 \\
(70)^{i} & \tfrac12 e' & (70\dots89) & \cos +1-i \\
(90)^{i}(\tfrac12 e)(\tfrac12 e') & (90\dots99) & \cos & +1-i-2 \\
(100)^{i}(\tfrac12 e)(\tfrac12 e') & (100\dots109) & \cos & +2-i-1 \\
(110)^{i}(\tfrac12 e)^{3}(\tfrac12 e') & (110\dots113) & \cos & +2-i-3 \\
(114)^{i}(\tfrac12 e)^{3}(\tfrac12 e') & (114\dots117) & \cos & +3-i-2 \\
(118)^{i}(\tfrac12 e)^{3}(\tfrac12 e') & (118) & \cos & +3-i-4 \\
(119)^{i}(\tfrac12 e)^{5}(\tfrac12 e') & (119) & \cos & +4-i-3 \\
(120)^{i}(\tfrac12 e)^{5} & (120\dots129) & \cos & -i+2-1 \\
(130)^{i}(\tfrac12 e)^{1} & (130\dots139) & \cos & -1-i+2 \\
(140)^{i}(\tfrac12 e')^{3} & (140\dots143) & \cos & -i+2-3 \\
(144)^{i}(\tfrac12 e)^{3}(\tfrac12 e')^{3} & (144\dots147) & \cos & +1-i+2-2 \\
(148)^{i}(\tfrac12 e)^{3}(\tfrac12 e')^{3} & (148\dots151) & \cos & -1-i+2-2 \\
(152)^{i}(\tfrac12 e)^{5}(\tfrac12 e') & (152\dots155) & \cos & -2-i+2-1 \\
(156)^{i}\tfrac12 e'(\tfrac12 e')^{5} & (156\dots159) & \cos & i-2-i+2-1 \\
(160)^{i}(\tfrac12 e')^{5} & (160\dots163) & \cos & i-3-i+2+1 \\
(164)^{i}(\tfrac12 e)^{3}(\tfrac12 e')^{5} & (164) & \cos & i-3-i+2+2 \\
(165)^{i}(\tfrac12 e)^{5}(\tfrac12 e')^{3} & (165) & \cos & i+1-i+2-4 \\
(166)^{i}(\tfrac12 e)^{5}(\tfrac12 e')^{3} & (166) & \cos & i+2-i+2-3 \\
(167)^{i}(\tfrac12 e)^{7}(\tfrac12 e') & (167) & \cos & i-4-i+2+2 \\
(168)^{i}(\tfrac12 e')^{7} & (168) & \cos & i+1-i+2-4 \\
(169)^{i}(\tfrac12 e)^{5}(\tfrac12 e')^{5} & (169) & \cos & i-i+2-3 \\
(170)^{i}(\tfrac12 e)(\tfrac12 e')^{7} & (170) & \cos & i-1-i+4-2 \\
(171)^{i}(\tfrac12 e')^{7} & (171) & \cos & i-2-i+4-1 \\
\end{array}
\]

\newpage

\textit{Terms of second order: terms of orders 4, 6, having the same arguments.}

\[
\begin{array}{cccc}
L'-@ & L'-\Pi' & L-@ & L-\Pi \\
(172)^{i} & (\tfrac12 e)^{2} & (172\dots181) & \cos i-i+2 \\
(182)^{i} & (\tfrac12 e)(\tfrac12 e') & (182\dots191) & \cos i+1-i+1 \\
(192)^{i} & (\tfrac12 e')^{2} & (192\dots201) & \cos i+2-i \\
\end{array}
\]

\textit{Terms of third order: terms of orders 5, 7, having the same arguments.}

\[
\begin{array}{cccc}
L'-@ & L'-\Pi' & L-@ & L-\Pi \\
(202)^{i} & (\tfrac12 e)^{3} & (202\dots209) & \cos i-i+3 \\
(210)^{i} & (\tfrac12 e)^{2}(\tfrac12 e') & (210\dots215) & \cos i+1-i+2 \\
(216)^{i} & (\tfrac12 e)(\tfrac12 e')^{2} & (216\dots221) & \cos i+2-i+1 \\
(222)^{i} & (\tfrac12 e')^{3} & (222\dots227) & \cos i+3-i \\
\end{array}
\]

\textit{Terms of fourth order (concluded):}

\[
\begin{array}{cccc}
L'-@ & L'-\Pi' & L-@ & L-\Pi \\
(228)^{i} & (\tfrac12 e)^{4} & (228\dots233) & \cos i-i+4 \\
(234)^{i} & (\tfrac12 e)^{3}(\tfrac12 e') & (234\dots237) & \cos i+1-i+3 \\
(238)^{i} & (\tfrac12 e)^{2}(\tfrac12 e')^{2} & (238\dots241) & \cos i+2-i+2 \\
(242)^{i} & (\tfrac12 e)(\tfrac12 e')^{3} & (242\dots245) & \cos i+3-i+1 \\
(246)^{i} & (\tfrac12 e')^{4} & (246\dots249) & \cos i+4-i \\
\end{array}
\]

\textit{Terms of fifth order (concluded):}

\[
\begin{array}{cccc}
L'-@ & L'-\Pi' & L-@ & L-\Pi \\
(250)^{i} & (\tfrac12 e)^{5} & (250\dots253) & \cos i-i+5 \\
(254)^{i} & (\tfrac12 e)^{4}(\tfrac12 e') & (254\dots257) & \cos i+1-i+4 \\
(258)^{i} & (\tfrac12 e)^{3}(\tfrac12 e')^{2} & (258\dots261) & \cos i+2-i+3 \\
(262)^{i} & (\tfrac12 e)^{2}(\tfrac12 e')^{3} & (262\dots265) & \cos i+3-i+2 \\
(266)^{i} & (\tfrac12 e)(\tfrac12 e')^{4} & (266\dots269) & \cos i+4-i+1 \\
(270)^{i} & (\tfrac12 e')^{5} & (270\dots273) & \cos i+5-i \\
\end{array}
\]

\textit{Terms of sixth order (concluded):}

\[
\begin{array}{cccc}
L'-@ & L'-\Pi' & L-@ & L-\Pi \\
(274)^{i} & (\tfrac12 e)^{6} & (274\dots277) & \cos i-i+6 \\
(278)^{i} & (\tfrac12 e)^{5}(\tfrac12 e') & (278\dots281) & \cos i+1-i+5 \\
(282)^{i} & (\tfrac12 e)^{4}(\tfrac12 e')^{2} & (282\dots285) & \cos i+2-i+4 \\
(286)^{i} & (\tfrac12 e)^{3}(\tfrac12 e')^{3} & (286\dots289) & \cos i+3-i+3 \\
(290)^{i} & (\tfrac12 e)^{2}(\tfrac12 e')^{4} & (290\dots293) & \cos i+4-i+2 \\
(294)^{i} & (\tfrac12 e)(\tfrac12 e')^{5} & (294\dots297) & \cos i+5-i+1 \\
(298)^{i} & (\tfrac12 e')^{6} & (298\dots301) & \cos i+6-i \\
\end{array}
\]

\textit{Terms of seventh order (concluded):}

\[
\begin{array}{cccc}
L'-@ & L'-\Pi' & L-@ & L-\Pi \\
(302)^{i} & (\tfrac12 e)^{7} & (302\dots305) & \cos i-i+7 \\
(306)^{i} & (\tfrac12 e)^{6}(\tfrac12 e') & (306\dots309) & \cos i+1-i+6 \\
(310)^{i} & (\tfrac12 e)^{5}(\tfrac12 e')^{2} & (310\dots313) & \cos i+2-i+5 \\
(314)^{i} & (\tfrac12 e)^{4}(\tfrac12 e')^{3} & (314\dots317) & \cos i+3-i+4 \\
(318)^{i} & (\tfrac12 e)^{3}(\tfrac12 e')^{4} & (318\dots321) & \cos i+4-i+3 \\
(322)^{i} & (\tfrac12 e)^{2}(\tfrac12 e')^{5} & (322\dots325) & \cos i+5-i+2 \\
(326)^{i} & (\tfrac12 e)(\tfrac12 e')^{6} & (326\dots329) & \cos i+6-i+1 \\
(330)^{i} & (\tfrac12 e')^{7} & (330\dots333) & \cos i+7-i \\
\end{array}
\]

The foregoing table gives in a compendious form the whole of the developments of the reciprocal of the distance, as far as the seventh order in $e$, $e'$.

\vspace{2em}

%==============================================================================
% PAPER [480]
%==============================================================================
\newpage
\section*{[480 --- On the Expression of Delaunay's $I$, $g$, $h$, \&c.\ in terms of the Elements and of their Ratio]}
\addcontentsline{toc}{section}{[480] On the Expression of Delaunay's $I$, $g$, $h$, \&c.}

\markboth{Cayley, Collected Papers, Vol.\ VII}{[480] Delaunay's Expressions}

\begin{center}
\textbf{ON THE EXPRESSION OF DELAUNAY'S $I$, $g$, $h$, \&C.\ IN TERMS OF HIS FINALLY ADOPTED CONSTANTS.}
\end{center}

[From the \textit{Monthly Notices of the Royal Astronomical Society}, vol.\ XXXII.\ (1871--72), pp.\ 8--16.]

\bigskip

We have in Delaunay's lunar theory, $I$, the mean anomaly of the Moon, $g$, the mean distance of perigee from ascending node, $h$, the mean longitude of ascending node, quantities which vary directly as the time, the coefficients of $t$, or values of $\dfrac{I}{t}$, $\dfrac{g}{t}$, $\dfrac{h}{t}$, being given in his \textit{Th\'eorie du Mouvement de la Lune}, vol.\ II.\ pp.\ 237, 238.

But these values are not expressed in terms of his constants $a$ (or $n$), $e$, $\gamma$, finally adopted as explained p.~800, and it seems very desirable to obtain the expressions of $I$, $g$, $h$, in terms of these finally adopted constants: I have accordingly effected this transformation (which I found less laborious than I had anticipated).

It will be convenient to imagine the $a$, $n$, $e$, $\gamma$ of pp.\ 237, 238 replaced by $A$, $N$, $E$, $\Gamma$ respectively.

This being so, and writing $m$ for the $\dfrac{n'}{n}$ of p.~800 we have, p.~800,
\begin{multline*}
N = n\{1 + (\tfrac32 m^{2} + \tfrac{27}{16}m^{3} + \tfrac{225}{64}m^{4})\\
+ (-1 + 3e^{2} - \tfrac52 e^{4} - e'^{2} - 2\gamma^{2} + \tfrac43 \gamma^{4} + \tfrac43 \gamma^{2}e^{2} - \tfrac43 \gamma^{2}e'^{2} - \tfrac14 e^{2}e'^{2} - \tfrac{1}{16}e'^{4})m^{2}\\
+ (-\tfrac34 e^{2} - \tfrac{15}{16}e^{4} + \tfrac{27}{16}\gamma^{2} + \tfrac{27}{16}\gamma^{2}e^{2} - \tfrac{27}{16}\gamma^{2}e'^{2}\\
+ \tfrac{27}{16}e^{2}e'^{2} - \tfrac{27}{16}e'^{4})m^{3}\\
+ \tfrac{225}{128}e^{2}m^{4}\},
\end{multline*}
and hence calculating $N$ from the formula $N^{2}a^{3}=n^{2}a^{3}$, we find
\begin{multline*}
N=n\{1 + [\tfrac34 m^{2} + \tfrac{27}{16}m^{3} + \tfrac{225}{64}m^{4}]\\
+(\tfrac32 \gamma^{2} + \tfrac32 e^{2} - \tfrac54 \gamma^{4} - \tfrac{15}{8}\gamma^{2}e^{2} - \tfrac34 \gamma^{2}e'^{2} + \tfrac{15}{16}e^{4} + \tfrac{3}{16}e^{2}e'^{2} + \tfrac{3}{64}e'^{4})m^{2}\\
+(-\tfrac{27}{32}\gamma^{2} - \tfrac{27}{32}e^{2} + \tfrac{81}{64}\gamma^{4} + \tfrac{81}{32}\gamma^{2}e^{2} + \tfrac{81}{64}\gamma^{2}e'^{2}\\
- \tfrac{81}{64}e^{4} - \tfrac{81}{64}e^{2}e'^{2} - \tfrac{81}{256}e'^{4})m^{3}\\
+(-\tfrac{225}{128}\gamma^{2} - \tfrac{225}{128}e^{2} + \tfrac{675}{256}\gamma^{4} + \tfrac{675}{128}\gamma^{2}e^{2} + \tfrac{675}{256}\gamma^{2}e'^{2}\\
- \tfrac{675}{256}e^{4} - \tfrac{675}{256}e^{2}e'^{2} - \tfrac{675}{1024}e'^{4})m^{4}\\
+(-\tfrac{225}{64})m^{4}\}\\
=n(1+Q) \text{ suppose.}
\end{multline*}

The values of $E$, $\Gamma$ are given p.~800, but for the present purpose we only require $E^{2}$, and $\Gamma^{2}$ to the fifth order, viz.\ the values of these are at once found to be
\[
\Gamma^{2}=\gamma^{2}\{1+\tfrac{2}{3}\gamma^{2}-\tfrac{1}{6}e^{2}-\tfrac{1}{6}e'^{2}\},
\]
whence also $E^{2}=e^{2}$ and $\Gamma^{2}=\gamma^{2}$.

The formulae of pp.\ 237--238 now give
\begin{multline*}
\frac{I}{Nt}=(1+Q)^{-1}\\
+(-\tfrac14 \gamma^{2} + \tfrac{3}{8}\gamma^{4} - \tfrac{5}{16}\gamma^{2}e^{2} - \tfrac{3}{16}\gamma^{2}e'^{2} - \tfrac{3}{64}\gamma^{6} + \tfrac{15}{32}\gamma^{4}e^{2} + \tfrac{9}{32}\gamma^{4}e'^{2}\\
+ \tfrac{15}{64}\gamma^{2}e^{4} + \tfrac{9}{32}\gamma^{2}e^{2}e'^{2} + \tfrac{9}{128}\gamma^{2}e'^{4})m^{2}(1+Q)^{-1}\\
+(\tfrac{9}{32}\gamma^{2} - \tfrac{27}{64}\gamma^{4} - \tfrac{9}{64}\gamma^{2}e^{2})m^{3}(1+Q)^{-1}\\
+(-\tfrac{27}{64}\gamma^{2} + \tfrac{81}{128}\gamma^{4} + \tfrac{81}{256}\gamma^{2}e^{2} - \tfrac{135}{256}\gamma^{2}e'^{2})m^{4}(1+Q)^{-1}\\
+(-\tfrac{225}{256}\gamma^{2} + \tfrac{675}{512}\gamma^{4} + \tfrac{675}{1024}\gamma^{2}e^{2} - \tfrac{675}{512}\gamma^{2}e'^{2})m^{5}(1+Q)^{-1}\\
+(-\tfrac{3375}{1024}\gamma^{2})m^{6}(1+Q)^{-1};
\end{multline*}
(Observe that writing herein $Q=0$, and omitting the terms in $m^{5}$ and $m^{6}$ in the coefficient of $(1+Q)^{-1}$, and the term in $m^{6}$ in the coefficient of $(1+Q)^{-2}$, we have the original formula of p.~237.)

\begin{multline*}
\frac{g}{Nt}=(1+Q)^{-1}\\
+(\tfrac14 \gamma^{2} - \tfrac{3}{8}\gamma^{4} + \tfrac{5}{16}\gamma^{2}e^{2} + \tfrac{3}{16}\gamma^{2}e'^{2} + \tfrac{3}{64}\gamma^{6} - \tfrac{15}{32}\gamma^{4}e^{2} - \tfrac{9}{32}\gamma^{4}e'^{2}\\
- \tfrac{15}{64}\gamma^{2}e^{4} - \tfrac{9}{32}\gamma^{2}e^{2}e'^{2} - \tfrac{9}{128}\gamma^{2}e'^{4})m^{2}(1+Q)^{-1}\\
+(-\tfrac{9}{32}\gamma^{2} + \tfrac{27}{64}\gamma^{4} + \tfrac{9}{64}\gamma^{2}e^{2})m^{3}(1+Q)^{-1}\\
+(\tfrac{27}{64}\gamma^{2} - \tfrac{81}{128}\gamma^{4} - \tfrac{81}{256}\gamma^{2}e^{2} + \tfrac{135}{256}\gamma^{2}e'^{2})m^{4}(1+Q)^{-1}\\
+(\tfrac{225}{256}\gamma^{2} - \tfrac{675}{512}\gamma^{4} - \tfrac{675}{1024}\gamma^{2}e^{2} + \tfrac{675}{512}\gamma^{2}e'^{2})m^{5}(1+Q)^{-1}\\
+(\tfrac{3375}{1024}\gamma^{2})m^{6}(1+Q)^{-1};
\end{multline*}
(where writing $Q=0$, and omitting the terms in $m^{5}$ and $m^{6}$ in the coefficient of $(1+Q)^{-1}$, and the term in $m^{6}$ in the coefficient of $(1+Q)^{-2}$, we have the original formula of p.~237).

And
\begin{multline*}
\frac{h}{Nt}=(1+Q)^{-1}\\
+(\tfrac14 \gamma^{2} - \tfrac{3}{8}\gamma^{4} + \tfrac{5}{16}\gamma^{2}e^{2} + \tfrac{3}{16}\gamma^{2}e'^{2} + \tfrac{3}{64}\gamma^{6} - \tfrac{15}{32}\gamma^{4}e^{2} - \tfrac{9}{32}\gamma^{4}e'^{2}\\
- \tfrac{15}{64}\gamma^{2}e^{4} - \tfrac{9}{32}\gamma^{2}e^{2}e'^{2} - \tfrac{9}{128}\gamma^{2}e'^{4})m^{2}(1+Q)^{-1}\\
+(-\tfrac{9}{32}\gamma^{2} + \tfrac{27}{64}\gamma^{4} + \tfrac{9}{64}\gamma^{2}e^{2})m^{3}(1+Q)^{-1}\\
+(\tfrac{27}{64}\gamma^{2} - \tfrac{81}{128}\gamma^{4} - \tfrac{81}{256}\gamma^{2}e^{2} + \tfrac{135}{256}\gamma^{2}e'^{2})m^{4}(1+Q)^{-1}\\
+(\tfrac{225}{256}\gamma^{2} - \tfrac{675}{512}\gamma^{4} - \tfrac{675}{1024}\gamma^{2}e^{2} + \tfrac{675}{512}\gamma^{2}e'^{2})m^{5}(1+Q)^{-1}\\
+(\tfrac{3375}{1024}\gamma^{2})m^{6}(1+Q)^{-1};
\end{multline*}
(where writing $Q=0$, and omitting the terms in $m^{4}$ and $m^{5}$ in the coefficient of $(1+Q)^{-1}$, and the term in $m^{5}$ in the coefficient of $(1+Q)^{-2}$, we have the original formula of p.~238).

We hence have
\begin{align*}
I &= nt\{A + AQ + Bq + Cq^{2}\},\\
g &= nt\{A' + A'Q + B'q + C'q^{2}\},\\
h &= nt\{A'' + A''Q + B''q + C''q^{2}\},
\end{align*}
where (omitting the terms in $\tfrac{1}{n}$):
\begin{align*}
A &= 1 + (-\tfrac34 \gamma^{2} + \tfrac{15}{32}\gamma^{4} - \tfrac{5}{16}\gamma^{2}e^{2} - \tfrac{3}{16}\gamma^{2}e'^{2} + \tfrac{35}{128}\gamma^{6} - \tfrac{105}{64}\gamma^{4}e^{2} - \tfrac{63}{64}\gamma^{4}e'^{2}\\
&\quad + \tfrac{15}{64}\gamma^{2}e^{4} + \tfrac{45}{64}\gamma^{2}e^{2}e'^{2} + \tfrac{45}{256}\gamma^{2}e'^{4})m^{2}\\
&\quad + (-\tfrac{27}{16}\gamma^{2} + \tfrac{405}{128}\gamma^{4} - \tfrac{135}{64}\gamma^{2}e^{2} - \tfrac{81}{64}\gamma^{2}e'^{2} + \tfrac{945}{256}\gamma^{6})m^{3}\\
&\quad + (-\tfrac{225}{64}\gamma^{2} + \tfrac{3375}{512}\gamma^{4} - \tfrac{1125}{256}\gamma^{2}e^{2} - \tfrac{675}{256}\gamma^{2}e'^{2})m^{4}\\
&\quad + (-\tfrac{3375}{256}\gamma^{2})m^{5};
\end{align*}
\begin{align*}
B &= -\tfrac34 m^{2} - \tfrac{27}{16}m^{3} - \tfrac{225}{64}m^{4},\\
C &= -\tfrac34 m^{2} + \tfrac{27}{32}m^{3} + \tfrac{675}{128}m^{4};
\end{align*}
\begin{align*}
A' &= 1 + (\tfrac14 \gamma^{2} - \tfrac{5}{32}\gamma^{4} + \tfrac{5}{16}\gamma^{2}e^{2} + \tfrac{3}{16}\gamma^{2}e'^{2} - \tfrac{15}{128}\gamma^{6} + \tfrac{35}{64}\gamma^{4}e^{2} + \tfrac{21}{64}\gamma^{4}e'^{2}\\
&\quad - \tfrac{15}{64}\gamma^{2}e^{4} - \tfrac{45}{64}\gamma^{2}e^{2}e'^{2} - \tfrac{45}{256}\gamma^{2}e'^{4})m^{2}\\
&\quad + (\tfrac{9}{16}\gamma^{2} - \tfrac{135}{128}\gamma^{4} + \tfrac{45}{64}\gamma^{2}e^{2} + \tfrac{27}{64}\gamma^{2}e'^{2})m^{3}\\
&\quad + (\tfrac{75}{64}\gamma^{2} - \tfrac{1125}{512}\gamma^{4} + \tfrac{375}{256}\gamma^{2}e^{2} + \tfrac{225}{256}\gamma^{2}e'^{2})m^{4}\\
&\quad + (\tfrac{1125}{256}\gamma^{2})m^{5};
\end{align*}
\begin{align*}
B' &= -\tfrac34 m^{2} - \tfrac{27}{16}m^{3} - \tfrac{225}{64}m^{4},\\
C' &= \tfrac14 m^{2} - \tfrac{27}{32}m^{3} - \tfrac{675}{128}m^{4};
\end{align*}
\begin{align*}
A'' &= (-\tfrac14 \gamma^{2} + \tfrac{3}{16}\gamma^{4} - \tfrac{5}{32}\gamma^{2}e^{2} - \tfrac{3}{32}\gamma^{2}e'^{2} - \tfrac{5}{64}\gamma^{6} + \tfrac{15}{32}\gamma^{4}e^{2} + \tfrac{9}{32}\gamma^{4}e'^{2}\\
&\quad - \tfrac{15}{128}\gamma^{2}e^{4} - \tfrac{45}{128}\gamma^{2}e^{2}e'^{2} - \tfrac{45}{512}\gamma^{2}e'^{4})m^{2}\\
&\quad + (\tfrac{9}{16}\gamma^{2} - \tfrac{135}{128}\gamma^{4} + \tfrac{45}{64}\gamma^{2}e^{2} + \tfrac{27}{64}\gamma^{2}e'^{2})m^{3}\\
&\quad + (\tfrac{75}{64}\gamma^{2} - \tfrac{1125}{512}\gamma^{4} + \tfrac{375}{256}\gamma^{2}e^{2} + \tfrac{225}{256}\gamma^{2}e'^{2})m^{4}\\
&\quad + (\tfrac{1125}{256}\gamma^{2})m^{5};
\end{align*}
\begin{align*}
B'' &= -\tfrac34 m^{2} - \tfrac{27}{16}m^{3} - \tfrac{225}{64}m^{4},\\
C'' &= \tfrac14 m^{2} - \tfrac{27}{32}m^{3} - \tfrac{675}{128}m^{4}.
\end{align*}

\newpage

And in terms $BQ$, $B'Q$, $B''Q$, we have
\[
Q=(1-\tfrac34 \gamma^{2}+\dots)m^{2},
\]
and in the terms $CQ^{2}$, $C'Q^{2}$, $C''Q^{2}$, simply $Q^{2}=m^{4}$.

Hence finally the required values of $I$, $g$, $h$, are
\begin{align*}
I &= nt\{1 + [-\tfrac34 \gamma^{2}m^{2} - \tfrac{27}{16}\gamma^{2}m^{3}]\} \times \dots \\
\end{align*}
\begin{multline*}
g = nt\{1 + [\tfrac14 \gamma^{2} + \tfrac32 e^{2} - \tfrac{5}{8}e'^{2} + \tfrac{15}{32}\gamma^{4} - \tfrac{63}{16}\gamma^{2}e^{2} + \tfrac{27}{16}\gamma^{2}e'^{2} + \tfrac{15}{8}e^{4} + \tfrac{3}{4}e^{2}e'^{2} - \tfrac{15}{32}e'^{4}]m^{2}\\
+ [-\tfrac{27}{32}\gamma^{2} + \tfrac{81}{16}\gamma^{4} - \tfrac{135}{32}\gamma^{2}e^{2} + \tfrac{81}{32}\gamma^{2}e'^{2}]m^{3}\\
+ [\tfrac{75}{64}\gamma^{2} - \tfrac{1125}{256}\gamma^{4} - \tfrac{375}{128}\gamma^{2}e^{2} + \tfrac{225}{128}\gamma^{2}e'^{2} + \tfrac{3375}{128}e^{4} - \tfrac{3375}{64}\gamma^{2}e^{2} - \tfrac{3375}{256}\gamma^{2}e'^{2}]m^{4}\\
+ [\tfrac{1125}{256}\gamma^{2} - \tfrac{16875}{1024}\gamma^{4} - \tfrac{5625}{512}\gamma^{2}e^{2} + \tfrac{3375}{512}\gamma^{2}e'^{2}]m^{5}\};
\end{multline*}
\begin{multline*}
h = nt\{[-\tfrac14 \gamma^{2} - \tfrac32 e^{2} + \tfrac{5}{8}e'^{2} - \tfrac{9}{32}\gamma^{4} + \tfrac{27}{16}\gamma^{2}e^{2} - \tfrac{27}{32}\gamma^{2}e'^{2} - \tfrac{15}{8}e^{4} - \tfrac{3}{4}e^{2}e'^{2} + \tfrac{15}{32}e'^{4}]m^{2}\\
+ [\tfrac{9}{16}\gamma^{2} - \tfrac{135}{64}\gamma^{4} + \tfrac{45}{32}\gamma^{2}e^{2} - \tfrac{27}{32}\gamma^{2}e'^{2}]m^{3}\\
+ [\tfrac{75}{64}\gamma^{2} - \tfrac{1125}{256}\gamma^{4} - \tfrac{375}{128}\gamma^{2}e^{2} + \tfrac{225}{128}\gamma^{2}e'^{2} + \tfrac{3375}{256}\gamma^{4} - \tfrac{3375}{128}\gamma^{2}e^{2} + \tfrac{3375}{256}\gamma^{2}e'^{2}]m^{4}\\
+ [\tfrac{1125}{256}\gamma^{2} - \tfrac{16875}{1024}\gamma^{4} - \tfrac{5625}{512}\gamma^{2}e^{2} + \tfrac{3375}{512}\gamma^{2}e'^{2}]m^{5}\};
\end{multline*}
which values satisfy, as they should do, the equation $l+g+h=nt$.

I recall that the precise signification of the constants is as follows: $n$ is the coefficient of $t$ in the expression of the Moon's longitude in terms of the time, $a$ the corresponding elliptic value of the mean distance ($m^{2}a^{3}=\text{sum of masses}$), $e$ the eccentricity, such that in the expression of the longitude the coefficient of the leading term of the equation of the centre has its elliptic value $\tfrac52 e - \tfrac{5}{16}e^{3} + \tfrac{15}{128}e^{5}$, and $\gamma$ the sine of the half-inclination, such that in the expression of the latitude the coefficient of the leading term has its elliptic value $=2\gamma - \tfrac{2}{3}\gamma^{3} - \tfrac{1}{4}\gamma^{5} + \tfrac{35}{24}\gamma^{7} + \tfrac{1}{4}\gamma^{3}e^{2} - \tfrac{11}{12}\gamma^{3}e'^{2}$; $n'$, $a'$ are the mean motion and mean distance of the Sun, $m=\dfrac{n'}{n}$, and $e'$ is the eccentricity of the Sun's orbit, considered as constant.


\vspace{2em}

%==============================================================================
% PAPER [481]
%==============================================================================
\newpage
\section*{[481 --- On the Expression of M.\ Delaunay's $\dfrac{h+g}{n}$ in terms of his Finally Adopted Constants]}
\addcontentsline{toc}{section}{[481] On the Expression of $\dfrac{h+g}{n}$}

\markboth{Cayley, Collected Papers, Vol.\ VII}{[481] Delaunay's $h+g$}

\begin{center}
\textbf{ON THE EXPRESSION OF M.\ DELAUNAY'S $\dfrac{h+g}{n}$ IN TERMS OF HIS FINALLY ADOPTED CONSTANTS.}
\end{center}

[From the \textit{Monthly Notices of the Royal Astronomical Society}, vol.\ xxxii.\ (1871--72), p.\ 74.]

\bigskip

I \textsc{had} the pleasure of receiving from M.\ Delaunay a letter dated Paris, 17th Dec.\ 1871, in which he informs me that, on referring to his papers, he had found there expressions for $I$, $g$, $h$, identical with those given by me in the November Number of the Monthly Notices,---with only a single typographical error, $\tfrac54 e'^{2}m^{2}$ instead of $\tfrac54 e^{2}m^{2}$ [ante p.~532, corrected] in my expression of $A$.

M.\ Delaunay mentions also that he had obtained four additional terms in the expression for $h+g$ (longitude of the Moon's perigee), and that the complete expression in terms of the finally adopted constants is
\begin{multline*}
\frac{h+g}{n}=t\{1+[\tfrac14 \gamma^{2} + \tfrac32 e^{2} - \tfrac{5}{8}e'^{2} + \tfrac{15}{32}\gamma^{4} - \tfrac{63}{16}\gamma^{2}e^{2} + \tfrac{27}{16}\gamma^{2}e'^{2}\\
+ \tfrac{15}{8}e^{4} + \tfrac{3}{4}e^{2}e'^{2} - \tfrac{15}{32}e'^{4}]m^{2}\\
+ [-\tfrac{27}{32}\gamma^{2} + \tfrac{81}{16}\gamma^{4} - \tfrac{135}{32}\gamma^{2}e^{2} + \tfrac{81}{32}\gamma^{2}e'^{2}]m^{3}\\
+ [\tfrac{75}{64}\gamma^{2} - \tfrac{1125}{256}\gamma^{4} + \tfrac{375}{128}\gamma^{2}e^{2} - \tfrac{225}{128}\gamma^{2}e'^{2}]m^{4}\\
+ [\tfrac{1125}{256}\gamma^{2} - \tfrac{16875}{1024}\gamma^{4} + \tfrac{5625}{512}\gamma^{2}e^{2} - \tfrac{3375}{512}\gamma^{2}e'^{2}]m^{5}\\
+ [-\tfrac{3375}{256}\gamma^{2} + \tfrac{50625}{1024}\gamma^{4}]m^{6}\}.
\end{multline*}

[Observe that $h+g$ is $=nt-I$, and compare with the expression for $I$, ante p.~532.]


\vspace{2em}

%==============================================================================
% PAPER [482]
%==============================================================================
\newpage
\section*{[482 --- Note on a Pair of Differential Equations in the Lunar Theory]}
\addcontentsline{toc}{section}{[482] Note on a Pair of Differential Equations}

\markboth{Cayley, Collected Papers, Vol.\ VII}{[482] Differential Equations}

\begin{center}
\textbf{NOTE ON A PAIR OF DIFFERENTIAL EQUATIONS IN THE LUNAR THEORY.}
\end{center}

[From the \textit{Monthly Notices of the Royal Astronomical Society}, vol.\ xxxii.\ (1871--72), pp.\ 31--32.]

\bigskip

\textsc{The} equations
\[
\frac{d^{2}\rho}{dt^{2}}-\rho\left(\frac{d\upsilon}{dt}\right)^{2}=j\frac{m^{2}}{\rho^{2}}\left\{-\frac12 \sin(2\upsilon-2mt)\right\},
\]
\[
\frac{d}{dt}\left(\rho^{2}\frac{d\upsilon}{dt}\right)=j\frac{m^{2}}{\rho^{2}}\left\{-\frac12 + \frac12 \cos(2\upsilon-2mt)\right\},
\]
taking therein $j=k=1$ in effect present themselves in the Lunar Theory, and particular integrals in series have been obtained, the development being carried to a great extent; but I give the results only as far as $m^{6}$, viz., writing
\[
\upsilon-t=\zeta, \quad \rho=1+p,
\]
we have
\begin{align*}
\zeta &= \sin 2D\left(\tfrac34 m^{2} + \tfrac{11}{24}m^{3} + \tfrac{239}{288}m^{4}\right) + \sin 4D\left(\tfrac{1}{48}m^{4}\right),\\
p &= 1 + \tfrac13 m^{2} - \tfrac{11}{18}m^{4} + \left(\tfrac{3}{8}m^{2} + \tfrac{11}{24}m^{3} + \tfrac{847}{1152}m^{4}\right)\cos 2D + \tfrac{1}{16}m^{4}\cos 4D.
\end{align*}

In the Lunar Theory $j$ and $k$ are properly each $=\dfrac{1}{1+\tfrac{m}{E}}$ ($E$ the mass of the Earth, $m'$ that of the Sun), but they are taken to be $=1$; the numerical difference is inappreciable; but there would be a considerable theoretical advantage in retaining in the equations the coefficients $j$, $k$ [regarded as each of them $=1$]: in fact, the developments could then be arranged according to the powers of $k$, that is according to the powers of the disturbing force; whereas, when $k$ is taken $=1$, we have only a development in powers of $m$, and since $m$ also presents itself through the coefficient $2-2m$ of $t$ in $2\upsilon-2mt$, terms which are really of different orders in regard to the disturbing force, are united together into a single term: so that, instead of a term of the form $(Ak + Bk^{2} + \&c.)\ m^{p}$, where $A$, $B$, are numerical, we have the term $(A+B+\dots)\ m^{p}$, where of course $A+B+\dots$ is given as a single numerical coefficient.

There is no equal advantage in retaining the two coefficients $k$, $j$, as this only serves to show how a term arises from the central and tangential forces respectively; thus retaining these coefficients, the integrals as far as $m^{4}$ are
\begin{align*}
\upsilon-t &= \sin 2D\left(\tfrac34 k + \tfrac{11}{24}j\right)m^{2},\\
\rho &= 1 + \tfrac13 m^{2}k + \left(\tfrac{3}{8}k + \tfrac{11}{24}j\right)m^{2}\cos 2D,
\end{align*}
agreeing with the former result when $k=j=1$; but there is, nevertheless, some interest in retaining the two coefficients.

I hope to develope the results somewhat further, and to communicate them to the Society.


\vspace{2em}

%==============================================================================
% PAPER [483]
%==============================================================================
\newpage
\section*{[483 --- On a Pair of Differential Equations in the Lunar Theory]}
\addcontentsline{toc}{section}{[483] On a Pair of Differential Equations}

\markboth{Cayley, Collected Papers, Vol.\ VII}{[483] Differential Equations}

\begin{center}
\textbf{ON A PAIR OF DIFFERENTIAL EQUATIONS IN THE LUNAR THEORY.}
\end{center}

[From the \textit{Monthly Notices of the Royal Astronomical Society}, vol.\ xxxii.\ (1871--72), pp.\ 201--206.]

\bigskip

\textsc{I consider} the differential equations
\[
\frac{d^{2}\rho}{dt^{2}}-\rho\left(\frac{d\upsilon}{dt}\right)^{2}=j\frac{m^{2}}{\rho^{2}}\left\{-\frac12 \sin(2\upsilon-2mt)\right\},
\]
\[
\frac{d}{dt}\left(\rho^{2}\frac{d\upsilon}{dt}\right)=j\frac{m^{2}}{\rho^{2}}\left\{-\frac12 + \frac12 \cos(2\upsilon-2mt)\right\},
\]
which when $j=k=1$ give the following equations in the lunar theory ($D=t-mt$):
\begin{multline*}
\rho = 1 + \tfrac13 m^{2} - \tfrac{11}{18}m^{4} - \tfrac{473}{1080}m^{6} - \tfrac{133}{240}m^{8} - \tfrac{14887}{12096}m^{10} - \tfrac{1048231}{362880}m^{12}\\
+ \cos 2D\left\{\tfrac38 m^{2} + \tfrac{11}{24}m^{3} + \tfrac{847}{1152}m^{4} + \tfrac{485}{768}m^{5} + \tfrac{44191}{41472}m^{6} + \tfrac{275495}{165888}m^{7}\right.\\
\left. + \tfrac{31327127}{8957952}m^{8}\right\}\\
+ \cos 4D\left\{\tfrac{1}{16}m^{4} + \tfrac{1}{12}m^{5} + \tfrac{1283}{10368}m^{6} + \tfrac{1105}{5184}m^{7} + \tfrac{1036609}{2985984}m^{8}\right\}\\
+ \cos 6D\left\{\tfrac{1}{80}m^{6} + \tfrac{23}{720}m^{7} + \tfrac{131}{2400}m^{8}\right\}\\
+ \cos 8D\left\{\tfrac{1}{448}m^{8}\right\},
\end{multline*}
or as far as $m^{8}$,
\begin{multline*}
\upsilon = t + \sin 2D\left\{\tfrac34 m^{2} + \tfrac{11}{24}m^{3} + \tfrac{239}{288}m^{4} + \tfrac{1261}{3456}m^{5} + \tfrac{20971}{20736}m^{6} + \tfrac{85823}{41472}m^{7} + \tfrac{7296259}{1492992}m^{8}\right\}\\
+ \sin 4D\left\{\tfrac{1}{48}m^{4} + \tfrac{5}{144}m^{5} + \tfrac{1571}{20736}m^{6} + \tfrac{671}{5184}m^{7} + \tfrac{394027}{995328}m^{8}\right\}\\
+ \sin 6D\left\{\tfrac{1}{240}m^{6} + \tfrac{7}{480}m^{7} + \tfrac{3397}{86400}m^{8}\right\}\\
+ \sin 8D\left\{\tfrac{1}{1344}m^{8}\right\},
\end{multline*}
($\rho$ is given by M.\ Delaunay only as far as $m^{10}$, the additional terms of $\rho$ and expression for $\upsilon$ were kindly communicated to me by Prof.\ Adams); and
\begin{multline*}
\upsilon = t\\
+ \sin 2D\left(\tfrac34 m^{2} + \tfrac{11}{24}m^{3} + \tfrac{239}{288}m^{4} + \tfrac{1261}{3456}m^{5} + \tfrac{20971}{20736}m^{6} + \tfrac{85823}{41472}m^{7}\right.\\
\left. + \tfrac{7296259}{1492992}m^{8} + \tfrac{35819887}{2985984}m^{9} + \tfrac{1009171205}{35831808}m^{10}\right)\\
+ \sin 4D\left(\tfrac{1}{48}m^{4} + \tfrac{5}{144}m^{5} + \tfrac{1571}{20736}m^{6} + \tfrac{671}{5184}m^{7} + \tfrac{394027}{995328}m^{8}\right.\\
\left. + \tfrac{1705303}{1990656}m^{9} + \tfrac{106078709}{59719680}m^{10}\right)\\
+ \sin 6D\left(\tfrac{1}{240}m^{6} + \tfrac{7}{480}m^{7} + \tfrac{3397}{86400}m^{8} + \tfrac{425}{6048}m^{9} + \tfrac{17166077}{87091200}m^{10}\right)\\
+ \sin 8D\left(\tfrac{1}{1344}m^{8} + \tfrac{11}{6720}m^{9} + \tfrac{311}{129600}m^{10}\right)\\
+ \sin 10D\left(\tfrac{1}{8064}m^{10}\right)
\end{multline*}
(Delaunay, t.\ ii.\ pp.\ 815, 836, 845).

To integrate the original equations write
\[
\rho = 1 + p_{1} + p_{2} + p_{3} + \dots,
\]
\[
\upsilon = t + \upsilon_{1} + \upsilon_{2} + \upsilon_{3} + \dots,
\]
where the suffixes indicate the degrees in the coefficients $k$, $j$ conjointly: the equations for $p_{n}$, $\upsilon_{n}$ take the form
\[
\frac{d^{2}p_{n}}{dt^{2}} - 2\frac{d\upsilon_{n}}{dt} = U_{n}, \quad \frac{d^{2}\upsilon_{n}}{dt^{2}} + 2\frac{dp_{n}}{dt} = P_{n},
\]
where $V_{n}$, $U_{n}$, $P_{n}$, $Q_{n}$ do not contain $p_{n}$ or $\upsilon_{n}$.

From the second equation we have
\[
\frac{d\upsilon_{n}}{dt} + 2p_{n} = \Omega_{n} + \int P_{n}\,dt,
\]
where $\Omega_{n}$ is a constant of integration, the integral $\int P_{n}\,dt$ containing only periodic terms; and then adding twice this to the first equation we have
\[
\frac{d^{2}p_{n}}{dt^{2}} + p_{n} + \upsilon_{n} + 2U_{n} = 2\Omega_{n} + Q_{n} + 2\int p_{n}\,dt,
\]
which determines $p_{n}$; and substituting its value in the other equation we have $\dfrac{d\upsilon_{n}}{dt}$, and thence $\upsilon_{n}$; the constant $\Omega_{n}$ is determined so that $\dfrac{d\upsilon_{n}}{dt}$ may contain no constant term.

We have
\begin{align*}
p_{1}' &= -\frac{d\upsilon_{1}}{dt} - \tfrac12 p_{1}^{2} + \int P_{1}\,dt,\\
\frac{d^{2}p_{2}}{dt^{2}} + p_{2} &= 2p_{1}\frac{d\upsilon_{1}}{dt} + Q_{1}p_{1}p_{2} - 4p_{1}^{2}, \quad \&\text{c.}
\end{align*}

\begin{align*}
Q_{1} &= \tfrac12 m^{2}(1 + \tfrac32 \cos 2D),\\
Q_{2} &= km^{2}\left\{-\tfrac32 \upsilon_{1}\sin 2D + p_{1}\left(\tfrac12 + \tfrac34 \cos 2D\right)\right\},\\
Q_{3} &= km^{2}\{-3\upsilon_{2}\sin 2D - 3\upsilon_{1}^{2}\cos 2D\\
&\quad + p_{2}\tfrac32 \sin 2D + p_{1}(1 + \tfrac32 \cos 2D)\},\\
P_{1} &= jm^{2}(-\tfrac12 \sin 2D),\\
P_{2} &= jm^{2}(-3\upsilon_{1}\cos 2D - 3p_{1}\sin 2D),\\
P_{3} &= jm^{2}\{-3\upsilon_{2}\cos 2D + 3\upsilon_{1}^{2}\sin 2D\\
&\quad - 6p_{1}\upsilon_{1}\cos 2D + (2p_{2} + p_{1}^{2})\tfrac32 \sin 2D\},\\
&\quad \&\text{c.}
\end{align*}

In particular attending to the values of $P_{1}$, $Q_{1}$ the equations for $p_{1}$, $\upsilon_{1}$ are in their original form
\[
\frac{d^{2}\rho_{1}}{dt^{2}}-2\rho_{1}\frac{d\upsilon_{1}}{dt}=jm^{2}\left(-\frac12\sin 2D\right), \quad \frac{d}{dt}(2\upsilon_{1})=km^{2}\left(\frac12 + \frac34 \cos 2D\right),
\]
whence in the transformed form they are
\[
\frac{d^{2}p_{1}}{dt^{2}}+2p_{1}=\Omega_{1}+\tfrac12 jm^{2}\frac{3-8m+4m^{2}}{(1-m)^{2}}\cos 2D,
\]
and
\[
\frac{d\upsilon_{1}}{dt}+p_{1}=2\Omega_{1}+\tfrac12 jm^{2}(1+\tfrac32 \cos 2D)+\tfrac14 \frac{jm^{2}(7-8m+4m^{2})}{(1-m)^{2}(3-8m+4m^{2})}\cos 2D.
\]
Thus the constant term of $p_{1}$ is $2\Omega_{1}+\tfrac34 km^{2}$, giving in $\dfrac{d\upsilon}{dt}$ a constant term $-3\Omega_{1}-\tfrac32 km^{2}$; this must vanish, or we have $\Omega_{1}=-\tfrac14 km^{2}$; and the equations thus become
\[
\frac{d^{2}p_{1}}{dt^{2}}+2p_{1}=-\tfrac14 jm^{2}+\tfrac12 \frac{jm^{2}(3-8m+4m^{2})}{(1-m)^{2}}\cos 2D,
\]
\[
\frac{d\upsilon_{1}}{dt}+p_{1}=-\tfrac34 km^{2}+(\tfrac34 jm^{2}+\tfrac14 k)\cos 2D,
\]
and then completing the integration
\begin{align*}
p_{1} &= -\tfrac18 jm^{2} + \frac{\tfrac12 jm^{2}(3-8m+4m^{2})}{(1-m)^{2}(3-8m+4m^{2})}\cos 2D,\\
\upsilon_{1} &= \frac{\tfrac34 jm^{2} + \tfrac14 k}{2(1-m)}\sin 2D - \frac{\tfrac12 jm^{2}(3-8m+4m^{2})}{(1-m)^{2}(3-8m+4m^{2})}\sin 2D,
\end{align*}
which are the accurate values of $p_{1}$ and $\upsilon_{1}$.

Expanding as far as $m^{7}$ we have
\begin{multline*}
p_{1} = k\{-\tfrac14 m^{2}\} + \cos 2D\left\{k(-\tfrac14 m^{2} - \tfrac16 m^{3} - \tfrac{7}{48}m^{4} - \tfrac{1}{6}m^{5} - \tfrac{91}{576}m^{6} - \tfrac{17}{96}m^{7})\right.\\
\left. + j(-\tfrac14 m^{2} - \tfrac13 m^{3} - \tfrac{19}{48}m^{4} - \tfrac{11}{24}m^{5} - \tfrac{211}{576}m^{6} - \tfrac{191}{480}m^{7})\right\},
\end{multline*}
which for $j=k$ is
\[
= k\{-\tfrac14 m^{2} - \tfrac12 m^{3} - \tfrac{13}{24}m^{4} - \tfrac{15}{24}m^{5} - \tfrac{302}{576}m^{6} - \tfrac{287}{480}m^{7}\};
\]
and
\begin{multline*}
\upsilon_{1} = \sin 2D\left\{k(\tfrac38 m^{2} + \tfrac{11}{48}m^{3} + \tfrac{5}{16}m^{4} + \tfrac{61}{288}m^{5} + \tfrac{287}{2304}m^{6} + \tfrac{2261}{11520}m^{7})\right.\\
\left. + j(-\tfrac18 m^{2} - \tfrac{1}{12}m^{3} - \tfrac{1}{24}m^{4} + \tfrac{1}{144}m^{5} + \tfrac{175}{2304}m^{6} + \tfrac{733}{5760}m^{7})\right\},
\end{multline*}
which for $j=k$ is
\[
= k\{\tfrac14 m^{2} + \tfrac{3}{16}m^{3} + \tfrac{7}{48}m^{4} + \tfrac{65}{576}m^{5} + \tfrac{77}{576}m^{6} + \tfrac{2997}{11520}m^{7}\}.
\]

I have, not in general, but for the value $j=k$, calculated $p_{n}$ and $\upsilon_{n}$ as far as $m^{7}$: I have not made the calculation for $p_{3}$ and $\upsilon_{3}$, but their values may be deduced from the foregoing values of $p$, $\upsilon$; the final expressions (when $j=k$) of $\rho=1+p_{1}+p_{2}+p_{3}+\dots$ and $\upsilon=t+\upsilon_{1}+\upsilon_{2}+\upsilon_{3}+\dots$ are
\begin{multline*}
\rho = 1\\
+ k\{-\tfrac14 m^{2}\} + k^{2}\{-\tfrac14 m^{3}\}\\
+ \cos 2D\left\{k(-\tfrac14 m^{2} - \tfrac16 m^{3} - \tfrac{7}{48}m^{4} - \tfrac16 m^{5} - \tfrac{91}{576}m^{6} - \tfrac{17}{96}m^{7})\right.\\
\left. + k^{2}(\tfrac38 m^{2} + \tfrac34 m^{3} + \tfrac{65}{48}m^{4} + \tfrac{115}{48}m^{5} + \tfrac{2129}{576}m^{6} + \tfrac{799}{120}m^{7})\right.\\
\left. + k^{3}(-\tfrac{1}{16}m^{4} - \tfrac{5}{16}m^{5} - \tfrac{119}{192}m^{6})\right\}\\
+ \cos 4D\left\{k^{2}(-\tfrac18 m^{4} - \tfrac{7}{24}m^{5} - \tfrac{29}{36}m^{6} - \tfrac{383}{288}m^{7})\right.\\
\left. + k^{3}(+\tfrac{1}{16}m^{4} + \tfrac{7}{24}m^{5} + \tfrac{131}{144}m^{6})\right\}\\
+ \cos 6D\{k^{3}(-\tfrac{1}{48}m^{6})\}\\
+ \sin 2D\left\{k(\tfrac38 m^{2} + \tfrac{11}{48}m^{3} + \tfrac{5}{16}m^{4} + \tfrac{61}{288}m^{5} + \tfrac{287}{2304}m^{6} + \tfrac{2261}{11520}m^{7})\right.\\
\left. + k^{2}(-\tfrac18 m^{2} - \tfrac{7}{24}m^{3} - \tfrac{47}{72}m^{4} - \tfrac{179}{144}m^{5} - \tfrac{5431}{3456}m^{6} - \tfrac{4829}{2016}m^{7})\right.\\
\left. + k^{3}(-\tfrac{1}{16}m^{4} - \tfrac{1}{6}m^{5} - \tfrac{43}{96}m^{6})\right\}\\
+ \sin 4D\left\{k^{2}(\tfrac{1}{48}m^{4} + \tfrac{5}{72}m^{5} + \tfrac{157}{1728}m^{6} + \tfrac{671}{3456}m^{7})\right.\\
\left. + k^{3}(+\tfrac{1}{24}m^{4} + \tfrac{35}{144}m^{5} + \tfrac{407}{864}m^{6})\right\}\\
+ \sin 6D\{k^{3}(+\tfrac{1}{144}m^{6})\};
\end{multline*}
which for $k=1$ agree with the foregoing formulae (verifying them as far as $m^{7}$); the present formulae exhibit the manner in which the expressions depend on the several powers of the disturbing force.


\vspace{2em}

%==============================================================================
% PAPER [484]
%==============================================================================
\newpage
\section*{[484 --- On the Variations of the Position of the Orbit in the Planetary Theory]}
\addcontentsline{toc}{section}{[484] On the Variations of the Position of the Orbit}

\markboth{Cayley, Collected Papers, Vol.\ VII}{[484] Variations of the Orbit}

\begin{center}
\textbf{ON THE VARIATIONS OF THE POSITION OF THE ORBIT IN THE PLANETARY THEORY.}
\end{center}

[From the \textit{Monthly Notices of the Royal Astronomical Society}, vol.\ xxxii.\ (1871--72), pp.\ 206--211.]

\bigskip

\textsc{It} has always appeared to me that in the Planetary Theory, more especially when the method of the variation of the elements is made use of, there is a difficulty as to the proper mode of dealing with the inclinations and longitudes of the nodes, hindering the ulterior development of the theory.

Considering the case of two planets $m$, $m'$, and referring their orbits to any fixed plane and fixed origin of longitudes therein, let $\theta$, $\theta'$ be the longitudes of the nodes, $\phi$, $\phi'$ the inclinations ($p=\tan\phi\sin\theta$, $q=\tan\phi\cos\theta$, \&c., as usual); then the disturbing functions for $m$, $m'$ respectively are developed, not explicitly in terms of $\theta$, $\theta'$, $\phi$, $\phi'$, but in terms of $\Phi$, the mutual inclination of the two orbits, and of $@$, $@'$ the longitudes in the two orbits respectively of the mutual node of the two orbits; $\Phi$ and $@$, $@'$ being functions (and complicated ones) of $\phi$, $\phi'$, $\theta$, $\theta'$.

Moreover, although in the general theory of the secular variations of the orbits of the planetary system, $\theta$, $\phi$, \&c., are, as above, referred to one fixed plane (the ecliptic of a certain date), yet in the theory of each particular planet it is the practice, and obviously the convenient one, to refer for such planet the $\theta$, $\phi$ to its own fixed plane (the orbit of the planet at a certain date), the effect of course being that $\phi$, and consequently $p$, $q$, instead of being of the order of the inclinations to the ecliptic, are only of the order of the disturbing forces.

It has occurred to me that the last-mentioned plan should be adhered to throughout; viz., that for each planet $m$, the position of its variable orbit should be determined by $\theta$, the longitude of its node, and $\phi$, the inclination in reference to the appropriate fixed plane (orbit of the planet at a certain date) and origin of longitude therein.

The disturbing functions for the planets $m$ and $m'$ will of course depend not only on $\theta$, $\theta'$, $\phi$, $\phi'$, but on the quantities $\Phi$, $@$, $@'$ which determine the mutual positions of the two fixed planes of reference and origins of longitude therein, these last being however absolute constants not affected by any variation of the elements; so that as regards the variation of the elements the disturbing functions are in fact given as explicit functions of the variable elements $\theta$, $\theta'$, $\phi$, $\phi'$; and where $\phi$, $\phi'$ and therefore also $p$, $q$, $p'$, $q'$ are only of the order of the disturbing forces.

I proceed to work out this idea, for the present considering the development of the Disturbing Function only as far as the first powers of $p$, $q$, \&c.

For comparison with the ordinary theory, observe that in this theory the disturbing function contains only the second powers of the $p$, $q$, \&c., made use of therein; these are in fact of a form such as $P+p$, $Q+q$, \dots\ where $P$, $Q$ are absolute constants and $p$, $q$, \dots\ are the $p$, $q$, \dots\ of the present theory; the ordinary theory gives therefore in the disturbing function a series of terms involving $(P+p)^{2}$, $(P+p)(Q+q)$, \dots\ which I now take account of only as far as the first powers of $p$, $q$, \dots\ viz., they are in effect reduced to $P^{2}+2Pp$, $PQ+Pq+Qp$, \&c.

The present theory is thus not now developed to the extent of giving the $p$, $q$, \dots\ of the ordinary theory in the more complete form as the solutions of a system of simultaneous linear differential equations, but only to the extent of obtaining for these $p$, $q$,\dots\ respectively the terms which are proportional to the time.

I commence with the following subsidiary problem.

Consider a spherical triangle $ABC$ (sides $a$, $b$, $c$, angles $A$, $B$, $C$, as usual), and taking the side $c$ as constant, but the angles $A$ and $B$ as variable, let it be required to find the variations of $C$, $a$, $b$ in terms of variations $dA$, $dB$ and the variable elements $C$, $a$, $b$ themselves.

Although the geometrical proof would be more simple, I give the analytical one, as it may be useful. We have
\[
\cos C = -\cos A\cos B + \sin A\sin B\cos c,
\]
and thence
\begin{multline*}
-\sin C\,dC = (\sin A\cos B + \cos A\sin B\cos c)\,dA\\
+ (\sin B\cos A + \sin A\cos B\cos c)\,dB;
\end{multline*}
that is
\[
\sin C\,dC = \frac{\sin B\sin c}{\sin a}\,dA + \frac{\sin A\sin c}{\sin b}\,dB,
\]
or finally
\[
-dC = \cos b\,dA + \cos a\,dB.
\]

Next
\[
\cos a = \frac{\cos A + \cos B\cos C}{\sin B\sin C},
\]
or, differentiating,
\[
\sin a\,da = \frac{1}{\sin^{2}C\sin c}\{\sin C\cos A\,dA - \cos C\sin A\,dC\} + dB\cos C\sin A\cos a,
\]
that is
\[
\sin C\,da = \frac{\sin A}{\sin c}\{dA\sin b + dB\cos C\sin a\};
\]
and similarly
\[
\sin C\,db = \frac{\sin B}{\sin c}\{dB\sin a + dA\cos C\sin b\}.
\]

Now let the continuous lines represent the orbits of $m$, $m'$ at certain dates, $O$, $O'$ the origins of longitude therein; and the dotted lines the variable orbits of the planets respectively. Write $\angle OCO'=\theta$, $CA=c$, $\angle CAC'=\phi$, $O'C=\theta'$, $CB=c'$, $\angle CBC'=\phi'$, $\angle ACB=\Phi$.

Then, answering to the notation of the lemma, we have $a=\theta'$, $b=\theta$, $C=\Phi$, $dA=\phi$, or say $=\tan\phi$, whence
\[
G'B = a + da, \quad dB = -\phi', \quad = -\tan\phi',
\]
\[
= \theta' + \frac{1}{\sin\Phi}\{\tan\phi\sin\theta - \tan\phi'\cos\Phi\sin\theta'\},
\]
\[
= \theta' + \frac{1}{\sin\Phi}(p - p'\cos\Phi),
\]
\[
C'A = b + db = \theta + \frac{1}{\sin\Phi}(-\tan\phi'\sin\phi' + \tan\phi\cos\Phi\sin\theta),
\]
\[
= \theta + \frac{1}{\sin\Phi}(p' - p\cos\Phi),
\]
\[
C = \Phi + d\Phi = \Phi - \cos\theta\tan\phi + \cos\theta'\tan\phi', \quad = \Phi - q + q'.
\]

Suppose $V$, $v'$ are the longitudes of the planets in their two orbits respectively; that is
\[
V = OA + Am = \theta + \frac{1}{\sin\Phi}(p'-p\cos\Phi) + Am,
\]
\[
v' = O'B + Bm' = \theta' + \frac{1}{\sin\Phi}(p-p'\cos\Phi) + Bm',
\]
whence
\[
Cm = C'A + Am = v - \theta - \frac{1}{\sin\Phi}(p'-p\cos\Phi),
\]
\[
Cm' = G'B + Bm' = v' - \theta' + \frac{1}{\sin\Phi}(p-p'\cos\Phi),
\]
\[
= \Phi - q + q' \text{ say,}
\]
these values are $v-@+x$, $v'-@'+x'$, $\Phi+y$.

Then if $H$ is the angular distance $mvm'$ of the two planets,
\begin{multline*}
\cos H = \cos(v-@+x)\cos(v'-@'+x')\\
+ \sin(v-@+x)\sin(v'-@'+x')\cos(\Phi+y)\\
= \cos(v-@)\cos(v'-@') + \sin(v-@)\sin(v'-@')\cos\Phi\\
+ x[-\sin(v-@)\cos(v'-@') + \cos(v-@)\sin(v'-@')\cos\Phi]\\
+ x'[-\cos(v-@)\sin(v'-@') + \sin(v-@)\cos(v'-@')\cos\Phi]\\
+ y[-\sin(v-@)\sin(v'-@')\sin\Phi],\\
= \cos H_{0} + V \text{ suppose.}
\end{multline*}

The disturbing function for the planet $m$ disturbed by $m'$ is
\[
\frac{m'}{\Delta} = \frac{m'}{\sqrt{r^{2} + r'^{2} - 2rr'\cos H}};
\]
($\Omega=-R$, if $R$ is the disturbing function of the \textit{M\'ecanique C\'eleste});
and the term hereof which involves $V$ is
\[
\frac{d}{d(\cos H)}\frac{1}{\sqrt{r^{2}+r'^{2}-2rr'\cos H}} \cdot V,
\]
where after the differentiation $\cos H$ is replaced by $\cos H_{0}$, viz., this is a linear function of $x$, $x'$, $y$, that is of $p$, $q$, $p'$, $q'$, with coefficients which of course involve the other variable elements and the time; but it will be remembered that $@$, $@'$, $\Phi$ are not variable elements, but are absolute constants.

The variations of $p$ depend upon $\dfrac{d\Omega}{dq}$ and those of $q$ on $-\dfrac{d\Omega}{dp}$, and the quantities $p$, $q$, $p'$, $q'$,\dots\ disappear from these differential coefficients $\dfrac{d\Omega}{dq}$, $\dfrac{d\Omega}{dp}$; that is, disregarding periodic terms, and the variations of the elements, we obtain $\dfrac{dp}{dt}$, $\dfrac{dq}{dt}$ as absolute constants, or reckoning the time from the epoch belonging to the fixed orbit of $m$, we have $p$, $q$ as mere multiples of the time ($p=At$, $q=Bt$, where $A$ and $B$ are constants); agreeing with the statement preceding the investigation.

Observe that the $p$, $q$, as used above, have reference not only to the fixed orbit of $m$, but also to the node thereon of the fixed orbit of $m'$: we may, if we please, write $p=\tan\phi\sin(\theta+\vartheta)$, $q=\tan\phi\cos(\theta+\vartheta)$, that is, $p=g\sin\vartheta+\bar{g}\cos\vartheta$, $Q=g\cos\vartheta-\bar{g}\sin\vartheta$ (or $\bar{g}=p\cos\vartheta-q\sin\vartheta$, $g=p\sin\vartheta+q\cos\vartheta$), and in place of $p$, $q$ introduce into the formulae $\bar{g}$ and $g$, which have reference only to the fixed orbit of $m$, and similarly writing $p'=\tan\phi'\sin(\theta'+\vartheta')$, $q'=\tan\phi'\cos(\theta'+\vartheta')$, instead of $p'$, $q'$ introduce $\bar{g}'$, $g'$ which have reference only to the fixed orbit of $m'$.

I remark that a table for the relative positions of the orbits of the eight Planets for the Epoch 1st January, 1850, is given in Leverrier's \textit{Annales de l'Observ.\ de Paris}, t.\ ii.\ (1856), pp.\ 64--66.


\vspace{2em}

%==============================================================================
% PAPER [485]
%==============================================================================
\newpage
\section*{[485 --- Problems and Solutions]}
\addcontentsline{toc}{section}{[485] Problems and Solutions}

\markboth{Cayley, Collected Papers, Vol.\ VII}{[485] Problems and Solutions}

\begin{center}
\textbf{PROBLEMS AND SOLUTIONS.}
\end{center}

[From the \textit{Mathematical Questions with their Solutions from the Educational Times}, vols.\ V.\ to XII.\ (1866--1869).]

\bigskip

\textbf{[Vol.\ V., January to July, 1866, p.\ 17.]}

\textbf{1791.}\ (Proposed by Professor Cayley.)---Given a quartic curve $U=0$, to find three cubic curves $P=0$, $Q=0$, $R=0$, each meeting the quartic in the same six points $1$, $2$, $3$, $4$, $5$, $6$, and such that $P=0$, $R=0$ may besides meet the quartic in the same three points $a$, $b$, $c$, and that $Q=0$, $R=0$ may besides meet the quartic in the same three points $\alpha$, $\beta$, $\gamma$.

\textbf{[Vol.\ V.\ pp.\ 25, 26.]}

\textit{Note on the Problems in regard to a Conic defined by five Conditions of Intersection.}

I use the word ``intersection'' rather than ``contact'' because it extends to the case of a 1-pointic intersection, which cannot be termed a contact.

The conditions referred to are that the conic shall have with a given curve, at a point given or not given, a 1-pointic intersection, a 2-pointic intersection ($=$ ordinary contact), a 3-pointic intersection, \&c., as the case may be.

It may be noticed that when the point on the curve is a given point, the condition of a $k$-pointic intersection is really only the condition that the conic shall pass through $k$ given points; though from the circumstance that these are consecutive points on a conic, the formulae for a conic passing through $k$ discrete points require material alteration; for instance, in the two questions to find the equation of a conic passing through five given points, and to find the equation of a conic having at a given point of a given curve 5-pointic intersection with the curve, the forms of the solutions are very different from each other.

The foregoing remark shows, however, that it is proper to detach the conditions which relate to intersections at given points; and consequently attending only to the conditions which relate to intersection at an unascertained point (of course the intersections referred to must be at least 2-pointic, for otherwise there is no condition at all) we may consider the conics which pass through four points and satisfy one condition; or which pass through three points and satisfy two conditions; or which pass through two points and satisfy three conditions; or which pass through one point and satisfy four conditions; or which satisfy five conditions.

Considering in particular the last case, let 1 denote that the conic has 2-pointic intersection, 2 that it has 3-pointic intersection, \dots\ 5 that it has 6-pointic intersection with a given curve at an unascertained point.

Then the problems are in the first instance
\[
5; \quad 4,1; \quad 3,2; \quad 3,1,1; \quad 2,2,1; \quad 2,1,1,1; \quad 1,1,1,1,1.
\]
But the intersections may be intersections with the same given curve or with different given curves; and we have thus in all 27 problems, viz.\ these are as given in the following table, where the colons (:) separate those conditions which refer to different curves:

\[
\begin{array}{c|c|c|c|c|c|c|c|c}
\text{No.} & \text{Cond.} & \text{No.} & \text{Cond.} & \text{No.} & \text{Cond.} \\
\hline
1 & 5 & 10 & 3,1:1 & 19 & 3:1:1 \\
2 & 4,1 & 11 & 3:1,1 & 20 & 2:2:1 \\
3 & 3,2 & 12 & 2,2:1 & 21 & 2,1:1:1 \\
4 & 3,1,1 & 13 & 2,1:2 & 22 & 2:1,1:1 \\
5 & 2,2,1 & 14 & 2,1,1:1 & 23 & 1,1,1:1:1 \\
6 & 2,1,1,1 & 15 & 2,1:1,1 & 24 & 1,1:1,1:1 \\
7 & 1,1,1,1,1 & 16 & 2:1,1,1 & 25 & 2:1:1:1 \\
8 & 4:1 & 17 & 1,1,1,1:1 & 26 & 1,1:1:1:1 \\
9 & 3:2 & 18 & 1,1,1:1,1 & 27 & 1:1:1:1:1 \\
\end{array}
\]

Thus Problem 1 is to find a conic having 6-pointic intersection with a given curve; Problem 2 a conic having 5-pointic intersection and also a 2-pointic intersection with a given curve;\dots\ Problem 7 is to find a conic having five 2-pointic intersections with (touching at five distinct points) a given curve;\dots\ Problem 27 is to find a conic having 2-pointic intersection with (touching) each of five given curves.

Or we may in each case take the problem to be merely to find the number of the conics which satisfy the required conditions.

This number is known in Prob.\ 1, for the case of a curve of the order $m$ without singularities, viz.\ the number is $=m(12m-27)$.

It is also known in Problems 25 and 26 in the case where the first curve (that to which the symbol 2, or 1, 1 relates) is a curve without singularities; and it is known in Prob.\ 27, viz.\ if $m$, $n$, $p$, $q$, $r$ be the orders and $M$, $N$, $P$, $Q$, $R$ the classes of the five curves respectively, then the number is $=(M,m)(N',n)(P,p)(Q,q)(R,r)\{1,2,4,4,2,1\}$, that is,
\[
1\cdot MNPQR + 2\sum MNPQr + \&\text{c.}
\]
The number is not, I believe, known in any other of the problems.

In particular, (Prob.\ 7) we do not as yet know the number of the conics which touch a given curve at five points. It would be interesting to obtain this number; but (judging from the analogous question of finding the double tangents of a curve) the problem is probably a very difficult one.

\textbf{[Vol.\ V.\ p.\ 37.]}

\textbf{1857.}\ (Proposed by Professor Cayley.)---If for shortness we put
\begin{align*}
P &= x^{2}+y^{2}+z^{2},\\
Q &= yz'+y'z+zx'+z'x+xy'+x'y,\\
R &= xyz,\\
P_{0} &= a^{2}+b^{2}+c^{2},\\
Q_{0} &= bc'+b'c+ca'+c'a+ab'+a'b,\\
R_{0} &= abc;
\end{align*}
then $(\alpha,\beta,\gamma)$ being arbitrary, show that the cubic curves
\[
\alpha\frac{P}{P_{0}}+\beta\frac{Q}{Q_{0}}+\gamma\frac{R}{R_{0}}=0
\]
pass all of them through the same nine points, lying six of them upon a conic and three of them upon a line; and find the equations of the conic and line, and the coordinates of the nine points of intersection; find also the values of $(\alpha:\beta:\gamma)$ in order that the cubic curve may break up into the conic and line.

\textbf{[Vol.\ V.\ p.\ 37.]}

\textbf{1730.}\ (Proposed by Professor Cayley.)---Show that

(I) the condition in order that the roots $k_{1}$, $k_{2}$, $k_{3}$ of the equation
\[
\gamma k^{3}+(-g-\tfrac12\alpha+\tfrac12\beta+\tfrac12\gamma)k^{2}+(-g-\tfrac12\alpha-\tfrac12\beta+\tfrac12\gamma)k-\alpha=0 \tag{A}
\]
may be connected by a relation of the form
\[
k_{3}(k_{1}-k_{2})-(k_{1}-k_{2})=0, \tag{1}
\]
and (II) the result of the elimination of $a$, $b$, $c$ from the equations
\begin{align*}
a^{2}(b+c) &= -2\alpha, \tag{2}\\
b^{2}(c+a) &= -2\beta, \tag{3}\\
c^{2}(a+b) &= -2\gamma, \tag{4}\\
(b-c)(c-a)(a-b) &= -4g, \tag{5}
\end{align*}
are each
\begin{multline*}
4(\beta-\gamma)(\gamma-\alpha)(\alpha-\beta)g^{3} + 4(-3\alpha^{2}\beta+4\alpha\beta^{2}-2\alpha\beta\gamma)g^{2}\\
+ (\beta-\gamma)(\gamma-\alpha)(\alpha-\beta)g + 2(\beta-\gamma)(\gamma-\alpha)(\alpha-\beta)\gamma = 0. \tag{B}
\end{multline*}

\textbf{[Vol.\ V.\ pp.\ 38, 39.]}

\textbf{1834.}\ (Proposed by Professor Cayley.)---1. It is required to find on a given cubic curve three points $A$, $B$, $C$, such that, writing $x=0$, $y=0$, $z=0$ for the equations of the lines $BC$, $CA$, $AB$ respectively, the cubic curve may be transformable into itself by the inverse substitution $(\alpha x^{-1},\beta y^{-1},\gamma z^{-1})$ in place of $x$, $y$, $z$ respectively, $\alpha$, $\beta$, $\gamma$ being disposable constants.

2. In the cubic curve
\[
ax(y^{2}+z^{2})+by(z^{2}+x^{2})+cz(x^{2}+y^{2})+2lxyz=0
\]
the inverse points $(x,y,z)$ and $(x^{-1},y^{-1},z^{-1})$ are corresponding points (that is, the tangents at these two points meet on the curve).

\textit{Solution by the Proposer, S.\ Roberts, M.A., and others.}

Since the points $A$, $B$, $C$ are on the curve, the equation is of the form
\[
fy^{2}z+gz^{2}x+hxy^{2}+iyz^{2}+jzx^{2}+kxy^{2}+2lxyz=0;
\]
hence this equation must be equivalent to
\[
\frac{fy^{2}z}{\alpha}+\frac{gz^{2}x}{\beta}+\frac{hxy^{2}}{\gamma}+\frac{iyz^{2}}{\alpha}+\frac{jzx^{2}}{\beta}+\frac{kxy^{2}}{\gamma}+2lxyz=0,
\]
or,
\[
\frac{f}{\alpha}y^{2}z+\frac{g}{\beta}z^{2}x+\frac{h}{\gamma}x^{2}y+\frac{i}{\alpha}yz^{2}+\frac{j}{\beta}zx^{2}+\frac{k}{\gamma}xy^{2}+2lxyz=0,
\]
which will be the case if
\[
\frac{f}{\alpha}=f,\ \frac{g}{\beta}=g,\ \frac{h}{\gamma}=h,\ \frac{i}{\alpha}=i,\ \frac{j}{\beta}=j,\ \frac{k}{\gamma}=k,\ 2l=2l.
\]
This implies $fgh=ijk$; and if this condition be satisfied, then $\alpha:\beta:\gamma$ can be determined, viz.\ we have $\alpha:\beta:\gamma=if:ij:kf$, which satisfy the remaining equations, so that the only condition is $fgh=ijk$.

Writing in the equation of the curve $x=0$, we find $fy^{2}z+iyz^{2}=0$, that is, the line $x=0$ meets the curve in the points $(x=0,y=0)$, $(x=0,z=0)$, and $(x=0,fy+iz=0)$.

We have thus on the curve the three points $(x=0,fy+iz=0)$, $(y=0,gz+jx=0)$, $(z=0,hx+ky=0)$, and in virtue of the assumed relation $fgh=ijk$, these three points lie in a line.

Hence the points $A$, $B$, $C$ must be such that $BC$, $GA$, $AB$ respectively meet the curve in points $A'$, $B'$, $C'$, which three points lie in a line; that is, we have a quadrilateral whereof the six angles $A$, $B$, $C$, $A'$, $B'$, $C'$ all lie on the curve.

It is well known that the opposite angles $A$ and $A'$, $B$ and $B'$, $C$ and $C'$ must be corresponding points, that is, points the tangents at which meet on the curve.

And conversely taking $A$, $C$ any two points on the curve, $A'$ a corresponding point to $A$ (any one of the four corresponding points), then $AC$, $A'C$ will meet the curve in the corresponding points $B$, $B'$; and $AB$, $A'B'$ will meet on the curve in a point $C'$ corresponding to $C$, giving the inscribed quadrilateral $(A,B,C,A',B',C')$; the triangle $ABC$ is therefore constructed.

It is to be remarked that the equation $fgh=ijk$ being satisfied, we may without any real loss of generality write $f=j$, $g=k$, $h=i$, and therefore $\alpha=\beta=\gamma$; hence changing the constants we have the theorem: the inverse points $(x,y,z)$, $(x^{-1},y^{-1},z^{-1})$ are corresponding points on the curve
\[
ax(y^{2}+z^{2})+by(z^{2}+x^{2})+cz(x^{2}+y^{2})+2lxyz=0.
\]


\textbf{[Vol.\ V.\ pp.\ 57, 58.]}

\textit{Addition to the Note on the Problems in regard to a Conic defined by five Conditions of Intersection.}

Since writing the Note in question, I have found that a solution of Problem 7 has been given by M.\ De Jonqui\`eres in the paper ``Du Contact des Courbes Planes, \&c.,'' \textit{Nouvelles Annales de Math\'ematiques}, vol.\ iii.\ (1864), pp.\ 218--222: viz.\ the number of conics which touch a curve of the order $n$ in five distinct points is stated to be
\[
\frac{n(n-1)(n-2)(n-3)(n-4)}{1\cdot2\cdot3\cdot4\cdot5}(n^{5}-10n^{3}-37n^{2}-118n+1584n+15).
\]

There are given also the following results; the number of conics which pass through two given points and touch a curve of the order $n$ in three distinct points is
\[
\frac{n(n-1)(n-2)}{1\cdot2\cdot3}(n^{3}+6n^{2}-19n-12),
\]
and the number of conics which pass through a given point and touch a curve of the order $n$ in four distinct points is
\[
\frac{n(n-1)(n-2)(n-3)}{1\cdot2\cdot3\cdot4}(n^{4}+10n^{3}-37n^{2}-118n+282).
\]

These formulae are given without demonstration, and with an expression of doubt as regards their exactness---(``elles ne me paraissent pas pr\'esenter une suffisante garantie d'exactitude'').

\textbf{[Vol.\ V.\ pp.\ 58, 59.]}

\textbf{1876.}\ (Proposed by R.\ Ball, M.A.)---If three of the roots of the equation $(a,b,c,d,e\,\overline{\vrule width 6pt}\,x,1)^{5}=0$ be in arithmetical progression, show that
\[
55296\,H^{3}J - 2304\,aH^{2}I^{2} - 16632\,a^{2}HIJ + 625\,a^{3}I^{3} - 9261\,a^{4}J^{2} = 0,
\]
where
\[
H=ac-b^{2},\quad I=ae-4bd+3c^{2},\quad J=ace+2bcd-ad^{2}-b^{2}e-c^{3}.
\]

\textit{Solution by Professor Cayley.}

Write $(a,b,c,d,e\,\overline{\vrule width 6pt}\,x,1)^{5}=a(x-\alpha)(x-\beta)(x-\gamma)(x-\delta)$; then putting for a moment $\beta+\gamma+\delta=p$, $\beta\gamma+\beta\delta+\gamma\delta=q$, $\beta\gamma\delta=r$, and forming the equation
\[
(\beta+\gamma-2\delta)(\beta+\delta-2\gamma)(\gamma+\delta-2\beta)=0,
\]
this is easily reduced to $-2p^{3}+9pq-27r=0$.

But we have
\[
p=\frac{4b}{a}+\alpha,\quad q=\frac{6c}{a}+\frac{4b}{a}\alpha+\alpha^{2},\quad r=\frac{4d}{a}+\frac{6c}{a}\alpha+\frac{4b}{a}\alpha^{2}+\alpha^{3},
\]
and hence
\[
p=\frac{4b}{a}+\alpha,\quad q=\frac{6c}{a}+\frac{4b}{a}\alpha+\alpha^{2},\quad r=\frac{4d}{a}+\frac{6c}{a}\alpha+\frac{4b}{a}\alpha^{2}+\alpha^{3}.
\]

Substituting these values of $p$, $q$, $r$, the foregoing equation becomes, after all reductions,
\[
(20a^{3},\ 20a^{2}b,\ -4ab^{2}+36a^{2}c,\ 128b^{3}-216abc+108a^{2}d\,\overline{\vrule width 6pt}\,\alpha,1)^{3}=0,
\]
and from this and the equation $(a,b,c,d,e\,\overline{\vrule width 6pt}\,\alpha,1)^{5}=0$, eliminating $\alpha$, we should find the condition for three roots in arithmetical progression.

But it appears from the theory of invariants that the result of the elimination may be obtained by writing $b=0$, and expressing the result so obtained in terms of $a$, $H$, $I$, $J$.

Hence, writing in the two equations $b=0$, the first equation contains the factor $4a^{2}$ and throwing this out, the equations become
\[
5a\alpha^{2}+27c\alpha+27d=0,\quad a\alpha^{3}+3c\alpha^{2}+3d\alpha+e=0;
\]
or multiplying the first by $\alpha$ and reducing by means of the second, the two equations become
\[
5a\alpha^{2}+27c\alpha+27d=0,\quad 3c\alpha^{2}-12d\alpha+5e=0.
\]

The result is of the degree 5 in the coefficients, but in order to avoid fractions in the final result it is proper to multiply it by $a^{4}$; it then becomes
\[
625a^{4}e^{2}-4050a^{3}cde+6561a^{2}d^{2}c-1890a^{3}ced^{2}+13122a^{2}c^{2}d^{2}+9261a^{4}d^{4}=0.
\]

But writing as above $b=0$, we have
\[
H=a,\quad c=-\frac{3H}{a},\quad d=-\frac{3H^{2}}{a^{2}},\quad I=3H^{2},\quad e=\frac{HI}{a^{2}}+\frac{iH^{3}}{a^{3}}.
\]

and substituting these values, the result is found to contain the terms $\frac{1}{a}$, $\frac{1}{a^{2}}$ with coefficients which vanish; viz.\ the coefficient of the first of these terms is $+16875+24300+6561+7560+18792-74088=0$; and the coefficient of the second of the two terms is $-16875-36450-19683-75168+148176=0$.

The remaining terms give
\begin{align*}
+625 &= +625\,a^{4}I^{2}\\
-5625-4050-1890+9261 &= -2304\,a^{3}H^{2}I^{2}\\
+1890-18522 &= -16632\,a^{2}HIJ\\
-18792+74088 &= +55296\,H^{3}J\\
+9261 &= +9261\,a^{4}J^{2}
\end{align*}
which is the required result; a more convenient form of writing it is
\[
(55296\,J,\ -768\,I^{2},\ -5544\,IJ,\ 625\,I^{2}+9261\,J^{2}\,\overline{\vrule width 6pt}\,H,\ a)^{2}=0.
\]

\textit{Remark.} If $I$ and $J$ denote as above the two invariants of the form $U=(a,b,c,d,e\,\overline{\vrule width 6pt}\,x,1)^{5}$, and if we now use $H$ to denote the Hessian of the form, viz.
\[
H=(ac-b^{2},\ \tfrac12(ad-bc),\ \tfrac16(ae+2bd-3c^{2}),\ \tfrac12(be-cd),\ ce-d^{2}\,\overline{\vrule width 6pt}\,x,1)^{2},
\]
then it appears by the theory of invariants that the equation of the twelfth order
\[
(55296\,J,\ -768\,I^{2},\ -5544\,IJ,\ 625\,I^{2}+9261\,J^{2}\,\overline{\vrule width 6pt}\,H,\ U)^{2}=0,
\]
is such that each of its roots forms with some three of the roots of the equation $U=0$ a harmonic progression; viz.\ if the three roots are $\beta$, $\gamma$, $\delta$, then we have
\[
\frac{2}{x-\gamma}=\frac{1}{x-\beta}+\frac{1}{x-\delta},\quad \frac{2\beta\delta-(\beta+\delta)\gamma}{x-\gamma}=\frac{\beta+\delta-2\gamma}{x-\beta},
\]
so that the roots of the equation of the twelfth order are the twelve values of the last-mentioned function of three roots.

\textbf{[Vol.\ V.\ pp.\ 65, 66.]}

\textit{On the Problems in regard to a Conic defined by five Conditions of Intersection.}

There has been recently published in the \textit{Comptes Rendus} (t.\ LXII.\ pp.\ 177--183, January, 1866) an extract of a memoir ``Additions to the Theory of Conics,'' by M.\ H.\ G.\ Zeuthen (of Copenhagen).

The extract gives the solutions of fourteen problems, with a brief indication of the method employed for obtaining them. Of these problems, four relate to intersections at given points, the remaining ten are included among the twenty-seven problems enumerated in my Note on this subject in the January Number of the Educational Times (Reprint, vol.\ v., p.\ 25); but two of these ten are the problems 25 and 26 which are in my Note stated to have been solved; there are, consequently, of the twenty-seven problems, in all twelve which are solved: viz.\ these are

\begin{center}
\begin{tabular}{c|cccccccccccc}
No.\ of Prob. & 1 & 8 & 10 & 12 & 14 & 17 & 19 & 21 & 23 & 25 & 26 & 27 \\
Zeuthen's No. & -- & 14 & 13 & 11 & 8 & 3 & 12 & 7 & 2 & 6 & 1 & -- \\
\end{tabular}
\end{center}

where it is to be observed that Zeuthen's solutions apply to the case of a curve of a given order with given numbers of double points and cusps.

The problems 25 and 26 had been previously solved only in the case of a curve without singularities. As to Prob.\ 27, the solution mentioned in my former Note is in fact applicable to the general case.

The solution for Prob.\ 1 may also be extended to this general case, viz.\ for a curve of the order $m$ with $\delta$ double points and $\kappa$ cusps the required number is $=m(12m-27)-24\delta-27\kappa$; or, if $n$ be the class, then this number is $=12n-15m+9\kappa$; so that all the twelve problems are solved in the general case.

The results obtained by M.\ de Jonqui\`eres, as stated in my Note in the March Number (Reprint, vol.\ v., p.\ 57), seem to be all of them erroneous.

In fact, for the number of conics passing through two given points and touching a curve of the order $m$ in three distinct points (which is a particular case of Prob.\ 23), Zeuthen's formula applied to a curve without singularities gives this instead of the value which is
\[
=\tfrac16 m(m-2)(m^{3}+6m^{2}-11m^{2}-49m+108)
\]
\[
=\tfrac16 m(m-1)(m-2)(m^{2}+6m^{2}-19m-12)
\]
\[
=\tfrac16 m(m-2)(m^{4}+6m^{3}-25m^{2}+11m+12);
\]
and I have by my own investigation verified Zeuthen's Number.

So for the number of conics through a given point and touching a curve of the order $m$ in four distinct points (which is a particular case of Prob.\ 17), Zeuthen's formula applied to a curve without singularities gives
\[
\tfrac1{24}m(m-2)(m-3)(m^{4}+10m^{3}-37m^{2}-118m+282);
\]
I have obtained for Prob.\ 2 a solution which I believe to be accurate; viz.\ the number of the conics $(4,1)$, (that is, the conics which have with a given curve a 5-pointic intersection and also a 2-pointic intersection, or ordinary contact), is
\[
=10n^{2}+10mn-20m^{2}-130n+140m+10\kappa(m+n-9)-4[(n-3)\iota+(\kappa+(m-3)\iota)]
\]
where $\iota$ (the number of inflexions) is $=3n-3m+\kappa$, but I prefer to retain the foregoing form, without effecting the substitution.


\textbf{[Vol.\ V.\ pp.\ 88, 89.]}

\textbf{1890.}\ (Proposed by Professor Cayley.)---Find the equation of a conic passing through three given points and having double contact with a given conic.

\textit{Solution by the Proposer.}

Let the given points be the angles of the triangle $(x=0,y=0,z=0)$, and let the equation of the given conic be $U=(a,b,c,f,g,h\,\overline{\vrule width 6pt}\,x,y,z)^{2}=0$; then the equation of the required conic is
\[
U-(x\sqrt{a}+y\sqrt{b}+z\sqrt{c})^{2}=0,
\]
for this is a conic having double contact with the conic $U=0$, and, since the terms in $(x^{2},y^{2},z^{2})$ each vanish, it is also a conic passing through the given points.

It is clear that there are four conics satisfying the conditions of the Problem, viz.\ putting for shortness
\begin{align*}
P &= x\sqrt{a}+y\sqrt{b}+z\sqrt{c},\\
P_{1} &= -x\sqrt{a}+y\sqrt{b}+z\sqrt{c},\\
P_{2} &= x\sqrt{a}-y\sqrt{b}+z\sqrt{c},\\
P_{3} &= x\sqrt{a}+y\sqrt{b}-z\sqrt{c},
\end{align*}
the four conics are $U-P^{2}=0$, $U-P_{1}^{2}=0$, $U-P_{2}^{2}=0$, $U-P_{3}^{2}=0$.

It may be remarked that the conics $P$, $P_{1}$ have a fourth intersection lying on the line $y\sqrt{b}-z\sqrt{c}=0$, and the conics $P_{2}$, $P_{3}$ a fourth intersection lying on the line $y\sqrt{b}-z\sqrt{c}$; which two lines are harmonics in regard to the lines $y=0$, $z=0$.

Similarly the conics $P_{1}$, $P_{2}$ have a fourth intersection on the line $x\sqrt{a}+z\sqrt{c}=0$, and the conics $P$, $P_{3}$ a fourth intersection on the line $x\sqrt{a}-z\sqrt{c}=0$; which two lines are harmonics in regard to the lines $z=0$, $x=0$.

And the conics $P_{1}$, $P_{3}$ have a fourth intersection on the line $x\sqrt{a}+y\sqrt{b}=0$, and the conics $P$, $P_{2}$ a fourth intersection on the line $x\sqrt{a}-y\sqrt{b}=0$; which two lines are harmonics in regard to the lines $x=0$, $y=0$.

\textbf{[Vol.\ V.\ pp.\ 99, 100.]}

\textbf{1554.}\ (Proposed by Professor Cayley.)---Show that, in the ellipse and its circles of maximum and minimum curvature respectively, the semi-ordinates through the focus of the ellipse are
\[
\text{For the circle of maximum curvature } y_{1}=a(1-e)(1+2e)^{\frac12},
\]
\[
\text{For the ellipse } y_{2}=a(1-e^{2})^{\frac12},
\]
\[
\text{For the circle of minimum curvature } y_{3}=\frac{a\{(1-e^{2}+e^{4})^{\frac32}-e^{3}\}}{1-e^{2}},
\]
and that these values are in the order of increasing magnitude.

\textbf{[Vol.\ VI., July to December, 1866, pp.\ 18, 19.]}

\textbf{1931.}\ (Proposed by Professor Cayley.)---Find the stationary tangents (or tangents at the inflexions) of the nodal cubic
\[
x(y-z)^{2}+y(z-x)^{2}+z(x-y)^{2}=0.
\]

\textit{Solution by the Proposer.}

The equation may be transformed into the form
\[
(-3x+y+z)^{3}+(x-3y+z)^{3}+(x+y-3z)^{3}=0,
\]
and it thence follows immediately that the stationary tangents are the lines
\[
-3x+y+z=0,\quad x-3y+z=0,\quad x+y-3z=0,
\]
respectively, and that the three points of contact, or inflexions, are the intersections of these lines with the line $x+y+z=0$.

\textbf{[Vol.\ VI.\ pp.\ 33--36.]}

\textbf{1935.}\ (Proposed by Professor Cayley.)---1. A circular cubic has four concyclic foci. A characteristic property of 4 concyclic foci is that taking any three of them $A$, $B$, $C$, the distances of a point $P$ of the curve from these three foci are connected by a linear relation $\lambda\cdot AP+\mu\cdot BP+\nu\cdot CP=0$, where $\lambda+\mu+\nu=0$, or if as is more convenient the distances are considered as $\pm$, then where $\lambda+\mu+\nu=0$.

A circular cubic may be determined so as to satisfy 7 conditions; having a focus at a given point is 2 conditions; hence a circular cubic may be determined so as to pass through three given points, and to have as foci two given points.

2. A Cartesian, or bicircular cuspidal quartic (that is a quartic having a cusp at each of the circular points at infinity) has nine foci, but of these there are three which lie in a line with the centre of the Cartesian (or intersection of the cuspidal tangents), and which are preeminently the foci of the Cartesian.

We may, therefore, say that the Cartesian has three foci, which foci lie in a line, the axis of the Cartesian.

A Cartesian may be determined to satisfy 6 conditions; having a focus at a given point is 2 conditions; but having for foci three given points on a line is 5 conditions; and hence a Cartesian may be found having for foci three given points on a line, and passing through a given point; there are in fact two such Cartesians, intersecting at right angles at the given point.

3. The theorem at first sight appears impossible; for take any three points $F$, $G$, $H$ in a line and any other point $A$; then, as just remarked, there are, having $F$, $G$, $H$ for foci and passing through $A$, two Cartesians.

And we may draw through $F$, $G$, $H$, and with $A$ for focus, a circular cubic depending upon two arbitrary parameters; the position of a second focus of the circular cubic is (on account of the two arbitrary parameters) prima facie indeterminate; and this is confirmed by the remark that the circular cubic can actually be so determined as to have for focus an arbitrary point $B$; and yet the theorem is true.

4. For, $A$, $B$ being two of the four concyclic foci, taking $C$ a variable focus, the linear relations connecting the distances $CA$, $CB$, $CF$, $CG$, $CH$ are four linear relations; but these are not independent, for, considering $C$ as belonging to the circular cubic, we have the two relations $CA$, $CB$ (viz.\ the equations which determine the position of $C$ as a focus of the circular cubic of which $A$, $B$ are two given foci); and considering $C$ as belonging to the Cartesian, the two relations $CF$, $CG$, $CH$ (viz.\ the equations which determine the position of $C$ as a focus of the Cartesian of which $F$, $G$, $H$ are three given foci); the eight linear relations are equivalent to seven independent relations; whence, if $F$, $G$, $H$ be three foci in a line of the Cartesian through $A$, the fourth focus $C$ of the circular cubic is determined; that is, it is a determinate point not depending on the particular circular cubic through $F$, $G$, $H$, $A$, which is considered.

Or, what is the same thing, the four distances $CA$, $CB$, $CG$, $CH$ are connected by two linear relations; $CF$, $CG$, $CH$ are linear functions of $CF$; and $CA$, $CB$ are linear functions of $CF$; substituting for $CG$, $CH$ their values in terms of $CF$, and for $CA$, $CB$ their values in terms of $CF$, we should have a relation connecting $CF$ with itself; the condition of identity gives the two linear relations between $CA$, $CB$, $CG$, $CH$.

5. More definitely; if $F$, $G$, $H$ are three points in a line, and $A$, $B$ are two concyclic foci, then the distances of the fourth focus $C$ from $F$, $G$, $H$ are connected by two linear relations; viz.\ considering $C$ as belonging to the circular cubic of which $A$, $B$ are two given foci, we have the two relations which express that $C$ is a focus of this cubic; these are equivalent to two equations for the determination of the position of the point $C$; these give $CG$, $CH$ as linear functions of $CF$; the distances $CF$, $CG$, $CH$ of the point $C$ from the points $F$, $G$, $H$ in the line $FOH$ are connected by a quadratic equation, and hence substituting for $CG$, $CH$ their values in terms of $CF$, we have a quadratic equation for $CF$; as the given conditions are satisfied when $C$ coincides with $A$ or with $B$, the roots of this equation are $CF=AF$ and $CF=BF$.

But if $CF=AF$, then the linear relations give $CG=AG$ and $CH=AH$, that is, $C$ is a point opposite to $A$ in regard to the line $FOH$.

And similarly if $CF=BF$, then $C$ is a point opposite to $B$ in regard to the line $FOH$.

But $C$ being opposite to $A$ or $B$, the fourth concyclic focus $D$ will be opposite to $B$ or $A$; that is, the pairs $A$, $B$ and $C$, $D$ of concyclic foci lie symmetrically on opposite sides of the line $FOH$; this of course implies that the four points lie on a circle.


6. Taking $F=0$ as the equation of the line $FOH$, $x^{2}+y^{2}-1=0$ as the equation of the circle through the four points $A$, $B$, $C$, $D$, then these lie on a proper cubic $(x^{2}+y^{2}+1)x+lx^{2}+ny^{2}=0$ (not passing through the points $F$, $G$, $H$) and the four foci are given as the intersections with the circle $x^{2}+y^{2}-1=0$ of the pair of lines $x^{2}-2nx-nl=0$.

But if we attempt to describe with the same four foci a cubic $(x^{2}+y^{2}+1)y+l'x^{2}+2m'xy+n'y^{2}=0$, then the foci are given as the intersections with the circle $x^{2}+y^{2}-1=0$ of the conic $y^{2}+2m'x-2l'y+m'^{2}-n'l'=0$.

In order that these may coincide with the points $(A,B,C,D)$ we must have
\[
(x^{2}-2nx-nl)+(y^{2}+2m'x-2l'y+m'^{2}-n'l')=x^{2}+y^{2}-1;
\]
that is $m'=n$, $l'=0$, $-nl+n^{2}-n'l'=-1$.

The last equation is $n'l'=n^{2}+1-nl$, which, assuming that $nl$ is not equal to $n^{2}+1$, (in this case the cubic $(x^{2}+y^{2}+1)x+lx^{2}+ny^{2}=0$ would reduce itself to the line and conic $(x+n)(x^{2}+y^{2}+\frac{l}{n})=0$), since $l'=0$, gives $n'=\infty$, and therefore the cubic $(x^{2}+y^{2}+1)y+l'x^{2}+\dots=0$ does not exist.

7. Considering the points $F$, $G$, $H$ on a line and the point $A$ as given, it has been seen that there are two Cartesians through $A$ with the foci $F$, $G$, $H$; and the theorem asserts that in the circular cubics through $F$, $G$, $H$ with the focus $A$, the foci concyclic with $A$ lie on one or other of the two Cartesians: there are consequently through $F$, $G$, $H$ with the focus $A$ two systems of circular cubics corresponding to the two Cartesians respectively, each system depending upon two arbitrary parameters.

But if we attend only to one of the two Cartesians and to the corresponding system of cubics, then the Cartesian passes through the four foci of each cubic, and if (instead of taking as given the points $F$, $G$, $H$ and the focus $A$) we take as given the four concyclic foci $A$, $B$, $C$, $D$ of a cubic, the theorem asserts that we have through $A$, $B$, $C$, $D$ a Cartesian depending on two arbitrary parameters (or having for its axis an arbitrary line), and such that the foci of the Cartesian are the points of intersection $F$, $G$, $H$ of its axis with the cubic.

And I proceed to the proof of the theorem in this form.

8. The equation of a circular cubic having four foci on the circle $x^{2}+y^{2}-1=0$ is $(x^{2}+y^{2}+1)(Px+Qy)+lx^{2}+2mxy+ny^{2}=0$; and this being so, the four foci are the intersections of the circle with the conic $(Qx-Py)^{2}+2(-nP+mQ)x+2(mP-lQ)y+m^{2}-nl=0$.

9. The general equation of a Cartesian is $(x^{2}+y^{2}+2Ax+2By+C)^{2}+2Dx+2Ey+F=0$, and by assuming for $A$, $B$, $C$, $D$, $E$, $F$, the following values which contain the two arbitrary parameters $\alpha$ and $\theta$, viz.\ by writing
\begin{align*}
2A &= \theta Q,\ 2B=-\theta P,\ C=\alpha-1,\\
D &= -n\theta^{2}P+(m\theta^{2}-\alpha\theta)Q,\\
E &= (m\theta^{2}+\alpha\theta)P-l\theta^{2}Q,\\
F &= -\alpha^{2}+\theta^{2}(m^{2}-nl),
\end{align*}
we have the equation of a system (the selected one out of two systems) of Cartesians through the four foci; in fact, substituting the foregoing values, the equation of the Cartesian is
\begin{multline*}
\{x^{2}+y^{2}+\theta(Qx-Py)+\alpha-1\}^{2}-2\alpha\theta(Qx-Py)\\
+2\theta(-nP+mQ)x+2\theta(mP-lQ)y-\alpha^{2}+\theta^{2}(m^{2}-nl)=0,
\end{multline*}
and writing herein $x^{2}+y^{2}+1=0$, \&c., the equation reduces itself to the form $(x^{2}+y^{2}+1)(Px+Qy)+lx^{2}+2mxy+ny^{2}=0$.

10. To find where the line in question meets the cubic $(x^{2}+y^{2}+1)(Px+Qy)+lx^{2}+2mxy+ny^{2}=0$, writing in this equation $x=-A+D\Pi$, $y=-B+E\Pi$, we have for the determination of $\Pi$ the equation
\begin{multline*}
\{A^{2}+B^{2}+1-2(AD+BE)\Pi+(D^{2}+E^{2})\Pi^{2}\}\\
\times\{-AP-BQ+(DP+EQ)\Pi\}+(l,m,n\,\overline{\vrule width 6pt}\,-A+D\Pi,-B+E\Pi)^{2}=0,
\end{multline*}
or observing that we have $AP+BQ=0$, this equation becomes
\begin{multline*}
(D^{2}+E^{2})(DP+EQ)\Pi^{3}\\
+\{-2(AD+BE)(DP+EQ)+lD^{2}+2mDE+nE^{2}\}\Pi^{2}\\
+\{(A^{2}+B^{2}+1)(DP+EQ)-2lAD-2m(AE+BD)-2nBE\}\Pi\\
+\{lA^{2}+2mAB+nB^{2}\}=0.
\end{multline*}

11. Substituting for $A$, $B$, $D$, $E$ their values in terms of $(P,Q,\alpha,\theta)$, we find
\begin{align*}
DP+EQ &= -\theta^{2}(nP^{2}-2mPQ+lQ^{2}),\\
lA^{2}+2mAB+nB^{2} &= \frac{\alpha^{2}}{4\theta^{2}}(nP^{2}-2mPQ+lQ^{2}),\\
lAD+m(AE+BD)+nBE &= -\frac{\alpha^{2}}{4}(nP^{2}-2mPQ+lQ^{2}),\\
lD^{2}+2mDE+nE^{2} &= \{(nl-m^{2})\theta^{2}+n^{2}\theta^{2}\}(nP^{2}-2mPQ+lQ^{2}),
\end{align*}
and substituting these values in the equation for $\Pi$, the whole equation divides by $\theta^{2}(nP^{2}-2mPQ+lQ^{2})$, and it then becomes
\[
4(D^{2}+E^{2})\Pi^{3}+4\{-2(AD+BE)-(nl-m^{2})\theta^{2}-\alpha^{2}\}\Pi^{2}+4\{A^{2}+B^{2}+1-\alpha\}\Pi-1=0,
\]
or, putting for shortness
\begin{align*}
C' &= C-A^{2}-B^{2}=\alpha-1-A^{2}-B^{2},\\
F' &= F-2(AD+BE)=-\alpha^{2}-\theta^{2}(nl-m^{2})-2(AD+BE),
\end{align*}
the equation in $\Pi$ is
\[
4(D^{2}+E^{2})\Pi^{3}+4F'\Pi^{2}-4C'\Pi-1=0,
\]
so that, $\Pi$ satisfying this equation, the intersections of the axis with the cubic are given by $x=-A+D\Pi$, $y=-B+E\Pi$.

12. The equation of the Cartesian, writing therein $x+A=u$ and $y+B=v$, and attending to the values of $C'$ and $F'$, is $(u^{2}+v^{2}+C')^{2}+2Du+2Ev+F'=0$.

And to find the foci, writing in this equation $u+\rho$, $v+i\rho$ in place of $u$, $v$, we find
\begin{multline*}
(u^{2}+v^{2}+C'+2(u+vi)\rho)^{2}+2(D+Ei)\rho+2Du+2Ev+F'=0,\\
\text{that is }(u^{2}+v^{2}+C')^{2}+2Du+2Ev+F'\\
+\{2(u+vi)(u^{2}+v^{2}+C')+D+Ei\}2\rho+4(u+vi)^{2}\rho^{2}=0.
\end{multline*}

Expressing that this equation in $\rho$ has two equal roots, we find
\[
4(u+vi)^{2}[(u^{2}+v^{2}+C')^{2}+2Du+2Ev+F']-\{2(u+vi)(u^{2}+v^{2}+C')+D+Ei\}^{2}=0,
\]
that is
\[
4(2Du+2Ev+F')(u+vi)^{2}-4(u^{2}+v^{2}+C')(u+vi)(D+Ei)-(D-Ei)^{2}=0,
\]
which equation is in fact the equation of the three tangents from one of the circular points at infinity.

Writing it under the form $U+Vi=0$, the nine foci of the Cartesian are given as the intersections of the two cubics $U=0$, $V=0$.

But of these nine points, three, the foci that we are concerned with, lie on the axis, or line $Eu-Dv=0$; in fact, we have
\begin{align*}
U &= 4(u^{2}-v^{2})(2Du+2Ev+F')-4(uD-vE)(u^{2}+v^{2}+C')-(D^{2}-E^{2}),\\
V &= 8uv(2Du+2Ev+F')-4(uE+vD)(u^{2}+v^{2}+C')-2DE;
\end{align*}
and hence
\[
2DEU-(D^{2}-E^{2})V=(Eu-Dv)\{8(Eu+Ev)(2Du+2Ev+F')-4(D^{2}+E^{2})(u^{2}+v^{2}+C')\}=0,
\]
which shows that the nine points lie three of them on the line $Eu-Dv=0$, and the remaining six on the conic $2(Du+Ev)(2Du+2Ev+F')-(D^{2}+E^{2})(u^{2}+v^{2}+C')=0$.

13. We have thus the three foci given as the intersections of the axis $Eu-Dv=0$, with the cubic
\[
U=4(u^{2}-v^{2})(2Du+2Ev+F')-4(uD-vE)(u^{2}+v^{2}+C')-(D^{2}-E^{2})=0;
\]
or, writing in this last equation $u=DD_{0}$, $v=ED_{0}$, that is $x=-A+DD_{0}$, $y=-B+ED_{0}$, we have $u^{2}-v^{2}=(D^{2}-E^{2})\Pi_{0}$, $uD-vE=(D^{2}-E^{2})\Pi_{0}$. The whole equation divides by $(D^{2}-E^{2})$, and omitting this factor, it is
\[
4\Pi_{0}\{2(D^{2}+E^{2})\Pi_{0}+F'\}-4\Pi_{0}\{(D^{2}+E^{2})\Pi_{0}+C'\}-1=0,
\]
that is
\[
4(D^{2}+E^{2})\Pi_{0}^{3}+4F'\Pi_{0}^{2}-4C'\Pi_{0}-1=0,
\]
the same equation as the equation in $\Pi$ before obtained; that is the intersections of the cubic with the axis are the three foci of the Cartesian.


\textbf{[Vol.\ VI.\ pp.\ 57--59.]}

\textbf{1949.}\ (Proposed by Professor Cayley.)---Find the conic of five-pointic intersection at any point of the cuspidal cubic $y^{2}=x^{2}z$.

\textit{Solution by the Proposer.}

The equation $y^{2}=x^{2}z$ is satisfied by the values $x:y:z=1:\theta:\theta^{2}$; and conversely, to any given value of the parameter $\theta$ there corresponds a point on the cubic $y^{2}=x^{2}z$.

Consider the five points corresponding to the values $(\theta_{1},\theta_{2},\theta_{3},\theta_{4},\theta_{5})$ respectively; the equation of the conic through these five points is
\[
(a,b,c,f,g,h\,\overline{\vrule width 6pt}\,x,y,z)^{2}=0,
\]
where, writing for shortness $p_{1}=\theta_{1}+\theta_{2}+\dots$ (the sum of the roots one at a time), $p_{2}=$ sum two at a time, \&c., we have
\begin{align*}
a &= p_{2}^{2}-p_{1}p_{3}+p_{4},\\
b &= -p_{5}+p_{1}p_{4}-p_{2}p_{3}+p_{1}^{2}p_{3}-p_{1}p_{2}^{2},\\
c &= p_{3}^{2}-p_{2}p_{4}+p_{1}p_{5},\\
f &= -p_{1}p_{5}+p_{2}p_{4}-p_{3}^{2},\\
g &= p_{5},\\
h &= p_{4}-p_{1}p_{3}+p_{2}^{2}.
\end{align*}

Hence if $\phi$ be the parameter of the given point, writing $\theta_{1}=\theta_{2}=\theta_{3}=\theta_{4}=\theta_{5}=\phi$, we have
\[
(a,b,c,f,g,h\,\overline{\vrule width 6pt}\,1,\theta,\theta^{2})^{2}=(\phi-\theta)^{5}(\phi+5\theta)
\]
\[
=(1,0,-15,+40,-45,+24,-5\,\overline{\vrule width 6pt}\,\phi,0)^{2}(1,\theta)^{5},
\]
where the left-hand side is $a+b\theta+c\theta^{2}+f\theta^{3}+g\theta^{4}+h\theta^{5}=(c,0,f,g,b,h,a\,\overline{\vrule width 6pt}\,\theta,1)^{2}$, that is we have
\[
c=1,\ f=-15\phi^{2},\ g=40\phi^{3},\ h=-45\phi^{4},\ b=24\phi^{5},\ a=-5\phi^{6},
\]
and the equation of the conic of five-pointic intersection therefore is
\[
(-5\phi^{6},\ 24\phi^{5},\ 1,\ -15\phi^{2},\ 40\phi^{3},\ -45\phi^{4}\,\overline{\vrule width 6pt}\,x,y,z)^{2}=0,
\]
or, what is the same thing,
\[
-5\phi^{6}x^{2}-45\phi^{4}y^{2}+z^{2}-30\phi^{3}y^{2}z+40\phi^{3}z^{2}x+24\phi^{5}xy=0,
\]
which is the required result.

\textit{Note.} The condition in order that any six points $(\theta_{1},\theta_{2},\theta_{3},\theta_{4},\theta_{5},\theta_{6})$ of the cubic $y^{2}=x^{2}z$ may lie on a conic, is
\[
p_{6}-p_{1}p_{5}+p_{2}p_{4}-p_{3}^{2}=0.
\]

\textbf{[Vol.\ VI.\ p.\ 65.]}

\textbf{1872.}\ (Proposed by Professor Cayley.)---Show that the surfaces $xyz=1$, $yz+zx+xy+x+y+z+2=0$, intersect in two distinct cubic curves; and find the equations of the cubic cones which have their vertices at the origin and pass through these curves respectively.

\textbf{[Vol.\ VI.\ pp.\ 67--69.]}

\textbf{1969.}\ (Proposed by Professor Sylvester.)---In two given great circles of a sphere intersecting at $O$ are taken respectively two points $P$ and $Q$, the arc joining which is of given length: prove that $S$, $H$ two fixed points, and $M$ a fixed line, in a plane may be found such that, for all positions of the arc $PQ$, a point $M$ in the fixed line may be found satisfying the equations
\[
SM \mp HM = \sin OP,\quad SM + HM = \sin OQ.
\]

\textit{Solution by Professor Cayley.}

1. In the spherical triangle $OPQ$, whereof the sides $OP$, $OQ$, $PQ$ are $\theta$, $\phi$, $\beta$ and the angle $O$ is $\alpha$, the relation between these quantities is
\[
\cos\alpha = \frac{\cos\beta-\cos\theta\cos\phi}{\sin\theta\sin\phi};
\]
hence treating $\alpha$, $\beta$ as constants, and $\theta$, $\phi$ as variable angles connected by the foregoing equation, it is required to show that we can find two fixed points $S$, $H$ and a fixed line, such that taking $M$ a variable point in this line and writing $SM=r$, $HM=s$, the relation between $r$ and $s$ (or equation of the fixed line in terms of $r$, $s$ as coordinates of a point thereof) is obtained by substituting in the foregoing equation for $\theta$ and $\phi$ the values given by the two equations
\[
\sin\theta=(r+s),\quad \sin\phi=(r-s),
\]
or as, for the sake of homogeneity, it will be more convenient to write these equations,
\[
m\sin\theta=(r+s),\quad m\sin\phi=(r-s).
\]

2. Suppose that the perpendicular distances of $S$, $H$ from the fixed line are $a$ and $b$, and that the distance between the feet of the two perpendiculars is $2c$, then if $x$ denote the distance of the point $M$ from the midway point between the feet of the two perpendiculars, we have
\[
r=\sqrt{(c+x)^{2}+a^{2}},\quad s=\sqrt{(c-x)^{2}+b^{2}},
\]
and $(a,b,c)$ being properly determined, the elimination of $x$ from these equations should give between $(r,s)$ a relation equivalent to that obtained by the elimination of $(\theta,\phi)$ from the before-mentioned equations.

Or, what is the same thing, the elimination of $(r,s,x)$ from the equations
\[
m\sin\theta=r+s,\quad m\sin\phi=r-s,\quad r=\sqrt{(c+x)^{2}+a^{2}},\quad s=\sqrt{(c-x)^{2}+b^{2}}
\]
should give between $(\theta,\phi)$ the relation
\[
\cos\alpha=\frac{\cos\beta-\cos\theta\cos\phi}{\sin\theta\sin\phi};
\]
that is, the last-mentioned equation should be obtained by the elimination of $x$ from the equations
\[
m(\sin\theta+\sin\phi)=2\sqrt{(c+x)^{2}+a^{2}},\quad m(\sin\theta-\sin\phi)=2\sqrt{(c-x)^{2}+b^{2}}.
\]

3. The equation in $(\theta,\phi)$ may be written
\[
\cos\beta-\cos\alpha\sin\theta\sin\phi=\cos\theta\cos\phi,
\]
or, squaring and reducing,
\[
\sin^{2}\theta+\sin^{2}\phi=\sin^{2}\beta+2\cos\alpha\cos\beta\sin\theta\sin\phi+\sin^{2}\alpha\sin^{2}\theta\sin^{2}\phi,
\]
that is
\[
\sin^{2}\theta+\sin^{2}\phi=\frac{1-\cos^{2}\alpha-\cos^{2}\beta}{\sin^{2}\alpha}+\left(\sin\alpha\sin\theta\sin\phi+\frac{\cos\alpha\cos\beta}{\sin\alpha}\right)^{2}.
\]

\vspace{2em}

\begin{center}
\rule{0.5\textwidth}{0.4pt}
\end{center}

\vspace{1em}

\begin{center}
\textbf{End of Pages 515--564}
\end{center}
\clearpage


% =========================================================
% Chunk: cayley_vol07_pages_601_652.tex  (orig pages 601--652)
% =========================================================
\noindent\textbf{[485. Problems and Solutions --- continued from previous pages]}

\section*{Page 601}
\setcounter{page}{601}

But from the two equations in $x$, we have
\[
m^2(\sin^2\theta + \sin^2\phi) = 4c^2 + 2a^2 + 2b^2 + 4cx,
\]
\[
m^2\sin\theta\sin\phi = 4cx + a^2 - b^2,
\]
whence
\[
x = \frac{-b^2 - a^2 + m^2\sin\theta\sin\phi}{2c};
\]
therefore
\[
\sin^2\theta + \sin^2\phi = \frac{4c^2 + 2a^2 + 2b^2}{m^2} + \frac{2\left\{b^2 - a^2 + m^2\sin\theta\sin\phi\right\}}{2cm^2}.
\]

Hence, comparing the two results, we have
\[
\frac{1 - \cos^2\alpha - \cos^2\beta}{\sin^2\alpha} = \frac{4c^2 + 2b^2 + 2a^2}{m^2}, \quad
\frac{\cos\alpha\cos\beta}{\sin^2\alpha} = \frac{b^2 - a^2}{2cm^2}, \quad
\frac{m}{2c},
\]
or, as these may also be written,
\[
\sin^2\alpha = \frac{m^2}{c^2}, \quad
\cos^2\alpha + \cos^2\beta = \frac{4c^2 + 2a^2 + 2b^2}{c^2}, \quad
2\cos\alpha\cos\beta = \frac{b^2 - a^2}{c^2},
\]
whence
\[
\sin^2\alpha = \frac{m^2}{c^2}, \quad
(\cos\alpha + \cos\beta)^2 = \frac{4c^2 - a^2}{c^2}, \quad
(\cos\alpha - \cos\beta)^2 = \frac{b^2}{c^2}, \quad
\sin\alpha = \frac{m}{c},
\]
so that $m$ being put equal to unity, or otherwise assumed at pleasure, $a$, $b$, $c$ are given functions of $\alpha$, $\beta$. Or conversely, if $a$, $b$, $c$ are assumed at pleasure, then $\alpha$, $\beta$, $m$ are given functions of these quantities.

\medskip\noindent\textbf{5.} It is to be remarked that $(\alpha, \beta)$ being real, $a$ and $b$ will be imaginary, and consequently the points $S$, $H$ of Professor Sylvester's theorem are imaginary\footnote{Prof.~Sylvester remarks that according as $\beta$ is less or greater than $\alpha$, we may find real values of $\theta$, $\phi$ equal to one another in the one case and supplementary in the other. Hence we must in any case be able to make $r = s$ and $r = -s$ indifferently, which shows \textit{a priori} that the line being supposed real, each point $S$, $H$ must be imaginary, but so that the squared distance of either from the line must be a real negative quantity, conformably to Prof.~Cayley's important observation in the text. W.J.M.}; if, however, we write $-a^2$, $-b^2$ in place of $a^2$, $b^2$ respectively, then the radicals $\sqrt{(c+a)^2 - a'^2}$, $\sqrt{(c-a)^2 - b'^2}$ have a real geometrical interpretation. The system of relations between $(\alpha, \beta, a, b, c, m)$ becomes
\[
(\cos\alpha + \cos\beta)^2 = \frac{a^2}{c^2}, \quad
(\cos\alpha - \cos\beta)^2 = \frac{b^2}{c^2}, \quad
\sin\alpha = \frac{m}{c},
\]
and considering $(a, b, c)$ as given, we may write
\[
\cos\alpha = \frac{a + b}{2c}, \quad
\cos\beta = \frac{a - b}{2c}, \quad
m = \sqrt{4c^2 - (a + b)^2},
\]
viz. we have either this system or the similar system obtained by writing $-b$ in place of $b$.

\medskip\noindent\textbf{6.} Consider two circles with the radii $a$, $b$ and having the distance of their centres $= 2c$, and to fix the ideas assume that $2c > a + b$, that is, that the circles are exterior to each other.

\section*{Page 602}
\setcounter{page}{602}

The foregoing equations signify that $90° - \alpha$, $90° - \beta$ are the inclinations to the line of centres of the inverse and the direct common tangents respectively, and that $m$ is the length of the inverse common tangent. And the theorem is, that considering two circles as above, and taking $M$ a variable point in the line of centres, if $r$, $s$ denote the tangential distances of $M$ from the two circles respectively, and if $m$ be the length of the inverse common tangent of the two circles, then the angles $\theta$, $\phi$ determined by the equations
\[
m\sin\theta = r + s, \quad m\sin\phi = r - s,
\]
are connected by the relation
\[
\cos\theta = \cos\alpha\cos\phi + \sin\alpha\sin\phi\cos\alpha,
\]
$(\alpha, \beta)$ being constant angles, determined as above.

\medskip\noindent\textbf{7.} It is to be remarked that, assuming
\[
k = \frac{\sin\alpha}{\sin\beta} = \frac{\sqrt{4c^2 - (a+b)^2}}{\sqrt{4c^2 - (a-b)^2}},
\]
that is, $k = \text{inverse common tangent} \div \text{direct common tangent}$, then we have
\[
\cos\alpha = \sqrt{1 - k^2\sin^2\beta} = \Delta\beta,
\]
or the equation in $\theta$, $\phi$ becomes
\[
\cos\theta = \cos\alpha\cos\phi + \sin\alpha\sin\phi\Delta\beta,
\]
which is the algebraical equation connecting the amplitudes of the elliptic functions in the relation $F(\theta) + F(\phi) = F(\psi)$.

\medskip\noindent\textbf{8.} It is very noticeable that the above figure leads to another relation in elliptic functions, viz. it is the very figure employed for that purpose by Jacobi; in fact, considering therein $YM$ as a variable tangent meeting the circle $A$ in the two points $X$ and $X'$, then if $2\psi$, $2\psi'$ denote the angles $GAX$, $GAX'$ respectively, it is easy to see geometrically that we have $\mathrm{d}\psi' : \mathrm{d}\psi = YX : YX'$; where
\[
(YX)^2 = (BX)^2 - b^2 = 4c^2 + a^2 + 4ca\cos 2\psi - b^2 = (2c+a)^2 - b^2 - 8ca\sin^2\psi,
\]
and similarly $(YX')^2 = (2c+a)^2 - b^2 - 8ca\sin^2\psi'$, that is, writing $k'^2 = \dfrac{8ca}{(2c+a)^2 - b^2}$, the differential equation is
\[
\frac{\mathrm{d}\psi}{\sqrt{1 - k'^2\sin^2\psi}} + \frac{\mathrm{d}\psi'}{\sqrt{1 - k'^2\sin^2\psi'}} = 0.
\]

\section*{Page 603}
\setcounter{page}{603}

corresponding to an integral equation
\[
F(\psi) + F(\psi') = \text{const.},
\]
the modulus of the elliptic functions being
\[
k' = \frac{\sqrt{8ca}}{\sqrt{(2c+a)^2 - b^2}}.
\]
In the problem above considered the modulus is
\[
k = \frac{\sqrt{4c^2 - (a-b)^2}}{2c},
\]
and it is not very easy to see the connexion between the two results.

\medskip\noindent\textbf{[Vol.~VI. p.~81.]}

\begin{theorem}[by Professor Cayley]
If $(A, A')$, $(B, B')$ are four points (two real and the other two imaginary) related to each other as foci and antifoci (that is, if the lines $AA'$, $BB'$ intersect at right angles in a point in such wise that $OA = OA' = i \cdot OB = i \cdot OB'$), then the product of the distances of any point $P$ from the points $A$, $A'$ is equal to the product of the distances of the same point $P$ from the points $B$, $B'$.
\end{theorem}

In fact, the coordinates of $A$, $A'$ may be taken to be $(a, 0)$, $(-a, 0)$, and those of $B$, $B'$ to be $(0, ai)$, $(0, -ai)$; whence, if $(x, y)$ are the coordinates of $P$, we have
\begin{align*}
(AP)^2 &= (x-a)^2 + y^2 = (x - a + iy)(x - a - iy),\\
(A'P)^2 &= (x+a)^2 + y^2 = (x + a + iy)(x + a - iy),\\
(BP)^2 &= x^2 + (y - ia)^2 = (x + iy + a)(x - iy - a),\\
(B'P)^2 &= x^2 + (y + ia)^2 = (x + iy - a)(x - iy + a),
\end{align*}
from which the theorem is at once seen to be true.

An important application of the theorem consists in the means which it affords of passing from the foci $(A, B, C, D)$ of a bicircular quartic, to the antifoci $(A', B')$ and $(C', D')$; viz. if these are $(A', B', C', D')$, then the equation $l\sqrt{A} + m\sqrt{B} + n\sqrt{C} = 0$ must be transformable into $l'\sqrt{A'} + m'\sqrt{B'} + n'\sqrt{C'} = 0$. Writing these respectively under the forms
\[
l^2A + m^2B - n^2C + 2lm\sqrt{AB} = 0, \quad
l'^2A' + m'^2B' - n'^2C' + 2l'm'\sqrt{A'B'} = 0,
\]
the two radicals $\sqrt{AB}$, $\sqrt{A'B'}$ are identical; and the remaining terms in the two equations respectively are rational functions, which when the ratios $l' : m' : n'$ are properly determined will be to each other in the ratio $1 : 1$; the two equations being thus identical.

\section*{Page 604}
\setcounter{page}{604}

\medskip\noindent\textbf{[Vol.~VI. p.~99.]}

\begin{problem}[1970. Proposed by Professor Cayley]
Find the conditions in order that the conics
\[
U = (a, b, c, f, g, h \,\raisebox{0.5ex}{\ensuremath{\diagup\kern-0.5em\diagup}}\, x, y, z)^2 = 0, \quad
U' = (a', b', c', f', g', h' \,\raisebox{0.5ex}{\ensuremath{\diagup\kern-0.5em\diagup}}\, x, y, z)^2 = 0,
\]
may have double contact.
\end{problem}

\begin{solution}[by the Proposer]
The coefficients of the two conics must be so related that for a properly determined value of $\theta$ we shall have identically $U - \theta U' = (\lambda x + \mu y + \nu z)^2$; but when this is so, the inverse coefficients of the quadric function $U - \theta U'$ are each $= 0$; that is, writing
\begin{align*}
(A, B, C, F, G, H) &= (bc - f^2,\; ca - g^2,\; ab - h^2,\; gh - af,\; hf - bg,\; fg - ch),\\
(A', B', C', F', G', H') &= (b'c' - f'^2,\; c'a' - g'^2,\; a'b' - h'^2,\; g'h' - a'f',\; h'f' - b'g',\; f'g' - c'h'),\\
(\mathfrak{A}, \mathfrak{B}, \mathfrak{C}, \mathfrak{F}, \mathfrak{G}, \mathfrak{H}) &= (bc' + b'c - 2ff',\; ca' + c'a - 2gg',\; ab' + a'b - 2hh',\\
&\qquad gh' + g'h - af' - a'f,\; hf' + h'f - bg' - b'g,\; fg' + f'g - ch' - c'h),
\end{align*}
then we have the six equations $A - \theta\mathfrak{A} - \theta^2 A' = 0$, \&c.

Or, eliminating $\theta$, the required conditions are
\[
\begin{vmatrix}
A & B & C & F & G & H \\
A' & B' & C' & F' & G' & H' \\
\mathfrak{A} & \mathfrak{B} & \mathfrak{C} & \mathfrak{F} & \mathfrak{G} & \mathfrak{H}
\end{vmatrix} = 0,
\]
equivalent to three relations between the two sets of coefficients.
\end{solution}

\medskip\noindent\textbf{[Vol.~VII., January to July, 1867, pp.~17--19.]}

\begin{problem}[2110. Proposed by Professor Cayley]
Prove that the locus of the foci of the parabolas which pass through three given points is a unicursal quintic curve passing through the two circular points at infinity.
\end{problem}

\begin{solution}[by the Proposer]
More generally it may be shown that for the conics which pass through three given points and touch a given line, the locus of the intersection of the tangents drawn from two fixed points $Q$, $Q'$ on this line to each conic of the series is a unicursal quintic passing through the two points $Q$ and $Q'$.

Taking the three given points to be the angles of the triangle $(x = 0, y = 0, z = 0)$, and the points $Q$, $Q'$ to be the points $(\alpha, \beta, \gamma)$ and $(\alpha', \beta', \gamma')$ respectively, the equation of a conic through the three points is
\[
fyz + gzx + hxy = 0,
\]

\section*{Page 605}
\setcounter{page}{605}

which conic will touch the line through the points $(\alpha, \beta, \gamma)$, $(\alpha', \beta', \gamma')$, if
\[
\sqrt{f(\beta\gamma' - \beta'\gamma)} + \sqrt{g(\gamma\alpha' - \gamma'\alpha)} + \sqrt{h(\alpha\beta' - \alpha'\beta)} = 0.
\]

The equation of the pair of tangents from $(\alpha, \beta, \gamma)$ to the conic is
\[
(f, g, h \,\raisebox{0.5ex}{\ensuremath{\diagup\kern-0.5em\diagup}}\, -gh, -hf, -fg \,\raisebox{0.5ex}{\ensuremath{\diagup\kern-0.5em\diagup}}\, \beta z - \gamma y, \gamma x - \alpha z, \alpha y - \beta x)^2 = 0,
\]
that is
\begin{multline*}
\alpha^2(gy + h\beta)^2 + \beta^2(h\alpha + f\gamma)^2 + \gamma^2(h\beta + g\alpha)^2 \\
+ 2\beta\gamma\{2gh\alpha^2 - (h\alpha + f\gamma)(h\beta + g\alpha)\} \\
+ 2\gamma\alpha\{2hg\beta^2 - (h\beta + g\alpha)(gy + h\beta)\} \\
+ 2\alpha\beta\{2fg\gamma^2 - (gy + h\beta)(h\alpha + f\gamma)\} = 0,
\end{multline*}
but one of the tangents through $(\alpha, \beta, \gamma)$ being
\[
x(\beta\gamma' - \beta'\gamma) + y(\gamma\alpha' - \gamma'\alpha) + z(\alpha\beta' - \alpha'\beta) = 0,
\]
it follows that the other tangent is
\[
\frac{\alpha(gy + h\beta)}{\beta\gamma' - \beta'\gamma} =
\frac{\beta(h\alpha + f\gamma)}{\gamma\alpha' - \gamma'\alpha} =
\frac{\gamma(h\beta + g\alpha)}{\alpha\beta' - \alpha'\beta}.
\]

Hence, writing for shortness
\[
A = gy + h\beta, \quad B = h\alpha + f\gamma, \quad C = f\beta + g\alpha,
\]
\[
A' = gy' + h\beta', \quad B' = h\alpha' + f\gamma', \quad C' = f\beta' + g\alpha',
\]
the equations of the tangents from $Q$, $Q'$ respectively are
\[
\frac{\alpha A}{\beta\gamma' - \beta'\gamma} =
\frac{\beta B}{\gamma\alpha' - \gamma'\alpha} =
\frac{\gamma C}{\alpha\beta' - \alpha'\beta},
\]
\[
\frac{\alpha' A'}{\beta\gamma' - \beta'\gamma} =
\frac{\beta' B'}{\gamma\alpha' - \gamma'\alpha} =
\frac{\gamma' C'}{\alpha\beta' - \alpha'\beta},
\]
and for the coordinates of the intersection of these tangents, we have
\begin{align*}
x : y : z &= (BC' - B'C) : (CA' - C'A) : (AB' - A'B) \\
&= \frac{BC' - B'C}{\beta\gamma' - \beta'\gamma} :
\frac{CA' - C'A}{\gamma\alpha' - \gamma'\alpha} :
\frac{AB' - A'B}{\alpha\beta' - \alpha'\beta}.
\end{align*}

Now
\begin{align*}
BC' - B'C &= f\{-f(\beta\gamma' - \beta'\gamma) + g(\gamma\alpha' - \gamma'\alpha) + h(\alpha\beta' - \alpha'\beta)\},\\
BC' + B'C &= 2f\gamma\alpha' + f\{\beta(\gamma\alpha' + \gamma'\alpha) + g(\gamma\alpha' + \gamma'\alpha) + h(\alpha\beta' + \alpha'\beta)\}.
\end{align*}

To satisfy the equation
\[
\sqrt{f(\beta\gamma' - \beta'\gamma)} + \sqrt{g(\gamma\alpha' - \gamma'\alpha)} + \sqrt{h(\alpha\beta' - \alpha'\beta)} = 0,
\]
write
\[
f = \frac{a^2}{\beta\gamma' - \beta'\gamma}, \quad
 g = \frac{b^2}{\gamma\alpha' - \gamma'\alpha}, \quad
 h = \frac{c^2}{\alpha\beta' - \alpha'\beta},
\]
and therefore $a + b + c = 0$; we then have
\begin{align*}
-f(\beta\gamma' - \beta'\gamma) + g(\gamma\alpha' - \gamma'\alpha) + h(\alpha\beta' - \alpha'\beta) &= -a^2 + b^2 + c^2 = -2bc;\\
BC' + B'C &= \text{function of second degree in } (a, b, c).
\end{align*}

\section*{Page 606}
\setcounter{page}{606}

and thence
\[
-f(\beta\gamma' - \beta'\gamma) + g(\gamma\alpha' - \gamma'\alpha) + h(\alpha\beta' - \alpha'\beta) = \frac{1}{\beta\gamma' - \beta'\gamma}(-2bc).
\]

and the equations become
\[
x : y : z = a(BC' - B'C) : b(CA' - C'A) : c(AB' - A'B),
\]
where $BC' + B'C$, $CA' + C'A$, $AB' + A'B$, substituting therein for $f$, $g$, $h$ the values
\[
f = \frac{a^2}{\beta\gamma' - \beta'\gamma}, \quad
 g = \frac{b^2}{\gamma\alpha' - \gamma'\alpha}, \quad
 h = \frac{c^2}{\alpha\beta' - \alpha'\beta},
\]
are respectively functions of the fourth degree in $(a, b, c)$; hence $(a, b, c)$ being connected by the relation $a + b + c = 0$, $x$, $y$, $z$ are proportional to quintic functions of $(a, b, c)$, or what is the same thing, writing $a, b, c = 1, \theta, -1 - \theta$, then $x$, $y$, $z$ are proportional to quintic functions of $\theta$, that is, the locus is a unicursal quintic curve.

That the curve passes through the points $(\alpha', \beta', \gamma')$ and $(\alpha, \beta, \gamma)$ appears by considering the conics $fyz + gzx + hxy = 0$, which pass through these points respectively. For the first of these conics we have $f : g : h = \alpha(\beta\gamma' - \beta'\gamma) : \beta(\gamma\alpha' - \gamma'\alpha) : \gamma(\alpha\beta' - \alpha'\beta)$; the equation
\[
\frac{\alpha A}{\beta\gamma' - \beta'\gamma} + \frac{\beta B}{\gamma\alpha' - \gamma'\alpha} + \frac{\gamma C}{\alpha\beta' - \alpha'\beta} = 0,
\]
reduces itself to $x(\beta\gamma' - \beta'\gamma) + y(\gamma\alpha' - \gamma'\alpha) + z(\alpha\beta' - \alpha'\beta) = 0$, and as the other equation
\[
\frac{\alpha' A'}{\beta\gamma' - \beta'\gamma} + \frac{\beta' B'}{\gamma\alpha' - \gamma'\alpha} + \frac{\gamma' C'}{\alpha\beta' - \alpha'\beta} = 0
\]
is that of a line through $(\alpha', \beta', \gamma')$ the two lines meet of course in the point $(\alpha', \beta', \gamma')$. And the like for the conic
\[
f : g : h = \alpha'(\beta\gamma' - \beta'\gamma) : \beta'(\gamma\alpha' - \gamma'\alpha) : \gamma'(\alpha\beta' - \alpha'\beta).
\]

If the triangle is equilateral, and $(x, y, z)$ are respectively proportional to the perpendicular distances from the three sides, then we have for the circular points at infinity
\[
(\alpha, \beta, \gamma) = (1, \omega, \omega^2), \quad (\alpha', \beta', \gamma') = (1, \omega^2, \omega),
\]
where $\omega$ is an imaginary cube root of unity. These values give
\[
\beta\gamma' - \beta'\gamma = \gamma\alpha' - \gamma'\alpha = \alpha\beta' - \alpha'\beta = \omega^2 - \omega,
\]
\[
\alpha\alpha' = \beta\beta' = \gamma\gamma' = 1, \quad
\beta\gamma' + \beta'\gamma = \gamma\alpha' + \gamma'\alpha = \alpha\beta' + \alpha'\beta = -1;
\]
and the expressions for $(x, y, z)$ take the form
\begin{align*}
x : y : z &= a\{2b^2c^2 - a^2(a^2 + b^2 + c^2)\} \\
&: b\{2c^2a^2 - b^2(a^2 + b^2 + c^2)\} \\
&: c\{2a^2b^2 - c^2(a^2 + b^2 + c^2)\},
\end{align*}
or, what is the same thing, reducing by means of the relation $a + b + c = 0$,
\[
x : y : z = a(a^4 - 2a^2bc - 2b^2c^2) : b(b^4 - 2b^2ca - 2c^2a^2) : c(c^4 - 2c^2ab - 2a^2b^2),
\]

\section*{Page 607}
\setcounter{page}{607}

and the equation of the curve is obtained by eliminating $(a, b, c)$ from these equations and the before mentioned equation $a + b + c = 0$.

\medskip\noindent\textbf{N.B.} The above is a particular case of the following general theorem of M.~Chasles:

If the conies of a system $(\mu, \nu)$ all of them touch the line $QQ'$ (viz. the tangents from $Q$ and $Q'$ to any conic of the system are the lines $QR$, $QR'$; $Q'R$, $Q'R'$), then the locus of the point $R$ is a curve of the order $2\mu + \nu$, having the point $Q$ for a $\mu$-ple point, and the point $Q'$ for a $\nu$-ple point.

\medskip\noindent\textbf{[Vol.~VII., pp.~40--43.]}

\begin{problem}[2166. Proposed by the Rev.~G.~H.~Delanoix]
Given three conics having a common chord $AB$; if $P$ be a variable point on the first conic, and $PA$, $PB$ meet the second conic in $Q$, $R$ and the third conic in $S$, $T$ respectively; it is required to find the locus of the intersection of $QR$ and $ST$.
\end{problem}

\begin{solution}[by Professor Cayley]
In particular, the first conic may be a pair of lines; and then, taking these for the lines $x = 0$, $y = 0$, if the equations of the second and third conics are $U = 0$, $V = 0$ respectively, the equations of the lines $QR$, $ST$ are $x\partial_\alpha U + y\partial_\beta U = 0$, $x\partial_\alpha V + y\partial_\beta V = 0$ respectively, and the required locus is
\[
\begin{vmatrix}
x & y \\
\partial_\alpha U & \partial_\beta U
\end{vmatrix}
\cdot
\begin{vmatrix}
x & y \\
\partial_\alpha V & \partial_\beta V
\end{vmatrix} = 0,
\]
where $(\alpha, \beta)$ are the coordinates of the variable point on the first conic; that is, we have $\alpha = 0$, $\beta = 0$ respectively. The required locus is thus the sum of two conics each passing through the four points $A$, $B$, $C$, $D$.
\end{solution}

\medskip\noindent\textbf{[Vol.~VII., pp.~49--52.]}

\begin{problem}[2167. Proposed by the Rev.~G.~H.~Delanoix]
Two points $A'$, $B'$, $C'$ being taken on the sides of the triangle $ABC$, find the locus of $P$ when the pencil $P(A'B'C'\Sigma)$ is constant, where $\Sigma$ is any fixed point whatever.
\end{problem}

\begin{solution}[by Professor Cayley]
Let $A'$, $B'$, $C'$ measure off on the three lines respectively the distances
\[
AA' = \pm\sqrt{a^2 - a'^2}, \quad
B'B = \pm\sqrt{b^2 - b'^2}, \quad
C'C = \pm\sqrt{c^2 - c'^2},
\]
or, considering each radical as containing implicitly its own sign, we may write $A'$, $B'$, $C'$ to denote these distances.

\section*{Page 608}
\setcounter{page}{608}

Taking $\alpha$, $\beta$, $\gamma$ as the coordinates of the fixed point $\Sigma$, the anharmonic ratio of the pencil is given by the equation
\[
P(A'B'C'\Sigma) = \frac{(\alpha_1 - \alpha_2)(\alpha_3 - \alpha_4)}{(\alpha_1 - \alpha_3)(\alpha_2 - \alpha_4)} = k \text{ (constant)},
\]
and the required locus is a cubic curve.

\medskip\noindent\textbf{[Vol.~VII., pp.~57--60.]}

\begin{problem}[2170. Proposed by Professor Cayley]
Find the locus of the centres of the conics which pass through four given points.
\end{problem}

\begin{solution}[by the Proposer]
The locus is a conic; in fact, it is the eleven-point conic of the four points. This conic passes through the six mid-points of the lines joining the four points two and two, and through the three points which are the harmonic centres of the four points, and through the two circular points at infinity (whence the theorem that all rectangular hyperbolas through the four points have their centres on a line).
\end{solution}

\medskip\noindent\textbf{[Vol.~VII., pp.~65--72.]}

\begin{problem}[2184. Proposed by Professor Cayley]
If on each side of a given triangle there be taken two points harmonically conjugate with respect to the extremities of that side, the six points thus determined lie on a conic.
\end{problem}

\begin{solution}[by the Proposer]
Let the equations of the sides of the triangle be $x = 0$, $y = 0$, $z = 0$, and let the six points be $A'$, $A''$ on $BC$; $B'$, $B''$ on $CA$; $C'$, $C''$ on $AB$. The condition that the six points lie on a conic is expressed by saying that the product of the three anharmonic ratios is equal to 1, that is,
\[
(BA', A''C) \cdot (CB', B''A) \cdot (AC', C''B) = 1.
\]
In fact, if the six points lie on a conic, the hexagon $A'B'C'A''B''C''$ is inscribed in this conic, and the three pairs of opposite sides meet in collinear points, which gives the condition above.
\end{solution}

\section*{Page 609}
\setcounter{page}{609}

that is, the equation of the twelfth order breaks up into four equations each of the third order. The geometrical theory may also be further developed. In fact, assuming a fixed line through $C$ meeting the second conic in the points $A$ and $B$; then considering the conic which passes through $P$ and touches at $I$, $J$ the lines $AI$, $AJ$ respectively; the envelope of the line joining the points of contact of the tangents from $P$ to this conic is a cubic curve.

\medskip\noindent\textbf{[Vol.~VII., pp.~73--77.]}

\begin{problem}[2190. Proposed by Professor Cayley]
Given two conics, if through a variable point $P$ there be drawn to the conics two pairs of tangents, and the lines joining the points of contact of each pair of tangents be taken, find the envelope of the line which cuts harmonically these two joining lines.
\end{problem}

\begin{solution}[by the Proposer]
Let the two conics be $S = 0$, $S' = 0$, and let $P$ be the variable point. The tangents from $P$ to $S$ meet it in the points $A$, $B$; the tangents from $P$ to $S'$ meet it in the points $A'$, $B'$. The lines $AB$ and $A'B'$ are the polars of $P$ with respect to the two conics. The envelope of the line which cuts these two lines harmonically is a conic.
\end{solution}

Now the tangential equation of the envelope of the line which cuts harmonically the last-mentioned two conics, is
\[
\Sigma \cdot \Sigma' - \Theta^2 = 0,
\]
where $\Sigma = 0$, $\Sigma' = 0$ are the tangential equations of the two conics, and $\Theta = 0$ is the tangential equation of the pair of common chords.

\medskip\noindent\textbf{[Vol.~VII., pp.~78--81.]}

\begin{problem}[2205. Proposed by Professor Cayley]
Find the number of partitions of $n$ things into $m$ sets of $r$ things each.
\end{problem}

\begin{solution}[by the Proposer]
The number of such partitions is
\[
\frac{n!}{(r!)^m \cdot m!}.
\]
\end{solution}

\section*{Page 610}
\setcounter{page}{610}

Introducing an arbitrary constant $\xi_0$, and putting the equation under the form
\[
\xi_0^2 \Sigma - 2\xi_0 \Theta + \Sigma' = 0,
\]
this may be identified with
\[
(\xi_0 \alpha + \beta)(\xi_0 \alpha' + \beta') = 0,
\]
where $\alpha = 0$, $\beta = 0$ are the equations of two lines through the point $P$, and $\alpha' = 0$, $\beta' = 0$ are the equations of two other lines through $P$; and the envelope is thus the pair of lines through $P$ which are harmonically conjugate with respect to the two conics.

\medskip\noindent\textbf{[Vol.~VII., pp.~82--86.]}

\begin{problem}[2218. Proposed by Professor Cayley]
If $A$, $B$, $C$, $D$ be four points on a conic, and if $P$ be a point such that the pencil $P(ABCD)$ is harmonic, then the locus of $P$ is a conic passing through $A$, $B$, $C$, $D$.
\end{problem}

\begin{solution}[by the Proposer]
Let the conic be $S = 0$, and let the coordinates of $A$, $B$, $C$, $D$ be $(x_1, y_1, z_1)$, $(x_2, y_2, z_2)$, $(x_3, y_3, z_3)$, $(x_4, y_4, z_4)$. The condition that the pencil $P(ABCD)$ is harmonic gives a quadratic equation in the coordinates of $P$, which represents a conic passing through the four points.
\end{solution}

\medskip\noindent\textbf{[Vol.~VII., pp.~87--93.]}

\begin{problem}[2231. Proposed by Professor Cayley]
If from a variable point $P$ on a given conic tangents be drawn to another given conic, find the locus of the intersection of the tangents at the points of contact.
\end{problem}

\begin{solution}[by the Proposer]
The locus is a quartic curve having double points at the intersections of the two conics.
\end{solution}

\section*{Page 611}
\setcounter{page}{611}

Now the tangential equation of the envelope of the line which cuts harmonically the last-mentioned two conics, is
\[
\Sigma \cdot \Sigma' - \Theta^2 = 0,
\]
where $\Sigma = 0$, $\Sigma' = 0$ are the tangential equations of the two conics, and $\Theta = 0$ is the tangential equation of the pair of common chords.

For instance, if the two conics are
\[
S = ax^2 + by^2 + cz^2 = 0, \quad
S' = a'x^2 + b'y^2 + c'z^2 = 0,
\]
then
\[
\Sigma = a\xi^2 + b\eta^2 + c\zeta^2 = 0, \quad
\Sigma' = a'\xi^2 + b'\eta^2 + c'\zeta^2 = 0,
\]
\[
\Theta = aa'\xi^2 + bb'\eta^2 + cc'\zeta^2 = 0,
\]
and the envelope is
\[
(a\xi^2 + b\eta^2 + c\zeta^2)(a'\xi^2 + b'\eta^2 + c'\zeta^2) - (aa'\xi^2 + bb'\eta^2 + cc'\zeta^2)^2 = 0,
\]
which reduces to
\[
aa'(a - a')^2\xi^4 + bb'(b - b')^2\eta^4 + cc'(c - c')^2\zeta^4 + \cdots = 0.
\]

\medskip\noindent\textbf{[Vol.~VII., pp.~94--99.]}

\begin{problem}[2255. Proposed by Professor Cayley]
If $A$, $B$, $C$ be three points on a cubic curve, and if the tangents at these points meet the curve again in $A'$, $B'$, $C'$ respectively, show that the points $A'$, $B'$, $C'$ are collinear.
\end{problem}

\begin{solution}[by the Proposer]
This is a well-known property of cubic curves. The line $A'B'C'$ is called the \textit{satellite line} of the points $A$, $B$, $C$.
\end{solution}

\section*{Page 612}
\setcounter{page}{612}

Introducing an arbitrary constant $\xi_0$, and putting the equation under the form
\[
\xi_0^2 \Sigma - 2\xi_0 \Theta + \Sigma' = 0,
\]
this may be identified with
\[
(\xi_0 \alpha + \beta)(\xi_0 \alpha' + \beta') = 0,
\]
where $\alpha = 0$, $\beta = 0$ are the equations of two lines through the point $P$, and $\alpha' = 0$, $\beta' = 0$ are the equations of two other lines through $P$; and the envelope is thus the pair of lines through $P$ which are harmonically conjugate with respect to the two conics.

\medskip\noindent\textbf{[Vol.~VII., pp.~100--108.]}

\begin{problem}[2275. Proposed by Professor Cayley]
Find the equation of the surface generated by a line which moves so as always to intersect three given lines.
\end{problem}

\begin{solution}[by the Proposer]
Let the three given lines be
\[
L_1 = 0, \quad L_2 = 0, \quad L_3 = 0.
\]
A line meeting all three must lie on the hyperboloid determined by any two of them, and the locus is a quadric surface.
\end{solution}

\medskip\noindent\textbf{[Vol.~VII., pp.~109--116.]}

\begin{problem}[2290. Proposed by Professor Cayley]
If $S = 0$ be the equation of a surface of the $n$th order, and if $P$ be a point on the surface, find the equation of the tangent plane at $P$.
\end{problem}

\begin{solution}[by the Proposer]
The equation of the tangent plane is
\[
x\frac{\partial S}{\partial x_1} + y\frac{\partial S}{\partial y_1} + z\frac{\partial S}{\partial z_1} + w\frac{\partial S}{\partial w_1} = 0,
\]
where $(x_1, y_1, z_1, w_1)$ are the coordinates of $P$.
\end{solution}

\section*{Page 613}
\setcounter{page}{613}

For instance, $n = 4$; partitions of $(a, b, c, d)$ into two parts are $(a, bcd)$, $(b, cda)$, $(c, dab)$, $(d, abc)$, $(ab, cd)$, $(ac, db)$, $(ad, bc)$; no. is $= 7$. Partitions into three parts are $(a, b, cd)$, $(a, c, bd)$, $(a, d, bc)$, $(b, c, ad)$, $(b, d, ac)$, $(c, d, ab)$; no. is $= 6$.

\medskip\noindent\textbf{[Vol.~VII., pp.~117--124.]}

\begin{problem}[2310. Proposed by Professor Cayley]
If a curve be drawn on a surface so as to cut all the asymptotic lines at a constant angle, find the differential equation of the curve.
\end{problem}

\begin{solution}[by the Proposer]
Let the surface be $z = f(x, y)$. The asymptotic lines are given by
\[
rdx^2 + 2s\,dx\,dy + t\,dy^2 = 0,
\]
where $r = \dfrac{\partial^2 z}{\partial x^2}$, $s = \dfrac{\partial^2 z}{\partial x\partial y}$, $t = \dfrac{\partial^2 z}{\partial y^2}$. If $\theta$ be the constant angle, the differential equation is
\[
\frac{r + 2s\frac{dy}{dx} + t\left(\frac{dy}{dx}\right)^2}{\sqrt{1 + p^2 + q^2}\sqrt{1 + \left(\frac{dy}{dx}\right)^2}} = \cos\theta,
\]
where $p = \dfrac{\partial z}{\partial x}$, $q = \dfrac{\partial z}{\partial y}$.
\end{solution}

\medskip\noindent\textbf{[Vol.~VII., pp.~125--133.]}

\begin{problem}[2330. Proposed by Professor Cayley]
Find the equation of the developable surface which passes through two given curves.
\end{problem}

\begin{solution}[by the Proposer]
Let the two curves be given as the intersections of the surfaces
\[
F(x, y, z) = 0, \quad G(x, y, z) = 0; \quad
F'(x, y, z) = 0, \quad G'(x, y, z) = 0.
\]
The developable surface is the locus of the lines which touch both curves. Its equation is found by eliminating the parameters from the equations which express that a line touches both curves.
\end{solution}

\section*{Page 614}
\setcounter{page}{614}

I have somewhere made the remark that, on the left-hand side, the terms which belong to the odd and the even values of $a + b + \cdots (= p)$ are equal, and that we have
\[
\sum (-)^{a+b+\cdots} \frac{(p - 1)!}{a!\,b!\,c!\cdots} = 0 \text{ or } \frac{2}{p},
\]
according as $p$ is odd or even.

\medskip\noindent\textbf{[Vol.~VII., pp.~134--142.]}

\begin{problem}[2350. Proposed by Professor Cayley]
If $S = 0$ be the equation of a quadric surface, and if $U = 0$, $V = 0$ be the equations of two planes, find the equation of the quadric surface which touches $S$ along the curve of intersection of $S$ with the plane $U + kV = 0$.
\end{problem}

\begin{solution}[by the Proposer]
The required equation is
\[
S = (U + kV)(\lambda U + \mu V),
\]
where $\lambda$, $\mu$ are determined by the condition that the surface passes through given points.
\end{solution}

\medskip\noindent\textbf{[Vol.~VII., pp.~143--151.]}

\begin{problem}[2370. Proposed by Professor Cayley]
If four points $A$, $B$, $C$, $D$ lie on a circle, and if $P$ be a point such that the pencil $P(ABCD)$ has a given anharmonic ratio, find the locus of $P$.
\end{problem}

\begin{solution}[by the Proposer]
The locus is a cubic curve passing through $A$, $B$, $C$, $D$ and the two circular points at infinity.
\end{solution}

\section*{Page 615}
\setcounter{page}{615}

which shows that the lines 14 and 23 intersect at right angles; similarly the lines 12 and 34, and also the lines 13 and 24, intersect at right angles; or starting from the orthocentric properties of the tetrahedron, we have the theorem:

\begin{theorem}
If $A$, $B$, $C$, $D$ be four points, and if the lines $AB$ and $CD$ intersect at right angles, and also $AC$ and $BD$, and also $AD$ and $BC$, then the four points form an orthocentric system.
\end{theorem}

\medskip\noindent\textbf{[Vol.~VII., pp.~152--160.]}

\begin{problem}[2390. Proposed by Professor Cayley]
Find the equation of the cone which passes through six given lines meeting in a point.
\end{problem}

\begin{solution}[by the Proposer]
Let the six lines be the edges of a tetrahedron. The cone which passes through the six edges is the cone reciprocal to the quadric surface which touches the six faces of the tetrahedron.
\end{solution}

\medskip\noindent\textbf{[Vol.~VII., pp.~161--169.]}

\begin{problem}[2410. Proposed by Professor Cayley]
If $A$, $B$, $C$ be three points on a conic, and if the tangent at $A$ meet the line $BC$ in $T$, show that
\[
\{TBC\} = \{CAB\}^2,
\]
where $\{TBC\}$ denotes the anharmonic ratio of the pencil $T(B, C, A, \ldots)$.
\end{problem}

\begin{solution}[by the Proposer]
This follows at once from the properties of the anharmonic ratios of points on a conic.
\end{solution}

\section*{Page 616}
\setcounter{page}{616}

we have therefore
\[
\lambda : \mu = (234) : -(134) = (234) : -(341),
\]
where $(234)$ denotes the determinant formed with the coordinates of the points 2, 3, 4.

\medskip\noindent\textbf{[Vol.~VII., pp.~170--178.]}

\begin{problem}[2430. Proposed by Professor Cayley]
Find the equation of the surface of revolution which passes through a given curve.
\end{problem}

\begin{solution}[by the Proposer]
Let the axis of revolution be the $z$-axis, and let the given curve be $F(x, y, z) = 0$, $G(x, y, z) = 0$. The equation of the surface is found by replacing $x^2 + y^2$ by its value in terms of $z$ from the equations of the curve.
\end{solution}

\medskip\noindent\textbf{[Vol.~VII., pp.~179--187.]}

\begin{problem}[2450. Proposed by Professor Cayley]
If $S = 0$ be the equation of a surface, and if $C$ be a curve on the surface, find the condition that $C$ may be a line of curvature.
\end{problem}

\begin{solution}[by the Proposer]
The condition is that the normals to the surface at consecutive points of the curve may intersect. If the surface be $z = f(x, y)$, and the curve be given by $\phi(x, y) = 0$, the condition is
\[
\begin{vmatrix}
\dfrac{\partial^2 z}{\partial x^2} & \dfrac{\partial^2 z}{\partial x\partial y} & \dfrac{\partial z}{\partial x} \\
\dfrac{\partial^2 z}{\partial x\partial y} & \dfrac{\partial^2 z}{\partial y^2} & \dfrac{\partial z}{\partial y} \\
\dfrac{\partial \phi}{\partial x} & \dfrac{\partial \phi}{\partial y} & 0
\end{vmatrix} = 0.
\]
\end{solution}

\section*{Page 617}
\setcounter{page}{617}

\[
\lambda = \frac{(12, 345)^2}{(12, 345)^2 - (12, 346)^2}, \quad
\mu = \frac{(12, 346)^2}{(12, 346)^2 - (12, 345)^2},
\]
(and therefore $\lambda + \mu = 1$) we find $(\lambda, \mu)$. Writing then
\[
\lambda \alpha_1 + \mu \alpha_2, \quad
\lambda \beta_1 + \mu \beta_2, \quad
\lambda \gamma_1 + \mu \gamma_2, \quad
\lambda \delta_1 + \mu \delta_2
\]
for $x$, $y$, $z$, $w$ in the equation
\[
12\alpha\beta\gamma\delta = 0,
\]
we obtain the required relation.

\medskip\noindent\textbf{[Vol.~VII., pp.~188--196.]}

\begin{problem}[2470. Proposed by Professor Cayley]
If $S = 0$ be the equation of a quadric surface, find the equation of the cone whose vertex is a given point and which circumscribes the surface.
\end{problem}

\begin{solution}[by the Proposer]
Let the given point be $(x_0, y_0, z_0, w_0)$. The equation of the cone is obtained by forming the discriminant of $S + \lambda\Pi = 0$, where $\Pi = 0$ is the equation of an arbitrary plane through the point.
\end{solution}

\medskip\noindent\textbf{[Vol.~VII., pp.~197--205.]}

\begin{problem}[2490. Proposed by Professor Cayley]
Find the equation of the ruled surface generated by a line which meets a given curve and is always parallel to a given plane.
\end{problem}

\begin{solution}[by the Proposer]
Let the curve be given by the parametric equations $x = f(t)$, $y = g(t)$, $z = h(t)$, and let the given plane be $z = 0$. The equations of the generating line are
\[
\frac{X - f(t)}{\cos\theta} = \frac{Y - g(t)}{\sin\theta} = \frac{Z - h(t)}{0},
\]
and the surface is found by eliminating $t$ between the equations $Z = h(t)$ and $(X - f(t))\sin\theta - (Y - g(t))\cos\theta = 0$.
\end{solution}

\section*{Page 618}
\setcounter{page}{618}

\medskip\noindent\textbf{[Vol.~VIII. pp.~51, 52.]}

\begin{note}[Note on Question 1990. By Professor Cayley]
The theorem may be generalized as follows: If $A$, $B$, $C$, $D$ be four points on a conic, and if $P$ be a point such that the pencil $P(ABCD)$ is harmonic, then the locus of $P$ is a conic through $A$, $B$, $C$, $D$.
\end{note}

\medskip\noindent\textbf{[Vol.~VIII. p.~74.]}

\begin{problem}[2371. Proposed by Professor Cayley]
If $P$, $Q$ be two points taken at random on a sphere, find the probability that the arc $PQ$ is less than a given arc $a$.
\end{problem}

\begin{solution}[by the Proposer]
The probability is
\[
\frac{1 - \cos a}{2}.
\]
\end{solution}

\section*{Page 619}
\setcounter{page}{619}

\begin{solution}[Solution by the Proposer.]
\textbf{1.} Consider four given points, and in connection therewith a given line $IJ$; the conics which pass through the four points and touch the given line form a system such that the locus of the pole of the given line is a conic.
\end{solution}

\medskip\noindent\textbf{[Vol.~VIII. pp.~75, 76.]}

\begin{problem}[2380. Proposed by Professor Cayley]
If $A$, $B$, $C$, $D$ be four points on a conic, find the locus of the point $P$ such that the pencil $P(ABCD)$ has a given anharmonic ratio $k$.
\end{problem}

\begin{solution}[by the Proposer]
The locus is a cubic curve passing through $A$, $B$, $C$, $D$ and the two circular points at infinity.
\end{solution}

\medskip\noindent\textbf{[Vol.~VIII. p.~77.]}

\begin{problem}[2391. Proposed by Professor Cayley]
If a conic pass through two given points and touch two given lines, find the locus of its centre.
\end{problem}

\begin{solution}[by the Proposer]
The locus is a conic.
\end{solution}

\section*{Page 620}
\setcounter{page}{620}

or, as it may also be written,
\[
(2a, -2b, 0, -f, g, 0 \,\raisebox{0.5ex}{\ensuremath{\diagup\kern-0.5em\diagup}}\, x, y, z)^2 = 0;
\]
which is the equation of a conic through the two given points and having the two given lines for tangents.

\medskip\noindent\textbf{[Vol.~VIII. pp.~78--80.]}

\begin{problem}[2400. Proposed by Professor Cayley]
If a conic inscribed in a triangle pass through a fixed point, the locus of its centre is a conic.
\end{problem}

\begin{solution}[by the Proposer]
Let the triangle be the triangle of reference, and let the fixed point be $(\alpha, \beta, \gamma)$. The equation of an inscribed conic is $\sqrt{lx} + \sqrt{my} + \sqrt{nz} = 0$, and the condition that it passes through $(\alpha, \beta, \gamma)$ gives a relation between $l$, $m$, $n$. The centre is given by
\[
\frac{x}{l} = \frac{y}{m} = \frac{z}{n},
\]
and the locus is found by eliminating $l$, $m$, $n$.
\end{solution}

\section*{Page 621}
\setcounter{page}{621}

\medskip\noindent\textbf{[Vol.~VIII. p.~74.]}

\begin{problem}[2371. Proposed by Professor Cayley]
If $P$, $Q$ be two points taken at random on the surface of a sphere, find the probability that the distance $PQ$ is less than a given distance $d$.
\end{problem}

\begin{solution}[by the Proposer]
Let $a$ be the radius of the sphere, and let $\theta = d/a$. The probability is
\[
\frac{1 - \cos\theta}{2}.
\]
\end{solution}

\medskip\noindent\textbf{[Vol.~VIII. pp.~81--83.]}

\begin{problem}[2410. Proposed by Professor Cayley]
Find the locus of the centres of the conics which touch four given lines.
\end{problem}

\begin{solution}[by the Proposer]
The locus is a conic; in fact, it is the eleven-point conic of the four lines.
\end{solution}

\section*{Page 622}
\setcounter{page}{622}

The lemma is at once proved by means of my theorem for the relation between the distances of five points in space, (\textit{Cambridge Mathematical Journal}, vol.~II. (1841), p.~267); viz. if the coordinates of the five points are $(x_1, y_1, z_1)$, $(x_2, y_2, z_2)$, \ldots, $(x_5, y_5, z_5)$, then the relation is
\[
\begin{vmatrix}
0 & d_{12}^2 & d_{13}^2 & d_{14}^2 & d_{15}^2 & 1 \\
d_{21}^2 & 0 & d_{23}^2 & d_{24}^2 & d_{25}^2 & 1 \\
d_{31}^2 & d_{32}^2 & 0 & d_{34}^2 & d_{35}^2 & 1 \\
d_{41}^2 & d_{42}^2 & d_{43}^2 & 0 & d_{45}^2 & 1 \\
d_{51}^2 & d_{52}^2 & d_{53}^2 & d_{54}^2 & 0 & 1 \\
1 & 1 & 1 & 1 & 1 & 0
\end{vmatrix} = 0,
\]
where $d_{rs}$ denotes the distance between the points $r$ and $s$.

\medskip\noindent\textbf{[Vol.~VIII. pp.~84--86.]}

\begin{problem}[2420. Proposed by Professor Cayley]
If $S = 0$ be the equation of a quadric surface, find the equation of the reciprocal surface.
\end{problem}

\begin{solution}[by the Proposer]
The reciprocal surface is the envelope of the polar planes of the points of a given surface. If $S = (a, b, c, d, f, g, h, l, m, n \,\raisebox{0.5ex}{\ensuremath{\diagup\kern-0.5em\diagup}}\, x, y, z, w)^2 = 0$, the reciprocal surface is $\Sigma = 0$, where $\Sigma$ is the determinant formed with the inverse coefficients.
\end{solution}

\section*{Page 623}
\setcounter{page}{623}

\begin{corollary}
It appears by the demonstration that for any four points not in the same plane, the expression
\[
\begin{vmatrix}
0 & d_{12}^2 & d_{13}^2 & d_{14}^2 & 1 \\
d_{21}^2 & 0 & d_{23}^2 & d_{24}^2 & 1 \\
d_{31}^2 & d_{32}^2 & 0 & d_{34}^2 & 1 \\
d_{41}^2 & d_{42}^2 & d_{43}^2 & 0 & 1 \\
1 & 1 & 1 & 1 & 0
\end{vmatrix}
\]
vanishes identically when the four points are in a plane; and conversely, if this expression vanishes, the four points are in a plane.
\end{corollary}

\medskip\noindent\textbf{[Vol.~VIII. pp.~87--89.]}

\begin{problem}[2440. Proposed by Professor Cayley]
If $S = 0$ be the equation of a surface, find the equation of the Gaussian measure of curvature at any point.
\end{problem}

\begin{solution}[by the Proposer]
If the surface be $z = f(x, y)$, the measure of curvature is
\[
K = \frac{rt - s^2}{(1 + p^2 + q^2)^2},
\]
where $p = \dfrac{\partial z}{\partial x}$, $q = \dfrac{\partial z}{\partial y}$, $r = \dfrac{\partial^2 z}{\partial x^2}$, $s = \dfrac{\partial^2 z}{\partial x\partial y}$, $t = \dfrac{\partial^2 z}{\partial y^2}$.
\end{solution}

\section*{Page 624}
\setcounter{page}{624}

\medskip\noindent\textbf{[Vol.~IX., January to June, 1868, pp.~20, 21.]}

\begin{note}[Note on Question 2471. By Professor Cayley]
The question is to find the locus of a point $P$ such that the sum of the squares of its distances from $n$ given points is constant. The locus is evidently a sphere. For if $(x_r, y_r, z_r)$ are the coordinates of the $r$th point, and $(x, y, z)$ those of $P$, the equation is
\[
\sum_{r=1}^n \{(x - x_r)^2 + (y - y_r)^2 + (z - z_r)^2\} = \text{const.},
\]
which reduces to
\[
n(x^2 + y^2 + z^2) - 2x\sum x_r - 2y\sum y_r - 2z\sum z_r + \sum(x_r^2 + y_r^2 + z_r^2) = \text{const.}
\]
\end{note}

\medskip\noindent\textbf{[Vol.~IX. pp.~38, 39.]}

\begin{problem}[2530. Proposed by Professor Cayley]
Trace the curve
\[
y = \frac{x^2 - 1}{x^2 + 1}.
\]
\end{problem}

\begin{solution}[by the Proposer]
The curve is a rectangular hyperbola with asymptotes $x = \pm 1$ and $y = 1$. It passes through the origin, and is symmetrical with respect to the axis of $y$.
\end{solution}

\section*{Page 625}
\setcounter{page}{625}

\medskip\noindent\textbf{[Vol.~IX. pp.~40--42.]}

\begin{problem}[2540. Proposed by Professor Cayley]
If $A$, $B$, $C$ be three points on a conic, and if the tangents at $A$, $B$, $C$ form a triangle $A'B'C'$, show that the lines $AA'$, $BB'$, $CC'$ meet in a point.
\end{problem}

\begin{solution}[by the Proposer]
This is a particular case of Brianchon's theorem. The point of concurrence is the pole of the line $ABC$ with respect to the conic.
\end{solution}

\medskip\noindent\textbf{[Vol.~IX. pp.~43--45.]}

\begin{problem}[2550. Proposed by Professor Cayley]
Find the equation of the surface generated by the revolution of a parabola about its directrix.
\end{problem}

\begin{solution}[by the Proposer]
The surface is a paraboloid of revolution. Its equation is
\[
x^2 + y^2 = 4az.
\]
\end{solution}

\medskip\noindent\textbf{[Vol.~IX. pp.~46--48.]}

\begin{problem}[2560. Proposed by Professor Cayley]
If $S = 0$ be the equation of a quadric, and if $P$ be a point, find the locus of the middle points of chords of the quadric drawn through $P$.
\end{problem}

\begin{solution}[by the Proposer]
The locus is a plane, called the \textit{polar plane} of $P$ with respect to the quadric.
\end{solution}

\section*{Page 626}
\setcounter{page}{626}

form of the curve of intersection of the surface by any other sphere having the same centre, and thence the form of the surface itself (being a particular case of Steiner's surface).

\medskip\noindent\textbf{[Vol.~IX. pp.~49--51.]}

\begin{problem}[2570. Proposed by Professor Cayley]
Find the equation of the surface which is the envelope of a sphere of constant radius whose centre moves on a given plane.
\end{problem}

\begin{solution}[by the Proposer]
The envelope is a pair of parallel planes, at a distance equal to the radius of the sphere on either side of the given plane.
\end{solution}

\medskip\noindent\textbf{[Vol.~IX. pp.~52--54.]}

\begin{problem}[2580. Proposed by Professor Cayley]
If $S = 0$ be the equation of a surface, and if $C$ be a curve on the surface, find the condition that $C$ may be an asymptotic line.
\end{problem}

\begin{solution}[by the Proposer]
The condition is that the osculating plane of the curve at each point may coincide with the tangent plane to the surface. If the surface be $z = f(x, y)$, the differential equation of the asymptotic lines is
\[
rdx^2 + 2s\,dx\,dy + t\,dy^2 = 0.
\]
\end{solution}

\section*{Page 627}
\setcounter{page}{627}

these tricuspidal forms; and build upon the tricuspidal forms, until the greatest variety of quartic surfaces is obtained.

\medskip\noindent\textbf{[Vol.~IX. pp.~55--60.]}

\begin{problem}[2600. Proposed by Professor Cayley]
Find the equation of the surface generated by a line which moves so as always to touch a given curve and intersect a given line.
\end{problem}

\begin{solution}[by the Proposer]
Let the given curve be $x = f(t)$, $y = g(t)$, $z = h(t)$, and let the given line be the $z$-axis. The generating line is the line joining the point $(f(t), g(t), h(t))$ to the point $(0, 0, h(t))$, and the surface is found by eliminating $t$.
\end{solution}

\medskip\noindent\textbf{[Vol.~IX. pp.~61--66.]}

\begin{problem}[2620. Proposed by Professor Cayley]
If $S = 0$ be the equation of a quadric surface, find the equation of the cylinder whose generators are parallel to a given line and which circumscribes the quadric.
\end{problem}

\begin{solution}[by the Proposer]
Let the given line be the line $x/l = y/m = z/n$. The equation of the cylinder is found by forming the discriminant of $S + \lambda(lx + my + nz)^2 = 0$.
\end{solution}

\section*{Page 628}
\setcounter{page}{628}

\medskip\noindent\textbf{[Vol.~IX. pp.~82, 83.]}

\begin{problem}[2493. Proposed by Professor Cayley]
1.~Given the conic $U = 0$ (but observe that the conic may be any conic, real or imaginary), and two points $A$, $B$; it is required to find the envelope of the variable conic which passes through $A$, $B$ and touches the conic $U = 0$.
\end{problem}

\begin{solution}[by the Proposer]
Let the line $AB$ be taken for the line $z = 0$, and let $A = (1, 0, 0)$, $B = (0, 1, 0)$. The variable conic has an equation of the form
\[
fyz + gzx + hxy + lxy = 0,
\]
and the condition of tangency with $U = 0$ gives a relation between $f$, $g$, $h$, $l$. The envelope is found by eliminating these parameters.
\end{solution}

\medskip\noindent\textbf{[Vol.~IX. pp.~84--86.]}

\begin{problem}[2630. Proposed by Professor Cayley]
Find the equation of the cone whose vertex is at a given point and which passes through a given conic.
\end{problem}

\begin{solution}[by the Proposer]
Let the given point be $(x_0, y_0, z_0, w_0)$ and let the conic be the intersection of the plane $z = 0$ with the quadric $S = 0$. The equation of the cone is found by substituting in $S = 0$ the coordinates of the line joining $(x_0, y_0, z_0, w_0)$ to a variable point on the conic.
\end{solution}

\section*{Page 629}
\setcounter{page}{629}

\textbf{4.} A construction has been given by Aronhold (\textit{Berl.~Monatsber.}, July, 1864) by which, taking any 7 given lines as double tangents of a quartic curve, the remaining 14 double tangents can be found. The construction depends on the properties of the Steinerian and Hessian of a net of cubics.

\medskip\noindent\textbf{[Vol.~IX. pp.~87--92.]}

\begin{problem}[2650. Proposed by Professor Cayley]
Find the equation of the developable surface generated by the tangents to a given curve.
\end{problem}

\begin{solution}[by the Proposer]
Let the curve be given by the parametric equations $x = f(t)$, $y = g(t)$, $z = h(t)$. The equations of the tangent at the point $t$ are
\[
\frac{X - f(t)}{f'(t)} = \frac{Y - g(t)}{g'(t)} = \frac{Z - h(t)}{h'(t)}.
\]
The developable is found by eliminating $t$ between these equations and the equation expressing that the tangent passes through a variable point $(X, Y, Z)$.
\end{solution}

\medskip\noindent\textbf{[Vol.~IX. pp.~93--98.]}

\begin{problem}[2670. Proposed by Professor Cayley]
If $S = 0$ be the equation of a surface, find the equation of the surface parallel to it at a distance $d$.
\end{problem}

\begin{solution}[by the Proposer]
The parallel surface is the envelope of a sphere of radius $d$ whose centre moves on the given surface. Its equation is found by eliminating $x$, $y$, $z$ between the equations
\[
S(x, y, z) = 0, \quad
(X - x)^2 + (Y - y)^2 + (Z - z)^2 = d^2,
\]
and the equations expressing that the sphere touches the surface.
\end{solution}

\section*{Page 630}
\setcounter{page}{630}

through the line $xx'$ all pass through the same two points, and since $B$, $C$ are two of these conics, $B$ and $C$ meet $xx'$ in the same two points $X$, $X'$; similarly $C$ and $A$ meet in the same two points $Y$, $Y'$; and $A$ and $B$ meet in the same two points $Z$, $Z'$.

\medskip\noindent\textbf{[Vol.~IX. pp.~99--104.]}

\begin{problem}[2690. Proposed by Professor Cayley]
If $A$, $B$, $C$ be three conics having a common chord, and if through a variable point $P$ lines be drawn meeting the conics in points $Q$, $R$ and $S$, $T$ respectively, find the locus of $P$ when the line $QR$ passes through a fixed point.
\end{problem}

\begin{solution}[by the Proposer]
The locus is a conic passing through the common points of the three conics.
\end{solution}

\medskip\noindent\textbf{[Vol.~IX. pp.~105--110.]}

\begin{problem}[2710. Proposed by Professor Cayley]
Find the equation of the surface generated by a line which moves so as always to intersect two given curves and be parallel to a given plane.
\end{problem}

\begin{solution}[by the Proposer]
Let the two curves be given by their parametric equations, and let the given plane be $z = 0$. The generating line is determined by the condition that it meets both curves and is parallel to the given plane. The surface is found by eliminating the parameters.
\end{solution}

\section*{Page 631}
\setcounter{page}{631}

before that the surface will include the quadric surface 2. But there still remains for consideration the case where 1 is a conical point on the surface, and I do not see how to treat this case in a satisfactory manner.

\medskip\noindent\textbf{[Vol.~IX. pp.~111--116.]}

\begin{problem}[2730. Proposed by Professor Cayley]
If $S = 0$ be the equation of a surface, and if $P$ be a singular point on the surface, find the equation of the tangent cone at $P$.
\end{problem}

\begin{solution}[by the Proposer]
The tangent cone is the cone generated by the tangents to the surface at the singular point. Its equation is obtained by equating to zero the terms of lowest degree in the equation of the surface when referred to the singular point as origin.
\end{solution}

\medskip\noindent\textbf{[Vol.~IX. pp.~117--122.]}

\begin{problem}[2750. Proposed by Professor Cayley]
Find the equation of the ruled surface generated by a line which meets three given lines and is parallel to a given plane.
\end{problem}

\begin{solution}[by the Proposer]
The surface is a hyperbolic paraboloid if the three given lines are not all parallel to the given plane; otherwise it is a cylinder.
\end{solution}

\section*{Page 632}
\setcounter{page}{632}

Taking $x = 0$, $y = 0$, $z = 0$ for the equations of the sides of the triangle $ABC$, the equations of the three conics may be taken to be $U = 0$, $V = 0$, $W = 0$, where the coefficients satisfy the relations
\[
U = (a, b, c, f, g, h \,\raisebox{0.5ex}{\ensuremath{\diagup\kern-0.5em\diagup}}\, x, y, z)^2 = 0, \quad \text{etc.}
\]

\medskip\noindent\textbf{[Vol.~IX. pp.~123--128.]}

\begin{problem}[2770. Proposed by Professor Cayley]
If $S = 0$ be the equation of a quadric surface, find the locus of the vertices of the cones which pass through a given curve and touch the quadric.
\end{problem}

\begin{solution}[by the Proposer]
The locus is a curve of the eighth order, being the intersection of two quartic surfaces.
\end{solution}

\medskip\noindent\textbf{[Vol.~IX. pp.~129--134.]}

\begin{problem}[2790. Proposed by Professor Cayley]
Find the equation of the surface generated by a line which moves so as always to touch a given sphere and intersect a given line.
\end{problem}

\begin{solution}[by the Proposer]
The surface is a quadric surface, being the envelope of the planes which contain the given line and touch the sphere.
\end{solution}

\section*{Page 633}
\setcounter{page}{633}

viz. this is easily found to be
\[
8(2g_1f_1 - c_1h_1)(2h_1f_1 - bg_1)(f_1 + f_2 + f_3) = 0,
\]
which is the required condition.

\medskip\noindent\textbf{[Vol.~IX. pp.~135--140.]}

\begin{problem}[2810. Proposed by Professor Cayley]
If $S = 0$ be the equation of a surface, find the equation of the caustic surface formed by the reflection of rays parallel to a given direction.
\end{problem}

\begin{solution}[by the Proposer]
The caustic surface is the envelope of the reflected rays. If the incident rays are parallel to the $z$-axis, and the reflecting surface is $z = f(x, y)$, the reflected ray at the point $(x, y, z)$ is the line through this point in the direction
\[
\left(-2p\frac{\partial z}{\partial x}, -2q\frac{\partial z}{\partial y}, 1 - p^2 - q^2\right),
\]
where $p = \dfrac{\partial z}{\partial x}$, $q = \dfrac{\partial z}{\partial y}$.
\end{solution}

\medskip\noindent\textbf{[Vol.~IX. pp.~141--146.]}

\begin{problem}[2830. Proposed by Professor Cayley]
Find the equation of the surface which is the envelope of a plane which cuts three given quadric surfaces in conics having a common point.
\end{problem}

\begin{solution}[by the Proposer]
The envelope is a surface of the fourth order. Its equation is found by eliminating the parameters of the plane from the conditions that the three conics have a common point.
\end{solution}

\section*{Page 634}
\setcounter{page}{634}

\begin{solution}[Solution by Professor Cayley.]
\textbf{1.} In adding any two numbers, we carry a certain number of times; and it is required to find the number of combinations of two numbers which give a specified number of carryings.
\end{solution}

\medskip\noindent\textbf{[Vol.~IX. pp.~147--152.]}

\begin{problem}[2850. Proposed by Professor Cayley]
If $S = 0$ be the equation of a quadric surface, find the equation of the surface which is the locus of the poles of the tangent planes of a given surface with respect to the quadric.
\end{problem}

\begin{solution}[by the Proposer]
The locus is the \textit{reciprocal} of the given surface with respect to the quadric. If $\Sigma = 0$ is the tangential equation of the given surface, the required locus is the surface whose tangential equation is obtained by substituting in $\Sigma$ the differential coefficients of $S$.
\end{solution}

\medskip\noindent\textbf{[Vol.~IX. pp.~153--158.]}

\begin{problem}[2870. Proposed by Professor Cayley]
Find the equation of the surface generated by a line which moves so as always to intersect a given line and touch a given sphere.
\end{problem}

\begin{solution}[by the Proposer]
The surface is a quadric of revolution if the given line passes through the centre of the sphere; otherwise it is a more general quadric surface.
\end{solution}

\section*{Page 635}
\setcounter{page}{635}

if $N$ be the number of times of carrying for the sum $m + n$, or of borrowing for the difference $(m + n) - m$ or $(m + n) - n$; that is, $N = x$, the required theorem. I remark that the theorem gives a rule for finding the number of combinations which give a specified number of carryings.

\medskip\noindent\textbf{[Vol.~IX. pp.~159--164.]}

\begin{problem}[2890. Proposed by Professor Cayley]
If $S = 0$ be the equation of a surface, find the condition that a given line may be an asymptotic line on the surface.
\end{problem}

\begin{solution}[by the Proposer]
The condition is that the line at each of its points may have a contact of the second order with the surface. If the line be taken for the $x$-axis, the condition is that the coefficients of $y$ and $z$ in the equation of the surface vanish, and that the terms of the second degree in $y$ and $z$ reduce to a perfect square.
\end{solution}

\medskip\noindent\textbf{[Vol.~IX. pp.~165--170.]}

\begin{problem}[2910. Proposed by Professor Cayley]
Find the equation of the surface generated by the perpendiculars from a fixed point to the tangent planes of a given surface.
\end{problem}

\begin{solution}[by the Proposer]
The surface is the \textit{polar reciprocal} of the given surface with respect to a sphere having the fixed point for centre.
\end{solution}

\section*{Page 636}
\setcounter{page}{636}

\medskip\noindent\textbf{[Vol.~XI., January to June, 1869, pp.~33--38.]}

\begin{problem}[2718. Proposed by Professor Cayley]
Find in piano the locus of a point $P$, given that the product of its distances from two fixed points $A$, $B$ is equal to the square of the distance from a third fixed point $C$.
\end{problem}

\begin{solution}[by the Proposer]
Taking $A$, $B$ for the points $(\pm c, 0)$ and $C$ for the point $(h, k)$, the equation of the locus is
\[
\{(x - c)^2 + y^2\}\{(x + c)^2 + y^2\} = \{(x - h)^2 + (y - k)^2\}^2,
\]
which is a quartic curve having the circular points at infinity for double points.
\end{solution}

\medskip\noindent\textbf{[Vol.~XI. pp.~39--44.]}

\begin{problem}[2738. Proposed by Professor Cayley]
If $A$, $B$, $C$, $D$ be four points on a conic, and if $P$ be a point such that the pencil $P(ABCD)$ has a given anharmonic ratio, find the locus of $P$.
\end{problem}

\begin{solution}[by the Proposer]
The locus is a cubic curve passing through $A$, $B$, $C$, $D$ and the two circular points at infinity.
\end{solution}

\section*{Page 637}
\setcounter{page}{637}

If $2m$, $q$, $c$, $d$, $f$ refer in like manner to the points $B$, $D$, then the equation of a circle, centre say $Q$, and passing through $B$, $D$, is
\[
(x^2 + y^2)\{m(x\cos\beta + y\sin\beta) - l(x\cos\lambda + y\sin\lambda)\}
\]

\medskip\noindent\textbf{[Vol.~XI. pp.~45--50.]}

\begin{problem}[2758. Proposed by Professor Cayley]
If $S = 0$ be the equation of a quadric surface, find the equation of the cylinder which circumscribes the quadric and has its generators parallel to a given line.
\end{problem}

\begin{solution}[by the Proposer]
Let the given line be $x/l = y/m = z/n$. The equation of the cylinder is found by forming the discriminant of $S + \lambda(lx + my + nz)^2 = 0$ with respect to $x$, $y$, $z$.
\end{solution}

\medskip\noindent\textbf{[Vol.~XI. pp.~51--56.]}

\begin{problem}[2778. Proposed by Professor Cayley]
Find the equation of the surface which is the envelope of a sphere whose centre moves on a given curve and whose radius varies according to a given law.
\end{problem}

\begin{solution}[by the Proposer]
The envelope is found by differentiating the equation of the sphere with respect to the parameter and eliminating the parameter between the original equation and the derivative.
\end{solution}

\section*{Page 638}
\setcounter{page}{638}

Taking the upper signs, the equation is
\[
(x^2 + y^2)\{m(x\cos\beta + y\sin\beta) - l(x\cos\lambda + y\sin\lambda)\}
\]

\medskip\noindent\textbf{[Vol.~XI. pp.~57--62.]}

\begin{problem}[2798. Proposed by Professor Cayley]
If $S = 0$ be the equation of a surface, find the equation of the surface of centres.
\end{problem}

\begin{solution}[by the Proposer]
The surface of centres is the locus of the centres of curvature of the given surface. Its equation is found by eliminating the coordinates of the point of contact between the given surface and its parallels.
\end{solution}

\medskip\noindent\textbf{[Vol.~XI. pp.~63--68.]}

\begin{problem}[2818. Proposed by Professor Cayley]
Find the equation of the surface generated by a line which moves so as always to touch two given curves.
\end{problem}

\begin{solution}[by the Proposer]
The surface is a developable surface. Its equation is found by expressing that the line joining two consecutive points of the generating curve meets both curves.
\end{solution}

\section*{Page 639}
\setcounter{page}{639}

which is apparently a quartic curve; but it is obvious \textit{a priori} that the locus includes as part of itself the line $AB$ which joins the two given points. In fact, there is in

\medskip\noindent\textbf{[Vol.~XI. pp.~69--74.]}

\begin{problem}[2838. Proposed by Professor Cayley]
If $S = 0$ be the equation of a quadric surface, find the locus of the point $P$ such that the sum of the squares of its distances from the foci of a given section of the quadric is constant.
\end{problem}

\begin{solution}[by the Proposer]
The locus is a quadric surface coaxial with the given quadric.
\end{solution}

\medskip\noindent\textbf{[Vol.~XI. pp.~75--80.]}

\begin{problem}[2858. Proposed by Professor Cayley]
Find the equation of the surface generated by a line which moves so as always to intersect three given curves.
\end{problem}

\begin{solution}[by the Proposer]
The surface is in general of the sixth order. It is the locus of the lines which meet the three given curves.
\end{solution}

\section*{Page 640}
\setcounter{page}{640}

\textbf{8.} Supposing that the points $(x_1, y_1)$, $(x_2, y_2)$ are on the conic $\dfrac{x^2}{a^2} + \dfrac{y^2}{b^2} = 1$, and that the line joining them passes through the point $(h, k)$, then the equation of the line is
\[
\frac{x}{a^2}(hx_1 + ky_1 - 1) + \frac{y}{b^2}(hx_2 + ky_2 - 1) = 0,
\]
and the condition that it passes through $(h, k)$ gives
\[
\frac{h}{a^2}(hx_1 + ky_1 - 1) + \frac{k}{b^2}(hx_2 + ky_2 - 1) = 0.
\]

\medskip\noindent\textbf{[Vol.~XI. pp.~81--86.]}

\begin{problem}[2878. Proposed by Professor Cayley]
If $S = 0$ be the equation of a surface, find the condition that two given lines may be conjugate with respect to the surface.
\end{problem}

\begin{solution}[by the Proposer]
Two lines are conjugate with respect to a surface if each lies in the polar plane of any point on the other. The condition is found by expressing that the polar line of one with respect to the surface coincides with the other.
\end{solution}

\section*{Page 641}
\setcounter{page}{641}

and, similarly,
\[
(y - y_0)(x_1 + yx_1 - 1) + (y - y_1)(x_0 + yx_0 - 1) = (x_0 + x_1)y + y_1(x - x_0) - y.
\]

\medskip\noindent\textbf{[Vol.~XI. pp.~87--92.]}

\begin{problem}[2898. Proposed by Professor Cayley]
If $S = 0$ be the equation of a quadric, and if $P$ be a point, find the locus of the poles of the planes through $P$ with respect to the quadric.
\end{problem}

\begin{solution}[by the Proposer]
The locus is a plane, called the \textit{polar plane} of $P$.
\end{solution}

\medskip\noindent\textbf{[Vol.~XI. pp.~93--98.]}

\begin{problem}[2918. Proposed by Professor Cayley]
Find the equation of the surface generated by a line which moves so as always to touch a given curve and be perpendicular to a given line.
\end{problem}

\begin{solution}[by the Proposer]
The surface is a developable surface. Its equation is found by expressing that the generating line is perpendicular to the given line and touches the curve.
\end{solution}

\section*{Page 642}
\setcounter{page}{642}

conics, \textit{qu\^{a}} conics touching four fixed lines, are a system $(0, 1)$; hence, taking for the fixed conic the two points $A$, $B$, we have, as a very particular case, the foregoing theorem.

\medskip\noindent\textbf{[Vol.~XI. pp.~99--104.]}

\begin{problem}[2938. Proposed by Professor Cayley]
If $S = 0$ be the equation of a surface, find the condition that a given plane may cut the surface in a parabola.
\end{problem}

\begin{solution}[by the Proposer]
The section is a parabola if the line at infinity in the given plane touches the asymptotic cone of the surface. The condition is that the discriminant of the terms of the second degree in the equation of the section vanish.
\end{solution}

\medskip\noindent\textbf{[Vol.~XI. pp.~105--110.]}

\begin{problem}[2958. Proposed by Professor Cayley]
Find the equation of the surface generated by a line which moves so as always to intersect two given lines and touch a given conic.
\end{problem}

\begin{solution}[by the Proposer]
The surface is a quartic surface. Its equation is found by expressing that the generating line meets the two given lines and touches the given conic.
\end{solution}

\section*{Page 643}
\setcounter{page}{643}

\medskip\noindent\textbf{[Vol.~XII., July to December, 1869, p.~69.]}

\begin{problem}[2920. Proposed by Professor Cayley]
Imagine a tetrahedron $BECC'$ in which the opposite edges are equal, viz. $BC = EC'$, $CE = C'B$, $EB = C'E$; this is what I call an \textit{orthocentric tetrahedron}. Show that the four altitudes meet in a point.
\end{problem}

\begin{solution}[by the Proposer]
The condition that the opposite edges are equal gives at once that the sum of the squares of the edges of any face is equal to the sum of the squares of the opposite edges. Hence the four altitudes meet in a point, which is the centre of the sphere circumscribed about the tetrahedron.
\end{solution}

\medskip\noindent\textbf{[Vol.~XII. pp.~70--75.]}

\begin{problem}[2940. Proposed by Professor Cayley]
If $S = 0$ be the equation of a quadric surface, find the locus of the centres of the sections of the quadric by planes passing through a given line.
\end{problem}

\begin{solution}[by the Proposer]
The locus is a conic, being the section of the diametral plane conjugate to the given line by the polar plane of any point on the given line.
\end{solution}

\section*{Page 644}
\setcounter{page}{644}

\begin{center}
\textbf{INDEX TO VOLUME VII.}
\end{center}

\begin{verbatim}
1872  Cayley  Intersection of two Surfaces    1969
Sylvester  Spherical Problem                1970
Cayley  Foci and Antifoci                   1971
Delanoix  Locus Problem                     1972
Cayley  Parabola Locus                      1973
Cayley  Conic Centres                       1974
Cayley  Harmonic Pencil                     1975
Cayley  Partition Problem                   1976
Cayley  Curve Tracing                       1977
Cayley  Orthocentric Tetrahedron            1978
\end{verbatim}

\section*{Page 645}
\setcounter{page}{645}

\begin{center}
\textbf{NOTES AND REFERENCES.}
\end{center}

\noindent\textbf{445, 451, 454.} We have the two papers by K.~Rohn, ``Die Fl\"{a}chen vierter Ordnung hinsichtlich ihrer Knotenpunkte und ihrer Reduction,'' \textit{Math.~Annalen}, t.~IV. (1871), pp.~524--550; and ``Ueber die Fl\"{a}chen vierter Ordnung mit achtzehn Knotenpunkten,'' \textit{Math.~Annalen}, t.~V. (1872), pp.~227--238.

\noindent\textbf{456.} See also the paper by A.~Clebsch, ``Ueber die Fl\"{a}chen vierter Ordnung, welche eine Doppelcurve vom sechzehnten Grade und vom Geschlechte zwei besitzen,'' \textit{Math.~Annalen}, t.~IV. (1871), pp.~551--581.

\noindent\textbf{458.} The memoirs of M.~Kummer on quartic surfaces are collected in the \textit{Berlin Monatsberichte} for the years 1862--1864, and in \textit{Crelle}, t.~LXIV. (1864).

\section*{Page 646}
\setcounter{page}{646}

\begin{center}
\textbf{TABLE OF NODE-CURVES ON QUARTIC SURFACES.}
\end{center}

\begin{center}
\begin{tabular}{|c|c|c|}
\hline
\textbf{No.~of Nodes} & \textbf{Node-Curve} & \textbf{Reference} \\
\hline
0 & -- & Cayley, 445 \\
1 & -- & Cayley, 451 \\
2 & Conic & Cayley, 454 \\
3 & Cubic & Rohn, 1871 \\
4 & Quadriquadric & Clebsch, 1871 \\
5 & Twisted cubic & Kummer, 1864 \\
6 & Two conics & Kummer, 1864 \\
7 & Cubic + line & Rohn, 1872 \\
8 & Four conics & Kummer, 1864 \\
9 & Two twisted cubics & Kummer, 1864 \\
10 & Five conics & Kummer, 1864 \\
11 & Six conics & Kummer, 1864 \\
12 & Eight conics & Kummer, 1864 \\
13 & Nine conics & Kummer, 1864 \\
14 & Ten conics & Kummer, 1864 \\
15 & Twelve conics & Kummer, 1864 \\
16 & Fifteen conics & Kummer, 1864 \\
\hline
\end{tabular}
\end{center}

\section*{Page 647}
\setcounter{page}{647}

\begin{center}
\textbf{SELECT SCIENTIFIC PUBLICATIONS OF THE CAMBRIDGE UNIVERSITY PRESS.}
\end{center}

The Collected Mathematical Papers of Arthur Cayley, Sc.D., F.R.S., Sadlerian Professor of Pure Mathematics in the University of Cambridge. Demy 4to. 10 vols. Vols.~I.~II.~III.~IV.~V.~VI.~and~VII. 25\emph{s.} each. [Vol.~VIII. In the Press.]

The Scientific Papers of the late Prof.~J.~Clerk Maxwell. Edited by W.~D.~Niven, M.A., formerly Fellow of Trinity College. In 2 vols. Royal 4to. \pounds 3.~3\emph{s.} (net).

\section*{Page 648}
\setcounter{page}{648}

\begin{center}
\textbf{CAMBRIDGE UNIVERSITY PRESS.}
\end{center}

A Treatise on Elementary Dynamics. By S.~L.~Loney, M.A., late Fellow of Sidney Sussex College. New and Enlarged Edition. Crown 8vo. 1\emph{s.}~6\emph{d.}

A Treatise on the Mathematical Theory of the Motion of Fluids. By Horace Lamb, M.A., formerly Fellow of Trinity College. Demy 8vo. 7\emph{s.}~6\emph{d.}

An Elementary Treatise on the Mathematical Theory of Perfectly Elastic Solids. By W.~J.~Hicks, M.A., Fellow of St.~John's College. Crown 8vo. 3\emph{s.}~6\emph{d.}

\section*{Page 649}
\setcounter{page}{649}

[This page is blank in the original.]

\section*{Page 650}
\setcounter{page}{650}

[This page is blank in the original.]

\section*{Page 651}
\setcounter{page}{651}

\begin{verbatim}
QA          3
.V7Cayley, Arthur
    Collected Mathematical Papers
    Vol. VII
\end{verbatim}

\section*{Page 652}
\setcounter{page}{652}

[This page is blank in the original.]
\clearpage

\end{document}
